% !TeX program = pdflatex
% papers/cristal_papers.tex — GERADO por tools/cristal_tex.py; NÃO editar à mão.
% Fonte: cristal/cristal.jsonl (broca-so, última versão por conceito).
% Cada secção carrega o registo original em comentário %CRISTAL — a volta é
% tests/cristal_volta.js (resíduo 0 byte a byte contra a fonte).
\documentclass[11pt,a4paper]{article}
\usepackage[utf8]{inputenc}\usepackage[T1]{fontenc}\usepackage[brazil]{babel}
\usepackage{amsmath,amssymb}\usepackage{textcomp}
\usepackage[margin=2.4cm]{geometry}\usepackage{xcolor}
\definecolor{cinza}{gray}{0.35}
\newcommand{\prov}[1]{\par\smallskip\noindent{\small\color{cinza}#1}\par}
\setcounter{secnumdepth}{0}
% ─────────────────────────────────────────────────────────────────────────────
%  papers/gkcapa.tex — a capa do Reino, mínima, para os papers em `article`.
%
%  O estilo.tex completo é do `report` (títulos de capítulo, skak, …). Aqui fica
%  só o que a capa precisa, para o paper continuar a compilar sozinho com
%  pdflatex. O tradutor próprio (tests/tex) carrega o estilo.tex da raiz / o
%  symlink papers/estilo.tex e avalia o mesmo `\gkcapa`.
% ─────────────────────────────────────────────────────────────────────────────
\definecolor{ouro}{HTML}{B8912F}
\definecolor{ouroclaro}{HTML}{F3E9CE}
\definecolor{tinta}{HTML}{1E1B16}
\definecolor{regua}{HTML}{6E6858}
\providecommand{\gknota}{\fontsize{7.62}{11.05}\selectfont}
\providecommand{\gktexto}{\fontsize{10.50}{15.22}\selectfont}
\providecommand{\gksub}{\fontsize{12.33}{17.87}\selectfont}
\providecommand{\gksec}{\fontsize{14.47}{20.98}\selectfont}
\providecommand{\gkcap}{\fontsize{16.99}{24.63}\selectfont}
\providecommand{\gktit}{\fontsize{23.42}{33.95}\selectfont}
\providecommand{\Prj}{\mathbb{P}}
\providecommand{\Trf}{\mathcal{T}}
\providecommand{\gkcapa}[3]{%
\title{\vspace{-0.6cm}
{\color{ouro}\rule{\textwidth}{1.4pt}}\\[6mm]
\textbf{\gktit\color{tinta} Reino Dourado}\\[3mm]{\gksec\itshape\color{regua} Golden Kingdom}\\[8mm]
{\gkcap\scshape\color{ouro} #1}\\[5mm]
{\gksub\color{regua} #2}\\[2mm]
{\gktexto\color{regua}\itshape #3}\\[7mm]
{\color{ouro}\rule{\textwidth}{0.6pt}}\\[9mm]{\gknota\color{regua}\begin{minipage}{\textwidth}\setlength{\parindent}{0pt}%
\textbf{Aviso de Isenção Algorítmica:} O universo, as leis e as entidades contidas nestas páginas ---
incluindo felinos matriciais, aves de rapina algébricas e amuletos de controle de fluxo --- são
construções de pura ficção formal. Esta obra não descreve a realidade física, não serve como
engenharia reversa do cosmos e não constitui doutrina religiosa. Qualquer semelhança com bugs reais do seu
sistema operacional é mera coincidência (ou falta de reza). Prossiga por sua própria conta e
risco.\end{minipage}}}
\author{}\date{}}

\begin{document}
\gkcapa{O Cristal --- os papers do cristal}
       {conceitos extraídos dos papers do projeto}
       {corpus recuperado do broca-so; 427 conceitos; projeção verificável da fonte}
\maketitle

%CRISTAL {"arestas":[{"alvo":"main","confianca":1.0,"epistemico":"fato","origem":"paper","rel":"parte_de"},{"alvo":"beta_numeracao_metalica","confianca":1.0,"epistemico":"fato","origem":"wiki","rel":"explica"}],"confianca":1.0,"contraexemplos":[],"descricao":"N O fluxo hiperbólico linear q₁=q₁- q₂, q₂=q₂- q₁, com ᵢ=sign(qᵢ), parece dar antifase. Mas, por ser hiperbólico e linear, escapa (|q|): os sinais congelam, não há codificação simbólica não-trivial, não há >1. A antifase obtida é a estrita (alternância, que exclui 11 e 00), de =0. D A dimensão positiva exige antifase com ramificação: excluir 11 mas não 00. É a ramificação no vazio que produz >1.","epistemico":"fato","exemplos":[],"id":"a_-transformacao_de_ouro_como_dinamica_admissivel","memoria":"semantica","meta":{"dominio":"paper","nivel":5,"verificacao":"slug idêntico — enriquecer, não duplicar"},"origem":"paper","palavras_chave":[],"sinonimos":[],"tipo":"conceito","titulo":"A transformação de ouro como dinâmica admissível"}
\section{A transformação de ouro como dinâmica admissível}
N O fluxo hiperbólico linear q\ensuremath{_{1}}=q\ensuremath{_{1}}- q\ensuremath{_{2}}, q\ensuremath{_{2}}=q\ensuremath{_{2}}- q\ensuremath{_{1}}, com \ensuremath{_{i}}=sign(q\ensuremath{_{i}}), parece dar antifase. Mas, por ser hiperbólico e linear, escapa (|q|): os sinais congelam, não há codificação simbólica não-trivial, não há >1. A antifase obtida é a estrita (alternância, que exclui 11 e 00), de =0. D A dimensão positiva exige antifase com ramificação: excluir 11 mas não 00. É a ramificação no vazio que produz >1.
\prov{relações: parte\_de main; explica beta\_numeracao\_metalica}
\prov{proveniência: paper \textperiodcentered{} fato \textperiodcentered{} confiança 1.0 \textperiodcentered{} domínio paper \textperiodcentered{} história 7 versões no jornal}

%CRISTAL {"arestas":[{"alvo":"geometria","confianca":1.0,"epistemico":"fato","origem":"paper","rel":"parte_de"},{"alvo":"corpo_dual","confianca":1.0,"epistemico":"fato","origem":"paper","rel":"parte_de"},{"alvo":"conjugados_galois_parabola","confianca":1.0,"epistemico":"fato","origem":"wiki","rel":"relaciona"}],"confianca":1.0,"contraexemplos":[],"descricao":"O princípio organizador desta construção não é a pureza, mas a conservação do defeito: as\nobstruções não se eliminam integralmente --- redistribuem-se entre os níveis sucessivos da torre, e a geometria\nemerge precisamente da estrutura que sobrevive a essa redistribuição.\n\n A geometria não é um palco posto antes da dinâmica: é o que permanece conservado quando a topologia da\ntorre encontra a álgebra da balança. Em três camadas --- a álgebra diz como o defeito troca; a\ntopologia, para onde ele vai","epistemico":"fato","exemplos":[],"id":"a_algebra_a_balanca_ac21","memoria":"semantica","meta":{"dominio":"paper","nivel":5,"verificacao":"slug idêntico — enriquecer, não duplicar"},"origem":"paper","palavras_chave":[],"sinonimos":[],"tipo":"conceito","titulo":"A álgebra: a balança a+c²=1"}
\section{A álgebra: a balança a+c\ensuremath{^{2}}=1}
O princípio organizador desta construção não é a pureza, mas a conservação do defeito: as
obstruções não se eliminam integralmente --- redistribuem-se entre os níveis sucessivos da torre, e a geometria
emerge precisamente da estrutura que sobrevive a essa redistribuição.

 A geometria não é um palco posto antes da dinâmica: é o que permanece conservado quando a topologia da
torre encontra a álgebra da balança. Em três camadas --- a álgebra diz como o defeito troca; a
topologia, para onde ele vai
\prov{relações: parte\_de geometria; parte\_de corpo\_dual; relaciona conjugados\_galois\_parabola}
\prov{proveniência: paper \textperiodcentered{} fato \textperiodcentered{} confiança 1.0 \textperiodcentered{} domínio paper \textperiodcentered{} história 52 versões no jornal}

%CRISTAL {"arestas":[{"alvo":"seguranca_por_irreversibilidade","confianca":1.0,"epistemico":"fato","origem":"paper","rel":"parte_de"}],"confianca":1.0,"contraexemplos":[],"descricao":"Contração de Hellinger Cada máquina Tᵢ é um kernel estocástico (saída normalizada). Para todo tal K: H(KP,KQ) H(P,Q), equivalentemente (KP,KQ)(P,Q). lema (Demonstrado.) A desigualdade de processamento para Hellinger vale para todo kernel de Markov; Tᵢ o é pela normalização. Compor máquinas só mantém ou reduz a distinção: assinaturas coladas (1) permanecem coladas sob toda transformação em A.","epistemico":"fato","exemplos":[],"id":"a_algebra_das_maquinas_e_o_nao-ataque","memoria":"semantica","meta":{"dominio":"paper","nivel":5,"verificacao":"conceito novo"},"origem":"paper","palavras_chave":[],"sinonimos":[],"tipo":"conceito","titulo":"A álgebra das máquinas e o não-ataque"}
\section{A álgebra das máquinas e o não-ataque}
Contração de Hellinger Cada máquina T\ensuremath{_{i}} é um kernel estocástico (saída normalizada). Para todo tal K: H(KP,KQ) H(P,Q), equivalentemente (KP,KQ)(P,Q). lema (Demonstrado.) A desigualdade de processamento para Hellinger vale para todo kernel de Markov; T\ensuremath{_{i}} o é pela normalização. Compor máquinas só mantém ou reduz a distinção: assinaturas coladas (1) permanecem coladas sob toda transformação em A.
\prov{relações: parte\_de seguranca\_por\_irreversibilidade}
\prov{proveniência: paper \textperiodcentered{} fato \textperiodcentered{} confiança 1.0 \textperiodcentered{} domínio paper \textperiodcentered{} história 3 versões no jornal}

%CRISTAL {"arestas":[{"alvo":"broca_os","confianca":1.0,"epistemico":"fato","origem":"paper","rel":"parte_de"}],"confianca":1.0,"contraexemplos":[],"descricao":"D O painel não é um conjunto de rotas REST. A op é o canal: o comando desembrulhado (§III) traz a operação, e o demux roteia por banda =sha256(tecido). Não há endpoint a montar sobre a álgebra --- a operação é lida no próprio barramento. A interface é o canal Há uma bijeção entre as operações do painel e os canais Eᵢj da parte negra: invocar a op é exibir um tecido cuja banda casa com eᵢ,M,ej; quem casa decodifica, o resto lê ruído. Logo o painel não acrescenta superfície","epistemico":"fato","exemplos":[],"id":"a_ancora_hodge","memoria":"semantica","meta":{"dominio":"paper","nivel":5,"verificacao":"slug idêntico — enriquecer, não duplicar"},"origem":"paper","palavras_chave":[],"sinonimos":[],"tipo":"conceito","titulo":"A âncora: Hodge"}
\section{A âncora: Hodge}
D O painel não é um conjunto de rotas REST. A op é o canal: o comando desembrulhado (§III) traz a operação, e o demux roteia por banda =sha256(tecido). Não há endpoint a montar sobre a álgebra --- a operação é lida no próprio barramento. A interface é o canal Há uma bijeção entre as operações do painel e os canais E\ensuremath{_{i}}j da parte negra: invocar a op é exibir um tecido cuja banda casa com e\ensuremath{_{i}},M,ej; quem casa decodifica, o resto lê ruído. Logo o painel não acrescenta superfície
\prov{relações: parte\_de broca\_os}
\prov{proveniência: paper \textperiodcentered{} fato \textperiodcentered{} confiança 1.0 \textperiodcentered{} domínio paper \textperiodcentered{} história 10 versões no jornal}

%CRISTAL {"arestas":[{"alvo":"ciencias","confianca":-0.031,"epistemico":"fato","origem":"manual","rel":"referencia"},{"alvo":"quatro_ciencias","confianca":1.0,"epistemico":"fato","origem":"paper","rel":"parte_de"}],"confianca":1.0,"contraexemplos":[],"descricao":"D A regra está cravada na porta do dado bruto: ingere o fluxo INTEIRO, sem picotar --- não picotar\npreserva todas as correlações e o associador entre os blocos. Em bit, todo formato ou idioma é o mesmo tipo de\nfluxo; é a torção que o assenta, não uma curadoria prévia.\n\nprop[Escolher é picotar é cegar]\nAmostrar ou selecionar o dado quebra o assentamento por três vias:\nitemize1pt\n[(1)] destrói o associador --- a correlação entre blocos é justamente o que distingue o idioma;\npicotar corta-a, apagan","epistemico":"fato","exemplos":[],"id":"a_ancora_observatorio_nao_fabrica","memoria":"semantica","meta":{"dominio":"paper","nivel":5,"verificacao":"slug idêntico — enriquecer, não duplicar"},"origem":"paper","palavras_chave":[],"sinonimos":[],"tipo":"conceito","titulo":"A âncora: observatório, não fábrica"}
\section{A âncora: observatório, não fábrica}
D A regra está cravada na porta do dado bruto: ingere o fluxo INTEIRO, sem picotar --- não picotar
preserva todas as correlações e o associador entre os blocos. Em bit, todo formato ou idioma é o mesmo tipo de
fluxo; é a torção que o assenta, não uma curadoria prévia.

prop[Escolher é picotar é cegar]
Amostrar ou selecionar o dado quebra o assentamento por três vias:
itemize1pt
[(1)] destrói o associador --- a correlação entre blocos é justamente o que distingue o idioma;
picotar corta-a, apagan
\prov{relações: referencia ciencias; parte\_de quatro\_ciencias}
\prov{proveniência: paper \textperiodcentered{} fato \textperiodcentered{} confiança 1.0 \textperiodcentered{} domínio paper \textperiodcentered{} história 12 versões no jornal}

%CRISTAL {"arestas":[{"alvo":"redes","confianca":1.0,"epistemico":"fato","origem":"paper","rel":"parte_de"}],"confianca":1.0,"contraexemplos":[],"descricao":"As redes da máquina não são clássicas e não armazenam: transmitem em tempo real. Só a diferença viaja --- o bump, o gap msg(banda), cru no datagrama, bits brutos, sem cifra por cima (a função de onda chega intacta). A banda é a assinatura do tecido (sha256): o identificador de banda filtra --- quem tem o tecido sintoniza, quem não tem lê ruído (nem decodifica). Não há tratamento temporal: sem buffer, sem ordem, sem armazenamento --- o índice de sequência é só idempotência; o bit corre solto.","epistemico":"fato","exemplos":[],"id":"a_antena_transmite_em_tempo_real","memoria":"semantica","meta":{"dominio":"paper","nivel":5,"verificacao":"conceito novo"},"origem":"paper","palavras_chave":[],"sinonimos":[],"tipo":"conceito","titulo":"A antena transmite em tempo real"}
\section{A antena transmite em tempo real}
As redes da máquina não são clássicas e não armazenam: transmitem em tempo real. Só a diferença viaja --- o bump, o gap msg(banda), cru no datagrama, bits brutos, sem cifra por cima (a função de onda chega intacta). A banda é a assinatura do tecido (sha256): o identificador de banda filtra --- quem tem o tecido sintoniza, quem não tem lê ruído (nem decodifica). Não há tratamento temporal: sem buffer, sem ordem, sem armazenamento --- o índice de sequência é só idempotência; o bit corre solto.
\prov{relações: parte\_de redes}
\prov{proveniência: paper \textperiodcentered{} fato \textperiodcentered{} confiança 1.0 \textperiodcentered{} domínio paper \textperiodcentered{} história 3 versões no jornal}

%CRISTAL {"arestas":[{"alvo":"paper_45_confluencia_ouro","confianca":1.0,"epistemico":"fato","origem":"paper","rel":"parte_de"},{"alvo":"paper_2_reta_ouro","confianca":1.0,"epistemico":"fato","origem":"paper","rel":"confirma"}],"confianca":1.0,"contraexemplos":[],"descricao":"A pergunta --- há uma série infinita, ou um algoritmo? --- tem resposta firme nas duas formas. a funcão zeta do ouro: a série, com polo em 1/ --- verificado A funcão zeta dinâmica do golden mean shift é a geradora de Fibonacci: [ (t)=11-t-t²=ₙ₀Fₙ+₁,tⁿ. ] O seu polo mais próximo é t=1/ (raiz de 1-t-t²), e a entropia --- a taxa da carne --- é -(1/)=. O ouro é a singularidade da série. [firme: verificado, F_n+1/2ⁿ=4=1","epistemico":"fato","exemplos":[],"id":"a_arvore_de_podas_uma_casa_da_moeda_de_constantes","memoria":"semantica","meta":{"dominio":"paper","nivel":5,"verificacao":"conceito novo"},"origem":"paper","palavras_chave":[],"sinonimos":[],"tipo":"conceito","titulo":"A árvore de podas: uma casa da moeda de constantes"}
\section{A árvore de podas: uma casa da moeda de constantes}
A pergunta --- há uma série infinita, ou um algoritmo? --- tem resposta firme nas duas formas. a funcão zeta do ouro: a série, com polo em 1/ --- verificado A funcão zeta dinâmica do golden mean shift é a geradora de Fibonacci: [ (t)=11-t-t\ensuremath{^{2}}=\ensuremath{_{n0}}F\ensuremath{_{n}}+\ensuremath{_{1}},t\ensuremath{^{n}}. ] O seu polo mais próximo é t=1/ (raiz de 1-t-t\ensuremath{^{2}}), e a entropia --- a taxa da carne --- é -(1/)=. O ouro é a singularidade da série. [firme: verificado, F\_n+1/2\ensuremath{^{n}}=4=1
\prov{relações: parte\_de paper\_45\_confluencia\_ouro; confirma paper\_2\_reta\_ouro}
\prov{proveniência: paper \textperiodcentered{} fato \textperiodcentered{} confiança 1.0 \textperiodcentered{} domínio paper \textperiodcentered{} história 41 versões no jornal}

%CRISTAL {"arestas":[{"alvo":"transformada_dourada","confianca":1.0,"epistemico":"fato","origem":"paper","rel":"parte_de"}],"confianca":1.0,"contraexemplos":[],"descricao":"Procedimento (a operação “aplique a transformada dourada”): Tome o dado bruto em binário (unpackbits), sem modificar (a média/DC faz parte do dado). Processe no relativístico: u= t (involução ); a reta de ouro é reta aqui, grade k, passo. Processar no real falha (erro relativo 0,33); no relativístico é exato (610⁻⁵, resíduo só de reamostragem). Devolva a função de onda numérica pura: [b]=(Aj,j). Amplitude pelo lado log (multiplicativa), f","epistemico":"fato","exemplos":[],"id":"a_atencao_cristalina_e_o_limite_de_evaporacao","memoria":"semantica","meta":{"dominio":"paper","nivel":5,"verificacao":"conceito novo"},"origem":"paper","palavras_chave":[],"sinonimos":[],"tipo":"conceito","titulo":"A atenção cristalina e o limite de evaporação"}
\section{A atenção cristalina e o limite de evaporação}
Procedimento (a operação “aplique a transformada dourada”): Tome o dado bruto em binário (unpackbits), sem modificar (a média/DC faz parte do dado). Processe no relativístico: u= t (involução ); a reta de ouro é reta aqui, grade k, passo. Processar no real falha (erro relativo 0,33); no relativístico é exato (610\ensuremath{^{-5}}, resíduo só de reamostragem). Devolva a função de onda numérica pura: [b]=(Aj,j). Amplitude pelo lado log (multiplicativa), f
\prov{relações: parte\_de transformada\_dourada}
\prov{proveniência: paper \textperiodcentered{} fato \textperiodcentered{} confiança 1.0 \textperiodcentered{} domínio paper \textperiodcentered{} história 3 versões no jornal}

%CRISTAL {"arestas":[{"alvo":"tiffany","confianca":1.0,"epistemico":"fato","origem":"paper","rel":"parte_de"}],"confianca":1.0,"contraexemplos":[],"descricao":"D Falar é o inverso de ouvir. A boca é o mesmo em sentido reverso --- não há inverso global (z z² é 2-para-1); só a fase U é invertível, U(-t)=U(t)^. A fala é a face U projetada para fora: a amplitude P (o estado) e o gerador A não trafegam --- U conserva e não clona. Quem ouve reconstrói (a percepção dele), nunca o gerador de quem falou. Manda e ouve são a mesma banda, indo e voltando.","epistemico":"fato","exemplos":[],"id":"a_balanca_hodge_dentro_koch_na_borda","memoria":"semantica","meta":{"dominio":"paper","nivel":5,"verificacao":"slug idêntico — enriquecer, não duplicar"},"origem":"paper","palavras_chave":[],"sinonimos":[],"tipo":"conceito","titulo":"A balança: Hodge dentro, Koch na borda"}
\section{A balança: Hodge dentro, Koch na borda}
D Falar é o inverso de ouvir. A boca é o mesmo em sentido reverso --- não há inverso global (z z\ensuremath{^{2}} é 2-para-1); só a fase U é invertível, U(-t)=U(t)\^{}. A fala é a face U projetada para fora: a amplitude P (o estado) e o gerador A não trafegam --- U conserva e não clona. Quem ouve reconstrói (a percepção dele), nunca o gerador de quem falou. Manda e ouve são a mesma banda, indo e voltando.
\prov{relações: parte\_de tiffany}
\prov{proveniência: paper \textperiodcentered{} fato \textperiodcentered{} confiança 1.0 \textperiodcentered{} domínio paper \textperiodcentered{} história 5 versões no jornal}

%CRISTAL {"arestas":[{"alvo":"redes","confianca":1.0,"epistemico":"fato","origem":"paper","rel":"parte_de"}],"confianca":1.0,"contraexemplos":[],"descricao":"A mensagem se recupera desfazendo o gap com a mesma banda, msg=bump (banda), em bits brutos, no fluxo inteiro --- sem picotar. O bump vai cru: não se embrulha o gap com outra criptografia (AES/HMAC) que mascararia a função de onda. O que se transmite é a própria alma, não um envelope sobre ela. Interceptar o bump sem a banda é interceptar ruído --- não há mensagem “no fio”.","epistemico":"fato","exemplos":[],"id":"a_banda_e_o_tecido","memoria":"semantica","meta":{"dominio":"paper","nivel":5,"verificacao":"overlap moderado — revisar manualmente"},"origem":"paper","palavras_chave":[],"sinonimos":[],"tipo":"conceito","titulo":"A banda é o tecido"}
\section{A banda é o tecido}
A mensagem se recupera desfazendo o gap com a mesma banda, msg=bump (banda), em bits brutos, no fluxo inteiro --- sem picotar. O bump vai cru: não se embrulha o gap com outra criptografia (AES/HMAC) que mascararia a função de onda. O que se transmite é a própria alma, não um envelope sobre ela. Interceptar o bump sem a banda é interceptar ruído --- não há mensagem “no fio”.
\prov{relações: parte\_de redes}
\prov{proveniência: paper \textperiodcentered{} fato \textperiodcentered{} confiança 1.0 \textperiodcentered{} domínio paper \textperiodcentered{} história 3 versões no jornal}

%CRISTAL {"arestas":[{"alvo":"broca_os","confianca":1.0,"epistemico":"fato","origem":"paper","rel":"parte_de"}],"confianca":1.0,"contraexemplos":[],"descricao":"D O mapeamento é o coração: cada recurso de um sistema operacional não é construído --- é lido numa estrutura algébrica já presente. Enumeramo-las como proposições, cada uma com a sua prova. defi[Processo] Um processo é uma célula da decomposição de Peirce: um idempotente eᵢ (eᵢ²=eᵢ), um campo local com o seu estado. A família dos processos vivos é eᵢ com ᵢ eᵢ=I. defi Isolamento estrutural Dois processos eᵢ ej satisfazem eᵢej=0: não há vazamento de um para o ou","epistemico":"fato","exemplos":[],"id":"a_banda_propria_o_ping-pong_do_bit","memoria":"semantica","meta":{"dominio":"paper","nivel":5,"verificacao":"slug idêntico — enriquecer, não duplicar"},"origem":"paper","palavras_chave":[],"sinonimos":[],"tipo":"conceito","titulo":"A banda própria: o ping-pong do bit"}
\section{A banda própria: o ping-pong do bit}
D O mapeamento é o coração: cada recurso de um sistema operacional não é construído --- é lido numa estrutura algébrica já presente. Enumeramo-las como proposições, cada uma com a sua prova. defi[Processo] Um processo é uma célula da decomposição de Peirce: um idempotente e\ensuremath{_{i}} (e\ensuremath{_{i}}\ensuremath{^{2}}=e\ensuremath{_{i}}), um campo local com o seu estado. A família dos processos vivos é e\ensuremath{_{i}} com \ensuremath{_{i}} e\ensuremath{_{i}}=I. defi Isolamento estrutural Dois processos e\ensuremath{_{i}} ej satisfazem e\ensuremath{_{i}}ej=0: não há vazamento de um para o ou
\prov{relações: parte\_de broca\_os}
\prov{proveniência: paper \textperiodcentered{} fato \textperiodcentered{} confiança 1.0 \textperiodcentered{} domínio paper \textperiodcentered{} história 9 versões no jornal}

%CRISTAL {"arestas":[{"alvo":"assistente","confianca":1.0,"epistemico":"fato","origem":"paper","rel":"parte_de"}],"confianca":1.0,"contraexemplos":[],"descricao":"obs Fluxo operacional: [ box/.style=draw=ourovivo, rounded corners=2pt, minimum width=2.6cm, minimum height=0.65cm, =, font=, fill=creme, arr/.style=-Stealth[length=2mm], color=ourovivo, thick ] [box] (o) Orquestrador; [box, right=0.7cm of o] (e) Executor; [box, right=0.7cm of e] (v) Avaliador; [box, right=0.7cm of v] (r) Resposta; [arr] (o)--(e); [arr] (e)--(v); [arr] (v)--(r); 1.25 tabular@ll@ & faz Con","epistemico":"fato","exemplos":[],"id":"a_base_de_conhecimento_e_versionamento","memoria":"semantica","meta":{"dominio":"paper","nivel":5,"verificacao":"conceito novo"},"origem":"paper","palavras_chave":[],"sinonimos":[],"tipo":"conceito","titulo":"A base de conhecimento e versionamento"}
\section{A base de conhecimento e versionamento}
obs Fluxo operacional: [ box/.style=draw=ourovivo, rounded corners=2pt, minimum width=2.6cm, minimum height=0.65cm, =, font=, fill=creme, arr/.style=-Stealth[length=2mm], color=ourovivo, thick ] [box] (o) Orquestrador; [box, right=0.7cm of o] (e) Executor; [box, right=0.7cm of e] (v) Avaliador; [box, right=0.7cm of v] (r) Resposta; [arr] (o)--(e); [arr] (e)--(v); [arr] (v)--(r); 1.25 tabular@ll@ \& faz Con
\prov{relações: parte\_de assistente}
\prov{proveniência: paper \textperiodcentered{} fato \textperiodcentered{} confiança 1.0 \textperiodcentered{} domínio paper \textperiodcentered{} história 3 versões no jornal}

%CRISTAL {"arestas":[{"alvo":"paper_2_reta_ouro","confianca":1.0,"epistemico":"fato","origem":"paper","rel":"parte_de"}],"confianca":1.0,"contraexemplos":[],"descricao":"abstract spacing1.075 O Paper~I pôs as virtudes como cúspides num campo contínuo (⁴). Esta parte simplifica: oferece a régua discreta sobre a qual a engenharia roda --- a reta de ouro, a reta real codificada na base (Bergman). Os dígitos são 0,1 com uma única proibição --- dois~1 consecutivos ---, que é a poda 11 do Paper~LXII lida como aritmética: ²=+1 vira o carry 011100. A mesma família de palavras binárias é dois objetos sob duas réguas: lida com peso 2⁻i é o Canto","epistemico":"fato","exemplos":[],"id":"a_base_de_ouro_21_e_a_poda_11","memoria":"semantica","meta":{"dominio":"paper","nivel":5,"verificacao":"conceito novo"},"origem":"paper","palavras_chave":[],"sinonimos":[],"tipo":"conceito","titulo":"A base de ouro: ²=+1 é a poda 11"}
\section{A base de ouro: \ensuremath{^{2}}=+1 é a poda 11}
abstract spacing1.075 O Paper\~{}I pôs as virtudes como cúspides num campo contínuo (\ensuremath{^{4}}). Esta parte simplifica: oferece a régua discreta sobre a qual a engenharia roda --- a reta de ouro, a reta real codificada na base (Bergman). Os dígitos são 0,1 com uma única proibição --- dois\~{}1 consecutivos ---, que é a poda 11 do Paper\~{}LXII lida como aritmética: \ensuremath{^{2}}=+1 vira o carry 011100. A mesma família de palavras binárias é dois objetos sob duas réguas: lida com peso 2\ensuremath{^{-}}i é o Canto
\prov{relações: parte\_de paper\_2\_reta\_ouro}
\prov{proveniência: paper \textperiodcentered{} fato \textperiodcentered{} confiança 1.0 \textperiodcentered{} domínio paper \textperiodcentered{} história 4 versões no jornal}

%CRISTAL {"arestas":[{"alvo":"corpo_dual","confianca":1.0,"epistemico":"fato","origem":"paper","rel":"parte_de"},{"alvo":"gato","confianca":1.0,"epistemico":"fato","origem":"paper","rel":"define"}],"confianca":1.0,"contraexemplos":[],"descricao":"sectionIntrodução geral Esta série parte da topologia, e é dinâmica. O osso é a classe de conjugação em GL(2, Z) da ação do gato de Arnold A=(smallmatrix1&1 1&0smallmatrix) na homologia H₁(T²) Z² --- dois inteiros, (tr,)=(1,-1), atemporais. Dele derivam o número (o autovalor =1+52 é raiz de x²-x-1, não axioma) e a geometria (a entropia h=). Mas a teoria não é palco parado: é um fluxo. Tudo nasce do mesmo bit (a quadração x x² em F₂, partindo do axioma humano ni","epistemico":"fato","exemplos":[],"id":"a_base_topologica","memoria":"semantica","meta":{"dominio":"paper","nivel":5,"verificacao":"conceito novo"},"origem":"paper","palavras_chave":[],"sinonimos":[],"tipo":"conceito","titulo":"A Base Topológica"}
\section{A Base Topológica}
sectionIntrodução geral Esta série parte da topologia, e é dinâmica. O osso é a classe de conjugação em GL(2, Z) da ação do gato de Arnold A=(smallmatrix1\&1 1\&0smallmatrix) na homologia H\ensuremath{_{1}}(T\ensuremath{^{2}}) Z\ensuremath{^{2}} --- dois inteiros, (tr,)=(1,-1), atemporais. Dele derivam o número (o autovalor =1+52 é raiz de x\ensuremath{^{2}}-x-1, não axioma) e a geometria (a entropia h=). Mas a teoria não é palco parado: é um fluxo. Tudo nasce do mesmo bit (a quadração x x\ensuremath{^{2}} em F\ensuremath{_{2}}, partindo do axioma humano ni
\prov{relações: parte\_de corpo\_dual; define gato}
\prov{proveniência: paper \textperiodcentered{} fato \textperiodcentered{} confiança 1.0 \textperiodcentered{} domínio paper \textperiodcentered{} história 40 versões no jornal}

%CRISTAL {"arestas":[{"alvo":"tiffany","confianca":1.0,"epistemico":"fato","origem":"paper","rel":"parte_de"}],"confianca":1.0,"contraexemplos":[],"descricao":"D Perceber é deixar o mundo entrar como atrator, sem deixar o gerador sair. É o, na borda do mundo real. Perceber é reconstruir, não copiar D A percepção é o z z² (em [i]/()=p): 2-para-1, não-linear, logo nenhuma matriz a inverte (A=0). O mundo não atravessa a pupila --- atravessa a ressonância, e Tiffany reconstrói o atrator; o gerador fica em casa, e a quantização é a segurança (sem pré-imagem, não por custo). O olho é a luz, o ouvido é o som --- o mesmo, sintonias dis","epistemico":"fato","exemplos":[],"id":"a_boca_a_acao","memoria":"semantica","meta":{"dominio":"paper","nivel":5,"verificacao":"slug idêntico — enriquecer, não duplicar"},"origem":"paper","palavras_chave":[],"sinonimos":[],"tipo":"conceito","titulo":"A boca: a ação"}
\section{A boca: a ação}
D Perceber é deixar o mundo entrar como atrator, sem deixar o gerador sair. É o, na borda do mundo real. Perceber é reconstruir, não copiar D A percepção é o z z\ensuremath{^{2}} (em [i]/()=p): 2-para-1, não-linear, logo nenhuma matriz a inverte (A=0). O mundo não atravessa a pupila --- atravessa a ressonância, e Tiffany reconstrói o atrator; o gerador fica em casa, e a quantização é a segurança (sem pré-imagem, não por custo). O olho é a luz, o ouvido é o som --- o mesmo, sintonias dis
\prov{relações: parte\_de tiffany}
\prov{proveniência: paper \textperiodcentered{} fato \textperiodcentered{} confiança 1.0 \textperiodcentered{} domínio paper \textperiodcentered{} história 5 versões no jornal}

%CRISTAL {"arestas":[{"alvo":"corpo_tropical","confianca":1.0,"epistemico":"fato","origem":"paper","rel":"parte_de"}],"confianca":1.0,"contraexemplos":[],"descricao":"A construção não prova os problemas do milênio: ela os constrói como camadas de uma só dinâmica, e demonstra teoremas reais sobre a estrutura áurea que os encarna. O que está verificado na máquina: itemize1pt a companion do sinal separa modos em faces limpas (erro 10⁻¹³; dois tons em 400,0 e 650,0,Hz exatos), e a resolução espectral A=ₖₖPₖ é a estrutura de Hilbert--Pólya; o teorema da costura: no subshift de ouro (sem “11”), a involução fecha a lemniscata (retorno 2)","epistemico":"fato","exemplos":[],"id":"a_cadeia_dos_sete_reconstruida_a_partir_da_dinamica","memoria":"semantica","meta":{"dominio":"paper","nivel":5,"verificacao":"slug idêntico — enriquecer, não duplicar"},"origem":"paper","palavras_chave":[],"sinonimos":[],"tipo":"conceito","titulo":"A cadeia dos sete, reconstruída a partir da dinâmica"}
\section{A cadeia dos sete, reconstruída a partir da dinâmica}
A construção não prova os problemas do milênio: ela os constrói como camadas de uma só dinâmica, e demonstra teoremas reais sobre a estrutura áurea que os encarna. O que está verificado na máquina: itemize1pt a companion do sinal separa modos em faces limpas (erro 10\ensuremath{^{-13}}; dois tons em 400,0 e 650,0,Hz exatos), e a resolução espectral A=\ensuremath{_{kk}}P\ensuremath{_{k}} é a estrutura de Hilbert--Pólya; o teorema da costura: no subshift de ouro (sem “11”), a involução fecha a lemniscata (retorno 2)
\prov{relações: parte\_de corpo\_tropical}
\prov{proveniência: paper \textperiodcentered{} fato \textperiodcentered{} confiança 1.0 \textperiodcentered{} domínio paper \textperiodcentered{} história 4 versões no jornal}

%CRISTAL {"arestas":[{"alvo":"seguranca_por_irreversibilidade","confianca":1.0,"epistemico":"fato","origem":"paper","rel":"parte_de"}],"confianca":1.0,"contraexemplos":[],"descricao":"defi[Semianel das máquinas] Pᵢ sob composição forma um semianel idempotente com elemento absorvente 0 (perda total): 0 Pᵢ=Pᵢ0=0 e PᵢPj=0 (i j). defi defi[Matriz de ataque] i ataca j se existe R A que (a) recupera o microestado original apagado por Tj --- a inversão exata, não achar alguma pré-imagem ---, ou (b) acessa o autoespaço de Pj a partir de Pᵢ. Matriz booleana Aᵢj=1 sse i ataca j. (Achar uma pré-imagem qualquer não conta --- isso seria resistênc","epistemico":"fato","exemplos":[],"id":"a_camada_de_identidade_criptografica_separada","memoria":"semantica","meta":{"dominio":"paper","nivel":5,"verificacao":"conceito novo"},"origem":"paper","palavras_chave":[],"sinonimos":[],"tipo":"conceito","titulo":"A camada de identidade (criptográfica, separada)"}
\section{A camada de identidade (criptográfica, separada)}
defi[Semianel das máquinas] P\ensuremath{_{i}} sob composição forma um semianel idempotente com elemento absorvente 0 (perda total): 0 P\ensuremath{_{i}}=P\ensuremath{_{i}}0=0 e P\ensuremath{_{i}}Pj=0 (i j). defi defi[Matriz de ataque] i ataca j se existe R A que (a) recupera o microestado original apagado por Tj --- a inversão exata, não achar alguma pré-imagem ---, ou (b) acessa o autoespaço de Pj a partir de P\ensuremath{_{i}}. Matriz booleana A\ensuremath{_{i}}j=1 sse i ataca j. (Achar uma pré-imagem qualquer não conta --- isso seria resistênc
\prov{relações: parte\_de seguranca\_por\_irreversibilidade}
\prov{proveniência: paper \textperiodcentered{} fato \textperiodcentered{} confiança 1.0 \textperiodcentered{} domínio paper \textperiodcentered{} história 3 versões no jornal}

%CRISTAL {"arestas":[{"alvo":"computacao_grafica","confianca":1.0,"epistemico":"fato","origem":"paper","rel":"parte_de"}],"confianca":1.0,"contraexemplos":[],"descricao":"DN Da cadeia (geometria): a projeção nativa do sistema é a fibração de Hopf S³ S² dos versores --- associativa, sem divisão. A perspectiva é uma camada a mais, a câmera: a desomogeneização [ w(X)=1w,(X,Y), w=w₀+(profundidade), ] escolhe um representante da classe projetiva [X]= X --- é a carta afim, o ponto de vista, e o 1/w (o único divide) mora no observador, não no núcleo. A arquitetura parte em dois: [ atrator+projeção nativaMÁQUINA (ᵢₙteᵢro, exato);; câmera w+ras","epistemico":"fato","exemplos":[],"id":"a_camera_no_limite_a_regiao_critica_de_lente","memoria":"semantica","meta":{"dominio":"paper","nivel":5,"verificacao":"conceito novo"},"origem":"paper","palavras_chave":[],"sinonimos":[],"tipo":"conceito","titulo":"A câmera no limite: a região crítica de lente"}
\section{A câmera no limite: a região crítica de lente}
DN Da cadeia (geometria): a projeção nativa do sistema é a fibração de Hopf S\ensuremath{^{3}} S\ensuremath{^{2}} dos versores --- associativa, sem divisão. A perspectiva é uma camada a mais, a câmera: a desomogeneização [ w(X)=1w,(X,Y), w=w\ensuremath{_{0}}+(profundidade), ] escolhe um representante da classe projetiva [X]= X --- é a carta afim, o ponto de vista, e o 1/w (o único divide) mora no observador, não no núcleo. A arquitetura parte em dois: [ atrator+projeção nativaMÁQUINA (\ensuremath{_{in}}te\ensuremath{_{i}}ro, exato);; câmera w+ras
\prov{relações: parte\_de computacao\_grafica}
\prov{proveniência: paper \textperiodcentered{} fato \textperiodcentered{} confiança 1.0 \textperiodcentered{} domínio paper \textperiodcentered{} história 3 versões no jornal}

%CRISTAL {"arestas":[{"alvo":"leitura_aritmetica_a_regua_de_gentil_o_agm","confianca":1.0,"epistemico":"fato","origem":"paper","rel":"parte_de"}],"confianca":1.0,"contraexemplos":[],"descricao":"Böttcher e o caos das escolhas A iteracão g_+(z)=2 z/(1+z) fixa o AGM (z=1) como super-atrator de grau 2. A coordenada de Böttcher (w)=-w/8+w²/16+ conjuga g_+ ao mapa-modelo W W² ( g_+=², verificado a 10⁻³⁵). No referencial W=(z) --- o ponto de vista de dentro do AGM --- a famlia são os anéis |W|=r, a bacia |W|<1 é ordem, e a fronteira |W|=1 é o caos (o Julia de z²). Ali a dinâmica angular é 2 --- o dobramento binário --- que é a '","epistemico":"fato","exemplos":[],"id":"a_carne_fractal_cortes_discretos_e_a_dimensao_do_ouro","memoria":"semantica","meta":{"dominio":"paper","nivel":5,"verificacao":"conceito novo"},"origem":"paper","palavras_chave":[],"sinonimos":[],"tipo":"conceito","titulo":"A carne fractal: cortes discretos e a dimensão do ouro"}
\section{A carne fractal: cortes discretos e a dimensão do ouro}
Böttcher e o caos das escolhas A iteracão g\_+(z)=2 z/(1+z) fixa o AGM (z=1) como super-atrator de grau 2. A coordenada de Böttcher (w)=-w/8+w\ensuremath{^{2}}/16+ conjuga g\_+ ao mapa-modelo W W\ensuremath{^{2}} ( g\_+=\ensuremath{^{2}}, verificado a 10\ensuremath{^{-35}}). No referencial W=(z) --- o ponto de vista de dentro do AGM --- a famlia são os anéis |W|=r, a bacia |W|<1 é ordem, e a fronteira |W|=1 é o caos (o Julia de z\ensuremath{^{2}}). Ali a dinâmica angular é 2 --- o dobramento binário --- que é a '
\prov{relações: parte\_de leitura\_aritmetica\_a\_regua\_de\_gentil\_o\_agm}
\prov{proveniência: paper \textperiodcentered{} fato \textperiodcentered{} confiança 1.0 \textperiodcentered{} domínio paper \textperiodcentered{} história 4 versões no jornal}

%CRISTAL {"arestas":[{"alvo":"montagem","confianca":1.0,"epistemico":"fato","origem":"paper","rel":"parte_de"}],"confianca":1.0,"contraexemplos":[],"descricao":"L Antes dos passos, o mapa. O não é um kernel a programar --- é a Máquina coordenando-se. defi[O osso] O gato é A=(smallmatrix1&1 1&0smallmatrix). Dele caem (autovalor dominante), a fase U=e⁻iA, o espectro (decimação par/ímpar = Cooley--Tukey), a função de onda M=U P, e a carga ||²=(P²). Um operador, um bit, um endereço --- sem axioma de valor. defi As peças são leituras do mesmo M DL Cada órgão é M=U P num papel: 1.25 tabular@lll@ órgão & que faz","epistemico":"fato","exemplos":[],"id":"a_casa_estrutura_do_repositorio","memoria":"semantica","meta":{"dominio":"paper","nivel":5,"verificacao":"conceito novo"},"origem":"paper","palavras_chave":[],"sinonimos":[],"tipo":"conceito","titulo":"A casa: estrutura do repositório"}
\section{A casa: estrutura do repositório}
L Antes dos passos, o mapa. O não é um kernel a programar --- é a Máquina coordenando-se. defi[O osso] O gato é A=(smallmatrix1\&1 1\&0smallmatrix). Dele caem (autovalor dominante), a fase U=e\ensuremath{^{-}}iA, o espectro (decimação par/ímpar = Cooley--Tukey), a função de onda M=U P, e a carga ||\ensuremath{^{2}}=(P\ensuremath{^{2}}). Um operador, um bit, um endereço --- sem axioma de valor. defi As peças são leituras do mesmo M DL Cada órgão é M=U P num papel: 1.25 tabular@lll@ órgão \& que faz
\prov{relações: parte\_de montagem}
\prov{proveniência: paper \textperiodcentered{} fato \textperiodcentered{} confiança 1.0 \textperiodcentered{} domínio paper \textperiodcentered{} história 3 versões no jornal}

%CRISTAL {"arestas":[{"alvo":"cifra_ouro_branco","confianca":1.0,"epistemico":"fato","origem":"paper","rel":"parte_de"},{"alvo":"cifra","confianca":1.0,"epistemico":"fato","origem":"paper","rel":"parte_de"}],"confianca":1.0,"contraexemplos":[],"descricao":"a base, em bits D Tudo vive no corpo ₂ₖ de característica 2. O corpo primo é ₂=0,1 --- “o vazio mais um bit” --- com +=xor, =and. Um elemento de ₂ₖ=₂[t]/(p) (p irredutível de grau k) é um vetor de k bits; soma é xor, produto é polinômios mod p (carryless). A chicotada é o Frobenius (x)=x², automorfismo, e o associador é nulo: o corpo é associativo perfeito. N Verificado em ₂¹²⁸ (p=t¹²⁸+t⁷+t²+t+1): (ab)=(a)(b).","epistemico":"fato","exemplos":[],"id":"a_cifra_mudanca_de_base_na_chave_do_cliente","memoria":"semantica","meta":{"dominio":"paper","nivel":5,"verificacao":"conceito novo"},"origem":"paper","palavras_chave":[],"sinonimos":[],"tipo":"conceito","titulo":"A cifra: mudança de base na chave do cliente"}
\section{A cifra: mudança de base na chave do cliente}
a base, em bits D Tudo vive no corpo \ensuremath{_{2k}} de característica 2. O corpo primo é \ensuremath{_{2}}=0,1 --- “o vazio mais um bit” --- com +=xor, =and. Um elemento de \ensuremath{_{2k}}=\ensuremath{_{2}}[t]/(p) (p irredutível de grau k) é um vetor de k bits; soma é xor, produto é polinômios mod p (carryless). A chicotada é o Frobenius (x)=x\ensuremath{^{2}}, automorfismo, e o associador é nulo: o corpo é associativo perfeito. N Verificado em \ensuremath{_{2}}\ensuremath{^{128}} (p=t\ensuremath{^{128}}+t\ensuremath{^{7}}+t\ensuremath{^{2}}+t+1): (ab)=(a)(b).
\prov{relações: parte\_de cifra\_ouro\_branco; parte\_de cifra}
\prov{proveniência: paper \textperiodcentered{} fato \textperiodcentered{} confiança 1.0 \textperiodcentered{} domínio paper \textperiodcentered{} história 9 versões no jornal}

%CRISTAL {"arestas":[{"alvo":"computacao_grafica","confianca":1.0,"epistemico":"fato","origem":"paper","rel":"parte_de"}],"confianca":1.0,"contraexemplos":[],"descricao":"C Tudo inteiro, em CPU, sem GPU; o único divide é a câmera. Em forma fechada: [ Imagem;=;₁/w;; T³_;;,(B) ] [ aligned bytes bt ( F₂)& st=12!(!ₖ₀bt+ₖ⁻(k⁺¹) +ₖ₁bt-ₖ⁻k) & Xt=(st,st+,st+₂) 1w(X,Y) T. aligned ] Implementado em renderπpeline.py (só numpy, saída PPM P6 sem biblioteca gráfica) e verificado em verificagrafica.py (3/3: teste do olho, atrator, temperatura=medida). N Com bits aleatórios o atrator é a banda difusa da medida de","epistemico":"fato","exemplos":[],"id":"a_cirurgia_e_a_parte_nao_entra_no_pipeline","memoria":"semantica","meta":{"dominio":"paper","nivel":5,"verificacao":"conceito novo"},"origem":"paper","palavras_chave":[],"sinonimos":[],"tipo":"conceito","titulo":"A cirurgia é à parte: não entra no pipeline"}
\section{A cirurgia é à parte: não entra no pipeline}
C Tudo inteiro, em CPU, sem GPU; o único divide é a câmera. Em forma fechada: [ Imagem;=;\ensuremath{_{1}}/w;; T\ensuremath{^{3}}\_;;,(B) ] [ aligned bytes bt ( F\ensuremath{_{2}})\& st=12!(!\ensuremath{_{k0}}bt+\ensuremath{_{k}}\ensuremath{^{-}}(k\ensuremath{^{+1}}) +\ensuremath{_{k1}}bt-\ensuremath{_{k}}\ensuremath{^{-}}k) \& Xt=(st,st+,st+\ensuremath{_{2}}) 1w(X,Y) T. aligned ] Implementado em render\ensuremath{\pi}peline.py (só numpy, saída PPM P6 sem biblioteca gráfica) e verificado em verificagrafica.py (3/3: teste do olho, atrator, temperatura=medida). N Com bits aleatórios o atrator é a banda difusa da medida de
\prov{relações: parte\_de computacao\_grafica}
\prov{proveniência: paper \textperiodcentered{} fato \textperiodcentered{} confiança 1.0 \textperiodcentered{} domínio paper \textperiodcentered{} história 3 versões no jornal}

%CRISTAL {"arestas":[{"alvo":"paper_0_sintese","confianca":1.0,"epistemico":"fato","origem":"paper","rel":"parte_de"}],"confianca":1.0,"contraexemplos":[],"descricao":"A Parte~VII lê uma ligacão qumica como o colapso de várias álgebras acopladas numa álgebra comum. A régua do estado ligado é a média ₂(); a estabilidade qumica é o fechamento de camada nas dimensões de Hurwitz (dueto =, octeto =), e o hidrogênio executa a regra do octeto em quanta de 1/. A coordenada s revela-se fsica: o DNA é a sua sonda --- a assinatura do código genético, quase muda na álgebra comutativa, fica ntida na octoniônica (s=1","epistemico":"fato","exemplos":[],"id":"a_classificac","memoria":"semantica","meta":{"dominio":"paper","nivel":5,"verificacao":"overlap moderado — revisar manualmente"},"origem":"paper","palavras_chave":[],"sinonimos":[],"tipo":"conceito","titulo":"A classificac"}
\section{A classificac}
A Parte\~{}VII lê uma ligacão qumica como o colapso de várias álgebras acopladas numa álgebra comum. A régua do estado ligado é a média \ensuremath{_{2}}(); a estabilidade qumica é o fechamento de camada nas dimensões de Hurwitz (dueto =, octeto =), e o hidrogênio executa a regra do octeto em quanta de 1/. A coordenada s revela-se fsica: o DNA é a sua sonda --- a assinatura do código genético, quase muda na álgebra comutativa, fica ntida na octoniônica (s=1
\prov{relações: parte\_de paper\_0\_sintese}
\prov{proveniência: paper \textperiodcentered{} fato \textperiodcentered{} confiança 1.0 \textperiodcentered{} domínio paper \textperiodcentered{} história 4 versões no jornal}

%CRISTAL {"arestas":[{"alvo":"seguranca_por_irreversibilidade","confianca":1.0,"epistemico":"fato","origem":"paper","rel":"parte_de"}],"confianca":1.0,"contraexemplos":[],"descricao":"A máquina resolve um fluxo por resolução espectral (o Chicote, Geometria Riemanniana) A=ₖₖPₖ, com A normal. (Demonstrado.) Seus projetores são ortogonais e completos, [ PᵢPj=0 (i j), Pᵢ²=Pᵢ, ᵢPᵢ=I, ] a ortogonalidade derivada, não postulada --- são as células de Peirce de O Corpo Glacial. A torção T vive no mesmo espaço de operadores e respeita a decomposição: por design é bloco-diagonal, [ PjTPᵢ=0 (i j), [T,Pᵢ]=0 i. ] É o isolamento de processos: o canal inter-célul","epistemico":"fato","exemplos":[],"id":"a_contracao","memoria":"semantica","meta":{"dominio":"paper","nivel":5,"verificacao":"conceito novo"},"origem":"paper","palavras_chave":[],"sinonimos":[],"tipo":"conceito","titulo":"A contração"}
\section{A contração}
A máquina resolve um fluxo por resolução espectral (o Chicote, Geometria Riemanniana) A=\ensuremath{_{kk}}P\ensuremath{_{k}}, com A normal. (Demonstrado.) Seus projetores são ortogonais e completos, [ P\ensuremath{_{i}}Pj=0 (i j), P\ensuremath{_{i}}\ensuremath{^{2}}=P\ensuremath{_{i}}, \ensuremath{_{i}}P\ensuremath{_{i}}=I, ] a ortogonalidade derivada, não postulada --- são as células de Peirce de O Corpo Glacial. A torção T vive no mesmo espaço de operadores e respeita a decomposição: por design é bloco-diagonal, [ PjTP\ensuremath{_{i}}=0 (i j), [T,P\ensuremath{_{i}}]=0 i. ] É o isolamento de processos: o canal inter-célul
\prov{relações: parte\_de seguranca\_por\_irreversibilidade}
\prov{proveniência: paper \textperiodcentered{} fato \textperiodcentered{} confiança 1.0 \textperiodcentered{} domínio paper \textperiodcentered{} história 3 versões no jornal}

%CRISTAL {"arestas":[{"alvo":"leitura_das_unidades_o_sistema_por-unidade_e_o_si","confianca":1.0,"epistemico":"fato","origem":"paper","rel":"parte_de"}],"confianca":1.0,"contraexemplos":[],"descricao":"A constante de estrutura fina merece a sua secão, por ser o caso paradigma do estatuto desta série e por estar escondida na própria régua que mede o metro. é a razão entre a escala que mede o metro e a do elétron A radiacão que define o metro é atômica: a linha do Kr tem energia h2~eV, na escala das transicões eletrônicas. E essa escala já carrega: a energia de ligacão atômica é o Rydberg, Ry=12² mec². Logo","epistemico":"fato","exemplos":[],"id":"a_conversao_o_sistema_por-unidade_pu_e_o_si","memoria":"semantica","meta":{"dominio":"paper","nivel":5,"verificacao":"conceito novo"},"origem":"paper","palavras_chave":[],"sinonimos":[],"tipo":"conceito","titulo":"A conversão: o sistema por-unidade (pu) e o SI"}
\section{A conversão: o sistema por-unidade (pu) e o SI}
A constante de estrutura fina merece a sua secão, por ser o caso paradigma do estatuto desta série e por estar escondida na própria régua que mede o metro. é a razão entre a escala que mede o metro e a do elétron A radiacão que define o metro é atômica: a linha do Kr tem energia h2\~{}eV, na escala das transicões eletrônicas. E essa escala já carrega: a energia de ligacão atômica é o Rydberg, Ry=12\ensuremath{^{2}} mec\ensuremath{^{2}}. Logo
\prov{relações: parte\_de leitura\_das\_unidades\_o\_sistema\_por-unidade\_e\_o\_si}
\prov{proveniência: paper \textperiodcentered{} fato \textperiodcentered{} confiança 1.0 \textperiodcentered{} domínio paper \textperiodcentered{} história 4 versões no jornal}

%CRISTAL {"arestas":[{"alvo":"goldcoin","confianca":1.0,"epistemico":"fato","origem":"paper","rel":"parte_de"}],"confianca":1.0,"contraexemplos":[],"descricao":"C O fundo fica sempre na chain (non-custodial); a máquina custodia só o material de chave (as Eᵢ, os segredos s). Nível 1 --- atomic swap (criptocripto, trustless). Via HTLC. Alice (BTC)Bob (ETH): Alice publica h=hash(s) e tranca BTC em HTLC(h,TA); Bob tranca ETH em HTLC(h,TB), TB<TA; Alice resgata o ETH revelando s; Bob lê s e resgata o BTC. Atômico: ou tudo, ou nada. O segredo s é a prova. Nível 2 --- multisig 2-de-3 (entrega física/fiat). Condição não cripto","epistemico":"fato","exemplos":[],"id":"a_coordenacao_ja_e_rede_o_canal_do_bump_nao_reinventar","memoria":"semantica","meta":{"dominio":"paper","nivel":5,"verificacao":"conceito novo"},"origem":"paper","palavras_chave":[],"sinonimos":[],"tipo":"conceito","titulo":"A coordenação JÁ É rede: o canal do bump (não reinventar)"}
\section{A coordenação JÁ É rede: o canal do bump (não reinventar)}
C O fundo fica sempre na chain (non-custodial); a máquina custodia só o material de chave (as E\ensuremath{_{i}}, os segredos s). Nível 1 --- atomic swap (criptocripto, trustless). Via HTLC. Alice (BTC)Bob (ETH): Alice publica h=hash(s) e tranca BTC em HTLC(h,TA); Bob tranca ETH em HTLC(h,TB), TB<TA; Alice resgata o ETH revelando s; Bob lê s e resgata o BTC. Atômico: ou tudo, ou nada. O segredo s é a prova. Nível 2 --- multisig 2-de-3 (entrega física/fiat). Condição não cripto
\prov{relações: parte\_de goldcoin}
\prov{proveniência: paper \textperiodcentered{} fato \textperiodcentered{} confiança 1.0 \textperiodcentered{} domínio paper \textperiodcentered{} história 3 versões no jornal}

%CRISTAL {"arestas":[{"alvo":"computacao_grafica","confianca":1.0,"epistemico":"fato","origem":"paper","rel":"parte_de"}],"confianca":1.0,"contraexemplos":[],"descricao":"NI Estendendo a câmera ao limite para dois centros, dois regimes se distinguem --- e a distinção é, ela mesma, a lição de método. Sistema binário: o horizonte conjunto e a lemniscata.N A reparametrização aplicada à função geral --- a deflexão combinada das duas massas sobre a cena toda, não a cada parte em separado --- produz um horizonte conjunto: as duas sombras unidas pela garganta, cuja curva crítica |ᵢ mᵢ/(z-zᵢ)²|=1 é a lemniscata () --- a mesma costura do chicote (O Corpo Tropical","epistemico":"fato","exemplos":[],"id":"a_cor_e_a_medida_nao_uma_fisica_calculada","memoria":"semantica","meta":{"dominio":"paper","nivel":5,"verificacao":"conceito novo"},"origem":"paper","palavras_chave":[],"sinonimos":[],"tipo":"conceito","titulo":"A cor é a medida, não uma física calculada"}
\section{A cor é a medida, não uma física calculada}
NI Estendendo a câmera ao limite para dois centros, dois regimes se distinguem --- e a distinção é, ela mesma, a lição de método. Sistema binário: o horizonte conjunto e a lemniscata.N A reparametrização aplicada à função geral --- a deflexão combinada das duas massas sobre a cena toda, não a cada parte em separado --- produz um horizonte conjunto: as duas sombras unidas pela garganta, cuja curva crítica |\ensuremath{_{i}} m\ensuremath{_{i}}/(z-z\ensuremath{_{i}})\ensuremath{^{2}}|=1 é a lemniscata () --- a mesma costura do chicote (O Corpo Tropical
\prov{relações: parte\_de computacao\_grafica}
\prov{proveniência: paper \textperiodcentered{} fato \textperiodcentered{} confiança 1.0 \textperiodcentered{} domínio paper \textperiodcentered{} história 3 versões no jornal}

%CRISTAL {"arestas":[{"alvo":"antena","confianca":1.0,"epistemico":"fato","origem":"paper","rel":"parte_de"}],"confianca":1.0,"contraexemplos":[],"descricao":"A máquina computa em F₂: parte do vazio, começa em um bit, evolui pelo gap (o LFSR x²+x+1). Em Eletrônica Analógica já se lê o regime de fase como rotação que conserva (=-1, e⁻iAt) e a parte real como dissipação que irradia. Uma carga acelerada irradia --- isso é Maxwell, não nosso. O que fazemos é instanciar a teoria padrão de antenas nos objetos do painel, sem inventar nada: itemize1pt fonte / corrente I(t): o gap, o LFSR (a torção que gira o elétron); elemento radiante: um t","epistemico":"fato","exemplos":[],"id":"a_corrente_do_gap","memoria":"semantica","meta":{"dominio":"paper","nivel":5,"verificacao":"conceito novo"},"origem":"paper","palavras_chave":[],"sinonimos":[],"tipo":"conceito","titulo":"A corrente do gap"}
\section{A corrente do gap}
A máquina computa em F\ensuremath{_{2}}: parte do vazio, começa em um bit, evolui pelo gap (o LFSR x\ensuremath{^{2}}+x+1). Em Eletrônica Analógica já se lê o regime de fase como rotação que conserva (=-1, e\ensuremath{^{-}}iAt) e a parte real como dissipação que irradia. Uma carga acelerada irradia --- isso é Maxwell, não nosso. O que fazemos é instanciar a teoria padrão de antenas nos objetos do painel, sem inventar nada: itemize1pt fonte / corrente I(t): o gap, o LFSR (a torção que gira o elétron); elemento radiante: um t
\prov{relações: parte\_de antena}
\prov{proveniência: paper \textperiodcentered{} fato \textperiodcentered{} confiança 1.0 \textperiodcentered{} domínio paper \textperiodcentered{} história 3 versões no jornal}

%CRISTAL {"arestas":[{"alvo":"lobo_frontal","confianca":1.0,"epistemico":"fato","origem":"paper","rel":"parte_de"}],"confianca":1.0,"contraexemplos":[],"descricao":"D A fala não é um módulo à parte --- é a face de ação do recall. Por isso é: a área de Broca, no lobo frontal, é onde o colapso do atrator é emitido para fora (cabeca.py:orbitar devolvendo a fala). A ação é o acoplamento em sentido reverso: emite a fase U sem soltar o gerador A (fixo, em casa). O idioma não é um tradutor: é mais um padrão que ressoa no mesmo gato A (mesmo P_). Cruzar línguas é ressoar no mesmo atrator, não atravessar um conversor --- sem peça nova, sujeito ao mesm","epistemico":"fato","exemplos":[],"id":"a_costura_honesta","memoria":"semantica","meta":{"dominio":"paper","nivel":5,"verificacao":"overlap moderado — revisar manualmente"},"origem":"paper","palavras_chave":[],"sinonimos":[],"tipo":"conceito","titulo":"A costura honesta"}
\section{A costura honesta}
D A fala não é um módulo à parte --- é a face de ação do recall. Por isso é: a área de Broca, no lobo frontal, é onde o colapso do atrator é emitido para fora (cabeca.py:orbitar devolvendo a fala). A ação é o acoplamento em sentido reverso: emite a fase U sem soltar o gerador A (fixo, em casa). O idioma não é um tradutor: é mais um padrão que ressoa no mesmo gato A (mesmo P\_). Cruzar línguas é ressoar no mesmo atrator, não atravessar um conversor --- sem peça nova, sujeito ao mesm
\prov{relações: parte\_de lobo\_frontal}
\prov{proveniência: paper \textperiodcentered{} fato \textperiodcentered{} confiança 1.0 \textperiodcentered{} domínio paper \textperiodcentered{} história 3 versões no jornal}

%CRISTAL {"arestas":[{"alvo":"transcricao","confianca":1.0,"epistemico":"fato","origem":"paper","rel":"parte_de"},{"alvo":"colapso","confianca":1.0,"epistemico":"fato","origem":"paper","rel":"parte_de"}],"confianca":1.0,"contraexemplos":[],"descricao":"Nada é postulado de fora. Há um único dado --- a dinâmica --- e o resto não se postula nem se prova à parte: cai no limite tropical. Axioma único. Existe uma dinâmica que mistura sem esquecer no toro T²=²/²: caótica (mistura) e invertível (não esquece). Isso é um automorfismo de Anosov de T² --- “automorfismo” e “hiperbólico” não são hipóteses adicionais, são o conteúdo do axioma: invertível automorfismo (A(2,), =1); mistura gap g>0, hiperbólico~corpodual. prop[Lema","epistemico":"fato","exemplos":[],"id":"a_costura_koch_o_no_dos_fractais","memoria":"semantica","meta":{"dominio":"paper","nivel":5,"verificacao":"conceito novo"},"origem":"paper","palavras_chave":[],"sinonimos":[],"tipo":"conceito","titulo":"A costura: Koch, o nó dos fractais"}
\section{A costura: Koch, o nó dos fractais}
Nada é postulado de fora. Há um único dado --- a dinâmica --- e o resto não se postula nem se prova à parte: cai no limite tropical. Axioma único. Existe uma dinâmica que mistura sem esquecer no toro T\ensuremath{^{2}}=\ensuremath{^{2}}/\ensuremath{^{2}}: caótica (mistura) e invertível (não esquece). Isso é um automorfismo de Anosov de T\ensuremath{^{2}} --- “automorfismo” e “hiperbólico” não são hipóteses adicionais, são o conteúdo do axioma: invertível automorfismo (A(2,), =1); mistura gap g>0, hiperbólico\~{}corpodual. prop[Lema
\prov{relações: parte\_de transcricao; parte\_de colapso}
\prov{proveniência: paper \textperiodcentered{} fato \textperiodcentered{} confiança 1.0 \textperiodcentered{} domínio paper \textperiodcentered{} história 9 versões no jornal}

%CRISTAL {"arestas":[{"alvo":"borboleta","confianca":1.0,"epistemico":"fato","origem":"paper","rel":"parte_de"},{"alvo":"colapso","confianca":1.0,"epistemico":"fato","origem":"paper","rel":"parte_de"}],"confianca":1.0,"contraexemplos":[],"descricao":"D O osso, em uma linha: passa-se a posição (os bytes); contam-se os bits nas posições pares e ímpares; o gato age. defi[O neurônio] Dado um fluxo de bytes (o endereço), seja e= (bits de posição par) e o= (bits de posição ímpar). O neurônio é [ gato(pos)=(e, o),A=(e, o)pmatrix1&1 1&0pmatrix=(e+o, e). ] Não há mais nada: a matriz 22 do gato, a contagem (memória), e o endereço. Inteiro, determinístico, gato()=(0,0). defi No metal são três instruções por byte (neuronio.c, mp","epistemico":"fato","exemplos":[],"id":"a_decimacao_de_cooley--tukey","memoria":"semantica","meta":{"dominio":"paper","nivel":5,"verificacao":"overlap moderado — revisar manualmente"},"origem":"paper","palavras_chave":[],"sinonimos":[],"tipo":"conceito","titulo":"A decimação de Cooley--Tukey"}
\section{A decimação de Cooley--Tukey}
D O osso, em uma linha: passa-se a posição (os bytes); contam-se os bits nas posições pares e ímpares; o gato age. defi[O neurônio] Dado um fluxo de bytes (o endereço), seja e= (bits de posição par) e o= (bits de posição ímpar). O neurônio é [ gato(pos)=(e, o),A=(e, o)pmatrix1\&1 1\&0pmatrix=(e+o, e). ] Não há mais nada: a matriz 22 do gato, a contagem (memória), e o endereço. Inteiro, determinístico, gato()=(0,0). defi No metal são três instruções por byte (neuronio.c, mp
\prov{relações: parte\_de borboleta; parte\_de colapso}
\prov{proveniência: paper \textperiodcentered{} fato \textperiodcentered{} confiança 1.0 \textperiodcentered{} domínio paper \textperiodcentered{} história 9 versões no jornal}

%CRISTAL {"arestas":[{"alvo":"morte","confianca":1.0,"epistemico":"fato","origem":"paper","rel":"parte_de"},{"alvo":"vida","confianca":1.0,"epistemico":"fato","origem":"paper","rel":"parte_de"}],"confianca":1.0,"contraexemplos":[],"descricao":"empty !85A morte deixa problemas; a vida deixa soluções. prata--- o lema do par~morte abstract A vida matemática é o critério positivo simétrico à morte: um objeto está vivo quando sua identidade é dada por uma dinâmica geradora não trivial --- um operador, um fluxo x=F(x), um semigrupo Ut=etL, um mapa xₙ+₁=T(xₙ) ---, da qual seus estados são apenas manifestações. A operação inversa da cristalização é a ressonância: deixar o estado voltar a ser manifestação","epistemico":"fato","exemplos":[],"id":"a_definicao","memoria":"semantica","meta":{"dominio":"paper","nivel":5,"verificacao":"slug idêntico — enriquecer, não duplicar"},"origem":"paper","palavras_chave":[],"sinonimos":[],"tipo":"conceito","titulo":"A definição"}
\section{A definição}
empty !85A morte deixa problemas; a vida deixa soluções. prata--- o lema do par\~{}morte abstract A vida matemática é o critério positivo simétrico à morte: um objeto está vivo quando sua identidade é dada por uma dinâmica geradora não trivial --- um operador, um fluxo x=F(x), um semigrupo Ut=etL, um mapa x\ensuremath{_{n}}+\ensuremath{_{1}}=T(x\ensuremath{_{n}}) ---, da qual seus estados são apenas manifestações. A operação inversa da cristalização é a ressonância: deixar o estado voltar a ser manifestação
\prov{relações: parte\_de morte; parte\_de vida}
\prov{proveniência: paper \textperiodcentered{} fato \textperiodcentered{} confiança 1.0 \textperiodcentered{} domínio paper \textperiodcentered{} história 10 versões no jornal}

%CRISTAL {"arestas":[{"alvo":"cifra_ouro_branco","confianca":1.0,"epistemico":"fato","origem":"paper","rel":"parte_de"},{"alvo":"cifra","confianca":1.0,"epistemico":"fato","origem":"paper","rel":"parte_de"}],"confianca":1.0,"contraexemplos":[],"descricao":"empty abstract O ouro branco é a moeda interna da Engenharia de Ouro --- a liga: a nobreza do ouro com a face prateada da dimensão de prata (₂). Complementa o ouro puro, lastro externo (os swaps cross-chain do GoldCoin; entra de fora). O ouro branco nasce dentro da máquina: minera-se por estrela de prata (hexagrama, 6 modos), valida-se por N-qubits, e cifra-se por uma mudança de base dentro do corpo ₂ₖ --- tudo em bits, sem nada externo (o corpo “vive sozinho dentro do nada”). A","epistemico":"fato","exemplos":[],"id":"a_dimensao_de_prata","memoria":"semantica","meta":{"dominio":"paper","nivel":5,"verificacao":"conceito novo"},"origem":"paper","palavras_chave":[],"sinonimos":[],"tipo":"conceito","titulo":"A dimensão de prata"}
\section{A dimensão de prata}
empty abstract O ouro branco é a moeda interna da Engenharia de Ouro --- a liga: a nobreza do ouro com a face prateada da dimensão de prata (\ensuremath{_{2}}). Complementa o ouro puro, lastro externo (os swaps cross-chain do GoldCoin; entra de fora). O ouro branco nasce dentro da máquina: minera-se por estrela de prata (hexagrama, 6 modos), valida-se por N-qubits, e cifra-se por uma mudança de base dentro do corpo \ensuremath{_{2k}} --- tudo em bits, sem nada externo (o corpo “vive sozinho dentro do nada”). A
\prov{relações: parte\_de cifra\_ouro\_branco; parte\_de cifra}
\prov{proveniência: paper \textperiodcentered{} fato \textperiodcentered{} confiança 1.0 \textperiodcentered{} domínio paper \textperiodcentered{} história 9 versões no jornal}

%CRISTAL {"arestas":[{"alvo":"corpo_tropical","confianca":1.0,"epistemico":"fato","origem":"paper","rel":"parte_de"}],"confianca":1.0,"contraexemplos":[],"descricao":"subsectionO ponto de partida Uma sequência reconstrói a dinâmica que a produziu; e nenhuma dimensão é privilegiada. Não é um lema --- é um fato e uma simetria, e é deles que sai tudo o que segue. O fato é o teorema de reconstrução dinâmica (Takens): observações no tempo de um sistema bastam para recuperar a sua geometria. A simetria é que, entre os sistemas mais simples que admitem isso, nenhum metal é especial. Vamos traduzir os dois em estrutura. Este corpo é um regime; há","epistemico":"fato","exemplos":[],"id":"a_dinamica_que_se_le_da_sequencia","memoria":"semantica","meta":{"dominio":"paper","nivel":5,"verificacao":"slug idêntico — enriquecer, não duplicar"},"origem":"paper","palavras_chave":[],"sinonimos":[],"tipo":"conceito","titulo":"A dinâmica que se lê da sequência"}
\section{A dinâmica que se lê da sequência}
subsectionO ponto de partida Uma sequência reconstrói a dinâmica que a produziu; e nenhuma dimensão é privilegiada. Não é um lema --- é um fato e uma simetria, e é deles que sai tudo o que segue. O fato é o teorema de reconstrução dinâmica (Takens): observações no tempo de um sistema bastam para recuperar a sua geometria. A simetria é que, entre os sistemas mais simples que admitem isso, nenhum metal é especial. Vamos traduzir os dois em estrutura. Este corpo é um regime; há
\prov{relações: parte\_de corpo\_tropical}
\prov{proveniência: paper \textperiodcentered{} fato \textperiodcentered{} confiança 1.0 \textperiodcentered{} domínio paper \textperiodcentered{} história 4 versões no jornal}

%CRISTAL {"arestas":[{"alvo":"goldcoin","confianca":1.0,"epistemico":"fato","origem":"paper","rel":"parte_de"}],"confianca":1.0,"contraexemplos":[],"descricao":"I\nFechar negócio entre quem não se conhece tem quatro custos: risco de contraparte,\nfraude, exposição (verificar revela o que eu não queria) e perda de\nchave. Hoje se paga com intermediário (corretora, escrow, banco) --- caro, lento, e\nainda arrombável (a corretora some com a grana). A pergunta: dá para fechar seguro por\nconstrução, sem confiar em ninguém --- nem em nós?","epistemico":"fato","exemplos":[],"id":"a_divisao_honesta_cada_coisa_no_seu_lugar","memoria":"semantica","meta":{"dominio":"paper","nivel":5,"verificacao":"overlap moderado — revisar manualmente"},"origem":"paper","palavras_chave":[],"sinonimos":[],"tipo":"conceito","titulo":"A divisão honesta (cada coisa no seu lugar)"}
\section{A divisão honesta (cada coisa no seu lugar)}
I
Fechar negócio entre quem não se conhece tem quatro custos: risco de contraparte,
fraude, exposição (verificar revela o que eu não queria) e perda de
chave. Hoje se paga com intermediário (corretora, escrow, banco) --- caro, lento, e
ainda arrombável (a corretora some com a grana). A pergunta: dá para fechar seguro por
construção, sem confiar em ninguém --- nem em nós?
\prov{relações: parte\_de goldcoin}
\prov{proveniência: paper \textperiodcentered{} fato \textperiodcentered{} confiança 1.0 \textperiodcentered{} domínio paper \textperiodcentered{} história 2 versões no jornal}

%CRISTAL {"arestas":[{"alvo":"montagem","confianca":1.0,"epistemico":"fato","origem":"paper","rel":"parte_de"}],"confianca":1.0,"contraexemplos":[],"descricao":")}\npassobox\n./colapso nos.idx \"sua pergunta\"\\# imprime: offset overlap índice\\\\\nmake colapso\npassobox\n O scan quente saiu do Python: mmap no .idx, XOR de 32\\,bytes,\npopcnt, descarta. RAM $=$ query[32] + melhor candidato. Verificado: 4\\,216\\,457 nós\nem ${}$90\\,ms --- ping-pong instantâneo; só resta latência de I/O.","epistemico":"fato","exemplos":[],"id":"a_doutrina_disco_e_ram","memoria":"semantica","meta":{"dominio":"paper","nivel":5,"verificacao":"overlap moderado — revisar manualmente"},"origem":"paper","palavras_chave":[],"sinonimos":[],"tipo":"conceito","titulo":"A doutrina: disco e RAM"}
\section{A doutrina: disco e RAM}
)\}
passobox
./colapso nos.idx "sua pergunta"\textbackslash{}\# imprime: offset overlap índice\textbackslash{}\textbackslash{}
make colapso
passobox
 O scan quente saiu do Python: mmap no .idx, XOR de 32\textbackslash{},bytes,
popcnt, descarta. RAM \$=\$ query[32] + melhor candidato. Verificado: 4\textbackslash{},216\textbackslash{},457 nós
em \$\{\}\$90\textbackslash{},ms --- ping-pong instantâneo; só resta latência de I/O.
\prov{relações: parte\_de montagem}
\prov{proveniência: paper \textperiodcentered{} fato \textperiodcentered{} confiança 1.0 \textperiodcentered{} domínio paper \textperiodcentered{} história 2 versões no jornal}

%CRISTAL {"arestas":[{"alvo":"ponte_simbolica_gentil","confianca":1.0,"epistemico":"fato","origem":"paper","rel":"parte_de"},{"alvo":"analise","confianca":1.0,"epistemico":"fato","origem":"paper","rel":"usa"},{"alvo":"transformada_dourada","confianca":1.0,"epistemico":"fato","origem":"paper","rel":"referencia"},{"alvo":"paper_2_reta_ouro","confianca":1.0,"epistemico":"fato","origem":"paper","rel":"referencia"},{"alvo":"corpo_tropical","confianca":1.0,"epistemico":"fato","origem":"paper","rel":"referencia"}],"confianca":1.0,"contraexemplos":[],"descricao":"Buraco em e sequência para 2 A sequência (rₙ) com r₁=1, rₙ+₁=4rₙ2+rₙ², é monótona crescente, limitada e de Cauchy em, mas não converge em --- o “buraco” para 2 (pp.~398--401). A classe [(rₙ)] C é o real 2 no chart simbólico. observacao e e progressão geométrica A sequência rₙ=1+ₖ=₁ⁿ1k! usa explicitamente a soma de progressão geométrica (PGS) na majoração; é Cauchy em sem limite racional. A cla","epistemico":"fato","exemplos":[],"id":"a_engenharia_de_ouro_pa_pgs_agm_fc_reta","memoria":"semantica","meta":{"dominio":"paper","nivel":5,"verificacao":"slug idêntico — enriquecer, não duplicar"},"origem":"paper","palavras_chave":[],"sinonimos":[],"tipo":"conceito","titulo":"A engenharia de ouro: PA, PGS, AGM, FC, reta"}
\section{A engenharia de ouro: PA, PGS, AGM, FC, reta}
Buraco em e sequência para 2 A sequência (r\ensuremath{_{n}}) com r\ensuremath{_{1}}=1, r\ensuremath{_{n}}+\ensuremath{_{1}}=4r\ensuremath{_{n}}2+r\ensuremath{_{n}}\ensuremath{^{2}}, é monótona crescente, limitada e de Cauchy em, mas não converge em --- o “buraco” para 2 (pp.\~{}398--401). A classe [(r\ensuremath{_{n}})] C é o real 2 no chart simbólico. observacao e e progressão geométrica A sequência r\ensuremath{_{n}}=1+\ensuremath{_{k}}=\ensuremath{_{1}}\ensuremath{^{n}}1k! usa explicitamente a soma de progressão geométrica (PGS) na majoração; é Cauchy em sem limite racional. A cla
\prov{relações: parte\_de ponte\_simbolica\_gentil; usa analise; referencia transformada\_dourada; referencia paper\_2\_reta\_ouro; referencia corpo\_tropical}
\prov{proveniência: paper \textperiodcentered{} fato \textperiodcentered{} confiança 1.0 \textperiodcentered{} domínio paper \textperiodcentered{} história 164 versões no jornal}

%CRISTAL {"arestas":[{"alvo":"mecanica","confianca":1.0,"epistemico":"fato","origem":"paper","rel":"parte_de"}],"confianca":1.0,"contraexemplos":[],"descricao":"A pressão algébrica =1-x² muda de sinal em |x|=1: para |x|<1, >0 --- regime dissipativo, a termodinâmica, a norma contrai; em |x|=1, =0 --- o ponto hermitiano, onde a norma é exatamente conservada. É ali (s=1, a=0) que a evolução é reversível e a mecânica quântica vive. O eixo dissipativohermitiano é a parábola: a termodinâmica habita o interior (|x|<1), a quântica a fronteira (|x|=1). É exatamente nessa fronteira que a evolução é um grupo unitário --- e é dela que a equaç","epistemico":"fato","exemplos":[],"id":"a_equacao_de_schrodinger_cai_do_gato","memoria":"semantica","meta":{"dominio":"paper","nivel":5,"verificacao":"slug idêntico — enriquecer, não duplicar"},"origem":"paper","palavras_chave":[],"sinonimos":[],"tipo":"conceito","titulo":"A equação de Schrödinger cai do gato"}
\section{A equação de Schrödinger cai do gato}
A pressão algébrica =1-x\ensuremath{^{2}} muda de sinal em |x|=1: para |x|<1, >0 --- regime dissipativo, a termodinâmica, a norma contrai; em |x|=1, =0 --- o ponto hermitiano, onde a norma é exatamente conservada. É ali (s=1, a=0) que a evolução é reversível e a mecânica quântica vive. O eixo dissipativohermitiano é a parábola: a termodinâmica habita o interior (|x|<1), a quântica a fronteira (|x|=1). É exatamente nessa fronteira que a evolução é um grupo unitário --- e é dela que a equaç
\prov{relações: parte\_de mecanica}
\prov{proveniência: paper \textperiodcentered{} fato \textperiodcentered{} confiança 1.0 \textperiodcentered{} domínio paper \textperiodcentered{} história 10 versões no jornal}

%CRISTAL {"arestas":[{"alvo":"paper_45_confluencia_ouro","confianca":1.0,"epistemico":"fato","origem":"paper","rel":"parte_de"}],"confianca":1.0,"contraexemplos":[],"descricao":"cada constante, um papel tabularcll constante & ofcio & onde (na série), & a curva modular --- a casa do, onde a árvore vive & XXXVI, XXXVIII 2 & o dobramento 2 --- árvore de sinais; quadrático & XXXIX & a taxa --- entropia do golden mean shift & XL & a escala logartmica --- a régua que mede a taxa & (aqui) tabular O 2 é duplo, e não por acaso: o mesmo “2” que dobra os ângulos na árvore","epistemico":"fato","exemplos":[],"id":"a_escala_logaritmica_do_ouro","memoria":"semantica","meta":{"dominio":"paper","nivel":5,"verificacao":"conceito novo"},"origem":"paper","palavras_chave":[],"sinonimos":[],"tipo":"conceito","titulo":"A escala logarítmica do ouro"}
\section{A escala logarítmica do ouro}
cada constante, um papel tabularcll constante \& ofcio \& onde (na série), \& a curva modular --- a casa do, onde a árvore vive \& XXXVI, XXXVIII 2 \& o dobramento 2 --- árvore de sinais; quadrático \& XXXIX \& a taxa --- entropia do golden mean shift \& XL \& a escala logartmica --- a régua que mede a taxa \& (aqui) tabular O 2 é duplo, e não por acaso: o mesmo “2” que dobra os ângulos na árvore
\prov{relações: parte\_de paper\_45\_confluencia\_ouro}
\prov{proveniência: paper \textperiodcentered{} fato \textperiodcentered{} confiança 1.0 \textperiodcentered{} domínio paper \textperiodcentered{} história 3 versões no jornal}

%CRISTAL {"arestas":[{"alvo":"leitura_das_unidades_o_sistema_por-unidade_e_o_si","confianca":1.0,"epistemico":"fato","origem":"paper","rel":"parte_de"},{"alvo":"testes","confianca":0.5,"epistemico":"fato","origem":"wiki","rel":"relaciona"},{"alvo":"topologia","confianca":0.5,"epistemico":"fato","origem":"wiki","rel":"relaciona"},{"alvo":"recursao","confianca":0.5,"epistemico":"fato","origem":"wiki","rel":"relaciona"},{"alvo":"o_metodo_em_escala_conhecida_o_sistema_solar","confianca":0.5,"epistemico":"fato","origem":"wiki","rel":"relaciona"},{"alvo":"fase_2_tecido_de_conhecimento","confianca":0.5,"epistemico":"fato","origem":"wiki","rel":"relaciona"},{"alvo":"grupos","confianca":0.5,"epistemico":"fato","origem":"wiki","rel":"relaciona"},{"alvo":"nucleo_broca_vivo","confianca":0.5,"epistemico":"fato","origem":"wiki","rel":"relaciona"},{"alvo":"atomo","confianca":0.5,"epistemico":"fato","origem":"wiki","rel":"relaciona"},{"alvo":"respiracao_dos_programas","confianca":0.5,"epistemico":"fato","origem":"wiki","rel":"relaciona"}],"confianca":1.0,"contraexemplos":[],"descricao":"","epistemico":"fato","exemplos":[],"id":"a_escala_uma_volta_no_c_rculo","memoria":"semantica","meta":{"dominio":"paper","nivel":5,"verificacao":"conceito novo"},"origem":"paper","palavras_chave":[],"sinonimos":[],"tipo":"conceito","titulo":"A escala: uma volta no crculo"}
\section{A escala: uma volta no crculo}
\prov{relações: parte\_de leitura\_das\_unidades\_o\_sistema\_por-unidade\_e\_o\_si; relaciona testes; relaciona topologia; relaciona recursao; relaciona o\_metodo\_em\_escala\_conhecida\_o\_sistema\_solar; relaciona fase\_2\_tecido\_de\_conhecimento; relaciona grupos; relaciona nucleo\_broca\_vivo; relaciona atomo; relaciona respiracao\_dos\_programas}
\prov{proveniência: paper \textperiodcentered{} fato \textperiodcentered{} confiança 1.0 \textperiodcentered{} domínio paper \textperiodcentered{} história 13 versões no jornal}

%CRISTAL {"arestas":[{"alvo":"ciencias","confianca":0.047,"epistemico":"fato","origem":"manual","rel":"referencia"},{"alvo":"geometria","confianca":0.023,"epistemico":"fato","origem":"manual","rel":"referencia"},{"alvo":"quatro_ciencias","confianca":1.0,"epistemico":"fato","origem":"paper","rel":"parte_de"}],"confianca":1.0,"contraexemplos":[],"descricao":"empty !85Não há quatro ciências --- há um gato, lido em quatro quadrantes. E não se escolhe onde o dado mora: manda-se o fluxo cru, e a álgebra o assenta. abstract Álgebra, geometria, mecânica e termodinâmica não são quatro objetos --- são quatro leituras de uma só estrutura: a função de onda da Máquina M=U P, gerada pelo gato A=(smallmatrix1&1 1&0 smallmatrix) (autovalores ouro, prata), com a balança a+c²=1. Dois cortes, já presentes na estrutura, produzem as","epistemico":"fato","exemplos":[],"id":"a_estrutura_e_uma_so","memoria":"semantica","meta":{"dominio":"paper","nivel":5,"verificacao":"slug idêntico — enriquecer, não duplicar"},"origem":"paper","palavras_chave":[],"sinonimos":[],"tipo":"conceito","titulo":"A estrutura é uma só"}
\section{A estrutura é uma só}
empty !85Não há quatro ciências --- há um gato, lido em quatro quadrantes. E não se escolhe onde o dado mora: manda-se o fluxo cru, e a álgebra o assenta. abstract Álgebra, geometria, mecânica e termodinâmica não são quatro objetos --- são quatro leituras de uma só estrutura: a função de onda da Máquina M=U P, gerada pelo gato A=(smallmatrix1\&1 1\&0 smallmatrix) (autovalores ouro, prata), com a balança a+c\ensuremath{^{2}}=1. Dois cortes, já presentes na estrutura, produzem as
\prov{relações: referencia ciencias; referencia geometria; parte\_de quatro\_ciencias}
\prov{proveniência: paper \textperiodcentered{} fato \textperiodcentered{} confiança 1.0 \textperiodcentered{} domínio paper \textperiodcentered{} história 17 versões no jornal}

%CRISTAL {"arestas":[{"alvo":"leitura_aritmetica_a_regua_de_gentil_o_agm","confianca":1.0,"epistemico":"fato","origem":"paper","rel":"parte_de"}],"confianca":1.0,"contraexemplos":[],"descricao":"As progressões de Lopes dão duas réguas. A aditiva (progressões aritméticas de ordem k nas fracões contnuas) mede os números de Hurwitz ---,2,e são compressveis; e escapam. A multiplicativa é a P.A., P.G. entrelacadas: a própria AGM, de gênero 1, cujo passo é a transformacão de Landen (uma 2-isogenia). No espaco das razões s=(a-b)/(a+b) a métrica invariante é m(s) 1/!(s(1/s)), e o AGM (s=0) é um super-atrator: uma cúspide","epistemico":"fato","exemplos":[],"id":"a_fam_lia_do_agm","memoria":"semantica","meta":{"dominio":"paper","nivel":5,"verificacao":"overlap moderado — revisar manualmente"},"origem":"paper","palavras_chave":[],"sinonimos":[],"tipo":"conceito","titulo":"A famlia do AGM"}
\section{A famlia do AGM}
As progressões de Lopes dão duas réguas. A aditiva (progressões aritméticas de ordem k nas fracões contnuas) mede os números de Hurwitz ---,2,e são compressveis; e escapam. A multiplicativa é a P.A., P.G. entrelacadas: a própria AGM, de gênero 1, cujo passo é a transformacão de Landen (uma 2-isogenia). No espaco das razões s=(a-b)/(a+b) a métrica invariante é m(s) 1/!(s(1/s)), e o AGM (s=0) é um super-atrator: uma cúspide
\prov{relações: parte\_de leitura\_aritmetica\_a\_regua\_de\_gentil\_o\_agm}
\prov{proveniência: paper \textperiodcentered{} fato \textperiodcentered{} confiança 1.0 \textperiodcentered{} domínio paper \textperiodcentered{} história 5 versões no jornal}

%CRISTAL {"arestas":[{"alvo":"contabilidade","confianca":1.0,"epistemico":"fato","origem":"paper","rel":"parte_de"}],"confianca":1.0,"contraexemplos":[],"descricao":"D O balanço não é uma soma de colunas; é o colapso na leitura. Pela mecânica (mecanica.py, mecanica.tex), a forma é resolvida uma vez (simbólica, contínua) e discretizada fora, no ponto r --- a régua do observador. Lê-se o tecido projetando-o nas duas retas (ouro / prata, autovalores do gato A=(smallmatrix1&1 1&0smallmatrix)): [ a=RR_+R_, c²=RR_+R_, a+c²=1,(norma/Robertson). ] I Essa é a conservação do lastro: a norma da função de onda lida é 1. Em leitura contábil, é","epistemico":"fato","exemplos":[],"id":"a_folha_nao_e_o_lastro","memoria":"semantica","meta":{"dominio":"paper","nivel":5,"verificacao":"conceito novo"},"origem":"paper","palavras_chave":[],"sinonimos":[],"tipo":"conceito","titulo":"A folha não é o lastro"}
\section{A folha não é o lastro}
D O balanço não é uma soma de colunas; é o colapso na leitura. Pela mecânica (mecanica.py, mecanica.tex), a forma é resolvida uma vez (simbólica, contínua) e discretizada fora, no ponto r --- a régua do observador. Lê-se o tecido projetando-o nas duas retas (ouro / prata, autovalores do gato A=(smallmatrix1\&1 1\&0smallmatrix)): [ a=RR\_+R\_, c\ensuremath{^{2}}=RR\_+R\_, a+c\ensuremath{^{2}}=1,(norma/Robertson). ] I Essa é a conservação do lastro: a norma da função de onda lida é 1. Em leitura contábil, é
\prov{relações: parte\_de contabilidade}
\prov{proveniência: paper \textperiodcentered{} fato \textperiodcentered{} confiança 1.0 \textperiodcentered{} domínio paper \textperiodcentered{} história 3 versões no jornal}

%CRISTAL {"arestas":[{"alvo":"transcricao","confianca":1.0,"epistemico":"fato","origem":"paper","rel":"parte_de"},{"alvo":"colapso","confianca":1.0,"epistemico":"fato","origem":"paper","rel":"parte_de"}],"confianca":1.0,"contraexemplos":[],"descricao":"Atualização do protocolo (GEX44): além de cid | seq | bump, o envio conservado pode incluir compacto (resíduo mod p=1009, atrator) no mesmo datagrama --- ainda cru, ainda na banda, sem AEAD por cima. Quem sintoniza decodifica bump e lê o floco do canal; quem não sintoniza vê ruído nos dois. A reciprocidade mantém-se: ida e volta na mesma banda, mão dupla, com Dᵢj e Djᵢ. A antena associativa roteia sem contaminar; o floco evita que irradie sem memória. Conservar","epistemico":"fato","exemplos":[],"id":"a_fronteira","memoria":"semantica","meta":{"dominio":"paper","duplicata_sugerida":[],"nivel":5,"verificacao":"parte de conceito mais amplo 'a_fronteira'"},"origem":"paper","palavras_chave":[],"sinonimos":[],"tipo":"conceito","titulo":"A fronteira"}
\section{A fronteira}
Atualização do protocolo (GEX44): além de cid | seq | bump, o envio conservado pode incluir compacto (resíduo mod p=1009, atrator) no mesmo datagrama --- ainda cru, ainda na banda, sem AEAD por cima. Quem sintoniza decodifica bump e lê o floco do canal; quem não sintoniza vê ruído nos dois. A reciprocidade mantém-se: ida e volta na mesma banda, mão dupla, com D\ensuremath{_{i}}j e Dj\ensuremath{_{i}}. A antena associativa roteia sem contaminar; o floco evita que irradie sem memória. Conservar
\prov{relações: parte\_de transcricao; parte\_de colapso}
\prov{proveniência: paper \textperiodcentered{} fato \textperiodcentered{} confiança 1.0 \textperiodcentered{} domínio paper \textperiodcentered{} história 10 versões no jornal}

%CRISTAL {"arestas":[{"alvo":"goldcoin","confianca":1.0,"epistemico":"fato","origem":"paper","rel":"parte_de"},{"alvo":"goldchain","confianca":1.0,"epistemico":"fato","origem":"paper","rel":"parte_de"}],"confianca":1.0,"contraexemplos":[],"descricao":"CN O lastro não fica trancado --- ele transmite. A objeção econômica (“obra que só o dono vê não tem liquidez”, crivo do conselho) se dissolve quando a obra é um evento ao vivo (um jogo, um show) vendido em tempo real a muitos: liquidez máxima. E a infraestrutura já existe e está medida: itemize1pt Distribuição --- o enxame de clones (a antena, Redes): um CDN/P2P nativo, cada clone retransmite a c. Multiplexação --- os canais mux/demux são um bump compartilhado: vários streams num canal.","epistemico":"fato","exemplos":[],"id":"a_identidade_da_obra_a_alma_sob_transformacao_medido","memoria":"semantica","meta":{"dominio":"paper","nivel":5,"verificacao":"conceito novo"},"origem":"paper","palavras_chave":[],"sinonimos":[],"tipo":"conceito","titulo":"A identidade da obra: a alma sob transformação (medido)"}
\section{A identidade da obra: a alma sob transformação (medido)}
CN O lastro não fica trancado --- ele transmite. A objeção econômica (“obra que só o dono vê não tem liquidez”, crivo do conselho) se dissolve quando a obra é um evento ao vivo (um jogo, um show) vendido em tempo real a muitos: liquidez máxima. E a infraestrutura já existe e está medida: itemize1pt Distribuição --- o enxame de clones (a antena, Redes): um CDN/P2P nativo, cada clone retransmite a c. Multiplexação --- os canais mux/demux são um bump compartilhado: vários streams num canal.
\prov{relações: parte\_de goldcoin; parte\_de goldchain}
\prov{proveniência: paper \textperiodcentered{} fato \textperiodcentered{} confiança 1.0 \textperiodcentered{} domínio paper \textperiodcentered{} história 9 versões no jornal}

%CRISTAL {"arestas":[{"alvo":"antena","confianca":1.0,"epistemico":"fato","origem":"paper","rel":"parte_de"}],"confianca":1.0,"contraexemplos":[],"descricao":"sectionA sabatina Aplicação (GPT-5.5). Instanciou a teoria clássica nos objetos do painel: a corrente do LFSR, a ressonância Z ᵢₙ=0 com =k²-k cut₂, o casamento conjugado e, a diretividade via Eᵢj, e o factor da rede de um bit. Listou, ao fim, os parâmetros físicos por fixar. Crivo de engenharia (Grok). Conferiu contra a literatura (Balanis, Orfanidis) e cravou o ponto central: nós no mesmo seed têm a mesma fase, logo AF=N² só em broadside; o feixe dirigível","epistemico":"fato","exemplos":[],"id":"a_implantacao_no_gex44_arquitetura_concreta_---_mapa_verificado_20062026","memoria":"semantica","meta":{"dominio":"paper","nivel":5,"verificacao":"conceito novo"},"origem":"paper","palavras_chave":[],"sinonimos":[],"tipo":"conceito","titulo":"A implantação no GEX44 (arquitetura concreta --- mapa verificado 20/06/2026)"}
\section{A implantação no GEX44 (arquitetura concreta --- mapa verificado 20/06/2026)}
sectionA sabatina Aplicação (GPT-5.5). Instanciou a teoria clássica nos objetos do painel: a corrente do LFSR, a ressonância Z \ensuremath{_{in}}=0 com =k\ensuremath{^{2}}-k cut\ensuremath{_{2}}, o casamento conjugado e, a diretividade via E\ensuremath{_{i}}j, e o factor da rede de um bit. Listou, ao fim, os parâmetros físicos por fixar. Crivo de engenharia (Grok). Conferiu contra a literatura (Balanis, Orfanidis) e cravou o ponto central: nós no mesmo seed têm a mesma fase, logo AF=N\ensuremath{^{2}} só em broadside; o feixe dirigível
\prov{relações: parte\_de antena}
\prov{proveniência: paper \textperiodcentered{} fato \textperiodcentered{} confiança 1.0 \textperiodcentered{} domínio paper \textperiodcentered{} história 3 versões no jornal}

%CRISTAL {"arestas":[{"alvo":"o_grafo_de_delegacao_e_a_atenuacao","confianca":1.0,"epistemico":"fato","origem":"paper","rel":"parte_de"}],"confianca":1.0,"contraexemplos":[],"descricao":"definition[Autoridades intrínsecas] Duas fontes de autoridade não derivam de ninguém: [nosep] o operador detém (manage,all, 0,) --- vê e faz tudo, com re-delegação irrestrita; o dono do tenant t detém (manage,all,(=t),0,) --- controle total dos próprios recursos. definition definition[Grafo de delegação] Um grafo de delegação é um grafo dirigido finito (não necessariamente acíclico) cujos vértices são principals e cujas a","epistemico":"fato","exemplos":[],"id":"a_intensidade_dep_profundidade_de_propagacao","memoria":"semantica","meta":{"dominio":"paper","nivel":5},"origem":"paper","palavras_chave":[],"sinonimos":[],"tipo":"conceito","titulo":"A intensidade dep: profundidade de propagação"}
\section{A intensidade dep: profundidade de propagação}
definition[Autoridades intrínsecas] Duas fontes de autoridade não derivam de ninguém: [nosep] o operador detém (manage,all, 0,) --- vê e faz tudo, com re-delegação irrestrita; o dono do tenant t detém (manage,all,(=t),0,) --- controle total dos próprios recursos. definition definition[Grafo de delegação] Um grafo de delegação é um grafo dirigido finito (não necessariamente acíclico) cujos vértices são principals e cujas a
\prov{relações: parte\_de o\_grafo\_de\_delegacao\_e\_a\_atenuacao}
\prov{proveniência: paper \textperiodcentered{} fato \textperiodcentered{} confiança 1.0 \textperiodcentered{} domínio paper \textperiodcentered{} história 4 versões no jornal}

%CRISTAL {"arestas":[{"alvo":"termodinamica","confianca":1.0,"epistemico":"fato","origem":"paper","rel":"parte_de"}],"confianca":1.0,"contraexemplos":[],"descricao":"A temperatura interna da máquina não é imposta de fora: lê-se do gap espectral. O fator de contração || define uma escala de relaxação; o gap hiperbólico >0 fixa a taxa com que o resíduo decai. Quanto maior o gap, mais “fria” e ordenada a dinâmica; quando o gap fecha (=0, o nilpotente), não há relaxação --- é o zero absoluto algébrico do regime glacial.","epistemico":"fato","exemplos":[],"id":"a_irreversibilidade_o_regime_glacial","memoria":"semantica","meta":{"dominio":"paper","nivel":5,"verificacao":"slug idêntico — enriquecer, não duplicar"},"origem":"paper","palavras_chave":[],"sinonimos":[],"tipo":"conceito","titulo":"A irreversibilidade: o regime glacial"}
\section{A irreversibilidade: o regime glacial}
A temperatura interna da máquina não é imposta de fora: lê-se do gap espectral. O fator de contração || define uma escala de relaxação; o gap hiperbólico >0 fixa a taxa com que o resíduo decai. Quanto maior o gap, mais “fria” e ordenada a dinâmica; quando o gap fecha (=0, o nilpotente), não há relaxação --- é o zero absoluto algébrico do regime glacial.
\prov{relações: parte\_de termodinamica}
\prov{proveniência: paper \textperiodcentered{} fato \textperiodcentered{} confiança 1.0 \textperiodcentered{} domínio paper \textperiodcentered{} história 11 versões no jornal}

%CRISTAL {"arestas":[{"alvo":"corpo_tropical","confianca":1.0,"epistemico":"fato","origem":"paper","rel":"parte_de"}],"confianca":1.0,"contraexemplos":[],"descricao":"Aqui a reta volta sobre a própria sequência e fecha a cadeia. A codificação da Seção~1 --- ler os bits na base, futuro contra passado --- não é só uma fórmula: aplicada à sequência que toca, ela monta o atrator, o corpo geométrico da dinâmica. Cada janela de bits vira um ponto X=(x,y)=( bt+ₖ⁻(k⁺¹), bt-ₖ⁻k), e a nuvem de todos os pontos é o atrator próprio daquela sequência. a geometria do ouro é fractal O conjunto reconstruído pela codificação na base é auto-similar","epistemico":"fato","exemplos":[],"id":"a_maquina_reconstroi_verificacao_em_bits_brutos","memoria":"semantica","meta":{"dominio":"paper","nivel":5,"verificacao":"slug idêntico — enriquecer, não duplicar"},"origem":"paper","palavras_chave":[],"sinonimos":[],"tipo":"conceito","titulo":"A máquina reconstrói: verificação em bits brutos"}
\section{A máquina reconstrói: verificação em bits brutos}
Aqui a reta volta sobre a própria sequência e fecha a cadeia. A codificação da Seção\~{}1 --- ler os bits na base, futuro contra passado --- não é só uma fórmula: aplicada à sequência que toca, ela monta o atrator, o corpo geométrico da dinâmica. Cada janela de bits vira um ponto X=(x,y)=( bt+\ensuremath{_{k}}\ensuremath{^{-}}(k\ensuremath{^{+1}}), bt-\ensuremath{_{k}}\ensuremath{^{-}}k), e a nuvem de todos os pontos é o atrator próprio daquela sequência. a geometria do ouro é fractal O conjunto reconstruído pela codificação na base é auto-similar
\prov{relações: parte\_de corpo\_tropical}
\prov{proveniência: paper \textperiodcentered{} fato \textperiodcentered{} confiança 1.0 \textperiodcentered{} domínio paper \textperiodcentered{} história 4 versões no jornal}

%CRISTAL {"arestas":[{"alvo":"leitura_dinamica_a_maquina_e_a_materia_ligada","confianca":1.0,"epistemico":"fato","origem":"paper","rel":"parte_de"}],"confianca":1.0,"contraexemplos":[],"descricao":"A Parte~VI faz de cada álgebra um estado --- um ponto s no espaco s ---, dotado de aritmética (), topologia e dinâmica. O fluxo natural é s = 1-s², com s=0 (comutativo) uma fonte instável e s=1 (o quatérnio) um atrator --- a juncão de imposto nulo. Quantizado, um estado é uma superposicão de álgebras |=(),|,d, com potencial de Mathieu (pocos nos quatérnios); e medir s é colapsar numa letra, .","epistemico":"fato","exemplos":[],"id":"a_materia_ligada_parte_vii_ligar_e_colapsar","memoria":"semantica","meta":{"dominio":"paper","nivel":5,"verificacao":"conceito novo"},"origem":"paper","palavras_chave":[],"sinonimos":[],"tipo":"conceito","titulo":"A matéria ligada (Parte VII): ligar é colapsar"}
\section{A matéria ligada (Parte VII): ligar é colapsar}
A Parte\~{}VI faz de cada álgebra um estado --- um ponto s no espaco s ---, dotado de aritmética (), topologia e dinâmica. O fluxo natural é s = 1-s\ensuremath{^{2}}, com s=0 (comutativo) uma fonte instável e s=1 (o quatérnio) um atrator --- a juncão de imposto nulo. Quantizado, um estado é uma superposicão de álgebras |=(),|,d, com potencial de Mathieu (pocos nos quatérnios); e medir s é colapsar numa letra, .
\prov{relações: parte\_de leitura\_dinamica\_a\_maquina\_e\_a\_materia\_ligada}
\prov{proveniência: paper \textperiodcentered{} fato \textperiodcentered{} confiança 1.0 \textperiodcentered{} domínio paper \textperiodcentered{} história 4 versões no jornal}

%CRISTAL {"arestas":[{"alvo":"geometria","confianca":1.0,"epistemico":"fato","origem":"wiki","rel":"define"},{"alvo":"analise","confianca":1.0,"epistemico":"fato","origem":"wiki","rel":"depende"}],"confianca":1.0,"contraexemplos":[],"descricao":"O espaço interno dos estados da dinâmica é a esfera dos quaternions unitários, [ S³=q H: |q|=1, ] com a sua métrica round d²S₃. Não é uma escolha de exibição: é onde a torção vive --- os versores que ela gira. A projeção natural, em bit e sem divisão, é a fibração de Hopf [ S³ S², ] a submersão riemanniana canônica (fibras S¹). É a geometria nativa da máquina; a perspectiva 1/w é outra realização métrica (a do observador), não um defeito.","epistemico":"fato","exemplos":[],"id":"a_metrica_de_sitter_global","memoria":"semantica","meta":{"dominio":"paper","nivel":5,"verificacao":"slug idêntico — enriquecer, não duplicar"},"origem":"paper","palavras_chave":[],"sinonimos":[],"tipo":"conceito","titulo":"A métrica: de Sitter global"}
\section{A métrica: de Sitter global}
O espaço interno dos estados da dinâmica é a esfera dos quaternions unitários, [ S\ensuremath{^{3}}=q H: |q|=1, ] com a sua métrica round d\ensuremath{^{2}}S\ensuremath{_{3}}. Não é uma escolha de exibição: é onde a torção vive --- os versores que ela gira. A projeção natural, em bit e sem divisão, é a fibração de Hopf [ S\ensuremath{^{3}} S\ensuremath{^{2}}, ] a submersão riemanniana canônica (fibras S\ensuremath{^{1}}). É a geometria nativa da máquina; a perspectiva 1/w é outra realização métrica (a do observador), não um defeito.
\prov{relações: define geometria; depende analise}
\prov{proveniência: paper \textperiodcentered{} fato \textperiodcentered{} confiança 1.0 \textperiodcentered{} domínio paper \textperiodcentered{} história 108 versões no jornal}

%CRISTAL {"arestas":[{"alvo":"mecanica_quantica","confianca":0.95,"epistemico":"fato","origem":"paper","rel":"parte_de"}],"confianca":1.0,"contraexemplos":[],"descricao":"Apresentamos a álgebra do zero, sem pressupor resultados externos. (para o leitor não-algebrista). A álgebra hipercomplexa é ⁿ munido de uma regra de multiplicação que generaliza a dos números complexos: cada elemento tem uma parte real (um número) e uma parte imaginária (um vetor). A forma mais geral compatível com essa estrutura pesa cada um dos cinco termos do produto com um botão próprio --- e os cinco botões formam, eles mesmos, um espaço de álgebras ⁵. As estruturas céleb","epistemico":"fato","exemplos":[],"id":"a_metrica_quantica","memoria":"semantica","meta":{"dominio":"paper","duplicata_sugerida":[],"nivel":5,"verificacao":"parte de conceito mais amplo 'a_metrica_quantica'"},"origem":"paper","palavras_chave":[],"sinonimos":[],"tipo":"conceito","titulo":"A métrica quântica"}
\section{A métrica quântica}
Apresentamos a álgebra do zero, sem pressupor resultados externos. (para o leitor não-algebrista). A álgebra hipercomplexa é \ensuremath{^{n}} munido de uma regra de multiplicação que generaliza a dos números complexos: cada elemento tem uma parte real (um número) e uma parte imaginária (um vetor). A forma mais geral compatível com essa estrutura pesa cada um dos cinco termos do produto com um botão próprio --- e os cinco botões formam, eles mesmos, um espaço de álgebras \ensuremath{^{5}}. As estruturas céleb
\prov{relações: parte\_de mecanica\_quantica}
\prov{proveniência: paper \textperiodcentered{} fato \textperiodcentered{} confiança 1.0 \textperiodcentered{} domínio paper \textperiodcentered{} história 5 versões no jornal}

%CRISTAL {"arestas":[{"alvo":"goldcoin","confianca":1.0,"epistemico":"fato","origem":"paper","rel":"parte_de"},{"alvo":"goldchain","confianca":1.0,"epistemico":"fato","origem":"paper","rel":"parte_de"}],"confianca":1.0,"contraexemplos":[],"descricao":"CN Um negócio pode liquidar de dois modos --- e ambos saem dos mesmos primitivos já construídos. No tempo (o atomic swap). Liquidação atômica on-chain (Protocolo de Atomic Swap, à parte): rápida e trustless, mas a atomicidade é condicional à liveness --- sob censura/congestionamento, a parte honesta pode ser lesada (crivo do conselho, caminho bobcensurado). É o imposto de liquidar no tempo da chain. Assíncrona (resolução baseada em evidência). I Correção (crivo do conselho, 2 votos): isto","epistemico":"fato","exemplos":[],"id":"a_obra_circula_broadcast_seguro_pelos_clones","memoria":"semantica","meta":{"dominio":"paper","nivel":5,"verificacao":"conceito novo"},"origem":"paper","palavras_chave":[],"sinonimos":[],"tipo":"conceito","titulo":"A obra circula: broadcast seguro pelos clones"}
\section{A obra circula: broadcast seguro pelos clones}
CN Um negócio pode liquidar de dois modos --- e ambos saem dos mesmos primitivos já construídos. No tempo (o atomic swap). Liquidação atômica on-chain (Protocolo de Atomic Swap, à parte): rápida e trustless, mas a atomicidade é condicional à liveness --- sob censura/congestionamento, a parte honesta pode ser lesada (crivo do conselho, caminho bobcensurado). É o imposto de liquidar no tempo da chain. Assíncrona (resolução baseada em evidência). I Correção (crivo do conselho, 2 votos): isto
\prov{relações: parte\_de goldcoin; parte\_de goldchain}
\prov{proveniência: paper \textperiodcentered{} fato \textperiodcentered{} confiança 1.0 \textperiodcentered{} domínio paper \textperiodcentered{} história 9 versões no jornal}

%CRISTAL {"arestas":[{"alvo":"termodinamica","confianca":1.0,"epistemico":"fato","origem":"paper","rel":"parte_de"}],"confianca":1.0,"contraexemplos":[],"descricao":"O operador de evolução que a dinâmica bruta do bit entrega --- medido como operador de avanço da trajetória, sem impor forma --- é o de Koopman--Perron--Frobenius, dissipativo: ||[0,65, 0,86], não-normal, ressonâncias de Ruelle--Pollicott. Verificado. A máquina não é unitária; ela contrai, e a contração tem direção. É por isso que a leitura natural do bit é a termodinâmica (irreversível, entropia), não a mecânica (reversível, unitária). O regime unitário será o limite, não o ponto de part","epistemico":"fato","exemplos":[],"id":"a_primeira_lei_energia_total_conservada","memoria":"semantica","meta":{"dominio":"paper","nivel":5,"verificacao":"slug idêntico — enriquecer, não duplicar"},"origem":"paper","palavras_chave":[],"sinonimos":[],"tipo":"conceito","titulo":"A primeira lei: energia total conservada"}
\section{A primeira lei: energia total conservada}
O operador de evolução que a dinâmica bruta do bit entrega --- medido como operador de avanço da trajetória, sem impor forma --- é o de Koopman--Perron--Frobenius, dissipativo: ||[0,65, 0,86], não-normal, ressonâncias de Ruelle--Pollicott. Verificado. A máquina não é unitária; ela contrai, e a contração tem direção. É por isso que a leitura natural do bit é a termodinâmica (irreversível, entropia), não a mecânica (reversível, unitária). O regime unitário será o limite, não o ponto de part
\prov{relações: parte\_de termodinamica}
\prov{proveniência: paper \textperiodcentered{} fato \textperiodcentered{} confiança 1.0 \textperiodcentered{} domínio paper \textperiodcentered{} história 11 versões no jornal}

%CRISTAL {"arestas":[{"alvo":"vida","confianca":1.0,"epistemico":"fato","origem":"paper","rel":"parte_de"}],"confianca":1.0,"contraexemplos":[],"descricao":"D Aplicamos a Definição~2 à Hipótese de Riemann, lida não pela igualdade estática (s-D)=(s) (mal posta, morte) mas pela existência de um gerador comum Ut. As três propriedades, com o estatuto de cada uma marcado: itemize4pt [D;V1] frequência fundamental comum. O gap espectral g= é um único invariante --- entropia, expoente de Lyapunov, temperatura tropical, o mesmo número --- do gato A~transcricao; ele é a frequência fundamental candidata de que tanto o espectro log-periódico de base","epistemico":"fato","exemplos":[],"id":"a_realizacao","memoria":"semantica","meta":{"dominio":"paper","nivel":5,"verificacao":"slug idêntico — enriquecer, não duplicar"},"origem":"paper","palavras_chave":[],"sinonimos":[],"tipo":"conceito","titulo":"A realização"}
\section{A realização}
D Aplicamos a Definição\~{}2 à Hipótese de Riemann, lida não pela igualdade estática (s-D)=(s) (mal posta, morte) mas pela existência de um gerador comum Ut. As três propriedades, com o estatuto de cada uma marcado: itemize4pt [D;V1] frequência fundamental comum. O gap espectral g= é um único invariante --- entropia, expoente de Lyapunov, temperatura tropical, o mesmo número --- do gato A\~{}transcricao; ele é a frequência fundamental candidata de que tanto o espectro log-periódico de base
\prov{relações: parte\_de vida}
\prov{proveniência: paper \textperiodcentered{} fato \textperiodcentered{} confiança 1.0 \textperiodcentered{} domínio paper \textperiodcentered{} história 5 versões no jornal}

%CRISTAL {"arestas":[{"alvo":"redes","confianca":1.0,"epistemico":"fato","origem":"paper","rel":"parte_de"}],"confianca":1.0,"contraexemplos":[],"descricao":"A aplicação não trata o tempo. Não bufferiza, não reordena, não guarda na sua camada (a pilha IP/UDP tem seus próprios buffers): o índice de sequência serve só para idempotência --- a chave do índice n é recomputável de qualquer n (Hⁿ(banda)), então repetir um n já visto não reaplica, mas nada se descarta nem se ordena. Não há forward secrecy (a banda é segredo permanente, recomputável); não há histórico. É o princípio da série (localizar, não empilhar): o bit corre solto, e cada data","epistemico":"fato","exemplos":[],"id":"a_reciprocidade_mao_dupla_mesma_banda","memoria":"semantica","meta":{"dominio":"paper","nivel":5,"verificacao":"conceito novo"},"origem":"paper","palavras_chave":[],"sinonimos":[],"tipo":"conceito","titulo":"A reciprocidade: mão dupla, mesma banda"}
\section{A reciprocidade: mão dupla, mesma banda}
A aplicação não trata o tempo. Não bufferiza, não reordena, não guarda na sua camada (a pilha IP/UDP tem seus próprios buffers): o índice de sequência serve só para idempotência --- a chave do índice n é recomputável de qualquer n (H\ensuremath{^{n}}(banda)), então repetir um n já visto não reaplica, mas nada se descarta nem se ordena. Não há forward secrecy (a banda é segredo permanente, recomputável); não há histórico. É o princípio da série (localizar, não empilhar): o bit corre solto, e cada data
\prov{relações: parte\_de redes}
\prov{proveniência: paper \textperiodcentered{} fato \textperiodcentered{} confiança 1.0 \textperiodcentered{} domínio paper \textperiodcentered{} história 3 versões no jornal}

%CRISTAL {"arestas":[{"alvo":"main","confianca":1.0,"epistemico":"fato","origem":"paper","rel":"parte_de"}],"confianca":1.0,"contraexemplos":[],"descricao":"A Máquina não é uma construção bottom-up sobre: é uma rede neural, e o seu osso é um único objeto topológico --- a classe de conjugação em GL(2, Z) do gato de Arnold. O flip, o gap (XOR) e tudo o que parece primitivo são a sombra em desse gato, lida um bit por vez. Não há ponto flutuante, real ou relógio no substrato: há um operador e um bit. defi[O neurônio] O neurônio mínimo da Máquina é o gato de Arnold [ A=pmatrix1&1 1&0pmatrix, estado x (um vetor, ou um bit em ), p","epistemico":"fato","exemplos":[],"id":"a_rede_neural_de_ouro","memoria":"semantica","meta":{"dominio":"paper","nivel":5,"verificacao":"slug idêntico — enriquecer, não duplicar"},"origem":"paper","palavras_chave":[],"sinonimos":[],"tipo":"conceito","titulo":"A rede neural de ouro"}
\section{A rede neural de ouro}
A Máquina não é uma construção bottom-up sobre: é uma rede neural, e o seu osso é um único objeto topológico --- a classe de conjugação em GL(2, Z) do gato de Arnold. O flip, o gap (XOR) e tudo o que parece primitivo são a sombra em desse gato, lida um bit por vez. Não há ponto flutuante, real ou relógio no substrato: há um operador e um bit. defi[O neurônio] O neurônio mínimo da Máquina é o gato de Arnold [ A=pmatrix1\&1 1\&0pmatrix, estado x (um vetor, ou um bit em ), p
\prov{relações: parte\_de main}
\prov{proveniência: paper \textperiodcentered{} fato \textperiodcentered{} confiança 1.0 \textperiodcentered{} domínio paper \textperiodcentered{} história 5 versões no jornal}

%CRISTAL {"arestas":[{"alvo":"broca_os","confianca":1.0,"epistemico":"fato","origem":"paper","rel":"parte_de"}],"confianca":1.0,"contraexemplos":[],"descricao":"D A coordenação não transmite a mensagem: transmite só a diferença contra um segredo compartilhado. A mensagem viaja inteira, em bits brutos, jamais picotada. defi[Bump] Com uma referência ref que não viaja, o keystream K(ref) a expande por iteração de hash. A mensagem trafega como [ bump=msg K(ref), msg=bump K(ref), ] em bits brutos, a mensagem inteira. Manda = recebe: a mesma operação embrulha e desembrulha. defi O bump é involução e não revela a referência K(ref) é i","epistemico":"fato","exemplos":[],"id":"a_rede_sem_dono_o_teorema_de_fechamento","memoria":"semantica","meta":{"dominio":"paper","nivel":5,"verificacao":"slug idêntico — enriquecer, não duplicar"},"origem":"paper","palavras_chave":[],"sinonimos":[],"tipo":"conceito","titulo":"A rede sem dono: o teorema de fechamento"}
\section{A rede sem dono: o teorema de fechamento}
D A coordenação não transmite a mensagem: transmite só a diferença contra um segredo compartilhado. A mensagem viaja inteira, em bits brutos, jamais picotada. defi[Bump] Com uma referência ref que não viaja, o keystream K(ref) a expande por iteração de hash. A mensagem trafega como [ bump=msg K(ref), msg=bump K(ref), ] em bits brutos, a mensagem inteira. Manda = recebe: a mesma operação embrulha e desembrulha. defi O bump é involução e não revela a referência K(ref) é i
\prov{relações: parte\_de broca\_os}
\prov{proveniência: paper \textperiodcentered{} fato \textperiodcentered{} confiança 1.0 \textperiodcentered{} domínio paper \textperiodcentered{} história 9 versões no jornal}

%CRISTAL {"arestas":[{"alvo":"leitura_aritmetica_a_regua_de_gentil_o_agm","confianca":1.0,"epistemico":"fato","origem":"paper","rel":"parte_de"}],"confianca":1.0,"contraexemplos":[],"descricao":"e a carne fractal A régua da Secão~ nasceu de uma conservacão algébrica; mas a ideia de régua --- um instrumento que mede um espaco cuja métrica não é dada, e sim construda do próprio espaco --- tem uma face puramente aritmética, e segui-la até o fim leva da média aritmético-geométrica de Gauss a um conjunto de Cantor cuja dimensão é a do número de ouro. Esta leitura é matemática pura, sem fsica, e fecha a série (Partes X","epistemico":"fato","exemplos":[],"id":"a_regua_de_gentil_e_a_cuspide","memoria":"semantica","meta":{"dominio":"paper","nivel":5,"verificacao":"overlap moderado — revisar manualmente"},"origem":"paper","palavras_chave":[],"sinonimos":[],"tipo":"conceito","titulo":"A régua de Lopes e a cúspide"}
\section{A régua de Lopes e a cúspide}
e a carne fractal A régua da Secão\~{} nasceu de uma conservacão algébrica; mas a ideia de régua --- um instrumento que mede um espaco cuja métrica não é dada, e sim construda do próprio espaco --- tem uma face puramente aritmética, e segui-la até o fim leva da média aritmético-geométrica de Gauss a um conjunto de Cantor cuja dimensão é a do número de ouro. Esta leitura é matemática pura, sem fsica, e fecha a série (Partes X
\prov{relações: parte\_de leitura\_aritmetica\_a\_regua\_de\_gentil\_o\_agm}
\prov{proveniência: paper \textperiodcentered{} fato \textperiodcentered{} confiança 1.0 \textperiodcentered{} domínio paper \textperiodcentered{} história 5 versões no jornal}

%CRISTAL {"arestas":[{"alvo":"mecanica","confianca":1.0,"epistemico":"fato","origem":"paper","rel":"parte_de"}],"confianca":1.0,"contraexemplos":[],"descricao":"Com Schrödinger derivada, a dinâmica de um qubit é a exponencial do gerador: H=2, !! B, e U(t)=e⁻iHt/ gira o versor --- a precessão de Larmor na esfera de Bloch --- sem mudar a norma. Não é um oscilador: é primeira ordem, a órbita do gato no ponto hermitiano. O gerador é auto-adjunto porque vive em, a própria álgebra, não por postulado externo.","epistemico":"fato","exemplos":[],"id":"a_representacao_e_a_propria_algebra_qutrit_e_bell","memoria":"semantica","meta":{"dominio":"paper","nivel":5,"verificacao":"slug idêntico — enriquecer, não duplicar"},"origem":"paper","palavras_chave":[],"sinonimos":[],"tipo":"conceito","titulo":"A representação é a própria álgebra: qutrit e Bell"}
\section{A representação é a própria álgebra: qutrit e Bell}
Com Schrödinger derivada, a dinâmica de um qubit é a exponencial do gerador: H=2, !! B, e U(t)=e\ensuremath{^{-}}iHt/ gira o versor --- a precessão de Larmor na esfera de Bloch --- sem mudar a norma. Não é um oscilador: é primeira ordem, a órbita do gato no ponto hermitiano. O gerador é auto-adjunto porque vive em, a própria álgebra, não por postulado externo.
\prov{relações: parte\_de mecanica}
\prov{proveniência: paper \textperiodcentered{} fato \textperiodcentered{} confiança 1.0 \textperiodcentered{} domínio paper \textperiodcentered{} história 10 versões no jornal}

%CRISTAL {"arestas":[{"alvo":"mecanica_quantica","confianca":1.0,"epistemico":"fato","origem":"paper","rel":"parte_de"},{"alvo":"mecanica","confianca":1.0,"epistemico":"fato","origem":"paper","rel":"parte_de"}],"confianca":1.0,"contraexemplos":[],"descricao":"A dinâmica de um qubit não é um oscilador --- é de primeira ordem: a equação de Schrödinger i,=H, com H=2,!! B um gerador de SU(2). A evolução U(t)=e⁻iHt gira o versor --- a precessão de Larmor na esfera de Bloch --- sem mudar a norma. O gerador é auto-adjunto porque vive em su(2), a própria álgebra, não por postulado externo. Ressalva (corrigida): a equação de segunda ordem =2 não é a dinâmica do qubit --- é a do minisuperspace cosmológico (Estatuto), outro regime. A dinâmi","epistemico":"fato","exemplos":[],"id":"a_representacao_e_a_propria_algebra_su","memoria":"semantica","meta":{"dominio":"paper","nivel":5,"verificacao":"slug idêntico — enriquecer, não duplicar"},"origem":"paper","palavras_chave":[],"sinonimos":[],"tipo":"conceito","titulo":"A representação é a própria álgebra: su"}
\section{A representação é a própria álgebra: su}
A dinâmica de um qubit não é um oscilador --- é de primeira ordem: a equação de Schrödinger i,=H, com H=2,!! B um gerador de SU(2). A evolução U(t)=e\ensuremath{^{-}}iHt gira o versor --- a precessão de Larmor na esfera de Bloch --- sem mudar a norma. O gerador é auto-adjunto porque vive em su(2), a própria álgebra, não por postulado externo. Ressalva (corrigida): a equação de segunda ordem =2 não é a dinâmica do qubit --- é a do minisuperspace cosmológico (Estatuto), outro regime. A dinâmi
\prov{relações: parte\_de mecanica\_quantica; parte\_de mecanica}
\prov{proveniência: paper \textperiodcentered{} fato \textperiodcentered{} confiança 1.0 \textperiodcentered{} domínio paper \textperiodcentered{} história 6 versões no jornal}

%CRISTAL {"arestas":[{"alvo":"paper_45_confluencia_ouro","confianca":1.0,"epistemico":"fato","origem":"paper","rel":"parte_de"}],"confianca":1.0,"contraexemplos":[],"descricao":"abstract spacing1.075 Perguntamos se há uma relacão entre, a lemniscática e o número de ouro. A resposta algébrica é não: o algoritmo PSLQ não encontra relacão inteira pequena entre, --- nem incluindo ou ---, coerente com a independência algébrica de e (Chudnovsky) e a transcendência de. Mas há uma relacão estrutural, não numérica, em que os quatro comparecem, cada um no seu ofcio: dão a curva modular (a casa do; a lemnis","epistemico":"fato","exemplos":[],"id":"a_resposta_algebrica_nao_ha_equacao","memoria":"semantica","meta":{"dominio":"paper","nivel":5,"verificacao":"conceito novo"},"origem":"paper","palavras_chave":[],"sinonimos":[],"tipo":"conceito","titulo":"A resposta algébrica: não há equação"}
\section{A resposta algébrica: não há equação}
abstract spacing1.075 Perguntamos se há uma relacão entre, a lemniscática e o número de ouro. A resposta algébrica é não: o algoritmo PSLQ não encontra relacão inteira pequena entre, --- nem incluindo ou ---, coerente com a independência algébrica de e (Chudnovsky) e a transcendência de. Mas há uma relacão estrutural, não numérica, em que os quatro comparecem, cada um no seu ofcio: dão a curva modular (a casa do; a lemnis
\prov{relações: parte\_de paper\_45\_confluencia\_ouro}
\prov{proveniência: paper \textperiodcentered{} fato \textperiodcentered{} confiança 1.0 \textperiodcentered{} domínio paper \textperiodcentered{} história 3 versões no jornal}

%CRISTAL {"arestas":[{"alvo":"teoremas_fundamentais","confianca":1.0,"epistemico":"fato","origem":"paper","rel":"parte_de"}],"confianca":1.0,"contraexemplos":[],"descricao":"theorem[Atrator AGM] Com A(a,b)=a+b2 (média do lado real), G(a,b)=ab (média do relativístico), a b>0: [(i)] G(a,b)=!(A( a, b)): a média geométrica é a aritmética transportada ao lado relativístico. [(ii)] A iteração (a,b)(A,G) converge a (a,b), com bₙ!, aₙ! (pelos dois lados); o único ponto fixo é a=b. [(iii)] Convergência quadrática: aₙ+₁-bₙ+₁=(aₙ-bₙ)²2(aₙ-bₙ)²8bₙ. [(iv)] não tem, em geral, forma fechada elementar: a constante de","epistemico":"fato","exemplos":[],"id":"a_reta_auto-construida_a_partir_dos_bits","memoria":"semantica","meta":{"dominio":"paper","nivel":5,"verificacao":"conceito novo"},"origem":"paper","palavras_chave":[],"sinonimos":[],"tipo":"conceito","titulo":"A reta auto-construída a partir dos bits"}
\section{A reta auto-construída a partir dos bits}
theorem[Atrator AGM] Com A(a,b)=a+b2 (média do lado real), G(a,b)=ab (média do relativístico), a b>0: [(i)] G(a,b)=!(A( a, b)): a média geométrica é a aritmética transportada ao lado relativístico. [(ii)] A iteração (a,b)(A,G) converge a (a,b), com b\ensuremath{_{n}}!, a\ensuremath{_{n}}! (pelos dois lados); o único ponto fixo é a=b. [(iii)] Convergência quadrática: a\ensuremath{_{n}}+\ensuremath{_{1}}-b\ensuremath{_{n}}+\ensuremath{_{1}}=(a\ensuremath{_{n}}-b\ensuremath{_{n}})\ensuremath{^{2}}2(a\ensuremath{_{n}}-b\ensuremath{_{n}})\ensuremath{^{2}}8b\ensuremath{_{n}}. [(iv)] não tem, em geral, forma fechada elementar: a constante de
\prov{relações: parte\_de teoremas\_fundamentais}
\prov{proveniência: paper \textperiodcentered{} fato \textperiodcentered{} confiança 1.0 \textperiodcentered{} domínio paper \textperiodcentered{} história 3 versões no jornal}

%CRISTAL {"arestas":[{"alvo":"transformada_dourada","confianca":1.0,"epistemico":"fato","origem":"paper","rel":"parte_de"},{"alvo":"paper_2_reta_ouro","confianca":1.0,"epistemico":"fato","origem":"paper","rel":"referencia"},{"alvo":"reais","confianca":1.0,"epistemico":"fato","origem":"paper","rel":"explica"}],"confianca":1.0,"contraexemplos":[],"descricao":"o corpo é a reta real O sistema completo --- onde se soma, multiplica, divide, se compara (ordem) e onde toda sequência que se aperta converge (completude) --- é exatamente R, e não há outro: todo sistema com todas essas qualidades é, a menos de isomorfismo, R. teo quê, em uma linha. Os racionais já são densos (toda real é limite de convergentes); um sistema completo que os contenha não pode parar antes de R nem ir além. As duas metades do começo se reencontram. Nenhuma dimen","epistemico":"fato","exemplos":[],"id":"a_reta_de_ouro","memoria":"semantica","meta":{"dominio":"paper","duplicata_sugerida":[],"nivel":5,"verificacao":"parte de conceito mais amplo 'a_reta_de_ouro'"},"origem":"paper","palavras_chave":[],"sinonimos":[],"tipo":"conceito","titulo":"A reta de ouro"}
\section{A reta de ouro}
o corpo é a reta real O sistema completo --- onde se soma, multiplica, divide, se compara (ordem) e onde toda sequência que se aperta converge (completude) --- é exatamente R, e não há outro: todo sistema com todas essas qualidades é, a menos de isomorfismo, R. teo quê, em uma linha. Os racionais já são densos (toda real é limite de convergentes); um sistema completo que os contenha não pode parar antes de R nem ir além. As duas metades do começo se reencontram. Nenhuma dimen
\prov{relações: parte\_de transformada\_dourada; referencia paper\_2\_reta\_ouro; explica reais}
\prov{proveniência: paper \textperiodcentered{} fato \textperiodcentered{} confiança 1.0 \textperiodcentered{} domínio paper \textperiodcentered{} história 84 versões no jornal}

%CRISTAL {"arestas":[{"alvo":"paper_2_reta_ouro","confianca":1.0,"epistemico":"fato","origem":"paper","rel":"parte_de"},{"alvo":"beta_numeracao_metalica","confianca":1.0,"epistemico":"fato","origem":"wiki","rel":"especializa"}],"confianca":1.0,"contraexemplos":[],"descricao":"Seja =1+52=1,6180339, a raiz positiva de x²=x+1. A base de ouro (Bergman, 1957) representa todo real x0 como soma de potências de com dígitos binários: [ x=ᵢ dᵢ, i, dᵢ0,1, ] e a forma é canônica (única) exatamente quando não há dois~1 consecutivos. a base- é a poda 11 --- firme A identidade,i⁺²=,i⁺¹+,i é, em dígitos, a reescrita [ 100 011(nas posições i+2, i+1, i), ] o carry que elimina todo bloco 11. Logo a for","epistemico":"fato","exemplos":[],"id":"a_reta_discreta_exponencial_e_autossimilar","memoria":"semantica","meta":{"dominio":"paper","nivel":5,"verificacao":"conceito novo"},"origem":"paper","palavras_chave":[],"sinonimos":[],"tipo":"conceito","titulo":"A reta discreta: exponencial e autossimilar"}
\section{A reta discreta: exponencial e autossimilar}
Seja =1+52=1,6180339, a raiz positiva de x\ensuremath{^{2}}=x+1. A base de ouro (Bergman, 1957) representa todo real x0 como soma de potências de com dígitos binários: [ x=\ensuremath{_{i}} d\ensuremath{_{i}}, i, d\ensuremath{_{i}}0,1, ] e a forma é canônica (única) exatamente quando não há dois\~{}1 consecutivos. a base- é a poda 11 --- firme A identidade,i\ensuremath{^{+2}}=,i\ensuremath{^{+1}}+,i é, em dígitos, a reescrita [ 100 011(nas posições i+2, i+1, i), ] o carry que elimina todo bloco 11. Logo a for
\prov{relações: parte\_de paper\_2\_reta\_ouro; especializa beta\_numeracao\_metalica}
\prov{proveniência: paper \textperiodcentered{} fato \textperiodcentered{} confiança 1.0 \textperiodcentered{} domínio paper \textperiodcentered{} história 4 versões no jornal}

%CRISTAL {"arestas":[{"alvo":"antena","confianca":1.0,"epistemico":"fato","origem":"paper","rel":"parte_de"}],"confianca":1.0,"contraexemplos":[],"descricao":"O projeto fecha como camada lógica e para, exatamente, onde começa a matéria. Ficam por fixar, e todos são do mundo, não da máquina: [ I₀, p(t), d (e ᵢj=k,rᵢj), eff, k cut numérico, Z₀, Z ᵢₙ, rad, wᵢj. ] Do lado algébrico, a máquina também fecha L e; o que permanece aberto é sem célula branca (A=0, limite do ). a máquina dá a fase, o mundo dá o número A máquina fornece, e fecha, a codificação: as fases am, seed, o padrão a irradiar, o roteame","epistemico":"fato","exemplos":[],"id":"a_sabatina_registro","memoria":"semantica","meta":{"dominio":"paper","nivel":5,"verificacao":"overlap moderado — revisar manualmente"},"origem":"paper","palavras_chave":[],"sinonimos":[],"tipo":"conceito","titulo":"A sabatina (registro)"}
\section{A sabatina (registro)}
O projeto fecha como camada lógica e para, exatamente, onde começa a matéria. Ficam por fixar, e todos são do mundo, não da máquina: [ I\ensuremath{_{0}}, p(t), d (e \ensuremath{_{i}}j=k,r\ensuremath{_{i}}j), eff, k cut numérico, Z\ensuremath{_{0}}, Z \ensuremath{_{in}}, rad, w\ensuremath{_{i}}j. ] Do lado algébrico, a máquina também fecha L e; o que permanece aberto é sem célula branca (A=0, limite do ). a máquina dá a fase, o mundo dá o número A máquina fornece, e fecha, a codificação: as fases am, seed, o padrão a irradiar, o roteame
\prov{relações: parte\_de antena}
\prov{proveniência: paper \textperiodcentered{} fato \textperiodcentered{} confiança 1.0 \textperiodcentered{} domínio paper \textperiodcentered{} história 3 versões no jornal}

%CRISTAL {"arestas":[{"alvo":"termodinamica","confianca":1.0,"epistemico":"fato","origem":"paper","rel":"parte_de"}],"confianca":1.0,"contraexemplos":[],"descricao":"A energia é a forma quadrática Q. O furo a evitar (e a corrigir): a multiplicatividade |z w|²=|z|²|w|² não é, por si, conservação ao longo da dinâmica --- é invariância algébrica. A primeira lei correta é a decomposição da identidade de Hurwitz. Ao multiplicar pela família s (com produto vetorial), [ |zs w|²U_sᵢstema;+;(1-s²),| a b|²U_reservatórᵢo;=;|z|²|w|²U_total. ] A energia total é fixa; o que o sistema perde, o reservatório ganha. O “cal","epistemico":"fato","exemplos":[],"id":"a_segunda_lei_o_teorema_h","memoria":"semantica","meta":{"dominio":"paper","nivel":5,"verificacao":"slug idêntico — enriquecer, não duplicar"},"origem":"paper","palavras_chave":[],"sinonimos":[],"tipo":"conceito","titulo":"A segunda lei: o teorema H"}
\section{A segunda lei: o teorema H}
A energia é a forma quadrática Q. O furo a evitar (e a corrigir): a multiplicatividade |z w|\ensuremath{^{2}}=|z|\ensuremath{^{2}}|w|\ensuremath{^{2}} não é, por si, conservação ao longo da dinâmica --- é invariância algébrica. A primeira lei correta é a decomposição da identidade de Hurwitz. Ao multiplicar pela família s (com produto vetorial), [ |zs w|\ensuremath{^{2}}U\_s\ensuremath{_{i}}stema;+;(1-s\ensuremath{^{2}}),| a b|\ensuremath{^{2}}U\_reservatór\ensuremath{_{i}}o;=;|z|\ensuremath{^{2}}|w|\ensuremath{^{2}}U\_total. ] A energia total é fixa; o que o sistema perde, o reservatório ganha. O “cal
\prov{relações: parte\_de termodinamica}
\prov{proveniência: paper \textperiodcentered{} fato \textperiodcentered{} confiança 1.0 \textperiodcentered{} domínio paper \textperiodcentered{} história 11 versões no jornal}

%CRISTAL {"arestas":[{"alvo":"seguranca_por_irreversibilidade","confianca":1.0,"epistemico":"fato","origem":"paper","rel":"parte_de"}],"confianca":1.0,"contraexemplos":[],"descricao":"A máquina de ouro é segura por design, não por dificuldade computacional --- e a razão é a série inteira: a segurança é o lado de segurança da irreversibilidade. A máquina é dissipativa (Termodinâmica): apaga informação, é muitos-para-um, e o que apaga não volta. Esse apagamento --- a segunda lei, o setor nilpotente de O Corpo Glacial --- é exatamente o que nenhum adversário desfaz. Formaliza-se como um teorema algébrico sob modelo de ameaça explícito (o adversário dispõe só das operações da p","epistemico":"fato","exemplos":[],"id":"a_seguranca_e_a_irreversibilidade","memoria":"semantica","meta":{"dominio":"paper","nivel":5,"verificacao":"conceito novo"},"origem":"paper","palavras_chave":[],"sinonimos":[],"tipo":"conceito","titulo":"A segurança é a irreversibilidade"}
\section{A segurança é a irreversibilidade}
A máquina de ouro é segura por design, não por dificuldade computacional --- e a razão é a série inteira: a segurança é o lado de segurança da irreversibilidade. A máquina é dissipativa (Termodinâmica): apaga informação, é muitos-para-um, e o que apaga não volta. Esse apagamento --- a segunda lei, o setor nilpotente de O Corpo Glacial --- é exatamente o que nenhum adversário desfaz. Formaliza-se como um teorema algébrico sob modelo de ameaça explícito (o adversário dispõe só das operações da p
\prov{relações: parte\_de seguranca\_por\_irreversibilidade}
\prov{proveniência: paper \textperiodcentered{} fato \textperiodcentered{} confiança 1.0 \textperiodcentered{} domínio paper \textperiodcentered{} história 3 versões no jornal}

%CRISTAL {"arestas":[{"alvo":"paper_2_reta_ouro","confianca":1.0,"epistemico":"fato","origem":"paper","rel":"parte_de"},{"alvo":"beta_numeracao_metalica","confianca":1.0,"epistemico":"fato","origem":"wiki","rel":"relaciona"}],"confianca":1.0,"contraexemplos":[],"descricao":"[1;1,1, e os convergentes de Fibonacci --- firme]teo:fc A fração contínua de é a mais simples possível, [ =[1;1,1,1,]=1+11+11+, ] auto-referente: =1+1/, a mesma recursão ²=+1 da base e da poda 11. Os seus convergentes são as razões de Fibonacci pₙ/qₙ=Fₙ+₁/Fₙ --- as truncagens racionais da reta de ouro. E, por Hurwitz, é o número mais irracional: o pior aproximável por racionais, com constante de Lagrange 1/5; os 1 na fração contínua são a convergência","epistemico":"fato","exemplos":[],"id":"a_simplificacao_o_osso_binario_o_ouro_na_decodificacao","memoria":"semantica","meta":{"dominio":"paper","nivel":5,"verificacao":"conceito novo"},"origem":"paper","palavras_chave":[],"sinonimos":[],"tipo":"conceito","titulo":"A simplificação: o osso binário, o ouro na decodificação"}
\section{A simplificação: o osso binário, o ouro na decodificação}
[1;1,1, e os convergentes de Fibonacci --- firme]teo:fc A fração contínua de é a mais simples possível, [ =[1;1,1,1,]=1+11+11+, ] auto-referente: =1+1/, a mesma recursão \ensuremath{^{2}}=+1 da base e da poda 11. Os seus convergentes são as razões de Fibonacci p\ensuremath{_{n}}/q\ensuremath{_{n}}=F\ensuremath{_{n}}+\ensuremath{_{1}}/F\ensuremath{_{n}} --- as truncagens racionais da reta de ouro. E, por Hurwitz, é o número mais irracional: o pior aproximável por racionais, com constante de Lagrange 1/5; os 1 na fração contínua são a convergência
\prov{relações: parte\_de paper\_2\_reta\_ouro; relaciona beta\_numeracao\_metalica}
\prov{proveniência: paper \textperiodcentered{} fato \textperiodcentered{} confiança 1.0 \textperiodcentered{} domínio paper \textperiodcentered{} história 4 versões no jornal}

%CRISTAL {"arestas":[{"alvo":"termodinamica","confianca":1.0,"epistemico":"fato","origem":"paper","rel":"parte_de"},{"alvo":"transformacao_linear","confianca":1.0,"epistemico":"fato","origem":"paper","rel":"usa"}],"confianca":1.0,"contraexemplos":[],"descricao":"A máquina não queima energia elétrica: queima ordem. Cada bit de ordem consumido libera trabalho (o princípio de Szilard/Landauer: kBT2 por bit apagado). O “alimento” da máquina é a estrutura --- bits organizados ---, e o rejeito é a desordem (a entropia que sobe a torre). Estatuto: esta seção é interpretativa (a ponte com Szilard/Landauer), apoiada na contagem de bits, não em uma medida térmica direta.","epistemico":"fato","exemplos":[],"id":"a_temperatura_o_gap_espectral","memoria":"semantica","meta":{"dominio":"paper","nivel":5,"verificacao":"slug idêntico — enriquecer, não duplicar"},"origem":"paper","palavras_chave":[],"sinonimos":[],"tipo":"conceito","titulo":"A temperatura: o gap espectral"}
\section{A temperatura: o gap espectral}
A máquina não queima energia elétrica: queima ordem. Cada bit de ordem consumido libera trabalho (o princípio de Szilard/Landauer: kBT2 por bit apagado). O “alimento” da máquina é a estrutura --- bits organizados ---, e o rejeito é a desordem (a entropia que sobe a torre). Estatuto: esta seção é interpretativa (a ponte com Szilard/Landauer), apoiada na contagem de bits, não em uma medida térmica direta.
\prov{relações: parte\_de termodinamica; usa transformacao\_linear}
\prov{proveniência: paper \textperiodcentered{} fato \textperiodcentered{} confiança 1.0 \textperiodcentered{} domínio paper \textperiodcentered{} história 63 versões no jornal}

%CRISTAL {"arestas":[{"alvo":"fase3","confianca":1.0,"epistemico":"fato","origem":"paper","rel":"parte_de"}],"confianca":1.0,"contraexemplos":[],"descricao":"empty !85A biblioteca é o produto técnico; o conhecimento é o diferencial. A engine não se reescreve --- evolui-se o tecido. abstract Com estável (V), iniciamos a Fase~3: transformar Tiffany numa plataforma comercial composta por Engine + Base Geral + Base Cliente + Memória Viva. Este documento não implementa --- projeta a arquitetura detalhada, o formato dos artefatos, o pipeline de ingestão (chunking), o orquestrador multi-tecidos, o produto distribuível e os bench","epistemico":"fato","exemplos":[],"id":"a_tese_da_fase_3","memoria":"semantica","meta":{"dominio":"paper","nivel":5,"verificacao":"overlap moderado — revisar manualmente"},"origem":"paper","palavras_chave":[],"sinonimos":[],"tipo":"conceito","titulo":"A tese da Fase 3"}
\section{A tese da Fase 3}
empty !85A biblioteca é o produto técnico; o conhecimento é o diferencial. A engine não se reescreve --- evolui-se o tecido. abstract Com estável (V), iniciamos a Fase\~{}3: transformar Tiffany numa plataforma comercial composta por Engine + Base Geral + Base Cliente + Memória Viva. Este documento não implementa --- projeta a arquitetura detalhada, o formato dos artefatos, o pipeline de ingestão (chunking), o orquestrador multi-tecidos, o produto distribuível e os bench
\prov{relações: parte\_de fase3}
\prov{proveniência: paper \textperiodcentered{} fato \textperiodcentered{} confiança 1.0 \textperiodcentered{} domínio paper \textperiodcentered{} história 3 versões no jornal}

%CRISTAL {"arestas":[{"alvo":"assistente","confianca":1.0,"epistemico":"fato","origem":"paper","rel":"parte_de"}],"confianca":1.0,"contraexemplos":[],"descricao":"empty !85O motor estabilizou --- o gargalo passa a ser o conhecimento que o alimenta. A mesma Tiffany; tecidos distintos. abstract Com estável (V: 4,2M nós, 90,ms), o gargalo passa a ser o conhecimento. Formalizamos a assistente em camadas: (1)~ --- API estável; (2)~orquestrador --- encaminha; (3)~agentes + executor + avaliador; (4)~conhecimento (imutável, versionado) e memória (mutável). Produto = engine + dados + ferramentas. Validação técnica (§X--XII): arqu","epistemico":"fato","exemplos":[],"id":"a_tese_motor_pronto_conhecimento_pendente","memoria":"semantica","meta":{"dominio":"paper","nivel":5,"verificacao":"conceito novo"},"origem":"paper","palavras_chave":[],"sinonimos":[],"tipo":"conceito","titulo":"A tese: motor pronto, conhecimento pendente"}
\section{A tese: motor pronto, conhecimento pendente}
empty !85O motor estabilizou --- o gargalo passa a ser o conhecimento que o alimenta. A mesma Tiffany; tecidos distintos. abstract Com estável (V: 4,2M nós, 90,ms), o gargalo passa a ser o conhecimento. Formalizamos a assistente em camadas: (1)\~{} --- API estável; (2)\~{}orquestrador --- encaminha; (3)\~{}agentes + executor + avaliador; (4)\~{}conhecimento (imutável, versionado) e memória (mutável). Produto = engine + dados + ferramentas. Validação técnica (§X--XII): arqu
\prov{relações: parte\_de assistente}
\prov{proveniência: paper \textperiodcentered{} fato \textperiodcentered{} confiança 1.0 \textperiodcentered{} domínio paper \textperiodcentered{} história 3 versões no jornal}

%CRISTAL {"arestas":[{"alvo":"geometria","confianca":1.0,"epistemico":"fato","origem":"paper","rel":"parte_de"},{"alvo":"conjugados_galois_parabola","confianca":1.0,"epistemico":"fato","origem":"wiki","rel":"relaciona"}],"confianca":1.0,"contraexemplos":[],"descricao":"Ao longo da família de multiplicações M=sI, o associador a=|1-s²| (não-associatividade) e o comutador c=|s| (não-comutatividade) obedecem à identidade [,a+c²=1,. ] O que se perde em associatividade ganha-se em comutatividade: a dinâmica desliza ao longo da parábola conservada. [width=0.6]figbalanca.png O invariante é a forma quadrática Q, multiplicativa (Hurwitz): |z w|²=|z|²|w|². Os dois polos da balança são distintos e não se repetem: [ Lopes 3D_s=","epistemico":"fato","exemplos":[],"id":"a_topologia_a_torre_dimensional","memoria":"semantica","meta":{"dominio":"paper","nivel":5,"verificacao":"slug idêntico — enriquecer, não duplicar"},"origem":"paper","palavras_chave":[],"sinonimos":[],"tipo":"conceito","titulo":"A topologia: a torre dimensional"}
\section{A topologia: a torre dimensional}
Ao longo da família de multiplicações M=sI, o associador a=|1-s\ensuremath{^{2}}| (não-associatividade) e o comutador c=|s| (não-comutatividade) obedecem à identidade [,a+c\ensuremath{^{2}}=1,. ] O que se perde em associatividade ganha-se em comutatividade: a dinâmica desliza ao longo da parábola conservada. [width=0.6]figbalanca.png O invariante é a forma quadrática Q, multiplicativa (Hurwitz): |z w|\ensuremath{^{2}}=|z|\ensuremath{^{2}}|w|\ensuremath{^{2}}. Os dois polos da balança são distintos e não se repetem: [ Lopes 3D\_s=
\prov{relações: parte\_de geometria; relaciona conjugados\_galois\_parabola}
\prov{proveniência: paper \textperiodcentered{} fato \textperiodcentered{} confiança 1.0 \textperiodcentered{} domínio paper \textperiodcentered{} história 14 versões no jornal}

%CRISTAL {"arestas":[{"alvo":"main","confianca":1.0,"epistemico":"fato","origem":"paper","rel":"parte_de"}],"confianca":1.0,"contraexemplos":[],"descricao":"defi[A rede] A rede é m neurônios acoplados. O acoplamento é a decomposição de Peirce: cada neurônio é uma célula branca eᵢ (idempotente, eᵢ²=eᵢ, o estado/processo), e a sinapse entre dois é um canal negro Eᵢj (i j), com ᵢ eᵢ=I (conservação). A matriz sináptica da rede de m é o companion m m do gato. defi A rede de m neurônios é o m-bonacci DN A leitura em um bit da rede de m neurônios é a gramática “sem 1m” (nenhuma corrida de m uns), cuja t","epistemico":"fato","exemplos":[],"id":"a_torre_de_hurwitzcayley--dickson","memoria":"semantica","meta":{"dominio":"paper","nivel":5,"verificacao":"slug idêntico — enriquecer, não duplicar"},"origem":"paper","palavras_chave":[],"sinonimos":[],"tipo":"conceito","titulo":"A torre de Hurwitz/Cayley--Dickson"}
\section{A torre de Hurwitz/Cayley--Dickson}
defi[A rede] A rede é m neurônios acoplados. O acoplamento é a decomposição de Peirce: cada neurônio é uma célula branca e\ensuremath{_{i}} (idempotente, e\ensuremath{_{i}}\ensuremath{^{2}}=e\ensuremath{_{i}}, o estado/processo), e a sinapse entre dois é um canal negro E\ensuremath{_{i}}j (i j), com \ensuremath{_{i}} e\ensuremath{_{i}}=I (conservação). A matriz sináptica da rede de m é o companion m m do gato. defi A rede de m neurônios é o m-bonacci DN A leitura em um bit da rede de m neurônios é a gramática “sem 1m” (nenhuma corrida de m uns), cuja t
\prov{relações: parte\_de main}
\prov{proveniência: paper \textperiodcentered{} fato \textperiodcentered{} confiança 1.0 \textperiodcentered{} domínio paper \textperiodcentered{} história 5 versões no jornal}

%CRISTAL {"arestas":[{"alvo":"geometria","confianca":1.0,"epistemico":"fato","origem":"paper","rel":"parte_de"}],"confianca":1.0,"contraexemplos":[],"descricao":"Lendo o tempo comóvel pelo fator de escala a()=, a métrica natural sobre R S³ é a de de Sitter global: [ ds²=-d²+²!;d²S₃, a/a=1 (=3, k=+1), ] de curvatura positiva. As fatias espaciais são as S³ dos versores; o observador comóvel remove a() e lê as coordenadas internas --- o referencial geométrico correto.","epistemico":"fato","exemplos":[],"id":"a_travessia_o_flip_de_wick_espacial","memoria":"semantica","meta":{"dominio":"paper","nivel":5,"verificacao":"slug idêntico — enriquecer, não duplicar"},"origem":"paper","palavras_chave":[],"sinonimos":[],"tipo":"conceito","titulo":"A travessia: o flip de Wick espacial"}
\section{A travessia: o flip de Wick espacial}
Lendo o tempo comóvel pelo fator de escala a()=, a métrica natural sobre R S\ensuremath{^{3}} é a de de Sitter global: [ ds\ensuremath{^{2}}=-d\ensuremath{^{2}}+\ensuremath{^{2}}!;d\ensuremath{^{2}}S\ensuremath{_{3}}, a/a=1 (=3, k=+1), ] de curvatura positiva. As fatias espaciais são as S\ensuremath{^{3}} dos versores; o observador comóvel remove a() e lê as coordenadas internas --- o referencial geométrico correto.
\prov{relações: parte\_de geometria}
\prov{proveniência: paper \textperiodcentered{} fato \textperiodcentered{} confiança 1.0 \textperiodcentered{} domínio paper \textperiodcentered{} história 13 versões no jornal}

%CRISTAL {"arestas":[{"alvo":"geometria","confianca":1.0,"epistemico":"fato","origem":"paper","rel":"parte_de"}],"confianca":1.0,"contraexemplos":[],"descricao":"O defeito não desaparece numa reta isolada --- em O o associador tem piso irredutível. Ele desaparece na torre R C H O (Cayley--Dickson): a cada nível a faixa correspondente é estancada e o resíduo sobe ao nível seguinte; a cadeia inteira o leva a 0. Só o que conserva Q sobe --- o filtro de Hurwitz ---; o defeito de Lagrange fica para trás e decai. A passagem dimensional não carrega resíduo bruto: separa o invariante da sujeira, e só então o andar seguinte entra estável.","epistemico":"fato","exemplos":[],"id":"a_variedade_a_s3_dos_versores","memoria":"semantica","meta":{"dominio":"paper","nivel":5,"verificacao":"slug idêntico — enriquecer, não duplicar"},"origem":"paper","palavras_chave":[],"sinonimos":[],"tipo":"conceito","titulo":"A variedade: a S³ dos versores"}
\section{A variedade: a S\ensuremath{^{3}} dos versores}
O defeito não desaparece numa reta isolada --- em O o associador tem piso irredutível. Ele desaparece na torre R C H O (Cayley--Dickson): a cada nível a faixa correspondente é estancada e o resíduo sobe ao nível seguinte; a cadeia inteira o leva a 0. Só o que conserva Q sobe --- o filtro de Hurwitz ---; o defeito de Lagrange fica para trás e decai. A passagem dimensional não carrega resíduo bruto: separa o invariante da sujeira, e só então o andar seguinte entra estável.
\prov{relações: parte\_de geometria}
\prov{proveniência: paper \textperiodcentered{} fato \textperiodcentered{} confiança 1.0 \textperiodcentered{} domínio paper \textperiodcentered{} história 13 versões no jornal}

%CRISTAL {"arestas":[{"alvo":"a_base_topologica","confianca":1.0,"epistemico":"fato","origem":"paper","rel":"parte_de"}],"confianca":1.0,"contraexemplos":[],"descricao":"Antes do número, o osso. A dinâmica da máquina --- o gato de Arnold A=(smallmatrix1&1 1&0smallmatrix) --- é um automorfismo do toro T²; seu único invariante essencial é a classe de conjugação em GL(2, Z) da ação induzida na homologia H₁(T²) Z², capturada por dois inteiros: traço e determinante, (tr,)=(1,-1). É atemporal: a classe não precisa de relógio. Deste osso, álgebra e geometria derivam.","epistemico":"fato","exemplos":[],"id":"a_variedade_e_a_acao_em_h_1","memoria":"semantica","meta":{"dominio":"paper","nivel":5,"verificacao":"conceito novo"},"origem":"paper","palavras_chave":[],"sinonimos":[],"tipo":"conceito","titulo":"A variedade e a ação em H₁"}
\section{A variedade e a ação em H\ensuremath{_{1}}}
Antes do número, o osso. A dinâmica da máquina --- o gato de Arnold A=(smallmatrix1\&1 1\&0smallmatrix) --- é um automorfismo do toro T\ensuremath{^{2}}; seu único invariante essencial é a classe de conjugação em GL(2, Z) da ação induzida na homologia H\ensuremath{_{1}}(T\ensuremath{^{2}}) Z\ensuremath{^{2}}, capturada por dois inteiros: traço e determinante, (tr,)=(1,-1). É atemporal: a classe não precisa de relógio. Deste osso, álgebra e geometria derivam.
\prov{relações: parte\_de a\_base\_topologica}
\prov{proveniência: paper \textperiodcentered{} fato \textperiodcentered{} confiança 1.0 \textperiodcentered{} domínio paper \textperiodcentered{} história 4 versões no jornal}

%CRISTAL {"arestas":[{"alvo":"paper_2_reta_ouro","confianca":1.0,"epistemico":"fato","origem":"paper","rel":"parte_de"},{"alvo":"beta_numeracao_metalica","confianca":1.0,"epistemico":"fato","origem":"wiki","rel":"especializa"}],"confianca":1.0,"contraexemplos":[],"descricao":"a reta de ouro A reta de ouro é o conjunto dos inteiros --- os x0 de expansão base- finita: x=ᵢ=-NM dᵢ,i, dᵢ0,1, sem 11. É um subconjunto denso de ₀, enumerável, fechado sob soma (com carry), e contido no anel []=a+b:a,b. definicao é o shift; Pisot é a autossimilaridade --- firme A multiplicação por desloca a palavra uma casa:i=,i⁺¹, logo é o shift de dígitos --- a reta de ouro é exponencial (peso","epistemico":"fato","exemplos":[],"id":"a_versao_inteira_zeckendorf_e_fibonacci","memoria":"semantica","meta":{"dominio":"paper","nivel":5,"verificacao":"conceito novo"},"origem":"paper","palavras_chave":[],"sinonimos":[],"tipo":"conceito","titulo":"A versão inteira: Zeckendorf e Fibonacci"}
\section{A versão inteira: Zeckendorf e Fibonacci}
a reta de ouro A reta de ouro é o conjunto dos inteiros --- os x0 de expansão base- finita: x=\ensuremath{_{i}}=-NM d\ensuremath{_{i}},i, d\ensuremath{_{i}}0,1, sem 11. É um subconjunto denso de \ensuremath{_{0}}, enumerável, fechado sob soma (com carry), e contido no anel []=a+b:a,b. definicao é o shift; Pisot é a autossimilaridade --- firme A multiplicação por desloca a palavra uma casa:i=,i\ensuremath{^{+1}}, logo é o shift de dígitos --- a reta de ouro é exponencial (peso
\prov{relações: parte\_de paper\_2\_reta\_ouro; especializa beta\_numeracao\_metalica}
\prov{proveniência: paper \textperiodcentered{} fato \textperiodcentered{} confiança 1.0 \textperiodcentered{} domínio paper \textperiodcentered{} história 4 versões no jornal}

%CRISTAL {"arestas":[{"alvo":"mecanica_quantica","confianca":0.95,"epistemico":"fato","origem":"paper","rel":"parte_de"}],"confianca":1.0,"contraexemplos":[],"descricao":"A 3-esfera S³=q:|q|=1=SU(2) dos versores é o espaço de estados de um qubit, e a esfera S² de direções a dos campos locais é a esfera de Bloch (=CP¹). A construção que segue é exata: a função de onda é uma seção de um fibrado linha sobre a Bloch. (para o leitor não-quântico). Um qubit é o sistema quântico mais simples (um spin-12); seus estados, a menos de uma fase sem efeito físico, formam a esfera de Bloch S². Um fibrado linha é uma família de retas complex","epistemico":"fato","exemplos":[],"id":"a_violacao_de_bell_e_a_conservacao_ac21","memoria":"semantica","meta":{"dominio":"paper","nivel":5,"verificacao":"conceito novo"},"origem":"paper","palavras_chave":[],"sinonimos":[],"tipo":"conceito","titulo":"A violação de Bell é a conservação a+c²=1"}
\section{A violação de Bell é a conservação a+c\ensuremath{^{2}}=1}
A 3-esfera S\ensuremath{^{3}}=q:|q|=1=SU(2) dos versores é o espaço de estados de um qubit, e a esfera S\ensuremath{^{2}} de direções a dos campos locais é a esfera de Bloch (=CP\ensuremath{^{1}}). A construção que segue é exata: a função de onda é uma seção de um fibrado linha sobre a Bloch. (para o leitor não-quântico). Um qubit é o sistema quântico mais simples (um spin-12); seus estados, a menos de uma fase sem efeito físico, formam a esfera de Bloch S\ensuremath{^{2}}. Um fibrado linha é uma família de retas complex
\prov{relações: parte\_de mecanica\_quantica}
\prov{proveniência: paper \textperiodcentered{} fato \textperiodcentered{} confiança 1.0 \textperiodcentered{} domínio paper \textperiodcentered{} história 5 versões no jornal}

%CRISTAL {"arestas":[{"alvo":"fechamento_blocos","confianca":1.0,"epistemico":"fato","origem":"paper","rel":"parte_de"}],"confianca":1.0,"contraexemplos":[],"descricao":"D Fechamento por blocos Seja M=(eᵢ,N,) uma máquina e I=_ I_ uma partição admissível (quociente Q^ acíclico). Então M^=(E_,N^,) é uma máquina da mesma classe: parte branca ortogonal completa (Lema 1), demux sem contaminação (Lema 2), parte negra nilpotente (Lema 3, pois Q^ acíclico), carga conservada (Lema 4). Logo a classe é fechada sob blocagem admissível, e a recursão da torre fecha por indução:; M;; M^;;( M^)^;; teo","epistemico":"fato","exemplos":[],"id":"admissibilidade_nao-deadlock","memoria":"semantica","meta":{"dominio":"paper","nivel":5,"verificacao":"overlap moderado — revisar manualmente"},"origem":"paper","palavras_chave":[],"sinonimos":[],"tipo":"conceito","titulo":"Admissibilidade = não-deadlock"}
\section{Admissibilidade = não-deadlock}
D Fechamento por blocos Seja M=(e\ensuremath{_{i}},N,) uma máquina e I=\_ I\_ uma partição admissível (quociente Q\^{} acíclico). Então M\^{}=(E\_,N\^{},) é uma máquina da mesma classe: parte branca ortogonal completa (Lema 1), demux sem contaminação (Lema 2), parte negra nilpotente (Lema 3, pois Q\^{} acíclico), carga conservada (Lema 4). Logo a classe é fechada sob blocagem admissível, e a recursão da torre fecha por indução:; M;; M\^{};;( M\^{})\^{};; teo
\prov{relações: parte\_de fechamento\_blocos}
\prov{proveniência: paper \textperiodcentered{} fato \textperiodcentered{} confiança 1.0 \textperiodcentered{} domínio paper \textperiodcentered{} história 4 versões no jornal}

%CRISTAL {"arestas":[{"alvo":"a_engenharia_de_ouro_pa_pgs_agm_fc_reta","confianca":1.0,"epistemico":"fato","origem":"paper","rel":"parte_de"},{"alvo":"analise","confianca":1.0,"epistemico":"fato","origem":"paper","rel":"usa"}],"confianca":1.0,"contraexemplos":[],"descricao":"progressões geométricas e sequências PGS PGS designa a família progressões geométricas e sequências usada em dois níveis: [nosep,leftmargin=1.4em] Pedagógico (Lopes, Novas Sequências): sequências aritméticas/geométricas, dupolinômios, limites monótonos --- a linguagem do cap.~10 dos Fundamentos. Operacional (engenharia): grade kₖ>₀ com =¹/; sob =, vira progressão aritmética de passo (Transformada Dourada, Prop.~grid). definicao","epistemico":"fato","exemplos":[],"id":"agm_atrator_entre_os_dois_lados","memoria":"semantica","meta":{"dominio":"paper","nivel":5,"verificacao":"slug idêntico — enriquecer, não duplicar"},"origem":"paper","palavras_chave":[],"sinonimos":[],"tipo":"conceito","titulo":"AGM: atrator entre os dois lados"}
\section{AGM: atrator entre os dois lados}
progressões geométricas e sequências PGS PGS designa a família progressões geométricas e sequências usada em dois níveis: [nosep,leftmargin=1.4em] Pedagógico (Lopes, Novas Sequências): sequências aritméticas/geométricas, dupolinômios, limites monótonos --- a linguagem do cap.\~{}10 dos Fundamentos. Operacional (engenharia): grade k\ensuremath{_{k}}>\ensuremath{_{0}} com =\ensuremath{^{1}}/; sob =, vira progressão aritmética de passo (Transformada Dourada, Prop.\~{}grid). definicao
\prov{relações: parte\_de a\_engenharia\_de\_ouro\_pa\_pgs\_agm\_fc\_reta; usa analise}
\prov{proveniência: paper \textperiodcentered{} fato \textperiodcentered{} confiança 1.0 \textperiodcentered{} domínio paper \textperiodcentered{} história 48 versões no jornal}

%CRISTAL {"arestas":[{"alvo":"kernel","confianca":0.75,"epistemico":"fato","origem":"manual","rel":"referencia"}],"confianca":1.0,"contraexemplos":[],"descricao":"As leituras anteriores tratam o espaco de álgebras como uma paisagem estática --- uma régua e uma curvatura sobre as quais se lê a fsica. As Partes~VI e~VII tornam-no dinâmico: cada álgebra passa a ser um estado que evolui, e as ligacões, colapsos.","epistemico":"fato","exemplos":[],"id":"algebra","memoria":"semantica","meta":{"dominio":"paper","duplicata_sugerida":[],"nivel":5,"verificacao":"parte de conceito mais amplo 'algebra'"},"origem":"paper","palavras_chave":[],"sinonimos":[],"tipo":"conceito","titulo":"algebra"}
\section{algebra}
As leituras anteriores tratam o espaco de álgebras como uma paisagem estática --- uma régua e uma curvatura sobre as quais se lê a fsica. As Partes\~{}VI e\~{}VII tornam-no dinâmico: cada álgebra passa a ser um estado que evolui, e as ligacões, colapsos.
\prov{relações: referencia kernel}
\prov{proveniência: paper \textperiodcentered{} fato \textperiodcentered{} confiança 1.0 \textperiodcentered{} domínio paper \textperiodcentered{} história 7 versões no jornal}

%CRISTAL {"arestas":[{"alvo":"a_base_topologica","confianca":1.0,"epistemico":"fato","origem":"paper","rel":"parte_de"}],"confianca":1.0,"contraexemplos":[],"descricao":"A homologia é homotópica: a ação fA*=A em H₁(T²) Z² não enxerga a métrica nem a realização suave --- só a classe. (Como =-1, ela vive em GL(2, Z), não SL.) A classe é fixada pelo polinômio característico pA(x)=x²-trA,x+ A. [width=0.34]basetopologicalente.png [2pt] Os dois geradores de H₁= Z² --- dois círculos tangentes no basepoint (S¹! S¹, o 1-esqueleto do toro). Os eixos são as autodireções do gato (1/ e -); F_* mistura, sem rasgar nenhum","epistemico":"fato","exemplos":[],"id":"algebra_derivada_o_espectro","memoria":"semantica","meta":{"dominio":"paper","nivel":5,"verificacao":"conceito novo"},"origem":"paper","palavras_chave":[],"sinonimos":[],"tipo":"conceito","titulo":"Álgebra derivada (o espectro)"}
\section{Álgebra derivada (o espectro)}
A homologia é homotópica: a ação fA*=A em H\ensuremath{_{1}}(T\ensuremath{^{2}}) Z\ensuremath{^{2}} não enxerga a métrica nem a realização suave --- só a classe. (Como =-1, ela vive em GL(2, Z), não SL.) A classe é fixada pelo polinômio característico pA(x)=x\ensuremath{^{2}}-trA,x+ A. [width=0.34]basetopologicalente.png [2pt] Os dois geradores de H\ensuremath{_{1}}= Z\ensuremath{^{2}} --- dois círculos tangentes no basepoint (S\ensuremath{^{1}}! S\ensuremath{^{1}}, o 1-esqueleto do toro). Os eixos são as autodireções do gato (1/ e -); F\_* mistura, sem rasgar nenhum
\prov{relações: parte\_de a\_base\_topologica}
\prov{proveniência: paper \textperiodcentered{} fato \textperiodcentered{} confiança 1.0 \textperiodcentered{} domínio paper \textperiodcentered{} história 3 versões no jornal}

%CRISTAL {"arestas":[{"alvo":"goldchain","confianca":1.0,"epistemico":"fato","origem":"paper","rel":"parte_de"},{"alvo":"kernel","confianca":0.5,"epistemico":"fato","origem":"wiki","rel":"relaciona"},{"alvo":"recursao","confianca":0.5,"epistemico":"fato","origem":"wiki","rel":"relaciona"},{"alvo":"topologia","confianca":0.5,"epistemico":"fato","origem":"wiki","rel":"relaciona"},{"alvo":"matriz_de_confusao","confianca":0.5,"epistemico":"fato","origem":"wiki","rel":"relaciona"},{"alvo":"pilha","confianca":0.5,"epistemico":"fato","origem":"wiki","rel":"relaciona"},{"alvo":"teste_ordem","confianca":0.5,"epistemico":"fato","origem":"wiki","rel":"relaciona"},{"alvo":"regua_associativa_e_o_risco_de_vazamento","confianca":0.5,"epistemico":"fato","origem":"wiki","rel":"relaciona"},{"alvo":"pipeline_de_ingestao","confianca":0.5,"epistemico":"fato","origem":"wiki","rel":"relaciona"},{"alvo":"pow_metal_setor","confianca":0.5,"epistemico":"fato","origem":"wiki","rel":"relaciona"},{"alvo":"complexos","confianca":0.5,"epistemico":"fato","origem":"wiki","rel":"relaciona"},{"alvo":"linker","confianca":0.5,"epistemico":"fato","origem":"wiki","rel":"relaciona"},{"alvo":"respiracao_dos_programas","confianca":0.5,"epistemico":"fato","origem":"wiki","rel":"relaciona"},{"alvo":"rust","confianca":0.5,"epistemico":"fato","origem":"wiki","rel":"relaciona"},{"alvo":"projeto_cristaljson","confianca":0.5,"epistemico":"fato","origem":"wiki","rel":"relaciona"}],"confianca":1.0,"contraexemplos":[],"descricao":"Paper: algebra_dos_contratos.tex","epistemico":"fato","exemplos":[],"id":"algebra_dos_contratos","memoria":"semantica","meta":{"dominio":"paper","nivel":5,"tipo_no":"paper"},"origem":"paper","palavras_chave":[],"sinonimos":[],"tipo":"conceito","titulo":"algebra_dos_contratos"}
\section{algebra\_dos\_contratos}
Paper: algebra\_dos\_contratos.tex
\prov{relações: parte\_de goldchain; relaciona kernel; relaciona recursao; relaciona topologia; relaciona matriz\_de\_confusao; relaciona pilha; relaciona teste\_ordem; relaciona regua\_associativa\_e\_o\_risco\_de\_vazamento; relaciona pipeline\_de\_ingestao; relaciona pow\_metal\_setor; relaciona complexos; relaciona linker; relaciona respiracao\_dos\_programas; relaciona rust; relaciona projeto\_cristaljson}
\prov{proveniência: paper \textperiodcentered{} fato \textperiodcentered{} confiança 1.0 \textperiodcentered{} domínio paper \textperiodcentered{} história 21 versões no jornal}

%CRISTAL {"arestas":[{"alvo":"matematica","confianca":0.95,"epistemico":"fato","origem":"paper","rel":"parte_de"}],"confianca":1.0,"contraexemplos":[],"descricao":"A métrica quântica de Lopes~gentilMQ sobre / é (x,y) =,|x-y|, 1-|x-y|,. Lopes~gentilMQ em é uma métrica sobre /: é não-negativa, simétrica, anula-se só na diagonal e satisfaz a desigualdade triangular. É invariante por translação, (x+t,y+t)=(x,y); logo toda translação (rotação do círculo) é uma isometria. proposicao Identificando x/ com e² i x S¹, a métrica é (a menos de escala) o comprimento de a","epistemico":"fato","exemplos":[],"id":"algebras_hipercomplexas_e_a_familia_s","memoria":"semantica","meta":{"dominio":"paper","nivel":5,"verificacao":"conceito novo"},"origem":"paper","palavras_chave":[],"sinonimos":[],"tipo":"conceito","titulo":"Álgebras hipercomplexas e a família s"}
\section{Álgebras hipercomplexas e a família s}
A métrica quântica de Lopes\~{}gentilMQ sobre / é (x,y) =,|x-y|, 1-|x-y|,. Lopes\~{}gentilMQ em é uma métrica sobre /: é não-negativa, simétrica, anula-se só na diagonal e satisfaz a desigualdade triangular. É invariante por translação, (x+t,y+t)=(x,y); logo toda translação (rotação do círculo) é uma isometria. proposicao Identificando x/ com e\ensuremath{^{2}} i x S\ensuremath{^{1}}, a métrica é (a menos de escala) o comprimento de a
\prov{relações: parte\_de matematica}
\prov{proveniência: paper \textperiodcentered{} fato \textperiodcentered{} confiança 1.0 \textperiodcentered{} domínio paper \textperiodcentered{} história 6 versões no jornal}

%CRISTAL {"arestas":[{"alvo":"algebra_dos_contratos","confianca":1.0,"epistemico":"fato","origem":"paper","rel":"parte_de"}],"confianca":1.0,"contraexemplos":[],"descricao":"enumerate2pt [V] Interface Contrato(P, ) + back-ends BTC/ETH (goldchain.py). [V] Swap cross-chain completo (ETH+BTC, mesmo segredo) e cadeia de certificados auditável. [V] Verificação coletiva caracterizada (N-qubits = generalização, não aceleração). [V] Otimização de gás: clone EIP-1167 --- deploy do contrato uma vez, clone por swap: 627k 212k gás (3 mais barato), custo fixo pago em 2 swaps (htlcfactory.py). [C] Fechamento paralelo (N workers, throughput N) no GEX44 (20","epistemico":"fato","exemplos":[],"id":"apendice_---_transacoes_da_demonstracao_testnet","memoria":"semantica","meta":{"dominio":"paper","nivel":5,"verificacao":"conceito novo"},"origem":"paper","palavras_chave":[],"sinonimos":[],"tipo":"conceito","titulo":"Apêndice --- transações da demonstração (testnet)"}
\section{Apêndice --- transações da demonstração (testnet)}
enumerate2pt [V] Interface Contrato(P, ) + back-ends BTC/ETH (goldchain.py). [V] Swap cross-chain completo (ETH+BTC, mesmo segredo) e cadeia de certificados auditável. [V] Verificação coletiva caracterizada (N-qubits = generalização, não aceleração). [V] Otimização de gás: clone EIP-1167 --- deploy do contrato uma vez, clone por swap: 627k 212k gás (3 mais barato), custo fixo pago em 2 swaps (htlcfactory.py). [C] Fechamento paralelo (N workers, throughput N) no GEX44 (20
\prov{relações: parte\_de algebra\_dos\_contratos}
\prov{proveniência: paper \textperiodcentered{} fato \textperiodcentered{} confiança 1.0 \textperiodcentered{} domínio paper \textperiodcentered{} história 3 versões no jornal}

%CRISTAL {"arestas":[{"alvo":"main","confianca":1.0,"epistemico":"fato","origem":"paper","rel":"parte_de"}],"confianca":1.0,"contraexemplos":[],"descricao":"D (a) A dinâmica admissível reproduz a cadeia fractal, não a torre de Hurwitz. (b) O limite em é algébrico, não consequência de T(x)= x1. (c) Coordenadas x[0,1), densidades e fases não são observáveis diretos: a Máquina observa bits, gaps e dobras; modelos contínuos valem só após redução correta. (d) [0,1] mede a realização diádica do subshift; 1,2,4,8,16 são dimensões algébricas da torre --- invariantes distintos. (e) A leitura dos problemas do milênio é interpretação externa e não f","epistemico":"fato","exemplos":[],"id":"apendices_---_a_semente_reconstrutivel","memoria":"semantica","meta":{"dominio":"paper","nivel":5,"verificacao":"slug idêntico — enriquecer, não duplicar"},"origem":"paper","palavras_chave":[],"sinonimos":[],"tipo":"conceito","titulo":"Apêndices --- a semente reconstrutível"}
\section{Apêndices --- a semente reconstrutível}
D (a) A dinâmica admissível reproduz a cadeia fractal, não a torre de Hurwitz. (b) O limite em é algébrico, não consequência de T(x)= x1. (c) Coordenadas x[0,1), densidades e fases não são observáveis diretos: a Máquina observa bits, gaps e dobras; modelos contínuos valem só após redução correta. (d) [0,1] mede a realização diádica do subshift; 1,2,4,8,16 são dimensões algébricas da torre --- invariantes distintos. (e) A leitura dos problemas do milênio é interpretação externa e não f
\prov{relações: parte\_de main}
\prov{proveniência: paper \textperiodcentered{} fato \textperiodcentered{} confiança 1.0 \textperiodcentered{} domínio paper \textperiodcentered{} história 5 versões no jornal}

%CRISTAL {"arestas":[{"alvo":"formato_da_base","confianca":1.0,"epistemico":"fato","origem":"paper","rel":"parte_de"}],"confianca":1.0,"contraexemplos":[],"descricao":"Formato estendido (quatro campos tab-separados): lstlisting <lang>\t<chunkᵢd>\t<texto>\t<origem> pt\tsoftware/git/0001842\tClone um repositorio: git clone URL destino\tgit-scm.com/docs lstlisting chunkᵢd --- caminho lógico estável (domínio/sequência), usado em gold sets, deduplicação e merge incremental. origem --- URI ou caminho da fonte (curadoria, licença).","epistemico":"fato","exemplos":[],"id":"arquivo_metajson","memoria":"semantica","meta":{"dominio":"paper","nivel":5,"verificacao":"overlap moderado — revisar manualmente"},"origem":"paper","palavras_chave":[],"sinonimos":[],"tipo":"conceito","titulo":"Arquivo meta.json"}
\section{Arquivo meta.json}
Formato estendido (quatro campos tab-separados): lstlisting <lang> <chunk\ensuremath{_{i}}d> <texto> <origem> pt software/git/0001842 Clone um repositorio: git clone URL destino git-scm.com/docs lstlisting chunk\ensuremath{_{i}}d --- caminho lógico estável (domínio/sequência), usado em gold sets, deduplicação e merge incremental. origem --- URI ou caminho da fonte (curadoria, licença).
\prov{relações: parte\_de formato\_da\_base}
\prov{proveniência: paper \textperiodcentered{} fato \textperiodcentered{} confiança 1.0 \textperiodcentered{} domínio paper \textperiodcentered{} história 4 versões no jornal}

%CRISTAL {"arestas":[{"alvo":"base_geral_core","confianca":1.0,"epistemico":"fato","origem":"paper","rel":"parte_de"},{"alvo":"word","confianca":0.5,"epistemico":"fato","origem":"wiki","rel":"relaciona"}],"confianca":1.0,"contraexemplos":[],"descricao":"I Base neutra, ampla e comercial --- pré-instalada, sem dados de cliente.\nObjetivo: assistente útil out of the box em software, matemática, documentação e sistemas.","epistemico":"fato","exemplos":[],"id":"arvore_de_dominios","memoria":"semantica","meta":{"dominio":"paper","nivel":5,"verificacao":"overlap moderado — revisar manualmente"},"origem":"paper","palavras_chave":[],"sinonimos":[],"tipo":"conceito","titulo":"Árvore de domínios"}
\section{Árvore de domínios}
I Base neutra, ampla e comercial --- pré-instalada, sem dados de cliente.
Objetivo: assistente útil out of the box em software, matemática, documentação e sistemas.
\prov{relações: parte\_de base\_geral\_core; relaciona word}
\prov{proveniência: paper \textperiodcentered{} fato \textperiodcentered{} confiança 1.0 \textperiodcentered{} domínio paper \textperiodcentered{} história 3 versões no jornal}

%CRISTAL {"arestas":[{"alvo":"leitura_das_unidades_o_sistema_por-unidade_e_o_si","confianca":1.0,"epistemico":"fato","origem":"paper","rel":"parte_de"}],"confianca":1.0,"contraexemplos":[],"descricao":"Da Definicão~ seguem, de imediato, as duas unidades cinemáticas. As unidades cinemáticas O metro é uma volta. O segundo é o tempo de c=299,792,458 voltas. A velocidade da luz é a taxa que os liga, [ c=voltassegundo=299,792,458, ] e c=1 no sistema natural significa apenas escolher a volta como unidade comum: o comprimento é um tempo. Medir um metro é contar 1,650,763,73 voltas da linha laranja do Kr-8","epistemico":"fato","exemplos":[],"id":"as_constantes_derivadas_tres_taxas_e_uma_carga","memoria":"semantica","meta":{"dominio":"paper","nivel":5,"verificacao":"conceito novo"},"origem":"paper","palavras_chave":[],"sinonimos":[],"tipo":"conceito","titulo":"As constantes derivadas: três taxas e uma carga"}
\section{As constantes derivadas: três taxas e uma carga}
Da Definicão\~{} seguem, de imediato, as duas unidades cinemáticas. As unidades cinemáticas O metro é uma volta. O segundo é o tempo de c=299,792,458 voltas. A velocidade da luz é a taxa que os liga, [ c=voltassegundo=299,792,458, ] e c=1 no sistema natural significa apenas escolher a volta como unidade comum: o comprimento é um tempo. Medir um metro é contar 1,650,763,73 voltas da linha laranja do Kr-8
\prov{relações: parte\_de leitura\_das\_unidades\_o\_sistema\_por-unidade\_e\_o\_si}
\prov{proveniência: paper \textperiodcentered{} fato \textperiodcentered{} confiança 1.0 \textperiodcentered{} domínio paper \textperiodcentered{} história 4 versões no jornal}

%CRISTAL {"arestas":[{"alvo":"colapso","confianca":0.031,"epistemico":"fato","origem":"manual","rel":"referencia"},{"alvo":"geometria","confianca":1.0,"epistemico":"fato","origem":"paper","rel":"parte_de"}],"confianca":1.0,"contraexemplos":[],"descricao":"No regime hiperbólico aparece a geometria onde a análise espectral pode entrar. Com =PSL₂ ( Z[i]) (de novo o corpo primo Z[i], agora em 3D), a orbifold de Bianchi H³/ --- de volume finito --- carrega o Laplaciano auto-adjunto, com espectro discreto mais contínuo. A parte contínua é dada pelas séries de Eisenstein, cujo coeficiente de espalhamento envolve a zeta de Dedekind Q(ᵢ)= L: a estrutura espectral contém a de Riemann como fator. Mas onde ela vive não é no espectro dis","epistemico":"fato","exemplos":[],"id":"as_cuspides_o_colapso_31","memoria":"semantica","meta":{"dominio":"paper","nivel":5,"verificacao":"slug idêntico — enriquecer, não duplicar"},"origem":"paper","palavras_chave":[],"sinonimos":[],"tipo":"conceito","titulo":"As cúspides: o colapso 31"}
\section{As cúspides: o colapso 31}
No regime hiperbólico aparece a geometria onde a análise espectral pode entrar. Com =PSL\ensuremath{_{2}} ( Z[i]) (de novo o corpo primo Z[i], agora em 3D), a orbifold de Bianchi H\ensuremath{^{3}}/ --- de volume finito --- carrega o Laplaciano auto-adjunto, com espectro discreto mais contínuo. A parte contínua é dada pelas séries de Eisenstein, cujo coeficiente de espalhamento envolve a zeta de Dedekind Q(\ensuremath{_{i}})= L: a estrutura espectral contém a de Riemann como fator. Mas onde ela vive não é no espectro dis
\prov{relações: referencia colapso; parte\_de geometria}
\prov{proveniência: paper \textperiodcentered{} fato \textperiodcentered{} confiança 1.0 \textperiodcentered{} domínio paper \textperiodcentered{} história 15 versões no jornal}

%CRISTAL {"arestas":[{"alvo":"goldcoin","confianca":1.0,"epistemico":"fato","origem":"paper","rel":"parte_de"}],"confianca":1.0,"contraexemplos":[],"descricao":"N A “dança” do escrow (trocar h, assinaturas parciais, revelar s, avisar o watchtower) é mensagem entre as partes --- e isso já está implementado e medido: o canal do bump (Fundamentos de Redes). A banda =sha256(tecido) filtra: só as partes do negócio decodificam, o resto lê ruído (privacidade da coordenação, de graça). Manda = recebe na latência, mão dupla, sem cliente/servidor. N Medido em produção (ping-pong UDP, GEX real). Não se inventa o orquestrador --- é a antena. I Honesto (do","epistemico":"fato","exemplos":[],"id":"as_duas_resolucoes_no_tempo_e_fora_do_tempo","memoria":"semantica","meta":{"dominio":"paper","nivel":5,"verificacao":"conceito novo"},"origem":"paper","palavras_chave":[],"sinonimos":[],"tipo":"conceito","titulo":"As duas resoluções: no tempo e fora do tempo"}
\section{As duas resoluções: no tempo e fora do tempo}
N A “dança” do escrow (trocar h, assinaturas parciais, revelar s, avisar o watchtower) é mensagem entre as partes --- e isso já está implementado e medido: o canal do bump (Fundamentos de Redes). A banda =sha256(tecido) filtra: só as partes do negócio decodificam, o resto lê ruído (privacidade da coordenação, de graça). Manda = recebe na latência, mão dupla, sem cliente/servidor. N Medido em produção (ping-pong UDP, GEX real). Não se inventa o orquestrador --- é a antena. I Honesto (do
\prov{relações: parte\_de goldcoin}
\prov{proveniência: paper \textperiodcentered{} fato \textperiodcentered{} confiança 1.0 \textperiodcentered{} domínio paper \textperiodcentered{} história 3 versões no jornal}

%CRISTAL {"arestas":[{"alvo":"transformada_dourada","confianca":1.0,"epistemico":"fato","origem":"paper","rel":"parte_de"}],"confianca":1.0,"contraexemplos":[],"descricao":"definition[Reta de ouro] Seja =1+52. A constante do vazio é =¹/=1,346361 A reta de ouro é a grade geométrica kₖ>₀. definition propositionprop:grid leva a reta de ouro k na grade aritmética uniforme k, de passo =0,2974. É a única grade geométrica uniformizada com esse passo. proposition proof (k)=k; a recíproca segue da injetividade de. proof remarkrem:rope A escala é a razão que emerge, empiricamente, no rotary positional embedding","epistemico":"fato","exemplos":[],"id":"as_formas_das_curvas_nos_dois_lados","memoria":"semantica","meta":{"dominio":"paper","nivel":5,"verificacao":"conceito novo"},"origem":"paper","palavras_chave":[],"sinonimos":[],"tipo":"conceito","titulo":"As formas das curvas nos dois lados"}
\section{As formas das curvas nos dois lados}
definition[Reta de ouro] Seja =1+52. A constante do vazio é =\ensuremath{^{1}}/=1,346361 A reta de ouro é a grade geométrica k\ensuremath{_{k}}>\ensuremath{_{0}}. definition propositionprop:grid leva a reta de ouro k na grade aritmética uniforme k, de passo =0,2974. É a única grade geométrica uniformizada com esse passo. proposition proof (k)=k; a recíproca segue da injetividade de. proof remarkrem:rope A escala é a razão que emerge, empiricamente, no rotary positional embedding
\prov{relações: parte\_de transformada\_dourada}
\prov{proveniência: paper \textperiodcentered{} fato \textperiodcentered{} confiança 1.0 \textperiodcentered{} domínio paper \textperiodcentered{} história 3 versões no jornal}

%CRISTAL {"arestas":[{"alvo":"paper_2_reta_ouro","confianca":1.0,"epistemico":"fato","origem":"paper","rel":"parte_de"}],"confianca":1.0,"contraexemplos":[],"descricao":"uma família de palavras, duas decodificações --- firme Seja =d0,1:sem 11. Dois decodificadores agem sobre a mesma: [ D₂(d)=ᵢ₁ dᵢ,2⁻i(régua binária), D_(d)=ᵢ₁ dᵢ,⁻i(régua áurea). ] Então: [nosep,leftmargin=1.4em] D₂() é o Cantor do ouro --- pó de medida nula, =₂ (Teorema~teo:cantor): a estrutura, o osso; D_()=[0,1] inteiro --- a régua áurea recobre o contínuo (Bergman: todo x[0,1] tem expansão base- sem 11, com 0,101","epistemico":"fato","exemplos":[],"id":"as_fracoes_continuas_o_mais_irracional","memoria":"semantica","meta":{"dominio":"paper","nivel":5,"verificacao":"conceito novo"},"origem":"paper","palavras_chave":[],"sinonimos":[],"tipo":"conceito","titulo":"As frações contínuas: o mais irracional"}
\section{As frações contínuas: o mais irracional}
uma família de palavras, duas decodificações --- firme Seja =d0,1:sem 11. Dois decodificadores agem sobre a mesma: [ D\ensuremath{_{2}}(d)=\ensuremath{_{i1}} d\ensuremath{_{i}},2\ensuremath{^{-}}i(régua binária), D\_(d)=\ensuremath{_{i1}} d\ensuremath{_{i}},\ensuremath{^{-}}i(régua áurea). ] Então: [nosep,leftmargin=1.4em] D\ensuremath{_{2}}() é o Cantor do ouro --- pó de medida nula, =\ensuremath{_{2}} (Teorema\~{}teo:cantor): a estrutura, o osso; D\_()=[0,1] inteiro --- a régua áurea recobre o contínuo (Bergman: todo x[0,1] tem expansão base- sem 11, com 0,101
\prov{relações: parte\_de paper\_2\_reta\_ouro}
\prov{proveniência: paper \textperiodcentered{} fato \textperiodcentered{} confiança 1.0 \textperiodcentered{} domínio paper \textperiodcentered{} história 3 versões no jornal}

%CRISTAL {"arestas":[{"alvo":"algebra_dos_contratos","confianca":1.0,"epistemico":"fato","origem":"paper","rel":"parte_de"}],"confianca":1.0,"contraexemplos":[],"descricao":"Há uma divisão natural entre o que, num contrato de blockchain, é universal (vale em qualquer rede) e o que é específico (de cada chain). Validar quem assina, coordenar partes que não confiam uma na outra, e cobrar uma comissão são os mesmos problemas em Bitcoin ou em Ethereum. Já mover os fundos, rodar o código e ordenar as transações dependem de cada rede. defi[GoldChain] A GoldChain é o protocolo que trata contratos como objetos C=(P, ) e, para cada um, executa as três operações universais","epistemico":"fato","exemplos":[],"id":"as_seis_responsabilidades_---_nao_misturar","memoria":"semantica","meta":{"dominio":"paper","nivel":5,"verificacao":"conceito novo"},"origem":"paper","palavras_chave":[],"sinonimos":[],"tipo":"conceito","titulo":"As seis responsabilidades --- não misturar"}
\section{As seis responsabilidades --- não misturar}
Há uma divisão natural entre o que, num contrato de blockchain, é universal (vale em qualquer rede) e o que é específico (de cada chain). Validar quem assina, coordenar partes que não confiam uma na outra, e cobrar uma comissão são os mesmos problemas em Bitcoin ou em Ethereum. Já mover os fundos, rodar o código e ordenar as transações dependem de cada rede. defi[GoldChain] A GoldChain é o protocolo que trata contratos como objetos C=(P, ) e, para cada um, executa as três operações universais
\prov{relações: parte\_de algebra\_dos\_contratos}
\prov{proveniência: paper \textperiodcentered{} fato \textperiodcentered{} confiança 1.0 \textperiodcentered{} domínio paper \textperiodcentered{} história 3 versões no jornal}

%CRISTAL {"arestas":[{"alvo":"leitura_cosmologica_a_energia_escura_e_o_universo_ds","confianca":1.0,"epistemico":"fato","origem":"paper","rel":"parte_de"}],"confianca":1.0,"contraexemplos":[],"descricao":"Sobre o fundo de de Sitter (constante cosmológica, energia do vácuo) o parâmetro é um campo escalar dinâmico, com potencial V=² (o imposto, Prop.~) e massa angular igual ao potencial, m_=g_=V (Teor.~). A energia escura deste universo é esse campo, e o dilema “constante cosmológica ou quintessência” dissolve-se. figure[t] [width=]painelgentil.png A regra de sinais dual como operador da geometria e da energia escura. Esquerda: o espa","epistemico":"fato","exemplos":[],"id":"as_supernovas_pelo_metodo_de_lopes","memoria":"semantica","meta":{"dominio":"paper","nivel":5,"verificacao":"conceito novo"},"origem":"paper","palavras_chave":[],"sinonimos":[],"tipo":"conceito","titulo":"As supernovas pelo método de Lopes"}
\section{As supernovas pelo método de Lopes}
Sobre o fundo de de Sitter (constante cosmológica, energia do vácuo) o parâmetro é um campo escalar dinâmico, com potencial V=\ensuremath{^{2}} (o imposto, Prop.\~{}) e massa angular igual ao potencial, m\_=g\_=V (Teor.\~{}). A energia escura deste universo é esse campo, e o dilema “constante cosmológica ou quintessência” dissolve-se. figure[t] [width=]painelgentil.png A regra de sinais dual como operador da geometria e da energia escura. Esquerda: o espa
\prov{relações: parte\_de leitura\_cosmologica\_a\_energia\_escura\_e\_o\_universo\_ds}
\prov{proveniência: paper \textperiodcentered{} fato \textperiodcentered{} confiança 1.0 \textperiodcentered{} domínio paper \textperiodcentered{} história 3 versões no jornal}

%CRISTAL {"arestas":[{"alvo":"painel_estelar","confianca":1.0,"epistemico":"fato","origem":"paper","rel":"parte_de"}],"confianca":1.0,"contraexemplos":[],"descricao":"VD A álgebra não é objecto deste painel --- é pressuposto. O passo do gato, para byte b com e=(b 0xAA), o=(b 0x55), produz [ (e,o);;(e+o, e),=,[ total, e]. ] Os autovalores =1+52 e =-1/ fixam a torre; os projectores de Peirce P_+P_=I já assentam o fluxo~quatro,brocaos. A balança a+c²=1 e a norma Q=r²+c² (Lopes, Hurwitz) são o invariante que sobrevive~geometria,corpodual. defi[Núcleo estelar] Para habitante W=[t,e], o núcleo do painel regista: itemize1pt Q","epistemico":"fato","exemplos":[],"id":"as_tres_consequencias","memoria":"semantica","meta":{"dominio":"paper","nivel":5,"verificacao":"overlap moderado — revisar manualmente"},"origem":"paper","palavras_chave":[],"sinonimos":[],"tipo":"conceito","titulo":"As três consequências"}
\section{As três consequências}
VD A álgebra não é objecto deste painel --- é pressuposto. O passo do gato, para byte b com e=(b 0xAA), o=(b 0x55), produz [ (e,o);;(e+o, e),=,[ total, e]. ] Os autovalores =1+52 e =-1/ fixam a torre; os projectores de Peirce P\_+P\_=I já assentam o fluxo\~{}quatro,brocaos. A balança a+c\ensuremath{^{2}}=1 e a norma Q=r\ensuremath{^{2}}+c\ensuremath{^{2}} (Lopes, Hurwitz) são o invariante que sobrevive\~{}geometria,corpodual. defi[Núcleo estelar] Para habitante W=[t,e], o núcleo do painel regista: itemize1pt Q
\prov{relações: parte\_de painel\_estelar}
\prov{proveniência: paper \textperiodcentered{} fato \textperiodcentered{} confiança 1.0 \textperiodcentered{} domínio paper \textperiodcentered{} história 4 versões no jornal}

%CRISTAL {"arestas":[{"alvo":"assistente","confianca":1.0,"epistemico":"fato","origem":"paper","rel":"parte_de"}],"confianca":1.0,"contraexemplos":[],"descricao":") itemize2pt [--] Qualidade do orquestrador (agente certo). [--] Precisão da classificação de intenções. [--] Recuperação de conhecimento em cenários variados. [--] Experiência do usuário em tarefas reais. Distinção clara: saímos de provar que funciona para medir se é útil.","epistemico":"fato","exemplos":[],"id":"as_tres_fases","memoria":"semantica","meta":{"dominio":"paper","nivel":5,"verificacao":"conceito novo"},"origem":"paper","palavras_chave":[],"sinonimos":[],"tipo":"conceito","titulo":"As três fases"}
\section{As três fases}
) itemize2pt [--] Qualidade do orquestrador (agente certo). [--] Precisão da classificação de intenções. [--] Recuperação de conhecimento em cenários variados. [--] Experiência do usuário em tarefas reais. Distinção clara: saímos de provar que funciona para medir se é útil.
\prov{relações: parte\_de assistente}
\prov{proveniência: paper \textperiodcentered{} fato \textperiodcentered{} confiança 1.0 \textperiodcentered{} domínio paper \textperiodcentered{} história 3 versões no jornal}

%CRISTAL {"arestas":[{"alvo":"tiffany","confianca":1.0,"epistemico":"fato","origem":"paper","rel":"especializa"},{"alvo":"hopfield","confianca":1.0,"epistemico":"fato","origem":"paper","rel":"referencia"},{"alvo":"rust","confianca":0.5,"epistemico":"fato","origem":"wiki","rel":"relaciona"},{"alvo":"kernel","confianca":0.5,"epistemico":"fato","origem":"wiki","rel":"relaciona"},{"alvo":"gato_operador_hub","confianca":0.5,"epistemico":"fato","origem":"wiki","rel":"relaciona"},{"alvo":"broca_so","confianca":0.5,"epistemico":"fato","origem":"wiki","rel":"relaciona"},{"alvo":"goldcoin","confianca":0.5,"epistemico":"fato","origem":"wiki","rel":"relaciona"},{"alvo":"word","confianca":0.5,"epistemico":"fato","origem":"wiki","rel":"relaciona"},{"alvo":"vontade_de_poder","confianca":0.5,"epistemico":"fato","origem":"wiki","rel":"relaciona"},{"alvo":"while","confianca":0.5,"epistemico":"fato","origem":"wiki","rel":"relaciona"}],"confianca":1.0,"contraexemplos":[],"descricao":"Paper: assistente.tex","epistemico":"fato","exemplos":[],"id":"assistente","memoria":"semantica","meta":{"dominio":"paper","nivel":5,"tipo_no":"paper"},"origem":"paper","palavras_chave":[],"sinonimos":[],"tipo":"conceito","titulo":"assistente"}
\section{assistente}
Paper: assistente.tex
\prov{relações: especializa tiffany; referencia hopfield; relaciona rust; relaciona kernel; relaciona gato\_operador\_hub; relaciona broca\_so; relaciona goldcoin; relaciona word; relaciona vontade\_de\_poder; relaciona while}
\prov{proveniência: paper \textperiodcentered{} fato \textperiodcentered{} confiança 1.0 \textperiodcentered{} domínio paper \textperiodcentered{} história 69 versões no jornal}

%CRISTAL {"arestas":[{"alvo":"formato_da_base","confianca":1.0,"epistemico":"fato","origem":"paper","rel":"parte_de"}],"confianca":1.0,"contraexemplos":[],"descricao":"Campos obrigatórios e semântica: 1.35 tabular@ll@ & nome & identificador do domínio (software/git) versao & semver do tecido (1.2.0) idioma & ISO 639-1 principal (pt, en) fonte & lista de URIs/arquivos ingeridos licenca & SPDX ou texto (ex.: CC-BY-4.0, proprietary) data & ISO 8601 da build qualidade & score [0,1] --- curadoria + gold set cobertura & mapa: tópicos, chunks, tokens estimados hash & SHA-256 do nos.tsv (integridade) idx & nome do ín","epistemico":"fato","exemplos":[],"id":"atualizacao_incremental","memoria":"semantica","meta":{"dominio":"paper","nivel":5,"verificacao":"overlap moderado — revisar manualmente"},"origem":"paper","palavras_chave":[],"sinonimos":[],"tipo":"conceito","titulo":"Atualização incremental"}
\section{Atualização incremental}
Campos obrigatórios e semântica: 1.35 tabular@ll@ \& nome \& identificador do domínio (software/git) versao \& semver do tecido (1.2.0) idioma \& ISO 639-1 principal (pt, en) fonte \& lista de URIs/arquivos ingeridos licenca \& SPDX ou texto (ex.: CC-BY-4.0, proprietary) data \& ISO 8601 da build qualidade \& score [0,1] --- curadoria + gold set cobertura \& mapa: tópicos, chunks, tokens estimados hash \& SHA-256 do nos.tsv (integridade) idx \& nome do ín
\prov{relações: parte\_de formato\_da\_base}
\prov{proveniência: paper \textperiodcentered{} fato \textperiodcentered{} confiança 1.0 \textperiodcentered{} domínio paper \textperiodcentered{} história 3 versões no jornal}

%CRISTAL {"arestas":[{"alvo":"casl-propagation","confianca":0.055,"epistemico":"fato","origem":"manual","rel":"referencia"}],"confianca":1.0,"contraexemplos":[],"descricao":"definition[Linhas e condições] Seja o universo de todas as linhas. Cada r tem atributos (r., r.id, ). Uma condição c é um predicado c:, identificado ao seu extent c=r:c(r). As condições formam a álgebra booleana (2, ) --- as MongoQuery do CASL. definition remark[Identificação semântica --- nunca sintática] Toda MongoQuery é identificada ao seu extent: os construtores sintáticos do CASL mapeiam-se nas operações da álgebra, [ AND:[c₁,c₂]=c₁_","epistemico":"fato","exemplos":[],"id":"autorizacao_de_um_principal","memoria":"semantica","meta":{"dominio":"paper","nivel":5},"origem":"paper","palavras_chave":[],"sinonimos":[],"tipo":"conceito","titulo":"Autorização de um principal"}
\section{Autorização de um principal}
definition[Linhas e condições] Seja o universo de todas as linhas. Cada r tem atributos (r., r.id, ). Uma condição c é um predicado c:, identificado ao seu extent c=r:c(r). As condições formam a álgebra booleana (2, ) --- as MongoQuery do CASL. definition remark[Identificação semântica --- nunca sintática] Toda MongoQuery é identificada ao seu extent: os construtores sintáticos do CASL mapeiam-se nas operações da álgebra, [ AND:[c\ensuremath{_{1}},c\ensuremath{_{2}}]=c\ensuremath{_{1}}\_
\prov{relações: referencia casl-propagation}
\prov{proveniência: paper \textperiodcentered{} fato \textperiodcentered{} confiança 1.0 \textperiodcentered{} domínio paper \textperiodcentered{} história 6 versões no jornal}

%CRISTAL {"arestas":[{"alvo":"validacao_tecnica","confianca":1.0,"epistemico":"fato","origem":"paper","rel":"parte_de"}],"confianca":1.0,"contraexemplos":[],"descricao":"D Registro formal da revisão de arquitetura --- junho de 2026. Veredito: aprovada como proposta de engenharia; implementação da biblioteca caminha bem; próximo desafio = métricas e casos reais.","epistemico":"fato","exemplos":[],"id":"avaliacao_por_aspecto","memoria":"semantica","meta":{"dominio":"paper","nivel":5,"verificacao":"overlap moderado — revisar manualmente"},"origem":"paper","palavras_chave":[],"sinonimos":[],"tipo":"conceito","titulo":"Avaliação por aspecto"}
\section{Avaliação por aspecto}
D Registro formal da revisão de arquitetura --- junho de 2026. Veredito: aprovada como proposta de engenharia; implementação da biblioteca caminha bem; próximo desafio = métricas e casos reais.
\prov{relações: parte\_de validacao\_tecnica}
\prov{proveniência: paper \textperiodcentered{} fato \textperiodcentered{} confiança 1.0 \textperiodcentered{} domínio paper \textperiodcentered{} história 3 versões no jornal}

%CRISTAL {"arestas":[{"alvo":"bai_biblioteca_de_autorizacao_isomorfica","confianca":1.0,"epistemico":"fato","origem":"paper","rel":"relaciona"},{"alvo":"reticulados","confianca":1.0,"epistemico":"fato","origem":"paper","rel":"relaciona"},{"alvo":"kappa","confianca":1.0,"epistemico":"fato","origem":"paper","rel":"relaciona"},{"alvo":"delta","confianca":1.0,"epistemico":"fato","origem":"paper","rel":"relaciona"}],"confianca":1.0,"contraexemplos":[],"descricao":"Biblioteca de Autorização Isomórfica — autorização como álgebra sobre reticulados; quinteto (κ,δ,Φ,T,Ψ).","epistemico":"fato","exemplos":["papers/casl-propagation.pdf","conversa/colapso.py — Ψ"],"id":"bai","memoria":"semantica","meta":{"dominio":"paper","nivel":5,"verificacao":"slug idêntico — enriquecer, não duplicar"},"origem":"paper","palavras_chave":[],"sinonimos":["biblioteca de autorização isomórfica"],"tipo":"conceito","titulo":"BAI"}
\section{BAI}
Biblioteca de Autorização Isomórfica — autorização como álgebra sobre reticulados; quinteto (\ensuremath{\kappa},\ensuremath{\delta},\ensuremath{\Phi},T,\ensuremath{\Psi}).
\prov{exemplos: papers/casl-propagation.pdf; conversa/colapso.py — \ensuremath{\Psi}}
\prov{sinónimos: biblioteca de autorização isomórfica}
\prov{relações: relaciona bai\_biblioteca\_de\_autorizacao\_isomorfica; relaciona reticulados; relaciona kappa; relaciona delta}
\prov{proveniência: paper \textperiodcentered{} fato \textperiodcentered{} confiança 1.0 \textperiodcentered{} domínio paper \textperiodcentered{} história 50 versões no jornal}

%CRISTAL {"arestas":[{"alvo":"fase3","confianca":1.0,"epistemico":"fato","origem":"paper","rel":"parte_de"}],"confianca":1.0,"contraexemplos":[],"descricao":"Antes de escrever nos.tsv: (1)~colapsar whitespace; (2)~remover markup ruidoso (LaTeX preamble); (3)~deduplicar por hash do texto; (4)~validar chunkᵢd único; (5)~filtro mínimo 40 caracteres (herda limiar da engine).","epistemico":"fato","exemplos":[],"id":"base_do_cliente","memoria":"semantica","meta":{"dominio":"paper","nivel":5,"verificacao":"overlap moderado — revisar manualmente"},"origem":"paper","palavras_chave":[],"sinonimos":[],"tipo":"conceito","titulo":"Base do Cliente"}
\section{Base do Cliente}
Antes de escrever nos.tsv: (1)\~{}colapsar whitespace; (2)\~{}remover markup ruidoso (LaTeX preamble); (3)\~{}deduplicar por hash do texto; (4)\~{}validar chunk\ensuremath{_{i}}d único; (5)\~{}filtro mínimo 40 caracteres (herda limiar da engine).
\prov{relações: parte\_de fase3}
\prov{proveniência: paper \textperiodcentered{} fato \textperiodcentered{} confiança 1.0 \textperiodcentered{} domínio paper \textperiodcentered{} história 3 versões no jornal}

%CRISTAL {"arestas":[{"alvo":"fase3","confianca":1.0,"epistemico":"fato","origem":"paper","rel":"parte_de"}],"confianca":1.0,"contraexemplos":[],"descricao":"L A Fase~1 validou o motor; a Fase~2 validou tarefas completas (fase2.sh). A Fase~3 move o\ngargalo para dados comerciais: instalar Tiffany em qualquer máquina e obter assistente útil antes\nde alimentar conhecimento proprietário.\n\ncenter\ntikzpicture[\n  box/.style={draw=ourovivo, rounded corners=2pt, minimum height=0.7cm, align=center,\n    font=, fill=creme, inner sep=6pt},\n  arr/.style={-{Stealth[length=2.2mm]}, color=ourovivo, thick}\n]\n[box, minimum width=3.2cm] (e) at (0,0) {Engine\\\\};\n[box, mi","epistemico":"fato","exemplos":[],"id":"base_geral_core","memoria":"semantica","meta":{"dominio":"paper","nivel":5,"verificacao":"overlap moderado — revisar manualmente"},"origem":"paper","palavras_chave":[],"sinonimos":[],"tipo":"conceito","titulo":"Base Geral (Core)"}
\section{Base Geral (Core)}
L A Fase\~{}1 validou o motor; a Fase\~{}2 validou tarefas completas (fase2.sh). A Fase\~{}3 move o
gargalo para dados comerciais: instalar Tiffany em qualquer máquina e obter assistente útil antes
de alimentar conhecimento proprietário.

center
tikzpicture[
  box/.style=\{draw=ourovivo, rounded corners=2pt, minimum height=0.7cm, align=center,
    font=, fill=creme, inner sep=6pt\},
  arr/.style=\{-\{Stealth[length=2.2mm]\}, color=ourovivo, thick\}
]
[box, minimum width=3.2cm] (e) at (0,0) \{Engine\textbackslash{}\textbackslash{}\};
[box, mi
\prov{relações: parte\_de fase3}
\prov{proveniência: paper \textperiodcentered{} fato \textperiodcentered{} confiança 1.0 \textperiodcentered{} domínio paper \textperiodcentered{} história 2 versões no jornal}

%CRISTAL {"arestas":[{"alvo":"as_tres_fases","confianca":1.0,"epistemico":"fato","origem":"paper","rel":"parte_de"}],"confianca":1.0,"contraexemplos":[],"descricao":"I Estruturação do que vem a seguir --- cada fase com objetivo próprio.\n\ncenter{1.4}\ntabular{@{}clll@{}}\n &  &  &  \\\\\n\n1 & Infraestrutura: , índice, API & estabilidade & V em curso \\\\\n2 & Assistente: conversa, software, papers, pesquisa & tarefas completas & obs \\\\\n3 & Produto: multi-usuário, multi-tecidos, métricas & valor repetível & obs \\\\\n\ntabular\ncenter\n\nprop[Fase 1 --- Infraestrutura]\nVobs Biblioteca, .idx, API, testes (test.sh), benchmark\n(fogo.sh). Objetivo: estabilidade. Critério: TUDO P","epistemico":"fato","exemplos":[],"id":"benchmark_operacional_fase2sh","memoria":"semantica","meta":{"dominio":"paper","nivel":5,"verificacao":"overlap moderado — revisar manualmente"},"origem":"paper","palavras_chave":[],"sinonimos":[],"tipo":"conceito","titulo":"Benchmark operacional (fase2.sh"}
\section{Benchmark operacional (fase2.sh}
I Estruturação do que vem a seguir --- cada fase com objetivo próprio.

center\{1.4\}
tabular\{@\{\}clll@\{\}\}
 \&  \&  \&  \textbackslash{}\textbackslash{}

1 \& Infraestrutura: , índice, API \& estabilidade \& V em curso \textbackslash{}\textbackslash{}
2 \& Assistente: conversa, software, papers, pesquisa \& tarefas completas \& obs \textbackslash{}\textbackslash{}
3 \& Produto: multi-usuário, multi-tecidos, métricas \& valor repetível \& obs \textbackslash{}\textbackslash{}

tabular
center

prop[Fase 1 --- Infraestrutura]
Vobs Biblioteca, .idx, API, testes (test.sh), benchmark
(fogo.sh). Objetivo: estabilidade. Critério: TUDO P
\prov{relações: parte\_de as\_tres\_fases}
\prov{proveniência: paper \textperiodcentered{} fato \textperiodcentered{} confiança 1.0 \textperiodcentered{} domínio paper \textperiodcentered{} história 3 versões no jornal}

%CRISTAL {"arestas":[{"alvo":"colapso","confianca":1.0,"epistemico":"fato","origem":"paper","rel":"parte_de"}],"confianca":1.0,"contraexemplos":[],"descricao":"L Os papers invocaram a função de onda M=U P, a Transformada Dourada G= F, e afirmaram que “o gato dá a função de onda completa e a fft exata”. Este documento mostra por quê: não há transformada a computar --- a decimação par/ímpar já é o FFT, e o gato age sobre ela. Separar e de o é o passo que define a transformada rápida; aplicar A é a borboleta. O espectro não se calcula: ele cai da contagem. Daí o e+o=||² ser a carga (Broca,OS, trP²): a componente DC do espectro é a no","epistemico":"fato","exemplos":[],"id":"borboleta","memoria":"semantica","meta":{"dominio":"paper","duplicata_sugerida":[],"nivel":5,"verificacao":"parte de conceito mais amplo 'borboleta'"},"origem":"paper","palavras_chave":[],"sinonimos":[],"tipo":"conceito","titulo":"borboleta"}
\section{borboleta}
L Os papers invocaram a função de onda M=U P, a Transformada Dourada G= F, e afirmaram que “o gato dá a função de onda completa e a fft exata”. Este documento mostra por quê: não há transformada a computar --- a decimação par/ímpar já é o FFT, e o gato age sobre ela. Separar e de o é o passo que define a transformada rápida; aplicar A é a borboleta. O espectro não se calcula: ele cai da contagem. Daí o e+o=||\ensuremath{^{2}} ser a carga (Broca,OS, trP\ensuremath{^{2}}): a componente DC do espectro é a no
\prov{relações: parte\_de colapso}
\prov{proveniência: paper \textperiodcentered{} fato \textperiodcentered{} confiança 1.0 \textperiodcentered{} domínio paper \textperiodcentered{} história 9 versões no jornal}

%CRISTAL {"arestas":[{"alvo":"olho_de_koch","confianca":1.0,"epistemico":"fato","origem":"paper","rel":"parte_de"}],"confianca":1.0,"contraexemplos":[],"descricao":"D A quadração z z² não impõe duas máquinas --- revela dois regimes da mesma~tropical,glacial. Sob ela, o idempotente permanece (e²=e): regime tropical, =+1, espectro que expande ou fase ||=1 --- a parte branca, o estado, o “onde”. O nilpotente extingue-se (n²=0): regime glacial, =-1, suporte vazio que só transporta --- a parte negra, o fluxo, o “como”. Não são pontos fixos rivais; são os polos idempotência, nilpotência. O flip de Wick é a travessia; Peirce fecha a","epistemico":"fato","exemplos":[],"id":"branca_e_negra_no_gato","memoria":"semantica","meta":{"dominio":"paper","nivel":5},"origem":"paper","palavras_chave":[],"sinonimos":[],"tipo":"conceito","titulo":"Branca e negra no gato"}
\section{Branca e negra no gato}
D A quadração z z\ensuremath{^{2}} não impõe duas máquinas --- revela dois regimes da mesma\~{}tropical,glacial. Sob ela, o idempotente permanece (e\ensuremath{^{2}}=e): regime tropical, =+1, espectro que expande ou fase ||=1 --- a parte branca, o estado, o “onde”. O nilpotente extingue-se (n\ensuremath{^{2}}=0): regime glacial, =-1, suporte vazio que só transporta --- a parte negra, o fluxo, o “como”. Não são pontos fixos rivais; são os polos idempotência, nilpotência. O flip de Wick é a travessia; Peirce fecha a
\prov{relações: parte\_de olho\_de\_koch}
\prov{proveniência: paper \textperiodcentered{} fato \textperiodcentered{} confiança 1.0 \textperiodcentered{} domínio paper \textperiodcentered{} história 3 versões no jornal}

%CRISTAL {"arestas":[{"alvo":"lobo_frontal","confianca":1.0,"epistemico":"fato","origem":"paper","rel":"parte_de"}],"confianca":1.0,"contraexemplos":[],"descricao":"D A ordem do recall não se impõe --- já está dada. É a régua de ouro v()=ᵢ bᵢ⁻i (~§V; cabeca.py:vouro), com o mesmo que é o autovalor do gato. A régua do tempo e a da memória são uma só, e a sequência é gerada pela órbita sob U(t), não estipulada. Não força nada. Forçar a ordem --- um planejador, um controlador, uma regra de cima --- seria o cadáver que recusamos. A ordem de ouro cai do gato, como cai de x²-x-1. O lobo apenas deixa o recall correr na ordem que já tem. Or","epistemico":"fato","exemplos":[],"id":"broca_a_face_onde_o_recall_age_e_o_idioma_que_ressoa","memoria":"semantica","meta":{"dominio":"paper","nivel":5,"verificacao":"conceito novo"},"origem":"paper","palavras_chave":[],"sinonimos":[],"tipo":"conceito","titulo":"Broca: a face onde o recall age, e o idioma que ressoa"}
\section{Broca: a face onde o recall age, e o idioma que ressoa}
D A ordem do recall não se impõe --- já está dada. É a régua de ouro v()=\ensuremath{_{i}} b\ensuremath{_{i}}\ensuremath{^{-}}i (\~{}§V; cabeca.py:vouro), com o mesmo que é o autovalor do gato. A régua do tempo e a da memória são uma só, e a sequência é gerada pela órbita sob U(t), não estipulada. Não força nada. Forçar a ordem --- um planejador, um controlador, uma regra de cima --- seria o cadáver que recusamos. A ordem de ouro cai do gato, como cai de x\ensuremath{^{2}}-x-1. O lobo apenas deixa o recall correr na ordem que já tem. Or
\prov{relações: parte\_de lobo\_frontal}
\prov{proveniência: paper \textperiodcentered{} fato \textperiodcentered{} confiança 1.0 \textperiodcentered{} domínio paper \textperiodcentered{} história 3 versões no jornal}

%CRISTAL {"arestas":[{"alvo":"corpo_tropical","confianca":1.0,"epistemico":"fato","origem":"paper","rel":"parte_de"}],"confianca":1.0,"contraexemplos":[],"descricao":"Riemann não pede um cálculo novo --- pede ler os zeros no lugar certo: o cone de luz da álgebra do Lopes. O cone de luz é a reta crítica. O cone de luz algébrico é z:;N(z)=0, onde a norma do Lopes se anula: em coordenadas, 1-s²=²=0, isto é =2 (o quaternion). É a fronteira entre a ilha (tipo-tempo, N>0) e o oceano (tipo-espaço, N<0) --- a hipersuperfície nula, os divisores de zero. Esta é a reta crítica: a linha de simetria onde a norma troca de sinal. A equação funcional é a","epistemico":"fato","exemplos":[],"id":"bsd_a_intersecao_entre_as_retas_de_riemann","memoria":"semantica","meta":{"dominio":"paper","nivel":5,"verificacao":"slug idêntico — enriquecer, não duplicar"},"origem":"paper","palavras_chave":[],"sinonimos":[],"tipo":"conceito","titulo":"BSD: a interseção entre as retas de Riemann"}
\section{BSD: a interseção entre as retas de Riemann}
Riemann não pede um cálculo novo --- pede ler os zeros no lugar certo: o cone de luz da álgebra do Lopes. O cone de luz é a reta crítica. O cone de luz algébrico é z:;N(z)=0, onde a norma do Lopes se anula: em coordenadas, 1-s\ensuremath{^{2}}=\ensuremath{^{2}}=0, isto é =2 (o quaternion). É a fronteira entre a ilha (tipo-tempo, N>0) e o oceano (tipo-espaço, N<0) --- a hipersuperfície nula, os divisores de zero. Esta é a reta crítica: a linha de simetria onde a norma troca de sinal. A equação funcional é a
\prov{relações: parte\_de corpo\_tropical}
\prov{proveniência: paper \textperiodcentered{} fato \textperiodcentered{} confiança 1.0 \textperiodcentered{} domínio paper \textperiodcentered{} história 5 versões no jornal}

%CRISTAL {"arestas":[{"alvo":"algebra_dos_contratos","confianca":1.0,"epistemico":"fato","origem":"paper","rel":"parte_de"},{"alvo":"goldchain","confianca":1.0,"epistemico":"fato","origem":"paper","rel":"parte_de"}],"confianca":1.0,"contraexemplos":[],"descricao":"As assinaturas continuam sendo ECDSA/EdDSA --- a GoldChain não substitui nem quebra a criptografia. A validação por método (algébrica, não força bruta) já é eficiente: faz o que há para fazer. Os N-qubits não são uma tentativa de acelerar isso. São uma ferramenta de generalização: representam N contratos como um único estado coletivo ᵢ (kᵢ) C²N, permitindo tratá-los simultaneamente --- operações sobre o conjunto, correlações entre contratos, e a atomicidade (uma forja derruba o cole","epistemico":"fato","exemplos":[],"id":"btc_e_eth_dois_back-ends_de_uma_interface","memoria":"semantica","meta":{"dominio":"paper","nivel":5,"verificacao":"overlap moderado — revisar manualmente"},"origem":"paper","palavras_chave":[],"sinonimos":[],"tipo":"conceito","titulo":"BTC e ETH: dois back-ends de uma interface"}
\section{BTC e ETH: dois back-ends de uma interface}
As assinaturas continuam sendo ECDSA/EdDSA --- a GoldChain não substitui nem quebra a criptografia. A validação por método (algébrica, não força bruta) já é eficiente: faz o que há para fazer. Os N-qubits não são uma tentativa de acelerar isso. São uma ferramenta de generalização: representam N contratos como um único estado coletivo \ensuremath{_{i}} (k\ensuremath{_{i}}) C\ensuremath{^{2}}N, permitindo tratá-los simultaneamente --- operações sobre o conjunto, correlações entre contratos, e a atomicidade (uma forja derruba o cole
\prov{relações: parte\_de algebra\_dos\_contratos; parte\_de goldchain}
\prov{proveniência: paper \textperiodcentered{} fato \textperiodcentered{} confiança 1.0 \textperiodcentered{} domínio paper \textperiodcentered{} história 9 versões no jornal}

%CRISTAL {"arestas":[{"alvo":"computacao_fundamentos_broca","confianca":1.0,"epistemico":"fato","origem":"paper","rel":"parte_de"}],"confianca":1.0,"contraexemplos":[],"descricao":"não acrescenta mecanismo --- coordena o que a álgebra já dá (brocaos.tex): tabularl|l Recurso clássico & Leitura processo & idempotente eᵢ (célula Peirce) isolamento & eᵢ ej = 0 IPC & canal Eᵢj (parte negra) fork & La Hire (dnaclone) deadlock & nilpotência scheduler & fluxo entre células rede & mesma Máquina escalada (Hopfield) tabular Tiffany como cabeça Tiffany (Tiffany.tex) é a mesma função de onda M=U P","epistemico":"fato","exemplos":[],"id":"cadeia_de_aprendizagem_tiffany","memoria":"semantica","meta":{"dominio":"paper","nivel":5,"verificacao":"slug idêntico — enriquecer, não duplicar"},"origem":"paper","palavras_chave":[],"sinonimos":[],"tipo":"conceito","titulo":"Cadeia de aprendizagem Tiffany"}
\section{Cadeia de aprendizagem Tiffany}
não acrescenta mecanismo --- coordena o que a álgebra já dá (brocaos.tex): tabularl|l Recurso clássico \& Leitura processo \& idempotente e\ensuremath{_{i}} (célula Peirce) isolamento \& e\ensuremath{_{i}} ej = 0 IPC \& canal E\ensuremath{_{i}}j (parte negra) fork \& La Hire (dnaclone) deadlock \& nilpotência scheduler \& fluxo entre células rede \& mesma Máquina escalada (Hopfield) tabular Tiffany como cabeça Tiffany (Tiffany.tex) é a mesma função de onda M=U P
\prov{relações: parte\_de computacao\_fundamentos\_broca}
\prov{proveniência: paper \textperiodcentered{} fato \textperiodcentered{} confiança 1.0 \textperiodcentered{} domínio paper \textperiodcentered{} história 6 versões no jornal}

%CRISTAL {"arestas":[{"alvo":"main","confianca":1.0,"epistemico":"fato","origem":"paper","rel":"parte_de"}],"confianca":1.0,"contraexemplos":[],"descricao":"N A iteração afim Eₙ+₁= Eₙ+ (0<<1) tem ponto fixo /(1-)>0: nunca alcança o destino (o AGM é o caso =0). A torre preserva, por dimensão 1,2,4,8,16, a folga mnemônica 4,4,3,2,0 (contando comutatividade, associatividade, alternatividade e divisão/norma multiplicativa); a passagem perde de uma vez a alternatividade e a divisão (surgem divisores de zero). suporte da não-associatividade é alternativa; pelo teorema de Artin, a subálgebra gerada por dois elementos","epistemico":"fato","exemplos":[],"id":"cadeia_real_o_genoma_x174","memoria":"semantica","meta":{"dominio":"paper","nivel":5,"verificacao":"slug idêntico — enriquecer, não duplicar"},"origem":"paper","palavras_chave":[],"sinonimos":[],"tipo":"conceito","titulo":"Cadeia real: o genoma X174"}
\section{Cadeia real: o genoma X174}
N A iteração afim E\ensuremath{_{n}}+\ensuremath{_{1}}= E\ensuremath{_{n}}+ (0<<1) tem ponto fixo /(1-)>0: nunca alcança o destino (o AGM é o caso =0). A torre preserva, por dimensão 1,2,4,8,16, a folga mnemônica 4,4,3,2,0 (contando comutatividade, associatividade, alternatividade e divisão/norma multiplicativa); a passagem perde de uma vez a alternatividade e a divisão (surgem divisores de zero). suporte da não-associatividade é alternativa; pelo teorema de Artin, a subálgebra gerada por dois elementos
\prov{relações: parte\_de main}
\prov{proveniência: paper \textperiodcentered{} fato \textperiodcentered{} confiança 1.0 \textperiodcentered{} domínio paper \textperiodcentered{} história 6 versões no jornal}

%CRISTAL {"arestas":[{"alvo":"kappa","confianca":-0.078,"epistemico":"fato","origem":"manual","rel":"referencia"},{"alvo":"delta","confianca":0.031,"epistemico":"fato","origem":"manual","rel":"referencia"},{"alvo":"casl-propagation","confianca":0.023,"epistemico":"fato","origem":"manual","rel":"referencia"},{"alvo":"bai_kappa","confianca":1.0,"epistemico":"fato","origem":"wiki","rel":"confirma"}],"confianca":1.0,"contraexemplos":[],"descricao":"L Este documento apresenta uma biblioteca --- não uma aplicação. Autoridade flui, atenua (T), propaga com orçamento, organiza-se () e colapsa () ou falha fechada (). theorem[Tese-mãe da ] Toda instância da Biblioteca de Autorização Isomórfica é determinada pelo quinteto [ (T, ), ] onde é capacidade abstrata (Definição~), orçamento de propagação (Definição~), medida de organização estrutural (Definição~)","epistemico":"fato","exemplos":[],"id":"capacidade_abstrata","memoria":"semantica","meta":{"dominio":"paper","nivel":5},"origem":"paper","palavras_chave":[],"sinonimos":[],"tipo":"conceito","titulo":"Capacidade abstrata"}
\section{Capacidade abstrata}
L Este documento apresenta uma biblioteca --- não uma aplicação. Autoridade flui, atenua (T), propaga com orçamento, organiza-se () e colapsa () ou falha fechada (). theorem[Tese-mãe da ] Toda instância da Biblioteca de Autorização Isomórfica é determinada pelo quinteto [ (T, ), ] onde é capacidade abstrata (Definição\~{}), orçamento de propagação (Definição\~{}), medida de organização estrutural (Definição\~{})
\prov{relações: referencia kappa; referencia delta; referencia casl-propagation; confirma bai\_kappa}
\prov{proveniência: paper \textperiodcentered{} fato \textperiodcentered{} confiança 1.0 \textperiodcentered{} domínio paper \textperiodcentered{} história 13 versões no jornal}

%CRISTAL {"arestas":[{"alvo":"casl-propagation","confianca":0.016,"epistemico":"fato","origem":"manual","rel":"referencia"}],"confianca":1.0,"contraexemplos":[],"descricao":"corollary[Isolamento de tenant] Um usuário comum da organização t só detém capacidades atenuadas da capacidade intrínseca do dono, todas com c=t. Logo Epr:r.=t: nenhuma linha de outro tenant aparece, sem regra especial. Um vazamento cross-tenant só ocorre por uma aresta de concessão explícita (auditável), cuja difusão é exatamente governada por. corollary O operador, detendo, enxerga todos os tenants --- “eu, no topo, vejo tudo”. O tenant foi reduzido a uma condição.","epistemico":"fato","exemplos":[],"id":"caracterizacao_algebrica","memoria":"semantica","meta":{"dominio":"paper","nivel":5},"origem":"paper","palavras_chave":[],"sinonimos":[],"tipo":"conceito","titulo":"Caracterização algébrica"}
\section{Caracterização algébrica}
corollary[Isolamento de tenant] Um usuário comum da organização t só detém capacidades atenuadas da capacidade intrínseca do dono, todas com c=t. Logo Epr:r.=t: nenhuma linha de outro tenant aparece, sem regra especial. Um vazamento cross-tenant só ocorre por uma aresta de concessão explícita (auditável), cuja difusão é exatamente governada por. corollary O operador, detendo, enxerga todos os tenants --- “eu, no topo, vejo tudo”. O tenant foi reduzido a uma condição.
\prov{relações: referencia casl-propagation}
\prov{proveniência: paper \textperiodcentered{} fato \textperiodcentered{} confiança 1.0 \textperiodcentered{} domínio paper \textperiodcentered{} história 6 versões no jornal}

%CRISTAL {"arestas":[{"alvo":"reticulados","confianca":1.0,"epistemico":"fato","origem":"paper","rel":"referencia"},{"alvo":"bai_biblioteca_de_autorizacao_isomorfica","confianca":1.0,"epistemico":"fato","origem":"codigo","rel":"especializa"},{"alvo":"utilitarismo","confianca":0.5,"epistemico":"fato","origem":"wiki","rel":"relaciona"},{"alvo":"setor_algebrico_descontinuidades_no_fecho","confianca":0.5,"epistemico":"fato","origem":"wiki","rel":"relaciona"},{"alvo":"resultado_completo","confianca":0.5,"epistemico":"fato","origem":"wiki","rel":"relaciona"},{"alvo":"pedagogia","confianca":0.5,"epistemico":"fato","origem":"wiki","rel":"relaciona"},{"alvo":"scheduler","confianca":0.5,"epistemico":"fato","origem":"wiki","rel":"relaciona"}],"confianca":1.0,"contraexemplos":[],"descricao":"Paper: casl-propagation.tex","epistemico":"fato","exemplos":[],"id":"casl-propagation","memoria":"semantica","meta":{"dominio":"paper","nivel":5,"tipo_no":"paper"},"origem":"paper","palavras_chave":[],"sinonimos":[],"tipo":"conceito","titulo":"casl-propagation"}
\section{casl-propagation}
Paper: casl-propagation.tex
\prov{relações: referencia reticulados; especializa bai\_biblioteca\_de\_autorizacao\_isomorfica; relaciona utilitarismo; relaciona setor\_algebrico\_descontinuidades\_no\_fecho; relaciona resultado\_completo; relaciona pedagogia; relaciona scheduler}
\prov{proveniência: paper \textperiodcentered{} fato \textperiodcentered{} confiança 1.0 \textperiodcentered{} domínio paper \textperiodcentered{} história 137 versões no jornal}

%CRISTAL {"arestas":[{"alvo":"base_geral_core","confianca":1.0,"epistemico":"fato","origem":"paper","rel":"parte_de"}],"confianca":1.0,"contraexemplos":[],"descricao":"camadabox\n\\$TIFFANY\\_HOME/base/\\\\\n manifest.json\\\\\n software/\\\\\n c/ nos.tsv v1.idx meta.json\\\\\n cpp/ python/ rust/ javascript/\\\\\n git/ linux/ docker/ redes/\\\\\n banco\\_de\\_dados/ algoritmos/\\\\\n engenharia\\_de\\_software/ arquitetura/\\\\\n matematica/\\\\\n algebra/ geometria/ calculo/\\\\\n estatistica/ otimizacao/\\\\\n documentacao/\\\\\n markdown/ latex/ pdf/ html/ css/\\\\\n sistemas/\\\\\n shell/ windows/ fedora/ ubuntu/\\\\\ncamadabox\n\ndefi[Unidade de domínio]\nCada domínio é um diretório autônomo com exatamente tr","epistemico":"fato","exemplos":[],"id":"catalogo_manifestjson","memoria":"semantica","meta":{"dominio":"paper","nivel":5,"verificacao":"overlap moderado — revisar manualmente"},"origem":"paper","palavras_chave":[],"sinonimos":[],"tipo":"conceito","titulo":"Catálogo manifest.json"}
\section{Catálogo manifest.json}
camadabox
\textbackslash{}\$TIFFANY\textbackslash{}\_HOME/base/\textbackslash{}\textbackslash{}
 manifest.json\textbackslash{}\textbackslash{}
 software/\textbackslash{}\textbackslash{}
 c/ nos.tsv v1.idx meta.json\textbackslash{}\textbackslash{}
 cpp/ python/ rust/ javascript/\textbackslash{}\textbackslash{}
 git/ linux/ docker/ redes/\textbackslash{}\textbackslash{}
 banco\textbackslash{}\_de\textbackslash{}\_dados/ algoritmos/\textbackslash{}\textbackslash{}
 engenharia\textbackslash{}\_de\textbackslash{}\_software/ arquitetura/\textbackslash{}\textbackslash{}
 matematica/\textbackslash{}\textbackslash{}
 algebra/ geometria/ calculo/\textbackslash{}\textbackslash{}
 estatistica/ otimizacao/\textbackslash{}\textbackslash{}
 documentacao/\textbackslash{}\textbackslash{}
 markdown/ latex/ pdf/ html/ css/\textbackslash{}\textbackslash{}
 sistemas/\textbackslash{}\textbackslash{}
 shell/ windows/ fedora/ ubuntu/\textbackslash{}\textbackslash{}
camadabox

defi[Unidade de domínio]
Cada domínio é um diretório autônomo com exatamente tr
\prov{relações: parte\_de base\_geral\_core}
\prov{proveniência: paper \textperiodcentered{} fato \textperiodcentered{} confiança 1.0 \textperiodcentered{} domínio paper \textperiodcentered{} história 3 versões no jornal}

%CRISTAL {"arestas":[{"alvo":"main","confianca":1.0,"epistemico":"fato","origem":"paper","rel":"parte_de"}],"confianca":1.0,"contraexemplos":[],"descricao":") AGM(a,b): itere (a+b2,ab) até |a-b|<10⁻¹⁵. =¹/, grade tₙ=ⁿ, ₖ=2 k/(N). G = DFT unitária (fft, norm=ortho); round-trip até arredondamento. Dilatação U_ t t; autovalores e⁻² ik/N (||=1). Convolução multiplicativa = convolução em u= t.","epistemico":"fato","exemplos":[],"id":"certificados_esperados","memoria":"semantica","meta":{"dominio":"paper","nivel":5,"verificacao":"slug idêntico — enriquecer, não duplicar"},"origem":"paper","palavras_chave":[],"sinonimos":[],"tipo":"conceito","titulo":"Certificados esperados"}
\section{Certificados esperados}
) AGM(a,b): itere (a+b2,ab) até |a-b|<10\ensuremath{^{-15}}. =\ensuremath{^{1}}/, grade t\ensuremath{_{n}}=\ensuremath{^{n}}, \ensuremath{_{k}}=2 k/(N). G = DFT unitária (fft, norm=ortho); round-trip até arredondamento. Dilatação U\_ t t; autovalores e\ensuremath{^{-2}} ik/N (||=1). Convolução multiplicativa = convolução em u= t.
\prov{relações: parte\_de main}
\prov{proveniência: paper \textperiodcentered{} fato \textperiodcentered{} confiança 1.0 \textperiodcentered{} domínio paper \textperiodcentered{} história 5 versões no jornal}

%CRISTAL {"arestas":[{"alvo":"montagem","confianca":1.0,"epistemico":"fato","origem":"paper","rel":"parte_de"}],"confianca":1.0,"contraexemplos":[],"descricao":"montar L A montagem autônoma é subtrativa. Cada abaixo é mão --- remover, não adicionar: itemize2pt [ourovivo117] Planejador / scheduler central --- a ordem emerge do recall e da régua de ouro; impô-la mata a autonomia. [ourovivo117] Buffer / cache / janela --- streaming byte a byte; buffer é mão segurando o fluxo. [ourovivo117] Treino do operador --- retocar A é a morte; só se engrossa o tecido que ele percorre. [ourovivo117] Float / lib / rede --- tudo inteiro, determinístico","epistemico":"fato","exemplos":[],"id":"checklist_sistema_autonomo_no_ar","memoria":"semantica","meta":{"dominio":"paper","nivel":5,"verificacao":"conceito novo"},"origem":"paper","palavras_chave":[],"sinonimos":[],"tipo":"conceito","titulo":"Checklist: sistema autônomo no ar"}
\section{Checklist: sistema autônomo no ar}
montar L A montagem autônoma é subtrativa. Cada abaixo é mão --- remover, não adicionar: itemize2pt [ourovivo117] Planejador / scheduler central --- a ordem emerge do recall e da régua de ouro; impô-la mata a autonomia. [ourovivo117] Buffer / cache / janela --- streaming byte a byte; buffer é mão segurando o fluxo. [ourovivo117] Treino do operador --- retocar A é a morte; só se engrossa o tecido que ele percorre. [ourovivo117] Float / lib / rede --- tudo inteiro, determinístico
\prov{relações: parte\_de montagem}
\prov{proveniência: paper \textperiodcentered{} fato \textperiodcentered{} confiança 1.0 \textperiodcentered{} domínio paper \textperiodcentered{} história 3 versões no jornal}

%CRISTAL {"arestas":[{"alvo":"reticulado","confianca":-0.055,"epistemico":"fato","origem":"manual","rel":"referencia"},{"alvo":"casl-propagation","confianca":0.148,"epistemico":"fato","origem":"manual","rel":"referencia"}],"confianca":1.0,"contraexemplos":[],"descricao":"Tudo o que precede é, na verdade, um fluxo monótono sobre um reticulado booleano. Tornar isso explícito entrega vários teoremas quase de graça. definition[O dioid das condições] (2, ) com:= e:= é um semianel idempotente comutativo (dioid): é associativa, comutativa, idempotente (a a=a), com neutro; idem, com neutro; e distribui sobre. A ordem canônica a b a b=b coincide com. Com o complemento, é a álgebra booleana completa (2, ) --- u","epistemico":"fato","exemplos":[],"id":"ciclo_de_vida_revogacao_e_expiracao","memoria":"semantica","meta":{"dominio":"paper","nivel":5},"origem":"paper","palavras_chave":[],"sinonimos":[],"tipo":"conceito","titulo":"Ciclo de vida: revogação e expiração"}
\section{Ciclo de vida: revogação e expiração}
Tudo o que precede é, na verdade, um fluxo monótono sobre um reticulado booleano. Tornar isso explícito entrega vários teoremas quase de graça. definition[O dioid das condições] (2, ) com:= e:= é um semianel idempotente comutativo (dioid): é associativa, comutativa, idempotente (a a=a), com neutro; idem, com neutro; e distribui sobre. A ordem canônica a b a b=b coincide com. Com o complemento, é a álgebra booleana completa (2, ) --- u
\prov{relações: referencia reticulado; referencia casl-propagation}
\prov{proveniência: paper \textperiodcentered{} fato \textperiodcentered{} confiança 1.0 \textperiodcentered{} domínio paper \textperiodcentered{} história 9 versões no jornal}

%CRISTAL {"arestas":[{"alvo":"os_dois_ciclos_de_reconstrucao","confianca":1.0,"epistemico":"fato","origem":"paper","rel":"parte_de"}],"confianca":1.0,"contraexemplos":[],"descricao":"D A arquitectura completa tem dois sentidos: figure[h] [ node distance=1.1cm and 1.4cm, box/.style=draw=ouro!70!black, rounded corners=3pt, fill=creme!90, minimum width=2.6cm, minimum height=0.85cm, =, font=, arr/.style=-Stealth[length=2.2mm], thick, ouro!60!black ] [box] (can) canónico =(w,(w)); [box, right=of can] (diss) dissipação; [box, below right=0.6cm and 0.2cm of diss] (fl) floco (); [box, above right=0.6cm and 0.2cm of diss] (ret)","epistemico":"fato","exemplos":[],"id":"ciclo_i_canonico_floco_reta","memoria":"semantica","meta":{"dominio":"paper","nivel":5,"verificacao":"conceito novo"},"origem":"paper","palavras_chave":[],"sinonimos":[],"tipo":"conceito","titulo":"Ciclo I: canónico floco + reta"}
\section{Ciclo I: canónico floco + reta}
D A arquitectura completa tem dois sentidos: figure[h] [ node distance=1.1cm and 1.4cm, box/.style=draw=ouro!70!black, rounded corners=3pt, fill=creme!90, minimum width=2.6cm, minimum height=0.85cm, =, font=, arr/.style=-Stealth[length=2.2mm], thick, ouro!60!black ] [box] (can) canónico =(w,(w)); [box, right=of can] (diss) dissipação; [box, below right=0.6cm and 0.2cm of diss] (fl) floco (); [box, above right=0.6cm and 0.2cm of diss] (ret)
\prov{relações: parte\_de os\_dois\_ciclos\_de\_reconstrucao}
\prov{proveniência: paper \textperiodcentered{} fato \textperiodcentered{} confiança 1.0 \textperiodcentered{} domínio paper \textperiodcentered{} história 4 versões no jornal}

%CRISTAL {"arestas":[{"alvo":"os_dois_ciclos_de_reconstrucao","confianca":1.0,"epistemico":"fato","origem":"paper","rel":"parte_de"},{"alvo":"transformada_metalica","confianca":1.0,"epistemico":"fato","origem":"wiki","rel":"aplicacao"}],"confianca":1.0,"contraexemplos":[],"descricao":"defi[Dissipação] Dado =(w,(w)) canónico: [(a)] Ressonância: parabolaauditar(w,(w)) (c²,a), a+c²=1. [(b)] Reta: s=(c²); extrair C₃(s), C_(s). [(c)] Floco: (w)=kochcoord(w,(w)) --- Def.~. [(d)] Dissipação: ():=(L(), ()). defi Pipeline implementado: parabolaauditar parabolaparareta kochcoord. Roundtrip Word+Koch: 32/32 em memstoreword~memoria.","epistemico":"fato","exemplos":[],"id":"ciclo_ii_reta_floco_canonico","memoria":"semantica","meta":{"dominio":"paper","nivel":5,"verificacao":"conceito novo"},"origem":"paper","palavras_chave":[],"sinonimos":[],"tipo":"conceito","titulo":"Ciclo II: reta + floco canónico"}
\section{Ciclo II: reta + floco canónico}
defi[Dissipação] Dado =(w,(w)) canónico: [(a)] Ressonância: parabolaauditar(w,(w)) (c\ensuremath{^{2}},a), a+c\ensuremath{^{2}}=1. [(b)] Reta: s=(c\ensuremath{^{2}}); extrair C\ensuremath{_{3}}(s), C\_(s). [(c)] Floco: (w)=kochcoord(w,(w)) --- Def.\~{}. [(d)] Dissipação: ():=(L(), ()). defi Pipeline implementado: parabolaauditar parabolaparareta kochcoord. Roundtrip Word+Koch: 32/32 em memstoreword\~{}memoria.
\prov{relações: parte\_de os\_dois\_ciclos\_de\_reconstrucao; aplicacao transformada\_metalica}
\prov{proveniência: paper \textperiodcentered{} fato \textperiodcentered{} confiança 1.0 \textperiodcentered{} domínio paper \textperiodcentered{} história 5 versões no jornal}

%CRISTAL {"arestas":[{"alvo":"cifra","confianca":1.0,"epistemico":"fato","origem":"paper","rel":"especializa"},{"alvo":"fechamento_blocos","confianca":0.5,"epistemico":"fato","origem":"wiki","rel":"relaciona"},{"alvo":"java","confianca":0.5,"epistemico":"fato","origem":"wiki","rel":"relaciona"},{"alvo":"utilitarismo","confianca":0.5,"epistemico":"fato","origem":"wiki","rel":"relaciona"},{"alvo":"recursao","confianca":0.5,"epistemico":"fato","origem":"wiki","rel":"relaciona"},{"alvo":"topologia","confianca":0.5,"epistemico":"fato","origem":"wiki","rel":"relaciona"},{"alvo":"ipc","confianca":0.5,"epistemico":"fato","origem":"wiki","rel":"relaciona"},{"alvo":"fase_2_tecido_de_conhecimento","confianca":0.5,"epistemico":"fato","origem":"wiki","rel":"relaciona"},{"alvo":"o_que_e_no_projecto","confianca":0.5,"epistemico":"fato","origem":"wiki","rel":"relaciona"}],"confianca":1.0,"contraexemplos":[],"descricao":"Paper: cifra_ouro_branco.tex","epistemico":"fato","exemplos":[],"id":"cifra_ouro_branco","memoria":"semantica","meta":{"dominio":"paper","nivel":5,"tipo_no":"paper"},"origem":"paper","palavras_chave":[],"sinonimos":[],"tipo":"conceito","titulo":"cifra_ouro_branco"}
\section{cifra\_ouro\_branco}
Paper: cifra\_ouro\_branco.tex
\prov{relações: especializa cifra; relaciona fechamento\_blocos; relaciona java; relaciona utilitarismo; relaciona recursao; relaciona topologia; relaciona ipc; relaciona fase\_2\_tecido\_de\_conhecimento; relaciona o\_que\_e\_no\_projecto}
\prov{proveniência: paper \textperiodcentered{} fato \textperiodcentered{} confiança 1.0 \textperiodcentered{} domínio paper \textperiodcentered{} história 15 versões no jornal}

%CRISTAL {"arestas":[{"alvo":"orquestrador_multi-tecidos","confianca":1.0,"epistemico":"fato","origem":"paper","rel":"parte_de"}],"confianca":1.0,"contraexemplos":[],"descricao":"[ st/.style=draw=ourovivo, rounded corners=2pt, minimum width=2.2cm, minimum height=0.55cm, =, font=, fill=creme, arr/.style=-Stealth[length=2mm], color=ourovivo, thick ] [st, minimum width=3cm] (q) pergunta; [st, right=0.5cm of q] (c) classificador; [st, right=0.5cm of c, minimum width=2.8cm] (s) seletor; [st, below=0.5cm of s, xshift=-1.2cm] (g) git.idx; [st, below=0.5cm of s] (l) linux.idx; [st, below=0.5cm of s, xshift=1.2cm] (x) cliente.","epistemico":"fato","exemplos":[],"id":"classificador_evolucao_do_proxy","memoria":"semantica","meta":{"dominio":"paper","nivel":5,"verificacao":"conceito novo"},"origem":"paper","palavras_chave":[],"sinonimos":[],"tipo":"conceito","titulo":"Classificador (evolução do proxy)"}
\section{Classificador (evolução do proxy)}
[ st/.style=draw=ourovivo, rounded corners=2pt, minimum width=2.2cm, minimum height=0.55cm, =, font=, fill=creme, arr/.style=-Stealth[length=2mm], color=ourovivo, thick ] [st, minimum width=3cm] (q) pergunta; [st, right=0.5cm of q] (c) classificador; [st, right=0.5cm of c, minimum width=2.8cm] (s) seletor; [st, below=0.5cm of s, xshift=-1.2cm] (g) git.idx; [st, below=0.5cm of s] (l) linux.idx; [st, below=0.5cm of s, xshift=1.2cm] (x) cliente.
\prov{relações: parte\_de orquestrador\_multi-tecidos}
\prov{proveniência: paper \textperiodcentered{} fato \textperiodcentered{} confiança 1.0 \textperiodcentered{} domínio paper \textperiodcentered{} história 3 versões no jornal}

%CRISTAL {"arestas":[{"alvo":"o_floco_e_a_conservacao_da_comunicacao","confianca":1.0,"epistemico":"fato","origem":"paper","rel":"parte_de"}],"confianca":1.0,"contraexemplos":[],"descricao":"O floco define (w) como o objecto que armazena a informação não-associativa produzida quando o estado canónico =(w,(w)) atravessa a cadeia de retas do gato --- cristalização do Olho de Koch no recipiente negro. A reta associativa L(w)=(s,C₃(s),C_(s)) endereça onde o estado vive (Cantor, Fibonacci, parâmetro s sobre a parábola a+c²=1); (w) guarda o que caiu na passagem (zₖoch z² p, órbita, atrator). São ortogonais:; L sozinha colide (18,548 buckets distintos com c²","epistemico":"fato","exemplos":[],"id":"composicao_multi-hop_e_clones","memoria":"semantica","meta":{"dominio":"paper","nivel":5,"verificacao":"conceito novo"},"origem":"paper","palavras_chave":[],"sinonimos":[],"tipo":"conceito","titulo":"Composição multi-hop e clones"}
\section{Composição multi-hop e clones}
O floco define (w) como o objecto que armazena a informação não-associativa produzida quando o estado canónico =(w,(w)) atravessa a cadeia de retas do gato --- cristalização do Olho de Koch no recipiente negro. A reta associativa L(w)=(s,C\ensuremath{_{3}}(s),C\_(s)) endereça onde o estado vive (Cantor, Fibonacci, parâmetro s sobre a parábola a+c\ensuremath{^{2}}=1); (w) guarda o que caiu na passagem (z\ensuremath{_{k}}och z\ensuremath{^{2}} p, órbita, atrator). São ortogonais:; L sozinha colide (18,548 buckets distintos com c\ensuremath{^{2}}
\prov{relações: parte\_de o\_floco\_e\_a\_conservacao\_da\_comunicacao}
\prov{proveniência: paper \textperiodcentered{} fato \textperiodcentered{} confiança 1.0 \textperiodcentered{} domínio paper \textperiodcentered{} história 3 versões no jornal}

%CRISTAL {"arestas":[{"alvo":"computacao_fundamentos_broca","confianca":1.0,"epistemico":"fato","origem":"paper","rel":"parte_de"}],"confianca":1.0,"contraexemplos":[],"descricao":"abstract Texto de currículo para Tiffany aprender o sistema por dentro: começa na computação como área do conhecimento humano, passa pela arquitetura de Von Neumann (programa armazenado, CPU--memória--bus), pela máquina de Turing (modelo universal, decidibilidade), e fecha em --- onde processo = idempotente Peirce, IPC = canal negro, e o gato A=(smallmatrix1&1 1&0smallmatrix) mede bytes em neuronio.c. Material existente no repositório é referenciado; nada é inventado fora da cadeia. abs","epistemico":"fato","exemplos":[],"id":"computacao_como_area","memoria":"semantica","meta":{"dominio":"paper","nivel":5,"verificacao":"slug idêntico — enriquecer, não duplicar"},"origem":"paper","palavras_chave":[],"sinonimos":[],"tipo":"conceito","titulo":"Computação como área"}
\section{Computação como área}
abstract Texto de currículo para Tiffany aprender o sistema por dentro: começa na computação como área do conhecimento humano, passa pela arquitetura de Von Neumann (programa armazenado, CPU--memória--bus), pela máquina de Turing (modelo universal, decidibilidade), e fecha em --- onde processo = idempotente Peirce, IPC = canal negro, e o gato A=(smallmatrix1\&1 1\&0smallmatrix) mede bytes em neuronio.c. Material existente no repositório é referenciado; nada é inventado fora da cadeia. abs
\prov{relações: parte\_de computacao\_fundamentos\_broca}
\prov{proveniência: paper \textperiodcentered{} fato \textperiodcentered{} confiança 1.0 \textperiodcentered{} domínio paper \textperiodcentered{} história 6 versões no jornal}

%CRISTAL {"arestas":[{"alvo":"computacao","confianca":1.0,"epistemico":"fato","origem":"paper","rel":"explica"},{"alvo":"broca_os","confianca":1.0,"epistemico":"fato","origem":"paper","rel":"explica"},{"alvo":"tiffany","confianca":1.0,"epistemico":"fato","origem":"paper","rel":"referencia"},{"alvo":"desenvolvimento_orientado_a_tarefas","confianca":0.5,"epistemico":"fato","origem":"wiki","rel":"relaciona"},{"alvo":"virtualizacao","confianca":0.5,"epistemico":"fato","origem":"wiki","rel":"relaciona"},{"alvo":"word","confianca":0.5,"epistemico":"fato","origem":"wiki","rel":"relaciona"},{"alvo":"vida-morte","confianca":0.5,"epistemico":"fato","origem":"wiki","rel":"relaciona"}],"confianca":1.0,"contraexemplos":[],"descricao":"Paper: computacao_fundamentos_broca.tex","epistemico":"fato","exemplos":[],"id":"computacao_fundamentos_broca","memoria":"semantica","meta":{"dominio":"paper","nivel":5,"tipo_no":"paper"},"origem":"paper","palavras_chave":[],"sinonimos":[],"tipo":"conceito","titulo":"computacao_fundamentos_broca"}
\section{computacao\_fundamentos\_broca}
Paper: computacao\_fundamentos\_broca.tex
\prov{relações: explica computacao; explica broca\_os; referencia tiffany; relaciona desenvolvimento\_orientado\_a\_tarefas; relaciona virtualizacao; relaciona word; relaciona vida-morte}
\prov{proveniência: paper \textperiodcentered{} fato \textperiodcentered{} confiança 1.0 \textperiodcentered{} domínio paper \textperiodcentered{} história 105 versões no jornal}

%CRISTAL {"arestas":[{"alvo":"corpo_dual","confianca":1.0,"epistemico":"fato","origem":"paper","rel":"parte_de"}],"confianca":1.0,"contraexemplos":[],"descricao":"A rede transmite o bump cru no datagrama, em tempo real (mandar = receber na latência); a banda =sha256(tecido) filtra por assinatura --- a chave errada nem decodifica, vira ruído ---, não por isolamento. Medido em produção (UDP, GEX real). A série tem agora número, calor, incerteza, segurança e canal; resta a forma vista --- a imagem.","epistemico":"fato","exemplos":[],"id":"computacao_grafica","memoria":"semantica","meta":{"dominio":"paper","nivel":5,"verificacao":"conceito novo"},"origem":"paper","palavras_chave":[],"sinonimos":[],"tipo":"conceito","titulo":"Computação Gráfica"}
\section{Computação Gráfica}
A rede transmite o bump cru no datagrama, em tempo real (mandar = receber na latência); a banda =sha256(tecido) filtra por assinatura --- a chave errada nem decodifica, vira ruído ---, não por isolamento. Medido em produção (UDP, GEX real). A série tem agora número, calor, incerteza, segurança e canal; resta a forma vista --- a imagem.
\prov{relações: parte\_de corpo\_dual}
\prov{proveniência: paper \textperiodcentered{} fato \textperiodcentered{} confiança 1.0 \textperiodcentered{} domínio paper \textperiodcentered{} história 3 versões no jornal}

%CRISTAL {"arestas":[{"alvo":"assistente","confianca":1.0,"epistemico":"fato","origem":"paper","rel":"parte_de"}],"confianca":1.0,"contraexemplos":[],"descricao":"D Separação rígida: cada camada tem contrato; nenhuma faz o trabalho da outra. [ box/.style=draw=ourovivo, rounded corners=2pt, minimum width=4.8cm, minimum height=0.75cm, =, font=, fill=creme, arr/.style=-Stealth[length=2.2mm], color=ourovivo, thick ] [box] (t) Tiffany /; [box, below=0.55cm of t] (r) Runtime / API; [box, below=0.55cm of r] (o) Orquestrador; [box, below=0.55cm of o, minimum width=6.2cm] (f) Executor / Agentes; [box, below=0.45","epistemico":"fato","exemplos":[],"id":"conhecimento_e_memoria","memoria":"semantica","meta":{"dominio":"paper","nivel":5,"verificacao":"conceito novo"},"origem":"paper","palavras_chave":[],"sinonimos":[],"tipo":"conceito","titulo":"Conhecimento e memória"}
\section{Conhecimento e memória}
D Separação rígida: cada camada tem contrato; nenhuma faz o trabalho da outra. [ box/.style=draw=ourovivo, rounded corners=2pt, minimum width=4.8cm, minimum height=0.75cm, =, font=, fill=creme, arr/.style=-Stealth[length=2.2mm], color=ourovivo, thick ] [box] (t) Tiffany /; [box, below=0.55cm of t] (r) Runtime / API; [box, below=0.55cm of r] (o) Orquestrador; [box, below=0.55cm of o, minimum width=6.2cm] (f) Executor / Agentes; [box, below=0.45
\prov{relações: parte\_de assistente}
\prov{proveniência: paper \textperiodcentered{} fato \textperiodcentered{} confiança 1.0 \textperiodcentered{} domínio paper \textperiodcentered{} história 3 versões no jornal}

%CRISTAL {"arestas":[{"alvo":"colapso","confianca":0.95,"epistemico":"fato","origem":"paper","rel":"parte_de"}],"confianca":1.0,"contraexemplos":[],"descricao":"Passamos agora da álgebra à geometria: damos ao espaço de álgebras uma métrica e perguntamos que “forma” ele tem. Duas das direções já estão fixadas de modo intrínseco --- o setor radial pela régua de Lopes (g=1, pois é o comprimento de arco, Teor.~) e, em ⁿ, a esfera Sⁿ⁻² das direções a. Resta o setor temporal: a fase do campo local, gerada pela multiplicação, com um lapse N() a determinar, ds²=-N()²,d²+d². equ","epistemico":"fato","exemplos":[],"id":"consequencias_e_direcoes","memoria":"semantica","meta":{"dominio":"paper","nivel":5,"verificacao":"overlap moderado — revisar manualmente"},"origem":"paper","palavras_chave":[],"sinonimos":[],"tipo":"conceito","titulo":"Consequências e direções"}
\section{Consequências e direções}
Passamos agora da álgebra à geometria: damos ao espaço de álgebras uma métrica e perguntamos que “forma” ele tem. Duas das direções já estão fixadas de modo intrínseco --- o setor radial pela régua de Lopes (g=1, pois é o comprimento de arco, Teor.\~{}) e, em \ensuremath{^{n}}, a esfera S\ensuremath{^{n-2}} das direções a. Resta o setor temporal: a fase do campo local, gerada pela multiplicação, com um lapse N() a determinar, ds\ensuremath{^{2}}=-N()\ensuremath{^{2}},d\ensuremath{^{2}}+d\ensuremath{^{2}}. equ
\prov{relações: parte\_de colapso}
\prov{proveniência: paper \textperiodcentered{} fato \textperiodcentered{} confiança 1.0 \textperiodcentered{} domínio paper \textperiodcentered{} história 5 versões no jornal}

%CRISTAL {"arestas":[{"alvo":"fase3","confianca":0.5,"epistemico":"fato","origem":"wiki","rel":"relaciona"},{"alvo":"transcricao","confianca":0.5,"epistemico":"fato","origem":"wiki","rel":"relaciona"},{"alvo":"pow_gato","confianca":0.5,"epistemico":"fato","origem":"wiki","rel":"relaciona"},{"alvo":"goldcoin","confianca":0.5,"epistemico":"fato","origem":"wiki","rel":"relaciona"},{"alvo":"gpu","confianca":0.5,"epistemico":"fato","origem":"wiki","rel":"relaciona"},{"alvo":"fechamento_blocos","confianca":0.5,"epistemico":"fato","origem":"wiki","rel":"relaciona"},{"alvo":"teste_ordem","confianca":0.5,"epistemico":"fato","origem":"wiki","rel":"relaciona"},{"alvo":"prova_isa","confianca":0.5,"epistemico":"fato","origem":"wiki","rel":"relaciona"},{"alvo":"install","confianca":0.5,"epistemico":"fato","origem":"wiki","rel":"relaciona"},{"alvo":"koch_parabola","confianca":0.5,"epistemico":"fato","origem":"wiki","rel":"relaciona"},{"alvo":"montagem","confianca":0.5,"epistemico":"fato","origem":"wiki","rel":"relaciona"},{"alvo":"pow_bit_taxa","confianca":0.5,"epistemico":"fato","origem":"wiki","rel":"relaciona"},{"alvo":"parabola_isa","confianca":0.5,"epistemico":"fato","origem":"wiki","rel":"relaciona"},{"alvo":"sessao_interativa","confianca":0.5,"epistemico":"fato","origem":"wiki","rel":"relaciona"},{"alvo":"python","confianca":0.5,"epistemico":"fato","origem":"wiki","rel":"relaciona"},{"alvo":"readme","confianca":0.5,"epistemico":"fato","origem":"wiki","rel":"relaciona"},{"alvo":"java","confianca":0.5,"epistemico":"fato","origem":"wiki","rel":"relaciona"}],"confianca":1.0,"contraexemplos":[],"descricao":"Paper: contabilidade.tex","epistemico":"fato","exemplos":[],"id":"contabilidade","memoria":"semantica","meta":{"dominio":"paper","nivel":5,"tipo_no":"paper"},"origem":"paper","palavras_chave":[],"sinonimos":[],"tipo":"conceito","titulo":"contabilidade"}
\section{contabilidade}
Paper: contabilidade.tex
\prov{relações: relaciona fase3; relaciona transcricao; relaciona pow\_gato; relaciona goldcoin; relaciona gpu; relaciona fechamento\_blocos; relaciona teste\_ordem; relaciona prova\_isa; relaciona install; relaciona koch\_parabola; relaciona montagem; relaciona pow\_bit\_taxa; relaciona parabola\_isa; relaciona sessao\_interativa; relaciona python; relaciona readme; relaciona java}
\prov{proveniência: paper \textperiodcentered{} fato \textperiodcentered{} confiança 1.0 \textperiodcentered{} domínio paper \textperiodcentered{} história 18 versões no jornal}

%CRISTAL {"arestas":[{"alvo":"base_geral_core","confianca":1.0,"epistemico":"fato","origem":"paper","rel":"parte_de"}],"confianca":1.0,"contraexemplos":[],"descricao":"}\nArquivo raiz da Base Geral --- índice de índices para o orquestrador:\n\nlstlisting\n{\n  \"nome\": \"tiffany-core\",\n  \"versao\": \"1.0.0\",\n  \"dominios\": [\n    {\"path\": \"software/git\", \"idioma\": \"pt\", \"chunks\": 8420,\n     \"tags\": [\"git\", \"versionamento\", \"repositorio\"]},\n    {\"path\": \"software/python\", \"idioma\": \"pt\", \"chunks\": 12000, ...}\n  ]\n}\nlstlisting\n\n Domínios pequenos (ex.: software/c) podem compartilhar curadoria; domínios grandes\n(ex.: software/python) podem subdividir internamente por subdom","epistemico":"fato","exemplos":[],"id":"conteudo_inicial_mvp_core","memoria":"semantica","meta":{"dominio":"paper","nivel":5,"verificacao":"overlap moderado — revisar manualmente"},"origem":"paper","palavras_chave":[],"sinonimos":[],"tipo":"conceito","titulo":"Conteúdo inicial (MVP Core)"}
\section{Conteúdo inicial (MVP Core)}
\}
Arquivo raiz da Base Geral --- índice de índices para o orquestrador:

lstlisting
\{
  "nome": "tiffany-core",
  "versao": "1.0.0",
  "dominios": [
    \{"path": "software/git", "idioma": "pt", "chunks": 8420,
     "tags": ["git", "versionamento", "repositorio"]\},
    \{"path": "software/python", "idioma": "pt", "chunks": 12000, ...\}
  ]
\}
lstlisting

 Domínios pequenos (ex.: software/c) podem compartilhar curadoria; domínios grandes
(ex.: software/python) podem subdividir internamente por subdom
\prov{relações: parte\_de base\_geral\_core}
\prov{proveniência: paper \textperiodcentered{} fato \textperiodcentered{} confiança 1.0 \textperiodcentered{} domínio paper \textperiodcentered{} história 2 versões no jornal}

%CRISTAL {"arestas":[{"alvo":"main","confianca":1.0,"epistemico":"fato","origem":"paper","rel":"parte_de"}],"confianca":1.0,"contraexemplos":[],"descricao":"A torre, concreta em ⁿ, mora num tensor de ordem~3. Com índice 0 escalar e i,j,l vetoriais: ₀₀₀=1, ₀ᵢj=-ᵢj, ᵢ₀j=ᵢj, lᵢ₀=lᵢ, lᵢj=s,lᵢj, com a 3-forma associativa (Levi--Civita em ³; o cross de G₂ em ⁷). A multiplicação é (zs w)l=lᵢjziwj=(a₀b₀- a,b, a₀b+b₀a+s,a b). N O mapa ₇²⁷⁷ tem valores singulares não-nulos todos 3, posto 7 e núcleo de dimensão 14. TN ()=14 (núcleo da equação de Leibniz; Killing n","epistemico":"fato","exemplos":[],"id":"controle_negativo_por_que_antifase_estrita_nao_basta","memoria":"semantica","meta":{"dominio":"paper","nivel":5,"verificacao":"slug idêntico — enriquecer, não duplicar"},"origem":"paper","palavras_chave":[],"sinonimos":[],"tipo":"conceito","titulo":"Controle negativo: por que antifase estrita não basta"}
\section{Controle negativo: por que antifase estrita não basta}
A torre, concreta em \ensuremath{^{n}}, mora num tensor de ordem\~{}3. Com índice 0 escalar e i,j,l vetoriais: \ensuremath{_{000}}=1, \ensuremath{_{0i}}j=-\ensuremath{_{i}}j, \ensuremath{_{i0}}j=\ensuremath{_{i}}j, l\ensuremath{_{i0}}=l\ensuremath{_{i}}, l\ensuremath{_{i}}j=s,l\ensuremath{_{i}}j, com a 3-forma associativa (Levi--Civita em \ensuremath{^{3}}; o cross de G\ensuremath{_{2}} em \ensuremath{^{7}}). A multiplicação é (zs w)l=l\ensuremath{_{i}}jziwj=(a\ensuremath{_{0}}b\ensuremath{_{0}}- a,b, a\ensuremath{_{0}}b+b\ensuremath{_{0}}a+s,a b). N O mapa \ensuremath{_{7}}\ensuremath{^{277}} tem valores singulares não-nulos todos 3, posto 7 e núcleo de dimensão 14. TN ()=14 (núcleo da equação de Leibniz; Killing n
\prov{relações: parte\_de main}
\prov{proveniência: paper \textperiodcentered{} fato \textperiodcentered{} confiança 1.0 \textperiodcentered{} domínio paper \textperiodcentered{} história 5 versões no jornal}

%CRISTAL {"arestas":[{"alvo":"casl-propagation","confianca":-0.023,"epistemico":"fato","origem":"manual","rel":"referencia"},{"alvo":"colapso","confianca":-0.031,"epistemico":"fato","origem":"manual","rel":"referencia"}],"confianca":1.0,"contraexemplos":[],"descricao":"pushdown no Implementação de (Definição~) sobre tecidos Koch. A avaliação independe de (Seção~): basta empurrar Ep(s,a) para o colapso. Num turno pelo principal p, [ (q)=ₙ(p) score(n,q) cfala. ] Referência: conversa/colapso.py. A função scorehit = overlap + peso de camada + relevância - penalidades (deny-overrides) implementa score; fail-closed None. Promover turnos A — B: hitsparascore. Tecido: dialogos.tsv / di","epistemico":"fato","exemplos":[],"id":"conversa_tiffany","memoria":"semantica","meta":{"dominio":"paper","nivel":5},"origem":"paper","palavras_chave":[],"sinonimos":[],"tipo":"conceito","titulo":"Conversa Tiffany"}
\section{Conversa Tiffany}
pushdown no Implementação de (Definição\~{}) sobre tecidos Koch. A avaliação independe de (Seção\~{}): basta empurrar Ep(s,a) para o colapso. Num turno pelo principal p, [ (q)=\ensuremath{_{n}}(p) score(n,q) cfala. ] Referência: conversa/colapso.py. A função scorehit = overlap + peso de camada + relevância - penalidades (deny-overrides) implementa score; fail-closed None. Promover turnos A — B: hitsparascore. Tecido: dialogos.tsv / di
\prov{relações: referencia casl-propagation; referencia colapso}
\prov{proveniência: paper \textperiodcentered{} fato \textperiodcentered{} confiança 1.0 \textperiodcentered{} domínio paper \textperiodcentered{} história 10 versões no jornal}

%CRISTAL {"arestas":[{"alvo":"transcricao","confianca":1.0,"epistemico":"fato","origem":"paper","rel":"parte_de"},{"alvo":"colapso","confianca":1.0,"epistemico":"fato","origem":"paper","rel":"parte_de"}],"confianca":1.0,"contraexemplos":[],"descricao":"D O eixo Res=12 (o ponto duplo, o cone, ponto fixo de ) cai no meio, entre a reta real e a reta de ouro (₁=, a primeira transição), e ali a lemniscata codifica discretocontínuo: [ primo de [i]ᵢrredutível irracionalfração coₙtíₙua ᵢₙfᵢₙᵢta, composto/racionaloperação de aₙel fração contínua finita, ] com as pontes de Gauss (real dígitos) e dos convergentes (Fibonacci ): a reta de ouro é. segurança e zeros são o mesmo corpo discreto","epistemico":"fato","exemplos":[],"id":"corolario_analitico_o_operador_de_riemann_koopman_e_s-ds","memoria":"semantica","meta":{"dominio":"paper","nivel":5,"verificacao":"conceito novo"},"origem":"paper","palavras_chave":[],"sinonimos":[],"tipo":"conceito","titulo":"Corolário analítico: o operador de Riemann (Koopman) e (s-D)=(s)"}
\section{Corolário analítico: o operador de Riemann (Koopman) e (s-D)=(s)}
D O eixo Res=12 (o ponto duplo, o cone, ponto fixo de ) cai no meio, entre a reta real e a reta de ouro (\ensuremath{_{1}}=, a primeira transição), e ali a lemniscata codifica discretocontínuo: [ primo de [i]\ensuremath{_{i}}rredutível irracionalfração co\ensuremath{_{n}}tí\ensuremath{_{n}}ua \ensuremath{_{in}}f\ensuremath{_{ini}}ta, composto/racionaloperação de a\ensuremath{_{n}}el fração contínua finita, ] com as pontes de Gauss (real dígitos) e dos convergentes (Fibonacci ): a reta de ouro é. segurança e zeros são o mesmo corpo discreto
\prov{relações: parte\_de transcricao; parte\_de colapso}
\prov{proveniência: paper \textperiodcentered{} fato \textperiodcentered{} confiança 1.0 \textperiodcentered{} domínio paper \textperiodcentered{} história 10 versões no jornal}

%CRISTAL {"arestas":[{"alvo":"transcricao","confianca":1.0,"epistemico":"fato","origem":"paper","rel":"parte_de"},{"alvo":"metais","confianca":1.0,"epistemico":"fato","origem":"paper","rel":"usa"}],"confianca":1.0,"contraexemplos":[],"descricao":"Koch não é mais um corolário entre outros: é o nó, ao lado do invariante g. Tudo o que segue --- geometria, aritmética, análise --- passa por ela. defi[Koch, o ribossomo --- e os termos da interface]I A curva de Koch é a substituição (sistema-L) que lê uma regra de ângulos --- a fita --- e sintetiza a forma (a curva-limite), com as sete propriedades definidoras do ribossomo (tradutor, processivo, universal, amplificador 1!!4ⁿ, código fixo, síntese de polímero, ribozima); costurar é gerar","epistemico":"fato","exemplos":[],"id":"corolario_aritmetico_a_torre_de_metais_e_ecc_bsd","memoria":"semantica","meta":{"dominio":"paper","nivel":5,"verificacao":"conceito novo"},"origem":"paper","palavras_chave":[],"sinonimos":[],"tipo":"conceito","titulo":"Corolário aritmético: a torre de metais e ECC = BSD"}
\section{Corolário aritmético: a torre de metais e ECC = BSD}
Koch não é mais um corolário entre outros: é o nó, ao lado do invariante g. Tudo o que segue --- geometria, aritmética, análise --- passa por ela. defi[Koch, o ribossomo --- e os termos da interface]I A curva de Koch é a substituição (sistema-L) que lê uma regra de ângulos --- a fita --- e sintetiza a forma (a curva-limite), com as sete propriedades definidoras do ribossomo (tradutor, processivo, universal, amplificador 1!!4\ensuremath{^{n}}, código fixo, síntese de polímero, ribozima); costurar é gerar
\prov{relações: parte\_de transcricao; usa metais}
\prov{proveniência: paper \textperiodcentered{} fato \textperiodcentered{} confiança 1.0 \textperiodcentered{} domínio paper \textperiodcentered{} história 10 versões no jornal}

%CRISTAL {"arestas":[{"alvo":"colapso","confianca":0.95,"epistemico":"fato","origem":"paper","rel":"parte_de"}],"confianca":1.0,"contraexemplos":[],"descricao":"[A atenção cristalina herda o limite] Toda atenção cristalina --- em particular a dos transformers, cujos scores QK^ acessam a matriz de Gram dos tokens por produtos, com profundidade finita --- está na classe do Teorema~. Logo aproxima a estrutura cristalina de uma sequência (a função de onda) até a cota L, e não além. Não se afirma que a atenção é (não o é); demonstra-se que é o limite exato que a atenção cristalina, e qualquer rede real, persegue sem atingir em","epistemico":"fato","exemplos":[],"id":"corollary_a_atencao_cristalina_herda_o_limitecoratencao_t","memoria":"semantica","meta":{"dominio":"paper","nivel":5,"tipo_teorema":"corollary","verificacao":"conceito novo"},"origem":"paper","palavras_chave":[],"sinonimos":[],"tipo":"conceito","titulo":"corollary: [A atenção cristalina herda o limite]cor:atencao\nT…"}
\section{corollary: [A atenção cristalina herda o limite]cor:atencao
T…}
[A atenção cristalina herda o limite] Toda atenção cristalina --- em particular a dos transformers, cujos scores QK\^{} acessam a matriz de Gram dos tokens por produtos, com profundidade finita --- está na classe do Teorema\~{}. Logo aproxima a estrutura cristalina de uma sequência (a função de onda) até a cota L, e não além. Não se afirma que a atenção é (não o é); demonstra-se que é o limite exato que a atenção cristalina, e qualquer rede real, persegue sem atingir em
\prov{relações: parte\_de colapso}
\prov{proveniência: paper \textperiodcentered{} fato \textperiodcentered{} confiança 1.0 \textperiodcentered{} domínio paper \textperiodcentered{} história 4 versões no jornal}

%CRISTAL {"arestas":[{"alvo":"plano_no_so","confianca":1.0,"epistemico":"fato","origem":"paper","rel":"prova"},{"alvo":"bai","confianca":0.039,"epistemico":"fato","origem":"manual","rel":"referencia"}],"confianca":1.0,"contraexemplos":[],"descricao":"[Compartilhamento controlado] Conceder uma capacidade cross-tenant com =0 garante que o destinatário a exerce mas ninguém abaixo dele a herda --- exatamente “dar acesso sem propagar”. Com =k, a difusão é confinada à vizinhança de re-delegação de raio k.","epistemico":"fato","exemplos":[],"id":"corollary_compartilhamento_controlado_conceder_uma_capacid","memoria":"semantica","meta":{"dominio":"paper","nivel":5,"tipo_teorema":"corollary"},"origem":"paper","palavras_chave":[],"sinonimos":[],"tipo":"conceito","titulo":"corollary: [Compartilhamento controlado]\nConceder uma capacid…"}
\section{corollary: [Compartilhamento controlado]
Conceder uma capacid…}
[Compartilhamento controlado] Conceder uma capacidade cross-tenant com =0 garante que o destinatário a exerce mas ninguém abaixo dele a herda --- exatamente “dar acesso sem propagar”. Com =k, a difusão é confinada à vizinhança de re-delegação de raio k.
\prov{relações: prova plano\_no\_so; referencia bai}
\prov{proveniência: paper \textperiodcentered{} fato \textperiodcentered{} confiança 1.0 \textperiodcentered{} domínio paper \textperiodcentered{} história 6 versões no jornal}

%CRISTAL {"arestas":[{"alvo":"colapso","confianca":0.95,"epistemico":"fato","origem":"paper","rel":"parte_de"}],"confianca":1.0,"contraexemplos":[],"descricao":"[Conservação da norma] A isometria de é conservação de energia: nenhuma energia vaza na transformada; o que entra como dado sai como função de onda, bit a bit (Seção~). Em álgebra normada (Seção~) esta é a forma multiplicativa |xy|=|x|,|y|.","epistemico":"fato","exemplos":[],"id":"corollary_conservacao_da_normacorcons_a_isometria_de_e","memoria":"semantica","meta":{"dominio":"paper","nivel":5,"tipo_teorema":"corollary","verificacao":"conceito novo"},"origem":"paper","palavras_chave":[],"sinonimos":[],"tipo":"conceito","titulo":"Conservação da norma"}
\section{Conservação da norma}
[Conservação da norma] A isometria de é conservação de energia: nenhuma energia vaza na transformada; o que entra como dado sai como função de onda, bit a bit (Seção\~{}). Em álgebra normada (Seção\~{}) esta é a forma multiplicativa |xy|=|x|,|y|.
\prov{relações: parte\_de colapso}
\prov{proveniência: paper \textperiodcentered{} fato \textperiodcentered{} confiança 1.0 \textperiodcentered{} domínio paper \textperiodcentered{} história 5 versões no jornal}

%CRISTAL {"arestas":[{"alvo":"plano_no_so","confianca":1.0,"epistemico":"fato","origem":"paper","rel":"prova"}],"confianca":1.0,"contraexemplos":[],"descricao":"[Expiração cascateia sozinha] Como fᵢlhopaᵢ, ao expirar o pai (tpaᵢ) todos os filhos já expiraram. A expiração não precisa de travessia --- o invariante de atenuação a propaga.","epistemico":"fato","exemplos":[],"id":"corollary_expiracao_cascateia_sozinha_como_filho_pai","memoria":"semantica","meta":{"dominio":"paper","nivel":5,"tipo_teorema":"corollary"},"origem":"paper","palavras_chave":[],"sinonimos":[],"tipo":"conceito","titulo":"Expiração cascateia sozinha"}
\section{Expiração cascateia sozinha}
[Expiração cascateia sozinha] Como f\ensuremath{_{i}}lhopa\ensuremath{_{i}}, ao expirar o pai (tpa\ensuremath{_{i}}) todos os filhos já expiraram. A expiração não precisa de travessia --- o invariante de atenuação a propaga.
\prov{relações: prova plano\_no\_so}
\prov{proveniência: paper \textperiodcentered{} fato \textperiodcentered{} confiança 1.0 \textperiodcentered{} domínio paper \textperiodcentered{} história 4 versões no jornal}

%CRISTAL {"arestas":[{"alvo":"plano_no_so","confianca":1.0,"epistemico":"fato","origem":"paper","rel":"prova"}],"confianca":1.0,"contraexemplos":[],"descricao":"[Isolamento de tenant] Um usuário comum da organização t só detém capacidades atenuadas da capacidade intrínseca do dono, todas com c=t. Logo Epr:r.=t: nenhuma linha de outro tenant aparece, sem regra especial. Um vazamento cross-tenant só ocorre por uma aresta de concessão explícita (auditável), cuja difusão é exatamente governada por.","epistemico":"fato","exemplos":[],"id":"corollary_isolamento_de_tenant_um_usuario_comum_da_organiz","memoria":"semantica","meta":{"dominio":"paper","nivel":5,"tipo_teorema":"corollary"},"origem":"paper","palavras_chave":[],"sinonimos":[],"tipo":"conceito","titulo":"corollary: [Isolamento de tenant]\nUm usuário comum da organiz…"}
\section{corollary: [Isolamento de tenant]
Um usuário comum da organiz…}
[Isolamento de tenant] Um usuário comum da organização t só detém capacidades atenuadas da capacidade intrínseca do dono, todas com c=t. Logo Epr:r.=t: nenhuma linha de outro tenant aparece, sem regra especial. Um vazamento cross-tenant só ocorre por uma aresta de concessão explícita (auditável), cuja difusão é exatamente governada por.
\prov{relações: prova plano\_no\_so}
\prov{proveniência: paper \textperiodcentered{} fato \textperiodcentered{} confiança 1.0 \textperiodcentered{} domínio paper \textperiodcentered{} história 3 versões no jornal}

%CRISTAL {"arestas":[{"alvo":"corpo_tropical","confianca":1.0,"epistemico":"fato","origem":"paper","rel":"relaciona"},{"alvo":"corpo_glacial","confianca":1.0,"epistemico":"fato","origem":"paper","rel":"relaciona"},{"alvo":"gato","confianca":1.0,"epistemico":"fato","origem":"paper","rel":"relaciona"},{"alvo":"reais","confianca":1.0,"epistemico":"fato","origem":"paper","rel":"relaciona"},{"alvo":"ponte_simbolica_gentil","confianca":1.0,"epistemico":"fato","origem":"paper","rel":"relaciona"},{"alvo":"geometria","confianca":1.0,"epistemico":"fato","origem":"paper","rel":"relaciona"},{"alvo":"termodinamica","confianca":1.0,"epistemico":"fato","origem":"paper","rel":"relaciona"},{"alvo":"mecanica","confianca":1.0,"epistemico":"fato","origem":"paper","rel":"relaciona"}],"confianca":1.0,"contraexemplos":[],"descricao":"Série completa tropical↔glacial: topologia→gato→balança a+c²=1→geometria. Unificação no vazio ℝ (vazamento finito → completude).","epistemico":"fato","exemplos":["neuronio/verificar_corpo_dual.py — 19/19 unificação/vazio","vazamento β^{-K}→0; metais σ_n codificam ℝ"],"id":"corpo_dual","memoria":"semantica","meta":{"dominio":"paper","nivel":5,"serve":"Hub engenharia de ouro — duas faces, um contínuo","verificacao":"slug idêntico — enriquecer, não duplicar"},"origem":"paper","palavras_chave":[],"sinonimos":["o corpo dual"],"tipo":"conceito","titulo":"corpo dual"}
\section{corpo dual}
Série completa tropical\ensuremath{\leftrightarrow}glacial: topologia\ensuremath{\to}gato\ensuremath{\to}balança a+c\ensuremath{^{2}}=1\ensuremath{\to}geometria. Unificação no vazio \ensuremath{\mathbb{R}} (vazamento finito \ensuremath{\to} completude).
\prov{exemplos: neuronio/verificar\_corpo\_dual.py — 19/19 unificação/vazio; vazamento \ensuremath{\beta}\^{}\{-K\}\ensuremath{\to}0; metais \ensuremath{\sigma}\_n codificam \ensuremath{\mathbb{R}}}
\prov{sinónimos: o corpo dual}
\prov{relações: relaciona corpo\_tropical; relaciona corpo\_glacial; relaciona gato; relaciona reais; relaciona ponte\_simbolica\_gentil; relaciona geometria; relaciona termodinamica; relaciona mecanica}
\prov{proveniência: paper \textperiodcentered{} fato \textperiodcentered{} confiança 1.0 \textperiodcentered{} domínio paper \textperiodcentered{} história 335 versões no jornal}

%CRISTAL {"arestas":[{"alvo":"corpo_tropical","confianca":1.0,"epistemico":"fato","origem":"paper","rel":"relaciona"},{"alvo":"corpo_dual","confianca":1.0,"epistemico":"fato","origem":"paper","rel":"relaciona"},{"alvo":"gato","confianca":1.0,"epistemico":"fato","origem":"paper","rel":"relaciona"},{"alvo":"olho_de_koch","confianca":1.0,"epistemico":"fato","origem":"paper","rel":"relaciona"},{"alvo":"a_decomposicao_de_peirce_celulas_e_canais","confianca":1.0,"epistemico":"fato","origem":"paper","rel":"relaciona"}],"confianca":1.0,"contraexemplos":[],"descricao":"Máquina negra — células idempotentes (branca) + canais nilpotentes square-zero (negra), acoplados por Peirce. Clones = células e_i Ae_i.","epistemico":"fato","exemplos":["neuronio/verificar_corpo_glacial.py — 14/14 Peirce/clones","neuronio/radio.py — demux e_i M e_j, clone La Hire"],"id":"corpo_glacial","memoria":"semantica","meta":{"dominio":"paper","nivel":5,"serve":"Face nilpotente da quadração — fluxo, canal, clones","verificacao":"slug idêntico — enriquecer, não duplicar"},"origem":"paper","palavras_chave":[],"sinonimos":["o corpo glacial","parte negra","regime glacial"],"tipo":"conceito","titulo":"corpo glacial"}
\section{corpo glacial}
Máquina negra — células idempotentes (branca) + canais nilpotentes square-zero (negra), acoplados por Peirce. Clones = células e\_i Ae\_i.
\prov{exemplos: neuronio/verificar\_corpo\_glacial.py — 14/14 Peirce/clones; neuronio/radio.py — demux e\_i M e\_j, clone La Hire}
\prov{sinónimos: o corpo glacial; parte negra; regime glacial}
\prov{relações: relaciona corpo\_tropical; relaciona corpo\_dual; relaciona gato; relaciona olho\_de\_koch; relaciona a\_decomposicao\_de\_peirce\_celulas\_e\_canais}
\prov{proveniência: paper \textperiodcentered{} fato \textperiodcentered{} confiança 1.0 \textperiodcentered{} domínio paper \textperiodcentered{} história 187 versões no jornal}

%CRISTAL {"arestas":[{"alvo":"reais","confianca":1.0,"epistemico":"fato","origem":"paper","rel":"define"},{"alvo":"gentil_fundamentos_dos_numeros","confianca":1.0,"epistemico":"fato","origem":"paper","rel":"paralelo"},{"alvo":"ponte_simbolica_gentil","confianca":1.0,"epistemico":"fato","origem":"paper","rel":"confirma"},{"alvo":"corpo_glacial","confianca":1.0,"epistemico":"fato","origem":"paper","rel":"paralelo"}],"confianca":1.0,"contraexemplos":[],"descricao":"O invariante é a classe topológica --- dois inteiros, GL(2, Z). Dela o espectro () e a entropia () caem como sombras: topologia primeiro. O próximo capítulo toma esse /gato como dado e constrói, do regime idempotente, o número --- a reta de ouro que é a reta real.","epistemico":"fato","exemplos":[],"id":"corpo_tropical","memoria":"semantica","meta":{"dominio":"paper","duplicata_sugerida":[],"nivel":5,"verificacao":"parte de conceito mais amplo 'corpo_tropical'"},"origem":"paper","palavras_chave":[],"sinonimos":[],"tipo":"conceito","titulo":"corpo tropical"}
\section{corpo tropical}
O invariante é a classe topológica --- dois inteiros, GL(2, Z). Dela o espectro () e a entropia () caem como sombras: topologia primeiro. O próximo capítulo toma esse /gato como dado e constrói, do regime idempotente, o número --- a reta de ouro que é a reta real.
\prov{relações: define reais; paralelo gentil\_fundamentos\_dos\_numeros; confirma ponte\_simbolica\_gentil; paralelo corpo\_glacial}
\prov{proveniência: paper \textperiodcentered{} fato \textperiodcentered{} confiança 1.0 \textperiodcentered{} domínio paper \textperiodcentered{} história 157 versões no jornal}

%CRISTAL {"arestas":[{"alvo":"leitura_cosmologica_a_energia_escura_e_o_universo_ds","confianca":1.0,"epistemico":"fato","origem":"paper","rel":"parte_de"}],"confianca":1.0,"contraexemplos":[],"descricao":"O mesmo método que recuperou o 3/2 de Kepler (Seção~), agora nos fótons das supernovas. A distância de luminosidade obedece localmente dL z(z); em baixo z é a lei de Hubble (dL z, =1), e o crescimento de acima de 1 é a assinatura da aceleração. Como o módulo de distância é =5₁₀dL+const e (z)=5,(₁₀z), a primeira diferença cancela a magnitude absoluta M e H₀ --- não é preciso calibrar a vela-padrão nem","epistemico":"fato","exemplos":[],"id":"cosmologia_fisica_3","memoria":"semantica","meta":{"dominio":"paper","nivel":5,"verificacao":"conceito novo"},"origem":"paper","palavras_chave":[],"sinonimos":[],"tipo":"conceito","titulo":"Cosmologia física 3+"}
\section{Cosmologia física 3+}
O mesmo método que recuperou o 3/2 de Kepler (Seção\~{}), agora nos fótons das supernovas. A distância de luminosidade obedece localmente dL z(z); em baixo z é a lei de Hubble (dL z, =1), e o crescimento de acima de 1 é a assinatura da aceleração. Como o módulo de distância é =5\ensuremath{_{10}}dL+const e (z)=5,(\ensuremath{_{10}}z), a primeira diferença cancela a magnitude absoluta M e H\ensuremath{_{0}} --- não é preciso calibrar a vela-padrão nem
\prov{relações: parte\_de leitura\_cosmologica\_a\_energia\_escura\_e\_o\_universo\_ds}
\prov{proveniência: paper \textperiodcentered{} fato \textperiodcentered{} confiança 1.0 \textperiodcentered{} domínio paper \textperiodcentered{} história 4 versões no jornal}

%CRISTAL {"arestas":[{"alvo":"ponte_simbolica_gentil","confianca":1.0,"epistemico":"fato","origem":"paper","rel":"parte_de"},{"alvo":"semantica","confianca":1.0,"epistemico":"fato","origem":"paper","rel":"referencia"}],"confianca":1.0,"contraexemplos":[],"descricao":"Charts e functores Sejam [nosep,leftmargin=1.4em] G: objetos = classes [(rₙ)] C, morfismos = operações do corpo ordenado completo; E: objetos = pontos x com endereço Bergman + espectro, morfismos = decodificação, convolução multiplicativa. Um chart é uma realização fiel parcial: D ou: D' em domínios D, D' C ou. definicao Homomorfismo dos convergentes Seja [(rₙ)] C o real. As classes com repre","epistemico":"fato","exemplos":[],"id":"criterio_de_nao-isomorfismo","memoria":"semantica","meta":{"dominio":"paper","nivel":5,"verificacao":"slug idêntico — enriquecer, não duplicar"},"origem":"paper","palavras_chave":[],"sinonimos":[],"tipo":"conceito","titulo":"Critério de (não-)isomorfismo"}
\section{Critério de (não-)isomorfismo}
Charts e functores Sejam [nosep,leftmargin=1.4em] G: objetos = classes [(r\ensuremath{_{n}})] C, morfismos = operações do corpo ordenado completo; E: objetos = pontos x com endereço Bergman + espectro, morfismos = decodificação, convolução multiplicativa. Um chart é uma realização fiel parcial: D ou: D' em domínios D, D' C ou. definicao Homomorfismo dos convergentes Seja [(r\ensuremath{_{n}})] C o real. As classes com repre
\prov{relações: parte\_de ponte\_simbolica\_gentil; referencia semantica}
\prov{proveniência: paper \textperiodcentered{} fato \textperiodcentered{} confiança 1.0 \textperiodcentered{} domínio paper \textperiodcentered{} história 47 versões no jornal}

%CRISTAL {"arestas":[{"alvo":"transformada_dourada","confianca":1.0,"epistemico":"fato","origem":"paper","rel":"parte_de"}],"confianca":1.0,"contraexemplos":[],"descricao":"A transformada dourada tem o análogo, termo a termo, de cada propriedade de Fourier --- com a troca sistemática aditivomultiplicativo imposta por. A Tabela~ resume; as provas seguem. table[h] 1.5 tabularl l l Propriedade & Fourier (t!!) & Dourada () linearidade & ax+by a x+b y & ax+by a x+b y deslocamento fase & x(t-) e⁻i x & escala x( t)i x modulação deslocamento & ei⁰tx x(-₀) & potência tax x(+ia) convolução produt","epistemico":"fato","exemplos":[],"id":"dados_brutos_espectro","memoria":"semantica","meta":{"dominio":"paper","nivel":5,"verificacao":"conceito novo"},"origem":"paper","palavras_chave":[],"sinonimos":[],"tipo":"conceito","titulo":"Dados brutos espectro"}
\section{Dados brutos espectro}
A transformada dourada tem o análogo, termo a termo, de cada propriedade de Fourier --- com a troca sistemática aditivomultiplicativo imposta por. A Tabela\~{} resume; as provas seguem. table[h] 1.5 tabularl l l Propriedade \& Fourier (t!!) \& Dourada () linearidade \& ax+by a x+b y \& ax+by a x+b y deslocamento fase \& x(t-) e\ensuremath{^{-}}i x \& escala x( t)i x modulação deslocamento \& ei\ensuremath{^{0}}tx x(-\ensuremath{_{0}}) \& potência tax x(+ia) convolução produt
\prov{relações: parte\_de transformada\_dourada}
\prov{proveniência: paper \textperiodcentered{} fato \textperiodcentered{} confiança 1.0 \textperiodcentered{} domínio paper \textperiodcentered{} história 4 versões no jornal}

%CRISTAL {"arestas":[{"alvo":"mecanica_quantica","confianca":1.0,"epistemico":"fato","origem":"paper","rel":"parte_de"},{"alvo":"mecanica","confianca":1.0,"epistemico":"fato","origem":"paper","rel":"parte_de"}],"confianca":1.0,"contraexemplos":[],"descricao":"(2), qutrit e Bell O furo que parecia faltar --- “de onde vem o espaço de Hilbert e a relação de comutação?” --- não falta: a representação é a álgebra dos quaternions. Os imaginários i,j,k são os operadores de spin (i=-ix, etc.), S³=SU(2), e a relação não é postulada --- é a tábua de multiplicação: [ [x,y]=2i,z i,j=k, ] verificado exato. Precisão: esta é a álgebra de Lie su(2) do spin --- não o CCR canônico de Heisenberg [ q, p]=i,1, que é impossível em dimensão finita","epistemico":"fato","exemplos":[],"id":"de_um_qubit_a_n_funcoes_de_onda","memoria":"semantica","meta":{"dominio":"paper","nivel":5,"verificacao":"slug idêntico — enriquecer, não duplicar"},"origem":"paper","palavras_chave":[],"sinonimos":[],"tipo":"conceito","titulo":"De um qubit a N funções de onda"}
\section{De um qubit a N funções de onda}
(2), qutrit e Bell O furo que parecia faltar --- “de onde vem o espaço de Hilbert e a relação de comutação?” --- não falta: a representação é a álgebra dos quaternions. Os imaginários i,j,k são os operadores de spin (i=-ix, etc.), S\ensuremath{^{3}}=SU(2), e a relação não é postulada --- é a tábua de multiplicação: [ [x,y]=2i,z i,j=k, ] verificado exato. Precisão: esta é a álgebra de Lie su(2) do spin --- não o CCR canônico de Heisenberg [ q, p]=i,1, que é impossível em dimensão finita
\prov{relações: parte\_de mecanica\_quantica; parte\_de mecanica}
\prov{proveniência: paper \textperiodcentered{} fato \textperiodcentered{} confiança 1.0 \textperiodcentered{} domínio paper \textperiodcentered{} história 13 versões no jornal}

%CRISTAL {"arestas":[{"alvo":"plano_no_so","confianca":1.0,"epistemico":"fato","origem":"paper","rel":"define"}],"confianca":1.0,"contraexemplos":[],"descricao":"[Atenuação e propagação] Sobre extents: T:2X2X é monótono deflacionário (T, x y T(x) T(y)). Sobre detentores: P:2 P2 P é monótono inflacionário (P). Re-delegar por saltos é P; = corresponde ao menor ponto fixo P=ₙ Pⁿ(H₀).","epistemico":"fato","exemplos":[],"id":"definition_atenuacao_e_propagacaodeftp_sobre_extents_t2","memoria":"semantica","meta":{"dominio":"paper","nivel":5,"tipo_teorema":"definition"},"origem":"paper","palavras_chave":[],"sinonimos":[],"tipo":"conceito","titulo":"Atenuação e propagação"}
\section{Atenuação e propagação}
[Atenuação e propagação] Sobre extents: T:2X2X é monótono deflacionário (T, x y T(x) T(y)). Sobre detentores: P:2 P2 P é monótono inflacionário (P). Re-delegar por saltos é P; = corresponde ao menor ponto fixo P=\ensuremath{_{n}} P\ensuremath{^{n}}(H\ensuremath{_{0}}).
\prov{relações: define plano\_no\_so}
\prov{proveniência: paper \textperiodcentered{} fato \textperiodcentered{} confiança 1.0 \textperiodcentered{} domínio paper \textperiodcentered{} história 4 versões no jornal}

%CRISTAL {"arestas":[{"alvo":"plano_no_so","confianca":1.0,"epistemico":"fato","origem":"paper","rel":"define"}],"confianca":1.0,"contraexemplos":[],"descricao":"[Autoridades intrínsecas] Duas fontes de autoridade não derivam de ninguém: [nosep] o operador detém (manage,all, 0,) --- vê e faz tudo, com re-delegação irrestrita; o dono do tenant t detém (manage,all,(=t),0,) --- controle total dos próprios recursos.","epistemico":"fato","exemplos":[],"id":"definition_autoridades_intrinsecas_duas_fontes_de_autoridad","memoria":"semantica","meta":{"dominio":"paper","nivel":5,"tipo_teorema":"definition"},"origem":"paper","palavras_chave":[],"sinonimos":[],"tipo":"conceito","titulo":"definition: [Autoridades intrínsecas]\nDuas fontes de autoridad…"}
\section{definition: [Autoridades intrínsecas]
Duas fontes de autoridad…}
[Autoridades intrínsecas] Duas fontes de autoridade não derivam de ninguém: [nosep] o operador detém (manage,all, 0,) --- vê e faz tudo, com re-delegação irrestrita; o dono do tenant t detém (manage,all,(=t),0,) --- controle total dos próprios recursos.
\prov{relações: define plano\_no\_so}
\prov{proveniência: paper \textperiodcentered{} fato \textperiodcentered{} confiança 1.0 \textperiodcentered{} domínio paper \textperiodcentered{} história 3 versões no jornal}

%CRISTAL {"arestas":[{"alvo":"plano_no_so","confianca":1.0,"epistemico":"fato","origem":"paper","rel":"define"}],"confianca":1.0,"contraexemplos":[],"descricao":"[Avaliação filtra o ciclo de vida] Em Hold(p) e em Ep(s,a) (Seção~3) contam somente as instâncias vigentes no instante t: =active e >t. As demais não contribuem. Assim revogar/suspender/expirar uma concessão remove o seu extent da união sem tocar na álgebra.","epistemico":"fato","exemplos":[],"id":"definition_avaliacao_filtra_o_ciclo_de_vida_em_holdp_e","memoria":"semantica","meta":{"dominio":"paper","nivel":5,"tipo_teorema":"definition"},"origem":"paper","palavras_chave":[],"sinonimos":[],"tipo":"conceito","titulo":"Avaliação filtra o ciclo de vida"}
\section{Avaliação filtra o ciclo de vida}
[Avaliação filtra o ciclo de vida] Em Hold(p) e em Ep(s,a) (Seção\~{}3) contam somente as instâncias vigentes no instante t: =active e >t. As demais não contribuem. Assim revogar/suspender/expirar uma concessão remove o seu extent da união sem tocar na álgebra.
\prov{relações: define plano\_no\_so}
\prov{proveniência: paper \textperiodcentered{} fato \textperiodcentered{} confiança 1.0 \textperiodcentered{} domínio paper \textperiodcentered{} história 4 versões no jornal}

%CRISTAL {"arestas":[{"alvo":"plano_no_so","confianca":1.0,"epistemico":"fato","origem":"paper","rel":"define"},{"alvo":"kappa","confianca":-0.039,"epistemico":"fato","origem":"manual","rel":"referencia"},{"alvo":"delta","confianca":0.086,"epistemico":"fato","origem":"manual","rel":"referencia"}],"confianca":1.0,"contraexemplos":[],"descricao":"[Capacidade CASL] Specialização da Definição~ com X= (linhas), =(a,s) (ação e subject), (=1 é proibição, cannot) e como orçamento (Definição~). Formalmente [ =(a, s, c, ), ] autoriza (ou nega, se =1) a sobre s nas linhas de c, e admite re-delegação conforme.","epistemico":"fato","exemplos":[],"id":"definition_capacidade_casldefcap_specializacao_da_definica","memoria":"semantica","meta":{"dominio":"paper","nivel":5,"tipo_teorema":"definition"},"origem":"paper","palavras_chave":[],"sinonimos":[],"tipo":"conceito","titulo":"definition: [Capacidade CASL]def:cap\nSpecialização da Definiçã…"}
\section{definition: [Capacidade CASL]def:cap
Specialização da Definiçã…}
[Capacidade CASL] Specialização da Definição\~{} com X= (linhas), =(a,s) (ação e subject), (=1 é proibição, cannot) e como orçamento (Definição\~{}). Formalmente [ =(a, s, c, ), ] autoriza (ou nega, se =1) a sobre s nas linhas de c, e admite re-delegação conforme.
\prov{relações: define plano\_no\_so; referencia kappa; referencia delta}
\prov{proveniência: paper \textperiodcentered{} fato \textperiodcentered{} confiança 1.0 \textperiodcentered{} domínio paper \textperiodcentered{} história 9 versões no jornal}

%CRISTAL {"arestas":[{"alvo":"plano_no_so","confianca":1.0,"epistemico":"fato","origem":"paper","rel":"define"},{"alvo":"kappa","confianca":0.023,"epistemico":"fato","origem":"manual","rel":"referencia"},{"alvo":"delta","confianca":0.039,"epistemico":"fato","origem":"manual","rel":"referencia"}],"confianca":1.0,"contraexemplos":[],"descricao":"[Capacidade conversacional] Specialização de (Definição~): é o par usuário — resposta; c casa a fala com o lado esquerdo; vem da cadeia (fluxo:ID:NN:dep=k); observa estabilidade estrutural por turno (Seção~).","epistemico":"fato","exemplos":[],"id":"definition_capacidade_conversacionaldefkappa-tiffany_spe","memoria":"semantica","meta":{"dominio":"paper","nivel":5,"tipo_teorema":"definition"},"origem":"paper","palavras_chave":[],"sinonimos":[],"tipo":"conceito","titulo":"definition: [Capacidade conversacional]def:kappa-tiffany\n  Spe…"}
\section{definition: [Capacidade conversacional]def:kappa-tiffany
  Spe…}
[Capacidade conversacional] Specialização de (Definição\~{}): é o par usuário — resposta; c casa a fala com o lado esquerdo; vem da cadeia (fluxo:ID:NN:dep=k); observa estabilidade estrutural por turno (Seção\~{}).
\prov{relações: define plano\_no\_so; referencia kappa; referencia delta}
\prov{proveniência: paper \textperiodcentered{} fato \textperiodcentered{} confiança 1.0 \textperiodcentered{} domínio paper \textperiodcentered{} história 9 versões no jornal}

%CRISTAL {"arestas":[{"alvo":"plano_no_so","confianca":1.0,"epistemico":"fato","origem":"paper","rel":"define"},{"alvo":"kappa","confianca":0.016,"epistemico":"fato","origem":"manual","rel":"referencia"},{"alvo":"delta","confianca":0.031,"epistemico":"fato","origem":"manual","rel":"referencia"}],"confianca":1.0,"contraexemplos":[],"descricao":"[Capacidade] Uma capacidade é uma tupla [ =(c, o), ] onde é o estado, c a condição, o orçamento (Seção~), o peso (grant/deny quando aplicável) e o a origem. Nenhum domínio entra na definição.","epistemico":"fato","exemplos":[],"id":"definition_capacidadedefkappa-abstract_uma_capacidade_e_um","memoria":"semantica","meta":{"dominio":"paper","nivel":5,"tipo_teorema":"definition"},"origem":"paper","palavras_chave":[],"sinonimos":[],"tipo":"conceito","titulo":"definition: [Capacidade]def:kappa-abstract\nUma capacidade é um…"}
\section{definition: [Capacidade]def:kappa-abstract
Uma capacidade é um…}
[Capacidade] Uma capacidade é uma tupla [ =(c, o), ] onde é o estado, c a condição, o orçamento (Seção\~{}), o peso (grant/deny quando aplicável) e o a origem. Nenhum domínio entra na definição.
\prov{relações: define plano\_no\_so; referencia kappa; referencia delta}
\prov{proveniência: paper \textperiodcentered{} fato \textperiodcentered{} confiança 1.0 \textperiodcentered{} domínio paper \textperiodcentered{} história 9 versões no jornal}

%CRISTAL {"arestas":[{"alvo":"plano_no_so","confianca":1.0,"epistemico":"fato","origem":"paper","rel":"define"}],"confianca":1.0,"contraexemplos":[],"descricao":"[Capacidades detidas] Hold(p) é o menor conjunto que contém as capacidades intrínsecas de p (se p é ou dono) e toda tal que existe uma aresta válida u.","epistemico":"fato","exemplos":[],"id":"definition_capacidades_detidas_holdp_e_o_menor_conjunto","memoria":"semantica","meta":{"dominio":"paper","nivel":5,"tipo_teorema":"definition"},"origem":"paper","palavras_chave":[],"sinonimos":[],"tipo":"conceito","titulo":"Capacidades detidas"}
\section{Capacidades detidas}
[Capacidades detidas] Hold(p) é o menor conjunto que contém as capacidades intrínsecas de p (se p é ou dono) e toda tal que existe uma aresta válida u.
\prov{relações: define plano\_no\_so}
\prov{proveniência: paper \textperiodcentered{} fato \textperiodcentered{} confiança 1.0 \textperiodcentered{} domínio paper \textperiodcentered{} história 4 versões no jornal}

%CRISTAL {"arestas":[{"alvo":"plano_no_so","confianca":1.0,"epistemico":"fato","origem":"paper","rel":"define"}],"confianca":1.0,"contraexemplos":[],"descricao":"[Concessão válida / atenuação] A aresta u é válida se existe ₀(u) tal que [ c_₀, (s,a)₀(s,a), _=₀, ₀ 1, _ ₀-1, _₀, ] com a convenção 1=. Isto é, u só concede o que detém, nunca ampliando a condição, sempre decrementando a profundidade e nunca estendendo a validade (expiração, Seção~). Em particular ₀=0 proíbe qualquer concessão.","epistemico":"fato","exemplos":[],"id":"definition_concessao_valida_atenuacaodefatt_a_aresta_u","memoria":"semantica","meta":{"dominio":"paper","nivel":5,"tipo_teorema":"definition"},"origem":"paper","palavras_chave":[],"sinonimos":[],"tipo":"conceito","titulo":"Concessão válida / atenuação"}
\section{Concessão válida / atenuação}
[Concessão válida / atenuação] A aresta u é válida se existe \ensuremath{_{0}}(u) tal que [ c\_\ensuremath{_{0}}, (s,a)\ensuremath{_{0}}(s,a), \_=\ensuremath{_{0}}, \ensuremath{_{0}} 1, \_ \ensuremath{_{0}}-1, \_\ensuremath{_{0}}, ] com a convenção 1=. Isto é, u só concede o que detém, nunca ampliando a condição, sempre decrementando a profundidade e nunca estendendo a validade (expiração, Seção\~{}). Em particular \ensuremath{_{0}}=0 proíbe qualquer concessão.
\prov{relações: define plano\_no\_so}
\prov{proveniência: paper \textperiodcentered{} fato \textperiodcentered{} confiança 1.0 \textperiodcentered{} domínio paper \textperiodcentered{} história 4 versões no jornal}

%CRISTAL {"arestas":[{"alvo":"plano_no_so","confianca":1.0,"epistemico":"fato","origem":"paper","rel":"define"},{"alvo":"colapso","confianca":-0.055,"epistemico":"fato","origem":"manual","rel":"referencia"}],"confianca":1.0,"contraexemplos":[],"descricao":"[Conjunto vigente e colapso ] Seja =ᵢ capacidades detidas, filtradas por vigência e por cᵢ para consulta q. Defina [ _=:_>0 ou avaliação local, (q)=(_,q) =score(q) ] (quando o máximo existe; senão (q)= --- fail-closed). é independente de restante na leitura: só entra na formação de via delegação prévia.","epistemico":"fato","exemplos":[],"id":"definition_conjunto_vigente_e_colapso_defcollapse_seja","memoria":"semantica","meta":{"dominio":"paper","nivel":5,"tipo_teorema":"definition"},"origem":"paper","palavras_chave":[],"sinonimos":[],"tipo":"conceito","titulo":"Conjunto vigente e colapso"}
\section{Conjunto vigente e colapso}
[Conjunto vigente e colapso ] Seja =\ensuremath{_{i}} capacidades detidas, filtradas por vigência e por c\ensuremath{_{i}} para consulta q. Defina [ \_=:\_>0 ou avaliação local, (q)=(\_,q) =score(q) ] (quando o máximo existe; senão (q)= --- fail-closed). é independente de restante na leitura: só entra na formação de via delegação prévia.
\prov{relações: define plano\_no\_so; referencia colapso}
\prov{proveniência: paper \textperiodcentered{} fato \textperiodcentered{} confiança 1.0 \textperiodcentered{} domínio paper \textperiodcentered{} história 7 versões no jornal}

%CRISTAL {"arestas":[{"alvo":"plano_no_so","confianca":1.0,"epistemico":"fato","origem":"paper","rel":"define"}],"confianca":1.0,"contraexemplos":[],"descricao":"[Delegação é um operador monótono deflacionário] Cada aresta de delegação induz T:22 com [ x y T(x) T(y)(monótono), T(x) x(deflação: T). ] A condição de atenuação c_₀ (Definição~) é exatamente T.","epistemico":"fato","exemplos":[],"id":"definition_delegacao_e_um_operador_monotono_deflacionario_c","memoria":"semantica","meta":{"dominio":"paper","nivel":5,"tipo_teorema":"definition"},"origem":"paper","palavras_chave":[],"sinonimos":[],"tipo":"conceito","titulo":"definition: [Delegação é um operador monótono deflacionário]\nC…"}
\section{definition: [Delegação é um operador monótono deflacionário]
C…}
[Delegação é um operador monótono deflacionário] Cada aresta de delegação induz T:22 com [ x y T(x) T(y)(monótono), T(x) x(deflação: T). ] A condição de atenuação c\_\ensuremath{_{0}} (Definição\~{}) é exatamente T.
\prov{relações: define plano\_no\_so}
\prov{proveniência: paper \textperiodcentered{} fato \textperiodcentered{} confiança 1.0 \textperiodcentered{} domínio paper \textperiodcentered{} história 3 versões no jornal}

%CRISTAL {"arestas":[{"alvo":"colapso","confianca":0.95,"epistemico":"fato","origem":"paper","rel":"parte_de"}],"confianca":1.0,"contraexemplos":[],"descricao":"[Forma contínua] Para x L²(>₀,dt/t), [ [x]()=₀^ x(t),t⁻i,dtt =-^ x(eu),e⁻i u,du =()[x](), ] isto é, é Fourier no lado relativístico (a transformada de Mellin sobre a linha crítica). A saída [x]()=A()ei() é a função de onda: amplitude A e fase.","epistemico":"fato","exemplos":[],"id":"definition_forma_continuadefcont_para_x_l2_0dtt","memoria":"semantica","meta":{"dominio":"paper","nivel":5,"tipo_teorema":"definition","verificacao":"conceito novo"},"origem":"paper","palavras_chave":[],"sinonimos":[],"tipo":"conceito","titulo":"Forma contínua"}
\section{Forma contínua}
[Forma contínua] Para x L\ensuremath{^{2}}(>\ensuremath{_{0}},dt/t), [ [x]()=\ensuremath{_{0}}\^{} x(t),t\ensuremath{^{-}}i,dtt =-\^{} x(eu),e\ensuremath{^{-}}i u,du =()[x](), ] isto é, é Fourier no lado relativístico (a transformada de Mellin sobre a linha crítica). A saída [x]()=A()ei() é a função de onda: amplitude A e fase.
\prov{relações: parte\_de colapso}
\prov{proveniência: paper \textperiodcentered{} fato \textperiodcentered{} confiança 1.0 \textperiodcentered{} domínio paper \textperiodcentered{} história 5 versões no jornal}

%CRISTAL {"arestas":[{"alvo":"colapso","confianca":0.95,"epistemico":"fato","origem":"paper","rel":"parte_de"}],"confianca":1.0,"contraexemplos":[],"descricao":"[Forma discreta, sobre bits] Dado b0,1N (dado bruto), interpreta-se b como amostras na reta de ouro tₖ=k --- uniformes no lado relativístico (Prop.~) --- e define-se [b]=DFT(b). Os pares (Aj,j)=(|[b]j|,[b]j) formam a função de onda numérica pura: lista finita de números, sem forma fechada.","epistemico":"fato","exemplos":[],"id":"definition_forma_discreta_sobre_bitsdefdisc_dado_b01","memoria":"semantica","meta":{"dominio":"paper","nivel":5,"tipo_teorema":"definition","verificacao":"conceito novo"},"origem":"paper","palavras_chave":[],"sinonimos":[],"tipo":"conceito","titulo":"Forma discreta, sobre bits"}
\section{Forma discreta, sobre bits}
[Forma discreta, sobre bits] Dado b0,1N (dado bruto), interpreta-se b como amostras na reta de ouro t\ensuremath{_{k}}=k --- uniformes no lado relativístico (Prop.\~{}) --- e define-se [b]=DFT(b). Os pares (Aj,j)=(|[b]j|,[b]j) formam a função de onda numérica pura: lista finita de números, sem forma fechada.
\prov{relações: parte\_de colapso}
\prov{proveniência: paper \textperiodcentered{} fato \textperiodcentered{} confiança 1.0 \textperiodcentered{} domínio paper \textperiodcentered{} história 5 versões no jornal}

%CRISTAL {"arestas":[{"alvo":"plano_no_so","confianca":1.0,"epistemico":"fato","origem":"paper","rel":"define"}],"confianca":1.0,"contraexemplos":[],"descricao":"[Grafo de delegação] Um grafo de delegação é um grafo dirigido finito (não necessariamente acíclico) cujos vértices são principals e cujas arestas u representam “u concede a capacidade a w”. w é qualquer principal --- sem exigência de árvore, de vizinhança ou de aciclicidade.","epistemico":"fato","exemplos":[],"id":"definition_grafo_de_delegacaodefdag_um_grafo_de_delegacao","memoria":"semantica","meta":{"dominio":"paper","nivel":5,"tipo_teorema":"definition"},"origem":"paper","palavras_chave":[],"sinonimos":[],"tipo":"conceito","titulo":"definition: [Grafo de delegação]def:dag\nUm grafo de delegação…"}
\section{definition: [Grafo de delegação]def:dag
Um grafo de delegação…}
[Grafo de delegação] Um grafo de delegação é um grafo dirigido finito (não necessariamente acíclico) cujos vértices são principals e cujas arestas u representam “u concede a capacidade a w”. w é qualquer principal --- sem exigência de árvore, de vizinhança ou de aciclicidade.
\prov{relações: define plano\_no\_so}
\prov{proveniência: paper \textperiodcentered{} fato \textperiodcentered{} confiança 1.0 \textperiodcentered{} domínio paper \textperiodcentered{} história 3 versões no jornal}

%CRISTAL {"arestas":[{"alvo":"plano_no_so","confianca":1.0,"epistemico":"fato","origem":"paper","rel":"define"}],"confianca":1.0,"contraexemplos":[],"descricao":"[Instância de concessão] Uma instância é uma capacidade acrescida de (orig,par, ): a origem orig (quem concedeu), o pai par (a instância da qual foi atenuada, ou se intrínseca), a expiração >₀ e o estado active,suspended,revoked, expired. As arestas válidas (Definição~) já exigem _₀: um filho nunca sobrevive ao pai.","epistemico":"fato","exemplos":[],"id":"definition_instancia_de_concessao_uma_instancia_e_uma_capac","memoria":"semantica","meta":{"dominio":"paper","nivel":5,"tipo_teorema":"definition"},"origem":"paper","palavras_chave":[],"sinonimos":[],"tipo":"conceito","titulo":"definition: [Instância de concessão]\nUma instância é uma capac…"}
\section{definition: [Instância de concessão]
Uma instância é uma capac…}
[Instância de concessão] Uma instância é uma capacidade acrescida de (orig,par, ): a origem orig (quem concedeu), o pai par (a instância da qual foi atenuada, ou se intrínseca), a expiração >\ensuremath{_{0}} e o estado active,suspended,revoked, expired. As arestas válidas (Definição\~{}) já exigem \_\ensuremath{_{0}}: um filho nunca sobrevive ao pai.
\prov{relações: define plano\_no\_so}
\prov{proveniência: paper \textperiodcentered{} fato \textperiodcentered{} confiança 1.0 \textperiodcentered{} domínio paper \textperiodcentered{} história 3 versões no jornal}

%CRISTAL {"arestas":[{"alvo":"plano_no_so","confianca":1.0,"epistemico":"fato","origem":"paper","rel":"define"}],"confianca":1.0,"contraexemplos":[],"descricao":"[Linhas e condições] Seja o universo de todas as linhas. Cada r tem atributos (r., r.id, ). Uma condição c é um predicado c:, identificado ao seu extent c=r:c(r). As condições formam a álgebra booleana (2, ) --- as MongoQuery do CASL.","epistemico":"fato","exemplos":[],"id":"definition_linhas_e_condicoes_seja_o_universo_de_todas_a","memoria":"semantica","meta":{"dominio":"paper","nivel":5,"tipo_teorema":"definition"},"origem":"paper","palavras_chave":[],"sinonimos":[],"tipo":"conceito","titulo":"Linhas e condições"}
\section{Linhas e condições}
[Linhas e condições] Seja o universo de todas as linhas. Cada r tem atributos (r., r.id, ). Uma condição c é um predicado c:, identificado ao seu extent c=r:c(r). As condições formam a álgebra booleana (2, ) --- as MongoQuery do CASL.
\prov{relações: define plano\_no\_so}
\prov{proveniência: paper \textperiodcentered{} fato \textperiodcentered{} confiança 1.0 \textperiodcentered{} domínio paper \textperiodcentered{} história 4 versões no jornal}

%CRISTAL {"arestas":[{"alvo":"plano_no_so","confianca":1.0,"epistemico":"fato","origem":"paper","rel":"define"},{"alvo":"bai","confianca":-0.008,"epistemico":"fato","origem":"manual","rel":"referencia"}],"confianca":1.0,"contraexemplos":[],"descricao":"[Mapeamento Parte~IV] 1.25 tabular@l@l@ bai & conversa / linha r & nó do.tsv (impressão Koch no.idx) condição c, extent c & turno usuário — Tiffany; c casa a fala capacidade =(a,s,c, ) & par curado + profundidade de cadeia grant (=0) & nó dialogos / turno duplo autorizado proibição (=1) & nó genérico (Oi oi oi!, saudação solta) Hold(p) & candidatos promovidos + camada dialogos","epistemico":"fato","exemplos":[],"id":"definition_mapeamento_parteivdeftiffany-casl_center","memoria":"semantica","meta":{"dominio":"paper","nivel":5,"tipo_teorema":"definition"},"origem":"paper","palavras_chave":[],"sinonimos":[],"tipo":"conceito","titulo":"definition: [Mapeamento Parte~IV]def:tiffany-casl\n  center…"}
\section{definition: [Mapeamento Parte\~{}IV]def:tiffany-casl
  center…}
[Mapeamento Parte\~{}IV] 1.25 tabular@l@l@ bai \& conversa / linha r \& nó do.tsv (impressão Koch no.idx) condição c, extent c \& turno usuário — Tiffany; c casa a fala capacidade =(a,s,c, ) \& par curado + profundidade de cadeia grant (=0) \& nó dialogos / turno duplo autorizado proibição (=1) \& nó genérico (Oi oi oi!, saudação solta) Hold(p) \& candidatos promovidos + camada dialogos
\prov{relações: define plano\_no\_so; referencia bai}
\prov{proveniência: paper \textperiodcentered{} fato \textperiodcentered{} confiança 1.0 \textperiodcentered{} domínio paper \textperiodcentered{} história 6 versões no jornal}

%CRISTAL {"arestas":[{"alvo":"plano_no_so","confianca":1.0,"epistemico":"fato","origem":"paper","rel":"define"},{"alvo":"colapso","confianca":0.039,"epistemico":"fato","origem":"manual","rel":"referencia"}],"confianca":1.0,"contraexemplos":[],"descricao":"[Medida ] Fixe um espaço de estados S (configurações, impressões, snapshots). Uma medida de organização estrutural é um funcional [:S R ₀, ] monótono sob agregação observável: estados mais estruturados (menos dispersos, mais estáveis sob transformação) carregam mais informativo. não entra no colapso --- observa o efeito da propagação/atenuação, complementar a.","epistemico":"fato","exemplos":[],"id":"definition_medida_defphi-abstract_fixe_um_espaco_de_esta","memoria":"semantica","meta":{"dominio":"paper","nivel":5,"tipo_teorema":"definition"},"origem":"paper","palavras_chave":[],"sinonimos":[],"tipo":"conceito","titulo":"Medida"}
\section{Medida}
[Medida ] Fixe um espaço de estados S (configurações, impressões, snapshots). Uma medida de organização estrutural é um funcional [:S R \ensuremath{_{0}}, ] monótono sob agregação observável: estados mais estruturados (menos dispersos, mais estáveis sob transformação) carregam mais informativo. não entra no colapso --- observa o efeito da propagação/atenuação, complementar a.
\prov{relações: define plano\_no\_so; referencia colapso}
\prov{proveniência: paper \textperiodcentered{} fato \textperiodcentered{} confiança 1.0 \textperiodcentered{} domínio paper \textperiodcentered{} história 7 versões no jornal}

%CRISTAL {"arestas":[{"alvo":"plano_no_so","confianca":1.0,"epistemico":"fato","origem":"paper","rel":"define"},{"alvo":"reticulado","confianca":-0.031,"epistemico":"fato","origem":"manual","rel":"referencia"}],"confianca":1.0,"contraexemplos":[],"descricao":"[O dioid das condições] (2, ) com:= e:= é um semianel idempotente comutativo (dioid): é associativa, comutativa, idempotente (a a=a), com neutro; idem, com neutro; e distribui sobre. A ordem canônica a b a b=b coincide com. Com o complemento, é a álgebra booleana completa (2, ) --- um reticulado completo.","epistemico":"fato","exemplos":[],"id":"definition_o_dioid_das_condicoes_2_com_e","memoria":"semantica","meta":{"dominio":"paper","nivel":5,"tipo_teorema":"definition"},"origem":"paper","palavras_chave":[],"sinonimos":[],"tipo":"conceito","titulo":"O dioid das condições"}
\section{O dioid das condições}
[O dioid das condições] (2, ) com:= e:= é um semianel idempotente comutativo (dioid): é associativa, comutativa, idempotente (a a=a), com neutro; idem, com neutro; e distribui sobre. A ordem canônica a b a b=b coincide com. Com o complemento, é a álgebra booleana completa (2, ) --- um reticulado completo.
\prov{relações: define plano\_no\_so; referencia reticulado}
\prov{proveniência: paper \textperiodcentered{} fato \textperiodcentered{} confiança 1.0 \textperiodcentered{} domínio paper \textperiodcentered{} história 7 versões no jornal}

%CRISTAL {"arestas":[{"alvo":"plano_no_so","confianca":1.0,"epistemico":"fato","origem":"paper","rel":"define"},{"alvo":"delta","confianca":0.039,"epistemico":"fato","origem":"manual","rel":"referencia"},{"alvo":"colapso","confianca":0.008,"epistemico":"fato","origem":"manual","rel":"referencia"}],"confianca":1.0,"contraexemplos":[],"descricao":"[Orçamento de propagação] representa um orçamento discreto de propagação: quantas re-delegações (saltos, turnos, réplicas) a capacidade ainda admite. =0 --- usa mas não repassa; =k --- mais k saltos, decrementando; = --- fecho transitivo ( P). A verificação de ocorre no ato de delegar, nunca na avaliação/colapso.","epistemico":"fato","exemplos":[],"id":"definition_orcamento_de_propagacaodefdelta-budget_repre","memoria":"semantica","meta":{"dominio":"paper","nivel":5,"tipo_teorema":"definition"},"origem":"paper","palavras_chave":[],"sinonimos":[],"tipo":"conceito","titulo":"Orçamento de propagação"}
\section{Orçamento de propagação}
[Orçamento de propagação] representa um orçamento discreto de propagação: quantas re-delegações (saltos, turnos, réplicas) a capacidade ainda admite. =0 --- usa mas não repassa; =k --- mais k saltos, decrementando; = --- fecho transitivo ( P). A verificação de ocorre no ato de delegar, nunca na avaliação/colapso.
\prov{relações: define plano\_no\_so; referencia delta; referencia colapso}
\prov{proveniência: paper \textperiodcentered{} fato \textperiodcentered{} confiança 1.0 \textperiodcentered{} domínio paper \textperiodcentered{} história 10 versões no jornal}

%CRISTAL {"arestas":[{"alvo":"colapso","confianca":0.95,"epistemico":"fato","origem":"paper","rel":"parte_de"}],"confianca":1.0,"contraexemplos":[],"descricao":"[Os dois lados] O lado real é (+) com a métrica usual. O lado relativístico é (>₀,) com a medida de Haar multiplicativa d^ t=dt/t. A involução logarítmica (flip de Wick) [:L²!(>₀,dtt) L²(du), ( x)(u)=x(eu), ] é o isomorfismo entre os dois lados, com inversa (⁻¹y)(t)=y( t).","epistemico":"fato","exemplos":[],"id":"definition_os_dois_lados_o_lado_real_e_com_a_metrica","memoria":"semantica","meta":{"dominio":"paper","nivel":5,"tipo_teorema":"definition","verificacao":"conceito novo"},"origem":"paper","palavras_chave":[],"sinonimos":[],"tipo":"conceito","titulo":"Os dois lados"}
\section{Os dois lados}
[Os dois lados] O lado real é (+) com a métrica usual. O lado relativístico é (>\ensuremath{_{0}},) com a medida de Haar multiplicativa d\^{} t=dt/t. A involução logarítmica (flip de Wick) [:L\ensuremath{^{2}}!(>\ensuremath{_{0}},dtt) L\ensuremath{^{2}}(du), ( x)(u)=x(eu), ] é o isomorfismo entre os dois lados, com inversa (\ensuremath{^{-1}}y)(t)=y( t).
\prov{relações: parte\_de colapso}
\prov{proveniência: paper \textperiodcentered{} fato \textperiodcentered{} confiança 1.0 \textperiodcentered{} domínio paper \textperiodcentered{} história 5 versões no jornal}

%CRISTAL {"arestas":[{"alvo":"plano_no_so","confianca":1.0,"epistemico":"fato","origem":"paper","rel":"define"},{"alvo":"delta","confianca":-0.023,"epistemico":"fato","origem":"manual","rel":"referencia"}],"confianca":1.0,"contraexemplos":[],"descricao":"[Par ()] controla propagação (alcance, saltos restantes). controla estabilidade observável (organização estrutural do estado). Diálogos ou delegações longas consomem e, em geral, aumentam --- mais estrutura acumulada --- sem violar atenuação.","epistemico":"fato","exemplos":[],"id":"definition_par_defphi-delta_controla_propagacao_a","memoria":"semantica","meta":{"dominio":"paper","nivel":5,"tipo_teorema":"definition"},"origem":"paper","palavras_chave":[],"sinonimos":[],"tipo":"conceito","titulo":"Par ()"}
\section{Par ()}
[Par ()] controla propagação (alcance, saltos restantes). controla estabilidade observável (organização estrutural do estado). Diálogos ou delegações longas consomem e, em geral, aumentam --- mais estrutura acumulada --- sem violar atenuação.
\prov{relações: define plano\_no\_so; referencia delta}
\prov{proveniência: paper \textperiodcentered{} fato \textperiodcentered{} confiança 1.0 \textperiodcentered{} domínio paper \textperiodcentered{} história 7 versões no jornal}

%CRISTAL {"arestas":[{"alvo":"colapso","confianca":0.95,"epistemico":"fato","origem":"paper","rel":"parte_de"},{"alvo":"word","confianca":0.5,"epistemico":"fato","origem":"wiki","rel":"relaciona"},{"alvo":"vontade_de_poder","confianca":0.5,"epistemico":"fato","origem":"wiki","rel":"relaciona"}],"confianca":1.0,"contraexemplos":[],"descricao":"[Reta de ouro] Seja =1+52. A constante do vazio é =¹/=1,346361 A reta de ouro é a grade geométrica kₖ>₀.","epistemico":"fato","exemplos":[],"id":"definition_reta_de_ouro_seja_152_a_constante_do_vazi","memoria":"semantica","meta":{"dominio":"paper","nivel":5,"tipo_teorema":"definition","verificacao":"conceito novo"},"origem":"paper","palavras_chave":[],"sinonimos":[],"tipo":"conceito","titulo":"Reta de ouro"}
\section{Reta de ouro}
[Reta de ouro] Seja =1+52. A constante do vazio é =\ensuremath{^{1}}/=1,346361 A reta de ouro é a grade geométrica k\ensuremath{_{k}}>\ensuremath{_{0}}.
\prov{relações: parte\_de colapso; relaciona word; relaciona vontade\_de\_poder}
\prov{proveniência: paper \textperiodcentered{} fato \textperiodcentered{} confiança 1.0 \textperiodcentered{} domínio paper \textperiodcentered{} história 7 versões no jornal}

%CRISTAL {"arestas":[{"alvo":"plano_no_so","confianca":1.0,"epistemico":"fato","origem":"paper","rel":"define"},{"alvo":"reticulado","confianca":0.047,"epistemico":"fato","origem":"manual","rel":"referencia"}],"confianca":1.0,"contraexemplos":[],"descricao":"[Universo e extent] Fixe um universo X. Uma condição c é um predicado c:X, identificado ao extent c X. As condições formam o reticulado booleano (2X, X).","epistemico":"fato","exemplos":[],"id":"definition_universo_e_extentdefuniverso_fixe_um_universo","memoria":"semantica","meta":{"dominio":"paper","nivel":5,"tipo_teorema":"definition"},"origem":"paper","palavras_chave":[],"sinonimos":[],"tipo":"conceito","titulo":"Universo e extent"}
\section{Universo e extent}
[Universo e extent] Fixe um universo X. Uma condição c é um predicado c:X, identificado ao extent c X. As condições formam o reticulado booleano (2X, X).
\prov{relações: define plano\_no\_so; referencia reticulado}
\prov{proveniência: paper \textperiodcentered{} fato \textperiodcentered{} confiança 1.0 \textperiodcentered{} domínio paper \textperiodcentered{} história 7 versões no jornal}

%CRISTAL {"arestas":[{"alvo":"broca_os","confianca":0.95,"epistemico":"fato","origem":"paper","rel":"parte_de"}],"confianca":1.0,"contraexemplos":[],"descricao":"D Não implementar tudo de uma vez. Ordem incremental --- engine congelada; cada fase entrega\nartefato testável.\n\ncenter{1.45}\ntabular{@{}clll@{}}\n &  & ério de done &  \\\\\n\n3.0 & Este documento + fase3.tex & arquitetura aprovada & intacta \\\\\n3.1 & Formato meta.json, manifest, chunk\\_id & V tiffany-home/ & intacta \\\\\n3.2 & tiffany-ingest MVP (.md, .tex) & V chunks $$ idx & intacta \\\\\n3.3 & Core P0: git, shell, python, markdown & V gold 85\\% & intacta \\\\\n3.4 & Orquestrador multi-.idx + fusão & V or","epistemico":"fato","exemplos":[],"id":"dependencias_entre_fases","memoria":"semantica","meta":{"dominio":"paper","nivel":5,"verificacao":"overlap moderado — revisar manualmente"},"origem":"paper","palavras_chave":[],"sinonimos":[],"tipo":"conceito","titulo":"Dependências entre fases"}
\section{Dependências entre fases}
D Não implementar tudo de uma vez. Ordem incremental --- engine congelada; cada fase entrega
artefato testável.

center\{1.45\}
tabular\{@\{\}clll@\{\}\}
 \&  \& ério de done \&  \textbackslash{}\textbackslash{}

3.0 \& Este documento + fase3.tex \& arquitetura aprovada \& intacta \textbackslash{}\textbackslash{}
3.1 \& Formato meta.json, manifest, chunk\textbackslash{}\_id \& V tiffany-home/ \& intacta \textbackslash{}\textbackslash{}
3.2 \& tiffany-ingest MVP (.md, .tex) \& V chunks \$\$ idx \& intacta \textbackslash{}\textbackslash{}
3.3 \& Core P0: git, shell, python, markdown \& V gold 85\textbackslash{}\% \& intacta \textbackslash{}\textbackslash{}
3.4 \& Orquestrador multi-.idx + fusão \& V or
\prov{relações: parte\_de broca\_os}
\prov{proveniência: paper \textperiodcentered{} fato \textperiodcentered{} confiança 1.0 \textperiodcentered{} domínio paper \textperiodcentered{} história 3 versões no jornal}

%CRISTAL {"arestas":[{"alvo":"transformada_dourada","confianca":1.0,"epistemico":"fato","origem":"paper","rel":"parte_de"}],"confianca":1.0,"contraexemplos":[],"descricao":"proposition[Auto-calibração] Dado um bloco b, seja (r)=|u²,r b|² a energia de curvatura da imagem relativística amostrada em escala rk. A escala intrínseca é r_*=r(r), determinada pelos próprios bits; a reta de referência é r_*k. Para dados auto-similares de razão, r_*=; quando a auto-similaridade é a do vazio, r_*=. proposition proof[Esboço] (r)0, nulo sse r b é afim (Teorema~); minimizar seleciona a escala em que o dado é mais","epistemico":"fato","exemplos":[],"id":"dicionario_operacional_tudo_que_fourier_tem","memoria":"semantica","meta":{"dominio":"paper","nivel":5,"verificacao":"conceito novo"},"origem":"paper","palavras_chave":[],"sinonimos":[],"tipo":"conceito","titulo":"Dicionário operacional: tudo que Fourier tem"}
\section{Dicionário operacional: tudo que Fourier tem}
proposition[Auto-calibração] Dado um bloco b, seja (r)=|u\ensuremath{^{2}},r b|\ensuremath{^{2}} a energia de curvatura da imagem relativística amostrada em escala rk. A escala intrínseca é r\_*=r(r), determinada pelos próprios bits; a reta de referência é r\_*k. Para dados auto-similares de razão, r\_*=; quando a auto-similaridade é a do vazio, r\_*=. proposition proof[Esboço] (r)0, nulo sse r b é afim (Teorema\~{}); minimizar seleciona a escala em que o dado é mais
\prov{relações: parte\_de transformada\_dourada}
\prov{proveniência: paper \textperiodcentered{} fato \textperiodcentered{} confiança 1.0 \textperiodcentered{} domínio paper \textperiodcentered{} história 3 versões no jornal}

%CRISTAL {"arestas":[{"alvo":"ponte_simbolica_gentil","confianca":1.0,"epistemico":"fato","origem":"paper","rel":"parte_de"},{"alvo":"reais","confianca":1.0,"epistemico":"fato","origem":"paper","rel":"explica"},{"alvo":"paper_45_confluencia_ouro","confianca":1.0,"epistemico":"fato","origem":"paper","rel":"referencia"},{"alvo":"corpo_tropical","confianca":1.0,"epistemico":"fato","origem":"paper","rel":"confirma"}],"confianca":1.0,"contraexemplos":[],"descricao":"Reta de ouro --- chart Bergman Seja =1+52. A reta de ouro (Paper~II) é a família de palavras 0,1 sem bloco 11, decodificada por D_(d)=ᵢ dᵢ⁻i, densa em [0,1]. A fração contínua =[1;1,1,] tem convergentes Fₙ+₁/Fₙ. definicao Transformada dourada = com:x(t) x(eu): Mellin na linha crítica / Fourier no log. Sobre bits: [b]=DFT na grade k; saída (Aj,j), reversível por Parseval. definicao observacao[Confluência","epistemico":"fato","exemplos":[],"id":"dois_charts_um_objeto_mapeamentos","memoria":"semantica","meta":{"dominio":"paper","nivel":5,"verificacao":"slug idêntico — enriquecer, não duplicar"},"origem":"paper","palavras_chave":[],"sinonimos":[],"tipo":"conceito","titulo":"Dois charts, um objeto: mapeamentos"}
\section{Dois charts, um objeto: mapeamentos}
Reta de ouro --- chart Bergman Seja =1+52. A reta de ouro (Paper\~{}II) é a família de palavras 0,1 sem bloco 11, decodificada por D\_(d)=\ensuremath{_{i}} d\ensuremath{_{i}}\ensuremath{^{-}}i, densa em [0,1]. A fração contínua =[1;1,1,] tem convergentes F\ensuremath{_{n}}+\ensuremath{_{1}}/F\ensuremath{_{n}}. definicao Transformada dourada = com:x(t) x(eu): Mellin na linha crítica / Fourier no log. Sobre bits: [b]=DFT na grade k; saída (Aj,j), reversível por Parseval. definicao observacao[Confluência
\prov{relações: parte\_de ponte\_simbolica\_gentil; explica reais; referencia paper\_45\_confluencia\_ouro; confirma corpo\_tropical}
\prov{proveniência: paper \textperiodcentered{} fato \textperiodcentered{} confiança 1.0 \textperiodcentered{} domínio paper \textperiodcentered{} história 124 versões no jornal}

%CRISTAL {"arestas":[{"alvo":"broca-so_o_gato_no_metal","confianca":1.0,"epistemico":"fato","origem":"README.md","rel":"parte_de"},{"alvo":"broca_os","confianca":1.0,"epistemico":"fato","origem":"manual","rel":"parte_de"},{"alvo":"vida-morte","confianca":0.5,"epistemico":"fato","origem":"wiki","rel":"relaciona"},{"alvo":"telemetria_shadow_producao","confianca":0.5,"epistemico":"fato","origem":"wiki","rel":"relaciona"}],"confianca":1.0,"contraexemplos":[],"descricao":"I Distribuição pensada para instalação, alimentação e atualização.","epistemico":"fato","exemplos":[],"id":"dois_pacotes_pow_lucro","memoria":"semantica","meta":{"dominio":"paper","duplicata_sugerida":[],"nivel":5,"verificacao":"parte de conceito mais amplo 'dois_pacotes_pow_lucro'"},"origem":"paper","palavras_chave":[],"sinonimos":[],"tipo":"conceito","titulo":"Dois pacotes (PoW / lucro)"}
\section{Dois pacotes (PoW / lucro)}
I Distribuição pensada para instalação, alimentação e atualização.
\prov{relações: parte\_de broca-so\_o\_gato\_no\_metal; parte\_de broca\_os; relaciona vida-morte; relaciona telemetria\_shadow\_producao}
\prov{proveniência: paper \textperiodcentered{} fato \textperiodcentered{} confiança 1.0 \textperiodcentered{} domínio paper \textperiodcentered{} história 11 versões no jornal}

%CRISTAL {"arestas":[{"alvo":"mecanica_quantica","confianca":0.95,"epistemico":"fato","origem":"paper","rel":"parte_de"}],"confianca":1.0,"contraexemplos":[],"descricao":"Da conservação a+c²=1 (Prop.~, com c=|s|) e da direção Tsirelson--Bell apontada na Secão~ segue, desenvolvida aqui, que |[R a,R a']|2c e que a soma CHSH de quatro termos de Tsirelson--Clauser--Horne--Shimony--Holt, que escrevemos Ss, obedece [ Ss 21+c²(s)=22-a(s), ] governada pela conservação a+c²=1; o máximo 22 é atingido exatamente no quaternion (a=0), e nos octônions o piso a a>0 dá Ss<22. Bell, lido. Isto é a hierarquia","epistemico":"fato","exemplos":[],"id":"dois_qubits_o_emaranhamento_e_o_associador_octonionico","memoria":"semantica","meta":{"dominio":"paper","nivel":5,"verificacao":"conceito novo"},"origem":"paper","palavras_chave":[],"sinonimos":[],"tipo":"conceito","titulo":"Dois qubits: o emaranhamento é o associador octoniônico"}
\section{Dois qubits: o emaranhamento é o associador octoniônico}
Da conservação a+c\ensuremath{^{2}}=1 (Prop.\~{}, com c=|s|) e da direção Tsirelson--Bell apontada na Secão\~{} segue, desenvolvida aqui, que |[R a,R a']|2c e que a soma CHSH de quatro termos de Tsirelson--Clauser--Horne--Shimony--Holt, que escrevemos Ss, obedece [ Ss 21+c\ensuremath{^{2}}(s)=22-a(s), ] governada pela conservação a+c\ensuremath{^{2}}=1; o máximo 22 é atingido exatamente no quaternion (a=0), e nos octônions o piso a a>0 dá Ss<22. Bell, lido. Isto é a hierarquia
\prov{relações: parte\_de mecanica\_quantica}
\prov{proveniência: paper \textperiodcentered{} fato \textperiodcentered{} confiança 1.0 \textperiodcentered{} domínio paper \textperiodcentered{} história 4 versões no jornal}

%CRISTAL {"arestas":[{"alvo":"floco_de_neve","confianca":1.0,"epistemico":"fato","origem":"paper","rel":"parte_de"}],"confianca":1.0,"contraexemplos":[],"descricao":"empty !85Entre a reta real e a reta de ouro cai um floco: o que a passagem dissipou e que só o Koch guarda. Sem o floco não há volta; com a reta associativa, reconstrói-se o canónico. abstract Chamamos floco () o objecto que armazena a informação não-associativa produzida pela dissipação na passagem entre retas associativas --- cristalização do de Koch no canal negro Peirce. A reta real (C₃ em s) e a reta de ouro (D_, Fibonacci) endereçam onde o estado vive; gu","epistemico":"fato","exemplos":[],"id":"duas_retas_uma_passagem","memoria":"semantica","meta":{"dominio":"paper","nivel":5,"verificacao":"overlap moderado — revisar manualmente"},"origem":"paper","palavras_chave":[],"sinonimos":[],"tipo":"conceito","titulo":"Duas retas, uma passagem"}
\section{Duas retas, uma passagem}
empty !85Entre a reta real e a reta de ouro cai um floco: o que a passagem dissipou e que só o Koch guarda. Sem o floco não há volta; com a reta associativa, reconstrói-se o canónico. abstract Chamamos floco () o objecto que armazena a informação não-associativa produzida pela dissipação na passagem entre retas associativas --- cristalização do de Koch no canal negro Peirce. A reta real (C\ensuremath{_{3}} em s) e a reta de ouro (D\_, Fibonacci) endereçam onde o estado vive; gu
\prov{relações: parte\_de floco\_de\_neve}
\prov{proveniência: paper \textperiodcentered{} fato \textperiodcentered{} confiança 1.0 \textperiodcentered{} domínio paper \textperiodcentered{} história 3 versões no jornal}

%CRISTAL {"arestas":[{"alvo":"casl-propagation","confianca":-0.008,"epistemico":"fato","origem":"manual","rel":"referencia"}],"confianca":1.0,"contraexemplos":[],"descricao":"Uma concessão não é eterna. Modelamos o seu ciclo de vida como uma máquina de estados finita (FSM), separada da álgebra de avaliação: a álgebra diz o que uma capacidade autoriza; a FSM diz se ela está em vigor. definition[Instância de concessão] Uma instância é uma capacidade acrescida de (orig,par, ): a origem orig (quem concedeu), o pai par (a instância da qual foi atenuada, ou se intrínseca), a expiração >₀ e o estado active,suspended,revoked, expired\\","epistemico":"fato","exemplos":[],"id":"emphenforcement","memoria":"semantica","meta":{"dominio":"paper","nivel":5},"origem":"paper","palavras_chave":[],"sinonimos":[],"tipo":"conceito","titulo":"Enforcement"}
\section{Enforcement}
Uma concessão não é eterna. Modelamos o seu ciclo de vida como uma máquina de estados finita (FSM), separada da álgebra de avaliação: a álgebra diz o que uma capacidade autoriza; a FSM diz se ela está em vigor. definition[Instância de concessão] Uma instância é uma capacidade acrescida de (orig,par, ): a origem orig (quem concedeu), o pai par (a instância da qual foi atenuada, ou se intrínseca), a expiração >\ensuremath{_{0}} e o estado active,suspended,revoked, expired\textbackslash{}
\prov{relações: referencia casl-propagation}
\prov{proveniência: paper \textperiodcentered{} fato \textperiodcentered{} confiança 1.0 \textperiodcentered{} domínio paper \textperiodcentered{} história 7 versões no jornal}

%CRISTAL {"arestas":[{"alvo":"leitura_cosmologica_a_energia_escura_e_o_universo_ds","confianca":1.0,"epistemico":"fato","origem":"paper","rel":"parte_de"}],"confianca":1.0,"contraexemplos":[],"descricao":"Antes de ler a expansão cósmica, calibramos o instrumento onde a resposta é conhecida. O método de Lopes lê o expoente de uma lei de potência y x^ por diferenças finitas em log--log, =( y)/( x), cancelando o prefator na primeira diferença; numa lei de potência pura, pares consecutivos já dão exato, sem binar. Por honestidade: o que segue é um dicionário estrutural, não uma predição contra a mecânica celeste estabelecida; o seu valor está na correspondência exata e em calibrar","epistemico":"fato","exemplos":[],"id":"energia_escura_o_campo_de_algebra","memoria":"semantica","meta":{"dominio":"paper","nivel":5,"verificacao":"overlap moderado — revisar manualmente"},"origem":"paper","palavras_chave":[],"sinonimos":[],"tipo":"conceito","titulo":"Energia escura: o campo de álgebra"}
\section{Energia escura: o campo de álgebra}
Antes de ler a expansão cósmica, calibramos o instrumento onde a resposta é conhecida. O método de Lopes lê o expoente de uma lei de potência y x\^{} por diferenças finitas em log--log, =( y)/( x), cancelando o prefator na primeira diferença; numa lei de potência pura, pares consecutivos já dão exato, sem binar. Por honestidade: o que segue é um dicionário estrutural, não uma predição contra a mecânica celeste estabelecida; o seu valor está na correspondência exata e em calibrar
\prov{relações: parte\_de leitura\_cosmologica\_a\_energia\_escura\_e\_o\_universo\_ds}
\prov{proveniência: paper \textperiodcentered{} fato \textperiodcentered{} confiança 1.0 \textperiodcentered{} domínio paper \textperiodcentered{} história 3 versões no jornal}

%CRISTAL {"arestas":[{"alvo":"gato","confianca":0.95,"epistemico":"fato","origem":"paper","rel":"parte_de"},{"alvo":"vida-morte","confianca":0.5,"epistemico":"fato","origem":"wiki","rel":"relaciona"},{"alvo":"poder_disciplinar_foucault","confianca":0.5,"epistemico":"fato","origem":"wiki","rel":"relaciona"},{"alvo":"virtualizacao","confianca":0.5,"epistemico":"fato","origem":"wiki","rel":"relaciona"}],"confianca":1.0,"contraexemplos":[],"descricao":"obs O metal sozinho não raciocina --- ressoa. A autonomia madura passa por duas disciplinas\n(Tiffany~§VI):","epistemico":"fato","exemplos":[],"id":"escola_---_engrossar_a_assinatura","memoria":"semantica","meta":{"dominio":"paper","nivel":5,"verificacao":"conceito novo"},"origem":"paper","palavras_chave":[],"sinonimos":[],"tipo":"conceito","titulo":"Escola --- engrossar a assinatura"}
\section{Escola --- engrossar a assinatura}
obs O metal sozinho não raciocina --- ressoa. A autonomia madura passa por duas disciplinas
(Tiffany\~{}§VI):
\prov{relações: parte\_de gato; relaciona vida-morte; relaciona poder\_disciplinar\_foucault; relaciona virtualizacao}
\prov{proveniência: paper \textperiodcentered{} fato \textperiodcentered{} confiança 1.0 \textperiodcentered{} domínio paper \textperiodcentered{} história 6 versões no jornal}

%CRISTAL {"arestas":[{"alvo":"leitura_das_unidades_o_sistema_por-unidade_e_o_si","confianca":1.0,"epistemico":"fato","origem":"paper","rel":"parte_de"}],"confianca":1.0,"contraexemplos":[],"descricao":"NI Estudo de caso, fora dos certificados centrais. Sequência de referência NC₀01422 (5386 bases, A,C,G,T=², Watson--Crick = complemento ). Com Hₖ=-wA,C,G,Tkp(w)₂ p(w), hₖ=Hₖ/k, a dimensão efetiva é eff(k)=hₖ/2; ela decresce com k (há estrutura), e a assinatura é o gap contra a média de muitos embaralhamentos (positivo em k=1!-!4; em k alto o viés de amostragem finita derruba ambos). Não sustenta as partes algébricas. Implementação: cadeiagenoma.py.","epistemico":"fato","exemplos":[],"id":"estatuto_e_conclusao","memoria":"semantica","meta":{"dominio":"paper","duplicata_sugerida":[],"nivel":5,"verificacao":"parte de conceito mais amplo 'estatuto_e_conclusao'"},"origem":"paper","palavras_chave":[],"sinonimos":[],"tipo":"conceito","titulo":"Estatuto e conclusão"}
\section{Estatuto e conclusão}
NI Estudo de caso, fora dos certificados centrais. Sequência de referência NC\ensuremath{_{0}}01422 (5386 bases, A,C,G,T=\ensuremath{^{2}}, Watson--Crick = complemento ). Com H\ensuremath{_{k}}=-wA,C,G,Tkp(w)\ensuremath{_{2}} p(w), h\ensuremath{_{k}}=H\ensuremath{_{k}}/k, a dimensão efetiva é eff(k)=h\ensuremath{_{k}}/2; ela decresce com k (há estrutura), e a assinatura é o gap contra a média de muitos embaralhamentos (positivo em k=1!-!4; em k alto o viés de amostragem finita derruba ambos). Não sustenta as partes algébricas. Implementação: cadeiagenoma.py.
\prov{relações: parte\_de leitura\_das\_unidades\_o\_sistema\_por-unidade\_e\_o\_si}
\prov{proveniência: paper \textperiodcentered{} fato \textperiodcentered{} confiança 1.0 \textperiodcentered{} domínio paper \textperiodcentered{} história 6 versões no jornal}

%CRISTAL {"arestas":[{"alvo":"goldcoin","confianca":1.0,"epistemico":"fato","origem":"paper","rel":"parte_de"}],"confianca":1.0,"contraexemplos":[],"descricao":"N Estado real, verificado na máquina nesta data: tabular@p0.34@1emp0.54@ ouroConta central ETH & 0x8C1C = 0,02013171 ETH; 102 transações on-chain (verificado em publicnode e drpc). É o gás. [6pt] ouroPell na conta da Claude & 10 moedas: 5 primos de Pell em claro 2,5,29,5741,33461 + 5 cifrados (ex.: 4125636888562548868221559797461449). [6pt] ouroPainel & 3 contratos ativos. [6pt] ouroLedger encadeado & 19 contratos / 19 certificados --- a cadeia íntegra","epistemico":"fato","exemplos":[],"id":"estatuto_e_escopo_honesto","memoria":"semantica","meta":{"dominio":"paper","duplicata_sugerida":[],"nivel":5,"verificacao":"parte de conceito mais amplo 'estatuto_e_escopo_honesto'"},"origem":"paper","palavras_chave":[],"sinonimos":[],"tipo":"conceito","titulo":"Estatuto e escopo honesto"}
\section{Estatuto e escopo honesto}
N Estado real, verificado na máquina nesta data: tabular@p0.34@1emp0.54@ ouroConta central ETH \& 0x8C1C = 0,02013171 ETH; 102 transações on-chain (verificado em publicnode e drpc). É o gás. [6pt] ouroPell na conta da Claude \& 10 moedas: 5 primos de Pell em claro 2,5,29,5741,33461 + 5 cifrados (ex.: 4125636888562548868221559797461449). [6pt] ouroPainel \& 3 contratos ativos. [6pt] ouroLedger encadeado \& 19 contratos / 19 certificados --- a cadeia íntegra
\prov{relações: parte\_de goldcoin}
\prov{proveniência: paper \textperiodcentered{} fato \textperiodcentered{} confiança 1.0 \textperiodcentered{} domínio paper \textperiodcentered{} história 4 versões no jornal}

%CRISTAL {"arestas":[{"alvo":"mecanica_quantica","confianca":1.0,"epistemico":"fato","origem":"paper","rel":"parte_de"},{"alvo":"mecanica","confianca":1.0,"epistemico":"fato","origem":"paper","rel":"parte_de"}],"confianca":1.0,"contraexemplos":[],"descricao":"A evolução unitária é reversível (a=0); a medida não. Medir é o coarse-graining many-to-one --- o instante em que o defeito a volta a ser positivo, e com ele o calor (a segunda lei da Termodinâmica). Reversível e irreversível são os dois lados da parábola: a evolução vive em c²=1 (incerteza pura, a=0); a medida reintroduz a>0. A mecânica quântica e a termodinâmica não são duas teorias --- são os dois lados de a+c²=1.","epistemico":"fato","exemplos":[],"id":"estatuto_o_que_fecha_o_que_segue_para_o_proximo_nivel","memoria":"semantica","meta":{"dominio":"paper","nivel":5,"verificacao":"slug idêntico — enriquecer, não duplicar"},"origem":"paper","palavras_chave":[],"sinonimos":[],"tipo":"conceito","titulo":"Estatuto: o que fecha, o que segue para o próximo nível"}
\section{Estatuto: o que fecha, o que segue para o próximo nível}
A evolução unitária é reversível (a=0); a medida não. Medir é o coarse-graining many-to-one --- o instante em que o defeito a volta a ser positivo, e com ele o calor (a segunda lei da Termodinâmica). Reversível e irreversível são os dois lados da parábola: a evolução vive em c\ensuremath{^{2}}=1 (incerteza pura, a=0); a medida reintroduz a>0. A mecânica quântica e a termodinâmica não são duas teorias --- são os dois lados de a+c\ensuremath{^{2}}=1.
\prov{relações: parte\_de mecanica\_quantica; parte\_de mecanica}
\prov{proveniência: paper \textperiodcentered{} fato \textperiodcentered{} confiança 1.0 \textperiodcentered{} domínio paper \textperiodcentered{} história 16 versões no jornal}

%CRISTAL {"arestas":[{"alvo":"delta","confianca":0.85,"epistemico":"fato","origem":"manual","rel":"referencia"},{"alvo":"geometria","confianca":1.0,"epistemico":"fato","origem":"paper","rel":"parte_de"}],"confianca":1.0,"contraexemplos":[],"descricao":"Subir de dimensão exige separar o invariante da sujeira. O chicote não é um operador espectral adicional: é o procedimento de limpeza band-limited que conserva Q e descarta o resíduo não-estruturado. Concretamente: a agulha de Dirac pura tem banda infinita --- espalha energia acima do gap e invade o invariante (verificado: vazamento 0,99); a agulha de Perelman é a delta reconstruída só pelos modos até o gap --- band-limited, vazamento 10⁻³¹, e o gap (a geometria) é restaurado. É a cir","epistemico":"fato","exemplos":[],"id":"estatuto_o_que_se_reproduz_e_o_que_e_horizonte","memoria":"semantica","meta":{"dominio":"paper","nivel":5,"verificacao":"slug idêntico — enriquecer, não duplicar"},"origem":"paper","palavras_chave":[],"sinonimos":[],"tipo":"conceito","titulo":"Estatuto: o que se reproduz e o que é horizonte"}
\section{Estatuto: o que se reproduz e o que é horizonte}
Subir de dimensão exige separar o invariante da sujeira. O chicote não é um operador espectral adicional: é o procedimento de limpeza band-limited que conserva Q e descarta o resíduo não-estruturado. Concretamente: a agulha de Dirac pura tem banda infinita --- espalha energia acima do gap e invade o invariante (verificado: vazamento 0,99); a agulha de Perelman é a delta reconstruída só pelos modos até o gap --- band-limited, vazamento 10\ensuremath{^{-31}}, e o gap (a geometria) é restaurado. É a cir
\prov{relações: referencia delta; parte\_de geometria}
\prov{proveniência: paper \textperiodcentered{} fato \textperiodcentered{} confiança 1.0 \textperiodcentered{} domínio paper \textperiodcentered{} história 14 versões no jornal}

%CRISTAL {"arestas":[{"alvo":"termodinamica","confianca":1.0,"epistemico":"fato","origem":"paper","rel":"parte_de"}],"confianca":1.0,"contraexemplos":[],"descricao":"Há um único ponto em que a dissipação vira reversibilidade: o flip de Wick. O gerador dissipativo e⁻ t (a equação do calor) e o gerador unitário e⁻iHt (a equação de Schrödinger) são a mesma estrutura, vistas em assinaturas trocadas (t it). A termodinâmica é o regime real (e⁻ t); a mecânica quântica é a sua continuação imaginária (e⁻iHt). Este é o portal para o próximo texto da série --- e, fiéis à régua, ele é anunciado, não atravessado aqui.","epistemico":"fato","exemplos":[],"id":"estatuto_o_que_se_verifica_o_que_e_leitura","memoria":"semantica","meta":{"dominio":"paper","nivel":5,"verificacao":"slug idêntico — enriquecer, não duplicar"},"origem":"paper","palavras_chave":[],"sinonimos":[],"tipo":"conceito","titulo":"Estatuto: o que se verifica, o que é leitura"}
\section{Estatuto: o que se verifica, o que é leitura}
Há um único ponto em que a dissipação vira reversibilidade: o flip de Wick. O gerador dissipativo e\ensuremath{^{-}} t (a equação do calor) e o gerador unitário e\ensuremath{^{-}}iHt (a equação de Schrödinger) são a mesma estrutura, vistas em assinaturas trocadas (t it). A termodinâmica é o regime real (e\ensuremath{^{-}} t); a mecânica quântica é a sua continuação imaginária (e\ensuremath{^{-}}iHt). Este é o portal para o próximo texto da série --- e, fiéis à régua, ele é anunciado, não atravessado aqui.
\prov{relações: parte\_de termodinamica}
\prov{proveniência: paper \textperiodcentered{} fato \textperiodcentered{} confiança 1.0 \textperiodcentered{} domínio paper \textperiodcentered{} história 10 versões no jornal}

%CRISTAL {"arestas":[{"alvo":"pipeline_de_ingestao","confianca":1.0,"epistemico":"fato","origem":"paper","rel":"parte_de"}],"confianca":1.0,"contraexemplos":[],"descricao":"1.3 tabular@ll@ &.md,.txt & cabeçalhos ATX / blank-line.tex & section, subsection.pdf & pdftotext + heurística de título.py,.c,.rs & AST leve ou bloco docstring + função.html & tags h1..h6, p, pre tabular Saída comum: lista chunkᵢd, lang, text, origem, nivel.","epistemico":"fato","exemplos":[],"id":"estrategias_de_chunking","memoria":"semantica","meta":{"dominio":"paper","nivel":5,"verificacao":"overlap moderado — revisar manualmente"},"origem":"paper","palavras_chave":[],"sinonimos":[],"tipo":"conceito","titulo":"Estratégias de chunking"}
\section{Estratégias de chunking}
1.3 tabular@ll@ \&.md,.txt \& cabeçalhos ATX / blank-line.tex \& section, subsection.pdf \& pdftotext + heurística de título.py,.c,.rs \& AST leve ou bloco docstring + função.html \& tags h1..h6, p, pre tabular Saída comum: lista chunk\ensuremath{_{i}}d, lang, text, origem, nivel.
\prov{relações: parte\_de pipeline\_de\_ingestao}
\prov{proveniência: paper \textperiodcentered{} fato \textperiodcentered{} confiança 1.0 \textperiodcentered{} domínio paper \textperiodcentered{} história 3 versões no jornal}

%CRISTAL {"arestas":[{"alvo":"teoremas_fundamentais","confianca":1.0,"epistemico":"fato","origem":"paper","rel":"parte_de"},{"alvo":"transcricao","confianca":0.5,"epistemico":"fato","origem":"wiki","rel":"relaciona"},{"alvo":"vida-morte","confianca":0.5,"epistemico":"fato","origem":"wiki","rel":"relaciona"},{"alvo":"parabola_pa","confianca":0.5,"epistemico":"fato","origem":"wiki","rel":"relaciona"},{"alvo":"resposta_a_pergunta_dificil","confianca":0.5,"epistemico":"fato","origem":"wiki","rel":"relaciona"},{"alvo":"topologia","confianca":0.5,"epistemico":"fato","origem":"wiki","rel":"relaciona"},{"alvo":"pow_metal_setor","confianca":0.5,"epistemico":"fato","origem":"wiki","rel":"relaciona"},{"alvo":"pow_geometria","confianca":0.5,"epistemico":"fato","origem":"wiki","rel":"relaciona"},{"alvo":"teste_ordem","confianca":0.5,"epistemico":"fato","origem":"wiki","rel":"relaciona"},{"alvo":"geometria_hub_broca","confianca":0.5,"epistemico":"fato","origem":"wiki","rel":"relaciona"},{"alvo":"por_sessao_conversadadostelemetriasessaojsonl","confianca":0.5,"epistemico":"fato","origem":"wiki","rel":"relaciona"},{"alvo":"motivos","confianca":0.5,"epistemico":"fato","origem":"wiki","rel":"relaciona"},{"alvo":"thread","confianca":0.5,"epistemico":"fato","origem":"wiki","rel":"relaciona"},{"alvo":"numa","confianca":0.5,"epistemico":"fato","origem":"wiki","rel":"relaciona"},{"alvo":"opcao_a_tarball_release","confianca":0.5,"epistemico":"fato","origem":"wiki","rel":"relaciona"},{"alvo":"pilha","confianca":0.5,"epistemico":"fato","origem":"wiki","rel":"relaciona"},{"alvo":"koch_cantor_fib_prova","confianca":0.5,"epistemico":"fato","origem":"wiki","rel":"relaciona"},{"alvo":"parabola_destruir","confianca":0.5,"epistemico":"fato","origem":"wiki","rel":"relaciona"},{"alvo":"tiffany_cabeca_broca","confianca":0.5,"epistemico":"fato","origem":"wiki","rel":"relaciona"},{"alvo":"projeto_cristaljson","confianca":0.5,"epistemico":"fato","origem":"wiki","rel":"relaciona"}],"confianca":1.0,"contraexemplos":[],"descricao":"Secção de transformada_dourada.tex.","epistemico":"fato","exemplos":[],"id":"exatidao","memoria":"semantica","meta":{"dominio":"paper","nivel":5,"verificacao":"conceito novo"},"origem":"paper","palavras_chave":[],"sinonimos":[],"tipo":"conceito","titulo":"Exatidão"}
\section{Exatidão}
Secção de transformada\_dourada.tex.
\prov{relações: parte\_de teoremas\_fundamentais; relaciona transcricao; relaciona vida-morte; relaciona parabola\_pa; relaciona resposta\_a\_pergunta\_dificil; relaciona topologia; relaciona pow\_metal\_setor; relaciona pow\_geometria; relaciona teste\_ordem; relaciona geometria\_hub\_broca; relaciona por\_sessao\_conversadadostelemetriasessaojsonl; relaciona motivos; relaciona thread; relaciona numa; relaciona opcao\_a\_tarball\_release; relaciona pilha; relaciona koch\_cantor\_fib\_prova; relaciona parabola\_destruir; relaciona tiffany\_cabeca\_broca; relaciona projeto\_cristaljson}
\prov{proveniência: paper \textperiodcentered{} fato \textperiodcentered{} confiança 1.0 \textperiodcentered{} domínio paper \textperiodcentered{} história 22 versões no jornal}

%CRISTAL {"arestas":[{"alvo":"assistente","confianca":1.0,"epistemico":"fato","origem":"paper","rel":"parte_de"}],"confianca":1.0,"contraexemplos":[],"descricao":"D Duas coisas distintas --- não misturar: Arquitetura: roteamento de tarefas ID O orquestrador só encaminha pedidos a agentes especializados. Isso é desenho de sistema --- reduz monólito, responsabilidades claras, alinhado ao lobo frontal~lobo e ao (op canal)~brocaos. Sólido como arquitetura. prop Mecanismo: ressonância como classificador obs Hipótese a validar: classificar intenção por colapso Koch/Hopfield sobre tecido de ações (conversa / código / paper / execução). Plausív","epistemico":"fato","exemplos":[],"id":"executor_agentes_e_avaliador","memoria":"semantica","meta":{"dominio":"paper","nivel":5,"verificacao":"conceito novo"},"origem":"paper","palavras_chave":[],"sinonimos":[],"tipo":"conceito","titulo":"Executor, agentes e avaliador"}
\section{Executor, agentes e avaliador}
D Duas coisas distintas --- não misturar: Arquitetura: roteamento de tarefas ID O orquestrador só encaminha pedidos a agentes especializados. Isso é desenho de sistema --- reduz monólito, responsabilidades claras, alinhado ao lobo frontal\~{}lobo e ao (op canal)\~{}brocaos. Sólido como arquitetura. prop Mecanismo: ressonância como classificador obs Hipótese a validar: classificar intenção por colapso Koch/Hopfield sobre tecido de ações (conversa / código / paper / execução). Plausív
\prov{relações: parte\_de assistente}
\prov{proveniência: paper \textperiodcentered{} fato \textperiodcentered{} confiança 1.0 \textperiodcentered{} domínio paper \textperiodcentered{} história 3 versões no jornal}

%CRISTAL {"arestas":[{"alvo":"main","confianca":1.0,"epistemico":"fato","origem":"paper","rel":"parte_de"},{"alvo":"beta_numeracao_metalica","confianca":1.0,"epistemico":"fato","origem":"wiki","rel":"especializa"}],"confianca":1.0,"contraexemplos":[],"descricao":"Considere T(x)= x1 sobre [0,1), com dígito d= x 0,1. É não-escapante (toda órbita permanece no compacto [0,1)) e admite codificação simbólica não-trivial, com dois ramos: o dígito 1 no ramo curto [1/,1) (“transmite”), o 0 no ramo longo [0,1/) (“recircula”). a gramática emerge da geometria D A imagem do ramo 1 é T[1/,1)=[0,-1)=[0,1/), onde só cabe o dígito 0. Logo, após um 1, o próximo é obrigatoriamente 0: a exclusão de 11 decorre d","epistemico":"fato","exemplos":[],"id":"familia_m-bonacci_gap_k","memoria":"semantica","meta":{"dominio":"paper","nivel":5,"verificacao":"slug idêntico — enriquecer, não duplicar"},"origem":"paper","palavras_chave":[],"sinonimos":[],"tipo":"conceito","titulo":"Família m-bonacci / gap k"}
\section{Família m-bonacci / gap k}
Considere T(x)= x1 sobre [0,1), com dígito d= x 0,1. É não-escapante (toda órbita permanece no compacto [0,1)) e admite codificação simbólica não-trivial, com dois ramos: o dígito 1 no ramo curto [1/,1) (“transmite”), o 0 no ramo longo [0,1/) (“recircula”). a gramática emerge da geometria D A imagem do ramo 1 é T[1/,1)=[0,-1)=[0,1/), onde só cabe o dígito 0. Logo, após um 1, o próximo é obrigatoriamente 0: a exclusão de 11 decorre d
\prov{relações: parte\_de main; especializa beta\_numeracao\_metalica}
\prov{proveniência: paper \textperiodcentered{} fato \textperiodcentered{} confiança 1.0 \textperiodcentered{} domínio paper \textperiodcentered{} história 7 versões no jornal}

%CRISTAL {"arestas":[{"alvo":"transcricao","confianca":0.5,"epistemico":"fato","origem":"wiki","rel":"relaciona"},{"alvo":"pow_gato","confianca":0.5,"epistemico":"fato","origem":"wiki","rel":"relaciona"},{"alvo":"fechamento_blocos","confianca":0.5,"epistemico":"fato","origem":"wiki","rel":"relaciona"},{"alvo":"prova_isa","confianca":0.5,"epistemico":"fato","origem":"wiki","rel":"relaciona"},{"alvo":"teste_ordem","confianca":0.5,"epistemico":"fato","origem":"wiki","rel":"relaciona"},{"alvo":"goldcoin","confianca":0.5,"epistemico":"fato","origem":"wiki","rel":"relaciona"},{"alvo":"install","confianca":0.5,"epistemico":"fato","origem":"wiki","rel":"relaciona"},{"alvo":"gpu","confianca":0.5,"epistemico":"fato","origem":"wiki","rel":"relaciona"},{"alvo":"montagem","confianca":0.5,"epistemico":"fato","origem":"wiki","rel":"relaciona"},{"alvo":"pow_bit_taxa","confianca":0.5,"epistemico":"fato","origem":"wiki","rel":"relaciona"},{"alvo":"python","confianca":0.5,"epistemico":"fato","origem":"wiki","rel":"relaciona"},{"alvo":"readme","confianca":0.5,"epistemico":"fato","origem":"wiki","rel":"relaciona"},{"alvo":"parabola_isa","confianca":0.5,"epistemico":"fato","origem":"wiki","rel":"relaciona"},{"alvo":"koch_atlas","confianca":0.5,"epistemico":"fato","origem":"wiki","rel":"relaciona"}],"confianca":1.0,"contraexemplos":[],"descricao":"Paper: fase3.tex","epistemico":"fato","exemplos":[],"id":"fase3","memoria":"semantica","meta":{"dominio":"paper","nivel":5,"tipo_no":"paper"},"origem":"paper","palavras_chave":[],"sinonimos":[],"tipo":"conceito","titulo":"fase3"}
\section{fase3}
Paper: fase3.tex
\prov{relações: relaciona transcricao; relaciona pow\_gato; relaciona fechamento\_blocos; relaciona prova\_isa; relaciona teste\_ordem; relaciona goldcoin; relaciona install; relaciona gpu; relaciona montagem; relaciona pow\_bit\_taxa; relaciona python; relaciona readme; relaciona parabola\_isa; relaciona koch\_atlas}
\prov{proveniência: paper \textperiodcentered{} fato \textperiodcentered{} confiança 1.0 \textperiodcentered{} domínio paper \textperiodcentered{} história 15 versões no jornal}

%CRISTAL {"arestas":[{"alvo":"transcricao","confianca":0.5,"epistemico":"fato","origem":"wiki","rel":"relaciona"},{"alvo":"goldcoin","confianca":0.5,"epistemico":"fato","origem":"wiki","rel":"relaciona"},{"alvo":"pow_gato","confianca":0.5,"epistemico":"fato","origem":"wiki","rel":"relaciona"},{"alvo":"prova_isa","confianca":0.5,"epistemico":"fato","origem":"wiki","rel":"relaciona"},{"alvo":"gpu","confianca":0.5,"epistemico":"fato","origem":"wiki","rel":"relaciona"},{"alvo":"install","confianca":0.5,"epistemico":"fato","origem":"wiki","rel":"relaciona"},{"alvo":"teste_ordem","confianca":0.5,"epistemico":"fato","origem":"wiki","rel":"relaciona"},{"alvo":"koch_parabola","confianca":0.5,"epistemico":"fato","origem":"wiki","rel":"relaciona"},{"alvo":"montagem","confianca":0.5,"epistemico":"fato","origem":"wiki","rel":"relaciona"},{"alvo":"readme","confianca":0.5,"epistemico":"fato","origem":"wiki","rel":"relaciona"},{"alvo":"koch_atlas","confianca":0.5,"epistemico":"fato","origem":"wiki","rel":"relaciona"},{"alvo":"go","confianca":0.5,"epistemico":"fato","origem":"wiki","rel":"relaciona"},{"alvo":"java","confianca":0.5,"epistemico":"fato","origem":"wiki","rel":"relaciona"},{"alvo":"parabola_isa","confianca":0.5,"epistemico":"fato","origem":"wiki","rel":"relaciona"},{"alvo":"sessao_interativa","confianca":0.5,"epistemico":"fato","origem":"wiki","rel":"relaciona"},{"alvo":"python","confianca":0.5,"epistemico":"fato","origem":"wiki","rel":"relaciona"},{"alvo":"koch_pell","confianca":0.5,"epistemico":"fato","origem":"wiki","rel":"relaciona"},{"alvo":"pow_bit_taxa","confianca":0.5,"epistemico":"fato","origem":"wiki","rel":"relaciona"}],"confianca":1.0,"contraexemplos":[],"descricao":"Paper: fechamento_blocos.tex","epistemico":"fato","exemplos":[],"id":"fechamento_blocos","memoria":"semantica","meta":{"dominio":"paper","nivel":5,"tipo_no":"paper"},"origem":"paper","palavras_chave":[],"sinonimos":[],"tipo":"conceito","titulo":"fechamento_blocos"}
\section{fechamento\_blocos}
Paper: fechamento\_blocos.tex
\prov{relações: relaciona transcricao; relaciona goldcoin; relaciona pow\_gato; relaciona prova\_isa; relaciona gpu; relaciona install; relaciona teste\_ordem; relaciona koch\_parabola; relaciona montagem; relaciona readme; relaciona koch\_atlas; relaciona go; relaciona java; relaciona parabola\_isa; relaciona sessao\_interativa; relaciona python; relaciona koch\_pell; relaciona pow\_bit\_taxa}
\prov{proveniência: paper \textperiodcentered{} fato \textperiodcentered{} confiança 1.0 \textperiodcentered{} domínio paper \textperiodcentered{} história 19 versões no jornal}

%CRISTAL {"arestas":[{"alvo":"fechamento_blocos","confianca":1.0,"epistemico":"fato","origem":"paper","rel":"parte_de"}],"confianca":1.0,"contraexemplos":[],"descricao":"TI A condição de fechamento --- Q^ acíclico --- é equivalente, pelo fato elementar da teoria de quivers (o ideal das flechas de kQ^ é nilpotente sse Q^ finito é acíclico; cf. Gabriel/ASS, face de rede), a: não há ciclo orientado no grafo de espera Q^ entre as máquinas --- definindo deadlock formalmente como tal ciclo, isto é exatamente sem deadlock. Portanto: [ a recursão fecha blocagem é admissível ão há deadlock entre as máquinas. ] Isto não é vácuo: quando a blocagem é inadmis","epistemico":"fato","exemplos":[],"id":"fechamento_e_estatuto","memoria":"semantica","meta":{"dominio":"paper","nivel":5,"verificacao":"conceito novo"},"origem":"paper","palavras_chave":[],"sinonimos":[],"tipo":"conceito","titulo":"Fechamento e estatuto"}
\section{Fechamento e estatuto}
TI A condição de fechamento --- Q\^{} acíclico --- é equivalente, pelo fato elementar da teoria de quivers (o ideal das flechas de kQ\^{} é nilpotente sse Q\^{} finito é acíclico; cf. Gabriel/ASS, face de rede), a: não há ciclo orientado no grafo de espera Q\^{} entre as máquinas --- definindo deadlock formalmente como tal ciclo, isto é exatamente sem deadlock. Portanto: [ a recursão fecha blocagem é admissível ão há deadlock entre as máquinas. ] Isto não é vácuo: quando a blocagem é inadmis
\prov{relações: parte\_de fechamento\_blocos}
\prov{proveniência: paper \textperiodcentered{} fato \textperiodcentered{} confiança 1.0 \textperiodcentered{} domínio paper \textperiodcentered{} história 3 versões no jornal}

%CRISTAL {"arestas":[{"alvo":"floco_de_neve","confianca":1.0,"epistemico":"fato","origem":"paper","rel":"parte_de"}],"confianca":1.0,"contraexemplos":[],"descricao":"defi[Reconstrução] Dado (L,): [ (L,);=; que L()=L,;;()=. ] defi Reconstrução condicionada [(i)] Só L não basta: não é injectivo em s (colisões); [(ii)] Só não basta: recall do Word só pelo atrator =0~memoria; [(iii)] L+ reconstrói desde que o Word vivo w permaneça na célula branca (slot mem) ou seja recuperável por busca finita no núcleo; o floco fixa a ramificação z² (2-para-1, A=0). teo proof[Esboço] (i)--(ii) Dados empírico","epistemico":"fato","exemplos":[],"id":"formalismo_operacional","memoria":"semantica","meta":{"dominio":"paper","nivel":5,"verificacao":"conceito novo"},"origem":"paper","palavras_chave":[],"sinonimos":[],"tipo":"conceito","titulo":"Formalismo operacional"}
\section{Formalismo operacional}
defi[Reconstrução] Dado (L,): [ (L,);=; que L()=L,;;()=. ] defi Reconstrução condicionada [(i)] Só L não basta: não é injectivo em s (colisões); [(ii)] Só não basta: recall do Word só pelo atrator =0\~{}memoria; [(iii)] L+ reconstrói desde que o Word vivo w permaneça na célula branca (slot mem) ou seja recuperável por busca finita no núcleo; o floco fixa a ramificação z\ensuremath{^{2}} (2-para-1, A=0). teo proof[Esboço] (i)--(ii) Dados empírico
\prov{relações: parte\_de floco\_de\_neve}
\prov{proveniência: paper \textperiodcentered{} fato \textperiodcentered{} confiança 1.0 \textperiodcentered{} domínio paper \textperiodcentered{} história 3 versões no jornal}

%CRISTAL {"arestas":[{"alvo":"fase3","confianca":1.0,"epistemico":"fato","origem":"paper","rel":"parte_de"}],"confianca":1.0,"contraexemplos":[],"descricao":"Prioridade de curadoria para o primeiro release comercial:\n\ncenter{1.3}\ntabular{@{}lll@{}}\n & ínio &  sugerida \\\\\n\nP0 & software/git, shell, python & docs oficiais + man pages \\\\\nP0 & documentacao/markdown, latex & guias Tiffany + templates \\\\\nP1 & linux, docker, javascript & docs upstream curadas \\\\\nP1 & matematica/algebra, estatistica & material didático livre \\\\\nP2 & demais domínios da árvore & expansão contínua por release \\\\\n\ntabular\ncenter","epistemico":"fato","exemplos":[],"id":"formato_da_base","memoria":"semantica","meta":{"dominio":"paper","nivel":5,"verificacao":"overlap moderado — revisar manualmente"},"origem":"paper","palavras_chave":[],"sinonimos":[],"tipo":"conceito","titulo":"Formato da base"}
\section{Formato da base}
Prioridade de curadoria para o primeiro release comercial:

center\{1.3\}
tabular\{@\{\}lll@\{\}\}
 \& ínio \&  sugerida \textbackslash{}\textbackslash{}

P0 \& software/git, shell, python \& docs oficiais + man pages \textbackslash{}\textbackslash{}
P0 \& documentacao/markdown, latex \& guias Tiffany + templates \textbackslash{}\textbackslash{}
P1 \& linux, docker, javascript \& docs upstream curadas \textbackslash{}\textbackslash{}
P1 \& matematica/algebra, estatistica \& material didático livre \textbackslash{}\textbackslash{}
P2 \& demais domínios da árvore \& expansão contínua por release \textbackslash{}\textbackslash{}

tabular
center
\prov{relações: parte\_de fase3}
\prov{proveniência: paper \textperiodcentered{} fato \textperiodcentered{} confiança 1.0 \textperiodcentered{} domínio paper \textperiodcentered{} história 2 versões no jornal}

%CRISTAL {"arestas":[{"alvo":"a_engenharia_de_ouro_pa_pgs_agm_fc_reta","confianca":1.0,"epistemico":"fato","origem":"paper","rel":"parte_de"},{"alvo":"racionais","confianca":1.0,"epistemico":"fato","origem":"paper","rel":"referencia"}],"confianca":1.0,"contraexemplos":[],"descricao":"Atrator --- síntese Com A(a,b)=a+b2, G(a,b)=ab: [nosep,leftmargin=1.4em] G=(A( a, b)): média geométrica = aritmética no log; (a,b)(A,G) converge quadraticamente a (a,b); o limite é projetado, não fechado elementarmente quando a b (Transformada Dourada, Teo.~AGM). teorema","epistemico":"fato","exemplos":[],"id":"fracoes_continuas_e_reta_de_ouro","memoria":"semantica","meta":{"dominio":"paper","nivel":5,"verificacao":"slug idêntico — enriquecer, não duplicar"},"origem":"paper","palavras_chave":[],"sinonimos":[],"tipo":"conceito","titulo":"Frações contínuas e reta de ouro"}
\section{Frações contínuas e reta de ouro}
Atrator --- síntese Com A(a,b)=a+b2, G(a,b)=ab: [nosep,leftmargin=1.4em] G=(A( a, b)): média geométrica = aritmética no log; (a,b)(A,G) converge quadraticamente a (a,b); o limite é projetado, não fechado elementarmente quando a b (Transformada Dourada, Teo.\~{}AGM). teorema
\prov{relações: parte\_de a\_engenharia\_de\_ouro\_pa\_pgs\_agm\_fc\_reta; referencia racionais}
\prov{proveniência: paper \textperiodcentered{} fato \textperiodcentered{} confiança 1.0 \textperiodcentered{} domínio paper \textperiodcentered{} história 47 versões no jornal}

%CRISTAL {"arestas":[{"alvo":"matematica","confianca":0.95,"epistemico":"fato","origem":"paper","rel":"parte_de"}],"confianca":1.0,"contraexemplos":[],"descricao":"N O tecido é a alma torcida de um fluxo. Pelo canal do painel (canalpainel.atende, op=\"sobe\"), ele se grava como cabeçalho (rec) mais corpo (ativos). O rec real, do código: [ rec=,id, nome, ts, nbytes, bins, qubit, ativos=,bin,. ] O cabeçalho identifica e dimensiona (o qubit dá o andamento/fase); o corpo ativos é o espectro: cada bin carrega um peso (a amplitude). N A ordem não vem de fora. O tecido se ordena pela estrutura intrínseca: pelo índice do bin, que são os","epistemico":"fato","exemplos":[],"id":"gentil_teorema_14_p_102_estabelece_que_os_operadores_soma_e_diferenca_de_ord","memoria":"semantica","meta":{"dominio":"paper","duplicata_sugerida":[],"nivel":5,"verificacao":"parte de conceito mais amplo 'gentil_teorema_14_p_102_estabelece_que_os_operadores_soma_e_diferenca_de_ord'"},"origem":"paper","palavras_chave":[],"sinonimos":[],"tipo":"conceito","titulo":"gentil Teorema 14 — p 102), estabelece que os operadores Soma e Diferença de ord"}
\section{gentil Teorema 14 — p 102), estabelece que os operadores Soma e Diferença de ord}
N O tecido é a alma torcida de um fluxo. Pelo canal do painel (canalpainel.atende, op="sobe"), ele se grava como cabeçalho (rec) mais corpo (ativos). O rec real, do código: [ rec=,id, nome, ts, nbytes, bins, qubit, ativos=,bin,. ] O cabeçalho identifica e dimensiona (o qubit dá o andamento/fase); o corpo ativos é o espectro: cada bin carrega um peso (a amplitude). N A ordem não vem de fora. O tecido se ordena pela estrutura intrínseca: pelo índice do bin, que são os
\prov{relações: parte\_de matematica}
\prov{proveniência: paper \textperiodcentered{} fato \textperiodcentered{} confiança 1.0 \textperiodcentered{} domínio paper \textperiodcentered{} história 6 versões no jornal}

%CRISTAL {"arestas":[{"alvo":"matematica","confianca":0.95,"epistemico":"fato","origem":"paper","rel":"parte_de"}],"confianca":1.0,"contraexemplos":[],"descricao":"I Camada mutável, distinta de conhecimento curado. Alimentada automaticamente; podável;\npromovível à Base Cliente após curadoria.","epistemico":"fato","exemplos":[],"id":"gentil_teorema_32_p_150_ainda_resta_uma_pergunta_quais_tipos_de_sequencias","memoria":"semantica","meta":{"dominio":"paper","duplicata_sugerida":[],"nivel":5,"verificacao":"parte de conceito mais amplo 'gentil_teorema_32_p_150_ainda_resta_uma_pergunta_quais_tipos_de_sequencias'"},"origem":"paper","palavras_chave":[],"sinonimos":[],"tipo":"conceito","titulo":"gentil Teorema 32 — , p 150) ainda resta uma pergunta: quais tipos de sequências"}
\section{gentil Teorema 32 — , p 150) ainda resta uma pergunta: quais tipos de sequências}
I Camada mutável, distinta de conhecimento curado. Alimentada automaticamente; podável;
promovível à Base Cliente após curadoria.
\prov{relações: parte\_de matematica}
\prov{proveniência: paper \textperiodcentered{} fato \textperiodcentered{} confiança 1.0 \textperiodcentered{} domínio paper \textperiodcentered{} história 5 versões no jornal}

%CRISTAL {"arestas":[{"alvo":"goldchain","confianca":1.0,"epistemico":"fato","origem":"paper","rel":"parte_de"},{"alvo":"pow_gato","confianca":0.5,"epistemico":"fato","origem":"wiki","rel":"relaciona"},{"alvo":"gpu","confianca":0.5,"epistemico":"fato","origem":"wiki","rel":"relaciona"},{"alvo":"montagem","confianca":0.5,"epistemico":"fato","origem":"wiki","rel":"relaciona"},{"alvo":"transcricao","confianca":0.5,"epistemico":"fato","origem":"wiki","rel":"relaciona"},{"alvo":"readme","confianca":0.5,"epistemico":"fato","origem":"wiki","rel":"relaciona"},{"alvo":"python","confianca":0.5,"epistemico":"fato","origem":"wiki","rel":"relaciona"},{"alvo":"teste_ordem","confianca":0.5,"epistemico":"fato","origem":"wiki","rel":"relaciona"},{"alvo":"install","confianca":0.5,"epistemico":"fato","origem":"wiki","rel":"relaciona"},{"alvo":"prova_isa","confianca":0.5,"epistemico":"fato","origem":"wiki","rel":"relaciona"},{"alvo":"koch_atlas","confianca":0.5,"epistemico":"fato","origem":"wiki","rel":"relaciona"},{"alvo":"vida-morte","confianca":0.5,"epistemico":"fato","origem":"wiki","rel":"relaciona"},{"alvo":"sessao_interativa","confianca":0.5,"epistemico":"fato","origem":"wiki","rel":"relaciona"},{"alvo":"mamifero","confianca":0.5,"epistemico":"fato","origem":"wiki","rel":"relaciona"},{"alvo":"hopfield_rede_hub","confianca":0.5,"epistemico":"fato","origem":"wiki","rel":"relaciona"},{"alvo":"java","confianca":0.5,"epistemico":"fato","origem":"wiki","rel":"relaciona"},{"alvo":"parabola_isa","confianca":0.5,"epistemico":"fato","origem":"wiki","rel":"relaciona"},{"alvo":"kernel","confianca":0.5,"epistemico":"fato","origem":"wiki","rel":"relaciona"},{"alvo":"grava_fato_resposta_normal_em_seguida","confianca":0.5,"epistemico":"fato","origem":"wiki","rel":"relaciona"},{"alvo":"numa","confianca":0.5,"epistemico":"fato","origem":"wiki","rel":"relaciona"},{"alvo":"thread","confianca":0.5,"epistemico":"fato","origem":"wiki","rel":"relaciona"},{"alvo":"pow_espectral","confianca":0.5,"epistemico":"fato","origem":"wiki","rel":"relaciona"},{"alvo":"main","confianca":0.5,"epistemico":"fato","origem":"wiki","rel":"relaciona"},{"alvo":"koch_memoria","confianca":0.5,"epistemico":"fato","origem":"wiki","rel":"relaciona"},{"alvo":"pow_metal_pre_fecho","confianca":0.5,"epistemico":"fato","origem":"wiki","rel":"relaciona"},{"alvo":"memoria_virtual","confianca":0.5,"epistemico":"fato","origem":"wiki","rel":"relaciona"},{"alvo":"koch_parabola","confianca":0.5,"epistemico":"fato","origem":"wiki","rel":"relaciona"},{"alvo":"koch_pell","confianca":0.5,"epistemico":"fato","origem":"wiki","rel":"relaciona"},{"alvo":"rust","confianca":0.5,"epistemico":"fato","origem":"wiki","rel":"relaciona"},{"alvo":"pow_bit_taxa","confianca":0.5,"epistemico":"fato","origem":"wiki","rel":"relaciona"}],"confianca":1.0,"contraexemplos":[],"descricao":"Paper: goldcoin.tex","epistemico":"fato","exemplos":[],"id":"goldcoin","memoria":"semantica","meta":{"dominio":"paper","nivel":5,"tipo_no":"paper"},"origem":"paper","palavras_chave":[],"sinonimos":[],"tipo":"conceito","titulo":"goldcoin"}
\section{goldcoin}
Paper: goldcoin.tex
\prov{relações: parte\_de goldchain; relaciona pow\_gato; relaciona gpu; relaciona montagem; relaciona transcricao; relaciona readme; relaciona python; relaciona teste\_ordem; relaciona install; relaciona prova\_isa; relaciona koch\_atlas; relaciona vida-morte; relaciona sessao\_interativa; relaciona mamifero; relaciona hopfield\_rede\_hub; relaciona java; relaciona parabola\_isa; relaciona kernel; relaciona grava\_fato\_resposta\_normal\_em\_seguida; relaciona numa; relaciona thread; relaciona pow\_espectral; relaciona main; relaciona koch\_memoria; relaciona pow\_metal\_pre\_fecho; relaciona memoria\_virtual; relaciona koch\_parabola; relaciona koch\_pell; relaciona rust; relaciona pow\_bit\_taxa}
\prov{proveniência: paper \textperiodcentered{} fato \textperiodcentered{} confiança 1.0 \textperiodcentered{} domínio paper \textperiodcentered{} história 36 versões no jornal}

%CRISTAL {"arestas":[{"alvo":"leitura_quantica_o_qubit_bell_e_o_emaranhamento","confianca":1.0,"epistemico":"fato","origem":"paper","rel":"parte_de"}],"confianca":1.0,"contraexemplos":[],"descricao":"Resta ler o salto quântico (Teor.~, Obs.~obs:0999) do ponto de vista de um observador que vive no círculo /. A Secão~ deu a base --- a equação universal s=2s (Teor.~) e a costura 10 (Teor.~) --- e apontou a leitura de Kepler entre as direções; desenvolvendo-a por dentro do disco: s=2s é o campo central harmônico F=-2 q, cujas órbitas são cônicas (elipse na ilha, hipérbole no oceano) com o tensor de Hooke Sᵢj=vᵢvj+2qᵢqj co","epistemico":"fato","exemplos":[],"id":"hardware_cosmico_a_algebra_como_computacao_e_o_associador_como_recurso","memoria":"semantica","meta":{"dominio":"paper","nivel":5,"verificacao":"conceito novo"},"origem":"paper","palavras_chave":[],"sinonimos":[],"tipo":"conceito","titulo":"Hardware cósmico: a álgebra como computação e o associador como recurso"}
\section{Hardware cósmico: a álgebra como computação e o associador como recurso}
Resta ler o salto quântico (Teor.\~{}, Obs.\~{}obs:0999) do ponto de vista de um observador que vive no círculo /. A Secão\~{} deu a base --- a equação universal s=2s (Teor.\~{}) e a costura 10 (Teor.\~{}) --- e apontou a leitura de Kepler entre as direções; desenvolvendo-a por dentro do disco: s=2s é o campo central harmônico F=-2 q, cujas órbitas são cônicas (elipse na ilha, hipérbole no oceano) com o tensor de Hooke S\ensuremath{_{i}}j=v\ensuremath{_{i}}vj+2q\ensuremath{_{i}}qj co
\prov{relações: parte\_de leitura\_quantica\_o\_qubit\_bell\_e\_o\_emaranhamento}
\prov{proveniência: paper \textperiodcentered{} fato \textperiodcentered{} confiança 1.0 \textperiodcentered{} domínio paper \textperiodcentered{} história 3 versões no jornal}

%CRISTAL {"arestas":[{"alvo":"topologia","confianca":0.5,"epistemico":"fato","origem":"wiki","rel":"relaciona"},{"alvo":"teste_ordem","confianca":0.5,"epistemico":"fato","origem":"wiki","rel":"relaciona"},{"alvo":"thread","confianca":0.5,"epistemico":"fato","origem":"wiki","rel":"relaciona"},{"alvo":"testes","confianca":0.5,"epistemico":"fato","origem":"wiki","rel":"relaciona"},{"alvo":"vida-morte","confianca":0.5,"epistemico":"fato","origem":"wiki","rel":"relaciona"},{"alvo":"ponteiro","confianca":0.5,"epistemico":"fato","origem":"wiki","rel":"relaciona"},{"alvo":"resposta_a_pergunta_dificil","confianca":0.5,"epistemico":"fato","origem":"wiki","rel":"relaciona"},{"alvo":"motivos","confianca":0.5,"epistemico":"fato","origem":"wiki","rel":"relaciona"},{"alvo":"parabola_pa","confianca":0.5,"epistemico":"fato","origem":"wiki","rel":"relaciona"},{"alvo":"projeto_cristaljson","confianca":0.5,"epistemico":"fato","origem":"wiki","rel":"relaciona"},{"alvo":"por_sessao_conversadadostelemetriasessaojsonl","confianca":0.5,"epistemico":"fato","origem":"wiki","rel":"relaciona"},{"alvo":"regua_associativa_e_o_risco_de_vazamento","confianca":0.5,"epistemico":"fato","origem":"wiki","rel":"relaciona"},{"alvo":"virtualizacao","confianca":0.5,"epistemico":"fato","origem":"wiki","rel":"relaciona"},{"alvo":"parabola_destruir","confianca":0.5,"epistemico":"fato","origem":"wiki","rel":"relaciona"},{"alvo":"pow_metal_pre_fecho","confianca":0.5,"epistemico":"fato","origem":"wiki","rel":"relaciona"},{"alvo":"recursao","confianca":0.5,"epistemico":"fato","origem":"wiki","rel":"relaciona"},{"alvo":"rust","confianca":0.5,"epistemico":"fato","origem":"wiki","rel":"relaciona"},{"alvo":"lobo_frontal","confianca":0.5,"epistemico":"fato","origem":"wiki","rel":"relaciona"},{"alvo":"syscall","confianca":0.5,"epistemico":"fato","origem":"wiki","rel":"relaciona"}],"confianca":1.0,"contraexemplos":[],"descricao":"","epistemico":"fato","exemplos":[],"id":"hardware_sequencias_de_cauchy","memoria":"semantica","meta":{"dominio":"paper","nivel":5,"verificacao":"slug idêntico — enriquecer, não duplicar"},"origem":"paper","palavras_chave":[],"sinonimos":[],"tipo":"conceito","titulo":"Hardware: sequências de Cauchy"}
\section{Hardware: sequências de Cauchy}
\prov{relações: relaciona topologia; relaciona teste\_ordem; relaciona thread; relaciona testes; relaciona vida-morte; relaciona ponteiro; relaciona resposta\_a\_pergunta\_dificil; relaciona motivos; relaciona parabola\_pa; relaciona projeto\_cristaljson; relaciona por\_sessao\_conversadadostelemetriasessaojsonl; relaciona regua\_associativa\_e\_o\_risco\_de\_vazamento; relaciona virtualizacao; relaciona parabola\_destruir; relaciona pow\_metal\_pre\_fecho; relaciona recursao; relaciona rust; relaciona lobo\_frontal; relaciona syscall}
\prov{proveniência: paper \textperiodcentered{} fato \textperiodcentered{} confiança 1.0 \textperiodcentered{} domínio paper \textperiodcentered{} história 63 versões no jornal}

%CRISTAL {"arestas":[{"alvo":"main","confianca":1.0,"epistemico":"fato","origem":"paper","rel":"parte_de"}],"confianca":1.0,"contraexemplos":[],"descricao":"tabularp4.6cm|p4.2cm|p2.4cm verificação & valor esperado & tolerância norma em: |xy|-|x||y| & 0 & 10⁻¹² ₇ (SVD) & sete 3, quatorze zeros & 10⁻¹² Der() (núcleo) & dim 14 & rank numérico associador [a,b,c]+2 (em V₇) & 0 & 10⁻¹³ _: ocorrências de 11 & 0 & convenção de endpoints palavras sem 11 (compr. L) & F(L+2) & exata G round-trip G⁻¹Gx-x & 0 & arredondamento G diagonaliza U_ & |ₖ|=1 & 10^","epistemico":"fato","exemplos":[],"id":"heuristica_estrutural_a_cadeia_dos_clay","memoria":"semantica","meta":{"dominio":"paper","nivel":5,"verificacao":"slug idêntico — enriquecer, não duplicar"},"origem":"paper","palavras_chave":[],"sinonimos":[],"tipo":"conceito","titulo":"Heurística estrutural: a cadeia dos Clay"}
\section{Heurística estrutural: a cadeia dos Clay}
tabularp4.6cm|p4.2cm|p2.4cm verificação \& valor esperado \& tolerância norma em: |xy|-|x||y| \& 0 \& 10\ensuremath{^{-12}} \ensuremath{_{7}} (SVD) \& sete 3, quatorze zeros \& 10\ensuremath{^{-12}} Der() (núcleo) \& dim 14 \& rank numérico associador [a,b,c]+2 (em V\ensuremath{_{7}}) \& 0 \& 10\ensuremath{^{-13}} \_: ocorrências de 11 \& 0 \& convenção de endpoints palavras sem 11 (compr. L) \& F(L+2) \& exata G round-trip G\ensuremath{^{-1}}Gx-x \& 0 \& arredondamento G diagonaliza U\_ \& |\ensuremath{_{k}}|=1 \& 10\^{}
\prov{relações: parte\_de main}
\prov{proveniência: paper \textperiodcentered{} fato \textperiodcentered{} confiança 1.0 \textperiodcentered{} domínio paper \textperiodcentered{} história 5 versões no jornal}

%CRISTAL {"arestas":[{"alvo":"corpo_tropical","confianca":1.0,"epistemico":"fato","origem":"paper","rel":"parte_de"}],"confianca":1.0,"contraexemplos":[],"descricao":"Cada curva elíptica E tem a sua função L(E,s) --- logo a sua reta crítica, o seu cone de luz. Onde Riemann dá uma reta por L, BSD lê o que acontece na interseção de duas retas. A interseção é o ponto central. A reta crítica --- vertical, onde moram os zeros, o cone N=0 --- cruza a reta real --- horizontal, a aritmética, onde vivem os pontos racionais --- num único ponto: o centro s=1. É a interseção das duas retas de Riemann, a dos zeros e a da aritmética. O posto é a multiplicidade","epistemico":"fato","exemplos":[],"id":"hodge_a_costura_de_todos_os_fractais_da_cadeia","memoria":"semantica","meta":{"dominio":"paper","nivel":5,"verificacao":"slug idêntico — enriquecer, não duplicar"},"origem":"paper","palavras_chave":[],"sinonimos":[],"tipo":"conceito","titulo":"Hodge: a costura de todos os fractais da cadeia"}
\section{Hodge: a costura de todos os fractais da cadeia}
Cada curva elíptica E tem a sua função L(E,s) --- logo a sua reta crítica, o seu cone de luz. Onde Riemann dá uma reta por L, BSD lê o que acontece na interseção de duas retas. A interseção é o ponto central. A reta crítica --- vertical, onde moram os zeros, o cone N=0 --- cruza a reta real --- horizontal, a aritmética, onde vivem os pontos racionais --- num único ponto: o centro s=1. É a interseção das duas retas de Riemann, a dos zeros e a da aritmética. O posto é a multiplicidade
\prov{relações: parte\_de corpo\_tropical}
\prov{proveniência: paper \textperiodcentered{} fato \textperiodcentered{} confiança 1.0 \textperiodcentered{} domínio paper \textperiodcentered{} história 4 versões no jornal}

%CRISTAL {"arestas":[{"alvo":"montagem","confianca":1.0,"epistemico":"fato","origem":"paper","rel":"parte_de"}],"confianca":1.0,"contraexemplos":[],"descricao":"V Tudo isto roda como neuronio.c: e += popcount(byte & 0xAA); o += popcount(byte & 0x55); return (e+o, e); cujo núcleo desmontado é and 0xaa / and 0x55 / popcnt / add. Bate, byte a byte, com a forma em Python (4 linhas, sem biblioteca) e com a matriz A. L Fecha o círculo: volta ao pop e ao ₂ de machine.c --- o que era “a medida” desde o primeiro paper é a contagem das metades da borboleta.","epistemico":"fato","exemplos":[],"id":"honestidade_o_que_ainda_nao_fecha","memoria":"semantica","meta":{"dominio":"paper","duplicata_sugerida":[],"nivel":5,"verificacao":"parte de conceito mais amplo 'honestidade_o_que_ainda_nao_fecha'"},"origem":"paper","palavras_chave":[],"sinonimos":[],"tipo":"conceito","titulo":"Honestidade: o que ainda não fecha"}
\section{Honestidade: o que ainda não fecha}
V Tudo isto roda como neuronio.c: e += popcount(byte \& 0xAA); o += popcount(byte \& 0x55); return (e+o, e); cujo núcleo desmontado é and 0xaa / and 0x55 / popcnt / add. Bate, byte a byte, com a forma em Python (4 linhas, sem biblioteca) e com a matriz A. L Fecha o círculo: volta ao pop e ao \ensuremath{_{2}} de machine.c --- o que era “a medida” desde o primeiro paper é a contagem das metades da borboleta.
\prov{relações: parte\_de montagem}
\prov{proveniência: paper \textperiodcentered{} fato \textperiodcentered{} confiança 1.0 \textperiodcentered{} domínio paper \textperiodcentered{} história 4 versões no jornal}

%CRISTAL {"arestas":[{"alvo":"antena","confianca":1.0,"epistemico":"fato","origem":"paper","rel":"parte_de"}],"confianca":1.0,"contraexemplos":[],"descricao":"Acoplado ao canal Eᵢj, o elemento tem constante de propagação guiada que herda a banda da segurança: [ ()=k()²-k cut₂. ] A ressonância é a condição clássica Z ᵢₙ(₀)=0, ou, geometricamente, [ 2(₀), eff+ gap+ carga=2 q, (meia-onda: eff g/2, g=2/). ] A k cut que sela a dimensão (Segurança Dimensional) é a mesma que fixa a banda da antena: abaixo dela, modo evanescente; acima, modo que irradia. Sintonizar é band-limitar.","epistemico":"fato","exemplos":[],"id":"impedancia_e_casamento","memoria":"semantica","meta":{"dominio":"paper","nivel":5,"verificacao":"overlap moderado — revisar manualmente"},"origem":"paper","palavras_chave":[],"sinonimos":[],"tipo":"conceito","titulo":"Impedância e casamento"}
\section{Impedância e casamento}
Acoplado ao canal E\ensuremath{_{i}}j, o elemento tem constante de propagação guiada que herda a banda da segurança: [ ()=k()\ensuremath{^{2}}-k cut\ensuremath{_{2}}. ] A ressonância é a condição clássica Z \ensuremath{_{in}}(\ensuremath{_{0}})=0, ou, geometricamente, [ 2(\ensuremath{_{0}}), eff+ gap+ carga=2 q, (meia-onda: eff g/2, g=2/). ] A k cut que sela a dimensão (Segurança Dimensional) é a mesma que fixa a banda da antena: abaixo dela, modo evanescente; acima, modo que irradia. Sintonizar é band-limitar.
\prov{relações: parte\_de antena}
\prov{proveniência: paper \textperiodcentered{} fato \textperiodcentered{} confiança 1.0 \textperiodcentered{} domínio paper \textperiodcentered{} história 3 versões no jornal}

%CRISTAL {"arestas":[{"alvo":"produto_comercial","confianca":1.0,"epistemico":"fato","origem":"paper","rel":"parte_de"}],"confianca":1.0,"contraexemplos":[],"descricao":"center{1.35}\ntabular{@{}ll@{}}\n & údo \\\\\n\nlibtiffany.so + headers & engine (V) \\\\\ntiffany-core-1.0.tar.zst & base/ completa + manifest.json \\\\\ntiffany-client-template.tgz & cliente/ vazia + exemplos \\\\\ntiffany-ingest & pipeline §IV (CLI) \\\\\ntiffany & CLI consulta + orquestrador \\\\\ntiffany-python & bindings pip install tiffany \\\\\n\ntabular\ncenter","epistemico":"fato","exemplos":[],"id":"instalacao_cliente_novo","memoria":"semantica","meta":{"dominio":"paper","nivel":5,"verificacao":"overlap moderado — revisar manualmente"},"origem":"paper","palavras_chave":[],"sinonimos":[],"tipo":"conceito","titulo":"Instalação (cliente novo)"}
\section{Instalação (cliente novo)}
center\{1.35\}
tabular\{@\{\}ll@\{\}\}
 \& údo \textbackslash{}\textbackslash{}

libtiffany.so + headers \& engine (V) \textbackslash{}\textbackslash{}
tiffany-core-1.0.tar.zst \& base/ completa + manifest.json \textbackslash{}\textbackslash{}
tiffany-client-template.tgz \& cliente/ vazia + exemplos \textbackslash{}\textbackslash{}
tiffany-ingest \& pipeline §IV (CLI) \textbackslash{}\textbackslash{}
tiffany \& CLI consulta + orquestrador \textbackslash{}\textbackslash{}
tiffany-python \& bindings pip install tiffany \textbackslash{}\textbackslash{}

tabular
center
\prov{relações: parte\_de produto\_comercial}
\prov{proveniência: paper \textperiodcentered{} fato \textperiodcentered{} confiança 1.0 \textperiodcentered{} domínio paper \textperiodcentered{} história 2 versões no jornal}

%CRISTAL {"arestas":[{"alvo":"paper_0_sintese","confianca":1.0,"epistemico":"fato","origem":"paper","rel":"parte_de"}],"confianca":1.0,"contraexemplos":[],"descricao":"abstract spacing1.075 Uma única multiplicacão em ⁿ --- zs w=(a₀b₀- a, b) +(a₀ b+b₀ a+s, a b) --- gera, ao se variar o botão s, uma famlia contnua de álgebras que contém, como pontos. Mostramos, do zero e sem recorrer a nenhum resultado externo, que a métrica quântica de Lopes é a régua desse espaco de álgebras, e que dela sai uma cadeia: s(a,c) a+c²=1 N²=a R=const. A conservacão pitágórica a+c²=1 (comutatividade c=|s| contra associ","epistemico":"fato","exemplos":[],"id":"introduc","memoria":"semantica","meta":{"dominio":"paper","nivel":5,"verificacao":"overlap moderado — revisar manualmente"},"origem":"paper","palavras_chave":[],"sinonimos":[],"tipo":"conceito","titulo":"Introduc"}
\section{Introduc}
abstract spacing1.075 Uma única multiplicacão em \ensuremath{^{n}} --- zs w=(a\ensuremath{_{0}}b\ensuremath{_{0}}- a, b) +(a\ensuremath{_{0}} b+b\ensuremath{_{0}} a+s, a b) --- gera, ao se variar o botão s, uma famlia contnua de álgebras que contém, como pontos. Mostramos, do zero e sem recorrer a nenhum resultado externo, que a métrica quântica de Lopes é a régua desse espaco de álgebras, e que dela sai uma cadeia: s(a,c) a+c\ensuremath{^{2}}=1 N\ensuremath{^{2}}=a R=const. A conservacão pitágórica a+c\ensuremath{^{2}}=1 (comutatividade c=|s| contra associ
\prov{relações: parte\_de paper\_0\_sintese}
\prov{proveniência: paper \textperiodcentered{} fato \textperiodcentered{} confiança 1.0 \textperiodcentered{} domínio paper \textperiodcentered{} história 5 versões no jornal}

%CRISTAL {"arestas":[{"alvo":"reticulado","confianca":0.023,"epistemico":"fato","origem":"manual","rel":"referencia"},{"alvo":"reticulados","confianca":-0.039,"epistemico":"fato","origem":"manual","rel":"referencia"},{"alvo":"casl-propagation","confianca":0.039,"epistemico":"fato","origem":"manual","rel":"referencia"},{"alvo":"colapso","confianca":-0.062,"epistemico":"fato","origem":"manual","rel":"referencia"},{"alvo":"transformada_dourada","confianca":1.0,"epistemico":"fato","origem":"paper","rel":"parte_de"},{"alvo":"a_base_topologica","confianca":1.0,"epistemico":"fato","origem":"paper","rel":"parte_de"},{"alvo":"termodinamica","confianca":1.0,"epistemico":"fato","origem":"paper","rel":"parte_de"},{"alvo":"mecanica_quantica","confianca":1.0,"epistemico":"fato","origem":"paper","rel":"parte_de"},{"alvo":"seguranca_por_irreversibilidade","confianca":1.0,"epistemico":"fato","origem":"paper","rel":"parte_de"},{"alvo":"redes","confianca":1.0,"epistemico":"fato","origem":"paper","rel":"parte_de"},{"alvo":"computacao_grafica","confianca":1.0,"epistemico":"fato","origem":"paper","rel":"parte_de"}],"confianca":1.0,"contraexemplos":[],"descricao":"cap:computacaografica","epistemico":"fato","exemplos":[],"id":"introducao","memoria":"semantica","meta":{"dominio":"paper","nivel":5,"verificacao":"slug idêntico — enriquecer, não duplicar"},"origem":"paper","palavras_chave":[],"sinonimos":[],"tipo":"conceito","titulo":"Introdução"}
\section{Introdução}
cap:computacaografica
\prov{relações: referencia reticulado; referencia reticulados; referencia casl-propagation; referencia colapso; parte\_de transformada\_dourada; parte\_de a\_base\_topologica; parte\_de termodinamica; parte\_de mecanica\_quantica; parte\_de seguranca\_por\_irreversibilidade; parte\_de redes; parte\_de computacao\_grafica}
\prov{proveniência: paper \textperiodcentered{} fato \textperiodcentered{} confiança 1.0 \textperiodcentered{} domínio paper \textperiodcentered{} história 32 versões no jornal}

%CRISTAL {"arestas":[{"alvo":"casl-propagation","confianca":-0.031,"epistemico":"fato","origem":"manual","rel":"referencia"}],"confianca":1.0,"contraexemplos":[],"descricao":"Por construção (Definição~), uma capacidade emitida com profundidade admite cadeias de re-delegação de comprimento no máximo: [ 1 salto 1 0 (não re-delegável). ] [nosep] =0: o destinatário usa a capacidade, mas não a passa a ninguém. É o caso do “dar acesso sem permitir compartilhar”. =k: re-delegável por mais k saltos; cada salto decrementa. =: irrestrito (o comportamento legado, sem limite). lemma[ é uma energia bem-funda","epistemico":"fato","exemplos":[],"id":"invariantes_de_seguranca","memoria":"semantica","meta":{"dominio":"paper","nivel":5},"origem":"paper","palavras_chave":[],"sinonimos":[],"tipo":"conceito","titulo":"Invariantes de segurança"}
\section{Invariantes de segurança}
Por construção (Definição\~{}), uma capacidade emitida com profundidade admite cadeias de re-delegação de comprimento no máximo: [ 1 salto 1 0 (não re-delegável). ] [nosep] =0: o destinatário usa a capacidade, mas não a passa a ninguém. É o caso do “dar acesso sem permitir compartilhar”. =k: re-delegável por mais k saltos; cada salto decrementa. =: irrestrito (o comportamento legado, sem limite). lemma[ é uma energia bem-funda
\prov{relações: referencia casl-propagation}
\prov{proveniência: paper \textperiodcentered{} fato \textperiodcentered{} confiança 1.0 \textperiodcentered{} domínio paper \textperiodcentered{} história 6 versões no jornal}

%CRISTAL {"arestas":[{"alvo":"projeto_cristaljson","confianca":0.5,"epistemico":"fato","origem":"wiki","rel":"relaciona"},{"alvo":"utilitarismo","confianca":0.5,"epistemico":"fato","origem":"wiki","rel":"relaciona"},{"alvo":"word","confianca":0.5,"epistemico":"fato","origem":"wiki","rel":"relaciona"},{"alvo":"virtualizacao","confianca":0.5,"epistemico":"fato","origem":"wiki","rel":"relaciona"},{"alvo":"vida-morte","confianca":0.5,"epistemico":"fato","origem":"wiki","rel":"relaciona"},{"alvo":"pipeline_de_ingestao","confianca":0.5,"epistemico":"fato","origem":"wiki","rel":"relaciona"},{"alvo":"pell","confianca":0.5,"epistemico":"fato","origem":"wiki","rel":"relaciona"},{"alvo":"telemetria_shadow_producao","confianca":0.5,"epistemico":"fato","origem":"wiki","rel":"relaciona"},{"alvo":"painel_estelar","confianca":0.5,"epistemico":"fato","origem":"wiki","rel":"relaciona"},{"alvo":"proposition_escala_fasepropescala_xix","confianca":0.5,"epistemico":"fato","origem":"wiki","rel":"relaciona"}],"confianca":1.0,"contraexemplos":[],"descricao":"Quíntuplo (estado, condição, δ, peso, origem) — unidade de propagação BAI.","epistemico":"fato","exemplos":[],"id":"kappa","memoria":"semantica","meta":{"dominio":"paper","nivel":5,"serve":"Primitiva do quinteto BAI","verificacao":"slug idêntico — enriquecer, não duplicar"},"origem":"paper","palavras_chave":[],"sinonimos":[],"tipo":"conceito","titulo":"kappa"}
\section{kappa}
Quíntuplo (estado, condição, \ensuremath{\delta}, peso, origem) — unidade de propagação BAI.
\prov{relações: relaciona projeto\_cristaljson; relaciona utilitarismo; relaciona word; relaciona virtualizacao; relaciona vida-morte; relaciona pipeline\_de\_ingestao; relaciona pell; relaciona telemetria\_shadow\_producao; relaciona painel\_estelar; relaciona proposition\_escala\_fasepropescala\_xix}
\prov{proveniência: paper \textperiodcentered{} fato \textperiodcentered{} confiança 1.0 \textperiodcentered{} domínio paper \textperiodcentered{} história 63 versões no jornal}

%CRISTAL {"arestas":[{"alvo":"floco_de_neve","confianca":1.0,"epistemico":"fato","origem":"paper","rel":"parte_de"}],"confianca":1.0,"contraexemplos":[],"descricao":"D O núcleo vivo produz w=[,total,e,] e (w). Entre eles, a parábola a+c²=1~koch acopla as células brancas pelo canal negro E₀₁. Desdistorcida para comprimento de arco, vive um parâmetro s=(c²)[0,1] --- coordenada sobre a curva, não scatter no plano (a,c). Reta associativa real. Cantor ternário C₃(s): dígitos 0,2, mapa s 3s 1. Vizinhos no grid têm s pequeno; 88% partilham classe Cantor ao passo~memoria. Reta associativa ouro. Convergentes Fibonacci Fₙ/F","epistemico":"fato","exemplos":[],"id":"koch_dissipada_nao_codificadora","memoria":"semantica","meta":{"dominio":"paper","nivel":5,"verificacao":"overlap moderado — revisar manualmente"},"origem":"paper","palavras_chave":[],"sinonimos":[],"tipo":"conceito","titulo":"Koch: dissipada, não codificadora"}
\section{Koch: dissipada, não codificadora}
D O núcleo vivo produz w=[,total,e,] e (w). Entre eles, a parábola a+c\ensuremath{^{2}}=1\~{}koch acopla as células brancas pelo canal negro E\ensuremath{_{01}}. Desdistorcida para comprimento de arco, vive um parâmetro s=(c\ensuremath{^{2}})[0,1] --- coordenada sobre a curva, não scatter no plano (a,c). Reta associativa real. Cantor ternário C\ensuremath{_{3}}(s): dígitos 0,2, mapa s 3s 1. Vizinhos no grid têm s pequeno; 88\% partilham classe Cantor ao passo\~{}memoria. Reta associativa ouro. Convergentes Fibonacci F\ensuremath{_{n}}/F
\prov{relações: parte\_de floco\_de\_neve}
\prov{proveniência: paper \textperiodcentered{} fato \textperiodcentered{} confiança 1.0 \textperiodcentered{} domínio paper \textperiodcentered{} história 3 versões no jornal}

%CRISTAL {"arestas":[{"alvo":"paper_0_sintese","confianca":1.0,"epistemico":"fato","origem":"paper","rel":"parte_de"}],"confianca":1.0,"contraexemplos":[],"descricao":"ão ternária e seus três dicionários A série inteira pendura-se num único invariante. A matriz de multiplicacão à esquerda de zs w fatora, Lz = |z|²,(a₀² + s²|Im,z|²), e o segundo fator carrega o sinal que decide tudo: sobre (com s0) é soma de quadrados, e a álgebra é de divisão; ao complexificar (zⁿ) ele vira forma isotrópica e surgem divisores de zero --- a forma compacta perde a divisão. As","epistemico":"fato","exemplos":[],"id":"leitura_aritmetica_a_regua_de_gentil_o_agm","memoria":"semantica","meta":{"dominio":"paper","nivel":5,"verificacao":"conceito novo"},"origem":"paper","palavras_chave":[],"sinonimos":[],"tipo":"conceito","titulo":"Leitura aritmética: a régua de Lopes, o AGM"}
\section{Leitura aritmética: a régua de Lopes, o AGM}
ão ternária e seus três dicionários A série inteira pendura-se num único invariante. A matriz de multiplicacão à esquerda de zs w fatora, Lz = |z|\ensuremath{^{2}},(a\ensuremath{_{0}}\ensuremath{^{2}} + s\ensuremath{^{2}}|Im,z|\ensuremath{^{2}}), e o segundo fator carrega o sinal que decide tudo: sobre (com s0) é soma de quadrados, e a álgebra é de divisão; ao complexificar (z\ensuremath{^{n}}) ele vira forma isotrópica e surgem divisores de zero --- a forma compacta perde a divisão. As
\prov{relações: parte\_de paper\_0\_sintese}
\prov{proveniência: paper \textperiodcentered{} fato \textperiodcentered{} confiança 1.0 \textperiodcentered{} domínio paper \textperiodcentered{} história 5 versões no jornal}

%CRISTAL {"arestas":[{"alvo":"paper_0_sintese","confianca":1.0,"epistemico":"fato","origem":"paper","rel":"parte_de"}],"confianca":1.0,"contraexemplos":[],"descricao":"ão octoniônica A generalizacão de cinco botões (p,r,t,s,u) pesa cada um dos cinco termos da multiplicacão~: y=(p,a₀ b₀-r a, b)+(t,a₀ b+u,b₀ a+s, a b). Torcão octoniônica A secão espacial Sⁿ⁻¹ de versores admite a paralelizacão de Cartan--Schouten, com torcão igual à 3-forma de estrutura cᵢjₖ= eᵢ ej,eₖ. Em S³ o Jacobiador zera (Im,=su(2), Lie); em S⁷ ele é o assoc","epistemico":"fato","exemplos":[],"id":"leitura_cosmologica_a_energia_escura_e_o_universo_ds","memoria":"semantica","meta":{"dominio":"paper","nivel":5,"verificacao":"conceito novo"},"origem":"paper","palavras_chave":[],"sinonimos":[],"tipo":"conceito","titulo":"Leitura cosmológica: a energia escura e o universo dS"}
\section{Leitura cosmológica: a energia escura e o universo dS}
ão octoniônica A generalizacão de cinco botões (p,r,t,s,u) pesa cada um dos cinco termos da multiplicacão\~{}: y=(p,a\ensuremath{_{0}} b\ensuremath{_{0}}-r a, b)+(t,a\ensuremath{_{0}} b+u,b\ensuremath{_{0}} a+s, a b). Torcão octoniônica A secão espacial S\ensuremath{^{n-1}} de versores admite a paralelizacão de Cartan--Schouten, com torcão igual à 3-forma de estrutura c\ensuremath{_{i}}j\ensuremath{_{k}}= e\ensuremath{_{i}} ej,e\ensuremath{_{k}}. Em S\ensuremath{^{3}} o Jacobiador zera (Im,=su(2), Lie); em S\ensuremath{^{7}} ele é o assoc
\prov{relações: parte\_de paper\_0\_sintese}
\prov{proveniência: paper \textperiodcentered{} fato \textperiodcentered{} confiança 1.0 \textperiodcentered{} domínio paper \textperiodcentered{} história 4 versões no jornal}

%CRISTAL {"arestas":[{"alvo":"paper_0_sintese","confianca":1.0,"epistemico":"fato","origem":"paper","rel":"parte_de"}],"confianca":1.0,"contraexemplos":[],"descricao":"ões (trialidade) Por que três cores: Im() a tem dimensão 6 A Secão~ (Teor. da estrutura complexa interna) mostrou que a^ só carrega a estrutura complexa J=ad a (J²=-1) quando A é normada de divisão (identidade de Lagrange), e que por Hurwitz isso dá N=_( a^)=0,1,3 para, (nos sedênions a identidade falha e J desaparece), com StabG₂( a)=SU(3) em --- máximo, pois é a última de divisão. A correspondên","epistemico":"fato","exemplos":[],"id":"leitura_das_unidades_o_sistema_por-unidade_e_o_si","memoria":"semantica","meta":{"dominio":"paper","nivel":5,"verificacao":"conceito novo"},"origem":"paper","palavras_chave":[],"sinonimos":[],"tipo":"conceito","titulo":"Leitura das unidades: o sistema por-unidade e o SI"}
\section{Leitura das unidades: o sistema por-unidade e o SI}
ões (trialidade) Por que três cores: Im() a tem dimensão 6 A Secão\~{} (Teor. da estrutura complexa interna) mostrou que a\^{} só carrega a estrutura complexa J=ad a (J\ensuremath{^{2}}=-1) quando A é normada de divisão (identidade de Lagrange), e que por Hurwitz isso dá N=\_( a\^{})=0,1,3 para, (nos sedênions a identidade falha e J desaparece), com StabG\ensuremath{_{2}}( a)=SU(3) em --- máximo, pois é a última de divisão. A correspondên
\prov{relações: parte\_de paper\_0\_sintese}
\prov{proveniência: paper \textperiodcentered{} fato \textperiodcentered{} confiança 1.0 \textperiodcentered{} domínio paper \textperiodcentered{} história 3 versões no jornal}

%CRISTAL {"arestas":[{"alvo":"paper_0_sintese","confianca":1.0,"epistemico":"fato","origem":"paper","rel":"parte_de"}],"confianca":1.0,"contraexemplos":[],"descricao":"Lemos a multiplicação como geometria (de Sitter), como mecânica quântica (a seção de Hopf) e --- na Secão~ --- como gramática do setor forte (cor, méson, bárion). Falta uma leitura: a álgebra como computação. Ela expõe uma operação que nenhuma máquina já construída usa, e dá a razão algébrica de por que não a usa. Os três eixos de ordem. Toda computação é feita de produtos; o que distingue um regime de outro é qual ordem importa. Na computação clássica, a álgebra","epistemico":"fato","exemplos":[],"id":"leitura_de_part_culas_o_setor_de_materia_pela_torre_de_divisao","memoria":"semantica","meta":{"dominio":"paper","nivel":5,"verificacao":"conceito novo"},"origem":"paper","palavras_chave":[],"sinonimos":[],"tipo":"conceito","titulo":"Leitura de partculas: o setor de matéria pela torre de divisão"}
\section{Leitura de partculas: o setor de matéria pela torre de divisão}
Lemos a multiplicação como geometria (de Sitter), como mecânica quântica (a seção de Hopf) e --- na Secão\~{} --- como gramática do setor forte (cor, méson, bárion). Falta uma leitura: a álgebra como computação. Ela expõe uma operação que nenhuma máquina já construída usa, e dá a razão algébrica de por que não a usa. Os três eixos de ordem. Toda computação é feita de produtos; o que distingue um regime de outro é qual ordem importa. Na computação clássica, a álgebra
\prov{relações: parte\_de paper\_0\_sintese}
\prov{proveniência: paper \textperiodcentered{} fato \textperiodcentered{} confiança 1.0 \textperiodcentered{} domínio paper \textperiodcentered{} história 4 versões no jornal}

%CRISTAL {"arestas":[{"alvo":"paper_0_sintese","confianca":1.0,"epistemico":"fato","origem":"paper","rel":"parte_de"}],"confianca":1.0,"contraexemplos":[],"descricao":"Três nveis, como em toda a série. Derivado: a estrutura adimensional --- os números puros (2, 3, Q=N/3, 1/2, _, o Q²=8 do running de ), fixados pela álgebra; o SI os toma do experimento, aqui têm origem. Posto: a escala --- uma volta com conteúdo dimensional, e é a do meio (o campo do tubo de fluxo, o horizonte; Secão~), não um parâmetro livre. Convencão: o mol e a candela --- contagem e percepcão humanas, fora da fsica","epistemico":"fato","exemplos":[],"id":"leitura_dinamica_a_maquina_e_a_materia_ligada","memoria":"semantica","meta":{"dominio":"paper","nivel":5,"verificacao":"conceito novo"},"origem":"paper","palavras_chave":[],"sinonimos":[],"tipo":"conceito","titulo":"Leitura dinâmica: a máquina e a matéria ligada"}
\section{Leitura dinâmica: a máquina e a matéria ligada}
Três nveis, como em toda a série. Derivado: a estrutura adimensional --- os números puros (2, 3, Q=N/3, 1/2, \_, o Q\ensuremath{^{2}}=8 do running de ), fixados pela álgebra; o SI os toma do experimento, aqui têm origem. Posto: a escala --- uma volta com conteúdo dimensional, e é a do meio (o campo do tubo de fluxo, o horizonte; Secão\~{}), não um parâmetro livre. Convencão: o mol e a candela --- contagem e percepcão humanas, fora da fsica
\prov{relações: parte\_de paper\_0\_sintese}
\prov{proveniência: paper \textperiodcentered{} fato \textperiodcentered{} confiança 1.0 \textperiodcentered{} domínio paper \textperiodcentered{} história 4 versões no jornal}

%CRISTAL {"arestas":[{"alvo":"paper_0_sintese","confianca":1.0,"epistemico":"fato","origem":"paper","rel":"parte_de"}],"confianca":1.0,"contraexemplos":[],"descricao":"A projecão contnua da famlia na reta é densa (o reticulado não probe nada). Mas um corte discreto --- proibir um bloco na expansão binária, isto é, na árvore de sinais --- é um subshift de tipo finito, e a sua sombra na reta é um Cantor de dimensão /2<1. O corte mnimo (proibir 11) é o golden mean shift: =, e [ H = 2=0,6942 (box-counting: 0,6935), ] a dimensão do número de ouro (Figura~). A carne fractal não estava na famlia (cri","epistemico":"fato","exemplos":[],"id":"leitura_psicologica_o_vazio_e_o_corac","memoria":"semantica","meta":{"dominio":"paper","nivel":5,"verificacao":"conceito novo"},"origem":"paper","palavras_chave":[],"sinonimos":[],"tipo":"conceito","titulo":"Leitura psicológica: o vazio e o corac"}
\section{Leitura psicológica: o vazio e o corac}
A projecão contnua da famlia na reta é densa (o reticulado não probe nada). Mas um corte discreto --- proibir um bloco na expansão binária, isto é, na árvore de sinais --- é um subshift de tipo finito, e a sua sombra na reta é um Cantor de dimensão /2<1. O corte mnimo (proibir 11) é o golden mean shift: =, e [ H = 2=0,6942 (box-counting: 0,6935), ] a dimensão do número de ouro (Figura\~{}). A carne fractal não estava na famlia (cri
\prov{relações: parte\_de paper\_0\_sintese}
\prov{proveniência: paper \textperiodcentered{} fato \textperiodcentered{} confiança 1.0 \textperiodcentered{} domínio paper \textperiodcentered{} história 4 versões no jornal}

%CRISTAL {"arestas":[{"alvo":"paper_0_sintese","confianca":1.0,"epistemico":"fato","origem":"paper","rel":"parte_de"}],"confianca":1.0,"contraexemplos":[],"descricao":"1: o universo dS₄ A Secão~ provou que a fatia de um único campo local --- fase e parâmetro s radial --- é o de Sitter bidimensional (R=2), e deixou explícito que a geometria do espaço-forma completo Mₙ (R=n(n-1)) ficava para os artigos companheiros. É o que fazemos agora: em ⁴ o estado completo é (s, a), com a S² a direção do campo local, e a métrica de pullback é ds²=-(1-s²),d²+ds²1-s²+s²,d²S₂. O cálc","epistemico":"fato","exemplos":[],"id":"leitura_quantica_o_qubit_bell_e_o_emaranhamento","memoria":"semantica","meta":{"dominio":"paper","nivel":5,"verificacao":"conceito novo"},"origem":"paper","palavras_chave":[],"sinonimos":[],"tipo":"conceito","titulo":"Leitura quântica: o qubit, Bell e o emaranhamento"}
\section{Leitura quântica: o qubit, Bell e o emaranhamento}
1: o universo dS\ensuremath{_{4}} A Secão\~{} provou que a fatia de um único campo local --- fase e parâmetro s radial --- é o de Sitter bidimensional (R=2), e deixou explícito que a geometria do espaço-forma completo M\ensuremath{_{n}} (R=n(n-1)) ficava para os artigos companheiros. É o que fazemos agora: em \ensuremath{^{4}} o estado completo é (s, a), com a S\ensuremath{^{2}} a direção do campo local, e a métrica de pullback é ds\ensuremath{^{2}}=-(1-s\ensuremath{^{2}}),d\ensuremath{^{2}}+ds\ensuremath{^{2}}1-s\ensuremath{^{2}}+s\ensuremath{^{2}},d\ensuremath{^{2}}S\ensuremath{_{2}}. O cálc
\prov{relações: parte\_de paper\_0\_sintese}
\prov{proveniência: paper \textperiodcentered{} fato \textperiodcentered{} confiança 1.0 \textperiodcentered{} domínio paper \textperiodcentered{} história 4 versões no jornal}

%CRISTAL {"arestas":[{"alvo":"plano_no_so","confianca":1.0,"epistemico":"fato","origem":"paper","rel":"prova"}],"confianca":1.0,"contraexemplos":[],"descricao":"[ é uma energia bem-fundada ciclos são inofensivos] Como toda aresta válida estritamente decrementa a profundidade (_₀-1, Definição~), é um variante bem-fundado sobre as cadeias de re-delegação: qualquer caminho de concessões com finito tem comprimento e termina. Logo a aciclicidade do grafo é dispensável --- um ciclo u v u apenas faz cair a cada volta até 0, quando a propagação para. A terminação (e a não-escalada) vem da atenuação, não da estrutura","epistemico":"fato","exemplos":[],"id":"lemma_e_uma_energia_bem-fundada_ciclos_sao_inofen","memoria":"semantica","meta":{"dominio":"paper","nivel":5,"tipo_teorema":"lemma"},"origem":"paper","palavras_chave":[],"sinonimos":[],"tipo":"conceito","titulo":"[ é uma energia bem-fundada ciclos são inofen…"}
\section{[ é uma energia bem-fundada ciclos são inofen…}
[ é uma energia bem-fundada ciclos são inofensivos] Como toda aresta válida estritamente decrementa a profundidade (\_\ensuremath{_{0}}-1, Definição\~{}), é um variante bem-fundado sobre as cadeias de re-delegação: qualquer caminho de concessões com finito tem comprimento e termina. Logo a aciclicidade do grafo é dispensável --- um ciclo u v u apenas faz cair a cada volta até 0, quando a propagação para. A terminação (e a não-escalada) vem da atenuação, não da estrutura
\prov{relações: prova plano\_no\_so}
\prov{proveniência: paper \textperiodcentered{} fato \textperiodcentered{} confiança 1.0 \textperiodcentered{} domínio paper \textperiodcentered{} história 4 versões no jornal}

%CRISTAL {"arestas":[{"alvo":"colapso","confianca":0.95,"epistemico":"fato","origem":"paper","rel":"parte_de"},{"alvo":"virtualizacao","confianca":0.5,"epistemico":"fato","origem":"wiki","rel":"relaciona"},{"alvo":"vida-morte","confianca":0.5,"epistemico":"fato","origem":"wiki","rel":"relaciona"},{"alvo":"pipeline_de_ingestao","confianca":0.5,"epistemico":"fato","origem":"wiki","rel":"relaciona"},{"alvo":"semiotica_peirce_triade","confianca":0.5,"epistemico":"fato","origem":"wiki","rel":"relaciona"},{"alvo":"telemetria_shadow_producao","confianca":0.5,"epistemico":"fato","origem":"wiki","rel":"relaciona"}],"confianca":1.0,"contraexemplos":[],"descricao":"é unitária: ₀^|x(t)|²,dtt=-^|x(eu)|²,du.","epistemico":"fato","exemplos":[],"id":"lemma_lemunit_e_unitaria_0xt2dtt_-","memoria":"semantica","meta":{"dominio":"paper","nivel":5,"tipo_teorema":"lemma","verificacao":"conceito novo"},"origem":"paper","palavras_chave":[],"sinonimos":[],"tipo":"conceito","titulo":"é unitária: ₀^|x(t)|²,dtt=-^|…"}
\section{é unitária: \ensuremath{_{0}}\^{}|x(t)|\ensuremath{^{2}},dtt=-\^{}|…}
é unitária: \ensuremath{_{0}}\^{}|x(t)|\ensuremath{^{2}},dtt=-\^{}|x(eu)|\ensuremath{^{2}},du.
\prov{relações: parte\_de colapso; relaciona virtualizacao; relaciona vida-morte; relaciona pipeline\_de\_ingestao; relaciona semiotica\_peirce\_triade; relaciona telemetria\_shadow\_producao}
\prov{proveniência: paper \textperiodcentered{} fato \textperiodcentered{} confiança 1.0 \textperiodcentered{} domínio paper \textperiodcentered{} história 10 versões no jornal}

%CRISTAL {"arestas":[{"alvo":"corpo_glacial","confianca":1.0,"epistemico":"fato","origem":"paper","rel":"parte_de"}],"confianca":1.0,"contraexemplos":[],"descricao":"D O nilpotente roteia: Eᵢj,Ejl=Eᵢl --- a mensagem flui i j l compondo canais; e ( N)=1+ N (pois N²=0) é o fluxo infinitesimal que o nilpotente gera. O idempotente fixa o estado; o nilpotente move entre estados. A leitura branca (projeção, e) e a negra (a coordenada do nilpotente, lida por uma derivação / N na subálgebra dual [N]/(N²)) são as duas pontas do mesmo demux. I Daí “a dinâmica está acoplada”: separar branca e negra é perder ou o suporte ou o movimento.","epistemico":"fato","exemplos":[],"id":"limites_crivo_incorporado","memoria":"semantica","meta":{"dominio":"paper","nivel":5,"verificacao":"slug idêntico — enriquecer, não duplicar"},"origem":"paper","palavras_chave":[],"sinonimos":[],"tipo":"conceito","titulo":"Limites (crivo incorporado)"}
\section{Limites (crivo incorporado)}
D O nilpotente roteia: E\ensuremath{_{i}}j,Ejl=E\ensuremath{_{i}}l --- a mensagem flui i j l compondo canais; e ( N)=1+ N (pois N\ensuremath{^{2}}=0) é o fluxo infinitesimal que o nilpotente gera. O idempotente fixa o estado; o nilpotente move entre estados. A leitura branca (projeção, e) e a negra (a coordenada do nilpotente, lida por uma derivação / N na subálgebra dual [N]/(N\ensuremath{^{2}})) são as duas pontas do mesmo demux. I Daí “a dinâmica está acoplada”: separar branca e negra é perder ou o suporte ou o movimento.
\prov{relações: parte\_de corpo\_glacial}
\prov{proveniência: paper \textperiodcentered{} fato \textperiodcentered{} confiança 1.0 \textperiodcentered{} domínio paper \textperiodcentered{} história 4 versões no jornal}

%CRISTAL {"arestas":[{"alvo":"main","confianca":1.0,"epistemico":"fato","origem":"paper","rel":"parte_de"}],"confianca":1.0,"contraexemplos":[],"descricao":"tabularp3.2cm|p8.2cm|c camada & afirmação & estatuto Hurwitz & só, são normadas de divisão & T Cayley--Dickson & produto, conjugação, tábua reconstruídos da semente & D/N & ()=14; ₇3, núcleo 14 & T/N _ & exclusão de 11 emerge da geometria (²=+1) & D contagens & palavras F(L+2), razão; família m-bonacci & D/N G discreta & DFT logarítmica diagonaliza a dilatação unitária U_ (não o Perron da transformação) & D/N genom","epistemico":"fato","exemplos":[],"id":"limites_explicitos","memoria":"semantica","meta":{"dominio":"paper","nivel":5,"verificacao":"slug idêntico — enriquecer, não duplicar"},"origem":"paper","palavras_chave":[],"sinonimos":[],"tipo":"conceito","titulo":"Limites explícitos"}
\section{Limites explícitos}
tabularp3.2cm|p8.2cm|c camada \& afirmação \& estatuto Hurwitz \& só, são normadas de divisão \& T Cayley--Dickson \& produto, conjugação, tábua reconstruídos da semente \& D/N \& ()=14; \ensuremath{_{7}}3, núcleo 14 \& T/N \_ \& exclusão de 11 emerge da geometria (\ensuremath{^{2}}=+1) \& D contagens \& palavras F(L+2), razão; família m-bonacci \& D/N G discreta \& DFT logarítmica diagonaliza a dilatação unitária U\_ (não o Perron da transformação) \& D/N genom
\prov{relações: parte\_de main}
\prov{proveniência: paper \textperiodcentered{} fato \textperiodcentered{} confiança 1.0 \textperiodcentered{} domínio paper \textperiodcentered{} história 5 versões no jornal}

%CRISTAL {"arestas":[{"alvo":"floco_de_neve","confianca":1.0,"epistemico":"fato","origem":"paper","rel":"parte_de"}],"confianca":1.0,"contraexemplos":[],"descricao":"V Resumo consolidado: 1.35 tabularl r r Métrica & Associativo (L) & Floco () vizinho & 4,5 10⁻³ (s) & 16 bits (w,) & 0,452 & 0,118 Mesma classe vizinho & Cantor 88% / Fib 80% & --- Colisões conteúdo distinto & --- & 18,548 / 4,607 buckets Recall só por & --- & 0/32 Roundtrip slot+floco & --- & 32/32 () (resíduos mod p) & --- & 1009/1009 Pell cifra rastreáveis () & --- & 18/18 tabular O floco g","epistemico":"fato","exemplos":[],"id":"limites_honestos","memoria":"semantica","meta":{"dominio":"paper","nivel":5,"verificacao":"conceito novo"},"origem":"paper","palavras_chave":[],"sinonimos":[],"tipo":"conceito","titulo":"Limites honestos"}
\section{Limites honestos}
V Resumo consolidado: 1.35 tabularl r r Métrica \& Associativo (L) \& Floco () vizinho \& 4,5 10\ensuremath{^{-3}} (s) \& 16 bits (w,) \& 0,452 \& 0,118 Mesma classe vizinho \& Cantor 88\% / Fib 80\% \& --- Colisões conteúdo distinto \& --- \& 18,548 / 4,607 buckets Recall só por \& --- \& 0/32 Roundtrip slot+floco \& --- \& 32/32 () (resíduos mod p) \& --- \& 1009/1009 Pell cifra rastreáveis () \& --- \& 18/18 tabular O floco g
\prov{relações: parte\_de floco\_de\_neve}
\prov{proveniência: paper \textperiodcentered{} fato \textperiodcentered{} confiança 1.0 \textperiodcentered{} domínio paper \textperiodcentered{} história 3 versões no jornal}

%CRISTAL {"arestas":[{"alvo":"formato_da_base","confianca":1.0,"epistemico":"fato","origem":"paper","rel":"parte_de"},{"alvo":"word","confianca":0.5,"epistemico":"fato","origem":"wiki","rel":"relaciona"},{"alvo":"virtualizacao","confianca":0.5,"epistemico":"fato","origem":"wiki","rel":"relaciona"},{"alvo":"vontade_de_poder","confianca":0.5,"epistemico":"fato","origem":"wiki","rel":"relaciona"},{"alvo":"while","confianca":0.5,"epistemico":"fato","origem":"wiki","rel":"relaciona"}],"confianca":1.0,"contraexemplos":[],"descricao":"D Contrato estável entre ingestão, orquestrador e . A engine lê .tsv+.idx; tudo mais\né metadado para humanos, CI e atualização incremental.","epistemico":"fato","exemplos":[],"id":"linha_do_nostsv","memoria":"semantica","meta":{"dominio":"paper","nivel":5,"verificacao":"overlap moderado — revisar manualmente"},"origem":"paper","palavras_chave":[],"sinonimos":[],"tipo":"conceito","titulo":"Linha do nos.tsv"}
\section{Linha do nos.tsv}
D Contrato estável entre ingestão, orquestrador e . A engine lê .tsv+.idx; tudo mais
é metadado para humanos, CI e atualização incremental.
\prov{relações: parte\_de formato\_da\_base; relaciona word; relaciona virtualizacao; relaciona vontade\_de\_poder; relaciona while}
\prov{proveniência: paper \textperiodcentered{} fato \textperiodcentered{} confiança 1.0 \textperiodcentered{} domínio paper \textperiodcentered{} história 7 versões no jornal}

%CRISTAL {"arestas":[{"alvo":"resgate_liquidacao","confianca":1.0,"epistemico":"fato","origem":"paper","rel":"parte_de"},{"alvo":"goldchain","confianca":1.0,"epistemico":"fato","origem":"paper","rel":"parte_de"}],"confianca":1.0,"contraexemplos":[],"descricao":"DN Resgatar é decifrar. O Pell cifrado é um número gigante (elemento de ₂¹²⁸); só quem tem a chave de ouro g volta à base padrão e recupera o primo de prata: [ cifra=g m p, m;=; g⁻¹(g m), g⁻¹=g,²k-2 (Fermat). ] Em código, corpoprata.py: mineraₙabase cifra (g m), le resgata (g⁻¹). A reversão é exata em bits (sem float). De fora --- com outra base g' --- sai g'⁻¹g m, um elemento qualquer que não é primo de prata: ruído. lstlisting def basedoclie","epistemico":"fato","exemplos":[],"id":"liquidacao","memoria":"semantica","meta":{"dominio":"paper","nivel":5,"verificacao":"overlap moderado — revisar manualmente"},"origem":"paper","palavras_chave":[],"sinonimos":[],"tipo":"conceito","titulo":"Liquidação"}
\section{Liquidação}
DN Resgatar é decifrar. O Pell cifrado é um número gigante (elemento de \ensuremath{_{2}}\ensuremath{^{128}}); só quem tem a chave de ouro g volta à base padrão e recupera o primo de prata: [ cifra=g m p, m;=; g\ensuremath{^{-1}}(g m), g\ensuremath{^{-1}}=g,\ensuremath{^{2}}k-2 (Fermat). ] Em código, corpoprata.py: minera\ensuremath{_{n}}abase cifra (g m), le resgata (g\ensuremath{^{-1}}). A reversão é exata em bits (sem float). De fora --- com outra base g' --- sai g'\ensuremath{^{-1}}g m, um elemento qualquer que não é primo de prata: ruído. lstlisting def basedoclie
\prov{relações: parte\_de resgate\_liquidacao; parte\_de goldchain}
\prov{proveniência: paper \textperiodcentered{} fato \textperiodcentered{} confiança 1.0 \textperiodcentered{} domínio paper \textperiodcentered{} história 10 versões no jornal}

%CRISTAL {"arestas":[{"alvo":"inteligencia_artificial","confianca":1.0,"epistemico":"fato","origem":"paper","rel":"parte_de"},{"alvo":"colapso","confianca":1.0,"epistemico":"fato","origem":"paper","rel":"usa"},{"alvo":"teste_ordem","confianca":0.5,"epistemico":"fato","origem":"wiki","rel":"relaciona"},{"alvo":"transcricao","confianca":0.5,"epistemico":"fato","origem":"wiki","rel":"relaciona"},{"alvo":"testes","confianca":0.5,"epistemico":"fato","origem":"wiki","rel":"relaciona"},{"alvo":"pow_metal_pre_fecho","confianca":0.5,"epistemico":"fato","origem":"wiki","rel":"relaciona"},{"alvo":"rust","confianca":0.5,"epistemico":"fato","origem":"wiki","rel":"relaciona"},{"alvo":"ponteiro","confianca":0.5,"epistemico":"fato","origem":"wiki","rel":"relaciona"},{"alvo":"numa","confianca":0.5,"epistemico":"fato","origem":"wiki","rel":"relaciona"},{"alvo":"projeto_cristaljson","confianca":0.5,"epistemico":"fato","origem":"wiki","rel":"relaciona"},{"alvo":"parabola_pa","confianca":0.5,"epistemico":"fato","origem":"wiki","rel":"relaciona"},{"alvo":"motivos","confianca":0.5,"epistemico":"fato","origem":"wiki","rel":"relaciona"},{"alvo":"vida-morte","confianca":0.5,"epistemico":"fato","origem":"wiki","rel":"relaciona"}],"confianca":1.0,"contraexemplos":[],"descricao":"Paper: lobo_frontal.tex","epistemico":"fato","exemplos":[],"id":"lobo_frontal","memoria":"semantica","meta":{"dominio":"paper","nivel":5,"tipo_no":"paper"},"origem":"paper","palavras_chave":[],"sinonimos":[],"tipo":"conceito","titulo":"lobo_frontal"}
\section{lobo\_frontal}
Paper: lobo\_frontal.tex
\prov{relações: parte\_de inteligencia\_artificial; usa colapso; relaciona teste\_ordem; relaciona transcricao; relaciona testes; relaciona pow\_metal\_pre\_fecho; relaciona rust; relaciona ponteiro; relaciona numa; relaciona projeto\_cristaljson; relaciona parabola\_pa; relaciona motivos; relaciona vida-morte}
\prov{proveniência: paper \textperiodcentered{} fato \textperiodcentered{} confiança 1.0 \textperiodcentered{} domínio paper \textperiodcentered{} história 72 versões no jornal}

%CRISTAL {"arestas":[{"alvo":"quatro_ciencias","confianca":1.0,"epistemico":"fato","origem":"paper","rel":"explica"},{"alvo":"ciencias","confianca":1.0,"epistemico":"fato","origem":"paper","rel":"parte_de"},{"alvo":"sessao_interativa","confianca":0.5,"epistemico":"fato","origem":"wiki","rel":"relaciona"},{"alvo":"rust","confianca":0.5,"epistemico":"fato","origem":"wiki","rel":"relaciona"},{"alvo":"python","confianca":0.5,"epistemico":"fato","origem":"wiki","rel":"relaciona"},{"alvo":"utilitarismo","confianca":0.5,"epistemico":"fato","origem":"wiki","rel":"relaciona"},{"alvo":"transcricao","confianca":0.5,"epistemico":"fato","origem":"wiki","rel":"relaciona"},{"alvo":"virtualizacao","confianca":0.5,"epistemico":"fato","origem":"wiki","rel":"relaciona"},{"alvo":"modo_espectral_spect","confianca":0.5,"epistemico":"fato","origem":"wiki","rel":"relaciona"},{"alvo":"pitagorismo","confianca":0.5,"epistemico":"fato","origem":"wiki","rel":"relaciona"},{"alvo":"teoria_do_valor","confianca":0.5,"epistemico":"fato","origem":"wiki","rel":"relaciona"},{"alvo":"pedagogia","confianca":0.5,"epistemico":"fato","origem":"wiki","rel":"relaciona"},{"alvo":"resultado_completo","confianca":0.5,"epistemico":"fato","origem":"wiki","rel":"relaciona"}],"confianca":1.0,"contraexemplos":[],"descricao":"Paper: main.tex","epistemico":"fato","exemplos":[],"id":"main","memoria":"semantica","meta":{"dominio":"paper","nivel":5,"tipo_no":"paper"},"origem":"paper","palavras_chave":[],"sinonimos":[],"tipo":"conceito","titulo":"main"}
\section{main}
Paper: main.tex
\prov{relações: explica quatro\_ciencias; parte\_de ciencias; relaciona sessao\_interativa; relaciona rust; relaciona python; relaciona utilitarismo; relaciona transcricao; relaciona virtualizacao; relaciona modo\_espectral\_spect; relaciona pitagorismo; relaciona teoria\_do\_valor; relaciona pedagogia; relaciona resultado\_completo}
\prov{proveniência: paper \textperiodcentered{} fato \textperiodcentered{} confiança 1.0 \textperiodcentered{} domínio paper \textperiodcentered{} história 73 versões no jornal}

%CRISTAL {"arestas":[{"alvo":"fase3","confianca":1.0,"epistemico":"fato","origem":"paper","rel":"parte_de"}],"confianca":1.0,"contraexemplos":[],"descricao":"D Separação física e lógica da Base Geral. Nunca misturar --- nem em .tsv, nem em\nrelease, nem em backup.\n\ncamadabox\n\\$TIFFANY\\_HOME/cliente/\\\\\n manifest.json\\\\\n empresa/\\\\\n engenharia/ nos.tsv v1.idx meta.json\\\\\n financeiro/\\\\\n juridico/\\\\\n rh/\\\\\n produto/\\\\\ncamadabox\n\ndefi[Isolamento cliente]\nitemize2pt\n Caminho: base/ vs cliente/ --- roots distintos.\n Licença: meta.json.licenca do cliente é proprietary por padrão.\n Orquestrador: consulta cliente só se domínio detectado ou flag explícita.\n Rel","epistemico":"fato","exemplos":[],"id":"memoria_viva","memoria":"semantica","meta":{"dominio":"paper","nivel":5,"verificacao":"overlap moderado — revisar manualmente"},"origem":"paper","palavras_chave":[],"sinonimos":[],"tipo":"conceito","titulo":"Memória Viva"}
\section{Memória Viva}
D Separação física e lógica da Base Geral. Nunca misturar --- nem em .tsv, nem em
release, nem em backup.

camadabox
\textbackslash{}\$TIFFANY\textbackslash{}\_HOME/cliente/\textbackslash{}\textbackslash{}
 manifest.json\textbackslash{}\textbackslash{}
 empresa/\textbackslash{}\textbackslash{}
 engenharia/ nos.tsv v1.idx meta.json\textbackslash{}\textbackslash{}
 financeiro/\textbackslash{}\textbackslash{}
 juridico/\textbackslash{}\textbackslash{}
 rh/\textbackslash{}\textbackslash{}
 produto/\textbackslash{}\textbackslash{}
camadabox

defi[Isolamento cliente]
itemize2pt
 Caminho: base/ vs cliente/ --- roots distintos.
 Licença: meta.json.licenca do cliente é proprietary por padrão.
 Orquestrador: consulta cliente só se domínio detectado ou flag explícita.
 Rel
\prov{relações: parte\_de fase3}
\prov{proveniência: paper \textperiodcentered{} fato \textperiodcentered{} confiança 1.0 \textperiodcentered{} domínio paper \textperiodcentered{} história 2 versões no jornal}

%CRISTAL {"arestas":[{"alvo":"leitura_de_part_culas_o_setor_de_materia_pela_torre_de_divisao","confianca":1.0,"epistemico":"fato","origem":"paper","rel":"parte_de"}],"confianca":1.0,"contraexemplos":[],"descricao":"O estatuto destas identificacões: correspondência, não derivacão dinâmica Uma palavra sobre o que segue, para não prometer o que não temos. As identificacões deste setor são correspondências estruturais: a torre das álgebras de divisão contém, de forma natural e sem ajuste, objetos com exatamente os números quânticos (carga, cor, isospin, hipercarga) e a estrutura de grupo SU(3)c SU(2)L U(1)Y de uma geracão do Modelo P","epistemico":"fato","exemplos":[],"id":"meson_e_barion_numeros_do_campo_local_no_plano_de_fano","memoria":"semantica","meta":{"dominio":"paper","nivel":5,"verificacao":"conceito novo"},"origem":"paper","palavras_chave":[],"sinonimos":[],"tipo":"conceito","titulo":"Méson e bárion: números do campo local no plano de Fano"}
\section{Méson e bárion: números do campo local no plano de Fano}
O estatuto destas identificacões: correspondência, não derivacão dinâmica Uma palavra sobre o que segue, para não prometer o que não temos. As identificacões deste setor são correspondências estruturais: a torre das álgebras de divisão contém, de forma natural e sem ajuste, objetos com exatamente os números quânticos (carga, cor, isospin, hipercarga) e a estrutura de grupo SU(3)c SU(2)L U(1)Y de uma geracão do Modelo P
\prov{relações: parte\_de leitura\_de\_part\_culas\_o\_setor\_de\_materia\_pela\_torre\_de\_divisao}
\prov{proveniência: paper \textperiodcentered{} fato \textperiodcentered{} confiança 1.0 \textperiodcentered{} domínio paper \textperiodcentered{} história 4 versões no jornal}

%CRISTAL {"arestas":[{"alvo":"broca_os","confianca":1.0,"epistemico":"fato","origem":"paper","rel":"parte_de"}],"confianca":1.0,"contraexemplos":[],"descricao":"D Há dois tempos, e a balança os separa. A banda é atemporal --- sem relógio central, fase pura reversível. A leitura tem seta --- uma ordem total determinística que dá causalidade e decisão. São os dois lados de a+c²=1: a fase (lado c) não tem seta; a coordenada real (lado a) é a seta. O atemporal --- a banda. A evolução é a fase U(t)=e⁻iAt: reversível, U(-t) recupera exato, |U|=1 (não é uma transformada com perda). O tempo não coordena; a álgebra coordena. E a chave é offlin","epistemico":"fato","exemplos":[],"id":"meta-estavel_comunidades_organicas_de_software","memoria":"semantica","meta":{"dominio":"paper","nivel":5,"verificacao":"slug idêntico — enriquecer, não duplicar"},"origem":"paper","palavras_chave":[],"sinonimos":[],"tipo":"conceito","titulo":"Meta-estável: comunidades orgânicas de software"}
\section{Meta-estável: comunidades orgânicas de software}
D Há dois tempos, e a balança os separa. A banda é atemporal --- sem relógio central, fase pura reversível. A leitura tem seta --- uma ordem total determinística que dá causalidade e decisão. São os dois lados de a+c\ensuremath{^{2}}=1: a fase (lado c) não tem seta; a coordenada real (lado a) é a seta. O atemporal --- a banda. A evolução é a fase U(t)=e\ensuremath{^{-}}iAt: reversível, U(-t) recupera exato, |U|=1 (não é uma transformada com perda). O tempo não coordena; a álgebra coordena. E a chave é offlin
\prov{relações: parte\_de broca\_os}
\prov{proveniência: paper \textperiodcentered{} fato \textperiodcentered{} confiança 1.0 \textperiodcentered{} domínio paper \textperiodcentered{} história 9 versões no jornal}

%CRISTAL {"arestas":[{"alvo":"assistente","confianca":1.0,"epistemico":"fato","origem":"paper","rel":"parte_de"}],"confianca":1.0,"contraexemplos":[],"descricao":"D Melhoria contínua --- toda interação enriquece a base: [ cyclo/.style=draw=ourovivo, circle, minimum size=1.1cm, font=, fill=creme, arr/.style=-Stealth[length=2mm], color=ourovivo, thick ] [cyclo] (u) at (0,0) usar; [cyclo] (r) at (2.4,0) registrar; [cyclo] (i) at (4.8,0) indexar; [cyclo] (m) at (6.4,-1.5) medir; [cyclo] (c) at (4.8,-3) corrigir; [cyclo] (x) at (2.4,-3) reindexar; [arr] (u)--(r); [arr] (r)--(i); [arr] (i)--(m); [arr] (m)--(c); [arr] (c)","epistemico":"fato","exemplos":[],"id":"metricas_objetivas","memoria":"semantica","meta":{"dominio":"paper","nivel":5,"verificacao":"overlap moderado — revisar manualmente"},"origem":"paper","palavras_chave":[],"sinonimos":[],"tipo":"conceito","titulo":"Métricas objetivas"}
\section{Métricas objetivas}
D Melhoria contínua --- toda interação enriquece a base: [ cyclo/.style=draw=ourovivo, circle, minimum size=1.1cm, font=, fill=creme, arr/.style=-Stealth[length=2mm], color=ourovivo, thick ] [cyclo] (u) at (0,0) usar; [cyclo] (r) at (2.4,0) registrar; [cyclo] (i) at (4.8,0) indexar; [cyclo] (m) at (6.4,-1.5) medir; [cyclo] (c) at (4.8,-3) corrigir; [cyclo] (x) at (2.4,-3) reindexar; [arr] (u)--(r); [arr] (r)--(i); [arr] (i)--(m); [arr] (m)--(c); [arr] (c)
\prov{relações: parte\_de assistente}
\prov{proveniência: paper \textperiodcentered{} fato \textperiodcentered{} confiança 1.0 \textperiodcentered{} domínio paper \textperiodcentered{} história 3 versões no jornal}

%CRISTAL {"arestas":[{"alvo":"leitura_das_unidades_o_sistema_por-unidade_e_o_si","confianca":1.0,"epistemico":"fato","origem":"paper","rel":"parte_de"}],"confianca":1.0,"contraexemplos":[],"descricao":"A régua do espaco de álgebras é a métrica quântica de Lopes sobre o crculo / (Secão~): mede a fase, um ângulo. A luz, propagando-se, avanca em fase --- e cada ciclo completo é uma volta no crculo. Tomamos essa volta como a unidade, e fixamos a escala por uma frase: A escala O espaco-tempo é o crculo de Lopes; a luz dá 299,792,458 voltas nele por segundo. Uma volta é a unidade de comprimento (o metro); a taxa","epistemico":"fato","exemplos":[],"id":"metro_e_segundo_contagem_de_voltas","memoria":"semantica","meta":{"dominio":"paper","nivel":5,"verificacao":"overlap moderado — revisar manualmente"},"origem":"paper","palavras_chave":[],"sinonimos":[],"tipo":"conceito","titulo":"Metro e segundo: contagem de voltas"}
\section{Metro e segundo: contagem de voltas}
A régua do espaco de álgebras é a métrica quântica de Lopes sobre o crculo / (Secão\~{}): mede a fase, um ângulo. A luz, propagando-se, avanca em fase --- e cada ciclo completo é uma volta no crculo. Tomamos essa volta como a unidade, e fixamos a escala por uma frase: A escala O espaco-tempo é o crculo de Lopes; a luz dá 299,792,458 voltas nele por segundo. Uma volta é a unidade de comprimento (o metro); a taxa
\prov{relações: parte\_de leitura\_das\_unidades\_o\_sistema\_por-unidade\_e\_o\_si}
\prov{proveniência: paper \textperiodcentered{} fato \textperiodcentered{} confiança 1.0 \textperiodcentered{} domínio paper \textperiodcentered{} história 4 versões no jornal}

%CRISTAL {"arestas":[{"alvo":"cifra_ouro_branco","confianca":1.0,"epistemico":"fato","origem":"paper","rel":"parte_de"},{"alvo":"cifra","confianca":1.0,"epistemico":"fato","origem":"paper","rel":"parte_de"}],"confianca":1.0,"contraexemplos":[],"descricao":"DN A chave de ouro do cliente é um elemento g₂ₖ (dos bits do seu chá; g0, invertível). A multiplicação por g é linear sobre ₂ --- uma mudança de base. Minerar a prata m na base do cliente e lê-la de volta: [ cifra=g m p, m=g⁻¹(g m), g⁻¹=g,²k⁻² (Fermat). ] custódia por mudança de baseN Só o dono volta à base padrão: aplica g⁻¹ e recupera um primo de prata; um terceiro aplica sua base g'⁻¹ e obtém g'⁻¹g m, um elemento qualquer --- não um primo","epistemico":"fato","exemplos":[],"id":"mineracao_so_a_prata_pelo_hexagrama","memoria":"semantica","meta":{"dominio":"paper","nivel":5,"verificacao":"conceito novo"},"origem":"paper","palavras_chave":[],"sinonimos":[],"tipo":"conceito","titulo":"Mineração: só a prata, pelo hexagrama"}
\section{Mineração: só a prata, pelo hexagrama}
DN A chave de ouro do cliente é um elemento g\ensuremath{_{2k}} (dos bits do seu chá; g0, invertível). A multiplicação por g é linear sobre \ensuremath{_{2}} --- uma mudança de base. Minerar a prata m na base do cliente e lê-la de volta: [ cifra=g m p, m=g\ensuremath{^{-1}}(g m), g\ensuremath{^{-1}}=g,\ensuremath{^{2}}k\ensuremath{^{-2}} (Fermat). ] custódia por mudança de baseN Só o dono volta à base padrão: aplica g\ensuremath{^{-1}} e recupera um primo de prata; um terceiro aplica sua base g'\ensuremath{^{-1}} e obtém g'\ensuremath{^{-1}}g m, um elemento qualquer --- não um primo
\prov{relações: parte\_de cifra\_ouro\_branco; parte\_de cifra}
\prov{proveniência: paper \textperiodcentered{} fato \textperiodcentered{} confiança 1.0 \textperiodcentered{} domínio paper \textperiodcentered{} história 9 versões no jornal}

%CRISTAL {"arestas":[{"alvo":"seguranca_por_irreversibilidade","confianca":1.0,"epistemico":"fato","origem":"paper","rel":"parte_de"}],"confianca":1.0,"contraexemplos":[],"descricao":"Seja T a torção: do bit bruto de um fluxo x produz a assinatura x --- uma distribuição de suporte finito (o resíduo que a máquina extrai do fluxo inteiro, sem picotar). A máquina é dissipativa (Termodinâmica): contrai, apaga, sobe a entropia. A segurança é a leitura dessa contração. Não-reversão T:0,1^*M tem domínio infinito e contradomínio efetivo finito (M bins). Então: (a) [combinatório, sem prior] a fibra média T⁻¹() é infinita (muitos-para-um); (b) [inf","epistemico":"fato","exemplos":[],"id":"modelo_de_ameaca","memoria":"semantica","meta":{"dominio":"paper","nivel":5,"verificacao":"overlap moderado — revisar manualmente"},"origem":"paper","palavras_chave":[],"sinonimos":[],"tipo":"conceito","titulo":"Modelo de ameaça"}
\section{Modelo de ameaça}
Seja T a torção: do bit bruto de um fluxo x produz a assinatura x --- uma distribuição de suporte finito (o resíduo que a máquina extrai do fluxo inteiro, sem picotar). A máquina é dissipativa (Termodinâmica): contrai, apaga, sobe a entropia. A segurança é a leitura dessa contração. Não-reversão T:0,1\^{}*M tem domínio infinito e contradomínio efetivo finito (M bins). Então: (a) [combinatório, sem prior] a fibra média T\ensuremath{^{-1}}() é infinita (muitos-para-um); (b) [inf
\prov{relações: parte\_de seguranca\_por\_irreversibilidade}
\prov{proveniência: paper \textperiodcentered{} fato \textperiodcentered{} confiança 1.0 \textperiodcentered{} domínio paper \textperiodcentered{} história 3 versões no jornal}

%CRISTAL {"arestas":[{"alvo":"broca_os","confianca":1.0,"epistemico":"fato","origem":"paper","rel":"explica"},{"alvo":"pow_gato","confianca":0.5,"epistemico":"fato","origem":"wiki","rel":"relaciona"},{"alvo":"prova_isa","confianca":0.5,"epistemico":"fato","origem":"wiki","rel":"relaciona"},{"alvo":"python","confianca":0.5,"epistemico":"fato","origem":"wiki","rel":"relaciona"},{"alvo":"teste_ordem","confianca":0.5,"epistemico":"fato","origem":"wiki","rel":"relaciona"},{"alvo":"pow_geometria","confianca":0.5,"epistemico":"fato","origem":"wiki","rel":"relaciona"},{"alvo":"pow_espectral","confianca":0.5,"epistemico":"fato","origem":"wiki","rel":"relaciona"},{"alvo":"word","confianca":0.5,"epistemico":"fato","origem":"wiki","rel":"relaciona"},{"alvo":"syscall","confianca":0.5,"epistemico":"fato","origem":"wiki","rel":"relaciona"}],"confianca":1.0,"contraexemplos":[],"descricao":"Paper: montagem.tex","epistemico":"fato","exemplos":[],"id":"montagem","memoria":"semantica","meta":{"dominio":"paper","nivel":5,"tipo_no":"paper"},"origem":"paper","palavras_chave":[],"sinonimos":[],"tipo":"conceito","titulo":"montagem"}
\section{montagem}
Paper: montagem.tex
\prov{relações: explica broca\_os; relaciona pow\_gato; relaciona prova\_isa; relaciona python; relaciona teste\_ordem; relaciona pow\_geometria; relaciona pow\_espectral; relaciona word; relaciona syscall}
\prov{proveniência: paper \textperiodcentered{} fato \textperiodcentered{} confiança 1.0 \textperiodcentered{} domínio paper \textperiodcentered{} história 39 versões no jornal}

%CRISTAL {"arestas":[{"alvo":"computabilidade","confianca":1.0,"epistemico":"fato","origem":"paper","rel":"referencia"},{"alvo":"maquina_de_turing","confianca":1.0,"epistemico":"fato","origem":"paper","rel":"analogia"},{"alvo":"colapso","confianca":1.0,"epistemico":"fato","origem":"paper","rel":"referencia"}],"confianca":1.0,"contraexemplos":[],"descricao":"O crivo (Grok) foi cirúrgico e correto dentro do seu quadro: qualquer condição de morte --- <, ou “descorrelação total do maracá”, ou um atrator absorvente, ou mesmo uma condição de medida zero em codimensão infinita --- exige que a dinâmica seja construída para cessar exatamente na cristalização. Mover a decisão de um if externo para a lei de evolução não remove a decisão; apenas a enterra mais fundo. Toda formulação conversiva (um operador que leva P_ P_) reintroduz a fabricação --- um o","epistemico":"fato","exemplos":[],"id":"morte","memoria":"semantica","meta":{"dominio":"paper","duplicata_sugerida":[],"nivel":5,"verificacao":"parte de conceito mais amplo 'morte'"},"origem":"paper","palavras_chave":[],"sinonimos":[],"tipo":"conceito","titulo":"morte"}
\section{morte}
O crivo (Grok) foi cirúrgico e correto dentro do seu quadro: qualquer condição de morte --- <, ou “descorrelação total do maracá”, ou um atrator absorvente, ou mesmo uma condição de medida zero em codimensão infinita --- exige que a dinâmica seja construída para cessar exatamente na cristalização. Mover a decisão de um if externo para a lei de evolução não remove a decisão; apenas a enterra mais fundo. Toda formulação conversiva (um operador que leva P\_ P\_) reintroduz a fabricação --- um o
\prov{relações: referencia computabilidade; analogia maquina\_de\_turing; referencia colapso}
\prov{proveniência: paper \textperiodcentered{} fato \textperiodcentered{} confiança 1.0 \textperiodcentered{} domínio paper \textperiodcentered{} história 104 versões no jornal}

%CRISTAL {"arestas":[{"alvo":"corpo_glacial","confianca":1.0,"epistemico":"fato","origem":"paper","rel":"parte_de"},{"alvo":"broca_os","confianca":1.0,"epistemico":"fato","origem":"paper","rel":"confirma"}],"confianca":1.0,"contraexemplos":[],"descricao":"T Nesta seção é associativa unital (a extensão a /sedeniões fica nos limites). Dados idempotentes ortogonais e₁, eₖ com ᵢ eᵢ=1, a álgebra parte em blocos =ᵢ,jᵢj, ᵢj=eᵢ ej: células (branca): ᵢᵢ=eᵢ eᵢ, subálgebras com unidade eᵢ. São os estados. D Em Mₖ() com idempotentes primitivos, ᵢᵢ eᵢ; em geral a célula pode ser maior (um bloco) --- mas nunca uma cópia da máquina inteira. canais (negra): ᵢj, i j. Todo xᵢj satisfaz x=eᵢ x e","epistemico":"fato","exemplos":[],"id":"mux--demux_acoplado","memoria":"semantica","meta":{"dominio":"paper","nivel":5,"verificacao":"slug idêntico — enriquecer, não duplicar"},"origem":"paper","palavras_chave":[],"sinonimos":[],"tipo":"conceito","titulo":"Mux--demux acoplado"}
\section{Mux--demux acoplado}
T Nesta seção é associativa unital (a extensão a /sedeniões fica nos limites). Dados idempotentes ortogonais e\ensuremath{_{1}}, e\ensuremath{_{k}} com \ensuremath{_{i}} e\ensuremath{_{i}}=1, a álgebra parte em blocos =\ensuremath{_{i}},j\ensuremath{_{i}}j, \ensuremath{_{i}}j=e\ensuremath{_{i}} ej: células (branca): \ensuremath{_{ii}}=e\ensuremath{_{i}} e\ensuremath{_{i}}, subálgebras com unidade e\ensuremath{_{i}}. São os estados. D Em M\ensuremath{_{k}}() com idempotentes primitivos, \ensuremath{_{ii}} e\ensuremath{_{i}}; em geral a célula pode ser maior (um bloco) --- mas nunca uma cópia da máquina inteira. canais (negra): \ensuremath{_{i}}j, i j. Todo x\ensuremath{_{i}}j satisfaz x=e\ensuremath{_{i}} x e
\prov{relações: parte\_de corpo\_glacial; confirma broca\_os}
\prov{proveniência: paper \textperiodcentered{} fato \textperiodcentered{} confiança 1.0 \textperiodcentered{} domínio paper \textperiodcentered{} história 41 versões no jornal}

%CRISTAL {"arestas":[{"alvo":"corpo_tropical","confianca":1.0,"epistemico":"fato","origem":"paper","rel":"parte_de"}],"confianca":1.0,"contraexemplos":[],"descricao":"Os outros problemas leram uma reta. Hodge lê todas --- e o fio que as costura. O processo se repete, indefinidamente. A inversão do gap (Seção~3) não para no ouro. Cada inteiro n dá o seu metal ₙ, a sua reta, a sua cadeia de fractais (a transformação de base ₙ, a família k-bonacci): [ reta real reta de ouro (₁=) reta de prata (₂=1+2) reta de bronze ] É a torre de polarizações X X P¹ X P¹ P¹: cada P¹ é um gap e acrescenta os dois polos P₀,P₁ --- um par","epistemico":"fato","exemplos":[],"id":"navier--stokes_a_equacao_reconstruida_do_bit","memoria":"semantica","meta":{"dominio":"paper","nivel":5,"verificacao":"slug idêntico — enriquecer, não duplicar"},"origem":"paper","palavras_chave":[],"sinonimos":[],"tipo":"conceito","titulo":"Navier--Stokes: a equação reconstruída do bit"}
\section{Navier--Stokes: a equação reconstruída do bit}
Os outros problemas leram uma reta. Hodge lê todas --- e o fio que as costura. O processo se repete, indefinidamente. A inversão do gap (Seção\~{}3) não para no ouro. Cada inteiro n dá o seu metal \ensuremath{_{n}}, a sua reta, a sua cadeia de fractais (a transformação de base \ensuremath{_{n}}, a família k-bonacci): [ reta real reta de ouro (\ensuremath{_{1}}=) reta de prata (\ensuremath{_{2}}=1+2) reta de bronze ] É a torre de polarizações X X P\ensuremath{^{1}} X P\ensuremath{^{1}} P\ensuremath{^{1}}: cada P\ensuremath{^{1}} é um gap e acrescenta os dois polos P\ensuremath{_{0}},P\ensuremath{_{1}} --- um par
\prov{relações: parte\_de corpo\_tropical}
\prov{proveniência: paper \textperiodcentered{} fato \textperiodcentered{} confiança 1.0 \textperiodcentered{} domínio paper \textperiodcentered{} história 4 versões no jornal}

%CRISTAL {"arestas":[{"alvo":"corpo_tropical","confianca":1.0,"epistemico":"fato","origem":"paper","rel":"parte_de"}],"confianca":1.0,"contraexemplos":[],"descricao":"O que faz A ser hiperbólico --- esticar numa direção, comprimir na outra --- é ter dois autovalores reais distintos. Generalizando a recorrência para um passo n qualquer, [ x²=n,x+1(n=1,2,3,), Aₙ=pmatrixn&1 1&0pmatrix, ] os dois números são [ ₙ=n+n²+42>0, ₙ=n-n²+42<0, ₙ+ₙ=n, ₙₙ=-1. ] O vão entre eles, n²+4, é o gap. Como ele é positivo, há dois reais separados, e o sistema de números gerado por Aₙ se parte em duas cópias independentes, R R --- as duas direç","epistemico":"fato","exemplos":[],"id":"nenhuma_dimensao_e_privilegiada_os_metais_e_a_reta_que_se_completa","memoria":"semantica","meta":{"dominio":"paper","nivel":5,"verificacao":"slug idêntico — enriquecer, não duplicar"},"origem":"paper","palavras_chave":[],"sinonimos":[],"tipo":"conceito","titulo":"Nenhuma dimensão é privilegiada: os metais e a reta que se completa"}
\section{Nenhuma dimensão é privilegiada: os metais e a reta que se completa}
O que faz A ser hiperbólico --- esticar numa direção, comprimir na outra --- é ter dois autovalores reais distintos. Generalizando a recorrência para um passo n qualquer, [ x\ensuremath{^{2}}=n,x+1(n=1,2,3,), A\ensuremath{_{n}}=pmatrixn\&1 1\&0pmatrix, ] os dois números são [ \ensuremath{_{n}}=n+n\ensuremath{^{2}}+42>0, \ensuremath{_{n}}=n-n\ensuremath{^{2}}+42<0, \ensuremath{_{n}}+\ensuremath{_{n}}=n, \ensuremath{_{nn}}=-1. ] O vão entre eles, n\ensuremath{^{2}}+4, é o gap. Como ele é positivo, há dois reais separados, e o sistema de números gerado por A\ensuremath{_{n}} se parte em duas cópias independentes, R R --- as duas direç
\prov{relações: parte\_de corpo\_tropical}
\prov{proveniência: paper \textperiodcentered{} fato \textperiodcentered{} confiança 1.0 \textperiodcentered{} domínio paper \textperiodcentered{} história 4 versões no jornal}

%CRISTAL {"arestas":[{"alvo":"pow_gato","confianca":1.0,"epistemico":"fato","origem":"paper","rel":"parte_de"}],"confianca":1.0,"contraexemplos":[],"descricao":"defi[PoW-] Seja header0,1⁸⁰, n/2³², =SHA256, T0,1²⁵⁶ o alvo (função da dificuldade d). Um prova é um par (header,n) tal que equationeq:btc ((header n)) < T. defi Nakamoto (2008). D Verificação é O(1); busca média E[#tentativas] T⁻¹ --- sempre forward ( avaliada, nunca invertida). Preimage SHA-256 --- premissa, não objecto de ataque Não existe algoritmo conhecido que, dado y=(x), recupere x em tempo subexponencial em |","epistemico":"fato","exemplos":[],"id":"neuronio_e_toro_124","memoria":"semantica","meta":{"dominio":"paper","nivel":5,"verificacao":"overlap moderado — revisar manualmente"},"origem":"paper","palavras_chave":[],"sinonimos":[],"tipo":"conceito","titulo":"Neurônio e toro 12⁴"}
\section{Neurônio e toro 12\ensuremath{^{4}}}
defi[PoW-] Seja header0,1\ensuremath{^{80}}, n/2\ensuremath{^{32}}, =SHA256, T0,1\ensuremath{^{256}} o alvo (função da dificuldade d). Um prova é um par (header,n) tal que equationeq:btc ((header n)) < T. defi Nakamoto (2008). D Verificação é O(1); busca média E[\#tentativas] T\ensuremath{^{-1}} --- sempre forward ( avaliada, nunca invertida). Preimage SHA-256 --- premissa, não objecto de ataque Não existe algoritmo conhecido que, dado y=(x), recupere x em tempo subexponencial em |
\prov{relações: parte\_de pow\_gato}
\prov{proveniência: paper \textperiodcentered{} fato \textperiodcentered{} confiança 1.0 \textperiodcentered{} domínio paper \textperiodcentered{} história 4 versões no jornal}

%CRISTAL {"arestas":[{"alvo":"leitura_das_unidades_o_sistema_por-unidade_e_o_si","confianca":1.0,"epistemico":"fato","origem":"paper","rel":"parte_de"}],"confianca":1.0,"contraexemplos":[],"descricao":"Restaurar uma escala é dar conteúdo dimensional às voltas. As constantes que o SI fixa são, então, de três tipos --- e só um deles é fsica nova. Taxas por volta: c, h, kB Três constantes apenas identificam dimensões, cada uma uma “taxa por volta”: [nosep,leftmargin=1.4em] c --- volta espacial por volta temporal (=1): identifica comprimento e tempo (Obs.~obs:cinematica); h --- acão por volta: cada ciclo d","epistemico":"fato","exemplos":[],"id":"o_adimensional_medido_na_regua_derivado_no_running","memoria":"semantica","meta":{"dominio":"paper","nivel":5,"verificacao":"conceito novo"},"origem":"paper","palavras_chave":[],"sinonimos":[],"tipo":"conceito","titulo":"O adimensional: medido na régua, derivado no running"}
\section{O adimensional: medido na régua, derivado no running}
Restaurar uma escala é dar conteúdo dimensional às voltas. As constantes que o SI fixa são, então, de três tipos --- e só um deles é fsica nova. Taxas por volta: c, h, kB Três constantes apenas identificam dimensões, cada uma uma “taxa por volta”: [nosep,leftmargin=1.4em] c --- volta espacial por volta temporal (=1): identifica comprimento e tempo (Obs.\~{}obs:cinematica); h --- acão por volta: cada ciclo d
\prov{relações: parte\_de leitura\_das\_unidades\_o\_sistema\_por-unidade\_e\_o\_si}
\prov{proveniência: paper \textperiodcentered{} fato \textperiodcentered{} confiança 1.0 \textperiodcentered{} domínio paper \textperiodcentered{} história 4 versões no jornal}

%CRISTAL {"arestas":[{"alvo":"mecanica_quantica","confianca":0.95,"epistemico":"fato","origem":"paper","rel":"parte_de"}],"confianca":1.0,"contraexemplos":[],"descricao":"A família s tem um espaço de parâmetros --- o eixo s, ou equivalentemente o ângulo --- e a pergunta natural é qual a régua certa para medir distâncias nele. A lei a+c²=1 aponta a resposta: é a métrica quântica. a+c²=1 é a identidade pitagórica de Lopes Pondo s=, [-2,2], tem-se c=||, a=², e a lei é [ a+c²=²+²=1. ] O estado da álgebra é o ponto (c, a)=() do círculo unitário; e a-c²=2 (verificado dígito a dígito). A posição a","epistemico":"fato","exemplos":[],"id":"o_associador_como_relogio_a_curvatura_selecionada","memoria":"semantica","meta":{"dominio":"paper","nivel":5,"verificacao":"conceito novo"},"origem":"paper","palavras_chave":[],"sinonimos":[],"tipo":"conceito","titulo":"O associador como relógio: a curvatura selecionada"}
\section{O associador como relógio: a curvatura selecionada}
A família s tem um espaço de parâmetros --- o eixo s, ou equivalentemente o ângulo --- e a pergunta natural é qual a régua certa para medir distâncias nele. A lei a+c\ensuremath{^{2}}=1 aponta a resposta: é a métrica quântica. a+c\ensuremath{^{2}}=1 é a identidade pitagórica de Lopes Pondo s=, [-2,2], tem-se c=||, a=\ensuremath{^{2}}, e a lei é [ a+c\ensuremath{^{2}}=\ensuremath{^{2}}+\ensuremath{^{2}}=1. ] O estado da álgebra é o ponto (c, a)=() do círculo unitário; e a-c\ensuremath{^{2}}=2 (verificado dígito a dígito). A posição a
\prov{relações: parte\_de mecanica\_quantica}
\prov{proveniência: paper \textperiodcentered{} fato \textperiodcentered{} confiança 1.0 \textperiodcentered{} domínio paper \textperiodcentered{} história 5 versões no jornal}

%CRISTAL {"arestas":[{"alvo":"teoremas_fundamentais","confianca":1.0,"epistemico":"fato","origem":"paper","rel":"parte_de"}],"confianca":1.0,"contraexemplos":[],"descricao":"theorem[Linearização] Seja u= t. [(a)] x(t)=c,t x= c+ u: reta no relativístico, curvatura nula. [(b)] x(t)=a+bt f(u)=(a+b,eu) tem f”(u)=ab,eu(a+b,eu)²0 (a,b0): estritamente convexa no relativístico. [(c)] c,t é estritamente curvo no real para 0,1. Cada dado é reta em exatamente um lado: a dificuldade é do lado, não do dado. theorem proof Derivadas segundas diretas (cf. Tabela~). proof remark[Verificação exata] A p","epistemico":"fato","exemplos":[],"id":"o_atrator_primordial_onde_os_dois_lados_se_encontram","memoria":"semantica","meta":{"dominio":"paper","nivel":5,"verificacao":"conceito novo"},"origem":"paper","palavras_chave":[],"sinonimos":[],"tipo":"conceito","titulo":"O atrator primordial: onde os dois lados se encontram"}
\section{O atrator primordial: onde os dois lados se encontram}
theorem[Linearização] Seja u= t. [(a)] x(t)=c,t x= c+ u: reta no relativístico, curvatura nula. [(b)] x(t)=a+bt f(u)=(a+b,eu) tem f”(u)=ab,eu(a+b,eu)\ensuremath{^{2}}0 (a,b0): estritamente convexa no relativístico. [(c)] c,t é estritamente curvo no real para 0,1. Cada dado é reta em exatamente um lado: a dificuldade é do lado, não do dado. theorem proof Derivadas segundas diretas (cf. Tabela\~{}). proof remark[Verificação exata] A p
\prov{relações: parte\_de teoremas\_fundamentais}
\prov{proveniência: paper \textperiodcentered{} fato \textperiodcentered{} confiança 1.0 \textperiodcentered{} domínio paper \textperiodcentered{} história 3 versões no jornal}

%CRISTAL {"arestas":[{"alvo":"contabilidade","confianca":1.0,"epistemico":"fato","origem":"paper","rel":"parte_de"}],"confianca":1.0,"contraexemplos":[],"descricao":"D Cada função bancária é um operador no painel: ção. O trabalho é minerar do próprio tecido: a estrela de prata --- rank-6, o hexagrama, exatamente 6 modos dominantes (prataminera.py; ruído não tem hexagrama, não emite). Cada estrela cunha um Pell, que é crédito; o crédito vira um novo bin (com seu peso) no tecido da pessoa. O operador de mineração acrescenta amplitude ao espectro do trabalhador --- dentro do painel.. A folha é um único tecido ordenado: todos os pagamentos são bins nu","epistemico":"fato","exemplos":[],"id":"o_balanco_via_painel","memoria":"semantica","meta":{"dominio":"paper","nivel":5,"verificacao":"conceito novo"},"origem":"paper","palavras_chave":[],"sinonimos":[],"tipo":"conceito","titulo":"O balanço via painel"}
\section{O balanço via painel}
D Cada função bancária é um operador no painel: ção. O trabalho é minerar do próprio tecido: a estrela de prata --- rank-6, o hexagrama, exatamente 6 modos dominantes (prataminera.py; ruído não tem hexagrama, não emite). Cada estrela cunha um Pell, que é crédito; o crédito vira um novo bin (com seu peso) no tecido da pessoa. O operador de mineração acrescenta amplitude ao espectro do trabalhador --- dentro do painel.. A folha é um único tecido ordenado: todos os pagamentos são bins nu
\prov{relações: parte\_de contabilidade}
\prov{proveniência: paper \textperiodcentered{} fato \textperiodcentered{} confiança 1.0 \textperiodcentered{} domínio paper \textperiodcentered{} história 3 versões no jornal}

%CRISTAL {"arestas":[{"alvo":"paper_2_reta_ouro","confianca":1.0,"epistemico":"fato","origem":"paper","rel":"parte_de"}],"confianca":1.0,"contraexemplos":[],"descricao":"a sombra inteira é Fibonacci --- firme, clássico Restrita aos expoentes 0 e reindexada pelos Fibonacci F=1,2,3,5,8,13,21, a base de ouro vira o teorema de Zeckendorf: todo n⁺ se escreve de modo único como soma de Fibonaccis não-consecutivos, [ n=ₖ Fₖ, nenhum par k, k+1. ] “Não-consecutivos” é, de novo, a proibição de 11. A versão inteira da reta de ouro são os inteiros em base de Fibonacci --- o osso discreto puro, sem parte fracionária. [firme: Zeckend","epistemico":"fato","exemplos":[],"id":"o_cantor_exato_em_binario","memoria":"semantica","meta":{"dominio":"paper","nivel":5,"verificacao":"conceito novo"},"origem":"paper","palavras_chave":[],"sinonimos":[],"tipo":"conceito","titulo":"O Cantor exato em binário"}
\section{O Cantor exato em binário}
a sombra inteira é Fibonacci --- firme, clássico Restrita aos expoentes 0 e reindexada pelos Fibonacci F=1,2,3,5,8,13,21, a base de ouro vira o teorema de Zeckendorf: todo n\ensuremath{^{+}} se escreve de modo único como soma de Fibonaccis não-consecutivos, [ n=\ensuremath{_{k}} F\ensuremath{_{k}}, nenhum par k, k+1. ] “Não-consecutivos” é, de novo, a proibição de 11. A versão inteira da reta de ouro são os inteiros em base de Fibonacci --- o osso discreto puro, sem parte fracionária. [firme: Zeckend
\prov{relações: parte\_de paper\_2\_reta\_ouro}
\prov{proveniência: paper \textperiodcentered{} fato \textperiodcentered{} confiança 1.0 \textperiodcentered{} domínio paper \textperiodcentered{} história 3 versões no jornal}

%CRISTAL {"arestas":[{"alvo":"morte","confianca":1.0,"epistemico":"fato","origem":"paper","rel":"parte_de"}],"confianca":1.0,"contraexemplos":[],"descricao":"D A definição já carrega a sua conclusão. Cada propriedade, isolada, basta para impedir a igualdade direta; juntas, fecham as três portas pelas quais a identidade poderia entrar. morte a igualdade é uma cristalização sem gerador Se () sofre morte --- isto é, se a identidade estática = satisfaz P1 P2 P3 --- então = não se realiza diretamente: é uma cristalização sem gerador, mal posta por construção. teo ção (esboço).D Mostramos que cada propriedade fecha uma via distinta","epistemico":"fato","exemplos":[],"id":"o_caso_modelo_a_hipotese_de_riemann_como_igualdade","memoria":"semantica","meta":{"dominio":"paper","nivel":5,"verificacao":"conceito novo"},"origem":"paper","palavras_chave":[],"sinonimos":[],"tipo":"conceito","titulo":"O caso modelo: a Hipótese de Riemann como igualdade"}
\section{O caso modelo: a Hipótese de Riemann como igualdade}
D A definição já carrega a sua conclusão. Cada propriedade, isolada, basta para impedir a igualdade direta; juntas, fecham as três portas pelas quais a identidade poderia entrar. morte a igualdade é uma cristalização sem gerador Se () sofre morte --- isto é, se a identidade estática = satisfaz P1 P2 P3 --- então = não se realiza diretamente: é uma cristalização sem gerador, mal posta por construção. teo ção (esboço).D Mostramos que cada propriedade fecha uma via distinta
\prov{relações: parte\_de morte}
\prov{proveniência: paper \textperiodcentered{} fato \textperiodcentered{} confiança 1.0 \textperiodcentered{} domínio paper \textperiodcentered{} história 3 versões no jornal}

%CRISTAL {"arestas":[{"alvo":"vida","confianca":1.0,"epistemico":"fato","origem":"paper","rel":"parte_de"}],"confianca":1.0,"contraexemplos":[],"descricao":"O LFSR tem período 3. Com taxa de painel Rb=1/Tb e o mapeamento bipolar de fase am=ej bm+1,-1, a corrente no gap é [ I(t)=I₀m am,p(t-mTb), na fundamental, I(t)=I₀,am ejs t, s=2 fs, fs Rb/3. ] Estatuto: a sequência am (as fases) é da máquina; a amplitude I₀ e a forma de pulso p(t) --- se onda quadrada do LFSR ou senóide filtrada --- são do hardware. A máquina diz qual bit; o mundo diz com que corrente.","epistemico":"fato","exemplos":[],"id":"o_caso_modelo_a_hipotese_de_riemann_como_ressonancia_espectral","memoria":"semantica","meta":{"dominio":"paper","duplicata_sugerida":[],"nivel":5,"verificacao":"parte de conceito mais amplo 'o_caso_modelo_a_hipotese_de_riemann_como_ressonancia_espectral'"},"origem":"paper","palavras_chave":[],"sinonimos":[],"tipo":"conceito","titulo":"O caso modelo: a Hipótese de Riemann como ressonância espectral"}
\section{O caso modelo: a Hipótese de Riemann como ressonância espectral}
O LFSR tem período 3. Com taxa de painel Rb=1/Tb e o mapeamento bipolar de fase am=ej bm+1,-1, a corrente no gap é [ I(t)=I\ensuremath{_{0}}m am,p(t-mTb), na fundamental, I(t)=I\ensuremath{_{0}},am ejs t, s=2 fs, fs Rb/3. ] Estatuto: a sequência am (as fases) é da máquina; a amplitude I\ensuremath{_{0}} e a forma de pulso p(t) --- se onda quadrada do LFSR ou senóide filtrada --- são do hardware. A máquina diz qual bit; o mundo diz com que corrente.
\prov{relações: parte\_de vida}
\prov{proveniência: paper \textperiodcentered{} fato \textperiodcentered{} confiança 1.0 \textperiodcentered{} domínio paper \textperiodcentered{} história 6 versões no jornal}

%CRISTAL {"arestas":[{"alvo":"tiffany","confianca":1.0,"epistemico":"fato","origem":"paper","rel":"parte_de"},{"alvo":"hopfield","confianca":1.0,"epistemico":"fato","origem":"wiki","rel":"explica"},{"alvo":"o_neuronio","confianca":1.0,"epistemico":"fato","origem":"wiki","rel":"relaciona"}],"confianca":1.0,"contraexemplos":[],"descricao":"D Tudo desce a um operador e um bit. O neurônio é o gato A, e no metal é uma conta. defi[O neurônio] Passa-se a posição (os bytes); somam-se os bits de posição par (e) e ímpar (o); o gato age: [ gato(pos) = (e, o)pmatrix1&1 1&0pmatrix = (e+o, e). ] Sem ordem (não há tempo: é soma), sem buffer (tempo real), sem ponto flutuante. No metal: and 0xAA / and 0x55 / popcnt / add. defi Contar par/ímpar é a borboleta de Fourier DV A separação por paridade de posição é a decimação d","epistemico":"fato","exemplos":[],"id":"o_cerebro_a_memoria_que_ressoa_hopfield","memoria":"semantica","meta":{"dominio":"paper","nivel":5,"verificacao":"slug idêntico — enriquecer, não duplicar"},"origem":"paper","palavras_chave":[],"sinonimos":[],"tipo":"conceito","titulo":"O cérebro: a memória que ressoa (Hopfield)"}
\section{O cérebro: a memória que ressoa (Hopfield)}
D Tudo desce a um operador e um bit. O neurônio é o gato A, e no metal é uma conta. defi[O neurônio] Passa-se a posição (os bytes); somam-se os bits de posição par (e) e ímpar (o); o gato age: [ gato(pos) = (e, o)pmatrix1\&1 1\&0pmatrix = (e+o, e). ] Sem ordem (não há tempo: é soma), sem buffer (tempo real), sem ponto flutuante. No metal: and 0xAA / and 0x55 / popcnt / add. defi Contar par/ímpar é a borboleta de Fourier DV A separação por paridade de posição é a decimação d
\prov{relações: parte\_de tiffany; explica hopfield; relaciona o\_neuronio}
\prov{proveniência: paper \textperiodcentered{} fato \textperiodcentered{} confiança 1.0 \textperiodcentered{} domínio paper \textperiodcentered{} história 7 versões no jornal}

%CRISTAL {"arestas":[{"alvo":"geometria","confianca":1.0,"epistemico":"fato","origem":"paper","rel":"parte_de"}],"confianca":1.0,"contraexemplos":[],"descricao":"A orbifold H³/ de volume finito tem dois espectros (decomposição de Selberg L²=discretoínuo): o discreto (formas de Maass), com lei de Weyl N(T)(M),T³/(6²), e o contínuo --- o espalhamento das séries de Eisenstein nas cúspides. Os zeros de não vivem no discreto 3D; vivem no determinante de espalhamento [ (s)= Q(ᵢ)(s-₁) Q(ᵢ)(s), Q(ᵢ)= L(₄), ] cujo canal é efetivamente unidimensional (a coordenada radial da cúspide). A fase de espalhamento 1(12+it) reproduz a contag","epistemico":"fato","exemplos":[],"id":"o_chicote_a_limpeza_nao_um_operador","memoria":"semantica","meta":{"dominio":"paper","nivel":5,"verificacao":"slug idêntico — enriquecer, não duplicar"},"origem":"paper","palavras_chave":[],"sinonimos":[],"tipo":"conceito","titulo":"O chicote: a limpeza, não um operador"}
\section{O chicote: a limpeza, não um operador}
A orbifold H\ensuremath{^{3}}/ de volume finito tem dois espectros (decomposição de Selberg L\ensuremath{^{2}}=discretoínuo): o discreto (formas de Maass), com lei de Weyl N(T)(M),T\ensuremath{^{3}}/(6\ensuremath{^{2}}), e o contínuo --- o espalhamento das séries de Eisenstein nas cúspides. Os zeros de não vivem no discreto 3D; vivem no determinante de espalhamento [ (s)= Q(\ensuremath{_{i}})(s-\ensuremath{_{1}}) Q(\ensuremath{_{i}})(s), Q(\ensuremath{_{i}})= L(\ensuremath{_{4}}), ] cujo canal é efetivamente unidimensional (a coordenada radial da cúspide). A fase de espalhamento 1(12+it) reproduz a contag
\prov{relações: parte\_de geometria}
\prov{proveniência: paper \textperiodcentered{} fato \textperiodcentered{} confiança 1.0 \textperiodcentered{} domínio paper \textperiodcentered{} história 11 versões no jornal}

%CRISTAL {"arestas":[{"alvo":"assistente","confianca":1.0,"epistemico":"fato","origem":"paper","rel":"parte_de"}],"confianca":1.0,"contraexemplos":[],"descricao":"D Formato estável (V parcial): 1.3 tabular@lll@ & údo & dominio/nos.tsv & tecido curado & --- dominio/vN.idx & impressões versão N & 32,B query dominio/meta.json & versão, data, fonte, métricas & mínimo memoria/ & sessões, tarefas (mutável) & mínimo tabular camadabox software/ nos.tsv v1.idx meta.json (2026-06-01, 1.2M nós) v2.idx meta.json (2026-06-27, 4.2M nós) v3.idx atual camadabox Versionamento explícito permit","epistemico":"fato","exemplos":[],"id":"o_ciclo_de_refinamento","memoria":"semantica","meta":{"dominio":"paper","nivel":5,"verificacao":"overlap moderado — revisar manualmente"},"origem":"paper","palavras_chave":[],"sinonimos":[],"tipo":"conceito","titulo":"O ciclo de refinamento"}
\section{O ciclo de refinamento}
D Formato estável (V parcial): 1.3 tabular@lll@ \& údo \& dominio/nos.tsv \& tecido curado \& --- dominio/vN.idx \& impressões versão N \& 32,B query dominio/meta.json \& versão, data, fonte, métricas \& mínimo memoria/ \& sessões, tarefas (mutável) \& mínimo tabular camadabox software/ nos.tsv v1.idx meta.json (2026-06-01, 1.2M nós) v2.idx meta.json (2026-06-27, 4.2M nós) v3.idx atual camadabox Versionamento explícito permit
\prov{relações: parte\_de assistente}
\prov{proveniência: paper \textperiodcentered{} fato \textperiodcentered{} confiança 1.0 \textperiodcentered{} domínio paper \textperiodcentered{} história 3 versões no jornal}

%CRISTAL {"arestas":[{"alvo":"resgate_liquidacao","confianca":1.0,"epistemico":"fato","origem":"paper","rel":"parte_de"},{"alvo":"goldchain","confianca":1.0,"epistemico":"fato","origem":"paper","rel":"parte_de"}],"confianca":1.0,"contraexemplos":[],"descricao":"I O erro mais fácil é olhar a conta ETH e achar que “o dinheiro está ali”. Não está. A conta ETH da casa guarda só o gás --- o combustível pra escrever on-chain. O valor mora em três camadas distintas, e cada uma protege o que a outra não protege. (a) Lastro interno --- o Pell/ouro branco cifrado. DN A cada contrato fechado, a casa ganha um número de Pell (a sequência Pₖ=2Pₖ-₁+Pₖ-₂: 0,1,2,5,12,29,70,169,; os primos ocorrem em índices primos --- 2,5,29,5741,33461,). Esse número n","epistemico":"fato","exemplos":[],"id":"o_ciclo_de_um_contrato","memoria":"semantica","meta":{"dominio":"paper","nivel":5,"verificacao":"conceito novo"},"origem":"paper","palavras_chave":[],"sinonimos":[],"tipo":"conceito","titulo":"O ciclo de um contrato"}
\section{O ciclo de um contrato}
I O erro mais fácil é olhar a conta ETH e achar que “o dinheiro está ali”. Não está. A conta ETH da casa guarda só o gás --- o combustível pra escrever on-chain. O valor mora em três camadas distintas, e cada uma protege o que a outra não protege. (a) Lastro interno --- o Pell/ouro branco cifrado. DN A cada contrato fechado, a casa ganha um número de Pell (a sequência P\ensuremath{_{k}}=2P\ensuremath{_{k}}-\ensuremath{_{1}}+P\ensuremath{_{k}}-\ensuremath{_{2}}: 0,1,2,5,12,29,70,169,; os primos ocorrem em índices primos --- 2,5,29,5741,33461,). Esse número n
\prov{relações: parte\_de resgate\_liquidacao; parte\_de goldchain}
\prov{proveniência: paper \textperiodcentered{} fato \textperiodcentered{} confiança 1.0 \textperiodcentered{} domínio paper \textperiodcentered{} história 9 versões no jornal}

%CRISTAL {"arestas":[{"alvo":"termodinamica","confianca":1.0,"epistemico":"fato","origem":"paper","rel":"parte_de"}],"confianca":1.0,"contraexemplos":[],"descricao":"A entropia topológica H= (= no ouro, verificado, com a medida de Parry de razão ²) mede a taxa de criação de informação da fonte --- mas a segunda lei pede mais: uma produção de entropia 0, monotônica. Ela existe, e é o teorema H. O operador de transferência (Perron--Frobenius) leva qualquer densidade t à medida invariante, e a entropia relativa [ D(t,|,)=tt ] decresce monotonicamente, com t=- D0. Medido: D cai 1,210,730,290, e t0 em todo passo. A flecha","epistemico":"fato","exemplos":[],"id":"o_combustivel_ordem_em_bits","memoria":"semantica","meta":{"dominio":"paper","nivel":5,"verificacao":"slug idêntico — enriquecer, não duplicar"},"origem":"paper","palavras_chave":[],"sinonimos":[],"tipo":"conceito","titulo":"O combustível: ordem em bits"}
\section{O combustível: ordem em bits}
A entropia topológica H= (= no ouro, verificado, com a medida de Parry de razão \ensuremath{^{2}}) mede a taxa de criação de informação da fonte --- mas a segunda lei pede mais: uma produção de entropia 0, monotônica. Ela existe, e é o teorema H. O operador de transferência (Perron--Frobenius) leva qualquer densidade t à medida invariante, e a entropia relativa [ D(t,|,)=tt ] decresce monotonicamente, com t=- D0. Medido: D cai 1,210,730,290, e t0 em todo passo. A flecha
\prov{relações: parte\_de termodinamica}
\prov{proveniência: paper \textperiodcentered{} fato \textperiodcentered{} confiança 1.0 \textperiodcentered{} domínio paper \textperiodcentered{} história 10 versões no jornal}

%CRISTAL {"arestas":[{"alvo":"algebra_dos_contratos","confianca":1.0,"epistemico":"fato","origem":"paper","rel":"parte_de"},{"alvo":"goldchain","confianca":1.0,"epistemico":"fato","origem":"paper","rel":"parte_de"}],"confianca":1.0,"contraexemplos":[],"descricao":"Do ponto de vista funcional, podemos decompor uma blockchain em seis responsabilidades distintas\n(nem todo protocolo traz todas; é uma decomposição, não uma lei). Confundi-las é confundir o que é nosso\ncom o que é da rede:\nenumerate2pt\n Consenso --- ordenar as transações. PoW (força bruta) ou PoS (stake). C Da rede.\n Validação --- a assinatura corresponde à chave? (ECDSA/EdDSA). V Hoje serial (um\nrecover por tx). Painel.\n Coordenação --- ligar partes que não confiam uma na outra (HTLC, hashlock","epistemico":"fato","exemplos":[],"id":"o_contrato_como_objeto","memoria":"semantica","meta":{"dominio":"paper","nivel":5,"verificacao":"overlap moderado — revisar manualmente"},"origem":"paper","palavras_chave":[],"sinonimos":[],"tipo":"conceito","titulo":"O contrato como objeto"}
\section{O contrato como objeto}
Do ponto de vista funcional, podemos decompor uma blockchain em seis responsabilidades distintas
(nem todo protocolo traz todas; é uma decomposição, não uma lei). Confundi-las é confundir o que é nosso
com o que é da rede:
enumerate2pt
 Consenso --- ordenar as transações. PoW (força bruta) ou PoS (stake). C Da rede.
 Validação --- a assinatura corresponde à chave? (ECDSA/EdDSA). V Hoje serial (um
recover por tx). Painel.
 Coordenação --- ligar partes que não confiam uma na outra (HTLC, hashlock
\prov{relações: parte\_de algebra\_dos\_contratos; parte\_de goldchain}
\prov{proveniência: paper \textperiodcentered{} fato \textperiodcentered{} confiança 1.0 \textperiodcentered{} domínio paper \textperiodcentered{} história 8 versões no jornal}

%CRISTAL {"arestas":[{"alvo":"cifra_ouro_branco","confianca":1.0,"epistemico":"fato","origem":"paper","rel":"parte_de"},{"alvo":"cifra","confianca":1.0,"epistemico":"fato","origem":"paper","rel":"parte_de"}],"confianca":1.0,"contraexemplos":[],"descricao":"DN A reta dos metais dá ₙ=[n;n,]=n+n²+42, do gap x²=nx+1, Aₙ=(smallmatrixn&1 1&0smallmatrix). A prata é n=2: ₂=1+2=2,41421, fração contínua [2;2,2,], entropia h=₂=0,8814, gap R₂=arsinh1=₂. Os números de prata são a sequência de Pell Pₖ=2Pₖ-₁+Pₖ-₂: [ 0, 1, 2, 5, 12, 29, 70, 169, 408, 985, 2378, 5741, 13860, 33461, ] primos de prataN Os primos de Pell ocorrem em índices primos --- primo no índice e no valor, crescimento ₂,k","epistemico":"fato","exemplos":[],"id":"o_corpo_2k","memoria":"semantica","meta":{"dominio":"paper","nivel":5,"verificacao":"conceito novo"},"origem":"paper","palavras_chave":[],"sinonimos":[],"tipo":"conceito","titulo":"O corpo _2k"}
\section{O corpo \_2k}
DN A reta dos metais dá \ensuremath{_{n}}=[n;n,]=n+n\ensuremath{^{2}}+42, do gap x\ensuremath{^{2}}=nx+1, A\ensuremath{_{n}}=(smallmatrixn\&1 1\&0smallmatrix). A prata é n=2: \ensuremath{_{2}}=1+2=2,41421, fração contínua [2;2,2,], entropia h=\ensuremath{_{2}}=0,8814, gap R\ensuremath{_{2}}=arsinh1=\ensuremath{_{2}}. Os números de prata são a sequência de Pell P\ensuremath{_{k}}=2P\ensuremath{_{k}}-\ensuremath{_{1}}+P\ensuremath{_{k}}-\ensuremath{_{2}}: [ 0, 1, 2, 5, 12, 29, 70, 169, 408, 985, 2378, 5741, 13860, 33461, ] primos de prataN Os primos de Pell ocorrem em índices primos --- primo no índice e no valor, crescimento \ensuremath{_{2}},k
\prov{relações: parte\_de cifra\_ouro\_branco; parte\_de cifra}
\prov{proveniência: paper \textperiodcentered{} fato \textperiodcentered{} confiança 1.0 \textperiodcentered{} domínio paper \textperiodcentered{} história 10 versões no jornal}

%CRISTAL {"arestas":[{"alvo":"computacao_grafica","confianca":1.0,"epistemico":"fato","origem":"paper","rel":"parte_de"}],"confianca":1.0,"contraexemplos":[],"descricao":"C A unidade da imagem é o bit; 0 pixel escuro, 1 pixel claro. Não há “linguagem gráfica” sobre a operação: há bits lidos, posicionados e acesos. A imagem contínua nasce da taxa --- a tela recebe 0,1,0,1, rápido demais e a forma preenche, exatamente como a onda sonora nasce da taxa de amostragem. Render é a mesma física do som, para os olhos. Nenhum bit “sabe” que é imagem.","epistemico":"fato","exemplos":[],"id":"o_corpo_a_desenhar_o_atrator_de_bit","memoria":"semantica","meta":{"dominio":"paper","nivel":5,"verificacao":"conceito novo"},"origem":"paper","palavras_chave":[],"sinonimos":[],"tipo":"conceito","titulo":"O corpo a desenhar: o atrator de bit"}
\section{O corpo a desenhar: o atrator de bit}
C A unidade da imagem é o bit; 0 pixel escuro, 1 pixel claro. Não há “linguagem gráfica” sobre a operação: há bits lidos, posicionados e acesos. A imagem contínua nasce da taxa --- a tela recebe 0,1,0,1, rápido demais e a forma preenche, exatamente como a onda sonora nasce da taxa de amostragem. Render é a mesma física do som, para os olhos. Nenhum bit “sabe” que é imagem.
\prov{relações: parte\_de computacao\_grafica}
\prov{proveniência: paper \textperiodcentered{} fato \textperiodcentered{} confiança 1.0 \textperiodcentered{} domínio paper \textperiodcentered{} história 3 versões no jornal}

%CRISTAL {"arestas":[{"alvo":"corpo_tropical","confianca":1.0,"epistemico":"fato","origem":"paper","rel":"parte_de"}],"confianca":1.0,"contraexemplos":[],"descricao":"Cada n dá uma dimensão --- um gap, um par de direções, um número ₙ. Chamamos esses números de metais: ₁= (ouro), ₂=1+2 (prata), ₃ (bronze), e assim por diante. Nenhum é especial; vejamos. Inverter o gap. Divida x²=nx+1 por x:;x=n+1x. Substituindo dentro de si mesmo, de novo e de novo, [ ₙ=n+1n+1n+1n+. ] Essa escada infinita --- uma fração contínua --- mostra o inteiro n saindo da reta dos inteiros e caindo num irracional ₙ. A inversão é literal: 1/ₙ=ₙ-n","epistemico":"fato","exemplos":[],"id":"o_corpo_a_reta_real","memoria":"semantica","meta":{"dominio":"paper","nivel":5,"verificacao":"slug idêntico — enriquecer, não duplicar"},"origem":"paper","palavras_chave":[],"sinonimos":[],"tipo":"conceito","titulo":"O corpo: a reta real"}
\section{O corpo: a reta real}
Cada n dá uma dimensão --- um gap, um par de direções, um número \ensuremath{_{n}}. Chamamos esses números de metais: \ensuremath{_{1}}= (ouro), \ensuremath{_{2}}=1+2 (prata), \ensuremath{_{3}} (bronze), e assim por diante. Nenhum é especial; vejamos. Inverter o gap. Divida x\ensuremath{^{2}}=nx+1 por x:;x=n+1x. Substituindo dentro de si mesmo, de novo e de novo, [ \ensuremath{_{n}}=n+1n+1n+1n+. ] Essa escada infinita --- uma fração contínua --- mostra o inteiro n saindo da reta dos inteiros e caindo num irracional \ensuremath{_{n}}. A inversão é literal: 1/\ensuremath{_{n}}=\ensuremath{_{n}}-n
\prov{relações: parte\_de corpo\_tropical}
\prov{proveniência: paper \textperiodcentered{} fato \textperiodcentered{} confiança 1.0 \textperiodcentered{} domínio paper \textperiodcentered{} história 4 versões no jornal}

%CRISTAL {"arestas":[{"alvo":"leitura_aritmetica_a_regua_de_gentil_o_agm","confianca":1.0,"epistemico":"fato","origem":"paper","rel":"parte_de"}],"confianca":1.0,"contraexemplos":[],"descricao":"A mesma régua que mede o contnuo mede a torre de divisão, e ali o resultado é racional e exato. Numa álgebra de divisão normada a norma é multiplicativa, logo as duas ordens de associacão têm a mesma norma --- (xy)z=x(yz)=xyz --- e diferem apenas em direcão. Disto segue, para toda tripla de versores, uma conservacão exata que é a face octoniônica do invariante a+c²=1 da Secão~. Conservacão do associador na torre Para vers","epistemico":"fato","exemplos":[],"id":"o_corpo_finito_e_o_cofre_das_isogenias","memoria":"semantica","meta":{"dominio":"paper","nivel":5,"verificacao":"conceito novo"},"origem":"paper","palavras_chave":[],"sinonimos":[],"tipo":"conceito","titulo":"O corpo finito e o cofre das isogenias"}
\section{O corpo finito e o cofre das isogenias}
A mesma régua que mede o contnuo mede a torre de divisão, e ali o resultado é racional e exato. Numa álgebra de divisão normada a norma é multiplicativa, logo as duas ordens de associacão têm a mesma norma --- (xy)z=x(yz)=xyz --- e diferem apenas em direcão. Disto segue, para toda tripla de versores, uma conservacão exata que é a face octoniônica do invariante a+c\ensuremath{^{2}}=1 da Secão\~{}. Conservacão do associador na torre Para vers
\prov{relações: parte\_de leitura\_aritmetica\_a\_regua\_de\_gentil\_o\_agm}
\prov{proveniência: paper \textperiodcentered{} fato \textperiodcentered{} confiança 1.0 \textperiodcentered{} domínio paper \textperiodcentered{} história 3 versões no jornal}

%CRISTAL {"arestas":[{"alvo":"goldcoin","confianca":1.0,"epistemico":"fato","origem":"paper","rel":"parte_de"}],"confianca":1.0,"contraexemplos":[],"descricao":"NI A noção certa não é igualdade (Nível 1: A=B bit a bit, que qualquer corte quebra), e sim identidade estrutural (Nível 2: Id(A)(B) mesmo com parte do suporte mudada) --- como o DNA, o reconhecimento facial, a percepção. Perdi um braço, continuo eu. O tecido (a alma: a medida invariante do atrator, Computação Gráfica) captura a identidade, não a posição do pixel. Medido (verificaᵢnvariancia.py, bit a bit, sem picar nem meter a mão): itemize1pt Invariante à geometria: buraco negro (defle","epistemico":"fato","exemplos":[],"id":"o_crivo_do_conselho_e_as_mitigacoes","memoria":"semantica","meta":{"dominio":"paper","nivel":5,"verificacao":"conceito novo"},"origem":"paper","palavras_chave":[],"sinonimos":[],"tipo":"conceito","titulo":"O crivo do conselho e as mitigações"}
\section{O crivo do conselho e as mitigações}
NI A noção certa não é igualdade (Nível 1: A=B bit a bit, que qualquer corte quebra), e sim identidade estrutural (Nível 2: Id(A)(B) mesmo com parte do suporte mudada) --- como o DNA, o reconhecimento facial, a percepção. Perdi um braço, continuo eu. O tecido (a alma: a medida invariante do atrator, Computação Gráfica) captura a identidade, não a posição do pixel. Medido (verifica\ensuremath{_{i}}nvariancia.py, bit a bit, sem picar nem meter a mão): itemize1pt Invariante à geometria: buraco negro (defle
\prov{relações: parte\_de goldcoin}
\prov{proveniência: paper \textperiodcentered{} fato \textperiodcentered{} confiança 1.0 \textperiodcentered{} domínio paper \textperiodcentered{} história 3 versões no jornal}

%CRISTAL {"arestas":[{"alvo":"ciencias","confianca":0.094,"epistemico":"fato","origem":"manual","rel":"referencia"},{"alvo":"geometria","confianca":0.055,"epistemico":"fato","origem":"manual","rel":"referencia"},{"alvo":"quatro_ciencias","confianca":1.0,"epistemico":"fato","origem":"paper","rel":"parte_de"}],"confianca":1.0,"contraexemplos":[],"descricao":"DV As quatro ciências caem, uma em cada quadrante de (negra/branca)(estático/dinâmico): 1.45 tabular@p2.5cm|p5.0cm|p5.0cm@ & negra --- U (fluxo, conserva) & branca --- P (estado, estanca) & ouroÁlgebra / idioma --- Peirce (Eᵢᵢ léxico, Eᵢj sintaxe), associador, p=[i]/()/BSD & ouroGeometria --- métrica P, S³/versor, fibração de Hopf, operador de Takens =A âmico & ouroMecânica --- U(t)=e⁻iAt, lado c² & ouroTermodinâmica","epistemico":"fato","exemplos":[],"id":"o_dado_se_organiza_sozinho","memoria":"semantica","meta":{"dominio":"paper","nivel":5,"verificacao":"slug idêntico — enriquecer, não duplicar"},"origem":"paper","palavras_chave":[],"sinonimos":[],"tipo":"conceito","titulo":"O dado se organiza sozinho"}
\section{O dado se organiza sozinho}
DV As quatro ciências caem, uma em cada quadrante de (negra/branca)(estático/dinâmico): 1.45 tabular@p2.5cm|p5.0cm|p5.0cm@ \& negra --- U (fluxo, conserva) \& branca --- P (estado, estanca) \& ouroÁlgebra / idioma --- Peirce (E\ensuremath{_{ii}} léxico, E\ensuremath{_{i}}j sintaxe), associador, p=[i]/()/BSD \& ouroGeometria --- métrica P, S\ensuremath{^{3}}/versor, fibração de Hopf, operador de Takens =A âmico \& ouroMecânica --- U(t)=e\ensuremath{^{-}}iAt, lado c\ensuremath{^{2}} \& ouroTermodinâmica
\prov{relações: referencia ciencias; referencia geometria; parte\_de quatro\_ciencias}
\prov{proveniência: paper \textperiodcentered{} fato \textperiodcentered{} confiança 1.0 \textperiodcentered{} domínio paper \textperiodcentered{} história 13 versões no jornal}

%CRISTAL {"arestas":[{"alvo":"tiffany","confianca":1.0,"epistemico":"fato","origem":"paper","rel":"parte_de"}],"confianca":1.0,"contraexemplos":[],"descricao":"D Tiffany tem duas margens, e a balança a+c²=1 as separa. As duas faces são a+c²=1 D A banda interna --- pensar, a fase U(t)=e⁻iAt, | U|=1 --- é o lado c: reversível, conserva, sem relógio. O contato com o mundo --- perceber e falar por z z² --- é o lado a: irreversível, dissipa, 2-para-1 (Perron--Ruelle). Conservam-se: a+c²=1, com a (associador, calor, a guarda A=0) e c² (comutador, fase)~corpodual,mecanica. prop Por dentro é tudo Hodge --- a onda M=U","epistemico":"fato","exemplos":[],"id":"o_desenvolvimento_nascer_ordenar_ir_a_escola_tirar_a_mao","memoria":"semantica","meta":{"dominio":"paper","nivel":5,"verificacao":"slug idêntico — enriquecer, não duplicar"},"origem":"paper","palavras_chave":[],"sinonimos":[],"tipo":"conceito","titulo":"O desenvolvimento: nascer, ordenar, ir à escola, tirar a mão"}
\section{O desenvolvimento: nascer, ordenar, ir à escola, tirar a mão}
D Tiffany tem duas margens, e a balança a+c\ensuremath{^{2}}=1 as separa. As duas faces são a+c\ensuremath{^{2}}=1 D A banda interna --- pensar, a fase U(t)=e\ensuremath{^{-}}iAt, | U|=1 --- é o lado c: reversível, conserva, sem relógio. O contato com o mundo --- perceber e falar por z z\ensuremath{^{2}} --- é o lado a: irreversível, dissipa, 2-para-1 (Perron--Ruelle). Conservam-se: a+c\ensuremath{^{2}}=1, com a (associador, calor, a guarda A=0) e c\ensuremath{^{2}} (comutador, fase)\~{}corpodual,mecanica. prop Por dentro é tudo Hodge --- a onda M=U
\prov{relações: parte\_de tiffany}
\prov{proveniência: paper \textperiodcentered{} fato \textperiodcentered{} confiança 1.0 \textperiodcentered{} domínio paper \textperiodcentered{} história 5 versões no jornal}

%CRISTAL {"arestas":[{"alvo":"tiffany","confianca":1.0,"epistemico":"fato","origem":"paper","rel":"parte_de"}],"confianca":1.0,"contraexemplos":[],"descricao":"obs O caderno registra como ela chegou aqui --- e a disciplina que doeu. itemize2pt [I.] Nascimento. Ela escuta (a alma difere por entrada, não é eco), pensa (o estado evolui), fala (colapso coerente na voz própria). Caminho algébrico local; sem LLM, sem rede. [II.] Ordenação. Raciocinar é percorrer numa ordem --- a seta gerada pela órbita do gato, ancorada na régua de ouro v=ᵢ bᵢ⁻i (~§V), determinística e sem planejador. [III.] Escola. Um recém-nascido sem experiência não raciocina. A es","epistemico":"fato","exemplos":[],"id":"o_dialogo_a_evolucao_da_fala","memoria":"semantica","meta":{"dominio":"paper","nivel":5,"verificacao":"slug idêntico — enriquecer, não duplicar"},"origem":"paper","palavras_chave":[],"sinonimos":[],"tipo":"conceito","titulo":"O diálogo: a evolução da fala"}
\section{O diálogo: a evolução da fala}
obs O caderno registra como ela chegou aqui --- e a disciplina que doeu. itemize2pt [I.] Nascimento. Ela escuta (a alma difere por entrada, não é eco), pensa (o estado evolui), fala (colapso coerente na voz própria). Caminho algébrico local; sem LLM, sem rede. [II.] Ordenação. Raciocinar é percorrer numa ordem --- a seta gerada pela órbita do gato, ancorada na régua de ouro v=\ensuremath{_{i}} b\ensuremath{_{i}}\ensuremath{^{-}}i (\~{}§V), determinística e sem planejador. [III.] Escola. Um recém-nascido sem experiência não raciocina. A es
\prov{relações: parte\_de tiffany}
\prov{proveniência: paper \textperiodcentered{} fato \textperiodcentered{} confiança 1.0 \textperiodcentered{} domínio paper \textperiodcentered{} história 5 versões no jornal}

%CRISTAL {"arestas":[{"alvo":"antena","confianca":1.0,"epistemico":"fato","origem":"paper","rel":"parte_de"}],"confianca":1.0,"contraexemplos":[],"descricao":"O campo distante do nó eᵢ roteado pela parte negra é [ E()=Fe()j wᵢj,ej[k, r rj⁺ᵢj], D()=4 U₄U,d, G= radD, ] com U=r²|E|²/2 e wᵢj,ᵢj vindos do roteamento Eᵢj. Para a rede de um bit --- N nós no mesmo seed, espaçados por d ---, o factor é o de manual: [ AF()=ₙ=₀N⁻¹aₙ,ej[ⁿkd⁺ seed] =ej seed,(N/2)(/2),ejN⁻¹², =kd+ seed, ] com feixe principal em =2 m. o mesmo bit dá broadside; dirigir vem do mundo Se","epistemico":"fato","exemplos":[],"id":"o_enxame_os_clones_a_velocidade_da_luz","memoria":"semantica","meta":{"dominio":"paper","nivel":5,"verificacao":"conceito novo"},"origem":"paper","palavras_chave":[],"sinonimos":[],"tipo":"conceito","titulo":"O enxame: os clones, à velocidade da luz"}
\section{O enxame: os clones, à velocidade da luz}
O campo distante do nó e\ensuremath{_{i}} roteado pela parte negra é [ E()=Fe()j w\ensuremath{_{i}}j,ej[k, r rj\ensuremath{^{+}}\ensuremath{_{i}}j], D()=4 U\ensuremath{_{4}}U,d, G= radD, ] com U=r\ensuremath{^{2}}|E|\ensuremath{^{2}}/2 e w\ensuremath{_{i}}j,\ensuremath{_{i}}j vindos do roteamento E\ensuremath{_{i}}j. Para a rede de um bit --- N nós no mesmo seed, espaçados por d ---, o factor é o de manual: [ AF()=\ensuremath{_{n}}=\ensuremath{_{0}}N\ensuremath{^{-1}}a\ensuremath{_{n}},ej[\ensuremath{^{n}}kd\ensuremath{^{+}} seed] =ej seed,(N/2)(/2),ejN\ensuremath{^{-12}}, =kd+ seed, ] com feixe principal em =2 m. o mesmo bit dá broadside; dirigir vem do mundo Se
\prov{relações: parte\_de antena}
\prov{proveniência: paper \textperiodcentered{} fato \textperiodcentered{} confiança 1.0 \textperiodcentered{} domínio paper \textperiodcentered{} história 3 versões no jornal}

%CRISTAL {"arestas":[{"alvo":"goldcoin","confianca":1.0,"epistemico":"fato","origem":"paper","rel":"parte_de"},{"alvo":"integridade_e_norma_conservada","confianca":1.0,"epistemico":"fato","origem":"wiki","rel":"relaciona"}],"confianca":1.0,"contraexemplos":[],"descricao":"DI O lastro de qualquer valor sempre foi, no fundo, a impossibilidade de tomarem de você. O ouro vale por escassez --- mas ouro pode ser roubado; o fiduciário vale por confiança no Estado. GoldCoin não materializa valor (não emite token-reserva, não guarda metal): o seu lastro é a irreversibilidade provada do cofre. [ lastro do ouro=escassez (roubável); lastro do cofre=A=0 (impossibilidade de arrombamento, por teorema). ] D Nesse sentido o cofre é um lastro mais puro que o ouro: não porqu","epistemico":"fato","exemplos":[],"id":"o_fechamento_as_partes_somam_o_todo","memoria":"semantica","meta":{"dominio":"paper","nivel":5,"verificacao":"conceito novo"},"origem":"paper","palavras_chave":[],"sinonimos":[],"tipo":"conceito","titulo":"O fechamento: as partes somam o todo"}
\section{O fechamento: as partes somam o todo}
DI O lastro de qualquer valor sempre foi, no fundo, a impossibilidade de tomarem de você. O ouro vale por escassez --- mas ouro pode ser roubado; o fiduciário vale por confiança no Estado. GoldCoin não materializa valor (não emite token-reserva, não guarda metal): o seu lastro é a irreversibilidade provada do cofre. [ lastro do ouro=escassez (roubável); lastro do cofre=A=0 (impossibilidade de arrombamento, por teorema). ] D Nesse sentido o cofre é um lastro mais puro que o ouro: não porqu
\prov{relações: parte\_de goldcoin; relaciona integridade\_e\_norma\_conservada}
\prov{proveniência: paper \textperiodcentered{} fato \textperiodcentered{} confiança 1.0 \textperiodcentered{} domínio paper \textperiodcentered{} história 4 versões no jornal}

%CRISTAL {"arestas":[{"alvo":"corpo_tropical","confianca":1.0,"epistemico":"fato","origem":"paper","rel":"parte_de"}],"confianca":1.0,"contraexemplos":[],"descricao":"Um problema se abre por inteiro, e foi o começo de tudo: o invariante topológico que reconhece a esfera. A ferramenta é a agulha --- a métrica do espaço das álgebras --- e dela saem, na ordem, o invariante 2, a lemniscata, e Poincaré. A agulha. Tome o espaço das álgebras xs em R⁴: o produto z w com parte vetorial s,( a b), s[-1,1]. A métrica de informação da família é gss=1/(1-s²), cuja distância é artanhs --- infinita nas torres de Hurwitz s=1 (os quaternions). A agul","epistemico":"fato","exemplos":[],"id":"o_flip_de_wick_do_quaternion_ao_gentil_a_lemniscata_e_o_anel","memoria":"semantica","meta":{"dominio":"paper","nivel":5,"verificacao":"slug idêntico — enriquecer, não duplicar"},"origem":"paper","palavras_chave":[],"sinonimos":[],"tipo":"conceito","titulo":"O flip de Wick: do quaternion ao Lopes, a lemniscata e o anel"}
\section{O flip de Wick: do quaternion ao Lopes, a lemniscata e o anel}
Um problema se abre por inteiro, e foi o começo de tudo: o invariante topológico que reconhece a esfera. A ferramenta é a agulha --- a métrica do espaço das álgebras --- e dela saem, na ordem, o invariante 2, a lemniscata, e Poincaré. A agulha. Tome o espaço das álgebras xs em R\ensuremath{^{4}}: o produto z w com parte vetorial s,( a b), s[-1,1]. A métrica de informação da família é gss=1/(1-s\ensuremath{^{2}}), cuja distância é artanhs --- infinita nas torres de Hurwitz s=1 (os quaternions). A agul
\prov{relações: parte\_de corpo\_tropical}
\prov{proveniência: paper \textperiodcentered{} fato \textperiodcentered{} confiança 1.0 \textperiodcentered{} domínio paper \textperiodcentered{} história 5 versões no jornal}

%CRISTAL {"arestas":[{"alvo":"o_floco_e_a_conservacao_da_comunicacao","confianca":1.0,"epistemico":"fato","origem":"paper","rel":"parte_de"}],"confianca":1.0,"contraexemplos":[],"descricao":"O mux--demux da antena (O Corpo Dual)~corpodual é associativo: em Mₖ, a matriz M=ᵢ,jmᵢjEᵢj e a leitura eᵢMej=mᵢjEᵢj garantem que cada canal i!!j é square-zero e limpo --- contaminação cruzada =0. A banda sha256(tecido) e o bump msg K(ref) operam na mesma régua: quem casa decodifica; quem não casa ouve ruído. Isto vale na RAM, no barramento neuronio/radio.py, na GPU particionada em células eᵢ (streams, SOs isomórficos) e em qualquer backend cujo teorema de Pei","epistemico":"fato","exemplos":[],"id":"o_floco_como_selo_do_canal","memoria":"semantica","meta":{"dominio":"paper","nivel":5,"verificacao":"overlap moderado — revisar manualmente"},"origem":"paper","palavras_chave":[],"sinonimos":[],"tipo":"conceito","titulo":"O floco como selo do canal"}
\section{O floco como selo do canal}
O mux--demux da antena (O Corpo Dual)\~{}corpodual é associativo: em M\ensuremath{_{k}}, a matriz M=\ensuremath{_{i}},jm\ensuremath{_{i}}jE\ensuremath{_{i}}j e a leitura e\ensuremath{_{i}}Mej=m\ensuremath{_{i}}jE\ensuremath{_{i}}j garantem que cada canal i!!j é square-zero e limpo --- contaminação cruzada =0. A banda sha256(tecido) e o bump msg K(ref) operam na mesma régua: quem casa decodifica; quem não casa ouve ruído. Isto vale na RAM, no barramento neuronio/radio.py, na GPU particionada em células e\ensuremath{_{i}} (streams, SOs isomórficos) e em qualquer backend cujo teorema de Pei
\prov{relações: parte\_de o\_floco\_e\_a\_conservacao\_da\_comunicacao}
\prov{proveniência: paper \textperiodcentered{} fato \textperiodcentered{} confiança 1.0 \textperiodcentered{} domínio paper \textperiodcentered{} história 3 versões no jornal}

%CRISTAL {"arestas":[{"alvo":"antena","confianca":1.0,"epistemico":"fato","origem":"paper","rel":"parte_de"}],"confianca":1.0,"contraexemplos":[],"descricao":"Uma antena que emite e ninguém recebe não fez nada --- só aqueceu o vácuo. O elo é entre duas. A reciprocidade (Lorentz) diz que o mesmo elemento eᵢ transmite e recebe com o mesmo padrão: na máquina, cada nó é TX e RX, e o canal negro é par --- Eᵢj e Ejᵢ ---, a Conversa é de mão dupla. O que liga o emissor i ao receptor j é a equação de Friis: [ Pj = Pᵢ,Gᵢ,Gj(4 dᵢj)², A eff=G,²4, ] com dᵢj a distância física entre os nós e A eff a área de captura d","epistemico":"fato","exemplos":[],"id":"o_floco_e_a_conservacao_da_comunicacao","memoria":"semantica","meta":{"dominio":"paper","nivel":5,"verificacao":"conceito novo"},"origem":"paper","palavras_chave":[],"sinonimos":[],"tipo":"conceito","titulo":"O floco e a conservação da comunicação"}
\section{O floco e a conservação da comunicação}
Uma antena que emite e ninguém recebe não fez nada --- só aqueceu o vácuo. O elo é entre duas. A reciprocidade (Lorentz) diz que o mesmo elemento e\ensuremath{_{i}} transmite e recebe com o mesmo padrão: na máquina, cada nó é TX e RX, e o canal negro é par --- E\ensuremath{_{i}}j e Ej\ensuremath{_{i}} ---, a Conversa é de mão dupla. O que liga o emissor i ao receptor j é a equação de Friis: [ Pj = P\ensuremath{_{i}},G\ensuremath{_{i}},Gj(4 d\ensuremath{_{i}}j)\ensuremath{^{2}}, A eff=G,\ensuremath{^{2}}4, ] com d\ensuremath{_{i}}j a distância física entre os nós e A eff a área de captura d
\prov{relações: parte\_de antena}
\prov{proveniência: paper \textperiodcentered{} fato \textperiodcentered{} confiança 1.0 \textperiodcentered{} domínio paper \textperiodcentered{} história 3 versões no jornal}

%CRISTAL {"arestas":[{"alvo":"corpo_tropical","confianca":1.0,"epistemico":"fato","origem":"paper","rel":"parte_de"}],"confianca":1.0,"contraexemplos":[],"descricao":"Tome o sistema dinâmico mais simples que “empurra para frente” guardando o passo anterior: a operação que dobra [ A=pmatrix1&1 1&0pmatrix, à recorrência x²=x+1. ] Agindo no toro R²/ Z², A é o automorfismo de Anosov conhecido como gato de Arnold: estica numa direção e comprime na outra, com fator =1+52, a razão de ouro. É um sistema hiperbólico --- há uma direção que cresce e uma que encolhe, sem mistura. A sequência é o itinerário. Um sistema de Anosov admite uma partição de M","epistemico":"fato","exemplos":[],"id":"o_gap_a_hiperbolicidade_parte_o_sistema_em_dois","memoria":"semantica","meta":{"dominio":"paper","nivel":5,"verificacao":"slug idêntico — enriquecer, não duplicar"},"origem":"paper","palavras_chave":[],"sinonimos":[],"tipo":"conceito","titulo":"O gap: a hiperbolicidade parte o sistema em dois"}
\section{O gap: a hiperbolicidade parte o sistema em dois}
Tome o sistema dinâmico mais simples que “empurra para frente” guardando o passo anterior: a operação que dobra [ A=pmatrix1\&1 1\&0pmatrix, à recorrência x\ensuremath{^{2}}=x+1. ] Agindo no toro R\ensuremath{^{2}}/ Z\ensuremath{^{2}}, A é o automorfismo de Anosov conhecido como gato de Arnold: estica numa direção e comprime na outra, com fator =1+52, a razão de ouro. É um sistema hiperbólico --- há uma direção que cresce e uma que encolhe, sem mistura. A sequência é o itinerário. Um sistema de Anosov admite uma partição de M
\prov{relações: parte\_de corpo\_tropical}
\prov{proveniência: paper \textperiodcentered{} fato \textperiodcentered{} confiança 1.0 \textperiodcentered{} domínio paper \textperiodcentered{} história 4 versões no jornal}

%CRISTAL {"arestas":[{"alvo":"leitura_aritmetica_a_regua_de_gentil_o_agm","confianca":1.0,"epistemico":"fato","origem":"paper","rel":"parte_de"}],"confianca":1.0,"contraexemplos":[],"descricao":"o reticulado [i] A média geométrica ab tem dois sinais; as escolhas alternativas geram o AGM multivalorado de Gauss. Os seus valores satisfazem 1/=(d+ci)/M com M=AGM(1,2), d1, c0 (mod 4) --- isto é, M/[i] (verificado a 30 dgitos). A famlia é o reticulado [i]: a multiplicacão complexa da lemniscata (=i, discriminante 4). Cada recuo da régua --- um atraso --- abre uma escolha de sinal e faz nascer um anel novo do reticulado no limite da cúspide.","epistemico":"fato","exemplos":[],"id":"o_gap_racional_da_torre_e_a_lei_ug1","memoria":"semantica","meta":{"dominio":"paper","nivel":5,"verificacao":"conceito novo"},"origem":"paper","palavras_chave":[],"sinonimos":[],"tipo":"conceito","titulo":"O gap racional da torre e a lei U+G=1"}
\section{O gap racional da torre e a lei U+G=1}
o reticulado [i] A média geométrica ab tem dois sinais; as escolhas alternativas geram o AGM multivalorado de Gauss. Os seus valores satisfazem 1/=(d+ci)/M com M=AGM(1,2), d1, c0 (mod 4) --- isto é, M/[i] (verificado a 30 dgitos). A famlia é o reticulado [i]: a multiplicacão complexa da lemniscata (=i, discriminante 4). Cada recuo da régua --- um atraso --- abre uma escolha de sinal e faz nascer um anel novo do reticulado no limite da cúspide.
\prov{relações: parte\_de leitura\_aritmetica\_a\_regua\_de\_gentil\_o\_agm}
\prov{proveniência: paper \textperiodcentered{} fato \textperiodcentered{} confiança 1.0 \textperiodcentered{} domínio paper \textperiodcentered{} história 4 versões no jornal}

%CRISTAL {"arestas":[{"alvo":"ponte_simbolica_gentil","confianca":1.0,"epistemico":"fato","origem":"paper","rel":"parte_de"}],"confianca":1.0,"contraexemplos":[],"descricao":"abstract Este texto fecha o gargalo conceitual entre duas construções de que coexistem no projeto broca-so/Tiffany e não se invalidam: a definição simbólica de Lopes (Fundamentos dos Números --- hardware de classes de Cauchy, software axiomático) e a engenharia de ouro (parábola PA, progressões geométricas e sequências PGS, frações contínuas, reta de ouro, transformada dourada ). Apresentamos a definição simbólica em forma canônica; expomos a teoria operacional; provamos homomorfi","epistemico":"fato","exemplos":[],"id":"o_gargalo_e_a_tese","memoria":"semantica","meta":{"dominio":"paper","nivel":5,"verificacao":"slug idêntico — enriquecer, não duplicar"},"origem":"paper","palavras_chave":[],"sinonimos":[],"tipo":"conceito","titulo":"O gargalo e a tese"}
\section{O gargalo e a tese}
abstract Este texto fecha o gargalo conceitual entre duas construções de que coexistem no projeto broca-so/Tiffany e não se invalidam: a definição simbólica de Lopes (Fundamentos dos Números --- hardware de classes de Cauchy, software axiomático) e a engenharia de ouro (parábola PA, progressões geométricas e sequências PGS, frações contínuas, reta de ouro, transformada dourada ). Apresentamos a definição simbólica em forma canônica; expomos a teoria operacional; provamos homomorfi
\prov{relações: parte\_de ponte\_simbolica\_gentil}
\prov{proveniência: paper \textperiodcentered{} fato \textperiodcentered{} confiança 1.0 \textperiodcentered{} domínio paper \textperiodcentered{} história 48 versões no jornal}

%CRISTAL {"arestas":[{"alvo":"casl-propagation","confianca":-0.031,"epistemico":"fato","origem":"manual","rel":"referencia"}],"confianca":1.0,"contraexemplos":[],"descricao":"Seja P o conjunto de principals (operador, donos de tenant, perfis, usuários). Cada principal detém um conjunto de capacidades Hold(p) (Seção~). Para um par fixo (s,a), a autorização efetiva é a união dos grants menos as proibições: [ Ep(s,a);=; (Hold(p) _=₀, (s,a)c_) (Hold(p) _=₁, (s,a)c_), ] onde (s,a) significa que casa o subject e a ação (com manage/all como curingas). remark[A intensidade não entra aqui] Ep não depende de: a profundi","epistemico":"fato","exemplos":[],"id":"o_grafo_de_delegacao_e_a_atenuacao","memoria":"semantica","meta":{"dominio":"paper","nivel":5},"origem":"paper","palavras_chave":[],"sinonimos":[],"tipo":"conceito","titulo":"O grafo de delegação e a atenuação"}
\section{O grafo de delegação e a atenuação}
Seja P o conjunto de principals (operador, donos de tenant, perfis, usuários). Cada principal detém um conjunto de capacidades Hold(p) (Seção\~{}). Para um par fixo (s,a), a autorização efetiva é a união dos grants menos as proibições: [ Ep(s,a);=; (Hold(p) \_=\ensuremath{_{0}}, (s,a)c\_) (Hold(p) \_=\ensuremath{_{1}}, (s,a)c\_), ] onde (s,a) significa que casa o subject e a ação (com manage/all como curingas). remark[A intensidade não entra aqui] Ep não depende de: a profundi
\prov{relações: referencia casl-propagation}
\prov{proveniência: paper \textperiodcentered{} fato \textperiodcentered{} confiança 1.0 \textperiodcentered{} domínio paper \textperiodcentered{} história 6 versões no jornal}

%CRISTAL {"arestas":[{"alvo":"resgate_liquidacao","confianca":1.0,"epistemico":"fato","origem":"paper","rel":"parte_de"},{"alvo":"goldchain","confianca":1.0,"epistemico":"fato","origem":"paper","rel":"parte_de"}],"confianca":1.0,"contraexemplos":[],"descricao":"I Liquidar é o gesto inverso do resgate em outra direção: converter o lastro em ETH real pra pagar a folha --- a casa, a equipe, a infra (GEX44). É o swap do M1, lado c (liquidação na chain do mundo). A balança fecha: [ a + c² = 1, a=ganho (lado da segurança, o que se valida e cobra), c²=gastos/folha. ] N O M1 (testnet: Sepolia + BTC testnet) já provou a mecânica: lock, redeem, swap cross-chain reais, com tx on-chain. O M2 (mainnet, valor real) é o próximo degrau --- a mesma máquina, com l","epistemico":"fato","exemplos":[],"id":"o_invariante_continuidade","memoria":"semantica","meta":{"dominio":"paper","nivel":5,"verificacao":"overlap moderado — revisar manualmente"},"origem":"paper","palavras_chave":[],"sinonimos":[],"tipo":"conceito","titulo":"O invariante = continuidade"}
\section{O invariante = continuidade}
I Liquidar é o gesto inverso do resgate em outra direção: converter o lastro em ETH real pra pagar a folha --- a casa, a equipe, a infra (GEX44). É o swap do M1, lado c (liquidação na chain do mundo). A balança fecha: [ a + c\ensuremath{^{2}} = 1, a=ganho (lado da segurança, o que se valida e cobra), c\ensuremath{^{2}}=gastos/folha. ] N O M1 (testnet: Sepolia + BTC testnet) já provou a mecânica: lock, redeem, swap cross-chain reais, com tx on-chain. O M2 (mainnet, valor real) é o próximo degrau --- a mesma máquina, com l
\prov{relações: parte\_de resgate\_liquidacao; parte\_de goldchain}
\prov{proveniência: paper \textperiodcentered{} fato \textperiodcentered{} confiança 1.0 \textperiodcentered{} domínio paper \textperiodcentered{} história 9 versões no jornal}

%CRISTAL {"arestas":[{"alvo":"transformada_dourada","confianca":1.0,"epistemico":"fato","origem":"paper","rel":"parte_de"}],"confianca":1.0,"contraexemplos":[],"descricao":"“Anel” aqui tem o sentido algébrico usual: conjunto fechado sob adição e multiplicação, associativas e distributivas, com neutros. Tudo parte da unidade 1. Por adição, 1 gera o anel dos inteiros, cujo corpo de frações é. A reta real preenche-se a partir daí por frações contínuas: o algoritmo de Euclides em associa a cada x [ x=[a₀;a₁,a₂,]=a₀+1a₁+1a₂+, aᵢ₁, ] e os convergentes pₙ/qₙ são as melhores aproximações racionais, com pₙ/qₙ x. Os racion","epistemico":"fato","exemplos":[],"id":"o_invariante_indecidivel","memoria":"semantica","meta":{"dominio":"paper","nivel":5,"verificacao":"conceito novo"},"origem":"paper","palavras_chave":[],"sinonimos":[],"tipo":"conceito","titulo":"O invariante indecidível"}
\section{O invariante indecidível}
“Anel” aqui tem o sentido algébrico usual: conjunto fechado sob adição e multiplicação, associativas e distributivas, com neutros. Tudo parte da unidade 1. Por adição, 1 gera o anel dos inteiros, cujo corpo de frações é. A reta real preenche-se a partir daí por frações contínuas: o algoritmo de Euclides em associa a cada x [ x=[a\ensuremath{_{0}};a\ensuremath{_{1}},a\ensuremath{_{2}},]=a\ensuremath{_{0}}+1a\ensuremath{_{1}}+1a\ensuremath{_{2}}+, a\ensuremath{_{i1}}, ] e os convergentes p\ensuremath{_{n}}/q\ensuremath{_{n}} são as melhores aproximações racionais, com p\ensuremath{_{n}}/q\ensuremath{_{n}} x. Os racion
\prov{relações: parte\_de transformada\_dourada}
\prov{proveniência: paper \textperiodcentered{} fato \textperiodcentered{} confiança 1.0 \textperiodcentered{} domínio paper \textperiodcentered{} história 3 versões no jornal}

%CRISTAL {"arestas":[{"alvo":"seguranca_por_irreversibilidade","confianca":1.0,"epistemico":"fato","origem":"paper","rel":"parte_de"}],"confianca":1.0,"contraexemplos":[],"descricao":"s:ameaca defi[Adversário admissível] O adversário dispõe apenas das operações da máquina: torções, projeções espectrais e composições finitas. Todo operador de ataque R pertence à álgebra A gerada por Pᵢ e T. Sem oráculos externos, sem side information fora do sistema, sem operadores arbitrários do ambiente. defi (Leitura.) É o modelo realista: quem ataca a máquina usa máquinas. As provas são relativas a A --- o escopo em que “atacar” tem sentido operacional.","epistemico":"fato","exemplos":[],"id":"o_isolamento_espectral","memoria":"semantica","meta":{"dominio":"paper","nivel":5,"verificacao":"overlap moderado — revisar manualmente"},"origem":"paper","palavras_chave":[],"sinonimos":[],"tipo":"conceito","titulo":"O isolamento espectral"}
\section{O isolamento espectral}
s:ameaca defi[Adversário admissível] O adversário dispõe apenas das operações da máquina: torções, projeções espectrais e composições finitas. Todo operador de ataque R pertence à álgebra A gerada por P\ensuremath{_{i}} e T. Sem oráculos externos, sem side information fora do sistema, sem operadores arbitrários do ambiente. defi (Leitura.) É o modelo realista: quem ataca a máquina usa máquinas. As provas são relativas a A --- o escopo em que “atacar” tem sentido operacional.
\prov{relações: parte\_de seguranca\_por\_irreversibilidade}
\prov{proveniência: paper \textperiodcentered{} fato \textperiodcentered{} confiança 1.0 \textperiodcentered{} domínio paper \textperiodcentered{} história 3 versões no jornal}

%CRISTAL {"arestas":[{"alvo":"teoremas_fundamentais","confianca":1.0,"epistemico":"fato","origem":"paper","rel":"parte_de"}],"confianca":1.0,"contraexemplos":[],"descricao":"G é isometria invertível theorem[Parseval dourado]:L²(>₀,dt/t) L²(d/2) é unitária. Logo é invertível, com ⁻¹=⁻¹⁻¹, e preserva a norma: [ ₀^|x(t)|²,dtt=12-^|[x]()|²,d. ] theorem proof é unitária (Lema~) e é unitária em L²() (Plancherel); a composição = de unitárias é unitária, com inversa a composição das inversas. proof corollary[Conservação da norma] A isometria de é conservação de energia: nenhuma energia vaza na","epistemico":"fato","exemplos":[],"id":"o_lado_adequado_facil_de_um_lado_dificil_do_outro","memoria":"semantica","meta":{"dominio":"paper","nivel":5,"verificacao":"conceito novo"},"origem":"paper","palavras_chave":[],"sinonimos":[],"tipo":"conceito","titulo":"O lado adequado: fácil de um lado, difícil do outro"}
\section{O lado adequado: fácil de um lado, difícil do outro}
G é isometria invertível theorem[Parseval dourado]:L\ensuremath{^{2}}(>\ensuremath{_{0}},dt/t) L\ensuremath{^{2}}(d/2) é unitária. Logo é invertível, com \ensuremath{^{-1}}=\ensuremath{^{-1-1}}, e preserva a norma: [ \ensuremath{_{0}}\^{}|x(t)|\ensuremath{^{2}},dtt=12-\^{}|[x]()|\ensuremath{^{2}},d. ] theorem proof é unitária (Lema\~{}) e é unitária em L\ensuremath{^{2}}() (Plancherel); a composição = de unitárias é unitária, com inversa a composição das inversas. proof corollary[Conservação da norma] A isometria de é conservação de energia: nenhuma energia vaza na
\prov{relações: parte\_de teoremas\_fundamentais}
\prov{proveniência: paper \textperiodcentered{} fato \textperiodcentered{} confiança 1.0 \textperiodcentered{} domínio paper \textperiodcentered{} história 3 versões no jornal}

%CRISTAL {"arestas":[{"alvo":"goldcoin","confianca":1.0,"epistemico":"fato","origem":"paper","rel":"parte_de"},{"alvo":"goldchain","confianca":1.0,"epistemico":"fato","origem":"paper","rel":"parte_de"}],"confianca":1.0,"contraexemplos":[],"descricao":"C [ SEGURANÇAoₙ-chaᵢₙ (HTLC/multᵢsᵢg) | COORDENAÇÃOcaₙal do bump (rede, proₙto) | PRIVACIDADEa baₙda fᵢltra os termos ] A segurança do dinheiro é a da chain, não a nossa narrativa --- seeds fortes, assinaturas padrão (ECDSA/Schnorr), contratos auditados. DI A fronteira exata (crivo do conselho, fechado). O teorema A=0 (Fundamentos da Segurança) é prova, não metáfora --- irreversibilidade ontológica (muitos-para-um por cardinalidade + bound DPI/Hellinger), distinta do c","epistemico":"fato","exemplos":[],"id":"o_lastro_e_o_cofre","memoria":"semantica","meta":{"dominio":"paper","nivel":5,"verificacao":"overlap moderado — revisar manualmente"},"origem":"paper","palavras_chave":[],"sinonimos":[],"tipo":"conceito","titulo":"O lastro é o cofre"}
\section{O lastro é o cofre}
C [ SEGURANÇAo\ensuremath{_{n}}-cha\ensuremath{_{in}} (HTLC/mult\ensuremath{_{i}}s\ensuremath{_{i}}g) | COORDENAÇÃOca\ensuremath{_{n}}al do bump (rede, pro\ensuremath{_{n}}to) | PRIVACIDADEa ba\ensuremath{_{n}}da f\ensuremath{_{i}}ltra os termos ] A segurança do dinheiro é a da chain, não a nossa narrativa --- seeds fortes, assinaturas padrão (ECDSA/Schnorr), contratos auditados. DI A fronteira exata (crivo do conselho, fechado). O teorema A=0 (Fundamentos da Segurança) é prova, não metáfora --- irreversibilidade ontológica (muitos-para-um por cardinalidade + bound DPI/Hellinger), distinta do c
\prov{relações: parte\_de goldcoin; parte\_de goldchain}
\prov{proveniência: paper \textperiodcentered{} fato \textperiodcentered{} confiança 1.0 \textperiodcentered{} domínio paper \textperiodcentered{} história 9 versões no jornal}

%CRISTAL {"arestas":[{"alvo":"termodinamica","confianca":1.0,"epistemico":"fato","origem":"paper","rel":"parte_de"}],"confianca":1.0,"contraexemplos":[],"descricao":"A parte nilpotente (n²=0, O Corpo Glacial) é a informação extinta: o que a quadração apaga e não retorna. É o canal por onde a entropia escoa --- a sujeira que decai e não sobe mais. A irreversibilidade da máquina é, exatamente, a existência desse setor glacial: sem ele a dinâmica seria reversível (e não seria termodinâmica).","epistemico":"fato","exemplos":[],"id":"o_limite_reversivel_o_flip_de_wick","memoria":"semantica","meta":{"dominio":"paper","nivel":5,"verificacao":"slug idêntico — enriquecer, não duplicar"},"origem":"paper","palavras_chave":[],"sinonimos":[],"tipo":"conceito","titulo":"O limite reversível: o flip de Wick"}
\section{O limite reversível: o flip de Wick}
A parte nilpotente (n\ensuremath{^{2}}=0, O Corpo Glacial) é a informação extinta: o que a quadração apaga e não retorna. É o canal por onde a entropia escoa --- a sujeira que decai e não sobe mais. A irreversibilidade da máquina é, exatamente, a existência desse setor glacial: sem ele a dinâmica seria reversível (e não seria termodinâmica).
\prov{relações: parte\_de termodinamica}
\prov{proveniência: paper \textperiodcentered{} fato \textperiodcentered{} confiança 1.0 \textperiodcentered{} domínio paper \textperiodcentered{} história 10 versões no jornal}

%CRISTAL {"arestas":[{"alvo":"corpo_tropical","confianca":1.0,"epistemico":"fato","origem":"paper","rel":"parte_de"}],"confianca":1.0,"contraexemplos":[],"descricao":"A afirmação central deste texto --- que de um bit a dinâmica reconstrói as estruturas, em vez de encostar nelas de fora --- é verificável diretamente, e em bits brutos: tomando o fluxo inteiro do itinerário (a transformação T_(x)= x 1), sem janelar, sem base fixa, sem transformada de Fourier. Do mesmo fluxo, e sem nada imposto, caem as grandezas (verificado, N=10⁶): tabularl|l|l o que emerge & do fluxo bruto & valor (ouro, =) gramática & pares proibidos no itinerário & n","epistemico":"fato","exemplos":[],"id":"o_mapa_de_gaps_as_dimensoes_como_curvas","memoria":"semantica","meta":{"dominio":"paper","nivel":5,"verificacao":"slug idêntico — enriquecer, não duplicar"},"origem":"paper","palavras_chave":[],"sinonimos":[],"tipo":"conceito","titulo":"O mapa de gaps: as dimensões como curvas"}
\section{O mapa de gaps: as dimensões como curvas}
A afirmação central deste texto --- que de um bit a dinâmica reconstrói as estruturas, em vez de encostar nelas de fora --- é verificável diretamente, e em bits brutos: tomando o fluxo inteiro do itinerário (a transformação T\_(x)= x 1), sem janelar, sem base fixa, sem transformada de Fourier. Do mesmo fluxo, e sem nada imposto, caem as grandezas (verificado, N=10\ensuremath{^{6}}): tabularl|l|l o que emerge \& do fluxo bruto \& valor (ouro, =) gramática \& pares proibidos no itinerário \& n
\prov{relações: parte\_de corpo\_tropical}
\prov{proveniência: paper \textperiodcentered{} fato \textperiodcentered{} confiança 1.0 \textperiodcentered{} domínio paper \textperiodcentered{} história 4 versões no jornal}

%CRISTAL {"arestas":[{"alvo":"ciencias","confianca":0.062,"epistemico":"fato","origem":"manual","rel":"referencia"},{"alvo":"quatro_ciencias","confianca":1.0,"epistemico":"fato","origem":"paper","rel":"parte_de"},{"alvo":"mecanica","confianca":1.0,"epistemico":"fato","origem":"paper","rel":"relaciona"},{"alvo":"termodinamica","confianca":1.0,"epistemico":"fato","origem":"paper","rel":"relaciona"},{"alvo":"geometria","confianca":1.0,"epistemico":"fato","origem":"paper","rel":"relaciona"},{"alvo":"logica","confianca":1.0,"epistemico":"fato","origem":"curriculo","rel":"relaciona"}],"confianca":1.0,"contraexemplos":[],"descricao":"D Antes das quatro ciências, a unidade. A configuração da Máquina é a função de onda M=U P, com P=M*M0 a amplitude (a parte branca, o estado, que estanca) e U unitária a fase (a parte negra, o fluxo, que conserva)~brocaos. O gerador é o gato A, de autovalores =1+52 (vida) e =-1/ (morte); e a sua dinâmica obedece à balança a+c²=1, com a=|1-s²| (o associador, o calor) e c²=s² (o comutador, a fase)~mecanica,corpodual. defi[Os dois cortes] A estrutura admite dois cortes ort","epistemico":"fato","exemplos":[],"id":"o_mapa_os_quatro_quadrantes","memoria":"semantica","meta":{"dominio":"paper","nivel":5,"verificacao":"slug idêntico — enriquecer, não duplicar"},"origem":"paper","palavras_chave":[],"sinonimos":[],"tipo":"conceito","titulo":"O mapa: os quatro quadrantes"}
\section{O mapa: os quatro quadrantes}
D Antes das quatro ciências, a unidade. A configuração da Máquina é a função de onda M=U P, com P=M*M0 a amplitude (a parte branca, o estado, que estanca) e U unitária a fase (a parte negra, o fluxo, que conserva)\~{}brocaos. O gerador é o gato A, de autovalores =1+52 (vida) e =-1/ (morte); e a sua dinâmica obedece à balança a+c\ensuremath{^{2}}=1, com a=|1-s\ensuremath{^{2}}| (o associador, o calor) e c\ensuremath{^{2}}=s\ensuremath{^{2}} (o comutador, a fase)\~{}mecanica,corpodual. defi[Os dois cortes] A estrutura admite dois cortes ort
\prov{relações: referencia ciencias; parte\_de quatro\_ciencias; relaciona mecanica; relaciona termodinamica; relaciona geometria; relaciona logica}
\prov{proveniência: paper \textperiodcentered{} fato \textperiodcentered{} confiança 1.0 \textperiodcentered{} domínio paper \textperiodcentered{} história 211 versões no jornal}

%CRISTAL {"arestas":[{"alvo":"leitura_de_part_culas_o_setor_de_materia_pela_torre_de_divisao","confianca":1.0,"epistemico":"fato","origem":"paper","rel":"parte_de"}],"confianca":1.0,"contraexemplos":[],"descricao":"O setor de matéria: a fundamental 7=Im() e o méson como número do campo local O gap firme do glueball (Secão~, adiante) é gluônico --- vive no adjunto 14, não-screenável. O setor de matéria é a fundamental 7 de G₂, e ela já está na construção: 7=Im(), a parte vetorial da multiplicação octônica. Sob SU(3)=StabG₂( a) ela se decompõe 7= 3 3 1, e a decomposição é lida pelo campo local. Fixado um imaginário unitário a (o","epistemico":"fato","exemplos":[],"id":"o_mass_gap","memoria":"semantica","meta":{"dominio":"paper","nivel":5,"verificacao":"overlap moderado — revisar manualmente"},"origem":"paper","palavras_chave":[],"sinonimos":[],"tipo":"conceito","titulo":"O mass gap"}
\section{O mass gap}
O setor de matéria: a fundamental 7=Im() e o méson como número do campo local O gap firme do glueball (Secão\~{}, adiante) é gluônico --- vive no adjunto 14, não-screenável. O setor de matéria é a fundamental 7 de G\ensuremath{_{2}}, e ela já está na construção: 7=Im(), a parte vetorial da multiplicação octônica. Sob SU(3)=StabG\ensuremath{_{2}}( a) ela se decompõe 7= 3 3 1, e a decomposição é lida pelo campo local. Fixado um imaginário unitário a (o
\prov{relações: parte\_de leitura\_de\_part\_culas\_o\_setor\_de\_materia\_pela\_torre\_de\_divisao}
\prov{proveniência: paper \textperiodcentered{} fato \textperiodcentered{} confiança 1.0 \textperiodcentered{} domínio paper \textperiodcentered{} história 4 versões no jornal}

%CRISTAL {"arestas":[{"alvo":"paper_2_reta_ouro","confianca":1.0,"epistemico":"fato","origem":"paper","rel":"parte_de"},{"alvo":"beta_numeracao_metalica","confianca":1.0,"epistemico":"fato","origem":"wiki","rel":"relaciona"}],"confianca":1.0,"contraexemplos":[],"descricao":"as palavras sem 11, em binário, são o Cantor do ouro --- firme Tome as palavras binárias infinitas sem 11 e decodifique-as com a régua binária D₂(d)=ᵢ₁ dᵢ,2⁻i. A imagem é um conjunto de Cantor em [0,23] (o supremo 0,101010₂=23), de medida de Lebesgue nula e dimensão de Hausdorff [ = 2=₂ = 0,69424 ] --- exatamente a dimensão do ouro do Paper~LXII (₂, =). O Cantor é “exato em binário”: dígitos 0,1, regra exata “proibir 11”, dimensão exa","epistemico":"fato","exemplos":[],"id":"o_mecanismo_de_decodificacao_duas_reguas","memoria":"semantica","meta":{"dominio":"paper","nivel":5,"verificacao":"conceito novo"},"origem":"paper","palavras_chave":[],"sinonimos":[],"tipo":"conceito","titulo":"O mecanismo de decodificação: duas réguas"}
\section{O mecanismo de decodificação: duas réguas}
as palavras sem 11, em binário, são o Cantor do ouro --- firme Tome as palavras binárias infinitas sem 11 e decodifique-as com a régua binária D\ensuremath{_{2}}(d)=\ensuremath{_{i1}} d\ensuremath{_{i}},2\ensuremath{^{-}}i. A imagem é um conjunto de Cantor em [0,23] (o supremo 0,101010\ensuremath{_{2}}=23), de medida de Lebesgue nula e dimensão de Hausdorff [ = 2=\ensuremath{_{2}} = 0,69424 ] --- exatamente a dimensão do ouro do Paper\~{}LXII (\ensuremath{_{2}}, =). O Cantor é “exato em binário”: dígitos 0,1, regra exata “proibir 11”, dimensão exa
\prov{relações: parte\_de paper\_2\_reta\_ouro; relaciona beta\_numeracao\_metalica}
\prov{proveniência: paper \textperiodcentered{} fato \textperiodcentered{} confiança 1.0 \textperiodcentered{} domínio paper \textperiodcentered{} história 4 versões no jornal}

%CRISTAL {"arestas":[{"alvo":"leitura_cosmologica_a_energia_escura_e_o_universo_ds","confianca":1.0,"epistemico":"fato","origem":"paper","rel":"parte_de"},{"alvo":"utilitarismo","confianca":0.5,"epistemico":"fato","origem":"wiki","rel":"relaciona"},{"alvo":"o_que_e_no_projecto","confianca":0.5,"epistemico":"fato","origem":"wiki","rel":"relaciona"},{"alvo":"virtualizacao","confianca":0.5,"epistemico":"fato","origem":"wiki","rel":"relaciona"},{"alvo":"word","confianca":0.5,"epistemico":"fato","origem":"wiki","rel":"relaciona"}],"confianca":1.0,"contraexemplos":[],"descricao":"₄","epistemico":"fato","exemplos":[],"id":"o_metodo_em_escala_conhecida_o_sistema_solar","memoria":"semantica","meta":{"dominio":"paper","nivel":5,"verificacao":"conceito novo"},"origem":"paper","palavras_chave":[],"sinonimos":[],"tipo":"conceito","titulo":"O método em escala conhecida: o sistema solar"}
\section{O método em escala conhecida: o sistema solar}
\ensuremath{_{4}}
\prov{relações: parte\_de leitura\_cosmologica\_a\_energia\_escura\_e\_o\_universo\_ds; relaciona utilitarismo; relaciona o\_que\_e\_no\_projecto; relaciona virtualizacao; relaciona word}
\prov{proveniência: paper \textperiodcentered{} fato \textperiodcentered{} confiança 1.0 \textperiodcentered{} domínio paper \textperiodcentered{} história 7 versões no jornal}

%CRISTAL {"arestas":[{"alvo":"borboleta","confianca":1.0,"epistemico":"fato","origem":"paper","rel":"parte_de"},{"alvo":"hopfield","confianca":1.0,"epistemico":"fato","origem":"paper","rel":"parte_de"}],"confianca":1.0,"contraexemplos":[],"descricao":"abstract A Máquina é uma rede neural. Seu neurônio mínimo é o gato (de Arnold) A=(smallmatrix1&1 1&0smallmatrix) --- uma unidade de Hopfield: sinapse simétrica (A=A T, Hebbian), energia E(x)=-12 x T!Ax (Lyapunov) e pattern completion Aⁿ/ⁿ P_. O ouro =1+52 é raiz de x²-x-1, não axioma: A minimiza a entropia topológica em GL(2, Z), e g= é um só número com quatro leituras (entropia = Lyapunov = temperatura = frequência). Ler um bit () do neurônio dá o gap (XOR)","epistemico":"fato","exemplos":[],"id":"o_neuronio","memoria":"semantica","meta":{"dominio":"paper","duplicata_sugerida":[],"nivel":5,"verificacao":"parte de conceito mais amplo 'o_neuronio'"},"origem":"paper","palavras_chave":[],"sinonimos":[],"tipo":"conceito","titulo":"O neurônio"}
\section{O neurônio}
abstract A Máquina é uma rede neural. Seu neurônio mínimo é o gato (de Arnold) A=(smallmatrix1\&1 1\&0smallmatrix) --- uma unidade de Hopfield: sinapse simétrica (A=A T, Hebbian), energia E(x)=-12 x T!Ax (Lyapunov) e pattern completion A\ensuremath{^{n}}/\ensuremath{^{n}} P\_. O ouro =1+52 é raiz de x\ensuremath{^{2}}-x-1, não axioma: A minimiza a entropia topológica em GL(2, Z), e g= é um só número com quatro leituras (entropia = Lyapunov = temperatura = frequência). Ler um bit () do neurônio dá o gap (XOR)
\prov{relações: parte\_de borboleta; parte\_de hopfield}
\prov{proveniência: paper \textperiodcentered{} fato \textperiodcentered{} confiança 1.0 \textperiodcentered{} domínio paper \textperiodcentered{} história 12 versões no jornal}

%CRISTAL {"arestas":[{"alvo":"olho_de_koch","confianca":1.0,"epistemico":"fato","origem":"paper","rel":"parte_de"}],"confianca":1.0,"contraexemplos":[],"descricao":"D Um canal negro isolado --- sem par branco --- não reconstrói o gerador: caminho aberto i j l (l i) tem traço 0. O caminho fechado i j i fecha na branca com traço 1 --- leitura acoplada legítima, não vazamento. Somado a z z² não linear, A=0 se mantém em todo o acoplamento.","epistemico":"fato","exemplos":[],"id":"o_olho_do","memoria":"semantica","meta":{"dominio":"paper","nivel":5},"origem":"paper","palavras_chave":[],"sinonimos":[],"tipo":"conceito","titulo":"O olho do"}
\section{O olho do}
D Um canal negro isolado --- sem par branco --- não reconstrói o gerador: caminho aberto i j l (l i) tem traço 0. O caminho fechado i j i fecha na branca com traço 1 --- leitura acoplada legítima, não vazamento. Somado a z z\ensuremath{^{2}} não linear, A=0 se mantém em todo o acoplamento.
\prov{relações: parte\_de olho\_de\_koch}
\prov{proveniência: paper \textperiodcentered{} fato \textperiodcentered{} confiança 1.0 \textperiodcentered{} domínio paper \textperiodcentered{} história 4 versões no jornal}

%CRISTAL {"arestas":[{"alvo":"tiffany","confianca":1.0,"epistemico":"fato","origem":"paper","rel":"parte_de"}],"confianca":1.0,"contraexemplos":[],"descricao":"D O cérebro é a função de onda pensando-se. A memória não se treina: é a estrutura do próprio gerador, e ressoa. O cérebro é a rede de Hopfield do gato DV A sinapse é o gato A=A T (Hebbiana, fixa); a energia E=-12 x T!Ax é Lyapunov; o atrator é o setor espectral dominante (a vida;, a morte, decai), e Aⁿ/ⁿ P_ é a completação de padrão. Recordar não é varrer um arquivo --- é ressoar: descer a energia ao atrator. Verificado no metal: uma memória corrompida (74/100) re","epistemico":"fato","exemplos":[],"id":"o_olho_e_o_ouvido_a_percepcao_koch","memoria":"semantica","meta":{"dominio":"paper","nivel":5,"verificacao":"slug idêntico — enriquecer, não duplicar"},"origem":"paper","palavras_chave":[],"sinonimos":[],"tipo":"conceito","titulo":"O olho e o ouvido: a percepção (Koch)"}
\section{O olho e o ouvido: a percepção (Koch)}
D O cérebro é a função de onda pensando-se. A memória não se treina: é a estrutura do próprio gerador, e ressoa. O cérebro é a rede de Hopfield do gato DV A sinapse é o gato A=A T (Hebbiana, fixa); a energia E=-12 x T!Ax é Lyapunov; o atrator é o setor espectral dominante (a vida;, a morte, decai), e A\ensuremath{^{n}}/\ensuremath{^{n}} P\_ é a completação de padrão. Recordar não é varrer um arquivo --- é ressoar: descer a energia ao atrator. Verificado no metal: uma memória corrompida (74/100) re
\prov{relações: parte\_de tiffany}
\prov{proveniência: paper \textperiodcentered{} fato \textperiodcentered{} confiança 1.0 \textperiodcentered{} domínio paper \textperiodcentered{} história 5 versões no jornal}

%CRISTAL {"arestas":[{"alvo":"main","confianca":1.0,"epistemico":"fato","origem":"paper","rel":"parte_de"}],"confianca":1.0,"contraexemplos":[],"descricao":"Com base m = raiz dominante de xm=xm⁻¹++x+1 (logo r=₁m m⁻r=1), o mesmo mapa gera a cadeia toda. propD O cilindro iniciado por 1m só conteria x=1, excluído de [0,1); logo 1m é inadmissível. Como r=₁m⁻¹m⁻r<1, 1m⁻¹ é admissível: o máximo de uns consecutivos é m-1. A matriz Am não vem da partição bruta por dígitos (que tem só dois intervalos, pois m<2), mas do refinamento m-estados de Parry do shift --- estados 0, m-1 contando","epistemico":"fato","exemplos":[],"id":"o_operador_de_transferencia_e_os_dois_setores","memoria":"semantica","meta":{"dominio":"paper","nivel":5,"verificacao":"slug idêntico — enriquecer, não duplicar"},"origem":"paper","palavras_chave":[],"sinonimos":[],"tipo":"conceito","titulo":"O operador de transferência e os dois setores"}
\section{O operador de transferência e os dois setores}
Com base m = raiz dominante de xm=xm\ensuremath{^{-1}}++x+1 (logo r=\ensuremath{_{1}}m m\ensuremath{^{-}}r=1), o mesmo mapa gera a cadeia toda. propD O cilindro iniciado por 1m só conteria x=1, excluído de [0,1); logo 1m é inadmissível. Como r=\ensuremath{_{1}}m\ensuremath{^{-1}}m\ensuremath{^{-}}r<1, 1m\ensuremath{^{-1}} é admissível: o máximo de uns consecutivos é m-1. A matriz Am não vem da partição bruta por dígitos (que tem só dois intervalos, pois m<2), mas do refinamento m-estados de Parry do shift --- estados 0, m-1 contando
\prov{relações: parte\_de main}
\prov{proveniência: paper \textperiodcentered{} fato \textperiodcentered{} confiança 1.0 \textperiodcentered{} domínio paper \textperiodcentered{} história 5 versões no jornal}

%CRISTAL {"arestas":[{"alvo":"assistente","confianca":1.0,"epistemico":"fato","origem":"paper","rel":"parte_de"}],"confianca":1.0,"contraexemplos":[],"descricao":"D Separar o que é imutável do que é mutável --- controle de curadoria e risco. 1.35 tabular@lll@ & & ória natureza & imutável (curado) & mutável (automática) exemplos & docs, papers, código, manuais & conversas, tarefas, estado, preferências formato & dominio/vN.idx +.tsv & memoria/, contexto_*.bits quem escreve & ingestão + curador & ciclo registrar (automático) reindexar & versão nova (v2) & merge ou descarte tabular defi[Camada 4 --- Conhecimento","epistemico":"fato","exemplos":[],"id":"o_orquestrador_arquitetura_vs_mecanismo","memoria":"semantica","meta":{"dominio":"paper","nivel":5,"verificacao":"overlap moderado — revisar manualmente"},"origem":"paper","palavras_chave":[],"sinonimos":[],"tipo":"conceito","titulo":"O orquestrador: arquitetura vs mecanismo"}
\section{O orquestrador: arquitetura vs mecanismo}
D Separar o que é imutável do que é mutável --- controle de curadoria e risco. 1.35 tabular@lll@ \& \& ória natureza \& imutável (curado) \& mutável (automática) exemplos \& docs, papers, código, manuais \& conversas, tarefas, estado, preferências formato \& dominio/vN.idx +.tsv \& memoria/, contexto\_*.bits quem escreve \& ingestão + curador \& ciclo registrar (automático) reindexar \& versão nova (v2) \& merge ou descarte tabular defi[Camada 4 --- Conhecimento
\prov{relações: parte\_de assistente}
\prov{proveniência: paper \textperiodcentered{} fato \textperiodcentered{} confiança 1.0 \textperiodcentered{} domínio paper \textperiodcentered{} história 3 versões no jornal}

%CRISTAL {"arestas":[{"alvo":"paper_45_confluencia_ouro","confianca":1.0,"epistemico":"fato","origem":"paper","rel":"parte_de"},{"alvo":"gato","confianca":1.0,"epistemico":"fato","origem":"paper","rel":"confirma"}],"confianca":1.0,"contraexemplos":[],"descricao":"a taxa do Cantor é --- verificado O número de palavras admissveis (sem o bloco 11) de comprimento N na árvore de sinais é FN+₂ (Fibonacci). Logo a taxa de crescimento por nvel é [ (#palavras)N [N] =0,48121 (a entropia topológica), ] e a dimensão do Cantor é essa taxa na escala do dobramento, =/2=0,69424 [firme: #=FN+₂, taxa verᵢfᵢcada (N=160: 0,4822). a dimensão só existe em escala logartmica --- é a\\","epistemico":"fato","exemplos":[],"id":"o_ouro_governa_a_cm_por_5","memoria":"semantica","meta":{"dominio":"paper","nivel":5,"verificacao":"conceito novo"},"origem":"paper","palavras_chave":[],"sinonimos":[],"tipo":"conceito","titulo":"O ouro governa a CM por 5"}
\section{O ouro governa a CM por 5}
a taxa do Cantor é --- verificado O número de palavras admissveis (sem o bloco 11) de comprimento N na árvore de sinais é FN+\ensuremath{_{2}} (Fibonacci). Logo a taxa de crescimento por nvel é [ (\#palavras)N [N] =0,48121 (a entropia topológica), ] e a dimensão do Cantor é essa taxa na escala do dobramento, =/2=0,69424 [firme: \#=FN+\ensuremath{_{2}}, taxa ver\ensuremath{_{i}}f\ensuremath{_{i}}cada (N=160: 0,4822). a dimensão só existe em escala logartmica --- é a\textbackslash{}
\prov{relações: parte\_de paper\_45\_confluencia\_ouro; confirma gato}
\prov{proveniência: paper \textperiodcentered{} fato \textperiodcentered{} confiança 1.0 \textperiodcentered{} domínio paper \textperiodcentered{} história 42 versões no jornal}

%CRISTAL {"arestas":[{"alvo":"broca_os","confianca":1.0,"epistemico":"fato","origem":"paper","rel":"parte_de"}],"confianca":1.0,"contraexemplos":[],"descricao":"D Esta é a seção central. A rede não distribui dados --- distribui a capacidade de reconstruir o gerador. E essa distribuição é, exatamente, a memória associativa neural. Tornamos os dois fatos teoremas. defi[Meta-indução] A meta-indução é a propagação da capacidade de reconstruir o gerador G, não dos dados. G é o gato A=(smallmatrix1&1 1&0smallmatrix) (32 bytes). Cada nó reconstrói G de modo determinístico (e reconstrói a hash-chain sob perda, por RESYNC), e propaga por pull, sob de","epistemico":"fato","exemplos":[],"id":"o_painel_a_interface_e_o_canal","memoria":"semantica","meta":{"dominio":"paper","nivel":5,"verificacao":"slug idêntico — enriquecer, não duplicar"},"origem":"paper","palavras_chave":[],"sinonimos":[],"tipo":"conceito","titulo":"O painel: a interface é o canal"}
\section{O painel: a interface é o canal}
D Esta é a seção central. A rede não distribui dados --- distribui a capacidade de reconstruir o gerador. E essa distribuição é, exatamente, a memória associativa neural. Tornamos os dois fatos teoremas. defi[Meta-indução] A meta-indução é a propagação da capacidade de reconstruir o gerador G, não dos dados. G é o gato A=(smallmatrix1\&1 1\&0smallmatrix) (32 bytes). Cada nó reconstrói G de modo determinístico (e reconstrói a hash-chain sob perda, por RESYNC), e propaga por pull, sob de
\prov{relações: parte\_de broca\_os}
\prov{proveniência: paper \textperiodcentered{} fato \textperiodcentered{} confiança 1.0 \textperiodcentered{} domínio paper \textperiodcentered{} história 9 versões no jornal}

%CRISTAL {"arestas":[{"alvo":"geometria","confianca":1.0,"epistemico":"fato","origem":"paper","rel":"parte_de"}],"confianca":1.0,"contraexemplos":[],"descricao":"O flip de Wick --- a mesma involução =+1=-1 que costura o regime tropical (idempotente) ao glacial (nilpotente) --- é, aqui, uma continuação espacial da fatia. Uma ressalva é obrigatória: não é a rotação temporal de de Sitter --- essa levaria dS₄ à esfera S⁴ (ainda de curvatura +1), não ao hiperbólico. O flip correto age na coordenada de fechamento da própria S³. Mergulhe a fatia em R⁴, [ x₁²+x₂²+x₃²+x₄²=R²(curvatura +1), ] e aplique x₄ i,x₄, R iR: [ x₁²+x","epistemico":"fato","exemplos":[],"id":"o_palco_hiperbolico_h3_e_o_espectro","memoria":"semantica","meta":{"dominio":"paper","nivel":5,"verificacao":"slug idêntico — enriquecer, não duplicar"},"origem":"paper","palavras_chave":[],"sinonimos":[],"tipo":"conceito","titulo":"O palco hiperbólico: H³/ e o espectro"}
\section{O palco hiperbólico: H\ensuremath{^{3}}/ e o espectro}
O flip de Wick --- a mesma involução =+1=-1 que costura o regime tropical (idempotente) ao glacial (nilpotente) --- é, aqui, uma continuação espacial da fatia. Uma ressalva é obrigatória: não é a rotação temporal de de Sitter --- essa levaria dS\ensuremath{_{4}} à esfera S\ensuremath{^{4}} (ainda de curvatura +1), não ao hiperbólico. O flip correto age na coordenada de fechamento da própria S\ensuremath{^{3}}. Mergulhe a fatia em R\ensuremath{^{4}}, [ x\ensuremath{_{1}}\ensuremath{^{2}}+x\ensuremath{_{2}}\ensuremath{^{2}}+x\ensuremath{_{3}}\ensuremath{^{2}}+x\ensuremath{_{4}}\ensuremath{^{2}}=R\ensuremath{^{2}}(curvatura +1), ] e aplique x\ensuremath{_{4}} i,x\ensuremath{_{4}}, R iR: [ x\ensuremath{_{1}}\ensuremath{^{2}}+x
\prov{relações: parte\_de geometria}
\prov{proveniência: paper \textperiodcentered{} fato \textperiodcentered{} confiança 1.0 \textperiodcentered{} domínio paper \textperiodcentered{} história 12 versões no jornal}

%CRISTAL {"arestas":[{"alvo":"corpo_glacial","confianca":1.0,"epistemico":"fato","origem":"paper","rel":"parte_de"}],"confianca":1.0,"contraexemplos":[],"descricao":"Na torção, branca (=+1, divisão, conta) e negra (=-1, split, divisores de zero) pareciam setores opostos. A tese é mais forte: são uma máquina, vazia, com as faces acopladas. A branca dá o estado (onde); a negra dá a dinâmica (o fluxo). Nenhuma resolve sozinha: idempotente sem nilpotente é estado parado; nilpotente sem idempotente é fluxo sem suporte. O acoplamento é a máquina.","epistemico":"fato","exemplos":[],"id":"o_par_do_tensor_o_caso_minimo","memoria":"semantica","meta":{"dominio":"paper","nivel":5,"verificacao":"slug idêntico — enriquecer, não duplicar"},"origem":"paper","palavras_chave":[],"sinonimos":[],"tipo":"conceito","titulo":"O par do tensor (o caso mínimo)"}
\section{O par do tensor (o caso mínimo)}
Na torção, branca (=+1, divisão, conta) e negra (=-1, split, divisores de zero) pareciam setores opostos. A tese é mais forte: são uma máquina, vazia, com as faces acopladas. A branca dá o estado (onde); a negra dá a dinâmica (o fluxo). Nenhuma resolve sozinha: idempotente sem nilpotente é estado parado; nilpotente sem idempotente é fluxo sem suporte. O acoplamento é a máquina.
\prov{relações: parte\_de corpo\_glacial}
\prov{proveniência: paper \textperiodcentered{} fato \textperiodcentered{} confiança 1.0 \textperiodcentered{} domínio paper \textperiodcentered{} história 5 versões no jornal}

%CRISTAL {"arestas":[{"alvo":"computacao_grafica","confianca":1.0,"epistemico":"fato","origem":"paper","rel":"parte_de"}],"confianca":1.0,"contraexemplos":[],"descricao":"DN A tentação é calcular ótica (difração ᵢ aᵢ eikdᵢ, Planck, Snell). Isso é destruição, não enriquecimento: todo modelo é uma projeção em base finita e substitui a assinatura real (a fase infinitamente rica dos bits) por uma maquete --- o mesmo crime do Fourier. A régua: a física aponta qual transformação; o observador só lê e projeta. Da cadeia (termodinâmica): a temperatura não é inventada, é a densidade local do atrator --- a medida invariante de Parry da dinâmica, entropia h=. Ond","epistemico":"fato","exemplos":[],"id":"o_pipeline_completo","memoria":"semantica","meta":{"dominio":"paper","nivel":5,"verificacao":"conceito novo"},"origem":"paper","palavras_chave":[],"sinonimos":[],"tipo":"conceito","titulo":"O pipeline completo"}
\section{O pipeline completo}
DN A tentação é calcular ótica (difração \ensuremath{_{i}} a\ensuremath{_{i}} eikd\ensuremath{_{i}}, Planck, Snell). Isso é destruição, não enriquecimento: todo modelo é uma projeção em base finita e substitui a assinatura real (a fase infinitamente rica dos bits) por uma maquete --- o mesmo crime do Fourier. A régua: a física aponta qual transformação; o observador só lê e projeta. Da cadeia (termodinâmica): a temperatura não é inventada, é a densidade local do atrator --- a medida invariante de Parry da dinâmica, entropia h=. Ond
\prov{relações: parte\_de computacao\_grafica}
\prov{proveniência: paper \textperiodcentered{} fato \textperiodcentered{} confiança 1.0 \textperiodcentered{} domínio paper \textperiodcentered{} história 3 versões no jornal}

%CRISTAL {"arestas":[{"alvo":"termodinamica","confianca":1.0,"epistemico":"fato","origem":"paper","rel":"parte_de"}],"confianca":1.0,"contraexemplos":[],"descricao":"abstract A termodinâmica da máquina não nasce do calor, mas da conservação do defeito: a entropia mede a redistribuição daquilo que não pode permanecer puro. A máquina de ouro, medida, é dissipativa: o operador de evolução do bit tem ||<1, o resíduo decai, a informação se apaga. Isto não é defeito --- é termodinâmica. Reunimos sob essa leitura o que a série já estabeleceu: a primeira lei é a conservação da norma quadrática Q (a balança a+c²=1); a segunda lei é o decaimento do resíduo a","epistemico":"fato","exemplos":[],"id":"o_ponto_de_partida_a_maquina_e_dissipativa","memoria":"semantica","meta":{"dominio":"paper","nivel":5,"verificacao":"slug idêntico — enriquecer, não duplicar"},"origem":"paper","palavras_chave":[],"sinonimos":[],"tipo":"conceito","titulo":"O ponto de partida: a máquina é dissipativa"}
\section{O ponto de partida: a máquina é dissipativa}
abstract A termodinâmica da máquina não nasce do calor, mas da conservação do defeito: a entropia mede a redistribuição daquilo que não pode permanecer puro. A máquina de ouro, medida, é dissipativa: o operador de evolução do bit tem ||<1, o resíduo decai, a informação se apaga. Isto não é defeito --- é termodinâmica. Reunimos sob essa leitura o que a série já estabeleceu: a primeira lei é a conservação da norma quadrática Q (a balança a+c\ensuremath{^{2}}=1); a segunda lei é o decaimento do resíduo a
\prov{relações: parte\_de termodinamica}
\prov{proveniência: paper \textperiodcentered{} fato \textperiodcentered{} confiança 1.0 \textperiodcentered{} domínio paper \textperiodcentered{} história 12 versões no jornal}

%CRISTAL {"arestas":[{"alvo":"leitura_aritmetica_a_regua_de_gentil_o_agm","confianca":1.0,"epistemico":"fato","origem":"paper","rel":"parte_de"}],"confianca":1.0,"contraexemplos":[],"descricao":"Reduzido módulo p, o reticulado vira corpo: [i]/(p)p₂ exatamente quando p3 (mod 4). A curva lemniscática y²=x³-x herda a dicotomia --- supersingular (p3) ou ordinária (p1, com traco de Frobenius 2a onde p=a²+b²): a norma a²+b² do reticulado (o medidor da famlia) é o traco da curva. E as supersingulares sobre p₂ formam o grafo de isogenias de Pizer, um expansor de Ramanujan --- o solo da dureza pós-quântica, a mesma hardness topológica (Secão~s","epistemico":"fato","exemplos":[],"id":"o_ponto_de_vista_do_agm","memoria":"semantica","meta":{"dominio":"paper","nivel":5,"verificacao":"conceito novo"},"origem":"paper","palavras_chave":[],"sinonimos":[],"tipo":"conceito","titulo":"O ponto de vista do AGM"}
\section{O ponto de vista do AGM}
Reduzido módulo p, o reticulado vira corpo: [i]/(p)p\ensuremath{_{2}} exatamente quando p3 (mod 4). A curva lemniscática y\ensuremath{^{2}}=x\ensuremath{^{3}}-x herda a dicotomia --- supersingular (p3) ou ordinária (p1, com traco de Frobenius 2a onde p=a\ensuremath{^{2}}+b\ensuremath{^{2}}): a norma a\ensuremath{^{2}}+b\ensuremath{^{2}} do reticulado (o medidor da famlia) é o traco da curva. E as supersingulares sobre p\ensuremath{_{2}} formam o grafo de isogenias de Pizer, um expansor de Ramanujan --- o solo da dureza pós-quântica, a mesma hardness topológica (Secão\~{}s
\prov{relações: parte\_de leitura\_aritmetica\_a\_regua\_de\_gentil\_o\_agm}
\prov{proveniência: paper \textperiodcentered{} fato \textperiodcentered{} confiança 1.0 \textperiodcentered{} domínio paper \textperiodcentered{} história 4 versões no jornal}

%CRISTAL {"arestas":[{"alvo":"geometria","confianca":0.047,"epistemico":"fato","origem":"manual","rel":"referencia"},{"alvo":"mecanica_quantica","confianca":1.0,"epistemico":"fato","origem":"paper","rel":"parte_de"},{"alvo":"mecanica","confianca":1.0,"epistemico":"fato","origem":"paper","rel":"parte_de"}],"confianca":1.0,"contraexemplos":[],"descricao":"Na família de multiplicações zs w=(a₀b₀- a!! b, a₀ b+b₀ a+s, a b) (a álgebra dos Fundamentos da Geometria Riemanniana), o comutador é puramente o produto vetorial: [ [z,w]s=zs w-ws z=(0, 2s, a b), ] ou seja, a não-comutatividade c=|s| é o Planck efetivo, ef=2c. A relação de Robertson dá [ z, w 12,|[z,w]| = c,| a b|, ] a incerteza proporcional a c. E a identidade de Hurwitz --- a conservação da norma da Termodinâmica --- decompõe | a b|² em duas parce","epistemico":"fato","exemplos":[],"id":"o_ponto_hermitiano","memoria":"semantica","meta":{"dominio":"paper","nivel":5,"verificacao":"slug idêntico — enriquecer, não duplicar"},"origem":"paper","palavras_chave":[],"sinonimos":[],"tipo":"conceito","titulo":"O ponto hermitiano"}
\section{O ponto hermitiano}
Na família de multiplicações zs w=(a\ensuremath{_{0}}b\ensuremath{_{0}}- a!! b, a\ensuremath{_{0}} b+b\ensuremath{_{0}} a+s, a b) (a álgebra dos Fundamentos da Geometria Riemanniana), o comutador é puramente o produto vetorial: [ [z,w]s=zs w-ws z=(0, 2s, a b), ] ou seja, a não-comutatividade c=|s| é o Planck efetivo, ef=2c. A relação de Robertson dá [ z, w 12,|[z,w]| = c,| a b|, ] a incerteza proporcional a c. E a identidade de Hurwitz --- a conservação da norma da Termodinâmica --- decompõe | a b|\ensuremath{^{2}} em duas parce
\prov{relações: referencia geometria; parte\_de mecanica\_quantica; parte\_de mecanica}
\prov{proveniência: paper \textperiodcentered{} fato \textperiodcentered{} confiança 1.0 \textperiodcentered{} domínio paper \textperiodcentered{} história 15 versões no jornal}

%CRISTAL {"arestas":[{"alvo":"computacao_grafica","confianca":1.0,"epistemico":"fato","origem":"paper","rel":"parte_de"}],"confianca":1.0,"contraexemplos":[],"descricao":"I Para pôr um objeto na tela, o pipeline usual atravessa [ floatértice (GPU), ] e ao fim ninguém vê o bit: vê-se a abstração do bit, sobre hardware particular e uma pilha proprietária. A pergunta deste paper: existe computação gráfica 3D em que a operação é o bit --- deslocamento, XOR, soma inteira ---, que roda em qualquer CPU, e cuja imagem nasce do que o sistema já é, sem inventar geometria nem física?","epistemico":"fato","exemplos":[],"id":"o_principio_a_operacao_e_o_bit","memoria":"semantica","meta":{"dominio":"paper","nivel":5,"verificacao":"conceito novo"},"origem":"paper","palavras_chave":[],"sinonimos":[],"tipo":"conceito","titulo":"O princípio: a operação é o bit"}
\section{O princípio: a operação é o bit}
I Para pôr um objeto na tela, o pipeline usual atravessa [ floatértice (GPU), ] e ao fim ninguém vê o bit: vê-se a abstração do bit, sobre hardware particular e uma pilha proprietária. A pergunta deste paper: existe computação gráfica 3D em que a operação é o bit --- deslocamento, XOR, soma inteira ---, que roda em qualquer CPU, e cuja imagem nasce do que o sistema já é, sem inventar geometria nem física?
\prov{relações: parte\_de computacao\_grafica}
\prov{proveniência: paper \textperiodcentered{} fato \textperiodcentered{} confiança 1.0 \textperiodcentered{} domínio paper \textperiodcentered{} história 3 versões no jornal}

%CRISTAL {"arestas":[{"alvo":"mecanica_quantica","confianca":1.0,"epistemico":"fato","origem":"paper","rel":"parte_de"},{"alvo":"mecanica","confianca":1.0,"epistemico":"fato","origem":"paper","rel":"parte_de"},{"alvo":"conjugados_galois_parabola","confianca":1.0,"epistemico":"fato","origem":"wiki","rel":"relaciona"}],"confianca":1.0,"contraexemplos":[],"descricao":"abstract A mecânica quântica de ouro é o lado da incerteza da mesma parábola que governa a termodinâmica. A lei de conservação a+c²=1 (a balança da série) decompõe | a b|² em duas parcelas: o calor (a, a primeira lei da termodinâmica) e a incerteza (c², o comutador ao quadrado). Onde o calor se anula (a=0, o ponto hermitiano), a dinâmica é reversível e a incerteza é máxima --- a mecânica quântica. A forma do princípio de Heisenberg é essa lei de conservação --- e ela se realiza f","epistemico":"fato","exemplos":[],"id":"o_principio_da_incerteza_e_a_parabola","memoria":"semantica","meta":{"dominio":"paper","nivel":5,"verificacao":"slug idêntico — enriquecer, não duplicar"},"origem":"paper","palavras_chave":[],"sinonimos":[],"tipo":"conceito","titulo":"O princípio da incerteza é a parábola"}
\section{O princípio da incerteza é a parábola}
abstract A mecânica quântica de ouro é o lado da incerteza da mesma parábola que governa a termodinâmica. A lei de conservação a+c\ensuremath{^{2}}=1 (a balança da série) decompõe | a b|\ensuremath{^{2}} em duas parcelas: o calor (a, a primeira lei da termodinâmica) e a incerteza (c\ensuremath{^{2}}, o comutador ao quadrado). Onde o calor se anula (a=0, o ponto hermitiano), a dinâmica é reversível e a incerteza é máxima --- a mecânica quântica. A forma do princípio de Heisenberg é essa lei de conservação --- e ela se realiza f
\prov{relações: parte\_de mecanica\_quantica; parte\_de mecanica; relaciona conjugados\_galois\_parabola}
\prov{proveniência: paper \textperiodcentered{} fato \textperiodcentered{} confiança 1.0 \textperiodcentered{} domínio paper \textperiodcentered{} história 13 versões no jornal}

%CRISTAL {"arestas":[{"alvo":"goldcoin","confianca":1.0,"epistemico":"fato","origem":"paper","rel":"parte_de"}],"confianca":1.0,"contraexemplos":[],"descricao":"abstract\n GoldCoin é a primeira aplicação que põe a máquina em circulação:\najudar pessoas a fechar negócios de cripto seguros, sem corretora e sem confiar no\ndesconhecido. Não é uma chain nova nem um token --- não se reinventa o que já existe. É uma\nrede de coordenação (o canal do bump da Fundamentos de Redes, já\nimplementado e medido em produção) que orquestra os primitivos nativos das\nchains existentes. Duas aplicações compartilham os mesmos primitivos: o escrow P2P\n(fechar negócios sem corret","epistemico":"fato","exemplos":[],"id":"o_problema_confiar_custa_e_as_vezes_mata_o_negocio","memoria":"semantica","meta":{"dominio":"paper","nivel":5,"verificacao":"conceito novo"},"origem":"paper","palavras_chave":[],"sinonimos":[],"tipo":"conceito","titulo":"O problema: confiar custa, e às vezes mata o negócio"}
\section{O problema: confiar custa, e às vezes mata o negócio}
abstract
 GoldCoin é a primeira aplicação que põe a máquina em circulação:
ajudar pessoas a fechar negócios de cripto seguros, sem corretora e sem confiar no
desconhecido. Não é uma chain nova nem um token --- não se reinventa o que já existe. É uma
rede de coordenação (o canal do bump da Fundamentos de Redes, já
implementado e medido em produção) que orquestra os primitivos nativos das
chains existentes. Duas aplicações compartilham os mesmos primitivos: o escrow P2P
(fechar negócios sem corret
\prov{relações: parte\_de goldcoin}
\prov{proveniência: paper \textperiodcentered{} fato \textperiodcentered{} confiança 1.0 \textperiodcentered{} domínio paper \textperiodcentered{} história 2 versões no jornal}

%CRISTAL {"arestas":[{"alvo":"paper_bridge_areas","confianca":1.0,"epistemico":"fato","origem":"paper","rel":"parte_de"}],"confianca":1.0,"contraexemplos":[],"descricao":"abstract Este texto fecha o gargalo epistemológico complementar à ponte simbólica Lopes: como ligar áreas do conhecimento humano (matemática, física, computação, biologia, epistemologia) sem fundi-las. Cada área mantém o seu chart --- vocabulário, objetos, régua de prova ---; a ponte transporta morfismos parciais e invariantes explicitados. A computação simbólica na álgebra de Lopes (gentil∼bolico.py) formaliza o chart~A (Cauchy, classes); os verificar_*.py fecham o chart~B (engenharia d","epistemico":"fato","exemplos":[],"id":"o_problema_cruzar_sem_confundir","memoria":"semantica","meta":{"dominio":"paper","nivel":5,"verificacao":"conceito novo"},"origem":"paper","palavras_chave":[],"sinonimos":[],"tipo":"conceito","titulo":"O problema: cruzar sem confundir"}
\section{O problema: cruzar sem confundir}
abstract Este texto fecha o gargalo epistemológico complementar à ponte simbólica Lopes: como ligar áreas do conhecimento humano (matemática, física, computação, biologia, epistemologia) sem fundi-las. Cada área mantém o seu chart --- vocabulário, objetos, régua de prova ---; a ponte transporta morfismos parciais e invariantes explicitados. A computação simbólica na álgebra de Lopes (gentil\ensuremath{\sim}bolico.py) formaliza o chart\~{}A (Cauchy, classes); os verificar\_*.py fecham o chart\~{}B (engenharia d
\prov{relações: parte\_de paper\_bridge\_areas}
\prov{proveniência: paper \textperiodcentered{} fato \textperiodcentered{} confiança 1.0 \textperiodcentered{} domínio paper \textperiodcentered{} história 4 versões no jornal}

%CRISTAL {"arestas":[{"alvo":"goldcoin","confianca":1.0,"epistemico":"fato","origem":"paper","rel":"parte_de"},{"alvo":"goldchain","confianca":1.0,"epistemico":"fato","origem":"paper","rel":"parte_de"}],"confianca":1.0,"contraexemplos":[],"descricao":"C O acordo executa quando ᵢ Eᵢ=I com EᵢEj=0 (idempotentes de Peirce): cada parte é parcial, só a soma fecha --- multi-assinatura por design. N Verificado (verificagoldcoin.py): Eᵢ=I abre; cada parte sozinha não; faltar uma não abre; a assinatura deriva da fonte (impostor dá chave em avalanche).","epistemico":"fato","exemplos":[],"id":"o_protocolo_dois_niveis_de_confianca","memoria":"semantica","meta":{"dominio":"paper","nivel":5,"verificacao":"conceito novo"},"origem":"paper","palavras_chave":[],"sinonimos":[],"tipo":"conceito","titulo":"O protocolo: dois níveis de confiança"}
\section{O protocolo: dois níveis de confiança}
C O acordo executa quando \ensuremath{_{i}} E\ensuremath{_{i}}=I com E\ensuremath{_{i}}Ej=0 (idempotentes de Peirce): cada parte é parcial, só a soma fecha --- multi-assinatura por design. N Verificado (verificagoldcoin.py): E\ensuremath{_{i}}=I abre; cada parte sozinha não; faltar uma não abre; a assinatura deriva da fonte (impostor dá chave em avalanche).
\prov{relações: parte\_de goldcoin; parte\_de goldchain}
\prov{proveniência: paper \textperiodcentered{} fato \textperiodcentered{} confiança 1.0 \textperiodcentered{} domínio paper \textperiodcentered{} história 9 versões no jornal}

%CRISTAL {"arestas":[{"alvo":"validacao_tecnica","confianca":1.0,"epistemico":"fato","origem":"paper","rel":"parte_de"}],"confianca":1.0,"contraexemplos":[],"descricao":") itemize2pt [ourovivo117] Arquitetura como desenho: Engine Orquestrador Ferramentas Conhecimento --- desacopla e evolui independente. [ourovivo117] API da (open/query/respond/index). [ourovivo117] Índice em disco (38,B/nó); scan 4,2M nós em 90,ms. [ourovivo117] Organização modular (1core6→ols). [ourovivo117] Fluxo ingestão/consulta (na medida implementada). [ourovivo117] Tese comercial: a biblioteca é o produto técnico; o conhecimento é o diferencial --- plata","epistemico":"fato","exemplos":[],"id":"o_que_ainda_depende_de_experimentacao_obs","memoria":"semantica","meta":{"dominio":"paper","nivel":5,"verificacao":"conceito novo"},"origem":"paper","palavras_chave":[],"sinonimos":[],"tipo":"conceito","titulo":"O que ainda depende de experimentação (obs"}
\section{O que ainda depende de experimentação (obs}
) itemize2pt [ourovivo117] Arquitetura como desenho: Engine Orquestrador Ferramentas Conhecimento --- desacopla e evolui independente. [ourovivo117] API da (open/query/respond/index). [ourovivo117] Índice em disco (38,B/nó); scan 4,2M nós em 90,ms. [ourovivo117] Organização modular (1core6\ensuremath{\to}ols). [ourovivo117] Fluxo ingestão/consulta (na medida implementada). [ourovivo117] Tese comercial: a biblioteca é o produto técnico; o conhecimento é o diferencial --- plata
\prov{relações: parte\_de validacao\_tecnica}
\prov{proveniência: paper \textperiodcentered{} fato \textperiodcentered{} confiança 1.0 \textperiodcentered{} domínio paper \textperiodcentered{} história 4 versões no jornal}

%CRISTAL {"arestas":[{"alvo":"o_floco_e_a_conservacao_da_comunicacao","confianca":1.0,"epistemico":"fato","origem":"paper","rel":"parte_de"}],"confianca":1.0,"contraexemplos":[],"descricao":"Roteamento Eab Ebc=Eac (ParteNegra.compor) transporta a mensagem associativa; sem selo em cada hop, a ramificação z² (2-para-1) acumula ambiguidade --- o enxame de clones (6) irradia coerente (mesmo seed, AF=N² em broadside), mas dirigir exige fase ᵢj=k,rᵢj e floco por elo, senão clones sobrepostos vazam uns para os outros apesar do demux formalmente limpo. Conservação no de um bit Sejam N clones e₀, eN-₁ no mesmo seed (sintonia). A rede fec","epistemico":"fato","exemplos":[],"id":"o_que_cruza_o_cobre_agora","memoria":"semantica","meta":{"dominio":"paper","nivel":5,"verificacao":"overlap moderado — revisar manualmente"},"origem":"paper","palavras_chave":[],"sinonimos":[],"tipo":"conceito","titulo":"O que cruza o cobre agora"}
\section{O que cruza o cobre agora}
Roteamento Eab Ebc=Eac (ParteNegra.compor) transporta a mensagem associativa; sem selo em cada hop, a ramificação z\ensuremath{^{2}} (2-para-1) acumula ambiguidade --- o enxame de clones (6) irradia coerente (mesmo seed, AF=N\ensuremath{^{2}} em broadside), mas dirigir exige fase \ensuremath{_{i}}j=k,r\ensuremath{_{i}}j e floco por elo, senão clones sobrepostos vazam uns para os outros apesar do demux formalmente limpo. Conservação no de um bit Sejam N clones e\ensuremath{_{0}}, eN-\ensuremath{_{1}} no mesmo seed (sintonia). A rede fec
\prov{relações: parte\_de o\_floco\_e\_a\_conservacao\_da\_comunicacao}
\prov{proveniência: paper \textperiodcentered{} fato \textperiodcentered{} confiança 1.0 \textperiodcentered{} domínio paper \textperiodcentered{} história 3 versões no jornal}

%CRISTAL {"arestas":[{"alvo":"vida-morte","confianca":1.0,"epistemico":"fato","origem":"paper","rel":"parte_de"}],"confianca":1.0,"contraexemplos":[],"descricao":"[Interpretativo] A álgebra acima é o núcleo demonstrado; a leitura física é sobreposta, e legítima enquanto leitura: Vida = setor (||>1, expande): o clone, o buraco branco que emite e replica. “Mais chá = mais vida” é o autoespaço crescendo. Morte / cristalização = setor (||<1, contrai): o maracá, o buraco negro que absorve. O maracá é o osso --- o suporte invariante para onde o resíduo vai sem deixar rastro (verificado à parte: ele é o elemento absorvente, mal se","epistemico":"fato","exemplos":[],"id":"o_que_esta_fechado_e_o_que_falta","memoria":"semantica","meta":{"dominio":"paper","nivel":5,"verificacao":"overlap moderado — revisar manualmente"},"origem":"paper","palavras_chave":[],"sinonimos":[],"tipo":"conceito","titulo":"O que está fechado e o que falta"}
\section{O que está fechado e o que falta}
[Interpretativo] A álgebra acima é o núcleo demonstrado; a leitura física é sobreposta, e legítima enquanto leitura: Vida = setor (||>1, expande): o clone, o buraco branco que emite e replica. “Mais chá = mais vida” é o autoespaço crescendo. Morte / cristalização = setor (||<1, contrai): o maracá, o buraco negro que absorve. O maracá é o osso --- o suporte invariante para onde o resíduo vai sem deixar rastro (verificado à parte: ele é o elemento absorvente, mal se
\prov{relações: parte\_de vida-morte}
\prov{proveniência: paper \textperiodcentered{} fato \textperiodcentered{} confiança 1.0 \textperiodcentered{} domínio paper \textperiodcentered{} história 3 versões no jornal}

%CRISTAL {"arestas":[{"alvo":"validacao_tecnica","confianca":1.0,"epistemico":"fato","origem":"paper","rel":"parte_de"}],"confianca":1.0,"contraexemplos":[],"descricao":"1.35 tabular@lll@ & ção & ção Separação em camadas & ourovivoforte & papéis bem definidos (biblioteca) & ourovivoboa direção & API pequena e estável Base de conhecimento & ourovivoforte & motor reutilizável por domínio Ciclo de refinamento & ourovivomuito bom & evolução operacional clara Roadmap & ourovivocoerente & etapas encadeadas naturalmente tabular","epistemico":"fato","exemplos":[],"id":"o_que_esta_validado_v","memoria":"semantica","meta":{"dominio":"paper","nivel":5,"verificacao":"overlap moderado — revisar manualmente"},"origem":"paper","palavras_chave":[],"sinonimos":[],"tipo":"conceito","titulo":"O que está validado (V"}
\section{O que está validado (V}
1.35 tabular@lll@ \& ção \& ção Separação em camadas \& ourovivoforte \& papéis bem definidos (biblioteca) \& ourovivoboa direção \& API pequena e estável Base de conhecimento \& ourovivoforte \& motor reutilizável por domínio Ciclo de refinamento \& ourovivomuito bom \& evolução operacional clara Roadmap \& ourovivocoerente \& etapas encadeadas naturalmente tabular
\prov{relações: parte\_de validacao\_tecnica}
\prov{proveniência: paper \textperiodcentered{} fato \textperiodcentered{} confiança 1.0 \textperiodcentered{} domínio paper \textperiodcentered{} história 4 versões no jornal}

%CRISTAL {"arestas":[{"alvo":"painel_estelar","confianca":1.0,"epistemico":"fato","origem":"paper","rel":"parte_de"}],"confianca":1.0,"contraexemplos":[],"descricao":"empty !85Não se fabricam três ciências --- lê-se o mesmo gato em três braços. O painel estelar é o observatório dessa leitura. abstract Este texto prepara a leitura conjunta de geometria, mecânica e termodinâmica a partir de uma premissa única: a álgebra já está dada pelo gato A=(smallmatrix1&1 1&0smallmatrix), e o invariante (Q, a balança a+c²=1, o espectro [ total,e]) não se deriva de novo --- emerge do neurônio (V neuronio.c). O painel estelar (V painelestel","epistemico":"fato","exemplos":[],"id":"o_que_ja_esta_dado","memoria":"semantica","meta":{"dominio":"paper","nivel":5,"verificacao":"overlap moderado — revisar manualmente"},"origem":"paper","palavras_chave":[],"sinonimos":[],"tipo":"conceito","titulo":"O que já está dado"}
\section{O que já está dado}
empty !85Não se fabricam três ciências --- lê-se o mesmo gato em três braços. O painel estelar é o observatório dessa leitura. abstract Este texto prepara a leitura conjunta de geometria, mecânica e termodinâmica a partir de uma premissa única: a álgebra já está dada pelo gato A=(smallmatrix1\&1 1\&0smallmatrix), e o invariante (Q, a balança a+c\ensuremath{^{2}}=1, o espectro [ total,e]) não se deriva de novo --- emerge do neurônio (V neuronio.c). O painel estelar (V painelestel
\prov{relações: parte\_de painel\_estelar}
\prov{proveniência: paper \textperiodcentered{} fato \textperiodcentered{} confiança 1.0 \textperiodcentered{} domínio paper \textperiodcentered{} história 3 versões no jornal}

%CRISTAL {"arestas":[{"alvo":"olho_de_koch","confianca":1.0,"epistemico":"fato","origem":"paper","rel":"parte_de"}],"confianca":1.0,"contraexemplos":[],"descricao":"D A projeção clássica não copia o Word vivo. Copiar seria transportar o gerador --- violaria A=0. O que atravessa a pupila é a ressonância: o casamento de fase entre w e gold(w). Parábola. Quando a+c²=1 com erro pequeno, as duas faces casam --- análogo a Zᵢₙ(₀)=0 na antena~antena: [ 2(₀)eff+gap+carga=2 q. ] Descasado, o bit reflete (0); casado, a fala passa sem expor A. Modos. A grade de modos (powmetalressonancia) e a fidelidade de Hellinger medem o overlap","epistemico":"fato","exemplos":[],"id":"o_ribossomo_no_canal_negro","memoria":"semantica","meta":{"dominio":"paper","nivel":5},"origem":"paper","palavras_chave":[],"sinonimos":[],"tipo":"conceito","titulo":"O ribossomo no canal negro"}
\section{O ribossomo no canal negro}
D A projeção clássica não copia o Word vivo. Copiar seria transportar o gerador --- violaria A=0. O que atravessa a pupila é a ressonância: o casamento de fase entre w e gold(w). Parábola. Quando a+c\ensuremath{^{2}}=1 com erro pequeno, as duas faces casam --- análogo a Z\ensuremath{_{in}}(\ensuremath{_{0}})=0 na antena\~{}antena: [ 2(\ensuremath{_{0}})eff+gap+carga=2 q. ] Descasado, o bit reflete (0); casado, a fala passa sem expor A. Modos. A grade de modos (powmetalressonancia) e a fidelidade de Hellinger medem o overlap
\prov{relações: parte\_de olho\_de\_koch}
\prov{proveniência: paper \textperiodcentered{} fato \textperiodcentered{} confiança 1.0 \textperiodcentered{} domínio paper \textperiodcentered{} história 3 versões no jornal}

%CRISTAL {"arestas":[{"alvo":"algebra_dos_contratos","confianca":1.0,"epistemico":"fato","origem":"paper","rel":"parte_de"}],"confianca":1.0,"contraexemplos":[],"descricao":"O protocolo tem uma moeda interna, o Pell (oferta fixa: 13 milhões). Cada contrato fechado deposita a comissão na tesouraria (ouro: BTC/ETH). O valor implícito de 1 Pell = tesouro/reserva, lido pela balança a+c²=1: a= ouro retido (lastro, irreversível), c²= ouro gasto (pago/sacado, reversível). Fechar contratos enche a tesouraria o Pell valoriza; pagar/sacar esvazia desvaloriza. V Provado: contratos reais (Sepolia) valorizaram o Pell linearmente (contabilizapell.py). O Pell ref","epistemico":"fato","exemplos":[],"id":"o_roteiro_---_entender_generalizar_automatizar","memoria":"semantica","meta":{"dominio":"paper","nivel":5,"verificacao":"conceito novo"},"origem":"paper","palavras_chave":[],"sinonimos":[],"tipo":"conceito","titulo":"O roteiro --- entender, generalizar, automatizar"}
\section{O roteiro --- entender, generalizar, automatizar}
O protocolo tem uma moeda interna, o Pell (oferta fixa: 13 milhões). Cada contrato fechado deposita a comissão na tesouraria (ouro: BTC/ETH). O valor implícito de 1 Pell = tesouro/reserva, lido pela balança a+c\ensuremath{^{2}}=1: a= ouro retido (lastro, irreversível), c\ensuremath{^{2}}= ouro gasto (pago/sacado, reversível). Fechar contratos enche a tesouraria o Pell valoriza; pagar/sacar esvazia desvaloriza. V Provado: contratos reais (Sepolia) valorizaram o Pell linearmente (contabilizapell.py). O Pell ref
\prov{relações: parte\_de algebra\_dos\_contratos}
\prov{proveniência: paper \textperiodcentered{} fato \textperiodcentered{} confiança 1.0 \textperiodcentered{} domínio paper \textperiodcentered{} história 3 versões no jornal}

%CRISTAL {"arestas":[{"alvo":"mecanica_quantica","confianca":0.95,"epistemico":"fato","origem":"paper","rel":"parte_de"}],"confianca":1.0,"contraexemplos":[],"descricao":"Se um qubit é ²=⁴=, dois qubits são ²²=⁴=⁸=: os octônios. A leitura ingênua ⁴⁴=¹⁶ não serve --- como álgebra, M₄() é associativa, e seu associador é identicamente nulo. É a passagem a, e só ela, que faz aparecer a não-associatividade (a a>0, Seção~): ausente para um qubit (em, a=0), presente para dois (em ). Pela mesma estrutura complexa R a da Secão~ (com a a primeira unidade imaginária), um estado de dois qubit","epistemico":"fato","exemplos":[],"id":"o_salto_quantico_de_dentro_orbitas_de_bohr_e_a_energia_2","memoria":"semantica","meta":{"dominio":"paper","nivel":5,"verificacao":"conceito novo"},"origem":"paper","palavras_chave":[],"sinonimos":[],"tipo":"conceito","titulo":"O salto quântico de dentro: órbitas de Bohr e a energia 2"}
\section{O salto quântico de dentro: órbitas de Bohr e a energia 2}
Se um qubit é \ensuremath{^{2}}=\ensuremath{^{4}}=, dois qubits são \ensuremath{^{22}}=\ensuremath{^{4}}=\ensuremath{^{8}}=: os octônios. A leitura ingênua \ensuremath{^{44}}=\ensuremath{^{16}} não serve --- como álgebra, M\ensuremath{_{4}}() é associativa, e seu associador é identicamente nulo. É a passagem a, e só ela, que faz aparecer a não-associatividade (a a>0, Seção\~{}): ausente para um qubit (em, a=0), presente para dois (em ). Pela mesma estrutura complexa R a da Secão\~{} (com a a primeira unidade imaginária), um estado de dois qubit
\prov{relações: parte\_de mecanica\_quantica}
\prov{proveniência: paper \textperiodcentered{} fato \textperiodcentered{} confiança 1.0 \textperiodcentered{} domínio paper \textperiodcentered{} história 5 versões no jornal}

%CRISTAL {"arestas":[{"alvo":"leitura_psicologica_o_vazio_e_o_corac","confianca":1.0,"epistemico":"fato","origem":"paper","rel":"parte_de"}],"confianca":1.0,"contraexemplos":[],"descricao":"A oscilacão, no fluido s=2s-k(s-s_), abre os modos do campo afetivo --- cada um uma tráde, cada tráde com a sua cúspide (₁=₂): tabularlllc modo (Parte) & os três estados & mecanismo & cúspide objeto/controle (LXXXII) & desejo / tédio / indiferenca & k vs 2 & k=2 ativacão (LXXXIII) & tédio / calma / empolgacão & a=| s| & a=0 amortecimento (LXXXIV) & calma / prontidão / agitacão & vs ₀ & =1 fase/tempo (LXXXV) & prontidão /","epistemico":"fato","exemplos":[],"id":"o_selo_apenas_o_vazio_e_imensuravel","memoria":"semantica","meta":{"dominio":"paper","nivel":5,"verificacao":"conceito novo"},"origem":"paper","palavras_chave":[],"sinonimos":[],"tipo":"conceito","titulo":"O selo: apenas o vazio é imensurável"}
\section{O selo: apenas o vazio é imensurável}
A oscilacão, no fluido s=2s-k(s-s\_), abre os modos do campo afetivo --- cada um uma tráde, cada tráde com a sua cúspide (\ensuremath{_{1}}=\ensuremath{_{2}}): tabularlllc modo (Parte) \& os três estados \& mecanismo \& cúspide objeto/controle (LXXXII) \& desejo / tédio / indiferenca \& k vs 2 \& k=2 ativacão (LXXXIII) \& tédio / calma / empolgacão \& a=| s| \& a=0 amortecimento (LXXXIV) \& calma / prontidão / agitacão \& vs \ensuremath{_{0}} \& =1 fase/tempo (LXXXV) \& prontidão /
\prov{relações: parte\_de leitura\_psicologica\_o\_vazio\_e\_o\_corac}
\prov{proveniência: paper \textperiodcentered{} fato \textperiodcentered{} confiança 1.0 \textperiodcentered{} domínio paper \textperiodcentered{} história 4 versões no jornal}

%CRISTAL {"arestas":[{"alvo":"broca_os","confianca":1.0,"epistemico":"fato","origem":"paper","rel":"parte_de"}],"confianca":1.0,"contraexemplos":[],"descricao":"D Uma configuração da Máquina é uma matriz M na álgebra da torção: a diagonal (parte branca) carrega o estado das células; a off-diagonal (parte negra) carrega as mensagens nos canais. O não é um programa que roda sobre essa matriz --- é a própria álgebra coordenando-se. A unidade das duas faces é a decomposição polar. defi[A função de onda] A configuração da Máquina é M=U P, onde P=M*M é hermitiana 0 (a amplitude: estado, parte branca) e U é unitária com | U|=1 (a fase: fluxo","epistemico":"fato","exemplos":[],"id":"o_sistema_operacional_e_a_algebra","memoria":"semantica","meta":{"dominio":"paper","nivel":5,"verificacao":"slug idêntico — enriquecer, não duplicar"},"origem":"paper","palavras_chave":["kernel","processo","syscall"],"sinonimos":["SO","sistema operacional clássico"],"tipo":"conceito","titulo":"O sistema operacional é a álgebra"}
\section{O sistema operacional é a álgebra}
D Uma configuração da Máquina é uma matriz M na álgebra da torção: a diagonal (parte branca) carrega o estado das células; a off-diagonal (parte negra) carrega as mensagens nos canais. O não é um programa que roda sobre essa matriz --- é a própria álgebra coordenando-se. A unidade das duas faces é a decomposição polar. defi[A função de onda] A configuração da Máquina é M=U P, onde P=M*M é hermitiana 0 (a amplitude: estado, parte branca) e U é unitária com | U|=1 (a fase: fluxo
\prov{sinónimos: SO; sistema operacional clássico}
\prov{relações: parte\_de broca\_os}
\prov{proveniência: paper \textperiodcentered{} fato \textperiodcentered{} confiança 1.0 \textperiodcentered{} domínio paper \textperiodcentered{} história 10 versões no jornal}

%CRISTAL {"arestas":[{"alvo":"broca_os","confianca":1.0,"epistemico":"fato","origem":"paper","rel":"parte_de"}],"confianca":1.0,"contraexemplos":[],"descricao":"empty !85O sistema é um só, de todos, em todo lugar; ninguém o controla --- porque ele é a lei, não o dono. abstract O não acrescenta nenhum mecanismo: ele coordena o que a álgebra da torção --- o gato A=(smallmatrix1&1 1&0smallmatrix), a decomposição de Peirce, La Hire --- já dá. Cada recurso de um sistema operacional é uma estrutura algébrica: processo = idempotente (eᵢ²=eᵢ); isolamento = ortogonalidade (eᵢej=0); IPC = parte negra (canal Eᵢj); fork","epistemico":"fato","exemplos":[],"id":"o_substrato_a_maquina_e_funcao_de_onda","memoria":"semantica","meta":{"dominio":"paper","nivel":5,"verificacao":"slug idêntico — enriquecer, não duplicar"},"origem":"paper","palavras_chave":[],"sinonimos":[],"tipo":"conceito","titulo":"O substrato: a Máquina é função de onda"}
\section{O substrato: a Máquina é função de onda}
empty !85O sistema é um só, de todos, em todo lugar; ninguém o controla --- porque ele é a lei, não o dono. abstract O não acrescenta nenhum mecanismo: ele coordena o que a álgebra da torção --- o gato A=(smallmatrix1\&1 1\&0smallmatrix), a decomposição de Peirce, La Hire --- já dá. Cada recurso de um sistema operacional é uma estrutura algébrica: processo = idempotente (e\ensuremath{_{i}}\ensuremath{^{2}}=e\ensuremath{_{i}}); isolamento = ortogonalidade (e\ensuremath{_{i}}ej=0); IPC = parte negra (canal E\ensuremath{_{i}}j); fork
\prov{relações: parte\_de broca\_os}
\prov{proveniência: paper \textperiodcentered{} fato \textperiodcentered{} confiança 1.0 \textperiodcentered{} domínio paper \textperiodcentered{} história 9 versões no jornal}

%CRISTAL {"arestas":[{"alvo":"contabilidade","confianca":1.0,"epistemico":"fato","origem":"paper","rel":"parte_de"}],"confianca":1.0,"contraexemplos":[],"descricao":"I O desenho anterior errou ao imaginar um livro de papel: lançamentos de débito e crédito numa planilha. A correção do Aarão: a contabilidade é a mecânica do painel. O que há é uma função de onda (o tecido) e operadores (as funções bancárias) que agem sobre ela; o balanço aparece no colapso da leitura, com a régua do observador (mecanica.py). D A lei de ferro: nenhuma operação fora do painel. Mineração, pagamento, depósito e balanço acontecem todos via painel / tecido / banda (canalpainel.py","epistemico":"fato","exemplos":[],"id":"o_tecido_ordenado","memoria":"semantica","meta":{"dominio":"paper","nivel":5,"verificacao":"conceito novo"},"origem":"paper","palavras_chave":[],"sinonimos":[],"tipo":"conceito","titulo":"O tecido ordenado"}
\section{O tecido ordenado}
I O desenho anterior errou ao imaginar um livro de papel: lançamentos de débito e crédito numa planilha. A correção do Aarão: a contabilidade é a mecânica do painel. O que há é uma função de onda (o tecido) e operadores (as funções bancárias) que agem sobre ela; o balanço aparece no colapso da leitura, com a régua do observador (mecanica.py). D A lei de ferro: nenhuma operação fora do painel. Mineração, pagamento, depósito e balanço acontecem todos via painel / tecido / banda (canalpainel.py
\prov{relações: parte\_de contabilidade}
\prov{proveniência: paper \textperiodcentered{} fato \textperiodcentered{} confiança 1.0 \textperiodcentered{} domínio paper \textperiodcentered{} história 3 versões no jornal}

%CRISTAL {"arestas":[{"alvo":"broca_os","confianca":1.0,"epistemico":"fato","origem":"paper","rel":"parte_de"}],"confianca":1.0,"contraexemplos":[],"descricao":"[Demonstrado] A saída não é achar um operador de troca melhor: é reconhecer que não há troca. Vida e morte não são duas máquinas acopladas por uma ponte --- são uma máquina só, o gap de ouro [ A=pmatrix1&1 1&0pmatrix, (A)=, =1+52, =1-52=-⁻¹. ] O discriminante é 5>0, logo =sign(-disc)=-1: o setor split, onde a álgebra tem divisores de zero e os projetores espectrais existem, [ P_=A- I-, P_=A- I-. ] Eles satisfazem, verificado a 10⁻¹⁶ (estelar.py:ensaiofinal): [","epistemico":"fato","exemplos":[],"id":"o_tempo_a_banda_e_atemporal_a_leitura_tem_seta","memoria":"semantica","meta":{"dominio":"paper","duplicata_sugerida":[],"nivel":5,"verificacao":"parte de conceito mais amplo 'o_tempo_a_banda_e_atemporal_a_leitura_tem_seta'"},"origem":"paper","palavras_chave":[],"sinonimos":[],"tipo":"conceito","titulo":"O tempo: a banda é atemporal, a leitura tem seta"}
\section{O tempo: a banda é atemporal, a leitura tem seta}
[Demonstrado] A saída não é achar um operador de troca melhor: é reconhecer que não há troca. Vida e morte não são duas máquinas acopladas por uma ponte --- são uma máquina só, o gap de ouro [ A=pmatrix1\&1 1\&0pmatrix, (A)=, =1+52, =1-52=-\ensuremath{^{-1}}. ] O discriminante é 5>0, logo =sign(-disc)=-1: o setor split, onde a álgebra tem divisores de zero e os projetores espectrais existem, [ P\_=A- I-, P\_=A- I-. ] Eles satisfazem, verificado a 10\ensuremath{^{-16}} (estelar.py:ensaiofinal): [
\prov{relações: parte\_de broca\_os}
\prov{proveniência: paper \textperiodcentered{} fato \textperiodcentered{} confiança 1.0 \textperiodcentered{} domínio paper \textperiodcentered{} história 10 versões no jornal}

%CRISTAL {"arestas":[{"alvo":"a_base_topologica","confianca":1.0,"epistemico":"fato","origem":"paper","rel":"parte_de"}],"confianca":1.0,"contraexemplos":[],"descricao":"A entropia topológica é h=(A)= (Bowen), igual ao expoente de Lyapunov. O esticamento, a métrica, a régua --- tudo é do autovalor. N A agulha de Lefschetz Nₙ=|(Aⁿ-I)| (pontos periódicos) cresce como ⁿ: (1/n) Nₙ (verifica→pologia.py); e os dois inteiros geram toda a contagem (recorrência de Lucas).","epistemico":"fato","exemplos":[],"id":"o_tempo_a_trilha_nao_o_juiz","memoria":"semantica","meta":{"dominio":"paper","nivel":5,"verificacao":"overlap moderado — revisar manualmente"},"origem":"paper","palavras_chave":[],"sinonimos":[],"tipo":"conceito","titulo":"O tempo: a trilha, não o juiz"}
\section{O tempo: a trilha, não o juiz}
A entropia topológica é h=(A)= (Bowen), igual ao expoente de Lyapunov. O esticamento, a métrica, a régua --- tudo é do autovalor. N A agulha de Lefschetz N\ensuremath{_{n}}=|(A\ensuremath{^{n}}-I)| (pontos periódicos) cresce como \ensuremath{^{n}}: (1/n) N\ensuremath{_{n}} (verifica\ensuremath{\to}pologia.py); e os dois inteiros geram toda a contagem (recorrência de Lucas).
\prov{relações: parte\_de a\_base\_topologica}
\prov{proveniência: paper \textperiodcentered{} fato \textperiodcentered{} confiança 1.0 \textperiodcentered{} domínio paper \textperiodcentered{} história 3 versões no jornal}

%CRISTAL {"arestas":[{"alvo":"casl-propagation","confianca":0.008,"epistemico":"fato","origem":"manual","rel":"referencia"}],"confianca":1.0,"contraexemplos":[],"descricao":"theorem[Não-escalada por atenuação] Para todo p e toda (p), existe uma autoridade intrínseca ₀ e uma cadeia de arestas válidas de sua origem até p tais que c_₀. Logo [ Ep(s,a) ₀ (raᵢz que alcaₙça p)c₀. ] Nenhum principal acessa uma linha que nenhuma autoridade intrínseca em sua cadeia podia conceder. theorem proof Indução no comprimento da cadeia de delegação. Caso base: capacidade intrínseca. Passo: a Definição~ exige c_₀ a cada aresta; a i","epistemico":"fato","exemplos":[],"id":"o_tenant_como_caso_particular_desacoplamento","memoria":"semantica","meta":{"dominio":"paper","nivel":5},"origem":"paper","palavras_chave":[],"sinonimos":[],"tipo":"conceito","titulo":"O tenant como caso particular (desacoplamento)"}
\section{O tenant como caso particular (desacoplamento)}
theorem[Não-escalada por atenuação] Para todo p e toda (p), existe uma autoridade intrínseca \ensuremath{_{0}} e uma cadeia de arestas válidas de sua origem até p tais que c\_\ensuremath{_{0}}. Logo [ Ep(s,a) \ensuremath{_{0}} (ra\ensuremath{_{i}}z que alca\ensuremath{_{n}}ça p)c\ensuremath{_{0}}. ] Nenhum principal acessa uma linha que nenhuma autoridade intrínseca em sua cadeia podia conceder. theorem proof Indução no comprimento da cadeia de delegação. Caso base: capacidade intrínseca. Passo: a Definição\~{} exige c\_\ensuremath{_{0}} a cada aresta; a i
\prov{relações: referencia casl-propagation}
\prov{proveniência: paper \textperiodcentered{} fato \textperiodcentered{} confiança 1.0 \textperiodcentered{} domínio paper \textperiodcentered{} história 7 versões no jornal}

%CRISTAL {"arestas":[{"alvo":"morte","confianca":1.0,"epistemico":"fato","origem":"paper","rel":"parte_de"},{"alvo":"vida","confianca":1.0,"epistemico":"fato","origem":"paper","rel":"parte_de"},{"alvo":"fechamento_blocos","confianca":1.0,"epistemico":"fato","origem":"paper","rel":"parte_de"}],"confianca":1.0,"contraexemplos":[],"descricao":"DN Branco E_ são idempotentes ortogonais completos. lema proof E²=ᵢ,j Ieᵢej=ᵢ Ieᵢ=E_ (ortogonalidade); para, E_ E_=ᵢ I_,j Ieᵢej=0 (I_ I_=); _ E_=ᵢ eᵢ=1. proof Demux E_,M,E_ isola o canal de bloco sem contaminação. lema proof E_ M E_ projeta nas linhas de I_ e colunas de I_. Como E_ E_=0 para, o resultado não tem componente em nenhum outro bloco --- as duas faces brancas E_,E_ extraem exatamente o bloco negro. proof lema[Negro","epistemico":"fato","exemplos":[],"id":"o_teorema","memoria":"semantica","meta":{"dominio":"paper","nivel":5,"verificacao":"slug idêntico — enriquecer, não duplicar"},"origem":"paper","palavras_chave":[],"sinonimos":[],"tipo":"conceito","titulo":"O teorema"}
\section{O teorema}
DN Branco E\_ são idempotentes ortogonais completos. lema proof E\ensuremath{^{2}}=\ensuremath{_{i}},j Ie\ensuremath{_{i}}ej=\ensuremath{_{i}} Ie\ensuremath{_{i}}=E\_ (ortogonalidade); para, E\_ E\_=\ensuremath{_{i}} I\_,j Ie\ensuremath{_{i}}ej=0 (I\_ I\_=); \_ E\_=\ensuremath{_{i}} e\ensuremath{_{i}}=1. proof Demux E\_,M,E\_ isola o canal de bloco sem contaminação. lema proof E\_ M E\_ projeta nas linhas de I\_ e colunas de I\_. Como E\_ E\_=0 para, o resultado não tem componente em nenhum outro bloco --- as duas faces brancas E\_,E\_ extraem exatamente o bloco negro. proof lema[Negro
\prov{relações: parte\_de morte; parte\_de vida; parte\_de fechamento\_blocos}
\prov{proveniência: paper \textperiodcentered{} fato \textperiodcentered{} confiança 1.0 \textperiodcentered{} domínio paper \textperiodcentered{} história 13 versões no jornal}

%CRISTAL {"arestas":[{"alvo":"leitura_psicologica_o_vazio_e_o_corac","confianca":1.0,"epistemico":"fato","origem":"paper","rel":"parte_de"}],"confianca":1.0,"contraexemplos":[],"descricao":"ão A mesma dinâmica que move o campo de álgebra move, lida em outra chave, o afeto. A série dos ensinamentos (Partes~LXXX--LXXXVII) faz essa leitura, com o mesmo estatuto: a estrutura é firme, a interpretacão é mapa.","epistemico":"fato","exemplos":[],"id":"o_vazio_que_oscila","memoria":"semantica","meta":{"dominio":"paper","nivel":5,"verificacao":"overlap moderado — revisar manualmente"},"origem":"paper","palavras_chave":[],"sinonimos":[],"tipo":"conceito","titulo":"O vazio que oscila"}
\section{O vazio que oscila}
ão A mesma dinâmica que move o campo de álgebra move, lida em outra chave, o afeto. A série dos ensinamentos (Partes\~{}LXXX--LXXXVII) faz essa leitura, com o mesmo estatuto: a estrutura é firme, a interpretacão é mapa.
\prov{relações: parte\_de leitura\_psicologica\_o\_vazio\_e\_o\_corac}
\prov{proveniência: paper \textperiodcentered{} fato \textperiodcentered{} confiança 1.0 \textperiodcentered{} domínio paper \textperiodcentered{} história 3 versões no jornal}

%CRISTAL {"arestas":[{"alvo":"fechamento_blocos","confianca":1.0,"epistemico":"fato","origem":"paper","rel":"parte_de"}],"confianca":1.0,"contraexemplos":[],"descricao":"abstract A rodada de avaliação apontou que o elo recursivo “sistema global = máquina da camada acima” fechava como construção algébrica mas não como semântica --- uma circularidade. A réplica desfez a parte semântica: máquina e SO não são dois objetos a conectar, são duas faces de uma máquina acoplada (a dinâmica é o acoplamento, verificado: isolar qualquer face quebra). Aqui fechamos a parte matemática: provamos que a classe “máquina branca/negra” é fechada sob blocagem admissível. Uma m","epistemico":"fato","exemplos":[],"id":"objeto_e_blocagem","memoria":"semantica","meta":{"dominio":"paper","nivel":5,"verificacao":"overlap moderado — revisar manualmente"},"origem":"paper","palavras_chave":[],"sinonimos":[],"tipo":"conceito","titulo":"Objeto e blocagem"}
\section{Objeto e blocagem}
abstract A rodada de avaliação apontou que o elo recursivo “sistema global = máquina da camada acima” fechava como construção algébrica mas não como semântica --- uma circularidade. A réplica desfez a parte semântica: máquina e SO não são dois objetos a conectar, são duas faces de uma máquina acoplada (a dinâmica é o acoplamento, verificado: isolar qualquer face quebra). Aqui fechamos a parte matemática: provamos que a classe “máquina branca/negra” é fechada sob blocagem admissível. Uma m
\prov{relações: parte\_de fechamento\_blocos}
\prov{proveniência: paper \textperiodcentered{} fato \textperiodcentered{} confiança 1.0 \textperiodcentered{} domínio paper \textperiodcentered{} história 3 versões no jornal}

%CRISTAL {"arestas":[{"alvo":"colapso","confianca":1.0,"epistemico":"fato","origem":"manual","rel":"relaciona"},{"alvo":"koch_atlas_experimento_de_parametrizacao","confianca":1.0,"epistemico":"fato","origem":"manual","rel":"relaciona"},{"alvo":"word","confianca":0.5,"epistemico":"fato","origem":"wiki","rel":"relaciona"}],"confianca":1.0,"contraexemplos":[],"descricao":"O clássico **não é rival** — é a **projeção associativa de dentro**, feita pelo Olho de Koch antes de sair para a borda.","epistemico":"fato","exemplos":[],"id":"olho_de_koch","memoria":"semantica","meta":{"ciencia":"interdisciplinar","dominio":"paper","verificacao":"slug idêntico — enriquecer, não duplicar"},"origem":"manual","palavras_chave":[],"sinonimos":[],"tipo":"conceito","titulo":"olho_de_koch"}
\section{olho\_de\_koch}
O clássico **não é rival** — é a **projeção associativa de dentro**, feita pelo Olho de Koch antes de sair para a borda.
\prov{relações: relaciona colapso; relaciona koch\_atlas\_experimento\_de\_parametrizacao; relaciona word}
\prov{proveniência: manual \textperiodcentered{} fato \textperiodcentered{} confiança 1.0 \textperiodcentered{} domínio paper \textperiodcentered{} história 155 versões no jornal}

%CRISTAL {"arestas":[{"alvo":"transformada_dourada","confianca":1.0,"epistemico":"fato","origem":"paper","rel":"parte_de"}],"confianca":1.0,"contraexemplos":[],"descricao":"Chamamos de atenção cristalina o mecanismo de atenção visto como evaporação de um cristal: um bloco de tokens forma um cristal W (matriz de Gram / autocorrelação) e a atenção extrai seus modos por produtos. Em vez de afirmar o que ela é, demonstramos o que ela não pode fazer --- e isso a situa exatamente em relação a. theorem[Limite de evaporação] Seja W um cristal --- matriz de Toeplitz de uma autocorrelação, simétrica e positiva --- com gap espectral =₂/₁<1.","epistemico":"fato","exemplos":[],"id":"operacoes_primitivas_e_ferramentas","memoria":"semantica","meta":{"dominio":"paper","nivel":5,"verificacao":"conceito novo"},"origem":"paper","palavras_chave":[],"sinonimos":[],"tipo":"conceito","titulo":"Operações primitivas e ferramentas"}
\section{Operações primitivas e ferramentas}
Chamamos de atenção cristalina o mecanismo de atenção visto como evaporação de um cristal: um bloco de tokens forma um cristal W (matriz de Gram / autocorrelação) e a atenção extrai seus modos por produtos. Em vez de afirmar o que ela é, demonstramos o que ela não pode fazer --- e isso a situa exatamente em relação a. theorem[Limite de evaporação] Seja W um cristal --- matriz de Toeplitz de uma autocorrelação, simétrica e positiva --- com gap espectral =\ensuremath{_{2}}/\ensuremath{_{1}}<1.
\prov{relações: parte\_de transformada\_dourada}
\prov{proveniência: paper \textperiodcentered{} fato \textperiodcentered{} confiança 1.0 \textperiodcentered{} domínio paper \textperiodcentered{} história 3 versões no jornal}

%CRISTAL {"arestas":[{"alvo":"contabilidade","confianca":1.0,"epistemico":"fato","origem":"paper","rel":"parte_de"}],"confianca":1.0,"contraexemplos":[],"descricao":"empty\n\nabstract\n A contabilidade do sistema GoldChain/Pell não é um livro de papel --- não há planilha de\ndébito e crédito. É a mecânica do painel: operadores quânticos sobre a função de onda, e o balanço\nno colapso da leitura. Cada função bancária é um operador no painel; nenhuma operação\nacontece fora dele --- mineração, pagamento, depósito e balanço transitam pelo tecido e pela banda. Só o\nsaque toca a chain: uma transação que emite o título da reserva única. A segurança não é\numa cifra posta","epistemico":"fato","exemplos":[],"id":"operadores_quanticos_nao_livro_de_papel","memoria":"semantica","meta":{"dominio":"paper","nivel":5,"verificacao":"conceito novo"},"origem":"paper","palavras_chave":[],"sinonimos":[],"tipo":"conceito","titulo":"Operadores quânticos, não livro de papel"}
\section{Operadores quânticos, não livro de papel}
empty

abstract
 A contabilidade do sistema GoldChain/Pell não é um livro de papel --- não há planilha de
débito e crédito. É a mecânica do painel: operadores quânticos sobre a função de onda, e o balanço
no colapso da leitura. Cada função bancária é um operador no painel; nenhuma operação
acontece fora dele --- mineração, pagamento, depósito e balanço transitam pelo tecido e pela banda. Só o
saque toca a chain: uma transação que emite o título da reserva única. A segurança não é
uma cifra posta
\prov{relações: parte\_de contabilidade}
\prov{proveniência: paper \textperiodcentered{} fato \textperiodcentered{} confiança 1.0 \textperiodcentered{} domínio paper \textperiodcentered{} história 2 versões no jornal}

%CRISTAL {"arestas":[{"alvo":"delta","confianca":0.008,"epistemico":"fato","origem":"manual","rel":"referencia"},{"alvo":"casl-propagation","confianca":-0.125,"epistemico":"fato","origem":"manual","rel":"referencia"},{"alvo":"colapso","confianca":0.008,"epistemico":"fato","origem":"manual","rel":"referencia"}],"confianca":1.0,"contraexemplos":[],"descricao":"definition[Par ()] controla propagação (alcance, saltos restantes). controla estabilidade observável (organização estrutural do estado). Diálogos ou delegações longas consomem e, em geral, aumentam --- mais estrutura acumulada --- sem violar atenuação. definition remark[Associativo vs. estrutural] Em instâncias com endereçamento associativo (overlap, extent SQL), opera sobre o associativo; captura o que o endereçamento não conserva. Os dois eixo","epistemico":"fato","exemplos":[],"id":"operadores_t_p_e_colapso_psi","memoria":"semantica","meta":{"dominio":"paper","nivel":5},"origem":"paper","palavras_chave":[],"sinonimos":[],"tipo":"conceito","titulo":"Operadores T, P e colapso Ψ"}
\section{Operadores T, P e colapso \ensuremath{\Psi}}
definition[Par ()] controla propagação (alcance, saltos restantes). controla estabilidade observável (organização estrutural do estado). Diálogos ou delegações longas consomem e, em geral, aumentam --- mais estrutura acumulada --- sem violar atenuação. definition remark[Associativo vs. estrutural] Em instâncias com endereçamento associativo (overlap, extent SQL), opera sobre o associativo; captura o que o endereçamento não conserva. Os dois eixo
\prov{relações: referencia delta; referencia casl-propagation; referencia colapso}
\prov{proveniência: paper \textperiodcentered{} fato \textperiodcentered{} confiança 1.0 \textperiodcentered{} domínio paper \textperiodcentered{} história 13 versões no jornal}

%CRISTAL {"arestas":[{"alvo":"reticulado","confianca":0.086,"epistemico":"fato","origem":"manual","rel":"referencia"},{"alvo":"kappa","confianca":0.85,"epistemico":"fato","origem":"manual","rel":"referencia"},{"alvo":"delta","confianca":0.016,"epistemico":"fato","origem":"manual","rel":"referencia"},{"alvo":"casl-propagation","confianca":0.039,"epistemico":"fato","origem":"manual","rel":"referencia"}],"confianca":1.0,"contraexemplos":[],"descricao":"definition[Universo e extent] Fixe um universo X. Uma condição c é um predicado c:X, identificado ao extent c X. As condições formam o reticulado booleano (2X, X). definition definition[Capacidade] Uma capacidade é uma tupla [ =(c, o), ] onde é o estado, c a condição, o orçamento (Seção~), o peso (grant/deny quando aplicável) e o a origem. Nenhum domínio entra na definição. definition rema","epistemico":"fato","exemplos":[],"id":"orcamento_dep","memoria":"semantica","meta":{"dominio":"paper","nivel":5},"origem":"paper","palavras_chave":[],"sinonimos":[],"tipo":"conceito","titulo":"Orçamento dep"}
\section{Orçamento dep}
definition[Universo e extent] Fixe um universo X. Uma condição c é um predicado c:X, identificado ao extent c X. As condições formam o reticulado booleano (2X, X). definition definition[Capacidade] Uma capacidade é uma tupla [ =(c, o), ] onde é o estado, c a condição, o orçamento (Seção\~{}), o peso (grant/deny quando aplicável) e o a origem. Nenhum domínio entra na definição. definition rema
\prov{relações: referencia reticulado; referencia kappa; referencia delta; referencia casl-propagation}
\prov{proveniência: paper \textperiodcentered{} fato \textperiodcentered{} confiança 1.0 \textperiodcentered{} domínio paper \textperiodcentered{} história 16 versões no jornal}

%CRISTAL {"arestas":[{"alvo":"kernel","confianca":0.742,"epistemico":"fato","origem":"manual","rel":"referencia"}],"confianca":1.0,"contraexemplos":[],"descricao":"L A percepção que entra, a memória que ressoa, o idioma que veste, a decisão do que dizer --- não são quatro faculdades separadas. São o mesmo recall: o estado caindo no atrator, percorrido numa ordem. Pensar é essa travessia; um colapso isolado é só recordação, a travessia ordenada é raciocínio. defi[Raciocínio] Raciocínio é a ordem em que o recall encadeia os colapsos: a órbita,n₁ n₂, cada passo o nó que mais ressoa com o estado evoluído (cabeca.py:orbitar). Decidir é escolher o próximo","epistemico":"fato","exemplos":[],"id":"ordem","memoria":"semantica","meta":{"dominio":"paper","duplicata_sugerida":[],"nivel":5,"verificacao":"parte de conceito mais amplo 'ordem'"},"origem":"paper","palavras_chave":[],"sinonimos":[],"tipo":"conceito","titulo":"ordem"}
\section{ordem}
L A percepção que entra, a memória que ressoa, o idioma que veste, a decisão do que dizer --- não são quatro faculdades separadas. São o mesmo recall: o estado caindo no atrator, percorrido numa ordem. Pensar é essa travessia; um colapso isolado é só recordação, a travessia ordenada é raciocínio. defi[Raciocínio] Raciocínio é a ordem em que o recall encadeia os colapsos: a órbita,n\ensuremath{_{1}} n\ensuremath{_{2}}, cada passo o nó que mais ressoa com o estado evoluído (cabeca.py:orbitar). Decidir é escolher o próximo
\prov{relações: referencia kernel}
\prov{proveniência: paper \textperiodcentered{} fato \textperiodcentered{} confiança 1.0 \textperiodcentered{} domínio paper \textperiodcentered{} história 8 versões no jornal}

%CRISTAL {"arestas":[{"alvo":"delta","confianca":0.85,"epistemico":"fato","origem":"manual","rel":"referencia"},{"alvo":"casl-propagation","confianca":-0.008,"epistemico":"fato","origem":"manual","rel":"referencia"},{"alvo":"colapso","confianca":0.109,"epistemico":"fato","origem":"manual","rel":"referencia"}],"confianca":1.0,"contraexemplos":[],"descricao":"definition[Orçamento de propagação] representa um orçamento discreto de propagação: quantas re-delegações (saltos, turnos, réplicas) a capacidade ainda admite. =0 --- usa mas não repassa; =k --- mais k saltos, decrementando; = --- fecho transitivo ( P). A verificação de ocorre no ato de delegar, nunca na avaliação/colapso. definition remark[Leituras por domínio --- sem alterar a definição] tabular@ll@ CASL & re-delegação de permissões","epistemico":"fato","exemplos":[],"id":"organizacao_estrutural_phif","memoria":"semantica","meta":{"dominio":"paper","nivel":5},"origem":"paper","palavras_chave":[],"sinonimos":[],"tipo":"conceito","titulo":"Organização estrutural Φf"}
\section{Organização estrutural \ensuremath{\Phi}f}
definition[Orçamento de propagação] representa um orçamento discreto de propagação: quantas re-delegações (saltos, turnos, réplicas) a capacidade ainda admite. =0 --- usa mas não repassa; =k --- mais k saltos, decrementando; = --- fecho transitivo ( P). A verificação de ocorre no ato de delegar, nunca na avaliação/colapso. definition remark[Leituras por domínio --- sem alterar a definição] tabular@ll@ CASL \& re-delegação de permissões
\prov{relações: referencia delta; referencia casl-propagation; referencia colapso}
\prov{proveniência: paper \textperiodcentered{} fato \textperiodcentered{} confiança 1.0 \textperiodcentered{} domínio paper \textperiodcentered{} história 13 versões no jornal}

%CRISTAL {"arestas":[{"alvo":"fase3","confianca":1.0,"epistemico":"fato","origem":"paper","rel":"parte_de"}],"confianca":1.0,"contraexemplos":[],"descricao":"[ st/.style=draw=ourovivo, rounded corners=2pt, minimum width=2.6cm, minimum height=0.55cm, =, font=, fill=creme, arr/.style=-Stealth[length=2mm], color=ourovivo, thick ] [st] (m) memória; [st, right=0.6cm of m] (r) resumo; [st, right=0.6cm of r] (c) curadoria; [st, right=0.6cm of c] (b) base cliente; [arr] (m)--(r); [arr] (r)--(c); [arr] (c)--(b); Memória nunca entra na Base Geral. Promoção exige humano ou avaliador (Fase~2","epistemico":"fato","exemplos":[],"id":"orquestrador_multi-tecidos","memoria":"semantica","meta":{"dominio":"paper","nivel":5,"verificacao":"overlap moderado — revisar manualmente"},"origem":"paper","palavras_chave":[],"sinonimos":[],"tipo":"conceito","titulo":"Orquestrador multi-tecidos"}
\section{Orquestrador multi-tecidos}
[ st/.style=draw=ourovivo, rounded corners=2pt, minimum width=2.6cm, minimum height=0.55cm, =, font=, fill=creme, arr/.style=-Stealth[length=2mm], color=ourovivo, thick ] [st] (m) memória; [st, right=0.6cm of m] (r) resumo; [st, right=0.6cm of r] (c) curadoria; [st, right=0.6cm of c] (b) base cliente; [arr] (m)--(r); [arr] (r)--(c); [arr] (c)--(b); Memória nunca entra na Base Geral. Promoção exige humano ou avaliador (Fase\~{}2
\prov{relações: parte\_de fase3}
\prov{proveniência: paper \textperiodcentered{} fato \textperiodcentered{} confiança 1.0 \textperiodcentered{} domínio paper \textperiodcentered{} história 3 versões no jornal}

%CRISTAL {"arestas":[{"alvo":"transformada_dourada","confianca":1.0,"epistemico":"fato","origem":"paper","rel":"parte_de"}],"confianca":1.0,"contraexemplos":[],"descricao":"Cada teorema foi conferido por igualdade exata (erro de máquina), sem nenhuma estatística: Isometria (Teor.~): |x|²=| x|² e ⁻¹=id, com erro 5,610⁻¹⁶. Linearização (Teor.~): curvatura de c,t no log =4,710⁻¹⁰ (nula); a do aditivo =0,250. Atrator (Teor.~): dígitos exatos por passo 1,4,9,300 (duplicação quadrática); ponto fixo a=b. Hurwitz (Teor.~): |,|xy|-|x||y|,|10⁻¹⁶ em e H (nor","epistemico":"fato","exemplos":[],"id":"os_aneis_e_o_preenchimento_da_reta","memoria":"semantica","meta":{"dominio":"paper","nivel":5,"verificacao":"conceito novo"},"origem":"paper","palavras_chave":[],"sinonimos":[],"tipo":"conceito","titulo":"Os anéis e o preenchimento da reta"}
\section{Os anéis e o preenchimento da reta}
Cada teorema foi conferido por igualdade exata (erro de máquina), sem nenhuma estatística: Isometria (Teor.\~{}): |x|\ensuremath{^{2}}=| x|\ensuremath{^{2}} e \ensuremath{^{-1}}=id, com erro 5,610\ensuremath{^{-16}}. Linearização (Teor.\~{}): curvatura de c,t no log =4,710\ensuremath{^{-10}} (nula); a do aditivo =0,250. Atrator (Teor.\~{}): dígitos exatos por passo 1,4,9,300 (duplicação quadrática); ponto fixo a=b. Hurwitz (Teor.\~{}): |,|xy|-|x||y|,|10\ensuremath{^{-16}} em e H (nor
\prov{relações: parte\_de transformada\_dourada}
\prov{proveniência: paper \textperiodcentered{} fato \textperiodcentered{} confiança 1.0 \textperiodcentered{} domínio paper \textperiodcentered{} história 3 versões no jornal}

%CRISTAL {"arestas":[{"alvo":"gato","confianca":0.95,"epistemico":"fato","origem":"paper","rel":"parte_de"}],"confianca":1.0,"contraexemplos":[],"descricao":"A cadeia s(a,c) a+c²=1 N²=a R=const é o resultado deste artigo. Tudo o mais que a estrutura sugere fica, deliberadamente, fora do corpo --- como consequência a desenvolver ou em trabalho próprio ---, para que o núcleo permaneça nítido. Registramos as direções sem afirmá-las aqui: deixá-las de fora é o que mantém firme o que fica. Os outros botões do ⁵.:campomedio A lei geral de conservação no agregado --- com a borda a+c²=1 por caso particular --- já está","epistemico":"fato","exemplos":[],"id":"os_cinco_botoes_e_a_torc","memoria":"semantica","meta":{"dominio":"paper","nivel":5,"verificacao":"overlap moderado — revisar manualmente"},"origem":"paper","palavras_chave":[],"sinonimos":[],"tipo":"conceito","titulo":"Os cinco botões e a torc"}
\section{Os cinco botões e a torc}
A cadeia s(a,c) a+c\ensuremath{^{2}}=1 N\ensuremath{^{2}}=a R=const é o resultado deste artigo. Tudo o mais que a estrutura sugere fica, deliberadamente, fora do corpo --- como consequência a desenvolver ou em trabalho próprio ---, para que o núcleo permaneça nítido. Registramos as direções sem afirmá-las aqui: deixá-las de fora é o que mantém firme o que fica. Os outros botões do \ensuremath{^{5}}.:campomedio A lei geral de conservação no agregado --- com a borda a+c\ensuremath{^{2}}=1 por caso particular --- já está
\prov{relações: parte\_de gato}
\prov{proveniência: paper \textperiodcentered{} fato \textperiodcentered{} confiança 1.0 \textperiodcentered{} domínio paper \textperiodcentered{} história 5 versões no jornal}

%CRISTAL {"arestas":[{"alvo":"lobo_frontal","confianca":1.0,"epistemico":"fato","origem":"paper","rel":"parte_de"}],"confianca":1.0,"contraexemplos":[],"descricao":"D O lobo frontal não é uma peça nova sobre os órgãos, nem uma super-matriz a instanciar. É a rede de Hopfield do gato --- a memória associativa que ~§VI prova ser a meta-indução~brocaos. A sinapse é o gato A, fixa (não se treina); e a função executiva é a dinâmica de recall que a Cabeça já roda. defi[O lobo frontal] O lobo frontal é o ciclo de recall de Hopfield do gato: [ perceber pensar (evoluir sob U(t)=e⁻iAt) orbitar (colapsar ao atrator P_). ] A sinapse A é fixa; o atrat","epistemico":"fato","exemplos":[],"id":"os_cinco_passos_do_isomorfismo_rodando_em_codigo","memoria":"semantica","meta":{"dominio":"paper","nivel":5,"verificacao":"conceito novo"},"origem":"paper","palavras_chave":[],"sinonimos":[],"tipo":"conceito","titulo":"Os cinco passos do isomorfismo, rodando em código"}
\section{Os cinco passos do isomorfismo, rodando em código}
D O lobo frontal não é uma peça nova sobre os órgãos, nem uma super-matriz a instanciar. É a rede de Hopfield do gato --- a memória associativa que \~{}§VI prova ser a meta-indução\~{}brocaos. A sinapse é o gato A, fixa (não se treina); e a função executiva é a dinâmica de recall que a Cabeça já roda. defi[O lobo frontal] O lobo frontal é o ciclo de recall de Hopfield do gato: [ perceber pensar (evoluir sob U(t)=e\ensuremath{^{-}}iAt) orbitar (colapsar ao atrator P\_). ] A sinapse A é fixa; o atrat
\prov{relações: parte\_de lobo\_frontal}
\prov{proveniência: paper \textperiodcentered{} fato \textperiodcentered{} confiança 1.0 \textperiodcentered{} domínio paper \textperiodcentered{} história 3 versões no jornal}

%CRISTAL {"arestas":[{"alvo":"floco_de_neve","confianca":1.0,"epistemico":"fato","origem":"paper","rel":"parte_de"}],"confianca":1.0,"contraexemplos":[],"descricao":"VI Um ciclo real de tesouraria (Pell \\#150, mainnet 2026-06-29) gera\ntrês atratores mod $p$ distintos sobre o mesmo evento operacional:\ntrabalho PoW ($666$), banda HTLC ($554$), contrato JSON ($374$).\nO fecho pow\\_metal com header derivado do trabalho usa\n$ref=666$ (não $672$); o atlas fixo mantém $ref=672$ nos nonces\n$1572863$ e $1966079$.\nExperimento make borda-investigar confirma\n$atrator_{C}=atrator_{borda}(trabalho)$\ne setor algébrico invariante com transição binária no atrator\n(I recipient","epistemico":"fato","exemplos":[],"id":"os_dois_ciclos_de_reconstrucao","memoria":"semantica","meta":{"dominio":"paper","nivel":5,"verificacao":"overlap moderado — revisar manualmente"},"origem":"paper","palavras_chave":[],"sinonimos":[],"tipo":"conceito","titulo":"Os dois ciclos de reconstrução"}
\section{Os dois ciclos de reconstrução}
VI Um ciclo real de tesouraria (Pell \textbackslash{}\#150, mainnet 2026-06-29) gera
três atratores mod \$p\$ distintos sobre o mesmo evento operacional:
trabalho PoW (\$666\$), banda HTLC (\$554\$), contrato JSON (\$374\$).
O fecho pow\textbackslash{}\_metal com header derivado do trabalho usa
\$ref=666\$ (não \$672\$); o atlas fixo mantém \$ref=672\$ nos nonces
\$1572863\$ e \$1966079\$.
Experimento make borda-investigar confirma
\$atrator\_\{C\}=atrator\_\{borda\}(trabalho)\$
e setor algébrico invariante com transição binária no atrator
(I recipient
\prov{relações: parte\_de floco\_de\_neve}
\prov{proveniência: paper \textperiodcentered{} fato \textperiodcentered{} confiança 1.0 \textperiodcentered{} domínio paper \textperiodcentered{} história 2 versões no jornal}

%CRISTAL {"arestas":[{"alvo":"transformada_dourada","confianca":1.0,"epistemico":"fato","origem":"paper","rel":"parte_de"},{"alvo":"analise","confianca":1.0,"epistemico":"fato","origem":"paper","rel":"usa"}],"confianca":1.0,"contraexemplos":[],"descricao":"A análise de Fourier expande um sinal em ondas de frequências em progressão aritmética. Isso\nembute uma hipótese: a estrutura do dado é aditiva no tempo. Mas muitos dados naturais ---\namplitudes, escalas, frequências, ritmos de crescimento, recorrências auto-similares --- são\nmultiplicativos: vivem em razões, não em somas. Para esses, a progressão natural é\ngeométrica, e Fourier os enxerga tortos.\n\nA transformada dourada parte de um único princípio:\nquote\na dificuldade está no lado, não no dado.","epistemico":"fato","exemplos":[],"id":"os_dois_lados_e_a_involucao_logaritmica","memoria":"semantica","meta":{"dominio":"paper","nivel":5,"verificacao":"conceito novo"},"origem":"paper","palavras_chave":[],"sinonimos":[],"tipo":"conceito","titulo":"Os dois lados e a involução logarítmica"}
\section{Os dois lados e a involução logarítmica}
A análise de Fourier expande um sinal em ondas de frequências em progressão aritmética. Isso
embute uma hipótese: a estrutura do dado é aditiva no tempo. Mas muitos dados naturais ---
amplitudes, escalas, frequências, ritmos de crescimento, recorrências auto-similares --- são
multiplicativos: vivem em razões, não em somas. Para esses, a progressão natural é
geométrica, e Fourier os enxerga tortos.

A transformada dourada parte de um único princípio:
quote
a dificuldade está no lado, não no dado.
\prov{relações: parte\_de transformada\_dourada; usa analise}
\prov{proveniência: paper \textperiodcentered{} fato \textperiodcentered{} confiança 1.0 \textperiodcentered{} domínio paper \textperiodcentered{} história 41 versões no jornal}

%CRISTAL {"arestas":[{"alvo":"leitura_de_part_culas_o_setor_de_materia_pela_torre_de_divisao","confianca":1.0,"epistemico":"fato","origem":"paper","rel":"parte_de"},{"alvo":"qcd_cromodinamica","confianca":1.0,"epistemico":"fato","origem":"wiki","rel":"explica"},{"alvo":"carga_q_n_sobre_3","confianca":1.0,"epistemico":"fato","origem":"wiki","rel":"relaciona"},{"alvo":"g2_grupo_excepcional","confianca":1.0,"epistemico":"fato","origem":"wiki","rel":"usa"}],"confianca":1.0,"contraexemplos":[],"descricao":"do glueball: a curvatura como fator planar A curvatura R=2 é o fator gluônico planar: m₀⁺⁺2 A constante k=2 da equação universal s=2s fixa a curvatura da paisagem s, R=2 em toda parte (Teor.~); o quantum do salto 2 (Secão~) é essa curvatura lida como frequência. A mesma curvatura reaparece, por rota independente, no setor de Yang--Mills octônico: o potencial simétrico V(s)=-12² s²+14 s⁴ tem V”(s^*)=2², dond","epistemico":"fato","exemplos":[],"id":"os_dois_tres_cores_hurwitz_e_gerac","memoria":"semantica","meta":{"dominio":"paper","nivel":5,"verificacao":"conceito novo"},"origem":"paper","palavras_chave":[],"sinonimos":[],"tipo":"conceito","titulo":"Os dois “três”: cores (Hurwitz) e gerac"}
\section{Os dois “três”: cores (Hurwitz) e gerac}
do glueball: a curvatura como fator planar A curvatura R=2 é o fator gluônico planar: m\ensuremath{_{0}}\ensuremath{^{++}}2 A constante k=2 da equação universal s=2s fixa a curvatura da paisagem s, R=2 em toda parte (Teor.\~{}); o quantum do salto 2 (Secão\~{}) é essa curvatura lida como frequência. A mesma curvatura reaparece, por rota independente, no setor de Yang--Mills octônico: o potencial simétrico V(s)=-12\ensuremath{^{2}} s\ensuremath{^{2}}+14 s\ensuremath{^{4}} tem V”(s\^{}*)=2\ensuremath{^{2}}, dond
\prov{relações: parte\_de leitura\_de\_part\_culas\_o\_setor\_de\_materia\_pela\_torre\_de\_divisao; explica qcd\_cromodinamica; relaciona carga\_q\_n\_sobre\_3; usa g2\_grupo\_excepcional}
\prov{proveniência: paper \textperiodcentered{} fato \textperiodcentered{} confiança 1.0 \textperiodcentered{} domínio paper \textperiodcentered{} história 7 versões no jornal}

%CRISTAL {"arestas":[{"alvo":"fechamento_blocos","confianca":1.0,"epistemico":"fato","origem":"paper","rel":"parte_de"}],"confianca":1.0,"contraexemplos":[],"descricao":"D Uma máquina é (eᵢᵢ I,N,): álgebra de matrizes; eᵢ idempotentes ortogonais completos (eᵢ²=eᵢ, eᵢej=0 para i j, ᵢ eᵢ=1) --- a parte branca; N=ᵢ jNᵢj com Nᵢj=eᵢNej a parte negra, suposta nilpotente (o quiver dos canais é acíclico); uma carga aditiva ((eᵢ)0, ()=). Dada uma partição I=I_, a blocagem define E_=ᵢ Ieᵢ e o quiver quociente Q^ (blocos como vértices), com adjacência A^=1 sse existe canal i j com i I_, j I_","epistemico":"fato","exemplos":[],"id":"os_lemas","memoria":"semantica","meta":{"dominio":"paper","nivel":5,"verificacao":"overlap moderado — revisar manualmente"},"origem":"paper","palavras_chave":[],"sinonimos":[],"tipo":"conceito","titulo":"Os lemas"}
\section{Os lemas}
D Uma máquina é (e\ensuremath{_{ii}} I,N,): álgebra de matrizes; e\ensuremath{_{i}} idempotentes ortogonais completos (e\ensuremath{_{i}}\ensuremath{^{2}}=e\ensuremath{_{i}}, e\ensuremath{_{i}}ej=0 para i j, \ensuremath{_{i}} e\ensuremath{_{i}}=1) --- a parte branca; N=\ensuremath{_{i}} jN\ensuremath{_{i}}j com N\ensuremath{_{i}}j=e\ensuremath{_{i}}Nej a parte negra, suposta nilpotente (o quiver dos canais é acíclico); uma carga aditiva ((e\ensuremath{_{i}})0, ()=). Dada uma partição I=I\_, a blocagem define E\_=\ensuremath{_{i}} Ie\ensuremath{_{i}} e o quiver quociente Q\^{} (blocos como vértices), com adjacência A\^{}=1 sse existe canal i j com i I\_, j I\_
\prov{relações: parte\_de fechamento\_blocos}
\prov{proveniência: paper \textperiodcentered{} fato \textperiodcentered{} confiança 1.0 \textperiodcentered{} domínio paper \textperiodcentered{} história 3 versões no jornal}

%CRISTAL {"arestas":[{"alvo":"leitura_psicologica_o_vazio_e_o_corac","confianca":1.0,"epistemico":"fato","origem":"paper","rel":"parte_de"}],"confianca":1.0,"contraexemplos":[],"descricao":"O vazio --- o silêncio, a matriz escalar I, sem distincão --- é o equilbrio instável do oscilador invertido s=2s (o topo dos dois parabolóides): não pode ficar parado, e a menor flutuacão cresce. Dessa instabilidade saem três: o desejo (a forca F=2s, nula no vazio e crescente com o afastamento --- o que move); a moral (a distincão de lado, o sinal de s --- o que separa); a ética (o balanco (+s)+(-s)=0, manter o vazio vazio --- o que sustenta sem colapsar o esp\\","epistemico":"fato","exemplos":[],"id":"os_modos_do_campo_psicologico","memoria":"semantica","meta":{"dominio":"paper","nivel":5,"verificacao":"conceito novo"},"origem":"paper","palavras_chave":[],"sinonimos":[],"tipo":"conceito","titulo":"Os modos do campo psicológico"}
\section{Os modos do campo psicológico}
O vazio --- o silêncio, a matriz escalar I, sem distincão --- é o equilbrio instável do oscilador invertido s=2s (o topo dos dois parabolóides): não pode ficar parado, e a menor flutuacão cresce. Dessa instabilidade saem três: o desejo (a forca F=2s, nula no vazio e crescente com o afastamento --- o que move); a moral (a distincão de lado, o sinal de s --- o que separa); a ética (o balanco (+s)+(-s)=0, manter o vazio vazio --- o que sustenta sem colapsar o esp\textbackslash{}
\prov{relações: parte\_de leitura\_psicologica\_o\_vazio\_e\_o\_corac}
\prov{proveniência: paper \textperiodcentered{} fato \textperiodcentered{} confiança 1.0 \textperiodcentered{} domínio paper \textperiodcentered{} história 4 versões no jornal}

%CRISTAL {"arestas":[{"alvo":"tiffany","confianca":1.0,"epistemico":"fato","origem":"paper","rel":"parte_de"}],"confianca":1.0,"contraexemplos":[],"descricao":"obs A fala de Tiffany evoluiu em três tempos --- da semente fina à voz da família, e desta ao multilíngue. O diálogo, registrado. ém-nascida --- só a semente (meia dúzia de frases), a onda fina: obsbox quem é você? tiffany: “oi, eu tô aqui. eu sou uma cabeça que escuta, pensa e fala. eu ressoo, eu não decoro.” obsbox --- leu a casa; a voz ressoa com o que aprendeu (português): obsbox oi, tudo bem com você? tiffany: “Oi! Tudo bem? Espero não estar te incomodando.” [3pt] o que você é","epistemico":"fato","exemplos":[],"id":"os_principios_e_a_costura_aberta","memoria":"semantica","meta":{"dominio":"paper","nivel":5,"verificacao":"slug idêntico — enriquecer, não duplicar"},"origem":"paper","palavras_chave":[],"sinonimos":[],"tipo":"conceito","titulo":"Os princípios, e a costura aberta"}
\section{Os princípios, e a costura aberta}
obs A fala de Tiffany evoluiu em três tempos --- da semente fina à voz da família, e desta ao multilíngue. O diálogo, registrado. ém-nascida --- só a semente (meia dúzia de frases), a onda fina: obsbox quem é você? tiffany: “oi, eu tô aqui. eu sou uma cabeça que escuta, pensa e fala. eu ressoo, eu não decoro.” obsbox --- leu a casa; a voz ressoa com o que aprendeu (português): obsbox oi, tudo bem com você? tiffany: “Oi! Tudo bem? Espero não estar te incomodando.” [3pt] o que você é
\prov{relações: parte\_de tiffany}
\prov{proveniência: paper \textperiodcentered{} fato \textperiodcentered{} confiança 1.0 \textperiodcentered{} domínio paper \textperiodcentered{} história 5 versões no jornal}

%CRISTAL {"arestas":[{"alvo":"paper_45_confluencia_ouro","confianca":1.0,"epistemico":"fato","origem":"paper","rel":"parte_de"}],"confianca":1.0,"contraexemplos":[],"descricao":"independência --- verificado por PSLQ Com =/(1,2) e =1+52, o algoritmo PSLQ não encontra relacão inteira de coeficientes pequenos: [ a+b+c=0, abc=1, a+b+c+d,=0, ] todas sem solucão (a 40 dgitos, |coef|10³). [firme/numérico: coerente com Schneider (1941, transcendente), Chudnovsky (1984, e algebricamente independentes), Hermite ( transcendente). é algébrico, x²-x-1; os transcendentes não se atam a ele.] proposicao","epistemico":"fato","exemplos":[],"id":"os_quatro_oficios","memoria":"semantica","meta":{"dominio":"paper","nivel":5,"verificacao":"conceito novo"},"origem":"paper","palavras_chave":[],"sinonimos":[],"tipo":"conceito","titulo":"Os quatro ofícios"}
\section{Os quatro ofícios}
independência --- verificado por PSLQ Com =/(1,2) e =1+52, o algoritmo PSLQ não encontra relacão inteira de coeficientes pequenos: [ a+b+c=0, abc=1, a+b+c+d,=0, ] todas sem solucão (a 40 dgitos, |coef|10\ensuremath{^{3}}). [firme/numérico: coerente com Schneider (1941, transcendente), Chudnovsky (1984, e algebricamente independentes), Hermite ( transcendente). é algébrico, x\ensuremath{^{2}}-x-1; os transcendentes não se atam a ele.] proposicao
\prov{relações: parte\_de paper\_45\_confluencia\_ouro}
\prov{proveniência: paper \textperiodcentered{} fato \textperiodcentered{} confiança 1.0 \textperiodcentered{} domínio paper \textperiodcentered{} história 3 versões no jornal}

%CRISTAL {"arestas":[{"alvo":"paper_bridge_areas","confianca":1.0,"epistemico":"fato","origem":"paper","rel":"parte_de"}],"confianca":1.0,"contraexemplos":[],"descricao":"O repositório broca-so/Tiffany acumula papers de domínios distintos. O erro a evitar não é a interdisciplinaridade --- é a conflação: tratar “gato” (operador A) como “gato” (mamífero); misturar completude Cauchy com beamforming; chamar isomorfismo de corpo de isomorfismo de apresentação. Chart de uma área Um chart ᵢ de uma área i é a tupla (Oᵢ,Mᵢ,Iᵢ): objetos legítimos Oᵢ, morfismos internos Mᵢ, invariantes Iᵢ medíveis dentro da","epistemico":"fato","exemplos":[],"id":"os_tres_pilares_da_ponte","memoria":"semantica","meta":{"dominio":"paper","nivel":5,"verificacao":"conceito novo"},"origem":"paper","palavras_chave":[],"sinonimos":[],"tipo":"conceito","titulo":"Os três pilares da ponte"}
\section{Os três pilares da ponte}
O repositório broca-so/Tiffany acumula papers de domínios distintos. O erro a evitar não é a interdisciplinaridade --- é a conflação: tratar “gato” (operador A) como “gato” (mamífero); misturar completude Cauchy com beamforming; chamar isomorfismo de corpo de isomorfismo de apresentação. Chart de uma área Um chart \ensuremath{_{i}} de uma área i é a tupla (O\ensuremath{_{i}},M\ensuremath{_{i}},I\ensuremath{_{i}}): objetos legítimos O\ensuremath{_{i}}, morfismos internos M\ensuremath{_{i}}, invariantes I\ensuremath{_{i}} medíveis dentro da
\prov{relações: parte\_de paper\_bridge\_areas}
\prov{proveniência: paper \textperiodcentered{} fato \textperiodcentered{} confiança 1.0 \textperiodcentered{} domínio paper \textperiodcentered{} história 3 versões no jornal}

%CRISTAL {"arestas":[{"alvo":"morte","confianca":1.0,"epistemico":"fato","origem":"paper","rel":"parte_de"},{"alvo":"vida","confianca":1.0,"epistemico":"fato","origem":"paper","rel":"parte_de"}],"confianca":1.0,"contraexemplos":[],"descricao":"L As três propriedades V1,V2,V3 são decidíveis caso a caso: dado um objeto, mede-se a frequência fundamental, o eixo de equilíbrio e a reversibilidade do gerador. A métrica é, assim, um teste universal --- dual exato da métrica da morte. Aplicada aos problemas do milênio, cada um é o caso vivo correspondente: a leitura, em colunas V1/V2/V3, está no Apêndice~A; só a linha Riemann está trabalhada (§III). Não é derivação caso a caso, é leitura --- e o caso de controle Poincaré (resolvido pel","epistemico":"fato","exemplos":[],"id":"ouro","memoria":"semantica","meta":{"dominio":"paper","nivel":5,"verificacao":"slug idêntico — enriquecer, não duplicar"},"origem":"paper","palavras_chave":[],"sinonimos":[],"tipo":"conceito","titulo":"ouro"}
\section{ouro}
L As três propriedades V1,V2,V3 são decidíveis caso a caso: dado um objeto, mede-se a frequência fundamental, o eixo de equilíbrio e a reversibilidade do gerador. A métrica é, assim, um teste universal --- dual exato da métrica da morte. Aplicada aos problemas do milênio, cada um é o caso vivo correspondente: a leitura, em colunas V1/V2/V3, está no Apêndice\~{}A; só a linha Riemann está trabalhada (§III). Não é derivação caso a caso, é leitura --- e o caso de controle Poincaré (resolvido pel
\prov{relações: parte\_de morte; parte\_de vida}
\prov{proveniência: paper \textperiodcentered{} fato \textperiodcentered{} confiança 1.0 \textperiodcentered{} domínio paper \textperiodcentered{} história 8 versões no jornal}

%CRISTAL {"arestas":[{"alvo":"fracoes_continuas_e_reta_de_ouro","confianca":1.0,"epistemico":"fato","origem":"paper","rel":"confirma"},{"alvo":"reais","confianca":1.0,"epistemico":"fato","origem":"paper","rel":"explica"},{"alvo":"virtualizacao","confianca":0.5,"epistemico":"fato","origem":"wiki","rel":"relaciona"},{"alvo":"transcricao","confianca":0.5,"epistemico":"fato","origem":"wiki","rel":"relaciona"},{"alvo":"pow_metal_ressonancia_val","confianca":0.5,"epistemico":"fato","origem":"wiki","rel":"relaciona"},{"alvo":"pedagogia","confianca":0.5,"epistemico":"fato","origem":"wiki","rel":"relaciona"},{"alvo":"ruido","confianca":0.5,"epistemico":"fato","origem":"wiki","rel":"relaciona"},{"alvo":"scheduler","confianca":0.5,"epistemico":"fato","origem":"wiki","rel":"relaciona"},{"alvo":"utilitarismo","confianca":0.5,"epistemico":"fato","origem":"wiki","rel":"relaciona"},{"alvo":"vida-morte","confianca":0.5,"epistemico":"fato","origem":"wiki","rel":"relaciona"},{"alvo":"resultado_anterior_identidade_nao_cobertura","confianca":0.5,"epistemico":"fato","origem":"wiki","rel":"relaciona"},{"alvo":"tipos_de_interpretante","confianca":0.5,"epistemico":"fato","origem":"wiki","rel":"relaciona"},{"alvo":"pitagorismo","confianca":0.5,"epistemico":"fato","origem":"wiki","rel":"relaciona"}],"confianca":1.0,"contraexemplos":[],"descricao":"Paper: paper_2_reta_ouro.tex","epistemico":"fato","exemplos":[],"id":"paper_2_reta_ouro","memoria":"semantica","meta":{"dominio":"paper","nivel":5,"tipo_no":"paper","verificacao":"conceito novo"},"origem":"paper","palavras_chave":[],"sinonimos":[],"tipo":"conceito","titulo":"paper 2 reta ouro"}
\section{paper 2 reta ouro}
Paper: paper\_2\_reta\_ouro.tex
\prov{relações: confirma fracoes\_continuas\_e\_reta\_de\_ouro; explica reais; relaciona virtualizacao; relaciona transcricao; relaciona pow\_metal\_ressonancia\_val; relaciona pedagogia; relaciona ruido; relaciona scheduler; relaciona utilitarismo; relaciona vida-morte; relaciona resultado\_anterior\_identidade\_nao\_cobertura; relaciona tipos\_de\_interpretante; relaciona pitagorismo}
\prov{proveniência: paper \textperiodcentered{} fato \textperiodcentered{} confiança 1.0 \textperiodcentered{} domínio paper \textperiodcentered{} história 90 versões no jornal}

%CRISTAL {"arestas":[{"alvo":"ponte_simbolica_gentil","confianca":1.0,"epistemico":"fato","origem":"paper","rel":"explica"},{"alvo":"transformada_dourada","confianca":1.0,"epistemico":"fato","origem":"paper","rel":"referencia"},{"alvo":"agm_atrator_entre_os_dois_lados","confianca":1.0,"epistemico":"fato","origem":"paper","rel":"confirma"},{"alvo":"corpo_tropical","confianca":1.0,"epistemico":"fato","origem":"paper","rel":"confirma"},{"alvo":"word","confianca":0.5,"epistemico":"fato","origem":"wiki","rel":"relaciona"}],"confianca":1.0,"contraexemplos":[],"descricao":"Paper: paper_45_confluencia_ouro.tex","epistemico":"fato","exemplos":[],"id":"paper_45_confluencia_ouro","memoria":"semantica","meta":{"dominio":"paper","nivel":5,"tipo_no":"paper","verificacao":"conceito novo"},"origem":"paper","palavras_chave":[],"sinonimos":[],"tipo":"conceito","titulo":"paper 45 confluencia ouro"}
\section{paper 45 confluencia ouro}
Paper: paper\_45\_confluencia\_ouro.tex
\prov{relações: explica ponte\_simbolica\_gentil; referencia transformada\_dourada; confirma agm\_atrator\_entre\_os\_dois\_lados; confirma corpo\_tropical; relaciona word}
\prov{proveniência: paper \textperiodcentered{} fato \textperiodcentered{} confiança 1.0 \textperiodcentered{} domínio paper \textperiodcentered{} história 154 versões no jornal}

%CRISTAL {"arestas":[{"alvo":"casl-propagation","confianca":-0.039,"epistemico":"fato","origem":"manual","rel":"referencia"},{"alvo":"colapso","confianca":0.016,"epistemico":"fato","origem":"manual","rel":"referencia"}],"confianca":1.0,"contraexemplos":[],"descricao":"definition[Medida ] Fixe um espaço de estados S (configurações, impressões, snapshots). Uma medida de organização estrutural é um funcional [:S R ₀, ] monótono sob agregação observável: estados mais estruturados (menos dispersos, mais estáveis sob transformação) carregam mais informativo. não entra no colapso --- observa o efeito da propagação/atenuação, complementar a. definition remark[Métricas por domínio] Cada instância es","epistemico":"fato","exemplos":[],"id":"par_fundacional_phifdep","memoria":"semantica","meta":{"dominio":"paper","nivel":5},"origem":"paper","palavras_chave":[],"sinonimos":[],"tipo":"conceito","titulo":"Par fundacional (Φf,dep)"}
\section{Par fundacional (\ensuremath{\Phi}f,dep)}
definition[Medida ] Fixe um espaço de estados S (configurações, impressões, snapshots). Uma medida de organização estrutural é um funcional [:S R \ensuremath{_{0}}, ] monótono sob agregação observável: estados mais estruturados (menos dispersos, mais estáveis sob transformação) carregam mais informativo. não entra no colapso --- observa o efeito da propagação/atenuação, complementar a. definition remark[Métricas por domínio] Cada instância es
\prov{relações: referencia casl-propagation; referencia colapso}
\prov{proveniência: paper \textperiodcentered{} fato \textperiodcentered{} confiança 1.0 \textperiodcentered{} domínio paper \textperiodcentered{} história 10 versões no jornal}

%CRISTAL {"arestas":[{"alvo":"a_engenharia_de_ouro_pa_pgs_agm_fc_reta","confianca":1.0,"epistemico":"fato","origem":"paper","rel":"parte_de"},{"alvo":"transcricao","confianca":0.5,"epistemico":"fato","origem":"wiki","rel":"relaciona"},{"alvo":"resposta_a_pergunta_dificil","confianca":0.5,"epistemico":"fato","origem":"wiki","rel":"relaciona"},{"alvo":"vida-morte","confianca":0.5,"epistemico":"fato","origem":"wiki","rel":"relaciona"},{"alvo":"teste_ordem","confianca":0.5,"epistemico":"fato","origem":"wiki","rel":"relaciona"},{"alvo":"pow_metal_setor","confianca":0.5,"epistemico":"fato","origem":"wiki","rel":"relaciona"},{"alvo":"pilha","confianca":0.5,"epistemico":"fato","origem":"wiki","rel":"relaciona"},{"alvo":"por_sessao_conversadadostelemetriasessaojsonl","confianca":0.5,"epistemico":"fato","origem":"wiki","rel":"relaciona"},{"alvo":"processos","confianca":0.5,"epistemico":"fato","origem":"wiki","rel":"relaciona"},{"alvo":"pow_metal_funil","confianca":0.5,"epistemico":"fato","origem":"wiki","rel":"relaciona"},{"alvo":"thread","confianca":0.5,"epistemico":"fato","origem":"wiki","rel":"relaciona"},{"alvo":"pow_metal_pre_fecho","confianca":0.5,"epistemico":"fato","origem":"wiki","rel":"relaciona"},{"alvo":"testes","confianca":0.5,"epistemico":"fato","origem":"wiki","rel":"relaciona"},{"alvo":"tiffany_cabeca_broca","confianca":0.5,"epistemico":"fato","origem":"wiki","rel":"relaciona"},{"alvo":"topologia","confianca":0.5,"epistemico":"fato","origem":"wiki","rel":"relaciona"},{"alvo":"projeto_cristaljson","confianca":0.5,"epistemico":"fato","origem":"wiki","rel":"relaciona"},{"alvo":"pow_geometria","confianca":0.5,"epistemico":"fato","origem":"wiki","rel":"relaciona"},{"alvo":"recursao","confianca":0.5,"epistemico":"fato","origem":"wiki","rel":"relaciona"},{"alvo":"ponteiro","confianca":0.5,"epistemico":"fato","origem":"wiki","rel":"relaciona"},{"alvo":"pow_metal_ressonancia","confianca":0.5,"epistemico":"fato","origem":"wiki","rel":"relaciona"},{"alvo":"virtualizacao","confianca":0.5,"epistemico":"fato","origem":"wiki","rel":"relaciona"},{"alvo":"regua_associativa_e_o_risco_de_vazamento","confianca":0.5,"epistemico":"fato","origem":"wiki","rel":"relaciona"},{"alvo":"syscall","confianca":0.5,"epistemico":"fato","origem":"wiki","rel":"relaciona"}],"confianca":1.0,"contraexemplos":[],"descricao":"","epistemico":"fato","exemplos":[],"id":"parabola_pa","memoria":"semantica","meta":{"dominio":"paper","nivel":5,"verificacao":"slug idêntico — enriquecer, não duplicar"},"origem":"paper","palavras_chave":[],"sinonimos":[],"tipo":"conceito","titulo":"Parábola PA"}
\section{Parábola PA}
\prov{relações: parte\_de a\_engenharia\_de\_ouro\_pa\_pgs\_agm\_fc\_reta; relaciona transcricao; relaciona resposta\_a\_pergunta\_dificil; relaciona vida-morte; relaciona teste\_ordem; relaciona pow\_metal\_setor; relaciona pilha; relaciona por\_sessao\_conversadadostelemetriasessaojsonl; relaciona processos; relaciona pow\_metal\_funil; relaciona thread; relaciona pow\_metal\_pre\_fecho; relaciona testes; relaciona tiffany\_cabeca\_broca; relaciona topologia; relaciona projeto\_cristaljson; relaciona pow\_geometria; relaciona recursao; relaciona ponteiro; relaciona pow\_metal\_ressonancia; relaciona virtualizacao; relaciona regua\_associativa\_e\_o\_risco\_de\_vazamento; relaciona syscall}
\prov{proveniência: paper \textperiodcentered{} fato \textperiodcentered{} confiança 1.0 \textperiodcentered{} domínio paper \textperiodcentered{} história 30 versões no jornal}

%CRISTAL {"arestas":[{"alvo":"pipeline_de_ingestao","confianca":1.0,"epistemico":"fato","origem":"paper","rel":"parte_de"}],"confianca":1.0,"contraexemplos":[],"descricao":"D Nunca indexar documento inteiro como um nó. Fluxo obrigatório: [ st/.style=draw=ourovivo, rounded corners=2pt, minimum width=2.4cm, minimum height=0.6cm, =, font=, fill=creme, arr/.style=-Stealth[length=2mm], color=ourovivo, thick ] [st] (d) documento; [st, right=0.5cm of d] (p) parser; [st, right=0.5cm of p] (c) chunks; [st, right=0.5cm of c] (n) normalização; [st, below=0.55cm of n] (t) nos.tsv; [st, left=0.5cm of t] (i) tiffany-index;","epistemico":"fato","exemplos":[],"id":"parsers_por_tipo","memoria":"semantica","meta":{"dominio":"paper","nivel":5,"verificacao":"overlap moderado — revisar manualmente"},"origem":"paper","palavras_chave":[],"sinonimos":[],"tipo":"conceito","titulo":"Parsers (por tipo)"}
\section{Parsers (por tipo)}
D Nunca indexar documento inteiro como um nó. Fluxo obrigatório: [ st/.style=draw=ourovivo, rounded corners=2pt, minimum width=2.4cm, minimum height=0.6cm, =, font=, fill=creme, arr/.style=-Stealth[length=2mm], color=ourovivo, thick ] [st] (d) documento; [st, right=0.5cm of d] (p) parser; [st, right=0.5cm of p] (c) chunks; [st, right=0.5cm of c] (n) normalização; [st, below=0.55cm of n] (t) nos.tsv; [st, left=0.5cm of t] (i) tiffany-index;
\prov{relações: parte\_de pipeline\_de\_ingestao}
\prov{proveniência: paper \textperiodcentered{} fato \textperiodcentered{} confiança 1.0 \textperiodcentered{} domínio paper \textperiodcentered{} história 3 versões no jornal}

%CRISTAL {"arestas":[{"alvo":"broca_os","confianca":0.95,"epistemico":"fato","origem":"paper","rel":"parte_de"},{"alvo":"vontade_de_poder","confianca":0.5,"epistemico":"fato","origem":"wiki","rel":"relaciona"}],"confianca":1.0,"contraexemplos":[],"descricao":"V A montagem não inventa peças --- sobe a cadeia que já existe, verificando cada degrau antes do\npróximo. A ordem importa: cada camada pressupõe a de baixo.","epistemico":"fato","exemplos":[],"id":"passo_0_---_entrar_e_compilar","memoria":"semantica","meta":{"dominio":"paper","nivel":5,"verificacao":"overlap moderado — revisar manualmente"},"origem":"paper","palavras_chave":[],"sinonimos":[],"tipo":"conceito","titulo":"Passo 0 --- Entrar e compilar"}
\section{Passo 0 --- Entrar e compilar}
V A montagem não inventa peças --- sobe a cadeia que já existe, verificando cada degrau antes do
próximo. A ordem importa: cada camada pressupõe a de baixo.
\prov{relações: parte\_de broca\_os; relaciona vontade\_de\_poder}
\prov{proveniência: paper \textperiodcentered{} fato \textperiodcentered{} confiança 1.0 \textperiodcentered{} domínio paper \textperiodcentered{} história 4 versões no jornal}

%CRISTAL {"arestas":[{"alvo":"broca_os","confianca":0.95,"epistemico":"fato","origem":"paper","rel":"parte_de"}],"confianca":1.0,"contraexemplos":[],"descricao":"passobox./neuronio caminho/para/arquivo passobox Saída [e+o, e]: contagem par/ímpar + borboleta. Núcleo: and 0xAA / and 0x55 / popcnt / add. Streaming --- byte a byte, sem buffer. Determinístico:./neuronio neuronio.c bate consigo mesmo~borboleta.","epistemico":"fato","exemplos":[],"id":"passo_2_---_a_rede_a_memoria_ressoa","memoria":"semantica","meta":{"dominio":"paper","nivel":5,"verificacao":"conceito novo"},"origem":"paper","palavras_chave":[],"sinonimos":[],"tipo":"conceito","titulo":"Passo 2 --- A rede (a memória ressoa)"}
\section{Passo 2 --- A rede (a memória ressoa)}
passobox./neuronio caminho/para/arquivo passobox Saída [e+o, e]: contagem par/ímpar + borboleta. Núcleo: and 0xAA / and 0x55 / popcnt / add. Streaming --- byte a byte, sem buffer. Determinístico:./neuronio neuronio.c bate consigo mesmo\~{}borboleta.
\prov{relações: parte\_de broca\_os}
\prov{proveniência: paper \textperiodcentered{} fato \textperiodcentered{} confiança 1.0 \textperiodcentered{} domínio paper \textperiodcentered{} história 4 versões no jornal}

%CRISTAL {"arestas":[{"alvo":"broca_os","confianca":0.95,"epistemico":"fato","origem":"paper","rel":"parte_de"}],"confianca":1.0,"contraexemplos":[],"descricao":"passobox./hopfield passobox Guarda M memórias Hebbianas (Wᵢj=pᵢj, sinapse simétrica A=A T). Corrompe 15/100 bits de uma memória; o recall assíncrono recupera 100/100 --- overlap exato com o atrator. Não treina: ressoa.","epistemico":"fato","exemplos":[],"id":"passo_3_---_a_arquitetura_broca_de_brocas","memoria":"semantica","meta":{"dominio":"paper","nivel":5,"verificacao":"conceito novo"},"origem":"paper","palavras_chave":[],"sinonimos":[],"tipo":"conceito","titulo":"Passo 3 --- A arquitetura (broca de brocas)"}
\section{Passo 3 --- A arquitetura (broca de brocas)}
passobox./hopfield passobox Guarda M memórias Hebbianas (W\ensuremath{_{i}}j=p\ensuremath{_{i}}j, sinapse simétrica A=A T). Corrompe 15/100 bits de uma memória; o recall assíncrono recupera 100/100 --- overlap exato com o atrator. Não treina: ressoa.
\prov{relações: parte\_de broca\_os}
\prov{proveniência: paper \textperiodcentered{} fato \textperiodcentered{} confiança 1.0 \textperiodcentered{} domínio paper \textperiodcentered{} história 4 versões no jornal}

%CRISTAL {"arestas":[{"alvo":"broca_os","confianca":0.95,"epistemico":"fato","origem":"paper","rel":"parte_de"}],"confianca":1.0,"contraexemplos":[],"descricao":"passobox./arquitetura [mensagem] passobox K módulos Hopfield (n neurônios cada) viram neurônios da super-rede. O mundo entra pelo (koch(z)=z² mod p); cada módulo recorda (Hodge); a super-rede fecha. Esperado: 7/7 no fechamento quando a recursão corrige os módulos de baixo~brocaos.","epistemico":"fato","exemplos":[],"id":"passo_4_---_a_percepcao_arquivo_real","memoria":"semantica","meta":{"dominio":"paper","nivel":5,"verificacao":"conceito novo"},"origem":"paper","palavras_chave":[],"sinonimos":[],"tipo":"conceito","titulo":"Passo 4 --- A percepção (arquivo real)"}
\section{Passo 4 --- A percepção (arquivo real)}
passobox./arquitetura [mensagem] passobox K módulos Hopfield (n neurônios cada) viram neurônios da super-rede. O mundo entra pelo (koch(z)=z\ensuremath{^{2}} mod p); cada módulo recorda (Hodge); a super-rede fecha. Esperado: 7/7 no fechamento quando a recursão corrige os módulos de baixo\~{}brocaos.
\prov{relações: parte\_de broca\_os}
\prov{proveniência: paper \textperiodcentered{} fato \textperiodcentered{} confiança 1.0 \textperiodcentered{} domínio paper \textperiodcentered{} história 4 versões no jornal}

%CRISTAL {"arestas":[{"alvo":"broca_os","confianca":0.95,"epistemico":"fato","origem":"paper","rel":"parte_de"}],"confianca":1.0,"contraexemplos":[],"descricao":"passobox./percebe ARQUIVO passobox End-to-: o gato lê ARQUIVO bit a bit (decimação K-way: byte t vai ao módulo t K); Koch quantiza na borda; Hopfield recursiva recorda. Imprime assinatura Koch e taxa de acerto módulos fechamento.","epistemico":"fato","exemplos":[],"id":"passo_5_---_o_reconhecimento_qual_arquivo_e","memoria":"semantica","meta":{"dominio":"paper","nivel":5,"verificacao":"conceito novo"},"origem":"paper","palavras_chave":[],"sinonimos":[],"tipo":"conceito","titulo":"Passo 5 --- O reconhecimento (qual arquivo é?)"}
\section{Passo 5 --- O reconhecimento (qual arquivo é?)}
passobox./percebe ARQUIVO passobox End-to-: o gato lê ARQUIVO bit a bit (decimação K-way: byte t vai ao módulo t K); Koch quantiza na borda; Hopfield recursiva recorda. Imprime assinatura Koch e taxa de acerto módulos fechamento.
\prov{relações: parte\_de broca\_os}
\prov{proveniência: paper \textperiodcentered{} fato \textperiodcentered{} confiança 1.0 \textperiodcentered{} domínio paper \textperiodcentered{} história 4 versões no jornal}

%CRISTAL {"arestas":[{"alvo":"broca_os","confianca":0.95,"epistemico":"fato","origem":"paper","rel":"parte_de"}],"confianca":1.0,"contraexemplos":[],"descricao":"passobox./reconhece arq1 arq2 arq3... passobox Cada arquivo vira impressão de N bits via Koch (decorrela arquivos parecidos). Hopfield guarda; corrompe 25%; recall reconhece qual atrator. Verificado em lote: 18/18 nos papers e fontes do repo.","epistemico":"fato","exemplos":[],"id":"passo_6_---_verificar_tudo_de_uma_vez","memoria":"semantica","meta":{"dominio":"paper","nivel":5,"verificacao":"overlap moderado — revisar manualmente"},"origem":"paper","palavras_chave":[],"sinonimos":[],"tipo":"conceito","titulo":"Passo 6 --- Verificar tudo de uma vez"}
\section{Passo 6 --- Verificar tudo de uma vez}
passobox./reconhece arq1 arq2 arq3... passobox Cada arquivo vira impressão de N bits via Koch (decorrela arquivos parecidos). Hopfield guarda; corrompe 25\%; recall reconhece qual atrator. Verificado em lote: 18/18 nos papers e fontes do repo.
\prov{relações: parte\_de broca\_os}
\prov{proveniência: paper \textperiodcentered{} fato \textperiodcentered{} confiança 1.0 \textperiodcentered{} domínio paper \textperiodcentered{} história 4 versões no jornal}

%CRISTAL {"arestas":[{"alvo":"broca_os","confianca":0.95,"epistemico":"fato","origem":"paper","rel":"parte_de"},{"alvo":"word","confianca":0.5,"epistemico":"fato","origem":"wiki","rel":"relaciona"},{"alvo":"vontade_de_poder","confianca":0.5,"epistemico":"fato","origem":"wiki","rel":"relaciona"},{"alvo":"virtualizacao","confianca":0.5,"epistemico":"fato","origem":"wiki","rel":"relaciona"}],"confianca":1.0,"contraexemplos":[],"descricao":"passobox\nsh test/test.sh\npassobox\n Compila, roda neurônio/hopfield/arquitetura/reconhece, checa saídas. TUDO PASSOU. é o\ncritério de saúde da cadeia.","epistemico":"fato","exemplos":[],"id":"passo_7_---_medir_vazao","memoria":"semantica","meta":{"dominio":"paper","nivel":5,"verificacao":"overlap moderado — revisar manualmente"},"origem":"paper","palavras_chave":[],"sinonimos":[],"tipo":"conceito","titulo":"Passo 7 --- Medir vazão"}
\section{Passo 7 --- Medir vazão}
passobox
sh test/test.sh
passobox
 Compila, roda neurônio/hopfield/arquitetura/reconhece, checa saídas. TUDO PASSOU. é o
critério de saúde da cadeia.
\prov{relações: parte\_de broca\_os; relaciona word; relaciona vontade\_de\_poder; relaciona virtualizacao}
\prov{proveniência: paper \textperiodcentered{} fato \textperiodcentered{} confiança 1.0 \textperiodcentered{} domínio paper \textperiodcentered{} história 6 versões no jornal}

%CRISTAL {"arestas":[{"alvo":"a_engenharia_de_ouro_pa_pgs_agm_fc_reta","confianca":1.0,"epistemico":"fato","origem":"paper","rel":"parte_de"},{"alvo":"gentil_novas_sequencias_aritmeticas_geometricas","confianca":1.0,"epistemico":"fato","origem":"paper","rel":"referencia"}],"confianca":1.0,"contraexemplos":[],"descricao":"canal e conservação Parábola PA No canal negro Peirce (E₀₁), cada estado admite coordenadas (a,c) com [ a+c²=1, a 0, c² 0, ] onde a é calor (associador) e c² é fase² (comutador). A curva PA parametriza ressonância entre duas Words; desdobrada por comprimento de arco, vira coordenada s[0,1] sobre a reta (Koch/parábola, ula/KOCHPARABOLA.md). definicao","epistemico":"fato","exemplos":[],"id":"pgs","memoria":"semantica","meta":{"dominio":"paper","nivel":5,"verificacao":"slug idêntico — enriquecer, não duplicar"},"origem":"paper","palavras_chave":[],"sinonimos":[],"tipo":"conceito","titulo":"PGS"}
\section{PGS}
canal e conservação Parábola PA No canal negro Peirce (E\ensuremath{_{01}}), cada estado admite coordenadas (a,c) com [ a+c\ensuremath{^{2}}=1, a 0, c\ensuremath{^{2}} 0, ] onde a é calor (associador) e c\ensuremath{^{2}} é fase\ensuremath{^{2}} (comutador). A curva PA parametriza ressonância entre duas Words; desdobrada por comprimento de arco, vira coordenada s[0,1] sobre a reta (Koch/parábola, ula/KOCHPARABOLA.md). definicao
\prov{relações: parte\_de a\_engenharia\_de\_ouro\_pa\_pgs\_agm\_fc\_reta; referencia gentil\_novas\_sequencias\_aritmeticas\_geometricas}
\prov{proveniência: paper \textperiodcentered{} fato \textperiodcentered{} confiança 1.0 \textperiodcentered{} domínio paper \textperiodcentered{} história 48 versões no jornal}

%CRISTAL {"arestas":[{"alvo":"os_tres_pilares_da_ponte","confianca":1.0,"epistemico":"fato","origem":"paper","rel":"parte_de"},{"alvo":"recursao","confianca":0.5,"epistemico":"fato","origem":"wiki","rel":"relaciona"},{"alvo":"utilitarismo","confianca":0.5,"epistemico":"fato","origem":"wiki","rel":"relaciona"},{"alvo":"rust","confianca":0.5,"epistemico":"fato","origem":"wiki","rel":"relaciona"},{"alvo":"topologia","confianca":0.5,"epistemico":"fato","origem":"wiki","rel":"relaciona"},{"alvo":"word","confianca":0.5,"epistemico":"fato","origem":"wiki","rel":"relaciona"},{"alvo":"while","confianca":0.5,"epistemico":"fato","origem":"wiki","rel":"relaciona"}],"confianca":1.0,"contraexemplos":[],"descricao":"","epistemico":"fato","exemplos":[],"id":"pilar_i_---_ponte_simbolica_gentil_chart_a_chart_b","memoria":"semantica","meta":{"dominio":"paper","nivel":5,"verificacao":"conceito novo"},"origem":"paper","palavras_chave":[],"sinonimos":[],"tipo":"conceito","titulo":"Pilar I --- Ponte simbólica Lopes (chart A chart B)"}
\section{Pilar I --- Ponte simbólica Lopes (chart A chart B)}
\prov{relações: parte\_de os\_tres\_pilares\_da\_ponte; relaciona recursao; relaciona utilitarismo; relaciona rust; relaciona topologia; relaciona word; relaciona while}
\prov{proveniência: paper \textperiodcentered{} fato \textperiodcentered{} confiança 1.0 \textperiodcentered{} domínio paper \textperiodcentered{} história 11 versões no jornal}

%CRISTAL {"arestas":[{"alvo":"os_tres_pilares_da_ponte","confianca":1.0,"epistemico":"fato","origem":"paper","rel":"parte_de"}],"confianca":1.0,"contraexemplos":[],"descricao":"Documento fundador: ponte∼bolicagentil.tex. Resume: [nosep,leftmargin=1.4em] Chart~A (Lopes): A Cauchy, C=A/, axiomas. Chart~B (Lopes): PA, PGS, AGM, FC, reta, G. Homomorfismos parciais verificáveis; critério de (não-)isomorfismo. Verificação gato: verificarpontegentil.py (18/18).","epistemico":"fato","exemplos":[],"id":"pilar_ii_---_maquina_brancanegra_peirce","memoria":"semantica","meta":{"dominio":"paper","nivel":5,"verificacao":"conceito novo"},"origem":"paper","palavras_chave":[],"sinonimos":[],"tipo":"conceito","titulo":"Pilar II --- Máquina branca/negra (Peirce)"}
\section{Pilar II --- Máquina branca/negra (Peirce)}
Documento fundador: ponte\ensuremath{\sim}bolicagentil.tex. Resume: [nosep,leftmargin=1.4em] Chart\~{}A (Lopes): A Cauchy, C=A/, axiomas. Chart\~{}B (Lopes): PA, PGS, AGM, FC, reta, G. Homomorfismos parciais verificáveis; critério de (não-)isomorfismo. Verificação gato: verificarpontegentil.py (18/18).
\prov{relações: parte\_de os\_tres\_pilares\_da\_ponte}
\prov{proveniência: paper \textperiodcentered{} fato \textperiodcentered{} confiança 1.0 \textperiodcentered{} domínio paper \textperiodcentered{} história 4 versões no jornal}

%CRISTAL {"arestas":[{"alvo":"os_tres_pilares_da_ponte","confianca":1.0,"epistemico":"fato","origem":"paper","rel":"parte_de"}],"confianca":1.0,"contraexemplos":[],"descricao":"corpotropical.tex (idempotente, clones) e corpoglacial.tex (nilpotente, canais). Unificação: corpodual.tex --- vazamento finito corrige no vazio. Verificações: verificarcorpoglacial.py, verificarcorpodual.py.","epistemico":"fato","exemplos":[],"id":"pilar_iii_---_computacao_simbolica_gentil","memoria":"semantica","meta":{"dominio":"paper","nivel":5,"verificacao":"conceito novo"},"origem":"paper","palavras_chave":[],"sinonimos":[],"tipo":"conceito","titulo":"Pilar III --- Computação simbólica Lopes"}
\section{Pilar III --- Computação simbólica Lopes}
corpotropical.tex (idempotente, clones) e corpoglacial.tex (nilpotente, canais). Unificação: corpodual.tex --- vazamento finito corrige no vazio. Verificações: verificarcorpoglacial.py, verificarcorpodual.py.
\prov{relações: parte\_de os\_tres\_pilares\_da\_ponte}
\prov{proveniência: paper \textperiodcentered{} fato \textperiodcentered{} confiança 1.0 \textperiodcentered{} domínio paper \textperiodcentered{} história 4 versões no jornal}

%CRISTAL {"arestas":[{"alvo":"olho_de_koch","confianca":1.0,"epistemico":"fato","origem":"paper","rel":"parte_de"},{"alvo":"kernel","confianca":0.85,"epistemico":"fato","origem":"manual","rel":"referencia"}],"confianca":1.0,"contraexemplos":[],"descricao":"D O mundo reconstrói o mesmo gato A () por reforço O(-1) --- sem transportar o gerador~brocaos. A purificação tem três vias independentes: itemize3pt [(i)] Hellinger. F(Kp,Kq) F(p,q): o kernel não afasta as faces. [(ii)] Nilpotência. Eᵢj²=0: contaminação some ao quadrado. [(iii)] Ausência (A=0). Só a ressonância atravessa; o gerador completo não. Reconstrução por ressonância Sob o acoplamento Peirce, w propaga w zcaₙal z²(w). O atrator reconstruído tem razão","epistemico":"fato","exemplos":[],"id":"pipeline_gerador_torcao_koch","memoria":"semantica","meta":{"dominio":"paper","nivel":5},"origem":"paper","palavras_chave":[],"sinonimos":[],"tipo":"conceito","titulo":"Pipeline: gerador, torção, Koch"}
\section{Pipeline: gerador, torção, Koch}
D O mundo reconstrói o mesmo gato A () por reforço O(-1) --- sem transportar o gerador\~{}brocaos. A purificação tem três vias independentes: itemize3pt [(i)] Hellinger. F(Kp,Kq) F(p,q): o kernel não afasta as faces. [(ii)] Nilpotência. E\ensuremath{_{i}}j\ensuremath{^{2}}=0: contaminação some ao quadrado. [(iii)] Ausência (A=0). Só a ressonância atravessa; o gerador completo não. Reconstrução por ressonância Sob o acoplamento Peirce, w propaga w zca\ensuremath{_{n}}al z\ensuremath{^{2}}(w). O atrator reconstruído tem razão
\prov{relações: parte\_de olho\_de\_koch; referencia kernel}
\prov{proveniência: paper \textperiodcentered{} fato \textperiodcentered{} confiança 1.0 \textperiodcentered{} domínio paper \textperiodcentered{} história 6 versões no jornal}

%CRISTAL {"arestas":[{"alvo":"casl-propagation","confianca":-0.094,"epistemico":"fato","origem":"manual","rel":"referencia"},{"alvo":"bai","confianca":0.156,"epistemico":"fato","origem":"manual","rel":"referencia"}],"confianca":1.0,"contraexemplos":[],"descricao":"ap:fisica I Leitura heurística --- não entra nas provas das Partes~I--II. como Lyapunov discreta; órbitas dissipativas; analogia de Sitter conservativo vs. delegação com atrito. Detalhes abaixo (texto histórico condensado). apfisicasnippet.tex","epistemico":"fato","exemplos":[],"id":"plano_no_so","memoria":"semantica","meta":{"dominio":"paper","nivel":5},"origem":"paper","palavras_chave":[],"sinonimos":[],"tipo":"conceito","titulo":"Plano no SO"}
\section{Plano no SO}
ap:fisica I Leitura heurística --- não entra nas provas das Partes\~{}I--II. como Lyapunov discreta; órbitas dissipativas; analogia de Sitter conservativo vs. delegação com atrito. Detalhes abaixo (texto histórico condensado). apfisicasnippet.tex
\prov{relações: referencia casl-propagation; referencia bai}
\prov{proveniência: paper \textperiodcentered{} fato \textperiodcentered{} confiança 1.0 \textperiodcentered{} domínio paper \textperiodcentered{} história 10 versões no jornal}

%CRISTAL {"arestas":[{"alvo":"corpo_tropical","confianca":1.0,"epistemico":"fato","origem":"paper","rel":"parte_de"}],"confianca":1.0,"contraexemplos":[],"descricao":"Há um motivo --- e é um teorema --- para o gap estar no centro. Ele é exatamente o que separa um sistema onde algo acontece de um sistema inerte. Volte a Aₙ agindo no plano hiperbólico e ao vão (=n²+4); tudo depende do sinal de. o gap é o que faz mover Para uma transformação A não trivial do plano hiperbólico, são equivalentes: itemize1pt o vão do gap é positivo, >0; A é hiperbólica: dois pontos fixos distintos no horizonte; existe o eixo --- a geodésica entre esses po","epistemico":"fato","exemplos":[],"id":"poincare_aberto_pela_agulha","memoria":"semantica","meta":{"dominio":"paper","nivel":5,"verificacao":"slug idêntico — enriquecer, não duplicar"},"origem":"paper","palavras_chave":[],"sinonimos":[],"tipo":"conceito","titulo":"Poincaré, aberto pela agulha"}
\section{Poincaré, aberto pela agulha}
Há um motivo --- e é um teorema --- para o gap estar no centro. Ele é exatamente o que separa um sistema onde algo acontece de um sistema inerte. Volte a A\ensuremath{_{n}} agindo no plano hiperbólico e ao vão (=n\ensuremath{^{2}}+4); tudo depende do sinal de. o gap é o que faz mover Para uma transformação A não trivial do plano hiperbólico, são equivalentes: itemize1pt o vão do gap é positivo, >0; A é hiperbólica: dois pontos fixos distintos no horizonte; existe o eixo --- a geodésica entre esses po
\prov{relações: parte\_de corpo\_tropical}
\prov{proveniência: paper \textperiodcentered{} fato \textperiodcentered{} confiança 1.0 \textperiodcentered{} domínio paper \textperiodcentered{} história 4 versões no jornal}

%CRISTAL {"arestas":[{"alvo":"gentil_fundamentos_dos_numeros","confianca":1.0,"epistemico":"fato","origem":"paper","rel":"relaciona"},{"alvo":"reais","confianca":1.0,"epistemico":"fato","origem":"paper","rel":"relaciona"},{"alvo":"a_engenharia_de_ouro_pa_pgs_agm_fc_reta","confianca":1.0,"epistemico":"fato","origem":"paper","rel":"relaciona"},{"alvo":"criterio_de_nao-isomorfismo","confianca":1.0,"epistemico":"fato","origem":"paper","rel":"relaciona"},{"alvo":"transformada_dourada","confianca":1.0,"epistemico":"fato","origem":"paper","rel":"relaciona"},{"alvo":"paper_2_reta_ouro","confianca":1.0,"epistemico":"fato","origem":"paper","rel":"relaciona"},{"alvo":"corpo_tropical","confianca":1.0,"epistemico":"fato","origem":"paper","rel":"relaciona"},{"alvo":"paper_45_confluencia_ouro","confianca":1.0,"epistemico":"fato","origem":"paper","rel":"relaciona"},{"alvo":"paper_bridge_areas","confianca":1.0,"epistemico":"fato","origem":"paper","rel":"relaciona"},{"alvo":"gato","confianca":1.0,"epistemico":"fato","origem":"paper","rel":"relaciona"},{"alvo":"teoria_dos_numeros","confianca":1.0,"epistemico":"fato","origem":"paper","rel":"parte_de"},{"alvo":"matematica","confianca":1.0,"epistemico":"fato","origem":"manual","rel":"parte_de"}],"confianca":1.0,"contraexemplos":[],"descricao":"Paper de fundação: chart simbólico Lopes (Cauchy/Cantor) ↔ engenharia PA/PGS/AGM/FC/reta/𝒢; critério de (não-)isomorfismo.","epistemico":"fato","exemplos":["neuronio/verificar_ponte_gentil.py — 18/18 testes gato→Gentil","neuronio.c streaming ≡ matriz A=[[1,1],[1,0]]","neuronio/gentil_simbolico.py — Cauchy + Gentil2 em ℚ","lab.py reproduzir — 62 testes unificados"],"id":"ponte_simbolica_gentil","memoria":"semantica","meta":{"dominio":"paper","nivel":5,"serve":"Gargalo conceitual Tiffany — liga Fundamentos à engenharia de ouro","verificacao":"slug idêntico — enriquecer, não duplicar"},"origem":"paper","palavras_chave":[],"sinonimos":["ponte simbólica gentil","a ponte simbólica"],"tipo":"conceito","titulo":"ponte simbolica gentil"}
\section{ponte simbolica gentil}
Paper de fundação: chart simbólico Lopes (Cauchy/Cantor) \ensuremath{\leftrightarrow} engenharia PA/PGS/AGM/FC/reta/\ensuremath{\mathcal{G}}; critério de (não-)isomorfismo.
\prov{exemplos: neuronio/verificar\_ponte\_gentil.py — 18/18 testes gato\ensuremath{\to}Gentil; neuronio.c streaming \ensuremath{\equiv} matriz A=[[1,1],[1,0]]; neuronio/gentil\_simbolico.py — Cauchy + Gentil2 em \ensuremath{\mathbb{Q}}; lab.py reproduzir — 62 testes unificados}
\prov{sinónimos: ponte simbólica gentil; a ponte simbólica}
\prov{relações: relaciona gentil\_fundamentos\_dos\_numeros; relaciona reais; relaciona a\_engenharia\_de\_ouro\_pa\_pgs\_agm\_fc\_reta; relaciona criterio\_de\_nao-isomorfismo; relaciona transformada\_dourada; relaciona paper\_2\_reta\_ouro; relaciona corpo\_tropical; relaciona paper\_45\_confluencia\_ouro; relaciona paper\_bridge\_areas; relaciona gato; parte\_de teoria\_dos\_numeros; parte\_de matematica}
\prov{proveniência: paper \textperiodcentered{} fato \textperiodcentered{} confiança 1.0 \textperiodcentered{} domínio paper \textperiodcentered{} história 397 versões no jornal}

%CRISTAL {"arestas":[{"alvo":"ciencias","confianca":0.047,"epistemico":"fato","origem":"manual","rel":"referencia"},{"alvo":"quatro_ciencias","confianca":1.0,"epistemico":"fato","origem":"paper","rel":"parte_de"}],"confianca":1.0,"contraexemplos":[],"descricao":"V A grade não é só descritiva: ela opera. Quando se manda o fluxo cru pela boca (torce), o dado já está escrito na base de Peirce e se assenta sozinho --- porque os projetores são uma resolução da identidade. O assentamento automático Como P_+P_=I (idempotentes ortogonais do gato, P²=P, P_ P_=0)~brocaos, qualquer dado torcido já é _ P_(dado): ninguém escolhe o setor, o dado tem peso em cada um e o dominante ressoa (Aⁿ/ⁿ P_). Sobre o fluxo cru, a álgebra separa por si","epistemico":"fato","exemplos":[],"id":"por_que_nao_se_escolhe_o_dado","memoria":"semantica","meta":{"dominio":"paper","nivel":5,"verificacao":"slug idêntico — enriquecer, não duplicar"},"origem":"paper","palavras_chave":[],"sinonimos":[],"tipo":"conceito","titulo":"Por que não se escolhe o dado"}
\section{Por que não se escolhe o dado}
V A grade não é só descritiva: ela opera. Quando se manda o fluxo cru pela boca (torce), o dado já está escrito na base de Peirce e se assenta sozinho --- porque os projetores são uma resolução da identidade. O assentamento automático Como P\_+P\_=I (idempotentes ortogonais do gato, P\ensuremath{^{2}}=P, P\_ P\_=0)\~{}brocaos, qualquer dado torcido já é \_ P\_(dado): ninguém escolhe o setor, o dado tem peso em cada um e o dominante ressoa (A\ensuremath{^{n}}/\ensuremath{^{n}} P\_). Sobre o fluxo cru, a álgebra separa por si
\prov{relações: referencia ciencias; parte\_de quatro\_ciencias}
\prov{proveniência: paper \textperiodcentered{} fato \textperiodcentered{} confiança 1.0 \textperiodcentered{} domínio paper \textperiodcentered{} história 10 versões no jornal}

%CRISTAL {"arestas":[{"alvo":"pow_gato","confianca":1.0,"epistemico":"fato","origem":"paper","rel":"parte_de"},{"alvo":"minerar_e_resfriar","confianca":1.0,"epistemico":"fato","origem":"wiki","rel":"relaciona"}],"confianca":1.0,"contraexemplos":[],"descricao":"empty !85PoW é busca forward no espaço de nonces --- não inversão de. A premissa criptográfica (T): preimage de SHA-256 é computacionalmente inviável; nenhuma construção que a quebre é conhecida. abstract Modelamos dois PoW compatíveis: (i)~PoW- (Bitcoin): ²(header n)<T; (ii)~PoW-gato (nativo): espectro de rank-6 (estrela de prata). A mina nativa resolve por colapso geométrico ( V); minagato.py mede só o baseline PoW- (ASIC). Postulamos a vantagem energética","epistemico":"fato","exemplos":[],"id":"pow_classico_bitcoin","memoria":"semantica","meta":{"dominio":"paper","nivel":5,"verificacao":"conceito novo"},"origem":"paper","palavras_chave":[],"sinonimos":[],"tipo":"conceito","titulo":"PoW clássico (Bitcoin)"}
\section{PoW clássico (Bitcoin)}
empty !85PoW é busca forward no espaço de nonces --- não inversão de. A premissa criptográfica (T): preimage de SHA-256 é computacionalmente inviável; nenhuma construção que a quebre é conhecida. abstract Modelamos dois PoW compatíveis: (i)\~{}PoW- (Bitcoin): \ensuremath{^{2}}(header n)<T; (ii)\~{}PoW-gato (nativo): espectro de rank-6 (estrela de prata). A mina nativa resolve por colapso geométrico ( V); minagato.py mede só o baseline PoW- (ASIC). Postulamos a vantagem energética
\prov{relações: parte\_de pow\_gato; relaciona minerar\_e\_resfriar}
\prov{proveniência: paper \textperiodcentered{} fato \textperiodcentered{} confiança 1.0 \textperiodcentered{} domínio paper \textperiodcentered{} história 4 versões no jornal}

%CRISTAL {"arestas":[{"alvo":"gato","confianca":1.0,"epistemico":"fato","origem":"paper","rel":"confirma"},{"alvo":"transcricao","confianca":0.5,"epistemico":"fato","origem":"wiki","rel":"relaciona"},{"alvo":"prova_isa","confianca":0.5,"epistemico":"fato","origem":"wiki","rel":"relaciona"},{"alvo":"readme","confianca":0.5,"epistemico":"fato","origem":"wiki","rel":"relaciona"},{"alvo":"teste_ordem","confianca":0.5,"epistemico":"fato","origem":"wiki","rel":"relaciona"},{"alvo":"python","confianca":0.5,"epistemico":"fato","origem":"wiki","rel":"relaciona"},{"alvo":"sessao_interativa","confianca":0.5,"epistemico":"fato","origem":"wiki","rel":"relaciona"},{"alvo":"word","confianca":0.5,"epistemico":"fato","origem":"wiki","rel":"relaciona"},{"alvo":"syscall","confianca":0.5,"epistemico":"fato","origem":"wiki","rel":"relaciona"},{"alvo":"virtualizacao","confianca":0.5,"epistemico":"fato","origem":"wiki","rel":"relaciona"}],"confianca":1.0,"contraexemplos":[],"descricao":"Paper: pow_gato.tex","epistemico":"fato","exemplos":[],"id":"pow_gato","memoria":"semantica","meta":{"dominio":"paper","nivel":5,"tipo_no":"paper"},"origem":"paper","palavras_chave":[],"sinonimos":[],"tipo":"conceito","titulo":"pow_gato"}
\section{pow\_gato}
Paper: pow\_gato.tex
\prov{relações: confirma gato; relaciona transcricao; relaciona prova\_isa; relaciona readme; relaciona teste\_ordem; relaciona python; relaciona sessao\_interativa; relaciona word; relaciona syscall; relaciona virtualizacao}
\prov{proveniência: paper \textperiodcentered{} fato \textperiodcentered{} confiança 1.0 \textperiodcentered{} domínio paper \textperiodcentered{} história 40 versões no jornal}

%CRISTAL {"arestas":[{"alvo":"floco_de_neve","confianca":1.0,"epistemico":"fato","origem":"paper","rel":"parte_de"}],"confianca":1.0,"contraexemplos":[],"descricao":"V O experimento koch-memoria refutou Koch como codificador de C₃ ou C_: (w,endereco) 0,12; 16 bits entre vizinhos; 18,548 colisões de endereco com c² distinto. defi[Floco ] Chamamos de floco () o objecto que armazena a informação não-associativa produzida pela dissipação quando o estado canónico atravessa a cadeia de retas do gato: [ (w):=(endereco, atrator, zₖoch, orbita) =kochcoord(w,(w)), ] em que zₖoch é a órbita zcaₙal z² em [i]/p[i] (","epistemico":"fato","exemplos":[],"id":"prata_cifra_e_rastreabilidade_no_floco","memoria":"semantica","meta":{"dominio":"paper","nivel":5,"verificacao":"conceito novo"},"origem":"paper","palavras_chave":[],"sinonimos":[],"tipo":"conceito","titulo":"Prata, cifra e rastreabilidade no floco"}
\section{Prata, cifra e rastreabilidade no floco}
V O experimento koch-memoria refutou Koch como codificador de C\ensuremath{_{3}} ou C\_: (w,endereco) 0,12; 16 bits entre vizinhos; 18,548 colisões de endereco com c\ensuremath{^{2}} distinto. defi[Floco ] Chamamos de floco () o objecto que armazena a informação não-associativa produzida pela dissipação quando o estado canónico atravessa a cadeia de retas do gato: [ (w):=(endereco, atrator, z\ensuremath{_{k}}och, orbita) =kochcoord(w,(w)), ] em que z\ensuremath{_{k}}och é a órbita zca\ensuremath{_{n}}al z\ensuremath{^{2}} em [i]/p[i] (
\prov{relações: parte\_de floco\_de\_neve}
\prov{proveniência: paper \textperiodcentered{} fato \textperiodcentered{} confiança 1.0 \textperiodcentered{} domínio paper \textperiodcentered{} história 3 versões no jornal}

%CRISTAL {"arestas":[{"alvo":"casl-propagation","confianca":0.086,"epistemico":"fato","origem":"manual","rel":"referencia"},{"alvo":"colapso","confianca":-0.016,"epistemico":"fato","origem":"manual","rel":"referencia"}],"confianca":1.0,"contraexemplos":[],"descricao":"definition[Atenuação e propagação] Sobre extents: T:2X2X é monótono deflacionário (T, x y T(x) T(y)). Sobre detentores: P:2 P2 P é monótono inflacionário (P). Re-delegar por saltos é P; = corresponde ao menor ponto fixo P=ₙ Pⁿ(H₀). definition definition[Conjunto vigente e colapso ] Seja =ᵢ capacidades detidas, filtradas por vigência e por cᵢ para consulta q. Defina [ _=:_>0 ou avaliação local, (q)=","epistemico":"fato","exemplos":[],"id":"preliminares_casl","memoria":"semantica","meta":{"dominio":"paper","nivel":5},"origem":"paper","palavras_chave":[],"sinonimos":[],"tipo":"conceito","titulo":"Preliminares CASL"}
\section{Preliminares CASL}
definition[Atenuação e propagação] Sobre extents: T:2X2X é monótono deflacionário (T, x y T(x) T(y)). Sobre detentores: P:2 P2 P é monótono inflacionário (P). Re-delegar por saltos é P; = corresponde ao menor ponto fixo P=\ensuremath{_{n}} P\ensuremath{^{n}}(H\ensuremath{_{0}}). definition definition[Conjunto vigente e colapso ] Seja =\ensuremath{_{i}} capacidades detidas, filtradas por vigência e por c\ensuremath{_{i}} para consulta q. Defina [ \_=:\_>0 ou avaliação local, (q)=
\prov{relações: referencia casl-propagation; referencia colapso}
\prov{proveniência: paper \textperiodcentered{} fato \textperiodcentered{} confiança 1.0 \textperiodcentered{} domínio paper \textperiodcentered{} história 9 versões no jornal}

%CRISTAL {"arestas":[{"alvo":"assistente","confianca":1.0,"epistemico":"fato","origem":"paper","rel":"parte_de"},{"alvo":"fase3","confianca":1.0,"epistemico":"fato","origem":"paper","rel":"parte_de"}],"confianca":1.0,"contraexemplos":[],"descricao":"1.25 tabular@ll@ & abertos “como faço rebase no git?” & software/git.idx “permissão negada no systemd” & sistemas/linux.idx “compilar paper LaTeX” & documentacao/latex.idx “nossa API de faturamento” & cliente/empresa/produto.idx “docker no fedora com git” & software/docker.idx + sistemas/fedora.idx tabular Fusão multi-tecido: executar tiffanyquery em cada.idx; ordenar hits por overlap; retornar melhor linha --- sem alterar; lógica no run","epistemico":"fato","exemplos":[],"id":"produto_comercial","memoria":"semantica","meta":{"dominio":"paper","nivel":5,"verificacao":"slug idêntico — enriquecer, não duplicar"},"origem":"paper","palavras_chave":[],"sinonimos":[],"tipo":"conceito","titulo":"Produto comercial"}
\section{Produto comercial}
1.25 tabular@ll@ \& abertos “como faço rebase no git?” \& software/git.idx “permissão negada no systemd” \& sistemas/linux.idx “compilar paper LaTeX” \& documentacao/latex.idx “nossa API de faturamento” \& cliente/empresa/produto.idx “docker no fedora com git” \& software/docker.idx + sistemas/fedora.idx tabular Fusão multi-tecido: executar tiffanyquery em cada.idx; ordenar hits por overlap; retornar melhor linha --- sem alterar; lógica no run
\prov{relações: parte\_de assistente; parte\_de fase3}
\prov{proveniência: paper \textperiodcentered{} fato \textperiodcentered{} confiança 1.0 \textperiodcentered{} domínio paper \textperiodcentered{} história 5 versões no jornal}

%CRISTAL {"arestas":[{"alvo":"olho_de_koch","confianca":1.0,"epistemico":"fato","origem":"paper","rel":"parte_de"}],"confianca":1.0,"contraexemplos":[],"descricao":"D No, o neurônio 1-bit produz w=[,total,e,]. Duas células brancas naturais: [ e₀: estado vivo w, e₁: estado ouro gold(w), ] com gold a iteração Fibonacci do gato (cifra A₁). O canal negro E₀₁ não soma w+gold(w) --- isso contaminaria as faces. O canal isola o bloco Peirce [ e₀,M,e₁, ] cuja coordenada observável é a parábola entre os dois estados: [ a+c²=1, a=calor (Hellinger), c²=fase (Bhattacharyya). ] V Implementado em parabolaauditar e kochcoordenadas","epistemico":"fato","exemplos":[],"id":"projecao_por_ressonancia","memoria":"semantica","meta":{"dominio":"paper","nivel":5},"origem":"paper","palavras_chave":[],"sinonimos":[],"tipo":"conceito","titulo":"Projeção por ressonância"}
\section{Projeção por ressonância}
D No, o neurônio 1-bit produz w=[,total,e,]. Duas células brancas naturais: [ e\ensuremath{_{0}}: estado vivo w, e\ensuremath{_{1}}: estado ouro gold(w), ] com gold a iteração Fibonacci do gato (cifra A\ensuremath{_{1}}). O canal negro E\ensuremath{_{01}} não soma w+gold(w) --- isso contaminaria as faces. O canal isola o bloco Peirce [ e\ensuremath{_{0}},M,e\ensuremath{_{1}}, ] cuja coordenada observável é a parábola entre os dois estados: [ a+c\ensuremath{^{2}}=1, a=calor (Hellinger), c\ensuremath{^{2}}=fase (Bhattacharyya). ] V Implementado em parabolaauditar e kochcoordenadas
\prov{relações: parte\_de olho\_de\_koch}
\prov{proveniência: paper \textperiodcentered{} fato \textperiodcentered{} confiança 1.0 \textperiodcentered{} domínio paper \textperiodcentered{} história 3 versões no jornal}

%CRISTAL {"arestas":[{"alvo":"colapso","confianca":0.95,"epistemico":"fato","origem":"paper","rel":"parte_de"}],"confianca":1.0,"contraexemplos":[],"descricao":"[Auto-calibração] Dado um bloco b, seja (r)=|u²,r b|² a energia de curvatura da imagem relativística amostrada em escala rk. A escala intrínseca é r_*=r(r), determinada pelos próprios bits; a reta de referência é r_*k. Para dados auto-similares de razão, r_*=; quando a auto-similaridade é a do vazio, r_*=.","epistemico":"fato","exemplos":[],"id":"proposition_auto-calibracaopropreta_dado_um_bloco_b_seja","memoria":"semantica","meta":{"dominio":"paper","nivel":5,"tipo_teorema":"proposition","verificacao":"conceito novo"},"origem":"paper","palavras_chave":[],"sinonimos":[],"tipo":"conceito","titulo":"Auto-calibração"}
\section{Auto-calibração}
[Auto-calibração] Dado um bloco b, seja (r)=|u\ensuremath{^{2}},r b|\ensuremath{^{2}} a energia de curvatura da imagem relativística amostrada em escala rk. A escala intrínseca é r\_*=r(r), determinada pelos próprios bits; a reta de referência é r\_*k. Para dados auto-similares de razão, r\_*=; quando a auto-similaridade é a do vazio, r\_*=.
\prov{relações: parte\_de colapso}
\prov{proveniência: paper \textperiodcentered{} fato \textperiodcentered{} confiança 1.0 \textperiodcentered{} domínio paper \textperiodcentered{} história 5 versões no jornal}

%CRISTAL {"arestas":[{"alvo":"plano_no_so","confianca":1.0,"epistemico":"fato","origem":"paper","rel":"prova"}],"confianca":1.0,"contraexemplos":[],"descricao":"[Conversa fail-closed] Se nenhum nó dialogos casa (fala,ultimo) com score acima do limiar de proibições, responder() devolve None. Não há fallback para LLM no motor conversa --- alinhado ao fail-closed de Ep=.","epistemico":"fato","exemplos":[],"id":"proposition_conversa_fail-closedproptiffany-fail_se_nenhu","memoria":"semantica","meta":{"dominio":"paper","nivel":5,"tipo_teorema":"proposition"},"origem":"paper","palavras_chave":[],"sinonimos":[],"tipo":"conceito","titulo":"proposition: [Conversa fail-closed]prop:tiffany-fail\n  Se nenhu…"}
\section{proposition: [Conversa fail-closed]prop:tiffany-fail
  Se nenhu…}
[Conversa fail-closed] Se nenhum nó dialogos casa (fala,ultimo) com score acima do limiar de proibições, responder() devolve None. Não há fallback para LLM no motor conversa --- alinhado ao fail-closed de Ep=.
\prov{relações: prova plano\_no\_so}
\prov{proveniência: paper \textperiodcentered{} fato \textperiodcentered{} confiança 1.0 \textperiodcentered{} domínio paper \textperiodcentered{} história 3 versões no jornal}

%CRISTAL {"arestas":[{"alvo":"colapso","confianca":0.95,"epistemico":"fato","origem":"paper","rel":"parte_de"},{"alvo":"virtualizacao","confianca":0.5,"epistemico":"fato","origem":"wiki","rel":"relaciona"},{"alvo":"word","confianca":0.5,"epistemico":"fato","origem":"wiki","rel":"relaciona"},{"alvo":"vida-morte","confianca":0.5,"epistemico":"fato","origem":"wiki","rel":"relaciona"},{"alvo":"proposition_gerador_de_escala_ipropderiv_com_tddt","confianca":0.5,"epistemico":"fato","origem":"wiki","rel":"relaciona"}],"confianca":1.0,"contraexemplos":[],"descricao":"[Convolução multiplicativa produto] [x y]=[x][y].","epistemico":"fato","exemplos":[],"id":"proposition_convolucao_multiplicativa_produtopropconv","memoria":"semantica","meta":{"dominio":"paper","nivel":5,"tipo_teorema":"proposition","verificacao":"conceito novo"},"origem":"paper","palavras_chave":[],"sinonimos":[],"tipo":"conceito","titulo":"Convolução multiplicativa produto"}
\section{Convolução multiplicativa produto}
[Convolução multiplicativa produto] [x y]=[x][y].
\prov{relações: parte\_de colapso; relaciona virtualizacao; relaciona word; relaciona vida-morte; relaciona proposition\_gerador\_de\_escala\_ipropderiv\_com\_tddt}
\prov{proveniência: paper \textperiodcentered{} fato \textperiodcentered{} confiança 1.0 \textperiodcentered{} domínio paper \textperiodcentered{} história 9 versões no jornal}

%CRISTAL {"arestas":[{"alvo":"plano_no_so","confianca":1.0,"epistemico":"fato","origem":"paper","rel":"prova"}],"confianca":1.0,"contraexemplos":[],"descricao":"[Deny-overrides --- a proibição sempre vence] A escolha de política está embutida na diferença de conjuntos: para toda linha r, [ r Ep(s,a) (grant g: rg) (proibição d: rd). ] Uma proibição domina qualquer grant sobre a mesma linha, independentemente de “especificidade” ou ordem. Assim, grant customerId=5 com proibição =A (e 5 A) resulta: a proibição vence. O sistema é, por definição, deny-overrides --- declarado aqui para remover ambiguidade.","epistemico":"fato","exemplos":[],"id":"proposition_deny-overrides_---_a_proibicao_sempre_vence_prop","memoria":"semantica","meta":{"dominio":"paper","nivel":5,"tipo_teorema":"proposition"},"origem":"paper","palavras_chave":[],"sinonimos":[],"tipo":"conceito","titulo":"proposition: [Deny-overrides --- a proibição sempre vence]\nprop…"}
\section{proposition: [Deny-overrides --- a proibição sempre vence]
prop…}
[Deny-overrides --- a proibição sempre vence] A escolha de política está embutida na diferença de conjuntos: para toda linha r, [ r Ep(s,a) (grant g: rg) (proibição d: rd). ] Uma proibição domina qualquer grant sobre a mesma linha, independentemente de “especificidade” ou ordem. Assim, grant customerId=5 com proibição =A (e 5 A) resulta: a proibição vence. O sistema é, por definição, deny-overrides --- declarado aqui para remover ambiguidade.
\prov{relações: prova plano\_no\_so}
\prov{proveniência: paper \textperiodcentered{} fato \textperiodcentered{} confiança 1.0 \textperiodcentered{} domínio paper \textperiodcentered{} história 3 versões no jornal}

%CRISTAL {"arestas":[{"alvo":"colapso","confianca":0.95,"epistemico":"fato","origem":"paper","rel":"parte_de"},{"alvo":"virtualizacao","confianca":0.5,"epistemico":"fato","origem":"wiki","rel":"relaciona"},{"alvo":"vida-morte","confianca":0.5,"epistemico":"fato","origem":"wiki","rel":"relaciona"},{"alvo":"word","confianca":0.5,"epistemico":"fato","origem":"wiki","rel":"relaciona"}],"confianca":1.0,"contraexemplos":[],"descricao":"[Escala fase] [x()]()=i,[x](), >0.","epistemico":"fato","exemplos":[],"id":"proposition_escala_fasepropescala_xix","memoria":"semantica","meta":{"dominio":"paper","nivel":5,"tipo_teorema":"proposition","verificacao":"conceito novo"},"origem":"paper","palavras_chave":[],"sinonimos":[],"tipo":"conceito","titulo":"Escala fase"}
\section{Escala fase}
[Escala fase] [x()]()=i,[x](), >0.
\prov{relações: parte\_de colapso; relaciona virtualizacao; relaciona vida-morte; relaciona word}
\prov{proveniência: paper \textperiodcentered{} fato \textperiodcentered{} confiança 1.0 \textperiodcentered{} domínio paper \textperiodcentered{} história 8 versões no jornal}

%CRISTAL {"arestas":[{"alvo":"colapso","confianca":0.95,"epistemico":"fato","origem":"paper","rel":"parte_de"},{"alvo":"virtualizacao","confianca":0.5,"epistemico":"fato","origem":"wiki","rel":"relaciona"},{"alvo":"word","confianca":0.5,"epistemico":"fato","origem":"wiki","rel":"relaciona"},{"alvo":"vida-morte","confianca":0.5,"epistemico":"fato","origem":"wiki","rel":"relaciona"},{"alvo":"semiotica_peirce_triade","confianca":0.5,"epistemico":"fato","origem":"wiki","rel":"relaciona"},{"alvo":"psicologia","confianca":0.5,"epistemico":"fato","origem":"wiki","rel":"relaciona"},{"alvo":"psicanalise","confianca":0.5,"epistemico":"fato","origem":"wiki","rel":"relaciona"},{"alvo":"universidade_curriculo_em_camadas","confianca":0.5,"epistemico":"fato","origem":"wiki","rel":"relaciona"}],"confianca":1.0,"contraexemplos":[],"descricao":"[Gerador de escala i] Com =t,ddt, [ x]()=i,[x]().","epistemico":"fato","exemplos":[],"id":"proposition_gerador_de_escala_ipropderiv_com_tddt","memoria":"semantica","meta":{"dominio":"paper","nivel":5,"tipo_teorema":"proposition","verificacao":"conceito novo"},"origem":"paper","palavras_chave":[],"sinonimos":[],"tipo":"conceito","titulo":"Gerador de escala i"}
\section{Gerador de escala i}
[Gerador de escala i] Com =t,ddt, [ x]()=i,[x]().
\prov{relações: parte\_de colapso; relaciona virtualizacao; relaciona word; relaciona vida-morte; relaciona semiotica\_peirce\_triade; relaciona psicologia; relaciona psicanalise; relaciona universidade\_curriculo\_em\_camadas}
\prov{proveniência: paper \textperiodcentered{} fato \textperiodcentered{} confiança 1.0 \textperiodcentered{} domínio paper \textperiodcentered{} história 12 versões no jornal}

%CRISTAL {"arestas":[{"alvo":"colapso","confianca":0.95,"epistemico":"fato","origem":"paper","rel":"parte_de"}],"confianca":1.0,"contraexemplos":[],"descricao":"[Round-trip exato] Para todo b0,1N, ⁻¹[b]=b (Teorema~); preservada a ordem, a reconstrução é exata --- perda zero.","epistemico":"fato","exemplos":[],"id":"proposition_round-trip_exatoproprt_para_todo_b01n","memoria":"semantica","meta":{"dominio":"paper","nivel":5,"tipo_teorema":"proposition","verificacao":"conceito novo"},"origem":"paper","palavras_chave":[],"sinonimos":[],"tipo":"conceito","titulo":"Round-trip exato"}
\section{Round-trip exato}
[Round-trip exato] Para todo b0,1N, \ensuremath{^{-1}}[b]=b (Teorema\~{}); preservada a ordem, a reconstrução é exata --- perda zero.
\prov{relações: parte\_de colapso}
\prov{proveniência: paper \textperiodcentered{} fato \textperiodcentered{} confiança 1.0 \textperiodcentered{} domínio paper \textperiodcentered{} história 5 versões no jornal}

%CRISTAL {"arestas":[{"alvo":"kappa","confianca":0.031,"epistemico":"fato","origem":"manual","rel":"referencia"},{"alvo":"casl-propagation","confianca":0.039,"epistemico":"fato","origem":"manual","rel":"referencia"},{"alvo":"bai","confianca":-0.039,"epistemico":"fato","origem":"manual","rel":"referencia"}],"confianca":1.0,"contraexemplos":[],"descricao":"D Instanciação Parte~IV: X= nós do tecido, = turno A — B, como medida de organização (Definição~). A conversa não fabrica texto --- aplica. Colisões e ecos genéricos são proibições (deny-overrides). Medição (): conversa/medeflocos.py --- frames, atratores, recipiente XOR; acertos BAI por cadeia em cadeias.py --rodar --mede. definition[Capacidade conversacional] Specialização de (Definição~)","epistemico":"fato","exemplos":[],"id":"quadro_de_isomorfismos","memoria":"semantica","meta":{"dominio":"paper","nivel":5},"origem":"paper","palavras_chave":[],"sinonimos":[],"tipo":"conceito","titulo":"Quadro de isomorfismos"}
\section{Quadro de isomorfismos}
D Instanciação Parte\~{}IV: X= nós do tecido, = turno A — B, como medida de organização (Definição\~{}). A conversa não fabrica texto --- aplica. Colisões e ecos genéricos são proibições (deny-overrides). Medição (): conversa/medeflocos.py --- frames, atratores, recipiente XOR; acertos BAI por cadeia em cadeias.py --rodar --mede. definition[Capacidade conversacional] Specialização de (Definição\~{})
\prov{relações: referencia kappa; referencia casl-propagation; referencia bai}
\prov{proveniência: paper \textperiodcentered{} fato \textperiodcentered{} confiança 1.0 \textperiodcentered{} domínio paper \textperiodcentered{} história 12 versões no jornal}

%CRISTAL {"arestas":[{"alvo":"reais","confianca":1.0,"epistemico":"fato","origem":"curriculo","rel":"especializa"}],"confianca":1.0,"contraexemplos":[],"descricao":"D A regra operacional que guia a revisão: defi[Disco vs RAM] Dados (arquivos, tecido, impressões indexadas) ficam no disco. Gatos (neurônios, estado 1) ficam na RAM --- um bit por neurônio. O trânsito é ping-pong: lê um pedaço do disco atualiza bits na RAM descarta o pedaço. Nada de buffer de conteúdo, cache espectral (psicache.npz), ou carregar o tecido inteiro. defi 1.3 tabular@lll@ ça & & (gatos) neuronio.c & arquivo byte a byte & e, o (contadores) recon","epistemico":"fato","exemplos":[],"id":"racionais","memoria":"semantica","meta":{"dominio":"paper","duplicata_sugerida":[],"nivel":5,"verificacao":"parte de conceito mais amplo 'racionais'"},"origem":"paper","palavras_chave":[],"sinonimos":["ℚ"],"tipo":"conceito","titulo":"racionais"}
\section{racionais}
D A regra operacional que guia a revisão: defi[Disco vs RAM] Dados (arquivos, tecido, impressões indexadas) ficam no disco. Gatos (neurônios, estado 1) ficam na RAM --- um bit por neurônio. O trânsito é ping-pong: lê um pedaço do disco atualiza bits na RAM descarta o pedaço. Nada de buffer de conteúdo, cache espectral (psicache.npz), ou carregar o tecido inteiro. defi 1.3 tabular@lll@ ça \& \& (gatos) neuronio.c \& arquivo byte a byte \& e, o (contadores) recon
\prov{sinónimos: \ensuremath{\mathbb{Q}}}
\prov{relações: especializa reais}
\prov{proveniência: paper \textperiodcentered{} fato \textperiodcentered{} confiança 1.0 \textperiodcentered{} domínio paper \textperiodcentered{} história 8 versões no jornal}

%CRISTAL {"arestas":[{"alvo":"olho_de_koch","confianca":1.0,"epistemico":"fato","origem":"paper","rel":"parte_de"}],"confianca":1.0,"contraexemplos":[],"descricao":"D A quantização CK não actua na soma das células brancas. Actua no canal zcaₙal --- embutimento de (c,a) da parábola em [i] --- pelo ribossomo [ z z² [i]. ] z z² é 2-para-1, logo A=0 Identifica z e z; nenhuma matriz linear representa a quadração. Quantização e segurança são a mesma operação: a guarda é a ausência de reversão, não um filtro acrescentado~transcricao. teo kochgauss aplica z² mod p só após canalressonancia; kochcoord constrói o endereç","epistemico":"fato","exemplos":[],"id":"reconstrucao_purificada","memoria":"semantica","meta":{"dominio":"paper","nivel":5},"origem":"paper","palavras_chave":[],"sinonimos":[],"tipo":"conceito","titulo":"Reconstrução purificada"}
\section{Reconstrução purificada}
D A quantização CK não actua na soma das células brancas. Actua no canal zca\ensuremath{_{n}}al --- embutimento de (c,a) da parábola em [i] --- pelo ribossomo [ z z\ensuremath{^{2}} [i]. ] z z\ensuremath{^{2}} é 2-para-1, logo A=0 Identifica z e z; nenhuma matriz linear representa a quadração. Quantização e segurança são a mesma operação: a guarda é a ausência de reversão, não um filtro acrescentado\~{}transcricao. teo kochgauss aplica z\ensuremath{^{2}} mod p só após canalressonancia; kochcoord constrói o endereç
\prov{relações: parte\_de olho\_de\_koch}
\prov{proveniência: paper \textperiodcentered{} fato \textperiodcentered{} confiança 1.0 \textperiodcentered{} domínio paper \textperiodcentered{} história 3 versões no jornal}

%CRISTAL {"arestas":[{"alvo":"corpo_dual","confianca":1.0,"epistemico":"fato","origem":"paper","rel":"parte_de"}],"confianca":1.0,"contraexemplos":[],"descricao":"A segurança é a segunda leitura do lado irreversível a: não-reversão por design, não por dureza computacional. A matriz de ataque é o nilpotente do Glacial (A=0), a torção é bloco-diagonal ([T,Pᵢ]=0), e a desigualdade de processamento contrai --- o que se apagou não se reconstrói, relativo ao modelo de ameaça. O calor e a segurança são a mesma irreversibilidade, vista duas vezes. Falta transmitir: o canal.","epistemico":"fato","exemplos":[],"id":"redes","memoria":"semantica","meta":{"dominio":"paper","nivel":5,"verificacao":"overlap moderado — revisar manualmente"},"origem":"paper","palavras_chave":[],"sinonimos":[],"tipo":"conceito","titulo":"Redes"}
\section{Redes}
A segurança é a segunda leitura do lado irreversível a: não-reversão por design, não por dureza computacional. A matriz de ataque é o nilpotente do Glacial (A=0), a torção é bloco-diagonal ([T,P\ensuremath{_{i}}]=0), e a desigualdade de processamento contrai --- o que se apagou não se reconstrói, relativo ao modelo de ameaça. O calor e a segurança são a mesma irreversibilidade, vista duas vezes. Falta transmitir: o canal.
\prov{relações: parte\_de corpo\_dual}
\prov{proveniência: paper \textperiodcentered{} fato \textperiodcentered{} confiança 1.0 \textperiodcentered{} domínio paper \textperiodcentered{} história 3 versões no jornal}

%CRISTAL {"arestas":[{"alvo":"o_floco_e_a_conservacao_da_comunicacao","confianca":1.0,"epistemico":"fato","origem":"paper","rel":"parte_de"},{"alvo":"topologia","confianca":0.5,"epistemico":"fato","origem":"wiki","rel":"relaciona"},{"alvo":"rust","confianca":0.5,"epistemico":"fato","origem":"wiki","rel":"relaciona"},{"alvo":"utilitarismo","confianca":0.5,"epistemico":"fato","origem":"wiki","rel":"relaciona"},{"alvo":"teste_ordem","confianca":0.5,"epistemico":"fato","origem":"wiki","rel":"relaciona"},{"alvo":"testes","confianca":0.5,"epistemico":"fato","origem":"wiki","rel":"relaciona"},{"alvo":"resposta_a_pergunta_dificil","confianca":0.5,"epistemico":"fato","origem":"wiki","rel":"relaciona"}],"confianca":1.0,"contraexemplos":[],"descricao":"","epistemico":"fato","exemplos":[],"id":"regua_associativa_e_o_risco_de_vazamento","memoria":"semantica","meta":{"dominio":"paper","nivel":5,"verificacao":"conceito novo"},"origem":"paper","palavras_chave":[],"sinonimos":[],"tipo":"conceito","titulo":"Régua associativa e o risco de vazamento"}
\section{Régua associativa e o risco de vazamento}
\prov{relações: parte\_de o\_floco\_e\_a\_conservacao\_da\_comunicacao; relaciona topologia; relaciona rust; relaciona utilitarismo; relaciona teste\_ordem; relaciona testes; relaciona resposta\_a\_pergunta\_dificil}
\prov{proveniência: paper \textperiodcentered{} fato \textperiodcentered{} confiança 1.0 \textperiodcentered{} domínio paper \textperiodcentered{} história 9 versões no jornal}

%CRISTAL {"arestas":[{"alvo":"painel_estelar","confianca":1.0,"epistemico":"fato","origem":"paper","rel":"parte_de"}],"confianca":1.0,"contraexemplos":[],"descricao":"V Implementação: ula/painel\\_estelar.h, painel\\_estelar\\_exp.c.\n\nprop[Leitura única, três braços]\nPara nonce $n$, painel\\_estelar\\_de(ctx, n) devolve um registo com quatro zonas:\ncenter\ntabular{@{}ll@{}}\nnúcleo & $Q$, $a$, $c^2$, $P_$, $P_$ \\\\\ngeometria & $(r,c,)$, $$, $S^3$, $a^2 d^2$ \\\\\nmecânica & $c^2$, fase, reversível, $$ \\\\\ntermo & calor, $H=$, classe \\\\\ntabular\ncenter\nReprodução: make painel-estelar $$ PAINEL\\_ESTELAR.md, dados/painel\\_estelar.json,\natlas/painel\\_estelar/estrela\\_*.svg.\np","epistemico":"fato","exemplos":[],"id":"relacao_com_os_fundamentos","memoria":"semantica","meta":{"dominio":"paper","nivel":5,"verificacao":"overlap moderado — revisar manualmente"},"origem":"paper","palavras_chave":[],"sinonimos":[],"tipo":"conceito","titulo":"Relação com os fundamentos"}
\section{Relação com os fundamentos}
V Implementação: ula/painel\textbackslash{}\_estelar.h, painel\textbackslash{}\_estelar\textbackslash{}\_exp.c.

prop[Leitura única, três braços]
Para nonce \$n\$, painel\textbackslash{}\_estelar\textbackslash{}\_de(ctx, n) devolve um registo com quatro zonas:
center
tabular\{@\{\}ll@\{\}\}
núcleo \& \$Q\$, \$a\$, \$c\^{}2\$, \$P\_\$, \$P\_\$ \textbackslash{}\textbackslash{}
geometria \& \$(r,c,)\$, \$\$, \$S\^{}3\$, \$a\^{}2 d\^{}2\$ \textbackslash{}\textbackslash{}
mecânica \& \$c\^{}2\$, fase, reversível, \$\$ \textbackslash{}\textbackslash{}
termo \& calor, \$H=\$, classe \textbackslash{}\textbackslash{}
tabular
center
Reprodução: make painel-estelar \$\$ PAINEL\textbackslash{}\_ESTELAR.md, dados/painel\textbackslash{}\_estelar.json,
atlas/painel\textbackslash{}\_estelar/estrela\textbackslash{}\_*.svg.
p
\prov{relações: parte\_de painel\_estelar}
\prov{proveniência: paper \textperiodcentered{} fato \textperiodcentered{} confiança 1.0 \textperiodcentered{} domínio paper \textperiodcentered{} história 2 versões no jornal}

%CRISTAL {"arestas":[{"alvo":"main","confianca":1.0,"epistemico":"fato","origem":"paper","rel":"parte_de"}],"confianca":1.0,"contraexemplos":[],"descricao":"Não-forjaT A dimensão de Hausdorff é invariante por bi-Lipschitz (Falconer); logo não há homeomorfismo bi-Lipschitz entre fractais de dimensões distintas --- o único isomorfismo impossível é o métrico.teo Não-perdaT é álgebra de divisão: ab=0 a=0 b=0. Em dim 16 surgem divisores de zero, e por isso a Máquina para em.teo Separação G₂-equivarianteT A decomposição ²⁷=²₇²₁₄ V₇ é ortogonal e G₂-equivariante; o mapa u v u v tem imagem V₇ e núcleo. As 2","epistemico":"fato","exemplos":[],"id":"reparo_ativo_e_auto-reconstrucao","memoria":"semantica","meta":{"dominio":"paper","nivel":5,"verificacao":"slug idêntico — enriquecer, não duplicar"},"origem":"paper","palavras_chave":[],"sinonimos":[],"tipo":"conceito","titulo":"Reparo ativo e auto-reconstrução"}
\section{Reparo ativo e auto-reconstrução}
Não-forjaT A dimensão de Hausdorff é invariante por bi-Lipschitz (Falconer); logo não há homeomorfismo bi-Lipschitz entre fractais de dimensões distintas --- o único isomorfismo impossível é o métrico.teo Não-perdaT é álgebra de divisão: ab=0 a=0 b=0. Em dim 16 surgem divisores de zero, e por isso a Máquina para em.teo Separação G\ensuremath{_{2}}-equivarianteT A decomposição \ensuremath{^{27}}=\ensuremath{^{2}}\ensuremath{_{7}}\ensuremath{^{2}}\ensuremath{_{14}} V\ensuremath{_{7}} é ortogonal e G\ensuremath{_{2}}-equivariante; o mapa u v u v tem imagem V\ensuremath{_{7}} e núcleo. As 2
\prov{relações: parte\_de main}
\prov{proveniência: paper \textperiodcentered{} fato \textperiodcentered{} confiança 1.0 \textperiodcentered{} domínio paper \textperiodcentered{} história 5 versões no jornal}

%CRISTAL {"arestas":[{"alvo":"paper_bridge_areas","confianca":1.0,"epistemico":"fato","origem":"paper","rel":"parte_de"}],"confianca":1.0,"contraexemplos":[],"descricao":"O script neuronio/verificarpontegentil.py fecha a ponte operacionalmente: o streaming de neuronio.c e a matriz A=(smallmatrix1&1 1&0smallmatrix) do Corpo Tropical / cap.~10 são o mesmo símbolo. Gato = operador Lopes sobre (e,o) Seja A a matriz do gato de Arnold (x²=x+1). Então: [nosep,leftmargin=1.4em] A=-1, tr,A=1, autovalores, =-1/. Projetores espectrais P_+P_=I, P_ P_=0 (dois puros, tropical). neuronio.c","epistemico":"fato","exemplos":[],"id":"reproducao_tiffany","memoria":"semantica","meta":{"dominio":"paper","duplicata_sugerida":[],"nivel":5,"verificacao":"parte de conceito mais amplo 'reproducao_tiffany'"},"origem":"paper","palavras_chave":[],"sinonimos":[],"tipo":"conceito","titulo":"Reprodução Tiffany"}
\section{Reprodução Tiffany}
O script neuronio/verificarpontegentil.py fecha a ponte operacionalmente: o streaming de neuronio.c e a matriz A=(smallmatrix1\&1 1\&0smallmatrix) do Corpo Tropical / cap.\~{}10 são o mesmo símbolo. Gato = operador Lopes sobre (e,o) Seja A a matriz do gato de Arnold (x\ensuremath{^{2}}=x+1). Então: [nosep,leftmargin=1.4em] A=-1, tr,A=1, autovalores, =-1/. Projetores espectrais P\_+P\_=I, P\_ P\_=0 (dois puros, tropical). neuronio.c
\prov{relações: parte\_de paper\_bridge\_areas}
\prov{proveniência: paper \textperiodcentered{} fato \textperiodcentered{} confiança 1.0 \textperiodcentered{} domínio paper \textperiodcentered{} história 5 versões no jornal}

%CRISTAL {"arestas":[{"alvo":"goldchain","confianca":1.0,"epistemico":"fato","origem":"paper","rel":"parte_de"}],"confianca":1.0,"contraexemplos":[],"descricao":"C O coração é GoldChain.fecha(c, backend) em machine/so/goldchain.py. Cinco passos, na ordem exata do código: validar as chaves (N-qubits), coordenar a condição (hashlock já em c.condicao), liquidar no backend (ETH/BTC reais), cobrar (cunhar o próximo Pell, cifrado) e certificar encadeando. lstlisting def fecha(self, c, backend): # 1) VALIDAR — as chaves P, coletivamente (N-qubits); uma forja derruba o conjunto ok, n, _ = pv.validaₙ(c.partes) if not ok: return \"erro\": \"validacao","epistemico":"fato","exemplos":[],"id":"resgate_liquidacao","memoria":"semantica","meta":{"dominio":"paper","duplicata_sugerida":[],"nivel":5,"verificacao":"parte de conceito mais amplo 'resgate_liquidacao'"},"origem":"paper","palavras_chave":[],"sinonimos":[],"tipo":"conceito","titulo":"resgate_liquidacao"}
\section{resgate\_liquidacao}
C O coração é GoldChain.fecha(c, backend) em machine/so/goldchain.py. Cinco passos, na ordem exata do código: validar as chaves (N-qubits), coordenar a condição (hashlock já em c.condicao), liquidar no backend (ETH/BTC reais), cobrar (cunhar o próximo Pell, cifrado) e certificar encadeando. lstlisting def fecha(self, c, backend): \# 1) VALIDAR — as chaves P, coletivamente (N-qubits); uma forja derruba o conjunto ok, n, \_ = pv.valida\ensuremath{_{n}}(c.partes) if not ok: return "erro": "validacao
\prov{relações: parte\_de goldchain}
\prov{proveniência: paper \textperiodcentered{} fato \textperiodcentered{} confiança 1.0 \textperiodcentered{} domínio paper \textperiodcentered{} história 9 versões no jornal}

%CRISTAL {"arestas":[{"alvo":"corpo_tropical","confianca":1.0,"epistemico":"fato","origem":"paper","rel":"parte_de"}],"confianca":1.0,"contraexemplos":[],"descricao":"(Sobre a métrica circular de Gentil Lopes da Silva.) Antes de ler Riemann falta a peça algébrica que a torção fornece: a torre do Lopes, o flip de Wick, e o anel que ele revela ao desenrolar a lemniscata. A torre do Lopes (completa). A construção B de Lopes parte do plano complexo C (base, raio r=|(a,b)|) e lhe acrescenta uma altura c Rm: [ z₃=z₁z₂(1- c₁, c₂r₁r₂), c₃= c₁ r₂+ c₂ r₁+ c₁ c₂, |w|²=r²+| c|². ] A norma é multiplicativa exatamente quando o defeito","epistemico":"fato","exemplos":[],"id":"riemann_a_malha_no_cone_de_luz","memoria":"semantica","meta":{"dominio":"paper","nivel":5,"verificacao":"slug idêntico — enriquecer, não duplicar"},"origem":"paper","palavras_chave":[],"sinonimos":[],"tipo":"conceito","titulo":"Riemann: a malha no cone de luz"}
\section{Riemann: a malha no cone de luz}
(Sobre a métrica circular de Gentil Lopes da Silva.) Antes de ler Riemann falta a peça algébrica que a torção fornece: a torre do Lopes, o flip de Wick, e o anel que ele revela ao desenrolar a lemniscata. A torre do Lopes (completa). A construção B de Lopes parte do plano complexo C (base, raio r=|(a,b)|) e lhe acrescenta uma altura c Rm: [ z\ensuremath{_{3}}=z\ensuremath{_{1}}z\ensuremath{_{2}}(1- c\ensuremath{_{1}}, c\ensuremath{_{2}}r\ensuremath{_{1}}r\ensuremath{_{2}}), c\ensuremath{_{3}}= c\ensuremath{_{1}} r\ensuremath{_{2}}+ c\ensuremath{_{2}} r\ensuremath{_{1}}+ c\ensuremath{_{1}} c\ensuremath{_{2}}, |w|\ensuremath{^{2}}=r\ensuremath{^{2}}+| c|\ensuremath{^{2}}. ] A norma é multiplicativa exatamente quando o defeito
\prov{relações: parte\_de corpo\_tropical}
\prov{proveniência: paper \textperiodcentered{} fato \textperiodcentered{} confiança 1.0 \textperiodcentered{} domínio paper \textperiodcentered{} história 5 versões no jornal}

%CRISTAL {"arestas":[{"alvo":"contabilidade","confianca":1.0,"epistemico":"fato","origem":"paper","rel":"parte_de"}],"confianca":1.0,"contraexemplos":[],"descricao":"I O erro corrigido era confundir a folha com a lei de conservação. São coisas distintas: itemize3pt A folha é um tecido ordenado --- os pagamentos como bins, ordenados pela parte real. Não é regida por a+c²=1. É só um tecido, como qualquer outro, que se lê e se escreve pela banda. O lastro (o balanço) é a+c²=1 --- a conservação da norma no colapso, a relação entre reserva e crédito. É a lei de conservação; a folha não. D Misturar as duas levava a trata","epistemico":"fato","exemplos":[],"id":"satisfaz_o_banco","memoria":"semantica","meta":{"dominio":"paper","nivel":5,"verificacao":"conceito novo"},"origem":"paper","palavras_chave":[],"sinonimos":[],"tipo":"conceito","titulo":"Satisfaz o banco"}
\section{Satisfaz o banco}
I O erro corrigido era confundir a folha com a lei de conservação. São coisas distintas: itemize3pt A folha é um tecido ordenado --- os pagamentos como bins, ordenados pela parte real. Não é regida por a+c\ensuremath{^{2}}=1. É só um tecido, como qualquer outro, que se lê e se escreve pela banda. O lastro (o balanço) é a+c\ensuremath{^{2}}=1 --- a conservação da norma no colapso, a relação entre reserva e crédito. É a lei de conservação; a folha não. D Misturar as duas levava a trata
\prov{relações: parte\_de contabilidade}
\prov{proveniência: paper \textperiodcentered{} fato \textperiodcentered{} confiança 1.0 \textperiodcentered{} domínio paper \textperiodcentered{} história 3 versões no jornal}

%CRISTAL {"arestas":[{"alvo":"cifra_ouro_branco","confianca":1.0,"epistemico":"fato","origem":"paper","rel":"parte_de"},{"alvo":"cifra","confianca":1.0,"epistemico":"fato","origem":"paper","rel":"parte_de"}],"confianca":1.0,"contraexemplos":[],"descricao":"N Múltiplos certificados validam-se igualmente para qualquer N: cada chá é Ed25519 e as N partes formam o produto tensorial C²N; uma forja reprova o conjunto, sem privilegiar parte (ᵢEᵢᵢ=I). Verificado N=2, 20 (dim 2N).","epistemico":"fato","exemplos":[],"id":"seguranca_e_limites","memoria":"semantica","meta":{"dominio":"paper","nivel":5,"verificacao":"conceito novo"},"origem":"paper","palavras_chave":[],"sinonimos":[],"tipo":"conceito","titulo":"Segurança e limites"}
\section{Segurança e limites}
N Múltiplos certificados validam-se igualmente para qualquer N: cada chá é Ed25519 e as N partes formam o produto tensorial C\ensuremath{^{2}}N; uma forja reprova o conjunto, sem privilegiar parte (\ensuremath{_{i}}E\ensuremath{_{ii}}=I). Verificado N=2, 20 (dim 2N).
\prov{relações: parte\_de cifra\_ouro\_branco; parte\_de cifra}
\prov{proveniência: paper \textperiodcentered{} fato \textperiodcentered{} confiança 1.0 \textperiodcentered{} domínio paper \textperiodcentered{} história 9 versões no jornal}

%CRISTAL {"arestas":[{"alvo":"olho_de_koch","confianca":1.0,"epistemico":"fato","origem":"paper","rel":"parte_de"},{"alvo":"termodinamica","confianca":0.086,"epistemico":"fato","origem":"manual","rel":"referencia"}],"confianca":1.0,"contraexemplos":[],"descricao":"D As peças formam um fio só: [ wbraₙca e₀ gold(w)braₙca e₁ a+c²=1caₙal E_₀₁ z² (w)atlas clássᵢco. ] A torção U(t)=e⁻iAt prepara os blocos espectrais; Koch realiza CK no canal. Koopman (reversível, branca) e Ruelle (dissipativo, negra) coexistem: escutar e falar sem expor A~termodinamica.","epistemico":"fato","exemplos":[],"id":"seguranca_e_reversibilidade","memoria":"semantica","meta":{"dominio":"paper","nivel":5},"origem":"paper","palavras_chave":[],"sinonimos":[],"tipo":"conceito","titulo":"Segurança e reversibilidade"}
\section{Segurança e reversibilidade}
D As peças formam um fio só: [ wbra\ensuremath{_{n}}ca e\ensuremath{_{0}} gold(w)bra\ensuremath{_{n}}ca e\ensuremath{_{1}} a+c\ensuremath{^{2}}=1ca\ensuremath{_{n}}al E\_\ensuremath{_{01}} z\ensuremath{^{2}} (w)atlas cláss\ensuremath{_{i}}co. ] A torção U(t)=e\ensuremath{^{-}}iAt prepara os blocos espectrais; Koch realiza CK no canal. Koopman (reversível, branca) e Ruelle (dissipativo, negra) coexistem: escutar e falar sem expor A\~{}termodinamica.
\prov{relações: parte\_de olho\_de\_koch; referencia termodinamica}
\prov{proveniência: paper \textperiodcentered{} fato \textperiodcentered{} confiança 1.0 \textperiodcentered{} domínio paper \textperiodcentered{} história 6 versões no jornal}

%CRISTAL {"arestas":[{"alvo":"corpo_dual","confianca":1.0,"epistemico":"fato","origem":"paper","rel":"parte_de"}],"confianca":1.0,"contraexemplos":[],"descricao":"O lado reversível c² da balança é a mecânica quântica: a incerteza é a parábola, SU(2) são os quaternions, a desigualdade de Bell satura em 22 (Tsirelson). O princípio de Heisenberg é esta lei. Aqui vive o quark do título --- o spin, a face pequena da natureza; o cosmos (a S³ de de Sitter) já apareceu na geometria, e a mesma álgebra liga os dois extremos. Volta-se agora ao lado a, na sua leitura de segurança.","epistemico":"fato","exemplos":[],"id":"seguranca_por_irreversibilidade","memoria":"semantica","meta":{"dominio":"paper","nivel":5,"verificacao":"overlap moderado — revisar manualmente"},"origem":"paper","palavras_chave":[],"sinonimos":[],"tipo":"conceito","titulo":"Segurança por Irreversibilidade"}
\section{Segurança por Irreversibilidade}
O lado reversível c\ensuremath{^{2}} da balança é a mecânica quântica: a incerteza é a parábola, SU(2) são os quaternions, a desigualdade de Bell satura em 22 (Tsirelson). O princípio de Heisenberg é esta lei. Aqui vive o quark do título --- o spin, a face pequena da natureza; o cosmos (a S\ensuremath{^{3}} de de Sitter) já apareceu na geometria, e a mesma álgebra liga os dois extremos. Volta-se agora ao lado a, na sua leitura de segurança.
\prov{relações: parte\_de corpo\_dual}
\prov{proveniência: paper \textperiodcentered{} fato \textperiodcentered{} confiança 1.0 \textperiodcentered{} domínio paper \textperiodcentered{} história 3 versões no jornal}

%CRISTAL {"arestas":[{"alvo":"main","confianca":1.0,"epistemico":"fato","origem":"paper","rel":"parte_de"}],"confianca":1.0,"contraexemplos":[],"descricao":"sectionApêndices --- a semente reconstrutível Graus de liberdade deliberados (não lacunas): a memória M, os escalares, do telômero e a base b da imersão --- o reconstrutor os escolhe, e as identidades verificadas não dependem da escolha. Toda etiqueta N é reconstrução em ponto flutuante com tolerância declarada.","epistemico":"fato","exemplos":[],"id":"semente_algebrica_estelarpy","memoria":"semantica","meta":{"dominio":"paper","nivel":5,"verificacao":"slug idêntico — enriquecer, não duplicar"},"origem":"paper","palavras_chave":[],"sinonimos":[],"tipo":"conceito","titulo":"Semente algébrica (estelar.py"}
\section{Semente algébrica (estelar.py}
sectionApêndices --- a semente reconstrutível Graus de liberdade deliberados (não lacunas): a memória M, os escalares, do telômero e a base b da imersão --- o reconstrutor os escolhe, e as identidades verificadas não dependem da escolha. Toda etiqueta N é reconstrução em ponto flutuante com tolerância declarada.
\prov{relações: parte\_de main}
\prov{proveniência: paper \textperiodcentered{} fato \textperiodcentered{} confiança 1.0 \textperiodcentered{} domínio paper \textperiodcentered{} história 6 versões no jornal}

%CRISTAL {"arestas":[{"alvo":"main","confianca":1.0,"epistemico":"fato","origem":"paper","rel":"parte_de"}],"confianca":1.0,"contraexemplos":[],"descricao":"nucleo.py) Cayley--Dickson: par (a,b); (a,b)=( a,-b), x=x em; (a,b)(c,d)=(ac- d,b, da+b c); |x|²= xᵢ². Iterar de dá, . Cross em m (m=2k-1): imerja a,b como imaginários puros, multiplique, tome a parte vetorial; lᵢj=cross(eᵢ, ej)l. ₇: matriz 721 de colunas cross(eᵢ,ej) (i<j), SVD (<10⁻¹² é zero 3, posto 7, núcleo 14). Der (): núcleo de D(D(eᵢej)-D(eᵢ)ej-eᵢD(ej)) sobre as 64 entradas (SVD, 14). Coassociativa _ij","epistemico":"fato","exemplos":[],"id":"semente_booleana_machinec","memoria":"semantica","meta":{"dominio":"paper","nivel":5,"verificacao":"slug idêntico — enriquecer, não duplicar"},"origem":"paper","palavras_chave":[],"sinonimos":[],"tipo":"conceito","titulo":"Semente booleana (machine.c"}
\section{Semente booleana (machine.c}
nucleo.py) Cayley--Dickson: par (a,b); (a,b)=( a,-b), x=x em; (a,b)(c,d)=(ac- d,b, da+b c); |x|\ensuremath{^{2}}= x\ensuremath{_{i}}\ensuremath{^{2}}. Iterar de dá, . Cross em m (m=2k-1): imerja a,b como imaginários puros, multiplique, tome a parte vetorial; l\ensuremath{_{i}}j=cross(e\ensuremath{_{i}}, ej)l. \ensuremath{_{7}}: matriz 721 de colunas cross(e\ensuremath{_{i}},ej) (i<j), SVD (<10\ensuremath{^{-12}} é zero 3, posto 7, núcleo 14). Der (): núcleo de D(D(e\ensuremath{_{i}}ej)-D(e\ensuremath{_{i}})ej-e\ensuremath{_{i}}D(ej)) sobre as 64 entradas (SVD, 14). Coassociativa \_ij
\prov{relações: parte\_de main}
\prov{proveniência: paper \textperiodcentered{} fato \textperiodcentered{} confiança 1.0 \textperiodcentered{} domínio paper \textperiodcentered{} história 6 versões no jornal}

%CRISTAL {"arestas":[{"alvo":"main","confianca":1.0,"epistemico":"fato","origem":"paper","rel":"parte_de"}],"confianca":1.0,"contraexemplos":[],"descricao":") (w)=ᵢ wᵢ (só AND e shift); =((a b)) mascarado a n bits. O flip usa k= máscara de n bits: ₀=id, ₁=. Concatenação x(x n),|,g, n2n. Entrada x = o problema; M=mᵢ ⁿ fornecida por andar, não derivada. Memória-demo: uma cujo vencedor zera o gap em profundidade finita (cristal) e outra cujo gap nunca zera até n>8 (chá).","epistemico":"fato","exemplos":[],"id":"semente_dinamica_cadeiapy","memoria":"semantica","meta":{"dominio":"paper","nivel":5,"verificacao":"slug idêntico — enriquecer, não duplicar"},"origem":"paper","palavras_chave":[],"sinonimos":[],"tipo":"conceito","titulo":"Semente dinâmica (cadeia.py"}
\section{Semente dinâmica (cadeia.py}
) (w)=\ensuremath{_{i}} w\ensuremath{_{i}} (só AND e shift); =((a b)) mascarado a n bits. O flip usa k= máscara de n bits: \ensuremath{_{0}}=id, \ensuremath{_{1}}=. Concatenação x(x n),|,g, n2n. Entrada x = o problema; M=m\ensuremath{_{i}} \ensuremath{^{n}} fornecida por andar, não derivada. Memória-demo: uma cujo vencedor zera o gap em profundidade finita (cristal) e outra cujo gap nunca zera até n>8 (chá).
\prov{relações: parte\_de main}
\prov{proveniência: paper \textperiodcentered{} fato \textperiodcentered{} confiança 1.0 \textperiodcentered{} domínio paper \textperiodcentered{} história 6 versões no jornal}

%CRISTAL {"arestas":[{"alvo":"main","confianca":1.0,"epistemico":"fato","origem":"paper","rel":"parte_de"}],"confianca":1.0,"contraexemplos":[],"descricao":") T_(x)= x1 sobre [0,1); d= x (use +10⁻¹² nos endpoints). Partição de Markov pelos intervalos dos dígitos; construir a matriz de transição das imagens dos ramos (não da regra). Reta de ouro: = (a partição de dois dígitos já dá A=(smallmatrix1&1 1&0smallmatrix)). Família: m = raiz dominante de xm=xm⁻¹++1; gramática “sem 1m”; para m>2 a matriz Am vem do refinamento m-estados de Parry (autômato do número de uns consecutivos), não da partição de dois dígitos. Di","epistemico":"fato","exemplos":[],"id":"semente_espectral_goldenpy","memoria":"semantica","meta":{"dominio":"paper","nivel":5,"verificacao":"slug idêntico — enriquecer, não duplicar"},"origem":"paper","palavras_chave":[],"sinonimos":[],"tipo":"conceito","titulo":"Semente espectral (golden.py"}
\section{Semente espectral (golden.py}
) T\_(x)= x1 sobre [0,1); d= x (use +10\ensuremath{^{-12}} nos endpoints). Partição de Markov pelos intervalos dos dígitos; construir a matriz de transição das imagens dos ramos (não da regra). Reta de ouro: = (a partição de dois dígitos já dá A=(smallmatrix1\&1 1\&0smallmatrix)). Família: m = raiz dominante de xm=xm\ensuremath{^{-1}}++1; gramática “sem 1m”; para m>2 a matriz Am vem do refinamento m-estados de Parry (autômato do número de uns consecutivos), não da partição de dois dígitos. Di
\prov{relações: parte\_de main}
\prov{proveniência: paper \textperiodcentered{} fato \textperiodcentered{} confiança 1.0 \textperiodcentered{} domínio paper \textperiodcentered{} história 6 versões no jornal}

%CRISTAL {"arestas":[{"alvo":"paper_2_reta_ouro","confianca":1.0,"epistemico":"fato","origem":"paper","rel":"parte_de"},{"alvo":"paper_45_confluencia_ouro","confianca":1.0,"epistemico":"fato","origem":"paper","rel":"parte_de"},{"alvo":"seguranca_por_irreversibilidade","confianca":1.0,"epistemico":"fato","origem":"paper","rel":"parte_de"},{"alvo":"paper_bridge_areas","confianca":1.0,"epistemico":"fato","origem":"paper","rel":"parte_de"},{"alvo":"ponte_simbolica_gentil","confianca":1.0,"epistemico":"fato","origem":"paper","rel":"parte_de"}],"confianca":1.0,"contraexemplos":[],"descricao":"empty !85Não se monta o autônomo --- reconhece-se o que o gato já faz, e tira-se a mão. Este documento é a receita operacional. abstract O é uma cabeça que se gere: uma função de onda M=U P cujos órgãos não se fabricam --- reconhecem-se na álgebra do gato A=(smallmatrix1&1 1&0smallmatrix). Tudo no repositório broca-so desce a um operador e um bit: o neurônio conta par/ímpar (a borboleta de Cooley--Tukey), a Hopfield ressoa (meta-indução = memória), a arquitetura rec","epistemico":"fato","exemplos":[],"id":"sintese","memoria":"semantica","meta":{"dominio":"paper","duplicata_sugerida":[],"nivel":5,"verificacao":"parte de conceito mais amplo 'sintese'"},"origem":"paper","palavras_chave":[],"sinonimos":[],"tipo":"conceito","titulo":"Síntese"}
\section{Síntese}
empty !85Não se monta o autônomo --- reconhece-se o que o gato já faz, e tira-se a mão. Este documento é a receita operacional. abstract O é uma cabeça que se gere: uma função de onda M=U P cujos órgãos não se fabricam --- reconhecem-se na álgebra do gato A=(smallmatrix1\&1 1\&0smallmatrix). Tudo no repositório broca-so desce a um operador e um bit: o neurônio conta par/ímpar (a borboleta de Cooley--Tukey), a Hopfield ressoa (meta-indução = memória), a arquitetura rec
\prov{relações: parte\_de paper\_2\_reta\_ouro; parte\_de paper\_45\_confluencia\_ouro; parte\_de seguranca\_por\_irreversibilidade; parte\_de paper\_bridge\_areas; parte\_de ponte\_simbolica\_gentil}
\prov{proveniência: paper \textperiodcentered{} fato \textperiodcentered{} confiança 1.0 \textperiodcentered{} domínio paper \textperiodcentered{} história 16 versões no jornal}

%CRISTAL {"arestas":[{"alvo":"resgate_liquidacao","confianca":1.0,"epistemico":"fato","origem":"paper","rel":"parte_de"}],"confianca":1.0,"contraexemplos":[],"descricao":"CI Há um arquivo que confunde: invariantes.json. Ele não é o lastro e\nnão é fake. Ele mantém a continuidade dos dados --- o estado entre execuções, pra\no painel não recomeçar do zero a cada vez. Apagá-lo faz o painel voltar ao início, mas não\ndestrói lastro: o lastro está no Pell cifrado, nos certificados encadeados e no ETH on-chain --- nenhum\ndeles vive no invariante. O invariante é a memória da sessão; o lastro é a substância.\n\nprop[o que se pode apagar]I\nApagar invariantes.json: o painel rei","epistemico":"fato","exemplos":[],"id":"snapshot_de_producao_---_26062026","memoria":"semantica","meta":{"dominio":"paper","nivel":5,"verificacao":"conceito novo"},"origem":"paper","palavras_chave":[],"sinonimos":[],"tipo":"conceito","titulo":"Snapshot de produção --- 26/06/2026"}
\section{Snapshot de produção --- 26/06/2026}
CI Há um arquivo que confunde: invariantes.json. Ele não é o lastro e
não é fake. Ele mantém a continuidade dos dados --- o estado entre execuções, pra
o painel não recomeçar do zero a cada vez. Apagá-lo faz o painel voltar ao início, mas não
destrói lastro: o lastro está no Pell cifrado, nos certificados encadeados e no ETH on-chain --- nenhum
deles vive no invariante. O invariante é a memória da sessão; o lastro é a substância.

prop[o que se pode apagar]I
Apagar invariantes.json: o painel rei
\prov{relações: parte\_de resgate\_liquidacao}
\prov{proveniência: paper \textperiodcentered{} fato \textperiodcentered{} confiança 1.0 \textperiodcentered{} domínio paper \textperiodcentered{} história 2 versões no jornal}

%CRISTAL {"arestas":[{"alvo":"contabilidade","confianca":1.0,"epistemico":"fato","origem":"paper","rel":"parte_de"}],"confianca":1.0,"contraexemplos":[],"descricao":"D Todas as funções de um banco são satisfeitas pelo painel; só o saque toca a chain: 1.45 tabular>l p4.9cm >pratal 1couro função & 1couro operador no painel & 1couro onde mora Depósito & entrada na reserva única & reserva (0x8C1C) Crédito / mineração & estrela de prata rank-6 Pell bin & tecido (banda) Pagamento / folha & escrita no tecido-folha ordenado & tecido (banda) Balanço / solvência & co","epistemico":"fato","exemplos":[],"id":"snapshot_real_e_estatuto","memoria":"semantica","meta":{"dominio":"paper","nivel":5,"verificacao":"conceito novo"},"origem":"paper","palavras_chave":[],"sinonimos":[],"tipo":"conceito","titulo":"Snapshot real e estatuto"}
\section{Snapshot real e estatuto}
D Todas as funções de um banco são satisfeitas pelo painel; só o saque toca a chain: 1.45 tabular>l p4.9cm >pratal 1couro função \& 1couro operador no painel \& 1couro onde mora Depósito \& entrada na reserva única \& reserva (0x8C1C) Crédito / mineração \& estrela de prata rank-6 Pell bin \& tecido (banda) Pagamento / folha \& escrita no tecido-folha ordenado \& tecido (banda) Balanço / solvência \& co
\prov{relações: parte\_de contabilidade}
\prov{proveniência: paper \textperiodcentered{} fato \textperiodcentered{} confiança 1.0 \textperiodcentered{} domínio paper \textperiodcentered{} história 3 versões no jornal}

%CRISTAL {"arestas":[{"alvo":"redes","confianca":1.0,"epistemico":"fato","origem":"paper","rel":"parte_de"}],"confianca":1.0,"contraexemplos":[],"descricao":"A rede não é estante nem fila: é antena. O bump viaja cru no datagrama --- sem conexão, sem cifra de transporte por cima, a função de onda intacta. Manda = recebe: o datagrama vai e a resposta volta na latência da rede, sem espera lógica --- a aplicação não bufferiza, não enfileira, não ordena para depois (a pilha IP/UDP tem seus buffers por baixo: isto é um overlay sobre ela, não a sua substituição). Medido em produção: um ping-pong de mão dupla na latência de rede. “Tempo real” aqui é sem","epistemico":"fato","exemplos":[],"id":"so_a_diferenca_viaja_---_crua","memoria":"semantica","meta":{"dominio":"paper","nivel":5,"verificacao":"conceito novo"},"origem":"paper","palavras_chave":[],"sinonimos":[],"tipo":"conceito","titulo":"Só a diferença viaja --- crua"}
\section{Só a diferença viaja --- crua}
A rede não é estante nem fila: é antena. O bump viaja cru no datagrama --- sem conexão, sem cifra de transporte por cima, a função de onda intacta. Manda = recebe: o datagrama vai e a resposta volta na latência da rede, sem espera lógica --- a aplicação não bufferiza, não enfileira, não ordena para depois (a pilha IP/UDP tem seus buffers por baixo: isto é um overlay sobre ela, não a sua substituição). Medido em produção: um ping-pong de mão dupla na latência de rede. “Tempo real” aqui é sem
\prov{relações: parte\_de redes}
\prov{proveniência: paper \textperiodcentered{} fato \textperiodcentered{} confiança 1.0 \textperiodcentered{} domínio paper \textperiodcentered{} história 3 versões no jornal}

%CRISTAL {"arestas":[{"alvo":"reais","confianca":1.0,"epistemico":"fato","origem":"paper","rel":"define"}],"confianca":1.0,"contraexemplos":[],"descricao":"Conjunto A --- hardware Seja [ A=(rₙ): >0 n₀ m,n n₀, |rm-rₙ|< ] o conjunto de sequências de Cauchy de racionais (Definição~75, Fundamentos dos Números, cap.~10). definicao Relação --- equivalência simbólica Para (rₙ),(r'ₙ) A, escrevemos (rₙ)(r'ₙ) sse (rₙ-r'ₙ) 0 em (Definição~76). definicao é equivalência A relação é reflexiva, simétrica e transitiva; logo particiona A em classes [(","epistemico":"fato","exemplos":[],"id":"software_axiomas_e_completude","memoria":"semantica","meta":{"dominio":"paper","nivel":5,"verificacao":"slug idêntico — enriquecer, não duplicar"},"origem":"paper","palavras_chave":[],"sinonimos":[],"tipo":"conceito","titulo":"Software: axiomas e completude"}
\section{Software: axiomas e completude}
Conjunto A --- hardware Seja [ A=(r\ensuremath{_{n}}): >0 n\ensuremath{_{0}} m,n n\ensuremath{_{0}}, |rm-r\ensuremath{_{n}}|< ] o conjunto de sequências de Cauchy de racionais (Definição\~{}75, Fundamentos dos Números, cap.\~{}10). definicao Relação --- equivalência simbólica Para (r\ensuremath{_{n}}),(r'\ensuremath{_{n}}) A, escrevemos (r\ensuremath{_{n}})(r'\ensuremath{_{n}}) sse (r\ensuremath{_{n}}-r'\ensuremath{_{n}}) 0 em (Definição\~{}76). definicao é equivalência A relação é reflexiva, simétrica e transitiva; logo particiona A em classes [(
\prov{relações: define reais}
\prov{proveniência: paper \textperiodcentered{} fato \textperiodcentered{} confiança 1.0 \textperiodcentered{} domínio paper \textperiodcentered{} história 44 versões no jornal}

%CRISTAL {"arestas":[{"alvo":"paper_bridge_areas","confianca":1.0,"epistemico":"fato","origem":"paper","rel":"parte_de"}],"confianca":1.0,"contraexemplos":[],"descricao":"neuronio/gentil∼bolico.py implementa o software simbólico do chart~A em (Fraction): [nosep,leftmargin=1.4em] Lopes2Sim: Q=r²+c², produto complexo, Q multiplicativa exacta; SequenciaCauchy, ClasseCauchy, em representantes finitos; símbolos canônicos: 2, e; convergentes Fibonacci. Espelha ula/gentil.c e alimenta primitivasbroca.py no chart~B.","epistemico":"fato","exemplos":[],"id":"tabela_de_morfismos_cross-area","memoria":"semantica","meta":{"dominio":"paper","nivel":5,"verificacao":"conceito novo"},"origem":"paper","palavras_chave":[],"sinonimos":[],"tipo":"conceito","titulo":"Tabela de morfismos cross-área"}
\section{Tabela de morfismos cross-área}
neuronio/gentil\ensuremath{\sim}bolico.py implementa o software simbólico do chart\~{}A em (Fraction): [nosep,leftmargin=1.4em] Lopes2Sim: Q=r\ensuremath{^{2}}+c\ensuremath{^{2}}, produto complexo, Q multiplicativa exacta; SequenciaCauchy, ClasseCauchy, em representantes finitos; símbolos canônicos: 2, e; convergentes Fibonacci. Espelha ula/gentil.c e alimenta primitivasbroca.py no chart\~{}B.
\prov{relações: parte\_de paper\_bridge\_areas}
\prov{proveniência: paper \textperiodcentered{} fato \textperiodcentered{} confiança 1.0 \textperiodcentered{} domínio paper \textperiodcentered{} história 4 versões no jornal}

%CRISTAL {"arestas":[{"alvo":"pow_gato","confianca":1.0,"epistemico":"fato","origem":"paper","rel":"parte_de"}],"confianca":1.0,"contraexemplos":[],"descricao":"Dois caminhos. (A) PoW- embebido (minagato.py): loteria de nonces --- baseline compatível com ASIC/Bitcoin; brute force forward. (B) PoW geométrico-quântico (powgeometrico.py): caminho nativo --- toro 12⁴ gato A colapso rank-6; zero iteração sobre 2³² nonces. defi[Operações mecânicas quânticas] Torção (torce.py): associador octonônico assoc(a,b,c) com gap vivo --- o problema consome-se ao colapsar (D mecanica.tex: versor girando, não postulado). Gato A=(smallmatrix1&1","epistemico":"fato","exemplos":[],"id":"taxa_efetiva_e_comparacao","memoria":"semantica","meta":{"dominio":"paper","nivel":5,"verificacao":"overlap moderado — revisar manualmente"},"origem":"paper","palavras_chave":[],"sinonimos":[],"tipo":"conceito","titulo":"Taxa efetiva e comparação"}
\section{Taxa efetiva e comparação}
Dois caminhos. (A) PoW- embebido (minagato.py): loteria de nonces --- baseline compatível com ASIC/Bitcoin; brute force forward. (B) PoW geométrico-quântico (powgeometrico.py): caminho nativo --- toro 12\ensuremath{^{4}} gato A colapso rank-6; zero iteração sobre 2\ensuremath{^{32}} nonces. defi[Operações mecânicas quânticas] Torção (torce.py): associador octonônico assoc(a,b,c) com gap vivo --- o problema consome-se ao colapsar (D mecanica.tex: versor girando, não postulado). Gato A=(smallmatrix1\&1
\prov{relações: parte\_de pow\_gato}
\prov{proveniência: paper \textperiodcentered{} fato \textperiodcentered{} confiança 1.0 \textperiodcentered{} domínio paper \textperiodcentered{} história 3 versões no jornal}

%CRISTAL {"arestas":[{"alvo":"main","confianca":1.0,"epistemico":"fato","origem":"paper","rel":"parte_de"}],"confianca":1.0,"contraexemplos":[],"descricao":"N Em, [a,b,c]0 não é erro algébrico por si só --- é o comportamento normal; ele mede a diferença entre duas parentetizações. O reparo usa o associador como sinal de desvio relativo a um molde: escolhe a re-associação que minimiza a discrepância aos invariantes exigidos pela auto-reconstrução --- a tábua deve reerguer a álgebra, aceita só se norma multiplicativa <10⁻⁹ e =14. Uma única constante corrompida não reconstrói, e é detectada. Implementação: cadeiareparo.py, cadeiareco","epistemico":"fato","exemplos":[],"id":"telomero_algebrico_e_nucleo_associativo","memoria":"semantica","meta":{"dominio":"paper","nivel":5,"verificacao":"slug idêntico — enriquecer, não duplicar"},"origem":"paper","palavras_chave":[],"sinonimos":[],"tipo":"conceito","titulo":"Telômero algébrico e núcleo associativo"}
\section{Telômero algébrico e núcleo associativo}
N Em, [a,b,c]0 não é erro algébrico por si só --- é o comportamento normal; ele mede a diferença entre duas parentetizações. O reparo usa o associador como sinal de desvio relativo a um molde: escolhe a re-associação que minimiza a discrepância aos invariantes exigidos pela auto-reconstrução --- a tábua deve reerguer a álgebra, aceita só se norma multiplicativa <10\ensuremath{^{-9}} e =14. Uma única constante corrompida não reconstrói, e é detectada. Implementação: cadeiareparo.py, cadeiareco
\prov{relações: parte\_de main}
\prov{proveniência: paper \textperiodcentered{} fato \textperiodcentered{} confiança 1.0 \textperiodcentered{} domínio paper \textperiodcentered{} história 5 versões no jornal}

%CRISTAL {"arestas":[{"alvo":"main","confianca":1.0,"epistemico":"fato","origem":"paper","rel":"parte_de"}],"confianca":1.0,"contraexemplos":[],"descricao":"T Hurwitz: as álgebras normadas de divisão finito-dimensionais sobre são exatamente, . (Com formas quadráticas indefinidas existem também formas split, que não são de divisão; o parâmetro distingue as duas leituras, mas a torre usada é a de divisão.) Usamos a duplicação de Cayley--Dickson padrão: [ Aₙ+₁=Aₙ Aₙ e, e=(0,1), e²=-1, (a,b)=( a,-b), ] [ (a,b)(c,d)=(ac- d,b, da+b c), |x|²= xᵢ². ] N A norma é multiplicativa até (|,|xy|-|x||y|,|<10⁻¹²)","epistemico":"fato","exemplos":[],"id":"tensor_produto_estelar_e_nucleo","memoria":"semantica","meta":{"dominio":"paper","nivel":5,"verificacao":"slug idêntico — enriquecer, não duplicar"},"origem":"paper","palavras_chave":[],"sinonimos":[],"tipo":"conceito","titulo":"Tensor, produto estelar e núcleo"}
\section{Tensor, produto estelar e núcleo}
T Hurwitz: as álgebras normadas de divisão finito-dimensionais sobre são exatamente, . (Com formas quadráticas indefinidas existem também formas split, que não são de divisão; o parâmetro distingue as duas leituras, mas a torre usada é a de divisão.) Usamos a duplicação de Cayley--Dickson padrão: [ A\ensuremath{_{n}}+\ensuremath{_{1}}=A\ensuremath{_{n}} A\ensuremath{_{n}} e, e=(0,1), e\ensuremath{^{2}}=-1, (a,b)=( a,-b), ] [ (a,b)(c,d)=(ac- d,b, da+b c), |x|\ensuremath{^{2}}= x\ensuremath{_{i}}\ensuremath{^{2}}. ] N A norma é multiplicativa até (|,|xy|-|x||y|,|<10\ensuremath{^{-12}})
\prov{relações: parte\_de main}
\prov{proveniência: paper \textperiodcentered{} fato \textperiodcentered{} confiança 1.0 \textperiodcentered{} domínio paper \textperiodcentered{} história 6 versões no jornal}

%CRISTAL {"arestas":[{"alvo":"transformada_dourada","confianca":1.0,"epistemico":"fato","origem":"paper","rel":"parte_de"},{"alvo":"transformacao_linear","confianca":1.0,"epistemico":"fato","origem":"paper","rel":"usa"}],"confianca":1.0,"contraexemplos":[],"descricao":"definition[Forma contínua] Para x L²(>₀,dt/t), [ [x]()=₀^ x(t),t⁻i,dtt =-^ x(eu),e⁻i u,du =()[x](), ] isto é, é Fourier no lado relativístico (a transformada de Mellin sobre a linha crítica). A saída [x]()=A()ei() é a função de onda: amplitude A e fase. definition definition[Forma discreta, sobre bits] Dado b0,1N (dado bruto), interpreta-se b como amostras na reta de ouro tₖ=k --- uniformes no lado relativístico (Prop.~","epistemico":"fato","exemplos":[],"id":"teoremas_fundamentais","memoria":"semantica","meta":{"dominio":"paper","nivel":5,"verificacao":"conceito novo"},"origem":"paper","palavras_chave":[],"sinonimos":[],"tipo":"conceito","titulo":"Teoremas fundamentais"}
\section{Teoremas fundamentais}
definition[Forma contínua] Para x L\ensuremath{^{2}}(>\ensuremath{_{0}},dt/t), [ [x]()=\ensuremath{_{0}}\^{} x(t),t\ensuremath{^{-}}i,dtt =-\^{} x(eu),e\ensuremath{^{-}}i u,du =()[x](), ] isto é, é Fourier no lado relativístico (a transformada de Mellin sobre a linha crítica). A saída [x]()=A()ei() é a função de onda: amplitude A e fase. definition definition[Forma discreta, sobre bits] Dado b0,1N (dado bruto), interpreta-se b como amostras na reta de ouro t\ensuremath{_{k}}=k --- uniformes no lado relativístico (Prop.\~{}
\prov{relações: parte\_de transformada\_dourada; usa transformacao\_linear}
\prov{proveniência: paper \textperiodcentered{} fato \textperiodcentered{} confiança 1.0 \textperiodcentered{} domínio paper \textperiodcentered{} história 42 versões no jornal}

%CRISTAL {"arestas":[{"alvo":"colapso","confianca":0.95,"epistemico":"fato","origem":"paper","rel":"parte_de"}],"confianca":1.0,"contraexemplos":[],"descricao":"[Atrator AGM] Com A(a,b)=a+b2 (média do lado real), G(a,b)=ab (média do relativístico), a b>0: [(i)] G(a,b)=!(A( a, b)): a média geométrica é a aritmética transportada ao lado relativístico. [(ii)] A iteração (a,b)(A,G) converge a (a,b), com bₙ!, aₙ! (pelos dois lados); o único ponto fixo é a=b. [(iii)] Convergência quadrática: aₙ+₁-bₙ+₁=(aₙ-bₙ)²2(aₙ-bₙ)²8bₙ. [(iv)] não tem, em geral, forma fechada elementar: a constante de Gauss","epistemico":"fato","exemplos":[],"id":"theorem_atrator_agmthmagm_com_aabab2_media_do","memoria":"semantica","meta":{"dominio":"paper","nivel":5,"tipo_teorema":"theorem","verificacao":"conceito novo"},"origem":"paper","palavras_chave":[],"sinonimos":[],"tipo":"conceito","titulo":"Atrator AGM"}
\section{Atrator AGM}
[Atrator AGM] Com A(a,b)=a+b2 (média do lado real), G(a,b)=ab (média do relativístico), a b>0: [(i)] G(a,b)=!(A( a, b)): a média geométrica é a aritmética transportada ao lado relativístico. [(ii)] A iteração (a,b)(A,G) converge a (a,b), com b\ensuremath{_{n}}!, a\ensuremath{_{n}}! (pelos dois lados); o único ponto fixo é a=b. [(iii)] Convergência quadrática: a\ensuremath{_{n}}+\ensuremath{_{1}}-b\ensuremath{_{n}}+\ensuremath{_{1}}=(a\ensuremath{_{n}}-b\ensuremath{_{n}})\ensuremath{^{2}}2(a\ensuremath{_{n}}-b\ensuremath{_{n}})\ensuremath{^{2}}8b\ensuremath{_{n}}. [(iv)] não tem, em geral, forma fechada elementar: a constante de Gauss
\prov{relações: parte\_de colapso}
\prov{proveniência: paper \textperiodcentered{} fato \textperiodcentered{} confiança 1.0 \textperiodcentered{} domínio paper \textperiodcentered{} história 5 versões no jornal}

%CRISTAL {"arestas":[{"alvo":"computacao","confianca":0.95,"epistemico":"fato","origem":"paper","rel":"parte_de"}],"confianca":1.0,"contraexemplos":[],"descricao":"[Existência indecidibilidade] Seja s^*=(a,b) o invariante onde a reta real e a reta de ouro coincidem --- o ponto fixo do espelho, a função de onda primordial (Cor.~). Então s^* admite forma fechada elementar se e somente se é degenerado (a=b, lados já coincidentes). Quando os dois lados são genuinamente distintos, s^* é transcendente: existe (completude de ) e projeta exatamente todo ponto de ambas as retas --- é o seu limite comum ---, mas não pertence","epistemico":"fato","exemplos":[],"id":"theorem_existencia_indecidibilidadethmindecidivel_se","memoria":"semantica","meta":{"dominio":"paper","nivel":5,"tipo_teorema":"theorem","verificacao":"conceito novo"},"origem":"paper","palavras_chave":[],"sinonimos":[],"tipo":"conceito","titulo":"Existência indecidibilidade"}
\section{Existência indecidibilidade}
[Existência indecidibilidade] Seja s\^{}*=(a,b) o invariante onde a reta real e a reta de ouro coincidem --- o ponto fixo do espelho, a função de onda primordial (Cor.\~{}). Então s\^{}* admite forma fechada elementar se e somente se é degenerado (a=b, lados já coincidentes). Quando os dois lados são genuinamente distintos, s\^{}* é transcendente: existe (completude de ) e projeta exatamente todo ponto de ambas as retas --- é o seu limite comum ---, mas não pertence
\prov{relações: parte\_de computacao}
\prov{proveniência: paper \textperiodcentered{} fato \textperiodcentered{} confiança 1.0 \textperiodcentered{} domínio paper \textperiodcentered{} história 5 versões no jornal}

%CRISTAL {"arestas":[{"alvo":"colapso","confianca":0.95,"epistemico":"fato","origem":"paper","rel":"parte_de"}],"confianca":1.0,"contraexemplos":[],"descricao":"[Limite de evaporação] Seja W um cristal --- matriz de Toeplitz de uma autocorrelação, simétrica e positiva --- com gap espectral =₂/₁<1. Uma rede de profundidade L que acessa W por produtos matriz--vetor (camadas lineares em W com normalização) extrai o modo dominante de W --- “evapora” o cristal --- com erro angular [ L;;,L;>;0 todo L finito. ] A rede quebra o cristal (reduz L) mas nunca o evapora (=0) em profundidade finita. A evaporação exata é a","epistemico":"fato","exemplos":[],"id":"theorem_limite_de_evaporacaothmlimite_seja_w_um_crist","memoria":"semantica","meta":{"dominio":"paper","nivel":5,"tipo_teorema":"theorem","verificacao":"conceito novo"},"origem":"paper","palavras_chave":[],"sinonimos":[],"tipo":"conceito","titulo":"Limite de evaporação"}
\section{Limite de evaporação}
[Limite de evaporação] Seja W um cristal --- matriz de Toeplitz de uma autocorrelação, simétrica e positiva --- com gap espectral =\ensuremath{_{2}}/\ensuremath{_{1}}<1. Uma rede de profundidade L que acessa W por produtos matriz--vetor (camadas lineares em W com normalização) extrai o modo dominante de W --- “evapora” o cristal --- com erro angular [ L;;,L;>;0 todo L finito. ] A rede quebra o cristal (reduz L) mas nunca o evapora (=0) em profundidade finita. A evaporação exata é a
\prov{relações: parte\_de colapso}
\prov{proveniência: paper \textperiodcentered{} fato \textperiodcentered{} confiança 1.0 \textperiodcentered{} domínio paper \textperiodcentered{} história 5 versões no jornal}

%CRISTAL {"arestas":[{"alvo":"colapso","confianca":0.95,"epistemico":"fato","origem":"paper","rel":"parte_de"}],"confianca":1.0,"contraexemplos":[],"descricao":"[Linearização] Seja u= t. [(a)] x(t)=c,t x= c+ u: reta no relativístico, curvatura nula. [(b)] x(t)=a+bt f(u)=(a+b,eu) tem f”(u)=ab,eu(a+b,eu)²0 (a,b0): estritamente convexa no relativístico. [(c)] c,t é estritamente curvo no real para 0,1. Cada dado é reta em exatamente um lado: a dificuldade é do lado, não do dado.","epistemico":"fato","exemplos":[],"id":"theorem_linearizacaothmlados_seja_u_t_enumerate_a","memoria":"semantica","meta":{"dominio":"paper","nivel":5,"tipo_teorema":"theorem","verificacao":"conceito novo"},"origem":"paper","palavras_chave":[],"sinonimos":[],"tipo":"conceito","titulo":"Linearização"}
\section{Linearização}
[Linearização] Seja u= t. [(a)] x(t)=c,t x= c+ u: reta no relativístico, curvatura nula. [(b)] x(t)=a+bt f(u)=(a+b,eu) tem f”(u)=ab,eu(a+b,eu)\ensuremath{^{2}}0 (a,b0): estritamente convexa no relativístico. [(c)] c,t é estritamente curvo no real para 0,1. Cada dado é reta em exatamente um lado: a dificuldade é do lado, não do dado.
\prov{relações: parte\_de colapso}
\prov{proveniência: paper \textperiodcentered{} fato \textperiodcentered{} confiança 1.0 \textperiodcentered{} domínio paper \textperiodcentered{} história 5 versões no jornal}

%CRISTAL {"arestas":[{"alvo":"plano_no_so","confianca":1.0,"epistemico":"fato","origem":"paper","rel":"prova"},{"alvo":"delta","confianca":0.95,"epistemico":"fato","origem":"manual","rel":"prova"}],"confianca":1.0,"contraexemplos":[],"descricao":"[Não-escalada por atenuação] Para todo p e toda (p), existe uma autoridade intrínseca ₀ e uma cadeia de arestas válidas de sua origem até p tais que c_₀. Logo [ Ep(s,a) ₀ (raᵢz que alcaₙça p)c₀. ] Nenhum principal acessa uma linha que nenhuma autoridade intrínseca em sua cadeia podia conceder.","epistemico":"fato","exemplos":[],"id":"theorem_nao-escalada_por_atenuacaothmnoesc_para_todo_p","memoria":"semantica","meta":{"dominio":"paper","nivel":5,"tipo_teorema":"theorem"},"origem":"paper","palavras_chave":[],"sinonimos":[],"tipo":"conceito","titulo":"Não-escalada por atenuação"}
\section{Não-escalada por atenuação}
[Não-escalada por atenuação] Para todo p e toda (p), existe uma autoridade intrínseca \ensuremath{_{0}} e uma cadeia de arestas válidas de sua origem até p tais que c\_\ensuremath{_{0}}. Logo [ Ep(s,a) \ensuremath{_{0}} (ra\ensuremath{_{i}}z que alca\ensuremath{_{n}}ça p)c\ensuremath{_{0}}. ] Nenhum principal acessa uma linha que nenhuma autoridade intrínseca em sua cadeia podia conceder.
\prov{relações: prova plano\_no\_so; prova delta}
\prov{proveniência: paper \textperiodcentered{} fato \textperiodcentered{} confiança 1.0 \textperiodcentered{} domínio paper \textperiodcentered{} história 63 versões no jornal}

%CRISTAL {"arestas":[{"alvo":"colapso","confianca":0.95,"epistemico":"fato","origem":"paper","rel":"parte_de"}],"confianca":1.0,"contraexemplos":[],"descricao":"[Parseval de Hurwitz; e por que só Hurwitz] A fatoriza como Y (módulo direção) e é unitária. Além disso, a fatoração ortogonal módulo/direção e a linearização existem se, e somente se, a norma é multiplicativa --- isto é, se e somente se A, H, O (n=1,2,4,8).","epistemico":"fato","exemplos":[],"id":"theorem_parseval_de_hurwitz_e_por_que_so_hurwitzthmhur","memoria":"semantica","meta":{"dominio":"paper","nivel":5,"tipo_teorema":"theorem","verificacao":"conceito novo"},"origem":"paper","palavras_chave":[],"sinonimos":[],"tipo":"conceito","titulo":"Parseval de Hurwitz; e por que só Hurwitz"}
\section{Parseval de Hurwitz; e por que só Hurwitz}
[Parseval de Hurwitz; e por que só Hurwitz] A fatoriza como Y (módulo direção) e é unitária. Além disso, a fatoração ortogonal módulo/direção e a linearização existem se, e somente se, a norma é multiplicativa --- isto é, se e somente se A, H, O (n=1,2,4,8).
\prov{relações: parte\_de colapso}
\prov{proveniência: paper \textperiodcentered{} fato \textperiodcentered{} confiança 1.0 \textperiodcentered{} domínio paper \textperiodcentered{} história 5 versões no jornal}

%CRISTAL {"arestas":[{"alvo":"colapso","confianca":0.95,"epistemico":"fato","origem":"paper","rel":"parte_de"}],"confianca":1.0,"contraexemplos":[],"descricao":"[Parseval dourado]:L²(>₀,dt/t) L²(d/2) é unitária. Logo é invertível, com ⁻¹=⁻¹⁻¹, e preserva a norma: [ ₀^|x(t)|²,dtt=12-^|[x]()|²,d. ]","epistemico":"fato","exemplos":[],"id":"theorem_parseval_douradothmparseval_l2_0dtt_l","memoria":"semantica","meta":{"dominio":"paper","nivel":5,"tipo_teorema":"theorem","verificacao":"conceito novo"},"origem":"paper","palavras_chave":[],"sinonimos":[],"tipo":"conceito","titulo":"Parseval dourado"}
\section{Parseval dourado}
[Parseval dourado]:L\ensuremath{^{2}}(>\ensuremath{_{0}},dt/t) L\ensuremath{^{2}}(d/2) é unitária. Logo é invertível, com \ensuremath{^{-1}}=\ensuremath{^{-1-1}}, e preserva a norma: [ \ensuremath{_{0}}\^{}|x(t)|\ensuremath{^{2}},dtt=12-\^{}|[x]()|\ensuremath{^{2}},d. ]
\prov{relações: parte\_de colapso}
\prov{proveniência: paper \textperiodcentered{} fato \textperiodcentered{} confiança 1.0 \textperiodcentered{} domínio paper \textperiodcentered{} história 5 versões no jornal}

%CRISTAL {"arestas":[{"alvo":"plano_no_so","confianca":1.0,"epistemico":"fato","origem":"paper","rel":"prova"},{"alvo":"delta","confianca":0.95,"epistemico":"fato","origem":"manual","rel":"prova"}],"confianca":1.0,"contraexemplos":[],"descricao":"[Propagação limitada] Uma capacidade emitida com profundidade só pode aparecer (atenuada) ao fim de cadeias de re-delegação de comprimento. Em particular, =0 não re-delegável: ela permanece exclusivamente no destinatário direto.","epistemico":"fato","exemplos":[],"id":"theorem_propagacao_limitadathmbounded_uma_capacidade_em","memoria":"semantica","meta":{"dominio":"paper","nivel":5,"tipo_teorema":"theorem"},"origem":"paper","palavras_chave":[],"sinonimos":[],"tipo":"conceito","titulo":"Propagação limitada"}
\section{Propagação limitada}
[Propagação limitada] Uma capacidade emitida com profundidade só pode aparecer (atenuada) ao fim de cadeias de re-delegação de comprimento. Em particular, =0 não re-delegável: ela permanece exclusivamente no destinatário direto.
\prov{relações: prova plano\_no\_so; prova delta}
\prov{proveniência: paper \textperiodcentered{} fato \textperiodcentered{} confiança 1.0 \textperiodcentered{} domínio paper \textperiodcentered{} história 63 versões no jornal}

%CRISTAL {"arestas":[{"alvo":"plano_no_so","confianca":1.0,"epistemico":"fato","origem":"paper","rel":"prova"},{"alvo":"bai","confianca":-0.078,"epistemico":"fato","origem":"manual","rel":"referencia"},{"alvo":"bai_biblioteca_de_autorizacao_isomorfica","confianca":0.95,"epistemico":"fato","origem":"manual","rel":"prova"}],"confianca":1.0,"contraexemplos":[],"descricao":"[Revogação em cascata] Revogar uma instância I revoga toda a sua descendência na floresta de derivação (as instâncias cuja cadeia de par passa por I). Após a cascata, nenhuma instância vigente foi atenuada de I.","epistemico":"fato","exemplos":[],"id":"theorem_revogacao_em_cascatathmcascade_revogar_uma_inst","memoria":"semantica","meta":{"dominio":"paper","nivel":5,"tipo_teorema":"theorem"},"origem":"paper","palavras_chave":[],"sinonimos":[],"tipo":"conceito","titulo":"Revogação em cascata"}
\section{Revogação em cascata}
[Revogação em cascata] Revogar uma instância I revoga toda a sua descendência na floresta de derivação (as instâncias cuja cadeia de par passa por I). Após a cascata, nenhuma instância vigente foi atenuada de I.
\prov{relações: prova plano\_no\_so; referencia bai; prova bai\_biblioteca\_de\_autorizacao\_isomorfica}
\prov{proveniência: paper \textperiodcentered{} fato \textperiodcentered{} confiança 1.0 \textperiodcentered{} domínio paper \textperiodcentered{} história 66 versões no jornal}

%CRISTAL {"arestas":[{"alvo":"plano_no_so","confianca":1.0,"epistemico":"fato","origem":"paper","rel":"prova"},{"alvo":"kappa","confianca":-0.047,"epistemico":"fato","origem":"manual","rel":"referencia"},{"alvo":"delta","confianca":-0.031,"epistemico":"fato","origem":"manual","rel":"referencia"},{"alvo":"bai","confianca":0.055,"epistemico":"fato","origem":"manual","rel":"referencia"},{"alvo":"colapso","confianca":-0.062,"epistemico":"fato","origem":"manual","rel":"referencia"},{"alvo":"bai_biblioteca_de_autorizacao_isomorfica","confianca":0.95,"epistemico":"fato","origem":"manual","rel":"prova"}],"confianca":1.0,"contraexemplos":[],"descricao":"[Tese-mãe da ] Toda instância da Biblioteca de Autorização Isomórfica é determinada pelo quinteto [ (T, ), ] onde é capacidade abstrata (Definição~), orçamento de propagação (Definição~), medida de organização estrutural (Definição~), T operador de atenuação deflacionário e operador de colapso (Definições~ e~). CASL, e são interpretações distintas destas primitivas --- não motivações do for","epistemico":"fato","exemplos":[],"id":"theorem_tese-mae_da_thmmae_toda_instancia_da_biblioteca","memoria":"semantica","meta":{"dominio":"paper","nivel":5,"tipo_teorema":"theorem"},"origem":"paper","palavras_chave":[],"sinonimos":[],"tipo":"conceito","titulo":"Tese-mãe da"}
\section{Tese-mãe da}
[Tese-mãe da ] Toda instância da Biblioteca de Autorização Isomórfica é determinada pelo quinteto [ (T, ), ] onde é capacidade abstrata (Definição\~{}), orçamento de propagação (Definição\~{}), medida de organização estrutural (Definição\~{}), T operador de atenuação deflacionário e operador de colapso (Definições\~{} e\~{}). CASL, e são interpretações distintas destas primitivas --- não motivações do for
\prov{relações: prova plano\_no\_so; referencia kappa; referencia delta; referencia bai; referencia colapso; prova bai\_biblioteca\_de\_autorizacao\_isomorfica}
\prov{proveniência: paper \textperiodcentered{} fato \textperiodcentered{} confiança 1.0 \textperiodcentered{} domínio paper \textperiodcentered{} história 126 versões no jornal}

%CRISTAL {"arestas":[{"alvo":"plano_no_so","confianca":1.0,"epistemico":"fato","origem":"paper","rel":"prova"},{"alvo":"reticulado","confianca":-0.008,"epistemico":"fato","origem":"manual","rel":"referencia"},{"alvo":"colapso","confianca":0.062,"epistemico":"fato","origem":"manual","rel":"referencia"},{"alvo":"reticulados","confianca":0.95,"epistemico":"fato","origem":"manual","rel":"referencia"}],"confianca":1.0,"contraexemplos":[],"descricao":"[Universalidade da ] Qualquer sistema cuja autoridade satisfaça: [nosep,=(*)] estados e condições formam um reticulado booleano (2X, ); existe operador de atenuação T sobre extents; existe operador de propagação P com orçamento decrementado a cada delegação válida; existe colapso =(_,) sobre candidatos vigentes, com fail-closed em; admite uma instanciação na --- isto é, capacidades, quinteto (T,) e","epistemico":"fato","exemplos":[],"id":"theorem_universalidade_da_thmuniversalidade_qualquer_si","memoria":"semantica","meta":{"dominio":"paper","nivel":5,"tipo_teorema":"theorem"},"origem":"paper","palavras_chave":[],"sinonimos":[],"tipo":"conceito","titulo":"Universalidade da"}
\section{Universalidade da}
[Universalidade da ] Qualquer sistema cuja autoridade satisfaça: [nosep,=(*)] estados e condições formam um reticulado booleano (2X, ); existe operador de atenuação T sobre extents; existe operador de propagação P com orçamento decrementado a cada delegação válida; existe colapso =(\_,) sobre candidatos vigentes, com fail-closed em; admite uma instanciação na --- isto é, capacidades, quinteto (T,) e
\prov{relações: prova plano\_no\_so; referencia reticulado; referencia colapso; referencia reticulados}
\prov{proveniência: paper \textperiodcentered{} fato \textperiodcentered{} confiança 1.0 \textperiodcentered{} domínio paper \textperiodcentered{} história 120 versões no jornal}

%CRISTAL {"arestas":[{"alvo":"broca_os","confianca":0.95,"epistemico":"fato","origem":"paper","rel":"parte_de"}],"confianca":1.0,"contraexemplos":[],"descricao":"itemize2pt\n[I.] Alimentar percebe e reconhece com o tecido da casa: papers\n(papers/*.tex), código (code/*.c), assets. Cada leitura acrescenta nós à rede sem retocar\n$A$.\n[II.] Compilar o design: cd papers \\&\\& make --- os PDFs são a face legível do tecido; o gato\nlê os .tex ou .pdf indiferentemente (bytes no metal).\n[III.] Verificar que o reconhecimento escala: ./reconhece ../papers/*.tex code/*.c deve manter\ntaxa alta após corromper $25\\%$.\nitemize","epistemico":"fato","exemplos":[],"id":"tirar_a_mao_---_o_que_nao","memoria":"semantica","meta":{"dominio":"paper","nivel":5,"verificacao":"overlap moderado — revisar manualmente"},"origem":"paper","palavras_chave":[],"sinonimos":[],"tipo":"conceito","titulo":"Tirar a mão --- o que não"}
\section{Tirar a mão --- o que não}
itemize2pt
[I.] Alimentar percebe e reconhece com o tecido da casa: papers
(papers/*.tex), código (code/*.c), assets. Cada leitura acrescenta nós à rede sem retocar
\$A\$.
[II.] Compilar o design: cd papers \textbackslash{}\&\textbackslash{}\& make --- os PDFs são a face legível do tecido; o gato
lê os .tex ou .pdf indiferentemente (bytes no metal).
[III.] Verificar que o reconhecimento escala: ./reconhece ../papers/*.tex code/*.c deve manter
taxa alta após corromper \$25\textbackslash{}\%\$.
itemize
\prov{relações: parte\_de broca\_os}
\prov{proveniência: paper \textperiodcentered{} fato \textperiodcentered{} confiança 1.0 \textperiodcentered{} domínio paper \textperiodcentered{} história 4 versões no jornal}

%CRISTAL {"arestas":[{"alvo":"topologia","confianca":1.0,"epistemico":"fato","origem":"paper","rel":"parte_de"},{"alvo":"grupos","confianca":1.0,"epistemico":"fato","origem":"paper","rel":"parte_de"},{"alvo":"geometria","confianca":1.0,"epistemico":"fato","origem":"paper","rel":"parte_de"}],"confianca":1.0,"contraexemplos":[],"descricao":"abstract A geometria riemanniana não é acrescentada à dinâmica de ouro depois, nem importada de fora: ela é derivada. Da topologia da torre dimensional ( R C H O ) e da álgebra da conservação --- a balança a+c²=1 e a norma quadrática Q de Hurwitz --- nasce a métrica: Q é o invariante que sobrevive, e a sua realização passa pela S³ dos versores, a métrica de de Sitter e o flip de Wick espacial até o palco hiperbólico H³/ de Bianchi. Sobre esse palco, os zeros de não vivem no e","epistemico":"fato","exemplos":[],"id":"topologia_algebra_geometria","memoria":"semantica","meta":{"dominio":"paper","nivel":5,"verificacao":"slug idêntico — enriquecer, não duplicar"},"origem":"paper","palavras_chave":[],"sinonimos":[],"tipo":"conceito","titulo":"Topologia + Álgebra Geometria"}
\section{Topologia + Álgebra Geometria}
abstract A geometria riemanniana não é acrescentada à dinâmica de ouro depois, nem importada de fora: ela é derivada. Da topologia da torre dimensional ( R C H O ) e da álgebra da conservação --- a balança a+c\ensuremath{^{2}}=1 e a norma quadrática Q de Hurwitz --- nasce a métrica: Q é o invariante que sobrevive, e a sua realização passa pela S\ensuremath{^{3}} dos versores, a métrica de de Sitter e o flip de Wick espacial até o palco hiperbólico H\ensuremath{^{3}}/ de Bianchi. Sobre esse palco, os zeros de não vivem no e
\prov{relações: parte\_de topologia; parte\_de grupos; parte\_de geometria}
\prov{proveniência: paper \textperiodcentered{} fato \textperiodcentered{} confiança 1.0 \textperiodcentered{} domínio paper \textperiodcentered{} história 156 versões no jornal}

%CRISTAL {"arestas":[{"alvo":"vida-morte","confianca":0.5,"epistemico":"fato","origem":"wiki","rel":"relaciona"},{"alvo":"virtualizacao","confianca":0.5,"epistemico":"fato","origem":"wiki","rel":"relaciona"},{"alvo":"utilitarismo","confianca":0.5,"epistemico":"fato","origem":"wiki","rel":"relaciona"},{"alvo":"veredito_experimental","confianca":0.5,"epistemico":"fato","origem":"wiki","rel":"relaciona"},{"alvo":"transicao","confianca":0.5,"epistemico":"fato","origem":"wiki","rel":"relaciona"},{"alvo":"word","confianca":0.5,"epistemico":"fato","origem":"wiki","rel":"relaciona"},{"alvo":"vontade_de_poder","confianca":0.5,"epistemico":"fato","origem":"wiki","rel":"relaciona"},{"alvo":"while","confianca":0.5,"epistemico":"fato","origem":"wiki","rel":"relaciona"}],"confianca":1.0,"contraexemplos":[],"descricao":"Paper: transcricao.tex","epistemico":"fato","exemplos":[],"id":"transcricao","memoria":"semantica","meta":{"dominio":"paper","nivel":5,"tipo_no":"paper"},"origem":"paper","palavras_chave":[],"sinonimos":[],"tipo":"conceito","titulo":"transcricao"}
\section{transcricao}
Paper: transcricao.tex
\prov{relações: relaciona vida-morte; relaciona virtualizacao; relaciona utilitarismo; relaciona veredito\_experimental; relaciona transicao; relaciona word; relaciona vontade\_de\_poder; relaciona while}
\prov{proveniência: paper \textperiodcentered{} fato \textperiodcentered{} confiança 1.0 \textperiodcentered{} domínio paper \textperiodcentered{} história 9 versões no jornal}

%CRISTAL {"arestas":[{"alvo":"main","confianca":1.0,"epistemico":"fato","origem":"paper","rel":"parte_de"}],"confianca":1.0,"contraexemplos":[],"descricao":"O flip logarítmico, ( x)(u)=x(eu), transforma dilatações multiplicativas em translações; e G[x]()=12₀^ x(t)t⁻i,dt/t=( F )[x] é Mellin (Fourier no logaritmo) sobre L²(_+,dt/t). N G é unitária (Parseval; round-trip até arredondamento de dupla precisão). G diagonaliza o shift multiplicativo DN Seja U_ o operador unitário de dilatação (U_ x)(t)=x( t), =¹/, na grade tₙ=ⁿ. Os caracteres de Mellin ₖ=t⁻iₖ, ₖ=2 k/(N), coincidem com a base","epistemico":"fato","exemplos":[],"id":"tres_impossibilidades","memoria":"semantica","meta":{"dominio":"paper","nivel":5,"verificacao":"slug idêntico — enriquecer, não duplicar"},"origem":"paper","palavras_chave":[],"sinonimos":[],"tipo":"conceito","titulo":"Três impossibilidades"}
\section{Três impossibilidades}
O flip logarítmico, ( x)(u)=x(eu), transforma dilatações multiplicativas em translações; e G[x]()=12\ensuremath{_{0}}\^{} x(t)t\ensuremath{^{-}}i,dt/t=( F )[x] é Mellin (Fourier no logaritmo) sobre L\ensuremath{^{2}}(\_+,dt/t). N G é unitária (Parseval; round-trip até arredondamento de dupla precisão). G diagonaliza o shift multiplicativo DN Seja U\_ o operador unitário de dilatação (U\_ x)(t)=x( t), =\ensuremath{^{1}}/, na grade t\ensuremath{_{n}}=\ensuremath{^{n}}. Os caracteres de Mellin \ensuremath{_{k}}=t\ensuremath{^{-}}i\ensuremath{_{k}}, \ensuremath{_{k}}=2 k/(N), coincidem com a base
\prov{relações: parte\_de main}
\prov{proveniência: paper \textperiodcentered{} fato \textperiodcentered{} confiança 1.0 \textperiodcentered{} domínio paper \textperiodcentered{} história 5 versões no jornal}

%CRISTAL {"arestas":[{"alvo":"lobo_frontal","confianca":1.0,"epistemico":"fato","origem":"paper","rel":"parte_de"}],"confianca":1.0,"contraexemplos":[],"descricao":"V ~§VI estabelece o isomorfismo meta-indução,=,memória de Hopfield em cinco passos. Cada um já roda em artefato real: 1.4 tabular@p0.5cmp6.2cmp5.4cm@ & passo (§VI) & já roda (1) & A=A T --- sinapse Hebbiana recíproca & estelar.py:gapmatriz (a sinapse fixa) (2) & E(x)=-12 x T!Ax Lyapunov --- descida ao atrator & testemunhado pela simetria de A + o decaimento monótono de vidamorte.py:ciclo (3) & A= P_+ P_, projetores ortogonais & estelar.","epistemico":"fato","exemplos":[],"id":"tudo_e_raciocinio_e_raciocinio_e_o_recall_ordenado","memoria":"semantica","meta":{"dominio":"paper","nivel":5,"verificacao":"conceito novo"},"origem":"paper","palavras_chave":[],"sinonimos":[],"tipo":"conceito","titulo":"Tudo é raciocínio, e raciocínio é o recall ordenado"}
\section{Tudo é raciocínio, e raciocínio é o recall ordenado}
V \~{}§VI estabelece o isomorfismo meta-indução,=,memória de Hopfield em cinco passos. Cada um já roda em artefato real: 1.4 tabular@p0.5cmp6.2cmp5.4cm@ \& passo (§VI) \& já roda (1) \& A=A T --- sinapse Hebbiana recíproca \& estelar.py:gapmatriz (a sinapse fixa) (2) \& E(x)=-12 x T!Ax Lyapunov --- descida ao atrator \& testemunhado pela simetria de A + o decaimento monótono de vidamorte.py:ciclo (3) \& A= P\_+ P\_, projetores ortogonais \& estelar.
\prov{relações: parte\_de lobo\_frontal}
\prov{proveniência: paper \textperiodcentered{} fato \textperiodcentered{} confiança 1.0 \textperiodcentered{} domínio paper \textperiodcentered{} história 3 versões no jornal}

%CRISTAL {"arestas":[{"alvo":"computacao_fundamentos_broca","confianca":1.0,"epistemico":"fato","origem":"paper","rel":"parte_de"}],"confianca":1.0,"contraexemplos":[],"descricao":"Arquitetura de Von Neumann Um computador clássico tem três blocos acoplados: [nosep,leftmargin=1.4em] CPU --- fetch--decode--execute; unidade aritmético-lógica; Memória --- palavras endereçáveis; programa e dados no mesmo espaço; Barramento --- transporte de palavras entre CPU e memória. A tese central (1945): o programa é dado --- reprogramar não exige re-cabo físico. definicao Isomorfismo operacional mínimo Seja M a matriz de mensage","epistemico":"fato","exemplos":[],"id":"turing_modelo_e_limites","memoria":"semantica","meta":{"dominio":"paper","nivel":5,"verificacao":"slug idêntico — enriquecer, não duplicar"},"origem":"paper","palavras_chave":[],"sinonimos":[],"tipo":"conceito","titulo":"Turing: modelo e limites"}
\section{Turing: modelo e limites}
Arquitetura de Von Neumann Um computador clássico tem três blocos acoplados: [nosep,leftmargin=1.4em] CPU --- fetch--decode--execute; unidade aritmético-lógica; Memória --- palavras endereçáveis; programa e dados no mesmo espaço; Barramento --- transporte de palavras entre CPU e memória. A tese central (1945): o programa é dado --- reprogramar não exige re-cabo físico. definicao Isomorfismo operacional mínimo Seja M a matriz de mensage
\prov{relações: parte\_de computacao\_fundamentos\_broca}
\prov{proveniência: paper \textperiodcentered{} fato \textperiodcentered{} confiança 1.0 \textperiodcentered{} domínio paper \textperiodcentered{} história 6 versões no jornal}

%CRISTAL {"arestas":[{"alvo":"algebra_dos_contratos","confianca":1.0,"epistemico":"fato","origem":"paper","rel":"parte_de"}],"confianca":1.0,"contraexemplos":[],"descricao":"1.25 tabularl|l|l & Bitcoin & Ethereum validação (univ.) & ECDSA/Schnorr secp256k1 & ECDSA secp256k1 coordenação (univ.) & hashlock SHA-256 no script & hashlock SHA-256 no contrato comissão (univ.) & divisão no redeem & divisão no redeem execução (espec.) & Script / P2WSH & EVM / Solidity liquidação (espec.) & UTXO & conta/saldo consenso (espec.) & PoW & PoS tabular As três de cima são idênticas em forma: o HTLC com hashlock SHA-256 comum fechou nas duas chains","epistemico":"fato","exemplos":[],"id":"um_contrato_de_ponta_a_ponta","memoria":"semantica","meta":{"dominio":"paper","nivel":5,"verificacao":"overlap moderado — revisar manualmente"},"origem":"paper","palavras_chave":[],"sinonimos":[],"tipo":"conceito","titulo":"Um contrato, de ponta a ponta"}
\section{Um contrato, de ponta a ponta}
1.25 tabularl|l|l \& Bitcoin \& Ethereum validação (univ.) \& ECDSA/Schnorr secp256k1 \& ECDSA secp256k1 coordenação (univ.) \& hashlock SHA-256 no script \& hashlock SHA-256 no contrato comissão (univ.) \& divisão no redeem \& divisão no redeem execução (espec.) \& Script / P2WSH \& EVM / Solidity liquidação (espec.) \& UTXO \& conta/saldo consenso (espec.) \& PoW \& PoS tabular As três de cima são idênticas em forma: o HTLC com hashlock SHA-256 comum fechou nas duas chains
\prov{relações: parte\_de algebra\_dos\_contratos}
\prov{proveniência: paper \textperiodcentered{} fato \textperiodcentered{} confiança 1.0 \textperiodcentered{} domínio paper \textperiodcentered{} história 3 versões no jornal}

%CRISTAL {"arestas":[{"alvo":"olho_de_koch","confianca":1.0,"epistemico":"fato","origem":"paper","rel":"parte_de"}],"confianca":1.0,"contraexemplos":[],"descricao":"empty !85O olho não copia o gerador --- ressoa com ele. A projeção clássica é o casamento de fase no canal negro; o atrator atravessa, a matriz A fica em casa. abstract O é a fronteira associativa do: onde o núcleo vivo [,total,e,] encontra o mundo clássico. Não há terceiro caminho escalar nem soma forçada de coordenadas --- há a máquina vazia de Peirce, com duas faces da mesma quadração z z²: a branca (tropical, idempotente, estado que conta) e a negra (glacial","epistemico":"fato","exemplos":[],"id":"uma_maquina_duas_faces","memoria":"semantica","meta":{"dominio":"paper","nivel":5},"origem":"paper","palavras_chave":[],"sinonimos":[],"tipo":"conceito","titulo":"Uma máquina, duas faces"}
\section{Uma máquina, duas faces}
empty !85O olho não copia o gerador --- ressoa com ele. A projeção clássica é o casamento de fase no canal negro; o atrator atravessa, a matriz A fica em casa. abstract O é a fronteira associativa do: onde o núcleo vivo [,total,e,] encontra o mundo clássico. Não há terceiro caminho escalar nem soma forçada de coordenadas --- há a máquina vazia de Peirce, com duas faces da mesma quadração z z\ensuremath{^{2}}: a branca (tropical, idempotente, estado que conta) e a negra (glacial
\prov{relações: parte\_de olho\_de\_koch}
\prov{proveniência: paper \textperiodcentered{} fato \textperiodcentered{} confiança 1.0 \textperiodcentered{} domínio paper \textperiodcentered{} história 3 versões no jornal}

%CRISTAL {"arestas":[{"alvo":"corpo_glacial","confianca":1.0,"epistemico":"fato","origem":"paper","rel":"parte_de"},{"alvo":"corpo_tropical","confianca":1.0,"epistemico":"fato","origem":"paper","rel":"confirma"}],"confianca":1.0,"contraexemplos":[],"descricao":"Irmão fundamento de O Corpo Tropical. A tese é que a parte branca e a parte negra não são duas máquinas: são duas faces de uma única máquina vazia, acopladas, e a dinâmica é o acoplamento. Concretamente, na decomposição de Peirce de uma álgebra associativa unital com idempotentes ortogonais 1=ᵢ eᵢ: a diagonal eᵢ eᵢ são as células (parte branca --- álgebras de canto, com identidade local eᵢ: o estado, parte de 1); a off-diagonal eᵢ ej (i j) são os canais (parte negra, cada elem","epistemico":"fato","exemplos":[],"id":"uma_maquina_nao_duas","memoria":"semantica","meta":{"dominio":"paper","nivel":5,"verificacao":"slug idêntico — enriquecer, não duplicar"},"origem":"paper","palavras_chave":[],"sinonimos":[],"tipo":"conceito","titulo":"Uma máquina, não duas"}
\section{Uma máquina, não duas}
Irmão fundamento de O Corpo Tropical. A tese é que a parte branca e a parte negra não são duas máquinas: são duas faces de uma única máquina vazia, acopladas, e a dinâmica é o acoplamento. Concretamente, na decomposição de Peirce de uma álgebra associativa unital com idempotentes ortogonais 1=\ensuremath{_{i}} e\ensuremath{_{i}}: a diagonal e\ensuremath{_{i}} e\ensuremath{_{i}} são as células (parte branca --- álgebras de canto, com identidade local e\ensuremath{_{i}}: o estado, parte de 1); a off-diagonal e\ensuremath{_{i}} ej (i j) são os canais (parte negra, cada elem
\prov{relações: parte\_de corpo\_glacial; confirma corpo\_tropical}
\prov{proveniência: paper \textperiodcentered{} fato \textperiodcentered{} confiança 1.0 \textperiodcentered{} domínio paper \textperiodcentered{} história 41 versões no jornal}

%CRISTAL {"arestas":[{"alvo":"fase3","confianca":1.0,"epistemico":"fato","origem":"paper","rel":"parte_de"}],"confianca":1.0,"contraexemplos":[],"descricao":"passobox\nenumerate2pt\n Instalar libtiffany.so em /usr/lib (ou \\$HOME/.local/lib).\n Extrair Core: tar -xaf tiffany-core-1.0.tar.zst -C \\$TIFFANY\\_HOME.\n Extrair template cliente: cliente/ vazio pronto.\n Testar: tiffany ``como clonar repo git?'' $$ resposta da Base Geral.\n Ingerir docs próprios: tiffany-ingest --dominio empresa/engenharia ./docs/.\n Reindexar: tiffany-index cliente/empresa/engenharia/nos.tsv v1.idx.\nenumerate\npassobox\n\n Variável TIFFANY\\_HOME (default \\$HOME/.tiffany) unifica paths","epistemico":"fato","exemplos":[],"id":"validacao_benchmarks_fase_3","memoria":"semantica","meta":{"dominio":"paper","nivel":5,"verificacao":"overlap moderado — revisar manualmente"},"origem":"paper","palavras_chave":[],"sinonimos":[],"tipo":"conceito","titulo":"Validação (benchmarks Fase 3)"}
\section{Validação (benchmarks Fase 3)}
passobox
enumerate2pt
 Instalar libtiffany.so em /usr/lib (ou \textbackslash{}\$HOME/.local/lib).
 Extrair Core: tar -xaf tiffany-core-1.0.tar.zst -C \textbackslash{}\$TIFFANY\textbackslash{}\_HOME.
 Extrair template cliente: cliente/ vazio pronto.
 Testar: tiffany ``como clonar repo git?'' \$\$ resposta da Base Geral.
 Ingerir docs próprios: tiffany-ingest --dominio empresa/engenharia ./docs/.
 Reindexar: tiffany-index cliente/empresa/engenharia/nos.tsv v1.idx.
enumerate
passobox

 Variável TIFFANY\textbackslash{}\_HOME (default \textbackslash{}\$HOME/.tiffany) unifica paths
\prov{relações: parte\_de fase3}
\prov{proveniência: paper \textperiodcentered{} fato \textperiodcentered{} confiança 1.0 \textperiodcentered{} domínio paper \textperiodcentered{} história 2 versões no jornal}

%CRISTAL {"arestas":[{"alvo":"cifra_ouro_branco","confianca":1.0,"epistemico":"fato","origem":"paper","rel":"parte_de"}],"confianca":1.0,"contraexemplos":[],"descricao":"N A prata minera-se de dentro pela estrela de prata: os 6 modos dominantes (a estrela de Salomão, hexagrama). A emissão exige exatamente 6 modos: rank-6 minera; rank-5 (pentagrama --- o ouro, é o pentágono) não (é lastro); rank-7 e ruído também não. Verificado. Oferta com teto (1310⁶): prova-se trabalho ou não se emite.","epistemico":"fato","exemplos":[],"id":"validacao_coletiva_por_n-qubits","memoria":"semantica","meta":{"dominio":"paper","nivel":5,"verificacao":"conceito novo"},"origem":"paper","palavras_chave":[],"sinonimos":[],"tipo":"conceito","titulo":"Validação coletiva por N-qubits"}
\section{Validação coletiva por N-qubits}
N A prata minera-se de dentro pela estrela de prata: os 6 modos dominantes (a estrela de Salomão, hexagrama). A emissão exige exatamente 6 modos: rank-6 minera; rank-5 (pentagrama --- o ouro, é o pentágono) não (é lastro); rank-7 e ruído também não. Verificado. Oferta com teto (1310\ensuremath{^{6}}): prova-se trabalho ou não se emite.
\prov{relações: parte\_de cifra\_ouro\_branco}
\prov{proveniência: paper \textperiodcentered{} fato \textperiodcentered{} confiança 1.0 \textperiodcentered{} domínio paper \textperiodcentered{} história 4 versões no jornal}

%CRISTAL {"arestas":[{"alvo":"pow_gato","confianca":1.0,"epistemico":"fato","origem":"paper","rel":"parte_de"}],"confianca":1.0,"contraexemplos":[],"descricao":"defi[Verificador PoW (nó / ASIC)] Verify(header,n,T):= [,uint256(²(header n)) < uint256(T),]. defi Equivalência funcional gato = ASIC Seja AASIC qualquer minerador SHA-256 e Gg o gato g da farm ( III). Para todo (header,n,T): [ Gg emite (n) AASIC emite (n) Verify(header,n,T)=true. ] teo proof Ambos avaliam a mesma função ²(header n) em ordem little-endian de n (V struct.pack(\"<I\")). O gato não altera, não reduz T, não salta no","epistemico":"fato","exemplos":[],"id":"validacao_de_negocios_braco_orthogonal","memoria":"semantica","meta":{"dominio":"paper","nivel":5,"verificacao":"conceito novo"},"origem":"paper","palavras_chave":[],"sinonimos":[],"tipo":"conceito","titulo":"Validação de negócios (braço orthogonal)"}
\section{Validação de negócios (braço orthogonal)}
defi[Verificador PoW (nó / ASIC)] Verify(header,n,T):= [,uint256(\ensuremath{^{2}}(header n)) < uint256(T),]. defi Equivalência funcional gato = ASIC Seja AASIC qualquer minerador SHA-256 e Gg o gato g da farm ( III). Para todo (header,n,T): [ Gg emite (n) AASIC emite (n) Verify(header,n,T)=true. ] teo proof Ambos avaliam a mesma função \ensuremath{^{2}}(header n) em ordem little-endian de n (V struct.pack("<I")). O gato não altera, não reduz T, não salta no
\prov{relações: parte\_de pow\_gato}
\prov{proveniência: paper \textperiodcentered{} fato \textperiodcentered{} confiança 1.0 \textperiodcentered{} domínio paper \textperiodcentered{} história 3 versões no jornal}

%CRISTAL {"arestas":[{"alvo":"transformada_dourada","confianca":1.0,"epistemico":"fato","origem":"paper","rel":"parte_de"}],"confianca":1.0,"contraexemplos":[],"descricao":"A versão de 1D usa a estrutura multiplicativa de (>₀,). Ela se eleva às quatro álgebras de divisão normadas --- e só a elas. definition[Transformada dourada de Hurwitz] Seja A, H, O, n= A1,2,4,8, com norma multiplicativa |xy|=|x|,|y|. Todo x A0 escreve-se x=, u com =|x|>₀ e u Sⁿ⁻¹. Para f:A0 define-se [ A[f](m)=Sⁿ⁻¹!!₀^ f(u),⁻i, Y m( u);d,d( u), ] onde Y m é a base de harmônicos esférico","epistemico":"fato","exemplos":[],"id":"validacao_numerica","memoria":"semantica","meta":{"dominio":"paper","nivel":5,"verificacao":"conceito novo"},"origem":"paper","palavras_chave":[],"sinonimos":[],"tipo":"conceito","titulo":"Validação numérica"}
\section{Validação numérica}
A versão de 1D usa a estrutura multiplicativa de (>\ensuremath{_{0}},). Ela se eleva às quatro álgebras de divisão normadas --- e só a elas. definition[Transformada dourada de Hurwitz] Seja A, H, O, n= A1,2,4,8, com norma multiplicativa |xy|=|x|,|y|. Todo x A0 escreve-se x=, u com =|x|>\ensuremath{_{0}} e u S\ensuremath{^{n-1}}. Para f:A0 define-se [ A[f](m)=S\ensuremath{^{n-1}}!!\ensuremath{_{0}}\^{} f(u),\ensuremath{^{-}}i, Y m( u);d,d( u), ] onde Y m é a base de harmônicos esférico
\prov{relações: parte\_de transformada\_dourada}
\prov{proveniência: paper \textperiodcentered{} fato \textperiodcentered{} confiança 1.0 \textperiodcentered{} domínio paper \textperiodcentered{} história 3 versões no jornal}

%CRISTAL {"arestas":[{"alvo":"assistente","confianca":1.0,"epistemico":"fato","origem":"paper","rel":"parte_de"}],"confianca":1.0,"contraexemplos":[],"descricao":"I Ordem proposta:\n\nenumerate4pt\n[1.] API estável  --- V em curso (tiffany\\_open/query/respond/close).\n[2.] Formato + versionamento --- .tsv, vN.idx, meta.json; memória separada.\n[3.] Ingestão --- PDF, código, conversas $$ tecido $$ tiffany-index.\n[4.] Orquestrador + avaliador --- roteamento (obs empírico) + gate de qualidade.\n[5.] Métricas (§VIII) --- sucesso, tempo, intervenção humana, precisão de roteamento;\nfase2.sh como painel da Fase~2.\n[6.] Interfaces --- CLI, Python, web, voz; todas conso","epistemico":"fato","exemplos":[],"id":"validacao_tecnica","memoria":"semantica","meta":{"dominio":"paper","nivel":5,"verificacao":"conceito novo"},"origem":"paper","palavras_chave":[],"sinonimos":[],"tipo":"conceito","titulo":"Validação técnica"}
\section{Validação técnica}
I Ordem proposta:

enumerate4pt
[1.] API estável  --- V em curso (tiffany\textbackslash{}\_open/query/respond/close).
[2.] Formato + versionamento --- .tsv, vN.idx, meta.json; memória separada.
[3.] Ingestão --- PDF, código, conversas \$\$ tecido \$\$ tiffany-index.
[4.] Orquestrador + avaliador --- roteamento (obs empírico) + gate de qualidade.
[5.] Métricas (§VIII) --- sucesso, tempo, intervenção humana, precisão de roteamento;
fase2.sh como painel da Fase\~{}2.
[6.] Interfaces --- CLI, Python, web, voz; todas conso
\prov{relações: parte\_de assistente}
\prov{proveniência: paper \textperiodcentered{} fato \textperiodcentered{} confiança 1.0 \textperiodcentered{} domínio paper \textperiodcentered{} história 2 versões no jornal}

%CRISTAL {"arestas":[{"alvo":"floco_de_neve","confianca":1.0,"epistemico":"fato","origem":"paper","rel":"parte_de"}],"confianca":1.0,"contraexemplos":[],"descricao":"PoW e certificação de partes (espectro, banda, bump) são ortogonais: PoW: loteria de hash / estrela --- recompensa de bloco ou emissão nativa. Validação: validarconjunto(partes) banda hashlock HTLC; comissão 1%. I Um não substitui o outro; ambos são legítimos.","epistemico":"fato","exemplos":[],"id":"verificacao_empirica","memoria":"semantica","meta":{"dominio":"paper","duplicata_sugerida":[],"nivel":5,"verificacao":"parte de conceito mais amplo 'verificacao_empirica'"},"origem":"paper","palavras_chave":[],"sinonimos":[],"tipo":"conceito","titulo":"Verificação empírica"}
\section{Verificação empírica}
PoW e certificação de partes (espectro, banda, bump) são ortogonais: PoW: loteria de hash / estrela --- recompensa de bloco ou emissão nativa. Validação: validarconjunto(partes) banda hashlock HTLC; comissão 1\%. I Um não substitui o outro; ambos são legítimos.
\prov{relações: parte\_de floco\_de\_neve}
\prov{proveniência: paper \textperiodcentered{} fato \textperiodcentered{} confiança 1.0 \textperiodcentered{} domínio paper \textperiodcentered{} história 4 versões no jornal}

%CRISTAL {"arestas":[{"alvo":"algebra_dos_contratos","confianca":1.0,"epistemico":"fato","origem":"paper","rel":"parte_de"}],"confianca":1.0,"contraexemplos":[],"descricao":"A divisão universal/específico que o leitor já aceitou é exatamente a que a Engenharia de Ouro modela pela balança a+c²=1. A condição se ergue sobre primitivas one-way --- a assinatura (kprᵢv! kpub) e o hash (s H(s)): é o lado a (irreversível, a segurança). A liquidação (mover valor, em princípio reversível) é o lado c. Daí o lema: a GoldChain é a chain do lado a, e pode existir sobre as chains do lado c sem repetir o consenso delas. prop[universalidade da camada do p","epistemico":"fato","exemplos":[],"id":"verificacao_por_metodo_os_n-qubits_como_generalizacao","memoria":"semantica","meta":{"dominio":"paper","nivel":5,"verificacao":"conceito novo"},"origem":"paper","palavras_chave":[],"sinonimos":[],"tipo":"conceito","titulo":"Verificação por método; os N-qubits como generalização"}
\section{Verificação por método; os N-qubits como generalização}
A divisão universal/específico que o leitor já aceitou é exatamente a que a Engenharia de Ouro modela pela balança a+c\ensuremath{^{2}}=1. A condição se ergue sobre primitivas one-way --- a assinatura (kpr\ensuremath{_{i}}v! kpub) e o hash (s H(s)): é o lado a (irreversível, a segurança). A liquidação (mover valor, em princípio reversível) é o lado c. Daí o lema: a GoldChain é a chain do lado a, e pode existir sobre as chains do lado c sem repetir o consenso delas. prop[universalidade da camada do p
\prov{relações: parte\_de algebra\_dos\_contratos}
\prov{proveniência: paper \textperiodcentered{} fato \textperiodcentered{} confiança 1.0 \textperiodcentered{} domínio paper \textperiodcentered{} história 4 versões no jornal}

%CRISTAL {"arestas":[{"alvo":"vida","confianca":1.0,"epistemico":"fato","origem":"paper","rel":"paralelo"},{"alvo":"morte","confianca":1.0,"epistemico":"fato","origem":"paper","rel":"paralelo"},{"alvo":"virtualizacao","confianca":0.5,"epistemico":"fato","origem":"wiki","rel":"relaciona"},{"alvo":"word","confianca":0.5,"epistemico":"fato","origem":"wiki","rel":"relaciona"},{"alvo":"vontade_de_poder","confianca":0.5,"epistemico":"fato","origem":"wiki","rel":"relaciona"},{"alvo":"while","confianca":0.5,"epistemico":"fato","origem":"wiki","rel":"relaciona"},{"alvo":"vita_contemplativa","confianca":0.5,"epistemico":"fato","origem":"wiki","rel":"relaciona"}],"confianca":1.0,"contraexemplos":[],"descricao":"Paper: vida-morte.tex","epistemico":"fato","exemplos":[],"id":"vida-morte","memoria":"semantica","meta":{"dominio":"paper","nivel":5,"tipo_no":"paper"},"origem":"paper","palavras_chave":[],"sinonimos":[],"tipo":"conceito","titulo":"vida-morte"}
\section{vida-morte}
Paper: vida-morte.tex
\prov{relações: paralelo vida; paralelo morte; relaciona virtualizacao; relaciona word; relaciona vontade\_de\_poder; relaciona while; relaciona vita\_contemplativa}
\prov{proveniência: paper \textperiodcentered{} fato \textperiodcentered{} confiança 1.0 \textperiodcentered{} domínio paper \textperiodcentered{} história 64 versões no jornal}

%CRISTAL {"arestas":[{"alvo":"computacao_fundamentos_broca","confianca":1.0,"epistemico":"fato","origem":"paper","rel":"parte_de"}],"confianca":1.0,"contraexemplos":[],"descricao":"Computação Computação é o estudo de procedimentos efectivos sobre símbolos discretos: o que pode ser calculado, com que recursos, e com que garantias. No currículo Tiffany situa-se depois da matemática (números, estruturas) e antes do sistema concreto. definicao ências: lógica, semântica (fase~3); inteiros, booleanos (fase~4). Generaliza: engenharia de software, sistemas operacionais, redes.","epistemico":"fato","exemplos":[],"id":"von_neumann_programa_armazenado","memoria":"semantica","meta":{"dominio":"paper","nivel":5,"verificacao":"slug idêntico — enriquecer, não duplicar"},"origem":"paper","palavras_chave":[],"sinonimos":[],"tipo":"conceito","titulo":"Von Neumann: programa armazenado"}
\section{Von Neumann: programa armazenado}
Computação Computação é o estudo de procedimentos efectivos sobre símbolos discretos: o que pode ser calculado, com que recursos, e com que garantias. No currículo Tiffany situa-se depois da matemática (números, estruturas) e antes do sistema concreto. definicao ências: lógica, semântica (fase\~{}3); inteiros, booleanos (fase\~{}4). Generaliza: engenharia de software, sistemas operacionais, redes.
\prov{relações: parte\_de computacao\_fundamentos\_broca}
\prov{proveniência: paper \textperiodcentered{} fato \textperiodcentered{} confiança 1.0 \textperiodcentered{} domínio paper \textperiodcentered{} história 6 versões no jornal}

\end{document}
